<<<<<<< HEAD
\setcounter{deso}{0}
\begin{name}
	{\tenchude}
	{TỔNG HỢP TRẢ LỜI NGẮN}
	{LỚP TOÁN THẦY PHÁT}
	{\lq\lq  It's not how much time you have, it's how you use it.\rq\rq}
\end{name}
\setcounter{ex}{0}
\begin{ex}%[0D0C2-3]%[KNTT - Lớp 12 - Dự án C - Đợt 2]%[Loc do]%[PB: Phan Anh]
	Bác Nghĩa đang giúp sắp xếp $16$ cuốn sách ôn thi vào một chiếc kệ có $5$ ngăn phân biệt, $16$ cuốn sách này thuộc $8$ môn học khác nhau: Toán, Lý, Hóa, Sinh, Sử, Địa, Văn, Anh. Mỗi môn học gồm đúng hai cuốn: một cuốn sách giáo khoa và một cuốn sách bài tập.  Để việc ôn tập đạt hiệu quả cao nhất theo từng khối thi, bác Nghĩa đặt ra các quy tắc khắc khe như sau: 
	\begin{itemize}
		\item Do ngăn kệ nhỏ, mỗi ngăn chỉ chứa được tối đa $5$ cuốn sách và không được để ngăn nào trống.
		\item Hai cuốn sách của cùng một môn học phải luôn nằm chung một ngăn với nhau.
		\item Các môn học trong cùng một tổ hợp môn thi phải nằm ở ba ngăn liên tiếp để thuận tiện cho việc tra cứu. Các tổ hợp bao gồm: (Văn, Sử, Địa); (Toán, Lý, Hóa); (Toán, Hóa, Sinh) và (Toán, Lý, Anh).
		\item Các cuốn sách trong mỗi ngăn được xếp theo hàng ngang với gáy sách quay ra ngoài ở mỗi ngăn, thứ tự từ trái qua phải.
	\end{itemize}
	Tổng số cách sắp xếp $16$ cuốn sách này vào $5$ kệ thỏa mãn điều kiện trên là $T$. Tính $\dfrac{T}{864} $.
	\par\shortans[0]{$6912$}
	\loigiai{
		Để xếp $16$ cuốn sách vào $5$ ngăn sao cho không có ngăn nào được để trống và các sách cùng môn sẽ được xếp chung ngăn thì sẽ xếp được $3$ ngăn chứa $4$ cuốn ($2$ môn) và $2$ ngăn chứa $2$ cuốn ($1$ môn).\\
		Số cách hoán vị sách trong các ngăn này (các sách cùng môn ở chung ngăn nhưng không nhất thiết phải ở cạnh nhau) là $(4!)^3 \cdot (2!)^2 $.\\
		Theo đề yêu cầu, ta đánh giá được như sau\\
		Ta thấy Toán tham gia vào $3$ tổ hợp liên tiếp nên cố định Toán ở ngăn $3$ để thỏa tính liền kề của ba tổ hợp môn.	Nhóm xã hội ( Văn, Sử, Địa) phải nằm ở ba ngăn liên tiếp nên hoán vị nội bộ của $3$ môn này là $3!$.\\
		Xét các trường hợp sau:\\
		\textbf{Trường hợp $1$}: Các môn tự nhiên xếp đều trên $5$ ngăn ( mỗi ngăn một môn tự nhiên). Khi đó có hai khả năng xảy ra là
		\begin{itemize}
			\item Sinh -- Hóa -- Toán -- Lý -- Anh.
			\item Anh -- Lý -- Toán -- Hóa -- Sinh.
		\end{itemize}
		Lúc này mỗi ngăn chỉ chứa $2$ cuốn của tự nhiên nên có thể đặt thêm ba môn Xã Hội vào bất kỳ $3$ ngăn liên tiếp trong $5$ ngăn đó là ($1,2,3$); ($2,3,4$) và ($3,4,5$).\\
		Vậy trường hợp này có: $2\cdot  3\cdot  3!=36$ (cách).\\
		\textbf{Trường hợp $2$}: Có $2$ môn tự nhiên nằm ở ngăn $2$. Khi đó có $4$ khả năng xảy ra là
		\begin{itemize}
			\item Anh ( ngăn $1$) -- Lý, Sinh ( ngăn $2$) -- Toán (ngăn $3$) -- Hóa (ngăn $4$);
			\item Hóa ( ngăn $1$) -- Lý, Sinh ( ngăn $2$) -- Toán (ngăn $3$) -- Anh (ngăn $4$);
			\item Sinh ( ngăn $1$) -- Hóa, Anh (ngăn $2$) -- Toán (ngăn $3$) -- Lý ( ngăn $4$)
			\item Lý (ngăn $1$) -- Hóa, Anh (ngăn $2$) -- Toán (ngăn $3$) -- Sinh (ngăn $4$)
		\end{itemize}
		Khi đó tổ hợp Xã Hội chỉ được xếp vào $3$ ngăn liên tiếp là ($3,4,5$).\\
		Vậy trường hợp này có: $4\cdot  1\cdot  3!=24$ (cách).\\
		\textbf{Trường hợp $3$}: Có một cặp tự nhiên nằm ở ngăn $4$. Trường hợp này là cách làm đối xứng của trường hợp $2$ nên có $24$ (cách).\\
		\textbf{Trường hợp $4$}: Có hai cặp môn tự nhiên cùng nằm ở ngăn $1$ và ngăn $2$. Khi đó có $2$ khả năng xảy ra là
		\begin{itemize}
			\item Lý, Sinh ( ngăn $1$) -- Hóa, Anh (ngăn $2$) -- Toán ( ngăn $3$);
			\item Hóa, Anh (ngăn $1$) -- Lý, Sinh ( ngăn $2$) -- Toán ( ngăn $3$)
		\end{itemize}
		Khi đó tổ hợp Xã Hội chỉ được xếp vào $3$ ngăn liên tiếp là ($3,4,5$).\\
		Vậy trường hợp này có $2\cdot  1\cdot  3!=12$ (cách).\\
		\textbf{Trường hợp $5$}: Có hai cặp môn tự nhiên cùng nằm ở ngăn $4$ và ngăn $5$. Đây là cách làm đối xứng với trường hợp $4$ nên có $12$ (cách).\\
		Suy ra tổng số xếp sách ở các trường hợp là $36+24+24+12+12=108$ ( cách).\\
		Vậy để xếp $18$ cuốn sách theo yêu cầu sẽ có $T=(4!)^3 \cdot  (2!)^2 \cdot  108=5\,971\,968$ (cách).\\
		Khi đó $\dfrac{T}{864} =6\,912$.}
\end{ex}

\begin{ex}%[0D0C2-5]%[Dự án C đề thi thử TN Môn Toán NH25-26 - Dot 4 - Phạm Đăng Đăng]%[THPT Lê Quý Đôn-Đống Đa-Hà Nội]%[GVPB - Phan Quốc Trí]
	Trong một trò chơi có thưởng bạn An được chơi một lần bằng cách lấy ngẫu nhiên ra $3$ quả cầu từ một hộp $150$ quả đánh số từ $1$ đến $150$. Bạn An sẽ nhận được phần thưởng nếu $3$ quả cầu bạn lấy ra thoả mãn đồng thời hai yêu cầu sau:\begin{itemize}
		\item Tổng các số được đánh số trên ba quả cầu  không vượt quả $300$.
		\item Các số trên $3$ quả cầu lấy được lập thành một cấp số cộng.
	\end{itemize}
	Xác suất để bạn An được phần thưởng là $p$, tính giá trị của $1\,000p$ (Kết quả là tròn đến hàng phần chục).
	\shortans{$7{,}8$}
	\loigiai{Chọn ngẫu nhiên $3$ quả cầu trong $150$ quả cầu nên $n(\Omega)=\mathrm{C}_{150}^3$.\\
		Giả sử 3 số chọn được là $a-d$; $a$; $a+d$ lập thành cấp số cộng với $a\ge 1$,  $d\ge 1$.\\
		Theo đề bài ta có
		$a-d+a+a+d\le 300 \Leftrightarrow 3a\le 300 \Leftrightarrow a\le 100
		$.\\
		Suy ra $\heva{&1\le a\le 100\\&a+d\le 150\\&a-d\ge 1}
		\Leftrightarrow\heva{&1\le a\le 100\\&d\le 150-a\\&d\le a-1.}$
		\begin{itemize}
			\item[\textbf{TH1:}] 
			$a-1\le 150-a \Leftrightarrow 2a\le 151 \Leftrightarrow a\le \dfrac{151}{2} \Leftrightarrow 1\le a\le 75.$\\
			Vậy có $75$ trường hợp xảy ra đối với $a$.\\
			Khi đó $d\le a-1 \Leftrightarrow 0< d\le 74$ nên có $74$ trường hợp xảy ra đối với $d$.\\
			Số trường hợp xảy ra là
			$\dfrac{75\cdot 74}{2}=2\,775$.
			\item[\textbf{TH2:}] $a-1\ge 150-a \Leftrightarrow 2a\ge 151 \Leftrightarrow a\ge \dfrac{151}{2} \Leftrightarrow 76\le a\le 100.$\\
			Khi đó $d\le 150-a \Leftrightarrow 150-100 \le d \le 150-76 \Leftrightarrow 50\le d\le 74 \ (\text{có }25\text{ số})$.
			Vậy số trường hợp xảy ra là $\dfrac{(50+74)\cdot 25}{2}=1\,550$.\\
			Tổng các trường hợp xảy ra là $2\,775+1\,550=4\,325$.
		\end{itemize}
		Vậy	xác suất cần tìm là $p=\dfrac{4\,325}{\mathrm{C}_{150}^3}\approx 0{,}0078469 \Rightarrow 1\,000\cdot p \approx 7{,}8.
		$}
\end{ex}

\begin{ex}%[Dự án đề TT Toán NH25-26-Dot 4- Hải Phụng]%[THCS-THPT LÊ THÁNH TÔNG - TP.HCM]%[GVPB - Viet Hoang Pham]%[0D0C2-6] 
	\immini[thm]{
		Trên một bàn cờ vua, Thông và Nghĩa cùng tham gia một trò chơi xác suất thú vị. Họ đặt một quân Mã tại ô D4 và một quân Vua tại ô F6. Quy ước rằng trong trò chơi này, hai quân cờ không ăn nhau. Trò chơi được tiến hành như sau:
		\begin{itemize}
			\item Thông: Thực hiện hành trình $4$ bước đi sao cho các ô đi qua trong $3$ bước đầu không trùng nhau và sau bước thứ $4$ Mã phải quay về lại đúng vị trí D4 ban đầu.
			\item Nghĩa: Thực hiện hành trình $3$ bước đi sao cho các ô đi qua trong $2$ bước đầu không trùng nhau và sau bước thứ $3$ Vua phải quay về lại đúng vị trí F6 ban đầu.
		\end{itemize}
	}
	{
		\begin{tikzpicture}[x=0.7cm, y=0.7cm]%Khai báo thêm gói chessfss để tạo quân cờ \usepackage{chessfss}
			%Thêm gói lệnh   \usepackage{chessfss}
			% 1. Định nghĩa màu sắc
			\definecolor{lightsq}{HTML}{F0D9B5} 
			\definecolor{darksq}{HTML}{B58863}  
			% 2. Vẽ bàn cờ
			\foreach \x in {0,...,7}
			\foreach \y in {0,...,7}
			{
				\pgfmathparse{mod(\x+\y,2) ? "darksq" : "lightsq"}
				\fill[\pgfmathresult] (\x,\y) rectangle (\x+1,\y+1);
			}
			% 3. Khung viền kép (Double border)
			\draw[thick] (0,0) rectangle (8,8);
			\draw[black, thin] (-0.8,-0.8) rectangle (8.8,8.8);
			\draw[black, thick] (-0.9,-0.9) rectangle (8.9,8.9);
			% 4. Tọa độ A-H
			\foreach \x [count=\i from 0] in {A,B,C,D,E,F,G,H} {
				\node at (\i+0.5, -0.4) {\small\textbf{\x}};
				\node at (\i+0.5, 8.4) {\small\textbf{\x}};
			}
			% 5. Tọa độ 1-8
			\foreach \y [count=\i from 0] in {1,2,3,4,5,6,7,8} {
				\node at (-0.4, \i+0.5) {\small\textbf{\y}};
				\node at (8.4, \i+0.5) {\small\textbf{\y}};
			}
			% 6. Đặt quân cờ bằng lệnh của gói skak
			% f6 tương ứng tọa độ (5.5, 5.5)
			% d4 tương ứng tọa độ (3.5, 3.5)
			% Quân Vua Đen (Black King)
			%\node at (5.5, 5.5) {\fontsize{35}{35}\selectfont \blackking}; 
			\node at (5.5, 5.5) {\fontsize{35}{35}\selectfont \BlackKingOnWhite};
			% Quân Mã Đen (Black Knight)
			\node at (3.5, 3.5) {\fontsize{35}{35}\selectfont \BlackKnightOnWhite};
			
		\end{tikzpicture}
	}
	\noindent Trình tự: Hai quân di chuyển xen kẽ nhau, bắt đầu bằng Mã.  Cụ thể: Mã đi bước $1$, Vua đi bước $1$, Mã đi bước $2$, Vua đi bước $2$,... cho đến khi cả hai hoàn tất hành trình. Xác suất để trong suốt quá trình di chuyển trên, có ít nhất một thời điểm Mã và Vua cùng đứng trên một ô cờ có dạng phân số tối giản  $\dfrac{a}{b}$ với $a, b \in \mathbb{N}^*$. Hãy tính giá trị của  $b-5a$.\\
	Cách di chuyển của quân Mã: Mã di chuyển theo đường chéo hình chữ nhật $2\times3$ ô vuông (hoặc $3\time2$ ô vuông).\\
	Cách di chuyển của quân Vua: Vua có thể di chuyển sang bất kỳ ô liền kề chung cạnh hoặc chung đỉnh xung quanh nó theo hướng ngang, dọc hoặc chéo.
	\shortans{$71$}
	\loigiai{
		\immini[thm]{
			Để dễ lập luận về sau, ta sẽ gọi hình vuông được tạo thành bởi vị trí ban đầu của Vua và các ô vuông xung quanh vị trí ban đầu này là hình vuông chứa Vua.\\
			Các bước đi của Mã và Vua sẽ được gọi là M1, M2, M3, M4, V1, V2, V3. Vị trí ban đầu của Mã và Vua gọi là M và V.\\
			Bây giờ ta sẽ tính số cách mà bộ Vua Mã có thể đi, đây chính là không gian mẫu.\\
			M1 có $8$ cách đi và tạo thành $2$ nhóm. Nhóm $1$ gồm $4$ vị trí mà cách lề bàn cờ nhiều hơn $1$ ô vuông, nhóm $2$ gồm $4$ vị trí cách ít nhất $1$ lề bàn cờ chỉ $1$ ô vuông.\\
			Để M4 về lại vị trí ban đầu thì bốn vị trí của Mã phải tạo thành hình thoi, khi đó M, M1, M2 phải không thẳng hàng.
		}
		{
			\begin{tikzpicture}[x=0.8cm, y=0.8cm]
				% 1. Định nghĩa màu sắc
				\definecolor{lightsq}{HTML}{F0D9B5} 
				\definecolor{darksq}{HTML}{B58863}  
				% 2. Vẽ bàn cờ
				\foreach \x in {0,...,7}
				\foreach \y in {0,...,7}
				{
					\pgfmathparse{mod(\x+\y,2) ? "darksq" : "lightsq"}
					\fill[\pgfmathresult] (\x,\y) rectangle (\x+1,\y+1);
				}
				% 3. Khung viền kép (Double border)
				\draw[thick] (0,0) rectangle (8,8);
				\draw[black, thin] (-0.8,-0.8) rectangle (8.8,8.8);
				\draw[black, thick] (-0.9,-0.9) rectangle (8.9,8.9);
				% 4. Tọa độ A-H
				\foreach \x [count=\i from 0] in {A,B,C,D,E,F,G,H} {
					\node at (\i+0.5, -0.4) {\small\textbf{\x}};
					\node at (\i+0.5, 8.4) {\small\textbf{\x}};
				}
				% 5. Tọa độ 1-8
				\foreach \y [count=\i from 0] in {1,2,3,4,5,6,7,8} {
					\node at (-0.4, \i+0.5) {\small\textbf{\y}};
					\node at (8.4, \i+0.5) {\small\textbf{\y}};
				}
				% Quân Vua Đen (Black King)
				\node at (5.5, 5.5) {\fontsize{35}{35}\selectfont \BlackKingOnWhite};
				% Khung bao từ góc dưới trái E5 (4,4) đến góc trên phải G7 (7,7)
				\draw[red, ultra thick] (4,4) rectangle (7,7);
				%Vẽ các quân cờ
				\node at (3.5, 3.5) {\fontsize{35}{35}\selectfont \BlackKnightOnWhite};
				\draw[solid, green!50!black, ultra thick] (3.5, 3.5) circle (0.42);
				% List tọa độ: B3(1,2), B5(1,4), C2(2,1), C6(2,5), E2(4,1), E6(4,5), F3(5,2), F5(5,4).
				% Cộng thêm 0.5 để vào tâm ô.
				\foreach \cx/\cy in {1.5/2.5, 1.5/4.5, 2.5/1.5, 2.5/5.5, 4.5/1.5, 4.5/5.5, 5.5/2.5, 5.5/4.5} {
					% 1. Vòng tròn nét đứt màu xám bao quanh ô đích
					\draw[dashed, black!50] (\cx, \cy) circle (0.35);
					% 2. Quân mã đen nhỏ trên ô tối
					\node at (\cx, \cy) {\fontsize{18}{18}\selectfont \BlackKnightOnBlack};
				}
			\end{tikzpicture}
		}
		\begin{itemize}
			\item  \textbf{Trường hợp 1:} M1 thuộc nhóm $1$. Khi đó M2 có $6$ cách đi do phải khác M và M, M1, M2 không thẳng hàng.
			\item  \textbf{Trường hợp 2:} M1 thuộc nhóm $2$. Khi đó M2 có $5$ cách đi do có $1$ vị trí nằm bên ngoài bàn cờ.
		\end{itemize}
		Vậy Mã có $4 \cdot 6 + 4 \cdot 5 = 44$ cách di chuyển.\\
		Để Vua có thể về vị trí đầu sau $3$ lần đi thì các bước đi của Vua phải nằm trong hình vuông chứa Vua đã nói ban đầu. Do đó, V1 có $8$ cách đi và chia thành $2$ nhóm.\\
		Nhóm $1$ gồm các ô vuông nhỏ chung đỉnh với hình vuông ban đầu và nhóm $2$ gồm các ô vuông nhỏ chung cạnh với hình vuông ban đầu.
		\begin{itemize}
			\item  \textbf{Trường hợp 1:} V1 thuộc nhóm $1$. Khi đó V2 có $2$ cách đi.
			\item  \textbf{Trường hợp 2: } V1 thuộc nhóm $2$. Khi đó V2 có $4$ cách đi.
		\end{itemize}
		Vậy Vua có $4 \cdot 2 + 4 \cdot 4 = 24$ cách.\\
		Do đó $n(\Omega) = 24 \cdot 44 = 1\,056$.\\
		Bây giờ ta sẽ tính số cách đi mà Mã và Vua có thể gặp nhau ít nhất $1$ lần.
		\begin{itemize}
			\item \textbf{Trường hợp 1:} M1 gặp V1.\\
			Khi đó M1 sẽ ở vị trí E6 hoặc F5. Với mỗi vị trí của M1 thì có $6$ cách chọn vị trí M2 nên Mã có $12$ cách di chuyển.\\
			Vua đi sau mã nên V1 sẽ phải đi vào ô M1. Với mỗi vị trí của V1 thì có $4$ cách chọn ví trí cho V2.\\
			Vậy số cách đi thoả mãn trường hợp này là $12 \cdot 4 =48$ cách.
			\item \textbf{Trường hợp 2:} M2 gặp V2.\\
			Để điều này xảy ra thì M2 phải nằm trong hình vuông chứa Vua nên M1 có $4$ cách chọn vị trí. Khi đó M2 có $2$ cách chọn vị trí tương ứng và điều nằm ở $4$ đỉnh hình vuông chứa Vua.\\
			Để V2 rơi vào ô M2 thì V1 phải nằm kề với ô mà M2 sẽ đi tới nên V1 có $2$ cách chọn vị trí.\\
			Vậy số cách đi thoả mãn trường hợp này là $4 \cdot 2 \cdot 2 = 16$ cách.
			\item  \textbf{Trường hợp 3:} M2 gặp V1.\\
			Cách di chuyển của M2 trong trường hợp này giống hệt \textbf{trường hợp 2}.\\
			Để V1 rơi vào ô M2 thì V1 sẽ đi vào ô mà M2 sẽ đi, khi đó V1 nằm ở  đỉnh hình vuông chứa Vua nên V2 sẽ có $2$ cách đi. \\
			Vậy có $4 \cdot 2 \cdot 2 = 16$ cách.
			\item  \textbf{Trường hợp 4:} M3 gặp V2.\\
			Do tính đối xứng nên trường hợp này thật ra là đi theo chiều ngược lại của trường hợp 1 nên cũng có $48$ cách.\\
			Tuy nhiên, có $2$ cách đi mà M1 gặp V1 thuộc trường hợp này.
		\end{itemize}
		Vậy số cách đi thoả mãn là $48+16+16+48-2=126$.\\
		Vậy xác suất để Mã và Vua gặp nhau ít nhất 1 lần là $\dfrac{126}{1\,056} = \dfrac{21}{176}$.\\
		Do đó $b-5a = 71$.
	}
\end{ex}

\begin{ex}%[0D0C2-7]%[Dự án C NH25-26 - Đợt 3 - Hector Tran]%[SGD-Tuyên Quang - Lần 1]%[GVPB: Thanh Phát]
	Gọi $S$ là tập hợp tất cả các số tự nhiên có $6$ chữ số khác nhau được lập từ các chữ số $1$, $2$, $3$, $4$, $5$, $6$. Lấy ngẫu nhiên một số thuộc $S$, gọi $T$ là xác suất số lấy được là số lẻ đồng thời tổng của ba chữ số đầu lớn hơn tổng của ba chữ số cuối một đơn vị. Giá trị của $230T$ bằng bao nhiêu?
	\par\shortans{23}
	\loigiai{
		Số phần tử của tập hợp $S$ tất cả các số tự nhiên có $6$ chữ số khác nhau được lập từ các chữ số $1$, $2$, $3$, $4$, $5$, $6$ là $n(\Omega) = 6! = 720$.\\
		Gọi $A$ là biến cố chọn từ $S$ được một số lẻ đồng thời tổng của ba chữ số đầu lớn hơn tổng của ba chữ số cuối một đơn vị.\\
		Giả sử số được chọn có dạng $\overline{a_1a_2a_3a_4a_5a_6}$. Theo giả thiết
		\begin{itemize}
			\item Số là số lẻ nên $a_6 \in \{1; 3; 5\}$.
			\item Tổng tất cả các chữ số là $1 + 2 + 3 + 4 + 5 + 6 = 21$.
			\item Gọi $S_1 = a_1 + a_2 + a_3$ là tổng ba chữ số đầu và $S_2 = a_4 + a_5 + a_6$ là tổng ba chữ số cuối.
		\end{itemize}
		Ta có hệ phương trình
		\[\heva{&S_1 + S_2 = 21 \\ &S_1 = S_2 + 1} \Leftrightarrow \heva{&S_1 = 11 \\ &S_2 = 10.}\]
		Các bộ gồm $3$ chữ số có tổng bằng $10$ từ tập $\{1; 2; 3; 4; 5; 6\}$ là $\{1; 3; 6\}$; $\{1; 4; 5\}$; $\{2; 3; 5\}$.
		\begin{itemize}
			\item Trường hợp $1$. $\{a_4; a_5; a_6\} = \{1; 3; 6\}$. Khi đó bộ còn lại là $\{a_1; a_2; a_3\} = \{2; 4; 5\}$.\\
			Vì số là số lẻ nên $a_6 \in \{1; 3\}$ ($2$\ cách chọn).\\
			Hai vị trí $a_4$; $a_5$ có $2!$ cách hoán vị.\\
			Ba vị trí $a_1$; $a_2$; $a_3$ có $3!$ cách hoán vị.\\
			Số các số thỏa mãn là $2 \cdot 2! \cdot 3! = 24$.
			\item Trường hợp $2$. $\{a_4; a_5; a_6\} = \{1; 4; 5\}$. Khi đó bộ còn lại là $\{a_1; a_2; a_3\} = \{2; 3; 6\}$.\\
			Vì số là số lẻ nên $a_6 \in \{1; 5\}$ ($2$\ cách chọn).\\ 
			Tương tự trường hợp $1$; có $2 \cdot 2! \cdot 3! = 24$ số.
			\item Trường hợp $3$. $\{a_4; a_5; a_6\} = \{2; 3; 5\}$. Khi đó bộ còn lại là $\{a_1; a_2; a_3\} = \{1; 4; 6\}$.\\
			Vì số là số lẻ nên $a_6 \in \{3; 5\}$ ($2$\ cách chọn).\\ 
			Tương tự trường hợp $1$; có $2 \cdot 2! \cdot 3! = 24$ số.
		\end{itemize}
		Tổng số kết quả thuận lợi cho biến cố $A$ là $n(A)=24 + 24 + 24 = 72$.\\
		Xác suất $T =\dfrac{n(A)}{n(\Omega)}= \dfrac{72}{720} = \dfrac{1}{10} = 0{,}1$.\\
		Vậy $230T = 230 \cdot 0{,}1 = 23$.
	}
\end{ex}

\begin{ex}%[0D0C2-9]%[Dự án C đề thi thử TN Môn Toán NH25-26 - Dot 4 - Phạm Hoàng Đăng]%[THPT Nguyễn Trãi-Hà Nội]%[GVPB - Khắc Thiên]
	\immini[thm]{Chọn ngẫu nhiên $5$ số phân biệt từ tập hợp $X=\{2; 4; 6; 8; 10; 12; 13; 14; 16; 18; 20\}$ để xếp vào $5$ ô $A$, $B$, $C$, $D$, $O$ như hình vẽ (mỗi ô xếp một số). Gọi $Y$ là biến cố \lq\lq Số $13$ xếp ở ô $O$ và tổng các số xếp ở các ô $A$, $O$, $C$ bằng tổng các số xếp ở các ô $B$, $O$, $D$\rq\rq. Biết xác suất của biến cố $Y$ là $\mathrm{P}(Y)=\dfrac{a}{b}$ (với $a$, $b \in \mathbb{N}$, $\dfrac{a}{b}$ là phân số tối giản). Khi đó giá trị $a+b$ bằng bao nhiêu?}{
		\begin{tikzpicture}[scale=1,>=stealth, line join=round, line cap=round, smooth]
			\def\R{1.5}
			\path (0,0) coordinate (O)
			(90:\R) coordinate (A)
			(180:\R) coordinate (B)
			(270:\R) coordinate (C)
			(0:\R) coordinate (D);
			\draw[blue!60!black] (O) -- (A) (O) -- (B) (O) -- (C) (O) -- (D);
			\foreach \p in {O,A,B,C,D} {
				\node[draw=green!60!black, circle, fill=white, inner sep=4pt, font=\large\itshape] at (\p) {\p};
			}
		\end{tikzpicture}
	}
	\par
	\shortans[oly]{698}
	\loigiai{
		Tập $X$ có $11$ phần tử.\\
		Chọn ngẫu nhiên $5$ phần tử từ tập $X$ và xếp vào $5$ vị trí.\\ Khi đó $n(\Omega)=\mathrm{A}_{11}^5=55\,440$.\\
		Xét biến cố $Y$: \lq\lq Số $13$ xếp ở ô $O$ và tổng các số xếp ở các ô $A$, $O$, $C$ bằng tổng các số xếp ở các ô $B$, $O$, $D$\rq\rq.
		Ta có:\\
		Số $13$ được xếp ở ô $O$ có $1$ (cách).\\
		Các số xếp ở các ô $A$, $B$, $C$, $D$ thỏa mãn $A+C=B+D=S$.\\ Ta chọn các bộ thỏa mãn:
		\begin{itemize}
			\item $S=10 \Rightarrow \{2;8\}$, $\{4;6\}$ có $1$ (bộ).
			\item $S=12 \Rightarrow \{2;10\}$, $\{4;8\}$ có $1$ (bộ).
			\item $S=14 \Rightarrow \{2;12\}$ và $\{4;10\}$, $\{2;12\}$ và $\{6;8\}$, $\{4;10\}$ và $\{6;8\}$ có $3$ (bộ).
			\item $S=16 \Rightarrow \{2;14\}$ và $\{6;10\}$, $\{2;14\}$ và $\{4;12\}$, $\{4;12\}$ và $\{6;10\}$ có $3$ (bộ).
			\item $S=18 \Rightarrow \{2;16\}$, $\{4;14\}$, $\{6;12\}$, $\{8;10\}$ suy ra chọn $2$ trong $4$ cặp có $\mathrm{C}_4^2=6$ (bộ).
			\item $S=20 \Rightarrow \{2;18\}$, $\{4;16\}$, $\{6;14\}$, $\{8;12\}$ suy ra chọn $2$ trong $4$ cặp có $\mathrm{C}_4^2=6$ (bộ).
			\item $S=22 \Rightarrow \{2;20\}$, $\{4;18\}$, $\{6;16\}$, $\{8;14\}$, $\{10;12\}$ suy ra chọn $2$ trong $5$ cặp có $\mathrm{C}_5^2=10$ (bộ).
			\item $S=24 \Rightarrow \{4;20\}$, $\{6;18\}$, $\{8;16\}$, $\{10;14\}$ suy ra chọn $2$ trong $4$ cặp có $\mathrm{C}_4^2=6$ (bộ).
			\item $S=26 \Rightarrow \{6;20\}$, $\{8;18\}$, $\{10;16\}$, $\{12;14\}$ suy ra chọn $2$ trong $4$ cặp có $\mathrm{C}_4^2=6$ (bộ).
			\item $S=28 \Rightarrow \{8;20\}$, $\{10;18\}$, $\{12;16\}$ suy ra chọn $2$ trong $3$ cặp có $\mathrm{C}_3^2=3$ (bộ).
			\item $S=30 \Rightarrow \{10;20\}$, $\{12;18\}$, $\{14;16\}$ suy ra chọn $2$ trong $3$ cặp có $\mathrm{C}_3^2=3$ (bộ).
			\item $S=32 \Rightarrow \{12;20\}$, $\{14;18\}$ có $1$ (bộ).
			\item $S=34 \Rightarrow \{14;20\}$, $\{16;18\}$ có $1$ (bộ).
		\end{itemize}
		Tổng số các bộ thỏa mãn là $1+1+3+3+6+6+10+6+6+3+3+1+1=50$ (bộ).\\
		Ứng với mỗi bộ ta chọn $1$ cặp xếp vào ô $A-C$, cặp còn lại cho ô $B-D$ có $2$ (cách).\\
		Hoán vị $2$ số trong cặp $A$, $C$ có $2$ (cách).\\
		Hoán vị $2$ số trong cặp $B$, $D$ có $2$ (cách).\\
		Suy ra mỗi bộ có $2 \cdot 2 \cdot 2=8$ (cách).\\
		Số phần tử thuận lợi của biến cố $Y$ là $n(Y)=50 \cdot 8=400$.\\
		Do đó $\mathrm{P}(Y)=\dfrac{n(Y)}{n(\Omega)}=\dfrac{400}{55\,440}=\dfrac{5}{693} \Rightarrow a+b=698$.
	}
\end{ex}

\begin{ex}%[0D0C2-9]%[Dự án C đợt 3 NH25-26- Phạm Đăng Đăng]%[TT-THPT-TranNhanTong-HaNoi-N25-26]%[GVPB - Phan Quốc Trí]
	Kết thúc học kì $1$, cô giáo chuẩn bị $8$ hộp bút bi, $7$ tập vở viết và $5$ bộ dụng cụ vẽ hình (thước kẻ, compa, êke, thước đo độ) để làm phần thưởng cho $10$ học sinh có kết quả xuất sắc nhất lớp. Mỗi học sinh nhận thưởng sẽ được $2$ phần thưởng khác loại. Trong số $10$ học sinh trên có $2$ học sinh tên Dũng và Nam. Tìm xác suất để Dũng và Nam có phần thưởng giống nhau.   
	\shortans{$0{,}31$}
	\loigiai{
		Giả sử số phần thưởng gồm $1$ hộp bút và $1$ tập vở là $x$ ($x \in \mathbb{N}$).\\
		Số phần thưởng gồm $1$ tập vở và $1$ bộ dụng cụ là $y$ ($y \in \mathbb{N}$).\\
		Số phần thưởng gồm $1$ hộp bút và $1$ bộ dụng cụ là $z$ ($z \in \mathbb{N}$).\\
		Vì cô có $8$ hộp bút bi, $7$ tập vở viết và $5$ bộ dụng cụ nên
		$
		\heva{&x + z = 8 \\ &x + y = 7 \\ &y + z = 5} \Leftrightarrow \heva{&x = 5 \\ &y = 2 \\ &z = 3.}
		$\\
		Tổng số cách tặng thưởng là $\mathrm{C}_{10}^5 \cdot \mathrm{C}_5^2 \cdot \mathrm{C}_3^3 = 2\,520$.\\
		Để các bạn nhận thưởng sao cho Dũng và Nam nhận thưởng như nhau, ta chia thành các trường hợp:
		\begin{itemize}
			\item \textbf{TH1:} Dũng và Nam nhận $1$ bút và $1$ vở, còn lại $3$ phần gồm $1$ bút và $1$ vở, $2$ phần $1$ vở và $1$ dụng cụ, $3$ phần $1$ bút và $1$ dụng cụ. Ta chọn $3$ phần gồm $1$ bút và $1$ vở cho $8$ bạn, chọn $2$ phần gồm $1$ tập vở và $1$ bộ dụng cụ cho $5$ bạn và chọn $3$ phần gồm $1$ hộp bút và $1$ bộ dụng cụ cho $3$  bạn còn lại.\\ Khi đó
			số cách chọn ở trường hợp này là $\mathrm{C}_8^3 \cdot \mathrm{C}_5^2 \cdot \mathrm{C}_3^3 = 560$.
			\item \textbf{TH2:} Dũng và Nam nhận $1$ vở và $1$ dụng cụ, còn lại $5$ phần gồm $1$ bút và $1$ vở, $3$ phần $1$ bút và $1$ dụng cụ nên có số cách phát $\mathrm{C}_8^5 \cdot \mathrm{C}_3^3 = 56$.
			\item 	\textbf{TH3:} Dũng và Nam nhận $1$ bút và $1$ dụng cụ, còn lại $5$ phần gồm $1$ bút và $1$ vở, $2$ phần $1$ vở và $1$ dụng cụ, 1 phần bút và dụng cụ nên có số cách phát $\mathrm{C}_8^5 \cdot \mathrm{C}_3^2 = 168$.
		\end{itemize}
		Suy ra tổng số cách phát phần thưởng là $560 + 56+168 = 784$.\\
		Vậy	xác suất cần tìm  $\dfrac{784}{2\,520} \approx 0{,}31$.}
\end{ex}

\begin{ex} %[0D0V2-5]%[TN-THPT-NguyenTrungThien-HaTinh-NH25-26]%[Loc Do]%[PB: Phan Anh]
	Một người môi giới bất động sản có $8$ chìa khoá để mở $8$ ngôi nhà mới. Mỗi chìa chỉ mở được đúng $1$ ngôi nhà. Biết có $3$ ngôi nhà thường không khoá cửa, người môi giới chọn ngẫu nhiên $3$ chìa khoá mang theo. Hỏi nếu người môi giới chọn ngẫu nhiên $1$ ngôi nhà để vào, thì xác suất để người môi giới này có thể vào được là bao nhiêu (\textit{kết quả làm tròn đến hàng phần trăm})?
	\shortans{0,61}
	\loigiai{
		Chọn $1$ nhà từ $8$ nhà có $8$ cách, sau đó chọn $3$ chìa từ $8$ chìa có $\mathrm{C}_{8}^{3}$ cách.\\
		Suy ra số phần tử của không gian mẫu là
		\[n(\Omega)=8\cdot \mathrm{C}_{8}^{3} =448.\]
		Gọi $A$ là biến cố: \lq\lq Để người môi giới này có thể vào được\rq\rq.
		\begin{itemize}
			\item Trường hợp $1$: Chọn đúng một trong ba nhà không khoá cửa, chọn $3$ chìa từ $8$ chìa bất kỳ (khi đó không cần quan tâm chìa chọn như thế nào). Nên có $3\cdot\mathrm{C}_{8}^{3} =168$ cách
			\item Trường hợp $2$: Chọn đúng nhà có khoá cửa.
			\begin{itemize}
				\item Chọn $1$ nhà từ $5$ nhà có $5$ cách;
				\item Chọn $1$ chìa đúng, và $2$ chìa còn lại từ $7$ chìa có: $1\cdot\mathrm{C}_{7}^{2} =21$ cách.
			\end{itemize}
		\end{itemize}
		Suy ra có $5\cdot 21=105$ cách.\\
		Ta có số phần tử của biến cố $A$ là $n(A)=168+105=273$.
		\[\Rightarrow \mathrm{P}(A)=\dfrac{n(A)}{n(\Omega)} =\dfrac{39}{64} \approx 0{,}61.\]}
\end{ex}

\begin{ex}%[0D0V2-7]%[Dự án đề thi thử NH25-26]%[THPT Chuyên Phan Bội Châu - Nghệ An]%[Sỹ Trường]
	Trong một hộp kín đựng $20$ viên bi được đánh số từ $1$ đến $20$. Bạn An chọn ngẫu nhiên $4$ viên bi trong hộp, xác suất để $4$ số trên $4$ viên bi được chọn lập thành một cấp số cộng là $\dfrac{1}{a}$. Tính $a$.
	\shortans{$85$}
	\loigiai{
		Số cách chọn $4$ viên bi từ hộp là $\mathrm{C}_{20}^4 = 4\,845$.\\
		Gọi cấp số cộng tạo thành là $a$, $b$, $c$, $d$. \\
		Ta có công sai $x = \dfrac{d-a}{3}$, do đó $a$, $d$ có cùng số dư khi chia cho $3$.\\
		Từ $1$ đến $20$ ta có $6$ số có dạng $3k$, có $7$ số có dạng $3k+1$ và $7$ số có dạng $3k+2$ với $k \in \mathbb{N}$.\\
		Mỗi cặp $a$, $d$ (với $a<d$) thỏa mãn tính chất trên sẽ xác định duy nhất hai số $b$, $c$ tạo thành cấp số cộng gồm $4$ phần tử trong tập hợp đã cho.\\
		Suy ra số cấp số cộng bằng số cách chọn cặp $a$, $d$ và bằng $\mathrm{C}_6^2+\mathrm{C}_7^2+\mathrm{C}_7^2 = 57$.\\
		Vậy xác suất là $\mathrm{P}(A) = \dfrac{57}{4\,845} = \dfrac{1}{85} \Rightarrow a=85$.
	}
\end{ex}

\begin{ex}%[0D0V2-8]%[Dự án C đề thi thử TN Môn Toán NH25-26 - Dot 4 - Vũ Hồng Toàn]%[THPT Mai Anh Tuấn - Thanh Hóa]%[GVPB - Lê Hữu Kiệt]
	Có $8$ bạn cùng ngồi xung quanh một cái bàn tròn, mỗi bạn cầm một đồng xu như sau. Tất cả $8$ bạn cùng tung đồng xu của mình, bạn có đồng xu ngửa thì đứng, bạn có đồng xu sấp thì ngồi. Xác suất để không có hai bạn liền kề cùng đứng là bao nhiêu {\it (kết quả được làm tròn đến hàng phần trăm)}?
	\par\shortans{0,18}
	\loigiai{
		Gọi $A$ là biến cố không có hai người liền kề cùng đứng.\\
		Số phần tử của không gian mẫu là $n(\Omega) = 2^8 = 256$.\\
		Rõ ràng nếu nhiều hơn $4$ đồng xu ngửa thì biến cố $A$ không xảy ra.        \\
		Để biến cố $A$ xảy ra có các trường hợp sau
		\begin{enumerate}[{\bf TH 1.}]
			\item Có nhiều nhất $1$ đồng xu ngửa.
			\begin{itemize}
				\item $0$ ngửa có $1$ cách.
				\item $1$ ngửa: chọn $1$ trong $8$ vị trí có $8$ cách.
			\end{itemize}
			\item Có $2$ đồng xu ngửa.\\
			Hai đồng xu ngửa kề nhau: có $8$ khả năng.\\
			Suy ra số kết quả của trường hợp này là $\mathrm{C}_8^2 - 8 = 20$.
			\item Có $3$ đồng xu ngửa.\\
			Cả $3$ đồng xu ngửa kề nhau: có $8$ kết quả.\\
			Trong $3$ đồng xu ngửa, có đúng một cặp kề nhau có $8\cdot 4 = 32$ cách.\\
			Suy ra số kết quả của trường hợp này là $\mathrm{C}_8^3 - 8 - 32 = 16$ cách.
			\item Có $4$ đồng xu ngửa.\\
			Trường hợp này có $2$ kết quả thỏa mãn biến cố $ A$ xảy ra.
		\end{enumerate}
		Suy ra $n(A) = 9 + 20 + 16 + 2 = 47$.\\
		Xác suất để không có hai bạn liền kề cùng đứng là $\mathrm{P}(A) = \dfrac{n(A)}{n(\Omega)} = \dfrac{47}{256} \approx 0{,}18$.
	}
\end{ex}

\begin{ex}%[0D0V2-9]%[TT-LienTruong-HaNoi-N25-26]%[Phạm Quốc Toàn - Lê Phúc]
	\immini[thm]{Cho một bảng ô vuông $3\times 3$ như hình vẽ. Điền ngẫu nhiên $9$ số thuộc tập hợp $X=\left\{0;1;2;3;4;5;6;7;8;9\right\}$ vào $9$ ô vuông trong bảng (mỗi ô điền một số khác nhau). Gọi $Y$ là biến cố \lq\lq Mỗi hàng, mỗi cột bất kì trong bảng đều có ít nhất một số lẻ\rq\rq. Biết xác suất $P(Y)=\dfrac{a}{b}$ (với $a$, $b\in \mathbb{N}$ và $\dfrac{a}{b}$ là phân số tối giản). Khi đó $a+b$ bằng bao nhiêu?}
	{	\begin{tikzpicture}[font=\footnotesize, line join=round, line cap=round, >=stealth,scale=0.7]
			%%%%%%%%%%%%%%%%%%%%%%%%%%%
			\draw[thick] (0,0) -- (3,0) -- (3,3) -- (0,3) -- (0,0) (1,0) -- (1,3) (2, 0) -- (2,3) (0, 1) -- (3, 1) (0,2) -- (3,2);
			%%%%%%%%%%%%%%%%%%%%%%%%%%%
	\end{tikzpicture}}
	\par\shortans[0]{43}
	\loigiai{
		Số phần tử của không gian mẫu là $n(\Omega)=\mathrm{A}_{10}^9$.\\
		Ta có $\overline{Y}$ là biến cố \lq\lq Có ít nhất một hàng hoặc một cột toàn chữ số chẵn\rq\rq.\\
		Tập $X$ gồm $5$ chữ số chẵn và $5$ chữ số lẻ nên ta có hai trường hợp sau.
		\begin{itemize}
			\item \textbf{Trường hợp $1$.} Bảng được xếp $5$ chữ số lẻ và $4$ chữ số chẵn\\
				Do chỉ có $4$ chữ số chẵn nên chỉ có tối đa đúng $1$ cột hoặc $1$ hàng toàn số chẵn.
			\begin{itemize}
				\item Chọn bất kì một hàng hoặc một cột để xếp toàn số chẵn có $6$ cách.
				\item Chọn $5$ trong $6$ ô còn lại và xếp $5$ số lẻ vào có $\mathrm{A}_6^5$ cách.
				\item Chọn $4$ trong $5$ số chẵn và xếp vào $4$ ô còn lại có $\mathrm{A}_5^4$ cách.
			\end{itemize}
			Trường hợp này có $6 \cdot \mathrm{A}_6^5 \cdot \mathrm{A}_5^4=518\,400$ cách xếp thỏa mãn yêu cầu đề bài.
			\item \textbf{Trường hợp $2$.} Bảng được xếp $4$ chữ số lẻ và $5$ chữ số chẵn
			\begin{itemize}
				\item Chọn bất kì một hàng hoặc một cột để xếp toàn số chẵn có $6$ cách.
				\item Chọn $4$ trong $6$ ô còn lại để xếp số lẻ vào có $\mathrm{C}_6^4$ cách.
				\item Chọn $4$ trong $5$ số lẻ và xếp vào $4$ ô đã chọn có $\mathrm{A}_5^4$ cách.
				\item Xếp nốt $5$ số chẵn vào $5$ ô còn lại có $5!$ cách.
			\end{itemize}
			Do $5$ số chẵn có thể xếp đầy một hàng và một cột (chung nhau một ô) nên ta đã đếm lặp trường hợp có cả một cột và một hàng toàn số chẵn, trường hợp này có $3 \cdot 3 \cdot \mathrm{A}_5^4 \cdot 5!$ cách.\\
			Suy ra trường hợp này có $6\cdot \mathrm{C}_6^4 \cdot \mathrm{A}_5^4\cdot 5!-3 \cdot 3\cdot \mathrm{A}_5^4 \cdot 5!=1\,166\,400$ cách.
		\end{itemize} 
		Vậy xác suất cần tính là $\mathrm{P}(Y)=1-\mathrm{P} \left(\overline{Y} \right)=1-\dfrac{518\,400+1\,166\,400}{\mathrm{A}_{10}^9}=\dfrac{15}{28}$.\\
		Khi đó $a+b=43$.
	}
\end{ex}

\begin{ex}%[0D1V3-5]%[Dự án C đề thi thử TN Môn Toán NH25-26 - Dot 4 - Hector Tran]%[Liên trường THPT cụm 09 - Đợt 2 - Hà Nội]%[GVPB - Thanh Phát]
	Lớp 10C có $45$ học sinh trong đó có $25$ em thích môn Văn, $20$ em thích môn Toán, $18$ em thích môn Sử, $6$ em không thích môn nào, $5$ em thích cả ba môn. Có bao nhiêu học sinh thích chỉ một trong ba môn học trên?
	\par\shortans{20}
	\loigiai{
		\immini {
			Gọi $a$, $b$, $c$ theo thứ tự là số học sinh chỉ thích môn Văn, Sử, Toán.\\
			Gọi $x$ là số học sinh chỉ thích hai môn là Văn và Toán.\\
			Gọi $y$ là số học sinh chỉ thích hai môn là Sử và Toán.\\
			Gọi $z$ là số học sinh chỉ thích hai môn là Văn và Sử.\\
			Ta có số em thích ít nhất một môn là $45 - 6 = 39$.\\
			Dựa vào biểu đồ Ven ta có hệ phương trình
			\[\heva{&a + x + z + 5 = 25&\quad (1) \\ &b + y + z + 5 = 18&\quad (2) \\ &c + x + y + 5 = 20&\quad (3) \\ &x + y + z + a + b + c + 5 = 39.& \quad (4)}\]
		}{
			\begin{tikzpicture}[scale=0.6, node/.style={scale=0.8}, declare function={
					xV=0; yV=0; rV=2.8;
					xT=2.6; yT=0.8; rT=1.9;
					xS=2.2; yS=-2.1; rS=2.3;
					nx25=-0.6; ny25=1.4;
					nxa=-0.6; nya=-0.2;
					nx20=3.2; ny20=1.5;
					nxc=3.6; nyc=0.6;
					nx18=2.2; ny18=-3.8;
					nxb=2.2; nyb=-2.1;
					nxx=1.4; nyx=1.1;
					nxy=3.1; nyy=-0.45;
					nxz=1; nyz=-1.3;
					nx5=1.9; ny5=-0.4;
				}]
				\tikzset{every node/.style={scale=0.8}}  
				\draw[line width=0.8pt] (xV, yV) circle (rV);
				\draw[line width=0.8pt] (xT, yT) circle (rT);
				\draw[line width=0.8pt] (xS, yS) circle (rS);				
				\node at (nx25, ny25) {$25(V)$};
				\node at (nxa, nya) {$a$};				
				\node at (nx20, ny20) {$20(T)$};
				\node at (nxc, nyc) {$c$};				
				\node at (nx18, ny18) {$18(S)$};
				\node at (nxb, nyb) {$b$};				
				\node at (nxx, nyx) {$x$};
				\node at (nxy, nyy) {$y$};
				\node at (nxz, nyz) {$z$};				
				\node at (nx5, ny5) {$5$};				
			\end{tikzpicture}
		}	\noindent
		Cộng vế với vế của phương trình $(1)$, $(2)$, $(3)$ ta có
		\begin{eqnarray*}
			&&(a + b + c) + 2(x + y + z) + 15 = 25 + 18 + 20 \\
			&\Leftrightarrow& (a + b + c) + 2(x + y + z) = 48. \quad (5)
		\end{eqnarray*}
		Từ $(4)$ ta có $x + y + z = 39 - 5 - (a + b + c) = 34 - (a + b + c)$.\\
		Thay $x + y + z = 34 - (a + b + c)$ vào $(5)$ ta được
		\begin{eqnarray*}
			&&(a + b + c) + 2[34 - (a + b + c)] = 48 \\
			&\Leftrightarrow& -(a + b + c) + 68 = 48 \\
			&\Leftrightarrow& a + b + c = 20.
		\end{eqnarray*}
		Vậy có $20$ em thích chỉ một trong ba môn trên.
	}
\end{ex}

\begin{ex}%[Đề thi thử môn Toán TRƯỜNG THPT Cổ Loa - Hà Nội năm học 2025-2026]%[Phạm Quốc Toàn - Phạm Hoàng Việt]%[0D2H1-3]
	Mỗi tuần, một cửa hàng bán điện thoại di động trung bình bán được $1\,000$ chiếc điện thoại $A$ với giá $14$ triệu đồng một chiếc. Biết rằng, nếu cứ giảm giá bán $500$ nghìn đồng một chiếc thì số lượng điện thoại $A$ bán ra sẽ tăng thêm $100$ chiếc mỗi tuần. Mặt khác, để bán được $n$ chiếc điện thoại $A$ mỗi tuần thì chi phí hàng tuần của cửa hàng là $C(n)=12\,000-3n$ (triệu đồng). Để lợi nhuận thu được lớn nhất, cửa hàng nên bán mỗi chiếc điện thoại $A$ với giá bao nhiêu triệu đồng?
	\par\shortans[0]{8}
	\loigiai{
		Gọi $x$ ($x\in \left[0;14\right]$; triệu đồng) là giá bán một chiếc điện thoại $A$.\\
		Đổi $500$ nghìn đồng bằng $0{,}5$ triệu đồng.\\
		Số lần giảm giá là $\dfrac{14-x}{0{,}5}$ .\\
		Số điện thoại bán ra là $1\,000+\dfrac{14-x}{0{,}5} \cdot 100=-200 x+3800=n$.\\
		Doanh thu của cửa hàng là $(-200x+3\,800) x$ (triệu đồng).\\
		Chi phí của cửa hàng là $12\,000-3(-200x+3\,800)=600x+600$ (triệu đồng).\\
		Lợi nhuận của cửa hàng là $f(x)=(-200x+3\,800)x-(600x+600)=-200x^2+3\,200x-600$. \\ 
		Ta có $f(x)$ là hàm số bậc hai nên đạt giá trị lớn nhất tại $x=-\dfrac{b}{2a}=8$.\\
		Để lợi nhuận thu được lớn nhất, cửa hàng nên bán mỗi chiếc điện thoại $A$ với giá là $8$ triệu đồng.
	}
\end{ex}

\begin{ex}%[0D2H2-3]%[Dự án C - Đợt 3 - Sở Giáo Dục Sơn La]%[Thy Nguyen Vo Diem - Nguyễn Tiến]
	Trong một cuộc thi làm bánh, mỗi đội chỉ được sử dụng tối đa $6$ kg bột mì và $4$ kg đậu để làm ra
	bánh loại I và bánh loại II. Để làm ra một chiếc bánh loại I cần $0{,}06$ kg bột mì và $0{,}08$ kg đậu, để làm ra một chiếc bánh loại II cần $0{,}08$ kg bột mì và $0{,}04$ kg đậu. Mỗi chiếc bánh loại I được $8$ điểm, mỗi chiếc bánh loại II được $10$ điểm. Số điểm tối đa có thể đạt được của mỗi đội bằng bao nhiêu?
	\shortans{760}
	\loigiai{
		Gọi $x$, $y$ lần lượt là số bánh loại I và loại II đội làm được ($x$, $y \ge 0$; $x$, $y \in \mathbb{Z}$).\\
		Theo giả thiết, ta có hệ bất phương trình $\heva{&x \ge 0\\&y \ge 0\\&0{,}06x+0{,}08y \le 6\\&0{,}08x+0{,}04y \le 4} \Leftrightarrow \heva{&x \ge 0\\&y \ge 0\\&3x+4y \le 300\\&2x+y \le 100.}$\\
		Số điểm đội đạt được là $P=8x+10y$.\\
		Miền nghiệm hệ bất phương trình là miền tứ giác $OABC$ với $O(0;0)$, $A(50;0)$, $B(20;60)$, $C(0;75)$.\\
		\begin{center}
			\begin{tikzpicture}[scale=0.7, font=\footnotesize, line join=round, line cap=round, >=stealth]
				\fill[blue!15] (0,0) -- (5,0) -- (2,6) -- (0,7.5) -- cycle;
				
				\draw[->] (-1,0) -- (11,0) node[right] {$x$};
				\draw[->] (0,-1) -- (0,11) node[above] {$y$};
				
				\draw[blue] (-1, 8.25) -- (11, -0.75) node[right] {$3x + 4y = 300$};
				
				\draw[red] (-0.5, 11) -- (5.5, -1) node[right] {$2x + y = 100$};
				
				\fill (0,0) circle (1pt) node[below left] {$O$};
				\fill (5,0) circle (1pt) node[below left] {$50$} node[above] {$A$};
				\fill (2,6) circle (1pt) (2,0) circle (1pt) (0,6) circle (1pt);
				\fill (0,7.5) circle (1pt) node[left] {$75$} node[right] {$C$};
				\fill (10,0) circle (1pt) node[below] {$100$};
				\fill (0,10) circle (1pt) node[left] {$100$};
				\draw[dashed] (2,0)node[below]{$20$}--(2,6) node[above right]{$B$}--(0,6) node[left]{$60$};
			\end{tikzpicture}
		\end{center}
		Giá trị của $P$ tại các đỉnh $P(0;0)=0$, $P(50;0)=400$, $P(20;60)=760$, $P(0;75)=750$.\\
		Vậy số điểm tối đa đạt được là $760$.
	}
\end{ex}

\begin{ex}%[0D2V2-3]%[TT-SGD-TuyenQuang-L1-NH25-26]%[Võ Thị Thùy Trang]%[GVPB: Thanh Phát]
	Để gây quỹ từ thiện, câu lạc bộ thiện nguyện của một trường THPT tổ chức hoạt động bán hàng với hai mặt hàng là trà sữa và bánh ngọt. Câu lạc bộ thiết kế hai thực đơn. Thực đơn $1$ có giá $40$ nghìn đồng, bao gồm hai ly trà sữa và một chiếc bánh ngọt. Thực đơn $2$ có giá $65$ nghìn đồng, bao gồm ba ly trà sữa và hai chiếc bánh ngọt. Biết rằng câu lạc bộ chỉ làm được không quá $180$ ly trà sữa và $110$ chiếc bánh ngọt. Số tiền lớn nhất mà câu lạc bộ có thể nhận được sau khi bán hết hàng bằng bao nhiêu nghìn đồng?
	\par\shortans{3800}
	\loigiai{
		Gọi $x$, $y$ lần lượt là số lượng thực đơn $1$ và thực đơn $2$ mà câu lạc bộ dự định bán ($x \ge 0$, $y \ge 0$).\\
		Theo đề bài, ta có
		\begin{itemize}
			\item Tổng số ly trà sữa cần dùng là $2x + 3y$. Do câu lạc bộ chỉ làm được không quá $180$ ly nên ta có $2x + 3y \le 180$.
			\item Tổng số chiếc bánh ngọt cần dùng là $x + 2y$. Do câu lạc bộ chỉ làm được không quá $110$ chiếc nên ta có $x + 2y \le 110$.
		\end{itemize}
		Từ đó ta có hệ bất phương trình sau $ \heva{&x \ge 0, y \ge 0\\&2x + 3y \le 180\\&x + 2y \le 110.}$\\
		Gọi $F(x,y)$ (nghìn đồng) là số tiền mà câu lạc bộ nhận được sau khi bán hết hàng.\\
		Tìm $x$, $y$ thỏa mãn hệ trên sao cho biểu thức $F(x, y) = 40x + 65y$ (nghìn đồng) đạt giá trị lớn nhất.
		\begin{center}
			\begin{tikzpicture}[>=stealth,line join=round,line cap=round,font=\footnotesize,scale=0.7]
				\draw[->] (-1,0) -- (12,0) node[above]{$x$};
				\draw (0,0) node[below left]{$O$};
				\fill[gray!20] (0,0) -- (9,0) -- (3,4) -- (0,5.5) -- cycle;
				\draw[domain=-0.5:10, samples=2] plot (\x, {6-2/3*\x}) node[right]{$2x+3y=180$};
				\draw[domain=-0.5:11.5, samples=2] plot (\x, {5.5-0.5*\x});
				\node at (9,1.1) [right]{$x+2y=110$};
				\fill (9,0) circle (1.5pt) node[below left] {$A(90;0)$};
				\fill (3,4) circle (1.5pt) node[above right] {$B(30;40)$};
				\fill (0,5.5) circle (1.5pt) node[below left] {$C(0;55)$};
				\draw[dashed] (3,0) node[below]{$30$} -- (3,4) -- (0,4) node[left]{$40$};
				\draw[->] (0,-1) -- (0,7) node[left]{$y$};
			\end{tikzpicture}
		\end{center}
		Miền nghiệm của hệ bất phương trình là miền tứ giác $OABC$ với các đỉnh $O(0;0)$, $A(90;0)$, $B(30;40)$ và $C(0;55)$.
		\begin{itemize}
			\item $F(0; 0) = 0$.
			\item $F(90; 0) = 40 \cdot 90 + 65 \cdot 0 = 3\,600$.
			\item $F(30; 40) = 40 \cdot 30 + 65 \cdot 40 = 3\,800$.
			\item $F(0; 55) = 40 \cdot 0 + 65 \cdot 55 = 3\,575$.
		\end{itemize}
		Ta thấy $F(x, y)$ đạt giá trị lớn nhất bằng $3\,800$ tại $x = 30$ và $y = 40$.\\
		Vậy số tiền lớn nhất mà câu lạc bộ có thể nhận được là $3\,800$ nghìn đồng.
	}
\end{ex}

\begin{ex}%[Đề thi thử môn Toán TRƯỜNG THPT Cổ Loa - Hà Nội năm học 2025-2026]%[Phạm Quốc Toàn - Phạm Hoàng Việt]%[0D2V2-3]
	Một cơ sở đóng thuyền thủ công cần $10$ giờ lao động để đóng một chiếc thuyền loại $A$ và $15$ giờ lao động để đóng một chiếc thuyền loại $B$. Mỗi tuần cơ sở bố trí được tối đa $120$ giờ lao động cho việc đóng hai loại thuyền này. Qua thực tế, người ta thấy mỗi tuần cơ sở bán được tối đa $6$ chiếc thuyền loại $A$ và tối thiểu $2$ chiếc thuyền loại $B$. Mỗi chiếc thuyền loại $A$ cho thuận lợi là $0{,}5$ triệu đồng và mỗi chiếc thuyền loại $B$ cho lợi nhuận là $0{,}7$ triệu đồng. Qua thời gian tìm hiểu, chủ cơ sở đã tìm được số lượng nên đóng thuyền mỗi loại để có thể thu được lợi nhuận cao nhất là $T$ (triệu đồng). Giá trị $T$ bằng bao nhiêu? 
	\par\shortans[0]{5{,}8}
	\loigiai{
		Gọi $x$, $y$ lần lượt là số thuyền loại $A$ và số thuyền loại $B$ được cơ sở đóng trong một tuần $\left(x, y\in \mathbb{N}\right)$. Theo đề bài ta có hệ bất phương trình sau 
		\[
		\heva{&10x+15y \le 120 \\&0 \le x \le 6 \\& y \ge 2} \Leftrightarrow \heva{&2x+3y\le 24 \\&0\le x\le 6 \\&y\ge 2.}.
		\]
		Hàm lợi nhuận là $T=0{,}5x+0{,}7y$.\\
		Ta có miền nghiệm của hệ bất phương trình trên là tứ giác $ABCD$ như hình vẽ 
		\begin{center}
			\begin{tikzpicture}[scale=0.8, font=\footnotesize, line join=round, line cap=round, >=stealth]
				\tikzset{every node/.style={scale=0.9}}
				\begin{scope}
					\clip (-1,-1) rectangle (13,9);
					\fill[black!35] (-2,9.33)--(14,9.33)--(14,-1.33)--cycle;
					\fill[black!35] (0,-1)--(-1,-1)--(-1,9)--(0,9)--cycle;
					\fill[black!35] (6,-1)--(13,-1)--(13,9)--(6,9)--cycle;
					\fill[black!35] (-1,2)--(-1,-1)--(13,-1)--(13,2)--cycle;
					\draw (-1.5,9)--(13.5,-1) node [pos=0.6, above, sloped] {$2x+3y-24=0$};
					\draw (6,-1)--(6,9) (-1,2)--(13,2) ;
				\end{scope}
				\draw[->] (-1,0)--(13,0) node[below]{$x$};
				\draw[->] (0,-1)--(0,9) node[left]{$y$};
				\draw (0,0) node[below left]{$O$};
				\foreach \x in {6,12}
				\draw[thin] (\x,1pt)--(\x,-1pt) node [below left] {$\x$};
				\foreach \y in {2,4,8}
				\draw[thin] (1pt,\y)--(-1pt,\y) node [below left] {$\y$};
				\draw(0,2)circle(1pt)node[above right]{$A$};
				\draw(6,2)circle(1pt)node[above left]{$B$};
				\draw(6,4)circle(1pt)node[below left]{$C$};
				\draw(0,8)circle(1pt)node[right]{$D$};
				\draw[dashed](0,4)--(6,4);
			\end{tikzpicture}
		\end{center}
		\begin{itemize}
			\item Tại $A(0;2)$, ta có $T=0{,}5 \cdot 0 +0{,}7 \cdot 2=1{,}4$.
			\item Tại $B(6;2)$, ta có $T=0{,}5\cdot 6+0{,}7 \cdot 2=4{,}4$. 
			\item Tại $C(6;4)$, ta có $T=0{,}5 \cdot 6+0{,}7 \cdot 4=5{,}8$.
			\item Tại $D(0;8)$, ta có $T=0{,}5 \cdot 0 +0{,}7 \cdot 8=5{,}6$.
		\end{itemize}
		Vậy lợi nhuận nhiều nhất là $5{,}8$ triệu đồng khi đóng $6$ thuyền loại $A$ và $4$ thuyền loại $B$ mỗi tuần.
	}
\end{ex}

\begin{ex}%[TT-THPT-HamRong-NH25-26 ]%[GV soạn: Nguyễn Quang Hiệp GVPB2 Đỗ Minh Vũ ]%[0D2V2-3]
	Một cửa hàng gạo tháng này cần nhập hai loại gạo là gạo loại I và gạo loại II. Số tiền chi cho việc nhập hai loại gạo này không quá $100$ triệu đồng. Gạo loại I mua vào với giá $30$ triệu đồng/tấn, lợi nhuận dự kiến là $2$ triệu đồng/tấn. Gạo loại II nhập vào với giá $40$ triệu đồng/tấn và lợi nhuận dự kiến là $2{,}5$ triệu đồng/tấn. Cửa hàng ước tính doanh số bán ra một tháng của hai loại gạo này không quá $3$ tấn. Hỏi cửa hàng có thể lãi nhiều nhất là bao nhiêu (triệu đồng) khi kinh doanh hai loại gạo này?
	\par\shortans[oly]{6,5}
	\loigiai{
		Gọi $x$, $y$ lần lượt là số tấn gạo loại I và loại II cần nhập trong tháng ($x \ge 0$, $y \ge 0$).\\
		Tiền lãi cửa hàng nhận được là $T = 2x + 2{,}5y$ (triệu đồng).\\
		Số tiền cửa hàng cần chi là $30x + 40y$ (triệu đồng).\\
		Tổng doanh số bán ra không quá $3$ tấn trong tháng nên ta có $x + y \le 3$.\\
		Do cửa hàng dự định chi không quá $100$ triệu đồng nên $30x + 40y \le 100 \Leftrightarrow 3x + 4y \le 10$.\\
		Bài toán đưa về tìm $x$, $y$ là nghiệm của hệ bất phương trình
		$ \heva{& x \ge 0 \\ & y \ge 0 \\ & x + y \le 3 \\ & 3x + 4y \le 10} $
		sao cho $T = 2x + 2{,}5y$ có giá trị lớn nhất.\\
		Miền nghiệm của hệ bất phương trình là miền tứ giác $OABC$ với $O(0;0)$, $A\left(0;\dfrac{5}{2}\right)$, $B(2;1)$, $C(3;0)$.
		\begin{center}
			\begin{tikzpicture}[line cap=round, line join=round, font=\footnotesize, >=stealth, scale=1.5]
				% Yêu cầu khai báo thư viện trong preamble: \usetikzlibrary{patterns}
				
				% Giới hạn của hệ trục (Bounding Box)
				\def\xmin{-1.0}
				\def\xmax{4.5}
				\def\ymin{-1.0}
				\def\ymax{4.0}
				
				% Vẽ lưới tọa độ (mỗi ô vuông 0.5 đơn vị)
				\draw[step=0.5cm, help lines, gray!40, thin] (\xmin,\ymin) grid (\xmax,\ymax);
				
				\begin{scope}
					\clip (\xmin,\ymin) rectangle (\xmax,\ymax);
					
					% === TÔ MIỀN KHÔNG NGHIỆM ===
					
					% 1. Miền KHÔNG thỏa: x >= 0 (Gạch bỏ x < 0)
					\fill[pattern=north west lines, pattern color=green!60!black]
					(\xmin,\ymin) rectangle (0,\ymax);
					\fill[opacity=0.05, green!60!black]
					(\xmin,\ymin) rectangle (0,\ymax);
					
					% 2. Miền KHÔNG thỏa: y >= 0 (Gạch bỏ y < 0)
					\fill[pattern=north east lines, pattern color=brown!80!red]
					(\xmin,\ymin) rectangle (\xmax,0);
					\fill[opacity=0.05, brown!80!red]
					(\xmin,\ymin) rectangle (\xmax,0);
					
					% 3. Miền KHÔNG thỏa: x + 2y <= 5 (Gạch bỏ x + 2y > 5 => y > 2.5 - 0.5x)
					\fill[pattern=north west lines, pattern color=blue!80]
					(\xmin, {2.5 - 0.5*\xmin}) -- (\xmax, {2.5 - 0.5*\xmax}) -- (\xmax, \ymax) -- (\xmin, \ymax) -- cycle;
					\fill[opacity=0.05, blue!80]
					(\xmin, {2.5 - 0.5*\xmin}) -- (\xmax, {2.5 - 0.5*\xmax}) -- (\xmax, \ymax) -- (\xmin, \ymax) -- cycle;
					
					% 4. Miền KHÔNG thỏa: x + y <= 3 (Gạch bỏ x + y > 3 => y > 3 - x)
					\fill[pattern=north east lines, pattern color=red!80]
					(\xmin, {3 - \xmin}) -- (\xmax, {3 - \xmax}) -- (\xmax, \ymax) -- (\xmin, \ymax) -- cycle;
					\fill[opacity=0.05, red!80]
					(\xmin, {3 - \xmin}) -- (\xmax, {3 - \xmax}) -- (\xmax, \ymax) -- (\xmin, \ymax) -- cycle;
					
					% === VẼ CÁC ĐƯỜNG THẲNG ===
					\draw[blue, thick] (\xmin, {2.5 - 0.5*\xmin}) -- (\xmax, {2.5 - 0.5*\xmax});
					\draw[red, thick] (\xmin, {3 - \xmin}) -- (\xmax, {3 - \xmax});
				\end{scope}
				
				% === TRỤC TỌA ĐỘ ===
				\draw[->, thick] (\xmin,0)--(\xmax,0) node[above left] {$x$};
				\draw[->, thick] (0,\ymin)--(0,\ymax) node[below right] {$y$};
				\node[below left] at (0,0) {$O$};
				
				% === CÁC ĐIỂM ĐÁNH DẤU TỌA ĐỘ ===
				% Trên trục Ox
				\draw[thick] (2, 0.05) -- (2, -0.05) node[below] {$2$};
				\draw[thick] (3, 0.05) -- (3, -0.05) node[below] {$3$};
				
				% Trên trục Oy
				\draw[thick] (0.05, 1) -- (-0.05, 1) node[left] {$1$};
				\draw[thick] (0.05, 2.5) -- (-0.05, 2.5) node[left] {$\dfrac{5}{2}$};
				
				% === VẼ ĐIỂM GIAO & DÁN NHÃN ===
				\filldraw[fill=yellow, draw=orange] (0, 2.5) circle (1.2pt) node[blue, above right] {$A$};
				\filldraw[fill=yellow, draw=orange] (1, 2) circle (1.2pt) node[blue, below left] {$B$};
				\filldraw[fill=yellow, draw=orange] (3, 0) circle (1.2pt) node[blue, above right] {$C$};
				
				% Khung bao viền (Bounding Box giống file xuất GeoGebra)
				\draw[thick, gray!80] (\xmin,\ymin) rectangle (\xmax,\ymax);
				
			\end{tikzpicture}
		\end{center}
		Giá trị biểu thức $T = 2x + 2{,}5y$ đạt giá trị lớn nhất tại một trong các đỉnh của tứ giác $OABC$.\\
		Ta có $T(O) = 0$, $T(A) = 2 \cdot 0 + 2{,}5 \cdot \dfrac{5}{2} = 6{,}25$, $T(B) = 2 \cdot 2 + 2{,}5 \cdot 1 = 6{,}5$, $T(C) = 2 \cdot 3 + 2{,}5 \cdot 0 = 6$.\\
		Suy ra $T$ đạt giá trị lớn nhất bằng $6{,}5$ khi $x = 2$, $y = 1$ ứng với tọa độ đỉnh $B$.\\
		Vậy cửa hàng lãi nhiều nhất là $6{,}5$ triệu đồng trong tháng khi kinh doanh hai loại gạo này.
	}
\end{ex}

\begin{ex}%[0D2V2-3]%[Dự án C đề thi thử TN Môn Toán NH25-26 - Dot 4 - Phạm Đăng Đăng]%[THPT Lê Quý Đôn-Đống Đa-Hà Nội]%[GVPB - Phan Quốc Trí]
	Một cơ sở rang xay cà phê để sản xuất và đóng gói cà phê  thành phẩm. Mỗi tháng cơ sở nhập về $300$ kg cà phê vối và $240$ kg cà phê chè để làm nguyên liệu. Hiện tại cơ sở sản xuất hai loại cà phê là loại $I$ và loại $II$.
	\begin{itemize}
		\item Loại $I$: Đóng gói $100$ g với tỉ lệ trộn $40$ g cà phê vối và $60$ g cà phê chè.
		\item Loại $II$: Đóng gói $100$ g với tỉ lệ $60$ g cà phê vối và $40$ g cà phê chè.
	\end{itemize}
	Biết rằng mỗi gói cà phê loại $I$ cơ sở lãi $10\,000$ đồng, mỗi gói cà phê loại $II$ cơ sở lãi $8\,000$ đồng. Nếu không phải lo khâu tiêu thụ thì cơ sở sản xuất thu được lợi nhuận mỗi tháng tối đa là bao nhiêu triệu đồng ( làm tròn đến hàng phần mười)?
	\shortans{$45{,}6$}
	\loigiai{Ta có $300$ kg = $300\,000$ g cà phê vối; $240$ kg = $240\,000$ g cà phê chè.\\
		Gọi $x$ là số gói loại $I$, $y$ là số gói loại $II$ ($x, y \geq 0$).\\
		Tổng cà phê vối: $40x + 60y \leq 300\,000 \Leftrightarrow 2x + 3y \leq 15\,000$.\\
		Tổng cà phê chè: $60x + 40y \leq 240\,000 \Leftrightarrow 3x + 2y \leq 12\,000$.\\
		Tổng tiền lãi của cơ sở $P = 10\,000x + 8\,000y$ (đơn vị: đồng).\\
		Bài toán trở thành tìm nghiệm của hệ bất phương trình 
		$\heva{ 
			&x \geq 0 \\ 
			&y \geq 0 \\ 
			&2x + 3y \leq 15\,000 \\ 
			&3x + 2y \leq 12\,000 
		}$ để biểu thức
		$P = 10\,000x + 8\,000y$ đạt giá trị lớn nhất.
		\begin{center}
			\begin{tikzpicture}[scale=0.7,font=\footnotesize,line join=round
				,line cap=round,>=stealth]
				\coordinate (A) at (0,5);
				\coordinate (B) at (1.2,4.2);
				\coordinate (C) at (4,0);
				\coordinate (O) at (0,0);
				\draw[fill=black] (A) circle (1pt) node[below,xshift=1.5mm,yshift=-4mm] {$A$};
				\draw[fill=black] (B) circle (1pt) node[below] {$B$};
				\draw[fill=black] (C) circle (1pt) node[above,xshift=-4mm]  {$C$};
				\draw[fill=black] (O) circle (1pt) node[ above right]  {$O$};
				\fill[pattern=crosshatch dots
				] 
				(-3,0)--(-3,-2)--(10,-2)--(10,0)--cycle
				(-3,0)--(-3,10)--(0,10)--(0,0)--cycle
				(0,10)--(A)--(B)--(C)--(10,0)--(10,10)--cycle;
					\draw[->] (-3,0)--(10,0) node[above right]{$x$};
				\draw[->] (0,-2)--(0,10) node[above right]{$y$};
				
				\draw[thick,blue,domain=9:-3,samples=2]
				plot (\x,{5-2/3*\x})
				node[below right,rotate=-33] {$2x+3y=15\,000$};

				\draw[thick,red,domain=5:-2.5,samples=2]
				plot (\x,{6-1.5*\x})
				node[above right,rotate=-55,] {$3x+2y=12\,000$};
				
			\end{tikzpicture}
		\end{center}
		Miền nghiệm của hệ bất phương trình là miền tứ giác $OABC$ với 
		$O(0; 0)$, $A(0; 5\,000)$, $C(4\,000; 0)$, $B(1\,200; 4\,200)$ (Miền không bị tô đậm).\\
		Tại $(0; 0): P = 0$\\
		Tại $(1\,200; 4\,200): P = 10\,000(1\,200) + 8\,000(4\,200)  = 45\,600\,000$ đồng.\\
		Tại $(0;5\,000): P = 8\,000(5\,000) = 40\,000\,000$ đồng.\\
		Tại $(4\,000; 0): P = 10\,000(4\,000) = 40\,000\,000$ đồng.\\
		Vậy lợi nhuận tối đa là $45{,}6$ triệu đồng/tháng.}
\end{ex}

\begin{ex}%[0D2V2-3]%[TN-THPT-NguyenTrungThien-HaTinh-NH25-26]%[Loc Do]%[PB: Phan Anh]
	Một cơ quan X chuẩn bị tổ chức cho $500$ người đi tham quan trải nghiệm. Để chuẩn bị cho chuyến đi, cơ quan cần vận chuyển tổng cộng $29$ tấn hàng hoá (bao gồm vật dụng, thực phẩm,...). Công ty vận tải báo giá cho thuê xe như sau: 
	\begin{itemize}
		\item Xe lớn: Có thể chở tối đa $50$ người và $2$ tấn hàng. Chi phí thuê là $10$ triệu đồng/xe. Công ty có $13$ xe loại này.
		\item Xe nhỏ: Có thể chở tối đa $30$ người và $3$ tấn hàng. Chi phí thuê là $7$ triệu đồng/xe. Công ty có $15$ xe loại này.
	\end{itemize}
	Sau khi tính toán, cơ quan X chọn phương án để chi phí thuê xe là thấp nhất. Số tiền thuê xe thấp nhất mà cơ quan X phải trả là bao nhiêu triệu đồng?
	\shortans{137}
	\loigiai{
		Giả sử, cơ quan X thuê số xe lớn và xe nhỏ lần lượt là $x$, $y$ (xe), ($x,y\in \mathbb{N}$).\\
		Khi đó, số tiền mà cơ quan phải trả là $F(x,y)=10x+7y$ (triệu đồng).\\
		Bất phương trình biểu thị số người là $50x+30y\ge 500\Leftrightarrow 5x+3y\ge 50$.\\
		Bất phương trình biểu thị khối lượng hàng hoá là  $2x+3y\ge 29$.\\
		Do đó, từ đề bài ta có hệ bất phương trình sau $\heva{& 0\le x\le 13\\ & 0\le y\le 15\\ & 5x+3y\ge 50\\ & 2x+3y\ge 29.}$ (*)\\
		Bài toán trở thành tìm giá trị nhỏ nhất của biểu thức $F(x,y)=10x+7y$ với $x,y$ thoả mãn (*).\\
		Biểu diễn miền nghiệm của hệ (*) bằng phần không bị tô đậm (kể cả bờ) được minh hoạ ở hình vẽ dưới đây:
		\begin{center}
			\begin{tikzpicture}[line join=round, line cap=round,>=stealth,thick,x=0.2cm,y=0.2cm]
				\tikzset{every node/.style={scale=0.9}}
				\begin{scope}
					\clip (-5,-5) rectangle (20,20);
					\fill[pattern=north east lines, pattern color=blue] (-5,-5) -- (0,-5) -- (0,20) -- (-5,20) -- cycle;
					\draw (0,-5)--(0,20)  ;
					\fill[pattern=north west lines, pattern color=green] (13,-5) -- (20,-5) -- (20,20) -- (13,20) -- cycle;
					\draw (13,-5)--(13,20)  ;
					\fill[pattern=vertical lines, pattern color=cyan] (-5,0) -- (-5,-5) -- (20,-5) -- (20,0) -- cycle;
					\draw (-5,0)--(20,0)  ;
					\fill[pattern=grid, pattern color=orange] (-5,15) -- (20,15) -- (20,20) -- (-5,20) -- cycle;
					\draw (-5,15)--(20,15)  ;
					\fill[pattern=dots, pattern color=magenta] (-5,-5) -- (13,-5) -- (-2,20) -- (-5,20) -- cycle;
					\draw (13,-5)--(-2,20)  ;
					\fill[pattern=crosshatch, pattern color=brown] (-5,13) -- (-5,-5) -- (20,-5) -- (20,-3.67) -- cycle;
					\draw (-5,13)--(20,-3.67)  ;
				\end{scope}
				\draw[->] (-5,0)--(22.2,0) node[below]{$x$};
				\draw[->] (0,-5)--(0,22.2) node[left]{$y$};
				\draw (0,0) node[below left]{$O$};
				\fill (1,15) circle (1.5pt) node[below right,scale=0.8] {$(1,15)$};
				\fill (13,15) circle (1.5pt) node[below left,scale=0.8] {$(13,15)$};
				\fill (13,1) circle (1.5pt) node[ right,scale=0.8] {$(13,1)$};
				\fill (7,5) circle (1.5pt) node[right,scale=0.8] {$(7,5)$};
				\foreach \x/\g in {5/-90,10/-90,15/-90}{
					\fill (\x,0) circle (1pt) +(\g:3mm) node[scale=0.8] {$\x$};
				}
				
				\foreach \y/\g in {5/180,10/180,15/160}{
					\fill (0,\y) circle (1pt) +(\g:3mm) node[scale=0.8] {$\y$};
				}
			\end{tikzpicture}
		\end{center}
		Trong đó, toạ độ các đỉnh được xác định như sau
		\begin{itemize}
			\item $A(1;15)$ là giao điểm của $2$ đường thẳng $y=15$ và $5x+3y=50$;
			\item $B(13;15)$ là giao điểm của $2$ đường thẳng $x=13$ và $y=15$;
			\item $C(7;5)$ là giao điểm của $2$ đường thẳng $5x+3y=50$ và $2x+3y=29$;
			\item $D(13; 1)$ là giao điểm của $2$ đường thẳng $x=13$ và $2x+3y=29$.
		\end{itemize}
		Kiểm tra giá trị của $F$ tại các đỉnh của tứ giác ta được
		\[F(1;15)=115, \,F(13;15)=235,\, F(7;5)=105,\, F(13;1)=137.\]
		Vậy số tiền ít nhất mà cơ quan X phải trả là $105$ triệu đồng (khi thuê $7$ xe lớn và $5$ xe nhỏ).
	}
\end{ex}

\begin{ex}%[0D2V2-3]%[Dự án C đợt 3 NH25-26- Phạm Đăng Đăng]%[TT-THPT-TranNhanTong-HaNoi-N25-26]
	Một xưởng mộc sản xuất hai loại sản phẩm: bàn gỗ và ghế gỗ. Xưởng có tối đa $240$ đơn vị gỗ và tổng thời gian nhân công là $100$ giờ. Để làm $1$ cái bàn cần $5$ đơn vị gỗ và $2$ giờ công. Để làm $1$ cái ghế cần $2$ đơn vị gỗ và $1$ giờ công. Mỗi cái bàn lãi $600\,000$ đồng, mỗi cái ghế lãi $300\,000$ đồng. Lợi nhuận của xưởng mộc lớn nhất là bao nhiêu triệu đồng?
	\shortans{$30$}
	\loigiai{
	\immini[thm]{	Gọi $x$, $y$ lần lượt là số bàn, ghế mà xưởng mộc sản xuất $(x,y>0)$.\\
		Theo đề bài, ta có hệ bất phương trình $$\heva{&5x+2y\le240\\
			&2x+y\le 100\\
			&x,y\ge0.}$$
	Miền nghiệm của hệ là miền tứ giác $OABC$ với $O(0;0)$, $A(0;100)$, $B(40;20)$ và $C(48;0)$.\\
			Lợi nhuận của xưởng thu được là \break $L(x;y)=600x+300y$ (đơn vị: nghìn đồng).\\
			Ta có $L(0;0)=0$, $L(0;100)=30\,000$,\\ $L(40;20)=30\,000$, $L(48;0)=28\,800$.\\Vậy lợi nhuận lớn nhất thu được là 30 triệu đồng.}{	\begin{tikzpicture}[scale=0.05, line cap=round, line join=round, font=\small,>=stealth]
				\clip (-50,-20) rectangle(100,200);
				% Trục tọa độ
				\draw[->] (-40,0)--(80,0) node[above right]{$x$};
				\draw[->] (0,-20)--(0,180) node[above right]{$y$};
				
				
				\draw[thick,blue,domain=110:-30,samples=2]
				plot (\x,{100-2*\x})
				node[below right,rotate=-65] {$2x+y=100$};
				
				% x + 4y = 200 <=> y = 50 - x/4
				\draw[thick,red,domain=220:-20,samples=2]
				plot (\x,{120-2.5*\x})
				node[above right,rotate=-65,] {$5x+2y=240$};
				
				
				
				% Các điểm A, B, C
				\coordinate (A) at (0,100);
				\coordinate (B) at (40,20);
				\coordinate (C) at (48,0);
				
				\draw[fill=black] (A) circle (20pt) node[below,xshift=1.5mm,yshift=-4mm] {$A$};
				\draw[fill=black] (B) circle (20pt) node[left] {$B$};
				\draw[fill=black] (C) circle (20pt) node[above,xshift=-4mm]  {$C$};
				\node[below left] at (0,0) {$O$};
				
				\node[below left] at (0,100) {$100$};
				\fill[pattern=horizontal lines] 
				(80,0)--(C)--(B)--(A)--(0,180)--(80,180)
				(80,0)--(80,-20)--(-40,-20)--(-40,0)--cycle
				(0,0)--(-40,0)--(-40,180)--(0,180)--cycle;
				%\fill[pattern=horizontal lines] plot[domain=-20:200/3]
				%(\x,{100-\x}) --
				%(220,-5) --(220,120)--cycle;
				%\fill[pattern=horizontal lines] plot[domain=90:20]
				%	(\x,{155-2*\x}) --
				%	(-20,120) --(-20,-20)--cycle;
		\end{tikzpicture}}
	}
\end{ex}

\begin{ex}%[Dự án đề TT Toán NH25-26-Dot 4- Hải Phụng]%[THCS-THPT LÊ THÁNH TÔNG - TP.HCM]%[GVPB - Viet Hoang Pham]%[0D3H2-6]
	Anh Trọng quyết định mua $2$ cây đào và $1$ cây mai tại một nhà vườn để trang trí nhà dịp Tết Nguyên Đán. Giá mỗi cây đào là $1\,000\,000$ đồng, giá mỗi cây mai là $2\,000\,000$ đồng. Để kích cầu, nhà vườn đưa ra chính sách: nếu khách hàng mua thêm $x$  cây cảnh nhỏ (với $x$  là số nguyên không âm), mỗi cây nhỏ có giá $50\,000$ đồng thì chi phí vận chuyển toàn bộ đơn hàng được tính theo hàm số $f(x)=x^2-100x+3\,000$  (đơn vị: nghìn đồng). Biết rằng số lượng cây mua thêm không vượt quá 40 cây để đảm bảo tải trọng xe $(0 \leq x \leq 40)$. Tổng chi phí (tiền mua cây và tiền vận chuyển) thấp nhất mà Anh Trọng phải bỏ ra để trang trí nhà dịp Tết là bao nhiêu nghìn đồng?
	\shortans{$6375$}
	\loigiai{
		Anh Trọng mua cố định:
		\begin{itemize}
			\item $2$ cây đào: \( 2 \cdot 1\,000\,000 = 2\,000\,000 \) đồng
			\item $1$ cây mai: \( 1 \cdot 2\,000\,000 = 2\,000\,000 \) đồng
		\end{itemize}
		Vậy tiền mua cây cố định là \( 4\,000\,000 \) đồng, tức là \( 4\,000 \) nghìn đồng.\\
		Tổng chi phí \( C(x) \) (đơn vị nghìn đồng) là
		\[
		C(x) = 4\,000 + 50x + (x^2 - 100x + 3\,000) = x^2 - 50x + 7\,000.
		\]
		Hàm \( C(x) = x^2 - 50x + 7\,000 \) là hàm bậc hai. Giá trị nhỏ nhất đạt tại
		\[
		x = -\frac{b}{2a} = \frac{50}{2 \cdot 1} = 25.
		\]
		Chi phí tại \( x = 25 \) là
		\[
		C(25) = 25^2 - 50 \cdot 25 + 7\,000 = 625 - 1250 + 7\,000 = 6\,375.
		\]
		Vậy chi phí nhỏ nhất là \( 6\,375 \) nghìn đồng khi mua thêm \( x = 25 \) cây cảnh nhỏ.
	}
\end{ex}

\begin{ex}%[0D8C1-3]%[Dự án đề thi thử đợt 3-Nguyễn Đức Lợi]%[SGD tỉnh Lạng Sơn]%[GVPB - ThanhTa]
	\immini[thm]{Một màn chơi của một trò chơi điện tử được thiết kế như sau: có năm vị trí $A$, $B$, $C$, $D$, $E$, được đặt ở năm đỉnh của một hình chóp tứ giác. Nhân vật sẽ được đặt ở một vị trí bất kỳ và có thể di chuyển tự do giữa các đỉnh với nhau, mỗi lần di chuyển đều phải từ đỉnh này di chuyển đến đỉnh khác. Giả sử khi bắt đầu, nhân vật được đặt ở vị trí $A$. Số cách di chuyển để sau sáu bước nhảy, nhân vật quay lại vị trí $A$ là bao nhiêu?
	}
	{
		\begin{tikzpicture}[font=\footnotesize, line join=round, line cap=round,>=stealth,scale=.8]
			
			\path (0,0) coordinate (B)
			(1,-1.5) coordinate (C)
			(3,-1.5) coordinate (D)
			(5,0) coordinate (E)
			(2,3) coordinate (A);
			\draw (A)--(B)--(C)--cycle (A)--(D)--(E)--cycle (C)--(D);
			\draw [dashed] (D)--(B)--(E)--(C);
			\foreach \t/\g in {A/90,B/180,C/-135,D/-45,E/0}{
				\fill (\t) circle (1pt) node[shift={(\g:10pt)}]{$ \t $};
			}
			\foreach \t/\g in {A,B,C,D,E}{
				\draw (\t) circle (5pt);}
		\end{tikzpicture}
	}
	\par
	\shortans{820}
	\loigiai{
		Từ một đỉnh bất kỳ, ta luôn có $4$ cách di chuyển đến $4$ đỉnh còn lại (các đỉnh $B$, $C$, $D$, $E$ có vai trò như nhau). Quá trình di chuyển gồm $6$ bước, bắt đầu từ $A$ và kết thúc tại $A$.\\
		Ta chia thành các trường hợp dựa vào số lần nhân vật quay trở lại $A$ ở các bước giữa (từ bước $1$ đến bước $5$).
		\begin{enumerate}[\it TH 1:]
			\item Trong $5$ bước ở giữa, không có lần nào quay lại $A$.\\
			Đường đi có dạng $A \to X_1 \to X_2 \to X_3 \to X_4 \to X_5 \to A$ (với $X_i \ne A$).
			\begin{itemize}
				\item Từ $A \to X_1$ có $4$ cách.
				\item Từ $X_1 \to X_2$:có $3$ cách (do $X_2 \ne A$ và $X_2 \ne X_1$).
				\item Tương tự, các bước $X_2 \to X_3$, $X_3 \to X_4$, $X_4 \to X_5$ mỗi bước đều có $3$ cách.
				\item Từ $X_5 \to A$ có $1$ cách.
			\end{itemize}
			Số cách di chuyển ở TH$_1$ là $4 \cdot 3 \cdot 3 \cdot 3 \cdot 3 \cdot 1=324$ (cách).
			\item Trong $5$ bước ở giữa, có đúng $1$ lần quay lại $A$.\\
			Do không thể đứng yên tại $A$ trong $2$ bước liên tiếp, lần quay lại $A$ không thể là bước $1$ và bước $5$, nên chỉ có thể rơi vào bước $2$, $3$ hoặc $4$.
			\begin{itemize}
				\item Về $A$ ở bước $2$ ($A \to X_1 \to A \to X_3 \to X_4 \to X_5 \to A$).\\
				Số cách là $4 \cdot 1 \cdot 4 \cdot 3 \cdot 3 \cdot 1=144$ (cách).
				\item Về $A$ ở bước $3$ ($A \to X_1 \to X_2 \to A \to X_4 \to X_5 \to A$).\\
				Số cách là $4 \cdot 3 \cdot 1 \cdot 4 \cdot 3 \cdot 1=144$ (cách).
				\item Về $A$ ở bước $4$ ($A \to X_1 \to X_2 \to X_3 \to A \to X_5 \to A$).\\
				Số cách là $4 \cdot 3 \cdot 3 \cdot 1 \cdot 4 \cdot 1=144$ (cách).
			\end{itemize}
			Số cách di chuyển ở TH$_2$ là $144 \cdot 3=432$ (cách).
			\item Trong $5$ bước ở giữa, có đúng $2$ lần quay lại $A$.\\
			Để $A$ không lặp lại kề nhau, $2$ lần quay lại $A$ bắt buộc phải ở bước $2$ và bước $4$.\\
			Đường đi có dạng $A \to X_1 \to A \to X_3 \to A \to X_5 \to A$.
			\begin{itemize}
				\item Bước $A \to X_1 \to A$ có $4 \cdot 1=4$ cách.
				\item Bước $A \to X_3 \to A$ có $4 \cdot 1=4$ cách.
				\item Bước $A \to X_5 \to A$ có $4 \cdot 1=4$ cách.
			\end{itemize}
			Số cách di chuyển ở TH$_3$ là $4 \cdot 4 \cdot 4=64$ (cách).
		\end{enumerate}
		(Không thể có $3$ lần quay lại $A$ vì sẽ dẫn đến $2$ bước liên tiếp đứng tại $A$).\\
		Vậy tổng số cách di chuyển thỏa mãn là $324 + 432 + 64 = 820$ (cách).
	}
\end{ex}

\begin{ex}%[0D8C2-5]%[Dự án C - Đợt 3 - Sở Giáo Dục Sơn La]%[Nguyễn Chiến - Nguyễn Tiến]
	\immini[thm]{
		Cho hình vẽ bên gồm $4$ tam giác. Người ta chọn $3$ số phân biệt từ tập hợp $S = \{1; 2; \ldots; 26\}$ để xếp vào $3$ tam giác ở $3$ góc. Sau đó, tính tổng bình phương của $3$ số đó rồi ghi kết quả vào tam giác còn lại ở giữa. Hỏi có bao nhiêu cách xếp sao cho số ghi ở tam giác giữa là một số chia hết cho $5$?
	}{
		\begin{tikzpicture}[scale=0.8, line join=round, line cap=round]
			\coordinate (A) at (0,0);
			\coordinate (B) at (3,0);
			\coordinate (C) at (1.5,{1.5*sqrt(3)});
			\coordinate (M) at ($(A)!0.5!(B)$);
			\coordinate (N) at ($(B)!0.5!(C)$);
			\coordinate (P) at ($(C)!0.5!(A)$);
			\draw[blue] (A) -- (B) -- (C) -- cycle;
			\draw[blue] (M) -- (N) -- (P) -- cycle;
		\end{tikzpicture}
	}
	\shortans[oly]{$3360$}
	\loigiai{
		Gọi $a$, $b$, $c$ là $3$ số phân biệt được chọn từ tập hợp $S$ để xếp vào $3$ góc.\\
		Theo bài ra, số ghi ở tam giác giữa là $a^2 + b^2 + c^2$.\\
		Ta cần tìm số cách chọn và xếp có thứ tự $3$ số $a$, $b$, $c$ sao cho tổng $a^2 + b^2 + c^2$ chia hết cho $5$.\\
		Với mọi số nguyên $x$, bình phương của nó chia cho $5$ chỉ có thể dư $0$, $1$ hoặc $4$.\\
		Gọi $R_a$, $R_b$, $R_c \in \{0; 1; 4\}$ lần lượt là số dư của $a^2$, $b^2$, $c^2$ khi chia cho $5$.\\
		Để $a^2 + b^2 + c^2$ chia hết cho $5$ thì tổng $R_a + R_b + R_c$ phải chia hết cho $5$.\\
		Ta chia tập $S = \{1; 2; \ldots; 26\}$ thành $3$ tập con theo số dư của bình phương khi chia cho $5$:
		\begin{itemize} 
			\item Tập $A_0$ gồm các số $x \in S$ thỏa mãn $x$ chia hết cho $5$.\\
			Suy ra $A_0 = \{5; 10; 15; 20; 25\}$ gồm $5$ phần tử.
			\item Tập $A_1$ gồm các số $x \in S$ thỏa mãn $x$ chia cho $5$ dư $1$ hoặc $x$ chia cho $5$ dư $4$.\\
			Suy ra $A_1 = \{1; 4; 6; 9; 11; 14; 16; 19; 21; 24; 26\}$ gồm $11$ phần tử.
			\item Tập $A_2$ gồm các số $x \in S$ thỏa mãn $x$ chia cho $5$ dư $2$ hoặc $x$ chia cho $5$ dư $3$.\\
			Suy ra $A_2 = \{2; 3; 7; 8; 12; 13; 17; 18; 22; 23\}$ gồm $10$ phần tử.
		\end{itemize}
		Trường hợp 1: Cả $3$ số $a$, $b$, $c$ đều có bình phương chia hết cho $5$.\\
		Ta chọn $3$ số phân biệt từ tập $A_0$ và xếp vào $3$ vị trí, số cách xếp là $\mathrm{A}_5^3 = 5 \cdot 4 \cdot 3 = 60$ (cách).\\
		Trường hợp 2: $3$ số được chọn có bình phương chia $5$ dư lần lượt là $0$, $1$ và $4$.\\
		Ta chọn $1$ số từ $A_0$, $1$ số từ $A_1$ và $1$ số từ $A_2$, sau đó hoán vị vào $3$ vị trí góc.\\
		Số cách thực hiện là $\mathrm{C}_5^1 \cdot \mathrm{C}_{11}^1 \cdot \mathrm{C}_{10}^1 \cdot 3! = 5 \cdot 11 \cdot 10 \cdot 6 = 3300$ (cách).\\
		Tổng số cách xếp thỏa mãn yêu cầu bài toán là $60 + 3300 = 3360$ (cách).
	}
\end{ex}

\begin{ex}%[0D8C2-7]%[Dự án C đề thi thử TN Môn Toán NH25-26 - Dot 4 - Phạm Hoàng Đăng]%[THPT Nguyễn Trãi-Hà Nội]%[GVPB - Khắc Thiên]
	Một kỹ sư tiến hành lắp ráp một rotor của động cơ phản lực. Rotor có $9$ khe cắm cánh quạt được đánh số cố định từ $1$ đến $9$ theo vòng tròn như hình vẽ. Khoảng cách giữa các khe đều nhau, tạo thành các đỉnh của một đa giác đều có $9$ cạnh. Do sai số chế tạo, $9$ cánh quạt có khối lượng thực tế là các số nguyên phân biệt từ $1$ gam đến $9$ gam. Để đảm bảo rotor cân bằng động học khi quay, kỹ sư lựa chọn phương án lắp đặt thỏa mãn đồng thời các điều kiện sau:
	\begin{itemize}
		\item Chia $9$ cánh quạt thành $3$ nhóm (mỗi nhóm $3$ cánh).
		\item Mỗi nhóm được lắp vào $3$ khe cắm tạo thành một tam giác đều (ba khe cắm tạo thành một tam giác đều khi và chỉ khi chúng cách nhau đúng $3$ khe theo vòng tròn).
		\item Tổng khối lượng của $3$ cánh quạt trong mỗi nhóm phải bằng nhau.
	\end{itemize}
	\begin{center}
		\begin{minipage}{0.45\textwidth}
			\begin{tikzpicture}[scale=1,>=stealth, line join=round, line cap=round, declare function={R=3; h=0.4; ecc=0.6;}]
				\path (0,0) coordinate (O);
				\fill[gray!30] (-R,0) arc (180:360:R and R*ecc) -- (R,-h) arc (0:-180:R and R*ecc) -- cycle;
				\filldraw[fill=white] (O) ellipse (R and R*ecc);
				\draw (-R,0) -- (-R,-h) arc (180:360:R and R*ecc) -- (R,0);
				\draw[fill=gray!10] (O) ellipse (0.5 and 0.5*ecc);
				\foreach \g in {0,40,...,320}{
					\filldraw[fill=gray!5] (\g:2.2 and 2.2*ecc) ellipse (0.35 and 0.35*ecc);
				}
			\end{tikzpicture}
		\end{minipage}
		\begin{minipage}{0.45\textwidth}
			\begin{tikzpicture}[scale=1,>=stealth, line join=round, line cap=round, smooth, samples=450]
				\def\R{3}
				\def\n{9}
				\def\ang{40}
				\def\xmin{-3.5} \def\xmax{3.5} \def\ymin{-3.5} \def\ymax{4}
				\path (0,0) coordinate (O)
				\foreach \i in {1,...,\n}{
					({90-(\i-1)*\ang}:\R) coordinate (P\i)
				};
				\draw (P1) \foreach \i in {2,...,\n}{-- (P\i)} -- cycle;
				\draw(O) circle (\R)
				(P1) \foreach \i in {2,...,\n}{-- (P\i)} -- cycle;
				\foreach \i/\g in {1/90,2/10,3/0,4/-10,5/-30,6/-120,7/-150,8/-180,9/140}
				\draw[fill=black] (P\i) circle (1pt) +(\g:.4) node{$\i$};
			\end{tikzpicture}
		\end{minipage}
	\end{center}
	Hỏi có bao nhiêu cách sắp xếp $9$ cánh quạt vào $9$ khe cắm thỏa mãn các điều kiện kỹ thuật trên (hai cách sắp xếp được coi là khác nhau nếu có ít nhất một cánh quạt ở một vị trí khe cắm khác nhau, không đồng nhất các cách lắp khác nhau bởi phép quay hay phép đối xứng của rotor)?
	\par
	\shortans[oly]{2\,592}
	\loigiai{
		Tổng khối lượng của các cánh quạt là $1+2+\ldots+9=45$ (gam).\\
		Ta chia $3$ nhóm cánh quạt sao cho tổng khối lượng của $3$ cánh quạt trong mỗi nhóm phải bằng nhau.\\ Suy ra, khối lượng của mỗi nhóm bằng $\dfrac{45}{3}=15$ (gam).\\
		Các bộ gồm $3$ số lấy từ tập $\{1;2;\ldots;9\}$ mà tổng của chúng bằng $15$ là\\
		\[\{1;5;9\},\{1;6;8\},\{2;4;9\},\{2;5;8\},\{2;6;7\},\{3;4;8\},\{3;5;7\},\{4;5;6\}.\]
		Chọn $3$ bộ trong số $8$ bộ trên sao cho các chữ số không trùng nhau để có được $3$ nhóm cánh quạt có tổng khối lượng thỏa mãn.
		\begin{itemize}
			\item \textbf{Trường hợp 1:} Ba bộ đó là $\{1;5;9\}$, $\{2;6;7\}$, $\{3;4;8\}$.\\
			Chọn ra $3$ tam giác đều thỏa đề bài có $3!$ (cách).\\
			Với mỗi tam giác đều, để lắp $3$ số vào $3$ đỉnh ta có $3!$ (cách).\\
			Vậy có tất cả $3$ tam giác đều.\\ 
			Do đó số cách xếp trong trường hợp này là $3!\cdot(3!)^3=1\,296$ (cách).
			\item \textbf{Trường hợp 2:} Ba bộ là $\{1;6;8\}$, $\{2;4;9\}$, $\{3;5;7\}$.\\
			Tương tự trường hợp $1$, có $1\,296$ (cách).
		\end{itemize}
		Vậy có tất cả $1\,296+1\,296=2\,592$ (cách).
	}
\end{ex}

\begin{ex}%[Đề thi thử môn Toán TRƯỜNG THPT Cổ Loa - Hà Nội năm học 2025-2026]%[Phạm Quốc Toàn - Phạm Hoàng Việt]%[0D8V2-5]
	Cho hình vuông $ABCD$ có tâm $O$. Gọi $M$, $N$, $P$, $Q$ lần lượt là trung điểm của các cạnh $AB$, $BC$, $CD$, $DA$. Kí hiệu $H_1$, $H_2$, $H_3$, $H_4$ lần lượt là các hình vuông $AMOQ$, $BMON$, $CPON$, $DPOQ$. Bạn An muốn tô màu mỗi hình vuông $H_1$, $H_2$, $H_3$, $H_4$ bởi một màu trong $5$ màu (đỏ, xanh, vàng, tím, nâu) sao cho hai hình vuông có chung cạnh sẽ được tô bởi hai màu khác nhau. Hỏi bạn An có bao nhiêu cách tô màu?
	\begin{center}
		\begin{tikzpicture}[scale=0.8, font=\footnotesize, line join=round, line cap=round, >=stealth]
			% Tọa độ các điểm
			\coordinate (A) at (0,4);
			\coordinate (B) at (4,4);
			\coordinate (C) at (4,0);
			\coordinate (D) at (0,0);
			%%%%%%%%%%%%
			\coordinate (M) at (2,4);
			\coordinate (N) at (4,2);
			\coordinate (P) at (2,0);
			\coordinate (Q) at (0,2);
			\coordinate (O) at (2,2);
			% Tô nền
			\fill[gray!20] (A) rectangle (C);
			% Khung ngoài
			\draw[thick] (A) -- (B) -- (C) -- (D) -- cycle;
			% Các đường chia
			\draw[thick] (M) -- (P);
			\draw[thick] (Q) -- (N);
			% Vẽ các điểm
			\foreach \p in {A,B,C,D,M,N,P,Q,O} \fill (\p) circle (2pt);
			% Nhãn các điểm
			\node[above left] at (A) {$A$};
			\node[above right] at (B) {$B$};
			\node[below right] at (C) {$C$};
			\node[below left] at (D) {$D$};
			%%%%%%%%%%%%%%%%%
			\node[above] at (M) {$M$};
			\node[right] at (N) {$N$};
			\node[below] at (P) {$P$};
			\node[left] at (Q) {$Q$};
			\node[below right] at (O) {$O$};
			% Nhãn các miền
			\node at (1,3) {$H_1$};
			\node at (3,3) {$H_2$};
			\node at (3,1) {$H_3$};
			\node at (1,1) {$H_4$};
		\end{tikzpicture}
	\end{center}
	\par\shortans[0]{260}
	\loigiai{
		\begin{itemize}
			\item 	\textbf{Trường hợp $1$:} Hình $H_1$, $H_3$ được tô cùng màu. 
			\begin{itemize}
				\item Có $5$ cách chọn màu cho hình $H_1$, $H_3$.
				\item Có $4$ cách chọn màu cho hình $H_2$.
				\item Có $4$ cách chọn màu cho hình $H_4$. \\ 
				Suy ra trường hợp $1$ có $5 \times 4 \times 4=80$ cách.
			\end{itemize}
		\item \textbf{Trường hợp $2$:} Hình $H_1$, $H_3$ được tô khác màu. 
		\begin{itemize}
			\item Có $5$ cách chọn màu cho hình $H_1$.
			\item Có $4$ cách chọn màu cho hình $H_3$.
			\item Có $3$ cách chọn màu cho hình $H_2$.
			\item Có $3$ cách chọn màu cho hình $H_4$.\\
			Suy ra trường hợp $2$ có $5 \times 4 \times 3 \times 3=180$ cách.
		\end{itemize}
	\end{itemize}
		Vậy tất cả có $80+180=260$ cách.
	}
\end{ex}

\begin{ex}%[Đề thi thử môn Toán TRƯỜNG THPT Cổ Loa - Hà Nội năm học 2025-2026]%[Phạm Quốc Toàn - Phạm Hoàng Việt]%[0D8V2-7]
	Sắp xếp ngẫu nhiên ba bạn nam $A$, $B$, $C$ và $5$ bạn nữ vào $8$ cái ghế được xếp theo hàng ngang (mỗi bạn ngồi một cái ghế). Gọi xác suất để $3$ bạn $A$, $B$, $C$ ngồi theo thứ tự đó từ trái qua phải, đồng thời giữa $A$ và $B$ có ít nhất một bạn nữ, giữa $B$ và $C$ có nhiều nhất một bạn nữ là $p$. Giá trị của $\dfrac{5}{p}$ bằng bao nhiêu?
	\par\shortans[0]{67{,}2}
	\loigiai{
		Số phần tử của không gian mẫu là $n(\Omega)=8!=40\,320$.
		\begin{itemize}
			\item \textbf{Trường hợp $1$:} Giữa $B$, $C$ không có bạn nữ nào.\\
			Ta gọi $x_1$, $x_2$, $x_3$ lần lượt là số bạn nữ xếp như sau 
			\[
			x_1 \quad A \quad x_2 \quad B\quad C \quad x_3\,\, \text{với}\,\, x_1\ge 0,\, x_2\ge 1,\, x_3\ge 0.
			\]
			Ta có $x_1+x_2+x_3=5 \Leftrightarrow x'_1+x_2+x'_3=7$ với $x'_1\ge 1$, $x_2\ge 1$, $x'_3 \ge 1$. \\
			Do vậy, ta có $\mathrm{C}_6^2 \cdot 5!$ cách.
			\item \textbf{Trường hợp $2$:} Giữa $B$, $C$ có một bạn nữ.\\
			Ta gọi $x_1$, $x_2$, $x_3$, $x_4$ lần lượt là số bạn nữ xếp như sau 
			\[
			x_1 \quad A \quad x_2 \quad B \quad x_3 \quad C \quad x_4\,\, \text{với}\,\, x_1\ge 0,\, x_2\ge 1,\, x_3=1,\, x_4\ge 0. 
			\]
			Ta có $x_1+x_2+1+x_4=5 \Leftrightarrow x_1+x_2+x_4=4\Leftrightarrow x'_1+x_2+x'_4=6$, với $x'_1\ge 1$, $x_2\ge 1$, $x'_4\ge 1$. \\
			Do vậy sẽ có $\mathrm{C}_5^2\cdot 5!$ cách.
		\end{itemize}
		Vậy có $\mathrm{C}_6^2 \cdot 5!+\mathrm{C}_5^2\cdot 5!=3\, 000$ (cách). \\ 
		Vậy xác suất là $p=\dfrac{3\,000}{8!}$. \\
		Do vậy $\dfrac{5}{p}=\dfrac{336}{5}=67{,}2$. 
	}
\end{ex}

\begin{ex}%[0H5V3-5]%[Dự án C-TT THPT NH25-26-Đợt 4- ĐỖ CHÍ TÂM]%[THPT Phan Bội Châu-Khánh Hòa]%[GVPB: Phạm Văn Long]
	Cho hình hộp $ABCD.A'B'C'D'$. Gọi $N$ là điểm trên đoạn $C'D$ sao cho $DN=\dfrac{1}{3}D{C}'$. Khi phân tích $\overrightarrow{BN}=x\cdot \overrightarrow{BA}+y\cdot\overrightarrow{BC}+z\cdot \overrightarrow{B{B}'}$ thì giá trị $3x+2y+3z$ bằng bao nhiêu?
	\shortans{$5$}
	\loigiai{
		\begin{center}
			\begin{tikzpicture}[scale=0.7]
				%%\draw[gray!20] (-3,-2) grid (6,5);
				\path (0,0) coordinate (A)  
				(-2,-2) coordinate (B)
				(5,0) coordinate (D)
				(3,-2) coordinate (C)
				(0,4) coordinate (A')
				(-2,2) coordinate (B')
				(3,2) coordinate (C')
				(5, 4) coordinate (D')
				($(D)!0.33!(C')$) coordinate(N);
				\draw[dashed] (A')--(A)--(B) (A)--(D) (B)--(N);
				\draw (B)--(C)--(D)--(D')--(C')--(B')--(A') (B')--(B) (C')--(C) (A')--(D') (C')--(D) (B)--(C');
				\foreach \p/\q in {A/180, B/-90, C/-90, D/0, A'/90, B'/90, C'/90, D'/90, N/45}
				\fill[blue] (\p) circle(2pt) node[shift={(\q:3mm)}]{\bfseries $\p$};
			\end{tikzpicture}
		\end{center}
		Ta có $\overrightarrow{BN}=\overrightarrow{B{C}'}+\overrightarrow{C'N}=\left(\overrightarrow{B{B}'}+\overrightarrow{BC} \right)+\dfrac{2}{3}\overrightarrow{C'D}$ $=\left(\overrightarrow{B{B}'}+\overrightarrow{BC} \right)+\dfrac{2}{3}\left(\overrightarrow{C'D'}+\overrightarrow{C'C} \right)$\\
		\hspace*{4.5cm}$=\left(\overrightarrow{B{B}'}+\overrightarrow{BC} \right)+\dfrac{2}{3}\left(\overrightarrow{BA}+\left(-\overrightarrow{B{B}'} \right) \right)$ $=\dfrac{2}{3}\overrightarrow{BA}+\overrightarrow{BC}+\dfrac{1}{3}\overrightarrow{B{B}'}$.\\
		Suy ra $x=\dfrac{2}{3}$; $y=1$; $z=\dfrac{1}{3}$.\\
		Vậy $3x+2y+3z=3\cdot \dfrac{2}{3}+2\cdot 1+3\cdot \dfrac{1}{3}=5$.
	}
\end{ex}

\begin{ex}%[0H9V4-7]%[KNTT - Lớp 12 - Dự án C - Đợt 2]%[Loc do]%[PB: Phan Anh]
	\immini[thm]{Miền Trung Việt Nam vốn luôn phải oằn mình chống chọi với thiên tai. Giả sử trong một đợt áp thấp nhiệt đới mạnh lên thành bão, Trung tâm Dự báo Khí tượng Thủy văn Quốc gia phát đi thông báo khẩn về cơn bão số $4$ có tên quốc tế là Sao La. Trên bản đồ quy hoạch phòng chống thiên tai được gắn hệ trục tọa độ Oxy với đơn vị đo là $10$ km (hướng Đông là trục $Ox$, hướng Bắc là trục $Oy$), vị trí ba thành phố trọng điểm là Hà Tĩnh, Vinh (Nghệ An) và Thanh Hóa được xác định lần lượt tại các điểm $T(0;0)$, $N(-2;3)$ và $H(-1;5)$. Tại thời điểm bản tin phát đi, tâm bão Sao La đang ở ngoài khơi Biển Đông, cách thành phố Hà Tĩnh $200$ km. Vị trí tâm bão nằm ở hướng Đông Nam so với Hà Tĩnh, tạo với hướng Đông một góc $\alpha $ sao cho $\cos \alpha =0{,}8$. }{\begin{tikzpicture}[>=stealth,line join=round,line cap=round,font=\footnotesize,scale=0.8]
			\draw[->] (-1,0)--(6,0) node[below]{$x$};
			\draw [->] (0,-4)--(0,4) node [right]{$y$};
			\path (0,0) coordinate (T)
			(5,-3) coordinate (S) (3.5,-1.2) coordinate (S1) 
			(3,0) coordinate (B)  (2,-3) coordinate (C) ;
			\fill (3.5,-1.2) circle (1pt) node[right] {$S$};
			\fill (-0.2,3) circle (1pt) node[left] {$H$};
			\fill (-0.5,2) circle (1pt) node[left] {$N$};
			\draw[dotted,green,thick] (S) circle (0.8cm);
			\draw[thick,red] (S1) circle (1cm);
			\coordinate (A) at ($(S1)+(180:1)$);
			\path pic[angle radius=5mm,fill=green,draw=blue,"$\alpha$",angle eccentricity=1.5] {angle = S--T--B};
			\path pic[angle radius=5mm,fill=green,draw=blue,"$60^\circ$",angle eccentricity=1.5] {angle = S1--S--C};
			\draw (T)--(S)--(S1)--(A) (S)--(C) ;
			\node at (3,-1) {$r(t)$};
			\node at (5.6,0.2) {Đông};
			\node at (-0.5,4) {Bắc};
			\foreach \t/\g in {T/-120,S/0}{
				\fill (\t) circle (1pt) node[shift={(\g:7pt)},font=\scriptsize]{$ \t $};
			}
			
	\end{tikzpicture}}
	Cơn bão di chuyển phức tạp theo hướng Tây Bắc lệch $60^\circ$ so với hướng Tây với vận tốc $20$ km/h. Sức tàn phá của bão rất lớn với vùng nguy hiểm là một hình tròn có bán kính ban đầu $20$ km và liên tục mở rộng thêm $10$ km mỗi giờ. Để lên phương án sơ tán dân cư đồng bộ, Ban chỉ đạo cần biết khoảng thời gian mà cả $3$ thành phố này cùng nằm trong vùng nguy hiểm của bão Sao La kéo dài bao nhiêu giờ? (làm tròn kết quả đến hàng phần chục).
	\par\shortans[0]{$10,8$}
	\loigiai{
		Ta có $T(0;0)$, $N(-2;3)$ và $H(-1;5)$.\\
		Tại thời điểm bản tin phát đi, tâm bão Sao La có tọa độ là $(x;y)$.\\
		Suy ra $\heva{&x=20\cdot \cos\alpha\\& y=-20\cdot \sin\alpha}\Rightarrow S(16;-12)$.\\
		Bão di chuyển theo hướng Tây Bắc lệch $60^\circ$ so với hướng Tây nên vectơ chỉ phương có tọa độ 
		\[\left(\cos120^\circ; \sin120^\circ\right)=\left(-\dfrac{1}{2}; -\dfrac{\sqrt{3}}{2}\right).\]
		Vì vận tốc $20$ km/h nên vectơ vận tốc \[\overrightarrow{v}=2\left(-\dfrac{1}{2}; -\dfrac{\sqrt{3}}{2}\right)=\left(-1;\sqrt{3}\right).\]
		Tọa độ tâm bão sau $t$ (giờ) là $S\left(16-t;\; -12+\sqrt{3} t\right)$.\\
		Bán kính vùng nguy hiểm sau $t$ (giờ) là $r=2+t$. Suy ra 
		\begin{itemize}
			\item $SH^2 =\left(17-t\right)^2 +\left(\sqrt{3} t-17\right)^2 $
			\item $ST^2 =\left(16-t\right)^2 +\left(\sqrt{3} t-12\right)^2 $
			\item $SN^2 =\left(18-t\right)^2 +\left(\sqrt{3} t-15\right)^2$.
		\end{itemize} 
		Thành phố Thanh Hóa, Hà Tĩnh, Vinh nằm trong vùng nguy hiểm khi khoảng cách từ tâm bão đến thành phố nhỏ hơn hoặc bằng bán kính vùng nguy hiểm do vậy ta được
		\begin{eqnarray*}
			\heva{& SH^2 \le r^2\\ & ST^2 \le r^2\\ & SN^2 \le r^2} &\Leftrightarrow& \heva{& 3t^2 -\left(38+34\sqrt{3}\right)t+574\le 0\\ & t^2 -\left(12+8\sqrt{3}\right)t+132\le 0\\ & 3t^2 -\left(40+30\sqrt{3}\right)t+545\le 0}\\
			&\Leftrightarrow& t\in \left[\dfrac{40+30\sqrt{3} -\sqrt{2400\sqrt{3} -2240}}{6} ;\; \dfrac{12+8\sqrt{3} +\sqrt{192\sqrt{3} -192}}{2} \right]
		\end{eqnarray*}
		Khoảng thời gian kéo dài là 
		\[\dfrac{12+8\sqrt{3} +\sqrt{192\sqrt{3} -192}}{2} -\dfrac{40+30\sqrt{3} -\sqrt{2400\sqrt{3} -2240}}{6} \approx 10{,}8\,\text{giờ}.\]
	}
\end{ex}

\begin{ex}%[1C2V3-1]%[Dự án C đề thi thử TN Môn Toán NH25-26 - Dot 4 - Hector Tran]%[Liên trường THPT cụm 09 - Đợt 2 - Hà Nội]%[GVPB - Thanh Phát]
	\immini [thm]{
		Hằng năm, trường THPT M tổ chức các đoàn đi thăm hỏi và chúc Tết tứ thân phụ mẫu cao tuổi của cán bộ, giáo viên và nhân viên trong trường. Đoàn số $1$ có nhiệm vụ đến thăm $4$ gia đình tại các địa điểm $A$, $B$, $C$, $D$. Đoàn xuất phát từ trường THPT M, đến thăm đủ $4$ gia đình (mỗi gia đình đúng một lần). Sau khi thăm xong gia đình cuối cùng, đoàn kết thúc hành trình tại đó (không quay về trường). Do điều kiện giao thông dịp cuối năm, đoạn đường $AB$ chỉ có thể đi theo chiều từ $A$ đến $B$, không đi được theo chiều ngược lại (\textit{hình vẽ}). Đoàn số $1$ đã chọn được lộ trình di chuyển có tổng quãng đường di chuyển ngắn nhất. Tổng quãng đường ngắn nhất đó dài bao nhiêu km?
	}{
		\begin{tikzpicture}[scale=0.5, line join=round, line cap=round, every node/.style={scale=0.8}, declare function={xM=-4; yM=1; xC=0; yC=4.5; xB=4.5; yB=0; xA=2.5; yA=-3.5; xD=-2; yD=-3; R=0.5;}]		
			\path
			(xM, yM) coordinate (M)
			(xC, yC) coordinate (C)
			(xB, yB) coordinate (B)
			(xA, yA) coordinate (A)
			(xD, yD) coordinate (D);		
			\draw[-stealth, shorten >=0.2cm, shorten <=0.2cm] (M) -- (C) node[pos=0.5, sloped, above] {$4$ km};
			\draw[stealth-stealth, shorten >=0.2cm, shorten <=0.2cm] (C) -- (B) node[pos=0.5, sloped, above] {$3$ km};
			\draw[-stealth, shorten >=0.2cm, shorten <=0.2cm] (A) -- (B) node[pos=0.5, sloped, below] {$3$ km};
			\draw[stealth-stealth, shorten >=0.2cm, shorten <=0.2cm] (D) -- (A) node[pos=0.5, sloped, below] {$5$ km};
			\draw[-stealth, shorten >=0.2cm, shorten <=0.2cm] (M) -- (D) node[pos=0.5, sloped, below] {$7$ km};		
			\draw[-stealth, shorten >=0.2cm, shorten <=0.2cm] (M) -- (B) node[pos=0.5, sloped, above] {$6$ km};
			\draw[-stealth, shorten >=0.2cm, shorten <=0.2cm] (M) -- (A) node[pos=0.3, sloped, above] {$8$ km};
			\draw[stealth-stealth, shorten >=0.2cm, shorten <=0.2cm] (D) -- (C) node[pos=0.7, sloped, above] {$5$ km};
			\draw[stealth-stealth, shorten >=0.2cm, shorten <=0.2cm] (D) -- (B) node[pos=0.7, sloped, above] {$4$ km};
			\draw[stealth-stealth, shorten >=0.2cm, shorten <=0.2cm] (C) -- (A) node[pos=0.25, sloped, above] {$6$ km};		
			\foreach \p/\pos in {M/135, C/90, B/0, A/-45, D/-135} {
				\draw (\p) circle (0.6pt);
				\draw (\p) circle (0.5cm) node[shift={(\pos:20pt)}] {$\p$};
			}
		\end{tikzpicture}
	}
	\par\shortans{16}
	\loigiai{
		- Gán nhãn cho điểm xuất phát $\ell(M)=0$ và coi đây là nhãn vĩnh viễn.\\ 		
		Các đỉnh $A$, $B$, $C$, $D$ lúc này có nhãn tạm thời dựa trên khoảng cách từ $M$ là $\ell(A)=8$, $\ell(B)=6$, $\ell(C)=4$, $\ell(D)=7$.\\
		- So sánh các nhãn tạm thời, ta thấy $\ell(C)=4$ là nhỏ nhất.\\ 
		Ta chọn $C$ làm điểm tiếp theo của hành trình.\\
		- Từ $C$, ta xem xét các gia đình chưa thăm ($A$, $B$, $D$).
		\begin{itemize}
			\item Khoảng cách $M$ đến $A$ là $\ell(C)+AC=4+6=10$ suy ra $\ell(A)=10$.
			\item Khoảng cách $M$ đến $B$ là $\ell(C)+CB=4+3=7$ suy ra $\ell(B)=7$.
			\item Khoảng cách $M$ đến $D$ là $\ell(C)+CD=4+5=9$ suy ra $\ell(D)=9$.
		\end{itemize}
		Hiện tại nhãn nhỏ nhất từ các lựa chọn là $\ell(B)=7$.\\ 		
		Ta chọn $B$ là điểm tiếp theo.\\
		- Từ $B$, ta còn $A$ và $D$.\\ 		
		Do tuyến đường từ $B$ đến $A$ bị chặn, nên đi đến $D$ trước, rồi tới $A$.\\ 
		Vậy lộ trình di chuyển có tổng quãng đường di chuyển ngắn nhất là 		
		\[MC+CB+BD+DA=\ell(B)+BD+DA=7+4+5=16\text { (km).}\]
	}
\end{ex}

\begin{ex}%[1D2C2-5]%[Dự án C đề thi thử TN Môn Toán NH25-26 - Dot 4 - Hector Tran]%[Liên trường THPT cụm 09 - Đợt 2 - Hà Nội]%[GVPB - Thanh Phát]
	Một hộp chứa $100$ cái thẻ được đánh số thứ tự liên tiếp từ $1$ đến $100$. Hai thẻ khác nhau thì đánh số thứ tự khác nhau. Chọn ngẫu nhiên $3$ thẻ trong hộp. Có bao nhiêu cách chọn $3$ thẻ có số thứ tự lập thành cấp số cộng, đồng thời có tổng không vượt quá $150$?
	\par\shortans{1225}
	\loigiai{
		Giả sử các số trên $3$ thẻ lập thành cấp số cộng với số hạng đầu là $a$, công sai $d$, với $a$, $d \in \mathbb{N}^*$.\\
		Tổng các số trên $3$ thẻ không vượt quá $150$ nên ta có
		\begin{eqnarray*}
			&&a + (a + d) + (a + 2d) \le 150 \\
			&\Leftrightarrow& 3a + 3d \le 150 \\
			&\Leftrightarrow& a + d \le 50.
		\end{eqnarray*}
		Vì $a$, $d \in \mathbb{N}^*$ nên
		\begin{itemize}
			\item Với $d = 1$ thì $a \le 49$. Với mỗi số tự nhiên $a \in \mathbb{N}^*$ ta có $1$ cấp số cộng. Do đó trường hợp này ta thu được $49$ cấp số cộng thỏa mãn.
			\item Với $d = 2$ thì $a \le 48$. Với mỗi số tự nhiên $a \in \mathbb{N}^*$ ta có $1$ cấp số cộng. Do đó trường hợp này ta thu được $48$ cấp số cộng thỏa mãn.
			\item \ldots
			\item Cứ tiếp tục như vậy với $d = 49$ thì $a \le 1$, trường hợp này có $1$ cấp số cộng thỏa mãn.
		\end{itemize}
		Do đó số cách chọn thỏa yêu cầu đề là $49 + 48 + \cdots + 1 = \dfrac{49 \cdot 50}{2} = 1\,225$.
	}
\end{ex}

\begin{ex}%[Dự án đề TT Toán NH25-26-Dot 4- Hải Phụng]%[THCS-THPT LÊ THÁNH TÔNG - TP.HCM]%[GVPB - Viet Hoang Pham]%[1D2V2-7]
	Trong dịp Tết Nguyên Đán, bạn Kiên nhận được tiền lì xì ngày mùng $1$ là $3\,000\,000$ đồng. Kể từ mùng $2$, mỗi ngày tiền lì xì giảm $300\,000$ đồng so với ngày trước đó. Nhưng đến ngày mùng $3$, mẹ bạn Kiên lại nói: \lq\lq Để mẹ giữ tiền cho sau này cưới vợ nhé\rq\rq. Nên bạn Kiên quyết định từ ngày hôm đó trở đi (từ mùng $3$ đến mùng $9$) mỗi ngày sẽ đưa cho mẹ $X$ triệu đồng. Biết rằng đến hết ngày mùng $9$, Kiên còn lại đúng $5\,000\,000$ đồng. Tính giá trị của $X$.
	\shortans{$1,6$}
	\loigiai{
		Gọi $u_n$ là số tiền lì xì Kiên nhận được vào ngày mùng $n$. \\
		Theo giả thiết, dãy số $(u_n)$ là một cấp số cộng có số hạng đầu $u_1 = 3\,000\,000$ và công sai $d = -300\,000$.\\
		Tổng số tiền lì xì Kiên nhận được từ mùng $1$ đến hết mùng $9$ là tổng của $9$ số hạng đầu tiên của cấp số cộng:
		$$S_9 = \dfrac{9}{2} \left[ 2u_1 + (9-1)d \right] = \dfrac{9}{2} \left[ 2 \cdot 3\,000\,000 + 8 \cdot (-300\,000) \right] = 16~200\,000 \text{ (đồng).}$$
		Từ ngày mùng $3$ đến ngày mùng $9$ có tất cả $9 - 3 + 1 = 7$ ngày.\\
		Mỗi ngày Kiên đưa cho mẹ $X$ triệu đồng (tức là $1\,000\,000X$ đồng).\\
		Tổng số tiền Kiên đưa cho mẹ là: $7 \cdot 1\,000\,000X = 7\,000\,000X$ (đồng).\\
		Số tiền còn lại của Kiên sau ngày mùng $9$ là
		{\allowdisplaybreaks
			\begin{eqnarray*}
				& & 16~200\,000 - 7\,000\,000X = 5\,000\,000  \\
				&\Leftrightarrow& 7\,000\,000X = 11~200\,000\\
				&\Leftrightarrow& X = \dfrac{11~200\,000}{7\,000\,000} = 1{,}6.
			\end{eqnarray*}
		}
	}
\end{ex}

\begin{ex}%[1D2V3-8]%[Dự án C đợt 3 NH25-26- Phạm Đăng Đăng]%[TT-THPT-TranNhanTong-HaNoi-N25-26]%[GVPB - Phan Quốc Trí]
	Bác An dự định gửi vào ngân hàng một số tiền với lãi suất không đổi là $7{,}5\%$ một năm. Biết rằng cứ sau mỗi năm số tiền lãi sẽ được nhập vào vốn ban đầu. Bác An cần gửi ít nhất vào ngân hàng bao nhiêu triệu đồng (kết quả làm tròn đến hàng đơn vị) để sau $5$ năm Bác rút được số tiền đủ mua $1$ chiếc xe máy điện trị giá $57$ triệu đồng.
	\shortans{$40$}
	\loigiai{Gọi $u_n$ là số tiền Bác An nhận được sau $n-1$ năm.\\ Khi đó $(u_n)$ là một cấp số nhân với công bội $q=1,075$.\\
		Ta có $u_6=u_1\cdot q^5\Leftrightarrow u_1=\dfrac{57}{1,075^5}\approx 40$ triệu đồng.
	}
\end{ex}

\begin{ex}%[1D5H1-3]%[Dự án C-TT THPT NH25-26-Đợt 4- ĐỖ CHÍ TÂM]%[THPT Phan Bội Châu-Khánh Hòa]%[GVPB: Phạm Văn Long]
	Thâm niên công tác của các công nhân hai nhà máy A và B được cho như trong bảng dưới đây.
	\begin{table}[h]
		\centering
		\begin{tabular}{|l|c|c|c|c|c|}
			\hline
			Thâm niên công tác (năm) & \textbf{$[0; 5)$} & \textbf{$[5; 10)$}\, & \textbf{$[10; 15)$} & \textbf{$[15; 20)$} & \textbf{$[20; 25)$} \\ \hline
			Số công nhân nhà máy $A$ & $30$ & $18$ & $20$ & $22$ & $10$ \\ \hline
			Số công nhân nhà máy $B$ & $27$ & $15$ & $30$ & $16$ & $12$ \\ \hline
		\end{tabular}
	\end{table}\\
	Tính tổng hai phương sai về thâm niên công tác của công nhân làm việc tại hai nhà máy trên (kết quả được qui tròn đến hàng phần chục).\shortans{$91{,}4$}
	\loigiai{
		\begin{table}[h]
			\centering
			\begin{tabular}{|l|c|c|c|c|c|}
				\hline
				{Thâm niên công tác (năm)} & \textbf{$[0; 5)$} & \textbf{$[5; 10)$}\, & \textbf{$[10; 15)$} & \textbf{$[15; 20)$} & \textbf{$[20; 25)$} \\ \hline
				Giá trị đại diện &$2{,}5$& $7{,}5$ & $12{,}5$ & $17{,}5$ & $22{,}5$ \\ \hline
			\end{tabular}
		\end{table}\\
		{\bf Nhà máy A:}\\
		Tổng số công nhân: $N_A=30+18+20+22+10=100$.\\
		Số trung bình:
		$\overline{x}_A=\dfrac{30\cdot2{,}5+18\cdot7{,}5+20\cdot12{,}5+22\cdot17{,}5+10\cdot22{,}5}{100}=10{,}7$.\\  
		Phương sai:\\ $s_A^2=\dfrac{30(2{,}5-10{,}7)^2+18(7{,}5-10{,}7)^2+20(12{,}5-10{,}7)^2+22(17{,}5-10{,}7)^2+10(22{,}5-10{,}7)^2}{100}$\\
		\hspace*{0.6cm}$=46{,}76$.\\
		{\bf Nhà máy B:}\\
		Tổng số công nhân: $N_B=27+15+30+16+12=100$.\\
		Số trung bình:
		$\overline{x}_B=\dfrac{27\cdot2{,}5+15\cdot7{,}5+30\cdot12{,}5+16\cdot17{,}5+12\cdot22{,}5}{100}=11{,}05$.\\ 
		Phương sai:\\
		$s_B^2=\dfrac{27(2{,}5-11{,}05)^2+15(7{,}5-11{,}05)^2+30(12{,}5-11{,}05)^2+16(17{,}5-11{,}05)^2+12(22{,}5-11{,}05)^2}{100}$\\
		\hspace*{0.6cm}$=44{,}6475$.\\  
		Tổng hai phương sai: $s_A^2+s_B^2=46{,}76+44{,}6475=91,4075\approx 91{,}4$.
	}
\end{ex}

\begin{ex}%[1D6H2-5]%[TT-LienTruong-HaNoi-N25-26]%[Phạm Quốc Toàn - Lê Phúc]
	Cường độ trận động đất (Richter) được cho bởi công thức $M = \log A - \log A_0$, với $A$ là biên độ rung chấn tối đa và $A_0$ là biên độ chuẩn (hằng số). Đầu thế kỷ $20$, một trận động đất ở San Francisco có cường độ $6{,}7$ độ Richter. Trong cùng năm đó, một trận động đất khác ở Nam Mỹ có biên độ mạnh hơn gấp $4$ lần. Cường độ của trận động đất ở Nam Mỹ là bao nhiêu? (\textit{không làm tròn kết quả các phép tính trung gian, chỉ làm tròn kết quả cuối cùng đến hàng phần mười}).
	\par\shortans[0]{7{,}3}
	\loigiai{
		Gọi $M_1$, $A_1$ lần lượt là cường độ và biên độ rung chấn tối đa trận động đất tại San Francisco.\\
		Ta có 
		\[
		M_1 = 6{,}7 = \log A_1 - \log A_0 \Leftrightarrow 6{,}7 = \log \dfrac{A_1}{A_0} \Leftrightarrow A_1=10^{6{,}7} \cdot A_0. 
		\]
		Gọi $M_2$, $A_2$ lần lượt là cường độ và biên độ rung chấn tối đa trận động đất tại Nam Mỹ.\\
		Ta có $A_2=4A_1=4\cdot 10^{6{,}7}\cdot A_0$. \\ 
		Vậy cường độ trận động đất tại Nam Mỹ là 
		\[
		M_2 = \log \dfrac{A_2}{A_0} =\log \dfrac{4\cdot 10^{6{,}7} \cdot A_0}{A_0} = \log \left(4 \cdot 10^{6{,}7}\right) \approx 7{,}3. 
		\].
	}
\end{ex}

\begin{ex}%[1D6H3-5]%[Dự án đề TT Toán NH25-26-Dot 3- Hải Phụng]%[Cụm 13-Hải Phòng]%[GVPB - Viet Hoang Pham]
Một người gửi tiết kiệm vào ngân hàng $200$ triệu đồng theo hình thức lãi kép (tức là tiền lãi được cộng vào vốn của kỳ kế tiếp). Ban đầu người đó gửi với kỳ hạn $3$ tháng, lãi suất $2{,}1 \%$/kỳ hạn, sau $2$ năm người đó thay đổi phương thức gửi, chuyển thành kỳ hạn $1$ tháng với lãi suất $0{,}65 \%$/tháng. Tính tổng số tiền lãi nhận được sau $5$ năm (đơn vị là triệu đồng và kết quả làm tròn đến hàng phần mười).
\shortans[oly]{$98{,}2$}
\loigiai{Xét $2$ năm đầu tiên, số tiền nhận được là
$$T_1 = 200 (1 + 2{,}1\%)^{2 \cdot \frac{ 12}{3}} \approx 236{,}176 \text{ (triệu đồng).}$$
Số tiền nhận được sau $5$ năm là
$$T_2 = T_1 (1 + 0{,}65\%)^{3 \cdot 12} = 298{,}217 \text{ (triệu đồng).}$$ 
Tổng số tiền lãi nhận được sau $5$ năm là $298{,}217 - 200 = 98{,}2$  (triệu đồng). 			
}
\end{ex}

\begin{ex}%[1D7V2-8]%[TT Cụm 5 - Ninh Bình-NH25-26]%[GVSB:Dương Công Tạo]%[GVPB1: Nguyễn Tấn Tài]
	Giả sử dân số Việt Nam được dự báo theo mô hình logistic, giai đoạn từ năm 2\,023 đến năm 2\,035 là hàm số $P(t) = \dfrac{120}{1+0{,}2\mathrm{e}^{-0{,}06t}}$ (triệu người), trong đó $t$ là số năm tính từ 2\,023. Chi phí an sinh xã hội bình quân theo đầu người được mô hình hóa bằng hàm số $C(t) = 25-20\mathrm{e}^{-0{,}05t}$ (triệu đồng/đầu người/năm). Tính tốc độ thay đổi của tổng chi phí an sinh xã hội toàn quốc (nghìn tỷ đồng/năm) vào đầu năm 2\,030 (\textit{làm tròn kết quả đến hàng đơn vị}).
	\par  \shortans{83}
	\loigiai{
		Gọi $T(t)$ là tổng chi phí an sinh xã hội toàn quốc tại thời điểm $t$ (tính từ đầu năm 2\,023). \\
		Khi đó $T(t) = P(t)\cdot C(t)=\dfrac{120}{1+0{,}2\mathrm{e}^{-0{,}06t}}\cdot \left(25-20\mathrm{e}^{-0{,}05t}\right)$. \\
		Ta có $T'(t) = \dfrac{1{,}44\mathrm{e}^{-0{,}06t}}{(1+0{,}2\mathrm{e}^{-0{,}06t})^2}(25-20\mathrm{e}^{-0{,}05t}) + \dfrac{120}{1+0{,}2\mathrm{e}^{-0{,}06t}}\cdot\mathrm{e}^{-0{,}05t}$ \\
		Tốc độ thay đổi của tổng chi phí an sinh xã hội toàn quốc (nghìn tỷ đồng/năm) vào đầu năm 2\,030 là $T'(7) \approx 83$ (nghìn tỷ đồng/năm).
	}
\end{ex}

\begin{ex}%[1D9V2-5]%[Dự án C - Đợt 3 - Sở Giáo Dục Sơn La]%[Lê Hùng Thắng - Nguyễn Tiến]% 
	Trong một cuộc thi đấu Robotics, sân đấu được thiết kế dạng lưới ô vuông như hình vẽ.
	\begin{center}
		\begin{tikzpicture}[scale=1,>=stealth, font=\footnotesize, line join=round, line cap=round]
			\draw (0,0) grid (6,4);
			\foreach \x in {0,1,...,6}{
				\foreach \y in {0,1,...,4}
				\fill (\x,\y) circle (2.5pt);}
			\draw (2,2) node [above right] {$M$};
			\draw (3,1) node [above right] {$N$};
			\draw (0,4) node [above left] {$A$};
			\draw (6,0) node [below right] {$B$};
		\end{tikzpicture}
	\end{center}
	Các robot xuất phát từ vị trí điểm ${A}$, di chuyển ngẫu nhiên theo cạnh của các ô vuông theo hướng xuống dưới hoặc sang phải đến vị trí điểm ${B}$. Tính xác suất robot đi từ ${A}$ đến ${B}$ mà không đi qua cả ${M}$ và ${N}$ (kết quả làm tròn đến chữ số hàng phần trăm).
	\shortans{$0,31$}
	\loigiai{
		\begin{center}
			\begin{tikzpicture}[scale=1,>=stealth, font=\footnotesize, line join=round, line cap=round]
				\draw (0,0) grid (6,4);
				\foreach \x in {0,1,...,6}{
					\foreach \y in {0,1,...,4}
					\fill (\x,\y) circle (2.5pt);}
				\draw (2,2) node [above right] {$M$};
				\draw (3,1) node [above right] {$N$};
				\draw (0,4) node [above left] {$A$};
				\draw (6,0) node [below right] {$B$};
			\end{tikzpicture}
		\end{center}
		Số cách di chuyển từ $A$ đến $B$ là $n(\Omega)=\mathrm{C}_{10}^4=210$. \\
		Gọi $E_M\colon$\lq\lq Robot đi qua điểm $M$\rq\rq.\\
		Ta có $n(E_M)=\mathrm{C}_5^2\cdot \mathrm{C}_5^2=10\cdot 10=100$.\\
		Gọi $E_N\colon$\lq\lq Robot đi qua điểm $N$\rq\rq.\\
		Ta có $n(E_N)=\mathrm{C}_7^3\cdot \mathrm{C}_3^1=35\cdot 3=105$.\\
		Số cách robot đi qua cả $M$ và $N$ là\\
		$n(E_M\cap E_N)=\mathrm{C}_5^2\cdot \mathrm{C}_2^1\cdot \mathrm{C}_3^1=10\cdot2\cdot3=60$.\\
		Số cách robot đi qua ít nhất một trong hai điểm $M$ hoặc $N$ là
		$$n(E_M\cup E_N)=n(E_M)+n(E_N)-n(E_M\cap E_N)=100+105-60=145>$$
		Số cách robot đi từ $A$ đến $B$ mà không đi qua $M$ và $N$ là
		$$n(X)=210-145=65.$$
		Xác suất cần tìm là $\mathrm{P}(X)=\dfrac{65}{210}\approx 0{,}31$.
	}
\end{ex}

\begin{ex}%[GV soạn: Nguyễn Quang Hiệp GVPB2 Đỗ Minh Vũ ]%[1D9V2-5]
	Trong thư viện có $3$ quyển sách toán, $3$ quyển sách lý, $3$ quyển sách hóa, $3$ quyển sách sinh. Biết các quyển sách cùng môn giống nhau, xếp $12$ quyển sách trên lên giá thành một hàng. Tính xác suất sao cho khi xếp không có $3$ quyển nào cùng môn đứng cạnh nhau (Kết quả làm tròn đến hàng phần trăm).
	\par\shortans[oly]{0{,}84}
	\loigiai{
		Có tất cả $12$ quyển sách, trong đó có $4$ nhóm (gồm toán, lý, hóa và sinh) các quyển sách cùng môn giống nhau. Số cách xếp là $n(\Omega) = \dfrac{12!}{3! \cdot 3! \cdot 3! \cdot 3!} = 369\,600$ (cách).\\
		Gọi $A_1, A_2, A_3, A_4$ lần lượt là tập hợp các cách xếp mà môn toán, lý, hóa, sinh các quyển sách cùng môn đứng cạnh nhau.\\
		\textbf{Trường hợp 1} Một nhóm $3$ quyển sách cạnh nhau (xem $3$ quyển cùng môn là một khối và $9$ quyển còn lại)
		\[n(A_i) = \dfrac{10!}{3! \cdot 3! \cdot 3! \cdot 1!} \cdot \mathrm{C}_4^1 = 67\,200.\]
		\textbf{Trường hợp 2} Hai nhóm mà mỗi nhóm gồm $3$ quyển sách cạnh nhau (xem $2$ khối và $6$ quyển còn lại)
		\[n(A_i \cap A_j) = \dfrac{8!}{3! \cdot 3! \cdot 1! \cdot 1!} \cdot \mathrm{C}_4^2 = 6\,720.\]
		\textbf{Trường hợp 3} Ba nhóm mà mỗi nhóm gồm $3$ quyển sách cạnh nhau (xem $3$ khối và $3$ quyển còn lại)
		\[n(A_i \cap A_j \cap A_k) = \dfrac{6!}{3! \cdot 1! \cdot 1! \cdot 1!} \cdot \mathrm{C}_4^3 = 480.\]
		\textbf{Trường hợp 4} Cả bốn nhóm mà mỗi nhóm gồm $3$ quyển sách cạnh nhau (xem $4$ khối)
		\[n(A_1 \cap A_2 \cap A_3 \cap A_4) = \dfrac{4!}{1! \cdot 1! \cdot 1! \cdot 1!} \cdot \mathrm{C}_4^4 = 24.\]
		Khi đó \allowdisplaybreaks
		\begin{eqnarray*}
			n(A_1 \cup A_2 \cup A_3 \cup A_4) &= &n(A_i) - n(A_i \cap A_j) + n(A_i \cap A_j \cap A_k) - n(A_1 \cap A_2 \cap A_3 \cap A_4)\\
			&=& 67\,200 - 6\,720 + 480 - 24 \\
			&=& 60\,936.\\
		\end{eqnarray*}
		Gọi $E$ là biến cố \lq\lq xếp không có $3$ quyển nào cùng môn đứng cạnh nhau\rq\rq\\
		$n(E) = 369\,600 - 60\,936 = 308\,664$.\\
		Xác suất sao cho khi xếp không có $3$ quyển nào cùng môn đứng cạnh nhau là
		\[\mathrm{P}(E) = \dfrac{n(E)}{n(\Omega)} = \dfrac{308\,664}{369\,600} \approx 0{,}84.\]
	}
\end{ex}

\begin{ex}%[1H8C1-4]%[Dự án C đề thi thử TN Môn Toán NH25-26 - Dot 4 - Phạm Đăng Đăng]%[THPT Lê Quý Đôn-Đống Đa-Hà Nội]%[GVPB - Phan Quốc Trí]		
	Để chuẩn bị chào đón năm mới xuân Bính Ngọ 2026, người ta cần trang trí dây đèn Led quanh một tòa nhà hình kim tự tháp là hình chóp tứ giác đều $S.ABCD$. Biết rằng cạnh bên của tòa nhà dài $40$ m, góc $\widehat{BSC}=15^\circ$. Dây đèn Led được trang trí theo đường $AMNPQRTUVS$ (tham khảo hình vẽ), trong đó $VS=3$ m. Hỏi cần dùng ít nhất bao nhiêu mét dây đèn Led để trang trí (kết quả làm tròn đến hàng phần chục)?
	
	\begin{center}
		\begin{tikzpicture}[scale=1.2,font=\footnotesize,line join=round
			,line cap=round,>=stealth]
			% 1. Định nghĩa các đỉnh chính của hình chóp S.ABCD
			\coordinate (S) at (2.2,4.5);
			\coordinate (A) at (0,0);
			\coordinate (B) at (1.6,-1.3);
			\coordinate (C) at (6.5,-0.7);
			\coordinate (D) at (4.5,0.8);
			
			% 2. Xác định các điểm trên các cạnh
			\coordinate (O) at ($(A)!0.45!(S)$);
			\coordinate (V) at ($(A)!0.8!(S)$);
			
			\coordinate (M) at ($(B)!0.15!(S)$); % M trên SB
			\coordinate (N) at ($(C)!0.4!(S)$);  % N trên SC
			
			\coordinate (R_SB) at ($(B)!0.52!(S)$); % R duy nhất trên SB
			\coordinate (T) at ($(C)!0.7!(S)$);      % T trên SC
			
			% P nằm trên SD
			\coordinate (P) at ($(D)!0.42!(S)$);    
			
			% SỬA ĐIỂM U: U là giao của VT và SD (cạnh khuất)
			\coordinate (U) at ($(D)!0.7!(S)$);
			
			% 3. Vẽ các nét khuất (Dashed)
			\draw[dashed, gray!80] (A) -- (D) -- (C);
			\draw[dashed, gray!80] (S) -- (D);
			\draw[dashed, gray!80] (A) -- (D);
			\draw[dashed, gray!80] (O) -- (P); 
			\draw[dashed, gray!80] (P) -- (N);
			\draw[dashed, gray!80] (V) -- (U); % Đoạn VU khuất vì U nằm trên SD
			\draw[dashed, gray!80] (U) -- (T); % Đoạn UT khuất
			% 4. Vẽ các nét thấy (Solid)
			\draw(R_SB) -- (T); 
			\draw (S) -- (A) -- (B) -- (C) -- (S);
			\draw (A)--(M);
			\draw (S) -- (B);
			\draw (M) -- (N);
			\draw (O) -- (R_SB);
			% 5. Vẽ ký hiệu góc 15 độ tại đỉnh S
			\begin{scope}
				\path[clip] (S) -- (D) -- (R_SB) -- cycle;
				\draw (S) circle (0.7);
				\node at ($(S)+(-72:0.85)$) {\footnotesize $15^\circ$};
			\end{scope}
			
			% 6. Gán nhãn và vẽ các điểm
			\foreach \p/\pos/\txt in {
				S/above/S, A/left/A, B/below left/B, C/right/C, D/below/D,
				M/left/M, N/right/N, P/right/P, O/left/O, V/above left/V,
				T/right/T, R_SB/right/R%
			} {
				\fill (\p) circle (1pt);
				\node[\pos] at (\p) {\small $\txt$};
			}
			
			% Vẽ điểm U đặc biệt (màu xám, nằm trên SD)
			\fill (U) circle (1pt);
			\draw (U) circle (1pt);
			\node[below left] at (U) {\small $U$};
			
		\end{tikzpicture}
	\end{center}
	\shortans{$44{,}6$}
	\loigiai{\begin{center}\begin{tikzpicture}[scale=1,font=\footnotesize,line join=round
				,line cap=round,>=stealth]
				% 1. Định nghĩa đỉnh S ở trên cùng làm gốc
				\coordinate (S) at (0,3); % Đặt S cao lên
				
				% 2. Tinh chỉnh các góc và bán kính để tạo độ lệch ABCD, xòe rộng và hướng xuống dưới
				% Các góc được tính toán trong khoảng từ 190 đến 350 độ (nửa mặt phẳng dưới S)
				\foreach \i/\ang/\r in {
					1/195/4.8,  % A trái cùng
					2/215/4.5,  % B
					3/235/4.2,  % C
					4/255/4.0,  % D
					5/275/3.8,  % A
					6/295/3.6,  % B
					7/315/3.7,  % C
					8/335/4.1,  % D
					9/355/4.6    % A phải cùng
				} {
					\coordinate (P\i) at ($(S) + (\ang:\r)$);
				}
				
				% 3. Vẽ các tia từ S xuống
				\foreach \i in {1,...,9} {
					\draw (S) -- (P\i);
				}
				
				% 4. Vẽ đường biên ngoài (đáy quạt, không đều)
				\draw (P1) -- (P2) -- (P3) -- (P4) -- (P5) -- (P6) -- (P7) -- (P8) -- (P9);
				
				% 5. Xác định các điểm nút bên trong (M, N, P, Q, R, T, U, V)
				% Giữ nguyên tỉ lệ lệch để các điểm này trông tự nhiên
				\coordinate (M) at ($(S)!0.9!(P2)$);
				\coordinate (N) at ($(S)!0.85!(P3)$);
				\coordinate (P) at ($(S)!0.8!(P4)$);
				\coordinate (Q) at ($(S)!0.7!(P5)$);
				\coordinate (R) at ($(S)!0.9!(P6)$);
				\coordinate (T) at ($(S)!0.8!(P7)$);
				\coordinate (U) at ($(S)!0.85!(P8)$);
				\coordinate (V) at ($(S)!0.3!(P9)$); % V nằm trên tia A (thứ 5)
				
				% Vẽ đường gấp khúc thiết diện bên trong
				\draw[thick] (P1)--(M) -- (N) -- (P) -- (Q) -- (R) -- (T) -- (U);
				\draw[thick] (U) -- (V)--(P1); 
				
				% 6. Gán nhãn các điểm ngoài (A, B, C, D...)
				\node[left] at (P1) {$A$};
				\node[below left] at (P2) {$B$};
				\node[below] at (P3) {$C$};
				\node[below] at (P4) {$D$};
				\node[below] at (P5) {$A$};
				\node[below right] at (P6) {$B$};
				\node[right] at (P7) {$C$};
				\node[right] at (P8) {$D$};
				\node[right] at (P9) {$A$};
				\node[above] at (S) {$S$}; % S ở trên cùng
				
				% 7. Gán nhãn các điểm nút bên trong
				\node[below right=-2pt] at (M) {$M$};
				\node[below right=-2pt] at (N) {$N$};
				\node[below right=-1pt] at (P) {$P$};
				\node[right] at (Q) {$Q$};
				\node[right] at (R) {$R$};
				\node[left] at (T) {$T$};
				\node[above] at (U) {$U$};
				\node[above] at (V) {$V$};
				
				% 8. Vẽ chấm điểm
				\foreach \p in {S, P1, P2, P3, P4, P5, P6, P7, P8, P9, M, N, P, Q, R, T, U, V} {
					\fill (\p) circle (1.2pt);
				}
			\end{tikzpicture}
		\end{center}
		Vì dây đèn led quấn đủ $2$ vòng quanh tòa nhà trước khi kết thúc tại $V$ và đi về $S$. Khi đó, chúng ta phải trải phẳng $4$ mặt tam giác của hình chóp đều ($8$ mặt là $8$ tam giác cân tại $S$).\\
		Vì mỗi mặt bên là một tam giác cân có góc ở đỉnh $S$ là $15^\circ$ nên tổng góc ở đỉnh $S$ trên hình trải phẳng sẽ là $15^\circ \cdot 8 = 120^\circ$.\\
		Xét tam giác $SAV$ trên mặt phẳng trải hình (với góc $\widehat{ASV} = 120^\circ$, $SA = 40$ m, $SV = 3$ m)
		$$AV = \sqrt{SA^2 + SV^2 - 2 \cdot SA \cdot SV \cdot \cos 120^\circ} = \sqrt{40^2 + 3^2 - 2 \cdot 40 \cdot 3 \cdot (-0,5)} = \sqrt{1729} \approx 41{,}6 \text{ m}.$$
		Để độ dài đường gấp khúc $AMNPQRTUVS$ ngắn nhất thì $AV$ ngắn nhất. Khi đó tất cả các điểm $A$, $M$, $N$, $P$, $Q$, $R$, $T$, $U$, $V$ thẳng hàng.\\
		Do đó tổng chiều dài dây đèn ngắn nhất
		$L = 41{,}6 + 3 = 44{,}6$ m.
	}
\end{ex}

\begin{ex}%[1H8C6-1]%[Dự án C đợt 3 NH25-26- Phạm Đăng Đăng]%[TT-THPT-TranNhanTong-HaNoi-N25-26]%[GVPB - Phan Quốc Trí]
	Cho hình chóp $S.ABC$ có đáy là tam giác cân tại $A$. Cạnh bên $SA$ vuông góc với đáy. Gọi $H$ là trực tâm tam giác $ABC$. Tính $\cos\alpha$ với $\alpha$ là góc giữa đường thẳng $CH$ và mặt phẳng $(SBC)$. Biết $AB=AC=3$, $SA=BC=4$.
	\shortans{$0{,}81$}
	\loigiai{
		\begin{center}
			\begin{tikzpicture}[scale=1, line cap=round, line join=round, font=\small,>=stealth]
				\def\a{4} %Khai báo cạnh
				\def\h{4}
				\path 	(0:0) coordinate (A)
				++(0:\a) coordinate (C)
				++(-150:4*\a/5) coordinate (B)
				($(A)+(90:\h)$) coordinate (S)
				($(B)!0.5!(C)$) coordinate (M)
				($(A)!2/3!(B)$) coordinate (F)
				($(S)!5/7!(M)$) coordinate (K)
				;
				
				\coordinate (H) at (intersection of A--M and C--F);
				\draw[thick] 	(A)--(B)--(C)
				(A)--(S)	(C)--(S)	(B)--(S)--(M) (K)--(C);
				\draw[dashed,thick] (C)--(F)	(A)--(C) (A)--(M) (H)--(K);
				\foreach \x /\goc in {A/180,B/270,C/0,S/90,M/270,F/180,H/250,K/90}
				\fill[black] (\x) circle (1.5pt)
				($(\x)+(\goc:3mm)$) node {$\x$};
				;%Theo chiều dương
			\end{tikzpicture}
		\end{center}
		Gọi $M$ là trung điểm $BC$, $F$ là giao điểm của $CH$ và $AB$. Kẻ $HK$ vuông góc với $SM$ tại $K$.\\
		Ta có $\heva{&BC\perp AM\\&BC\perp SA}\Rightarrow BC \perp (SAM)\Rightarrow BC\perp HK$.\\Ta lại có $HK\perp SM$ nên ta suy ra $HK\perp (SBC)$. Do đó $\left(CH,(SBC)\right)=(CH,CK)=\widehat{HCK}=\alpha$.\\
		Ta có $AM=\sqrt{AB^2-BM^2}=\sqrt{5}$, $\tan\widehat{ BAM}=\dfrac{BM}{AM}=\dfrac{2}{\sqrt{5}}$.\\
		Ta có $\widehat {BAM}=\widehat{BCF}$ (cùng phụ $\widehat{ABC}$) nên $\tan \widehat{BCF}=\dfrac{FB}{FC}=\dfrac{HM}{MC}$. Suy ra $HM=\dfrac{4}{\sqrt{5}}$.\\
		$SM=\sqrt{SA^2+AM^2}=\sqrt{21}$.\\
		Lại có $\Delta SAM \sim \Delta HKM$, suy ra $\dfrac{SA}{HK}=\dfrac{MS}{MH}$, do đó $HK=\dfrac{SA\cdot MH}{MS}=\dfrac{16}{\sqrt{105}}$.\\
		$HC=\sqrt{HM^2+MC^2}=\dfrac{6}{\sqrt{5}}$ và $KC=\sqrt{HC^2-HK^2}=\dfrac{10}{\sqrt{21}}$. \\Do đó $\cos\alpha=\cos\widehat{ HCK}=\dfrac{KC}{HC}\approx 0{,}81$.}
\end{ex}

\begin{ex}%[1H8C6-6]%[Dự án đề TT Toán NH25-26-Dot 3- Hải Phụng]%[Cụm 13-Hải Phòng]%[GVPB - Viet Hoang Pham]
Cho hình chóp $S.ABC$ có đáy $ABC$ là tam giác vuông tại $A$, $AB=1$ cm, $AC=\sqrt{2}$ cm; $\widehat{SBA}=\widehat{SCA}=90^\circ$, góc giữa $BC$ và mặt phẳng $(SAB)$ bằng $45^\circ$. Tính khoảng cách giữa hai đường thẳng $SA$ và $BC$, với đơn vị là cm và kết quả làm tròn đến hàng phần trăm.
\shortans[oly]{$0{,}68$}
\loigiai{
\begin{center}
	\begin{tikzpicture}[line join=round, line cap=round, scale=0.9]
		%\draw[dashed,cyan] (-5,-5) grid (5,5);
		\def\d{5} \def\r{1.8} \def\h{4} \def\l{2.2}
		
		\coordinate[label={below left}:$B$] (B) at (-3,-3);
		\coordinate[label={below right}:$A$] (A) at ($(B)+(\d,0)$);
		\coordinate[label={above left}:$H$] (H) at ($(B)+(\l,\r)$);
		\coordinate[label={above right}:$C$] (C) at ($(H)+(\d,0)$);
		\coordinate[label={above}:$S$] (S) at ($(H)+(0,\h)$);
		\coordinate[label={above left}:{$K$}] (K) at ($(S)!0.6!(B)$);
		\coordinate[label={right}:{$O$}] (O) at ($(H)!0.5!(A)$);
		\coordinate[label={above right}:{$I$}] (I) at ($(S)!0.5!(H)$);
		\coordinate[label={below left}:{$M$}] (M) at ($(B)!0.8!(O)$);
		\coordinate[label={right}:{$N$}] (N) at ($(I)!0.5!(M)$);
		\draw (S)--(B)--(A)--(C)--(S)--(A);
		\draw[dashed] (S)--(H)--(B) (H)--(C)--(B) (H)--(A) (H)--(K) (O)--(I)--(C) (B)--(I) (H)--(M)--(I) (H)--(N);
		\draw 
		pic[draw, angle radius=2mm]{right angle=S--B--A}
		pic[draw, angle radius=2mm]{right angle=S--C--A}
		pic[draw, angle radius=2mm]{right angle=S--H--B}
		pic[draw, angle radius=2mm]{right angle=B--A--C}
		pic[draw, angle radius=2mm]{right angle=H--K--S}
		pic[draw, angle radius=2mm]{right angle=H--M--B}
		pic[draw, angle radius=2mm]{right angle=H--N--I}
		; 
		
		\foreach \x in {A,B,C,H,K,S,O,I,M,N} 
		\draw[fill=black] (\x) circle (1.5pt); 
		
	\end{tikzpicture}
\end{center}
Kẻ $SH\perp (ABC)$. Khi đó $AB\perp SH$ và $AB\perp SB$ nên $AB\perp (SBH)$.\\
Tương tự $AC\perp (SCH)$.\\
Do đó $AB\perp BH$ và $AC\perp CH$. Nghĩa là $ACHB$ là hình chữ nhật.\\
Vì góc giữa $BC$ và mặt phẳng $(SAB)$ bằng $45^\circ$ nên
$$\sin 45^\circ = \dfrac{\mathrm{d}\left[ C,(SAB)\right]}{BC}\Rightarrow \dfrac{\sqrt{2}}{2} = \dfrac{\mathrm{d}\left[ C,(SAB)\right]}{\sqrt{3}} \Rightarrow \mathrm{d}\left[ C,(SAB)\right] = \dfrac{\sqrt{6}}{2}. $$
Vì $CH\parallel (SAB)$ nên $\mathrm{d}\left[ H,(SAB)\right] =\mathrm{d}\left[ C,(SAB)\right] = \dfrac{\sqrt{6}}{2}$.\\
Kẻ $HK\perp SB$ tại $K$. Khi đó $HK\perp SB$ và $HK\perp AB$ nên $HK\perp (SAB)$. Suy ra $\mathrm{d}\left[ H,(SAB)\right] = HK =\dfrac{\sqrt{6}}{2}$.\\
Mà $HK=\dfrac{HS\cdot HB}{\sqrt{SH^2 + HB^2}}\Rightarrow \dfrac{\sqrt{6}}{2} = \dfrac{HS\cdot \sqrt{2}}{\sqrt{SH^2+2}}\Rightarrow SH=\sqrt{6}$.
\\
Gọi $I$ là trung điểm $SH$ và $O=HA\cap BC$. Khi đó $(BIC)\parallel SA$.\\
Suy ra $\mathrm{d}(SA,BC) = \mathrm{d}(SA,(BIC)) = \mathrm{d}(A,(BIC)) = \mathrm{d}(H,(BIC))$. \hfill $(1)$\\
Kẻ $HM\perp BC$ tại $M$, $HN\perp IM$ tại $N$. Khi đó $\mathrm{d}(H,(BIC)) = HN = \dfrac{HM\cdot HI}{\sqrt{HM^2+HI^2}}$.\\
Mà $HM=\dfrac{HC\cdot HB}{\sqrt{HC^2+HB^2}} = \dfrac{1\cdot \sqrt{2}}{\sqrt{1^2+\sqrt{2}^2}} =\dfrac{\sqrt{6}}{3} $ và $HI=\dfrac{SH}{2} = \dfrac{\sqrt{6}}{2}$.\\
Suy ra $HN= \dfrac{\sqrt{78}}{13}$. \hfill $(2)$\\
Từ $(1)$ và $(2)$ suy ra $\mathrm{d}(SA,BC) = \dfrac{\sqrt{78}}{13}\approx 0{,}68$.
}
\end{ex}

\begin{ex}%[Dự án đề TT Toán NH25-26-Dot 4- Hải Phụng]%[THCS-THPT LÊ THÁNH TÔNG - TP.HCM]%[GVPB - Viet Hoang Pham]%[1H8C7-9]
	\immini[thm]{
		Một tấm pin năng lượng mặt trời hình chữ nhật $ABCD$ có kích thước $AB = 2$ m và $BC = 3$ m. Tấm pin được đặt nghiêng sao cho cạnh $AB$ nằm sát trên mặt đất phẳng. Một bóng đèn (xem như một điểm) được đặt tại vị trí $S$ cao $4$ m có hình chiếu vuông góc lên mặt đất trùng với trung điểm $I$ của cạnh $AB$. Vào buổi tối khi bật đèn lên bóng của tấm pin trên mặt đất tạo thành một hình thang cân $ABC'D'$ (với $C', D'$ lần lượt là bóng của $CD$). Biết rằng hình thang cân $ABC'D'$ có chiều cao bằng $4{,}5$ m. Hãy tính diện tích lớn nhất của hình thang cân $ABC'D'$ là bóng của tấm pin trên mặt đất? (Đơn vị: m$^2$, làm tròn kết quả đến hàng phần chục).
	}{
		\begin{tikzpicture}[line join=round, line cap=round, scale=1]
			\def\d{5} \def\r{1.8} \def\h{4} \def\l{1}
			
			\coordinate[label={below left}:$B$] (B) at (-3,-3);
			\coordinate[label={below right}:$C'$] (C') at ($(B)+(\d,-0.3)$);
			\coordinate[label={below left}:$A$] (A) at ($(B)+(-\l,\r)$);
			\coordinate[label={above right}:$D'$] (D') at ($(A)+(\d+0.5,0.7)$);
			\coordinate[label={below left}:$I$] (I) at ($(A)!0.5!(B)$);
			\coordinate[label={right,xshift=1mm, yshift=2mm}:$S$] (S) at ($(I)+(0,\h)$);
			\coordinate[label={above right}:{$C$}] (C) at ($(S)!0.7!(C')$);
			\coordinate[label={above right}:{$D$}] (D) at ($(S)!0.7!(D')$);
			
			\draw (S)--(B)--(C')--(D')--(S)--(C') (S)--(A)--(B) (S)--(I) (B)--(C)--(D);
			\draw[dashed] (A)--(D') (D)--(A);
			\draw 
			pic[draw, angle radius=2mm]{right angle=S--I--B}
			; 
			\draw (S)  node[above,yshift=-1mm]{\Large \faLightbulbO};
			\draw[draw=none, fill=gray, opacity=0.3] (A)--(B)--(C)--(D)--cycle;
			
		\end{tikzpicture}
	}
	\shortans{$15{,}8$}
	\loigiai{
		\immini[thm]{
			Gọi $H$, $K$ lần lượt là trung điểm của $CD$ và $C'D'$.\\
			Vì $ABC'D'$ là hình thang cân, mà $C'D'>CD=AB$ nên đáy hình thang này bắt buộc phải là $AB$ và $C'D'$ (vì hai cạnh bên hình thang cân bằng nhau). Từ đó suy ra $IK$ là đường cao hình thang $ABC'D'$ nên  $IK=4{,}5$.\\
			Diện tích hình thang $ABC'D'$ là
			$$S=\dfrac{1}{2}\cdot (AB+C'D')\cdot IK = \dfrac{1}{2}\cdot (2+C'D')\cdot 4{,}5.$$
			Như vậy $S_{\max}$ khi và chỉ khi $C'D'$ lớn nhất. \\
			Dễ thấy $AB$ vuông góc với các đường thẳng $SI$, $IH$ và $IK$. \\
			Ta có $SK=\sqrt{SI^2+IK^2} = \dfrac{\sqrt{145}}{2}$. Khi đó
		}{%Bổ sung gói vẽ bóng đèn \usepackage{fontawesome5}
			\begin{tikzpicture}[line join=round, line cap=round, scale=1]
				\def\d{5} \def\r{1.8} \def\h{4} \def\l{1}
				
				\coordinate[label={below left}:$B$] (B) at (-3,-3);
				\coordinate[label={below right}:$C'$] (C') at ($(B)+(\d,-0.3)$);
				\coordinate[label={below left}:$A$] (A) at ($(B)+(-\l,\r)$);
				\coordinate[label={above right}:$D'$] (D') at ($(A)+(\d+0.5,0.7)$);
				\coordinate[label={below left}:$I$] (I) at ($(A)!0.5!(B)$);
				\coordinate[label={right,xshift=1mm, yshift=2mm}:$S$] (S) at ($(I)+(0,\h)$);
				\coordinate[label={above right}:{$C$}] (C) at ($(S)!0.7!(C')$);
				\coordinate[label={above right}:{$D$}] (D) at ($(S)!0.7!(D')$);
				\coordinate[label={right}:$H$] (H) at ($(C)!0.5!(D)$);
				\coordinate[label={right}:$K$] (K) at ($(C')!0.5!(D')$);
				
				\draw (S)--(B)--(C')--(D')--(S)--(C') (S)--(A)--(B) (S)--(I) (B)--(C)--(D) (S)--(K);
				\draw[dashed] (A)--(D') (D)--(A) (K)--(I)--(H);
				\draw 
				pic[draw, angle radius=2mm]{right angle=S--I--B}
				; 
				\draw 
				pic[draw, angle radius=2mm]{right angle=B--I--K}
				; 
				\draw 
				pic[draw, angle radius=2mm]{right angle=C'--K--S}
				; 
				\draw 
				pic[draw, angle radius=2mm]{right angle=C'--K--I}
				; 
				\draw (S)  node[above,yshift=-1mm]{\Large \faLightbulbO};
				\draw[draw=none, fill=gray, opacity=0.3] (A)--(B)--(C)--(D)--cycle;
			\end{tikzpicture}
		}
		\begin{eqnarray*}
			\dfrac{CD}{C'D'}=\dfrac{SH}{SK} & \Rightarrow& \dfrac{2}{C'D'}=\dfrac{2\cdot SH}{\sqrt{145}}\\
			& \Rightarrow& C'D'=\dfrac{\sqrt{145}}{SH}.
		\end{eqnarray*}
		Từ đó, để $C'D'$ lớn nhất thì $SH$ nhỏ nhất.\\
		Ta có
		{\allowdisplaybreaks
			\begin{eqnarray*}
				IH^2=SI^2+SH^2-2SI\cdot SH\cdot \cos \widehat{ISH} &\Leftrightarrow& 3^2=4^2+SH^2 -2\cdot 4\cdot SH\cdot \dfrac{8}{\sqrt{145}}\\
				&\Leftrightarrow& SH=2{,}4.
			\end{eqnarray*}
		}
		\noindent Suy ra $C'D'=5$ (chọn $SH$ nhỏ nhất để $C'D'$ lớn nhất).\\
		Vậy $S_{\max} = \dfrac{1}{2}\cdot (2+5)\cdot 4{,}5 = 15{,}75\approx 15{,}8$.
	}
\end{ex}

\begin{ex}%[1H8C7-9]%[Dự án C đề thi thử TN Môn Toán NH25-26 - Dot 4 - Phạm Hoàng Đăng]%[THPT Nguyễn Trãi-Hà Nội]%[GVPB - Khắc Thiên]
	\immini[thm]{Cho một tòa nhà đồ chơi có dạng hình hộp chữ nhật $ABCD.A'B'C'D'$ với đáy là hình vuông cạnh $20$ (cm) và chiều cao $80$ (cm), $A'B'C'D'$ là mặt dưới, $ABCD$ là mặt trên. Một con kiến xuất phát từ điểm $A'$ và bò lên bề mặt xung quanh của tòa nhà để đến đích là điểm $A$ (hình vẽ minh họa). Đường đi của con kiến là một đường gấp khúc liên tục, nằm hoàn toàn trên bốn mặt bên (không bò trên mặt trên và mặt dưới), không trùng với các cạnh của hình hộp và lần lượt gặp các cạnh bên theo đúng thứ tự $BB'$, $CC'$, $BB'$, $CC'$, $DD'$. Biết rằng con kiến chạm vào cạnh $BB'$ ở các điểm cách $B'$ một khoảng không nhỏ hơn $16$ (cm) và chạm vào cạnh $CC'$ ở các điểm cách $C$ một khoảng không lớn hơn $16$ (cm). Tính độ dài ngắn nhất của quãng đường mà con kiến đi (\textit{làm tròn kết quả đến hàng đơn vị, đơn vị độ dài là cm}).}{
		\begin{tikzpicture}[line cap=round,line join=round,scale=.7,font=\footnotesize]
			\path
			(0,0) coordinate (A)
			(4,0) coordinate (B)
			(2,2) coordinate (D)
			($(B)+(D)-(A)$) coordinate (C)
			($(A)!-6cm!90:(B)$) coordinate (A')
			($(A')+(D)-(A)$) coordinate (D')
			($(A')+(B)-(A)$) coordinate (B')
			($(B')+(D')-(A')$) coordinate (C');
			\draw (A)--(B)--(C)--(D)--cycle (A')--(B')--(C')
			(B)--(B') (C)--(C')(A)--(A')
			(A')--($(B')!1cm!(B)$)--($(C')!2cm!(C)$)--($(B')!3cm!(B)$)--($(C')!4cm!(C)$);
			\draw[dashed] (C')--(D')--(A') (D)--(D') ($(C')!4cm!(C)$)--($(D')!5cm!(D)$)--(A)
			;
			\foreach \i/\g in {A'/-150,B'/-40,C'/0,D'/140,A/180,B/-140,C/45,D/90}
			\draw[fill=black] (\i) circle (1pt) +(\g:.4) node{$\i$};
	\end{tikzpicture}}
	\shortans[oly]{157}
	\loigiai{
		Trải phẳng các mặt bên của hình hộp chữ nhật theo đường đi của con kiến ta được hình phẳng như hình vẽ.
		\begin{center}
			\begin{tikzpicture}[scale=1,>=stealth, line join=round, line cap=round, smooth]
				\path (0,0) coordinate (Ap1)
				(1,0) coordinate (Bp1)
				(2,0) coordinate (Cp1)
				(3,0) coordinate (Bp2)
				(4,0) coordinate (Cp2)
				(5,0) coordinate (Dp1)
				(6,0) coordinate (Ap2)
				(0,4) coordinate (A1)
				(1,4) coordinate (B1)
				(2,4) coordinate (C1)
				(3,4) coordinate (B2)
				(4,4) coordinate (C2)
				(5,4) coordinate (D1)
				(6,4) coordinate (A2)
				(2, 3.25) coordinate (H)
				(1, 1.625) coordinate (K)
				(1, 0.666) coordinate (E)
				(0,0) coordinate (O);
				\foreach \x in {1,2,3,4,5} \draw (\x,0) -- (\x,4);
				\draw (Ap1) -- (H) -- (A2)
				(Ap1) -- (A2)
				(A1) -- (A2) (Ap1) -- (Ap2) (0,0) -- (0,4) (6,0) -- (6,4);
				\foreach \x/\y/\z in {A1/135/A, B1/90/B, C1/90/C, B2/90/B, C2/90/C, D1/90/D, A2/45/A, Ap1/-140/A', Bp1/-90/B', Cp1/-90/C', Bp2/-90/B', Cp2/-90/C', Dp1/-90/D', Ap2/-40/A', H/160/H, K/135/K, E/135/E}
				\path[fill=black, draw] (\x) circle (.035)+(\y:.3) node{$\z$};
			\end{tikzpicture}
		\end{center}
		Xét trường hợp con kiến đi thẳng từ $A'$ đến $A$.\\ Gọi $E$ là điểm con kiến chạm vào cạnh $BB'$.\\
		Khi đó
		\[\dfrac{EB'}{AA'}=\dfrac{20}{120} \Rightarrow EB'=\dfrac{80\cdot 20}{120} \approx 13{,}3 \text{ (cm)}<16 \text{ (cm)} \text{ (không thỏa mãn).}\]
		Xét trường hợp con kiến chạm vào cạnh $CC'$ tại $H$ (như hình vẽ), với $0<CH \leq 16$.\\
		Khi đó $A'H$ đạt giá trị nhỏ nhất khi $CH=16$.\\
		Suy ra $C'H=64$ và $A'H=\sqrt{40^2+64^2}=8 \sqrt{89}$.\\
		Gọi $K$ là điểm con kiếm chạm cạnh $BB'$ trong trường hợp này.\\
		Ta có $\dfrac{KB'}{HC'}=\dfrac{20}{40} \Rightarrow KB'=\dfrac{64\cdot 20}{40}=32 \text{ (cm)}>16 \text{ (cm)}$ (thỏa mãn yêu cầu).\\
		Con kiến tiếp tục đi thẳng từ $H$ đến $A$, $AH=\sqrt{16^2+80^2}=16\sqrt{26}$.\\
		Vậy độ dài ngắn nhất của quãng đường mà con kiến đi là
		\[A'H+HA=8\sqrt{89}+16\sqrt{26} \approx 157{,}056 \text{ (cm)}.\]
	}
\end{ex}

\begin{ex}%[Đề thi thử môn Toán TRƯỜNG THPT Cổ Loa - Hà Nội năm học 2025-2026]%[Phạm Quốc Toàn - Phạm Hoàng Việt]%[1H8H5-4]
	Cho hình chóp $S\cdot ABCD$ có đáy $ABCD$ là hình vuông tâm $O$, cạnh bằng $1$. Cạnh bên $SA$ vuông góc với đáy và góc $\widehat{SBD}=60^{\circ}$. Khoảng cách giữa hai đường thẳng $AB$ và $SO$ bằng bao nhiêu? (không làm tròn kết quả các phép tính trung gian, chỉ làm tròn kết quả cuối cùng đến hàng phần trăm)
	\par\shortans[0]{0{,}45}
	\loigiai{
		\begin{center}
			\begin{tikzpicture}[font=\footnotesize, line join=round, line cap=round, >=stealth,scale=.7]
				\coordinate (S) at (3,8);
				\coordinate (A) at (3,3);
				\coordinate (B) at (0,0);
				\coordinate (C) at (6,0);
				\coordinate (D) at (9,3);
				\coordinate (O) at (4.5,1.5);
				\coordinate (E) at (6,3);
				\coordinate (H) at (5.1,4.5);
				%%%%%%%%%%%%%%%%%%%%%%%%%%%%%%%%
				\draw (S) -- (B) -- (C) -- (D) -- (S) -- (C);
				\draw[dashed] (S) -- (A) -- (B) -- (D) -- (A) (A) -- (C) (O) -- (E) -- (S) (A) -- (H);
				\foreach \i/\g in {S/90,A/180,B/180,C/0, D/0, O/-90, E/90, H/0}{\draw(\i) circle (1pt) ($(\i)+(\g:3mm)$) node[scale=1]{$\i$};}
				\pic[draw,thin,angle radius=2mm] {right angle=S--A--D}
				pic[draw,thin,angle radius=2mm] {right angle=A--H--E};
				%%%%%%%%%%%%%%%%%%%%%%%%%%%%%%%%
			\end{tikzpicture}
		\end{center}
		Vì $ABCD$ là hình vuông có cạnh bằng $1$ nên $BD=\sqrt{2}$. \\ 
		$SA \perp (ABCD) \Rightarrow \heva{& SA \perp AB \\ & SA \perp AD} \Rightarrow \Delta SAB=\Delta SAD \Rightarrow SB=SD$. \\ 
		Tam giác $SBD$ có $SB=SD$ và $\widehat{SBD}=60^{\circ} \Rightarrow \Delta SBD$ đều $\Rightarrow SB=BD=\sqrt{2}$. \\ 
		Tam giác $SAB$ có $\widehat{A}=90^{\circ} \Rightarrow SA=\sqrt{SB^2-AB^2}=\sqrt{\left(\sqrt{2}\right)^2-1^2}=1$. \\ 
		Gọi $E$ là trung điểm của $AD \Rightarrow OE$ song song $AB$, $OE \perp AD$, $AE=\dfrac{1}{2}$. \\
		Lại có $SA\perp OE \Rightarrow OE \perp (SAD)$. \\ 
		Kẻ $AH \perp SE \Rightarrow OE\perp AH\Rightarrow AH\perp (SEO)$. \\
		Do $AB$ song song $OE \Rightarrow AB$ song song $(SEO)$. \\
		Từ đó, ta có $\mathrm{d} \left(AB, SO\right)=\mathrm{d} \big(AB, (SEO)\big)=\mathrm{d}\big(A, (SEO)\big)=AH$. \\ 
		Ta có $AH=\dfrac{SA\cdot AE}{\sqrt{SA^2+AE^2}}=\dfrac{1\cdot \dfrac{1}{2}}{\sqrt{1^2+\left(\dfrac{1}{2}\right)^2}}=\dfrac{1}{\sqrt{5}}\approx 0{,}45$.
	}
\end{ex}

\begin{ex}%[1H8V5-3]%[Dự án C - Đợt 3 - Sở Giáo Dục Sơn La]%[Lê Hùng Thắng - Nguyễn Tiến]% 
	Cho hình chóp $S.ABCD$ có đáy $ABCD$ là hình thoi tâm $O$, cạnh bằng $1$, góc $\widehat{ABC} = 60^{\circ}$. Biết rằng $SO \perp (ABCD)$, $SO = \dfrac{3}{2}$. Tính khoảng cách từ $A$ đến mặt phẳng $(SCD)$ (không làm tròn kết quả ở các bước trung gian, làm tròn kết quả ở bước cuối cùng đến chữ số thập phân hàng phần trăm).
	\shortans{0,83}
	\loigiai{
		\immini[thm]{
			Ta có $AO \cap(SCD)=C$ nên
			\begin{eqnarray*}
				&& \dfrac{\mathrm{d}(A, (SCD))}{\mathrm{d}(O, (SCD))} = \dfrac{AC}{OC} = 2  \\
				&\Rightarrow & \mathrm{d}(A, (SCD)) = 2\mathrm{d}(O, (SCD)).
			\end{eqnarray*}
			Tam giác $ACD$ cân tại $D$ có $\widehat{ADC}=60^{\circ}$ nên $ACD$ là tam giác đều cạnh bằng $1$.\\
			Gọi $M$ là trung điểm $CD$ và $H$ là trung điểm của $CM$, khi đó\\
			$AM \perp CD \Rightarrow OH \perp CD$ tại $H$  (vì $OH \parallel AM$).\\
			Kẻ $OK \perp SH$ tại $K$, suy ra $OK \perp (SCD)$.\\
			Khi đó $\mathrm{d}(O, (SCD)) = OK$.
			
		}
		{\begin{tikzpicture}[scale=1, font=\footnotesize, line join=round, line cap=round, >=stealth]
				\path (0,0) coordinate (A)
				(-1.5,-1.5) coordinate (B)
				(2.5,-1.5) coordinate (C) 
				($(A)+(C)-(B)$) coordinate (D)
				($(A)!0.5!(C)$) coordinate (O)
				($(O)+(0,4)$) coordinate (S)
				($(C)!0.5!(D)$) coordinate (M)
				($(C)!0.25!(D)$) coordinate (H)
				($(S)!0.65!(H)$) coordinate (K);
				\draw (H)--(S)--(B)--(C)--(D)--(S)--(C);
				\draw[dashed] (S)--(A)--(B) (A)--(D) (A)--(C) (B)--(D) (S)--(O) (O)--(H) (O)--(K) (A)--(M);
				\foreach \p/\r in {S/90,A/135,B/-135,M/0,C/-45,D/45,O/-90,H/-10,K/30}
				\fill (\p) circle (1.5pt) node[shift={(\r:3mm)}]{$\p$};
			\end{tikzpicture}
		}
		\noindent
		$AM$ là đường cao của tam giác đều cạnh bằng $1$ nên đường cao $AM = \dfrac{\sqrt{3}}{2}$.	\\
		Ta có $OH=\dfrac{1}{2}AM = \dfrac{\sqrt{3}}{4}$.\\
		Xét tam giác $SOH$ vuông tại $O$ có $OK$ là đường cao nên ta có
		\begin{eqnarray*}
			& & \dfrac{1}{OK^2} = \dfrac{1}{SO^2} + \dfrac{1}{OH^2} =\dfrac{1}{\left(\dfrac{3}{2}\right)^2} + \dfrac{1}{\left(\dfrac{\sqrt{3}}{4}\right)^2}= \dfrac{4}{9} + \dfrac{16}{3} = \dfrac{52}{9} \\
			&\Rightarrow & OK = \sqrt{\dfrac{9}{52}} = \dfrac{3}{2\sqrt{13}}.
		\end{eqnarray*}
		Vậy $\mathrm{d}(A, (SCD)) = 2 \cdot \dfrac{3}{2\sqrt{13}} = \dfrac{3}{\sqrt{13}} \approx 0,83$.
	}
\end{ex}

\begin{ex}%[Đề thi thử môn Toán Sở giáo dục tỉnh Tuyên Quang năm học 2025-2026]%[Đoàn Minh Tâm]%[1H8V5-3]%[GVPB: Thanh Phát]
	Cho hình lăng trụ đứng $ABC.A'B'C'$ có đáy $ABC$ là tam giác vuông tại $A$, $BC=2\sqrt{3}$, $AB=3$. Khoảng cách giữa đường thẳng $AA'$ và mặt phẳng $(BCC'B')$ bằng bao nhiêu?
	\shortans[oly]{1,5}
	\loigiai{
		\immini[thm]{
			Vì $ABC.A'B'C'$ là hình lăng trụ đứng nên cạnh bên $AA'$ song song với cạnh bên $BB'$.\\
			Suy ra $AA' \parallel (BCC'B')$.\\
			Do đó $\mathrm{d}\left(AA',(BCC'B')\right)=\mathrm{d}\left(A,(BCC'B')\right)$.\\
			Trong mặt phẳng $(ABC)$, kẻ $AH \perp BC$ tại $H$.\\
			Ta có lăng trụ đứng nên $BB' \perp (ABC) \Rightarrow BB' \perp AH$.\\
			Từ $\heva{&AH \perp BC \\ &AH \perp BB'}\Rightarrow AH \perp (BCC'B')$.\\
			Vậy khoảng cách cần tìm chính là độ dài đoạn thẳng $AH$.\\
			Xét tam giác $ABC$ vuông tại $A$, áp dụng định lí Pythagore, ta có
			\[ AC = \sqrt{BC^2 - AB^2} = \sqrt{\left(2\sqrt{3}\right)^2 - 3^2} = \sqrt{12 - 9} = \sqrt{3}. \]}{
			\begin{tikzpicture}[line join=round,line cap=round,line width=.6pt,font=\footnotesize,scale=1,>=stealth]
				\coordinate[label=left:$A$] (A) at (0,0);
				\coordinate[label=below left:$B$] (B) at (1,-1);
				\coordinate[label=right:$C$] (C) at (4,0);
				\coordinate[label=left:$A'$] (A1) at ($(A)+(90:4)$);
				\coordinate[label=below left:$B'$] (B1) at ($(B)-(A)+(A1)$);
				\coordinate[label=right:$C'$] (C1) at ($(C)-(A)+(A1)$);
				\coordinate[label=below right:$H$] (H) at ($(B)!1/3!(C)$);
				\draw (A1)--(A)--(B)--(C)--(C1)--(A1)--(B1)--(C1) (B)--(B1);
				\draw[dashed] (A)--(C) (A)--(H);
				\draw ($(H)!5pt!(A)$)--($(H)!5pt!(A)+(H)!5pt!(B)-(H)$)--($(H)!5pt!(B)$)--(H)--cycle;
				\fill (A)circle(1pt) (B)circle(1pt) (C)circle(1pt) (A1)circle(1pt) (B1)circle(1pt) (C1)circle(1pt) (H)circle(1pt);
			\end{tikzpicture}
		}
		\noindent Áp dụng hệ thức lượng trong tam giác vuông $ABC$ với đường cao $AH$ là
		\[ AH \cdot BC = AB \cdot AC \Rightarrow AH = \dfrac{AB \cdot AC}{BC} = \dfrac{3 \cdot \sqrt{3}}{2\sqrt{3}} = \dfrac{3}{2} = 1{,}5. \]
		Vậy $\mathrm{d}\left(AA',(BCC'B')\right)=1{,}5$.
	}
\end{ex}

\begin{ex}%[1H8V5-4]%[Dự án đề thi thử NH25-26]%[THPT Chuyên Phan Bội Châu - Nghệ An]%[Phạm Hoàng Đăng]
	\immini[thm]{Cho hình lăng trụ đều $ABC.A'B'C'$. Biết $AA'=6$, góc giữa đường thẳng $A'B$ và mặt phẳng $(ABC)$ bằng $60^\circ$. Tính khoảng cách giữa hai đường thẳng $AA'$ và $BC$.}{
		\begin{tikzpicture}[declare function={a=3; b=0.55*a; h=0.9*a;}]
			\path (0:0) coordinate (A)
			(-43:b) coordinate (B)
			(a,0) coordinate (C)
			\foreach \x in {A,B,C}{(\x)++(90:h) coordinate (\x1)};
			\draw[dashed] (A)--(C);
			\draw (B1)--(A1)--(C1)--cycle
			(A1)--(A)--(B)--(B1)
			(C1)--(C)--(B);
			\foreach \x/\goc/\t in {A/180/A,B/-90/B,C/0/C,A1/100/A',B1/85/B',C1/80/C'}
			\fill (\x) circle (1pt) node[shift={(\goc:9pt)},font=\scriptsize]{$\t$};
		\end{tikzpicture}
	}
	\shortans{$3$}
	\loigiai{
		Vì $ABC.A'B'C'$ là hình lăng trụ đều nên cạnh bên $AA' \perp (ABC)$.\\
		Hình chiếu vuông góc của đường thẳng $A'B$ lên mặt phẳng $(ABC)$ là $AB$.\\
		Do đó, góc giữa $A'B$ và mặt phẳng $(ABC)$ là $\left(A'B, AB\right) = \widehat{A'BA} = 60^\circ$.\\
		\begin{center}
			\begin{tikzpicture}[declare function={a=3; b=0.55*a; h=0.9*a;}]
				\path (0:0) coordinate (A)
				(-43:b) coordinate (B)
				(a,0) coordinate (C)
				\foreach \x in {A,B,C}{(\x)++(90:h) coordinate (\x1)};
				\path (B) -- (C) coordinate[midway] (H);
				\draw[dashed] (A)--(C)
				(A)--(H);
				\draw (B1)--(A1)--(C1)--cycle
				(A1)--(A)--(B)--(B1)
				(C1)--(C)--(B)
				(A1)--(B);
				\foreach \x/\y/\z in {A/H/C}{
					\path (\y) pic[draw,angle radius=5pt]{right angle = \x--\y--\z};
				}
				\foreach \x/\goc/\t in {A/180/A,B/-90/B,C/0/C,A1/100/A',B1/85/B',C1/80/C',H/-45/H}
				\fill (\x) circle (1pt) node[shift={(\goc:9pt)},font=\scriptsize]{$\t$};
			\end{tikzpicture}
		\end{center}
		Xét tam giác vuông $A'AB$ tại $A$, ta có $AB = \dfrac{AA'}{\tan 60^\circ} = \dfrac{6}{\sqrt{3}} = 2\sqrt{3}$.\\
		Vì $AA' \parallel BB'$ nên $AA' \parallel (BCC'B')$.\\ 
		Khoảng cách cần tìm là
		\[\mathrm{d}(AA', BC) = \mathrm{d}\big(AA', (BCC'B')\big) = \mathrm{d}\big(A, (BCC'B')\big).\]
		Gọi $H$ là trung điểm của $BC$, vì tam giác đáy $ABC$ là tam giác đều cạnh $2\sqrt{3}$ nên $AH \perp BC$.\\
		Mặt khác, $AH \perp BB'$ (do $BB' \perp (ABC)$).\\ Suy ra $AH \perp (BCC'B')$.\\
		Vậy khoảng cách $\mathrm{d}\big(A, (BCC'B')\big) = AH = AB \cdot \dfrac{\sqrt{3}}{2} = 2\sqrt{3} \cdot \dfrac{\sqrt{3}}{2} = 3$.
	}
\end{ex}

\begin{ex}%[1H8V5-4]%[Dự án C đề thi thử TN Môn Toán NH25-26 - Dot 4 - Vũ Hồng Toàn]%[THPT Mai Anh Tuấn - Thanh Hóa]%[GVPB - Lê Hữu Kiệt]
	Cho hình lăng trụ đứng $ABC.A'B'C'$ có $AB = 5$, $BC = 5$, $CA = 7$. Khoảng cách giữa hai đường thẳng $AA'$ và $BC$ bằng bao nhiêu {\it (làm tròn kết quả đến hàng phần chục)}?
	\par\shortans{5,0}
	\loigiai{
		\begin{center}
			\begin{tikzpicture}[scale=1, font=\footnotesize, line join=round, line cap=round, >=stealth]
				\tikzset{declare function={a=3.5;b=0.65*a;goc=-40;h=0.75*a;}}
				\path (0:0) coordinate (A)
				(a,0) coordinate (C)
				(goc:b) coordinate (B)
				\foreach \p in {A,B,C} {($(\p)+(90:h)$) coordinate (\p')}
				($(C)!0.5!(B)$)coordinate (H);
				\draw[dashed]  (H)--(A)--(C);
				\draw (B')--(B)--(A)--(A')--cycle--(C')--(C)--(B) (A')--(C');
				\foreach \x /\gN in {A/180,B/-90,C/-20,A'/90,B'/80,C'/90,H/0}{
					\fill (\x) circle (1pt)($(\x)+(\gN:2.5mm)$) node {$\x$};
				}
			\end{tikzpicture}
		\end{center}
		Trong tam giác $ABC$, dựng $AH \perp BC$ với $H \in BC$.\\
		Khi đó $\heva{AH \perp BC\\ AH \perp BB'} \Rightarrow AH \perp (BCC'B')$.\\
		Ta có $p = \dfrac{a+b+c}{2} = \dfrac{5+5+7}{2} = \dfrac{17}{2}$.\\
		$S_{ABC} = \sqrt{\dfrac{17}{2}\left(\dfrac{17}{2}-5\right)\left(\dfrac{17}{2}-5\right)\left(\dfrac{17}{2}-7\right)} = \dfrac{7\sqrt{51}}{4}$.\\
		Mặt khác $S_{ABC} = \dfrac{1}{2}\cdot AH\cdot BC \Rightarrow AH = \dfrac{2S_{ABC}}{BC} = \dfrac{2 \cdot \dfrac{7\sqrt{51}}{4}}{5} \approx 5,0$.\\
		Vậy $\mathrm{d}(AA', BC) = \mathrm{d}\left(A,(BCC'B')\right) = AH \approx 5{,}0$.
	}
\end{ex}

\begin{ex}%[1H8V6-2]%[TT Cụm 5 - Ninh Bình-NH25-26]%[GVSB:Dương Công Tạo]%[GVPB1: Nguyễn Tấn Tài]
	\immini[thm]{Một mái nhà được tạo bởi hai nửa lục giác đều $ABCD$, $ABC'D'$ và hai tam giác bằng nhau $ADD'$, $BCC'$. Biết $CDD'C'$ là hình chữ nhật và $AB\parallel CD\parallel C'D'$, $CD=C'D'=2AB=6$ m, $DD'=4$ m. Tìm số đo góc nhị diện $[D', AD, C]$ (\textit{làm tròn kết quả đến hàng đơn vị của độ}).}
	{\begin{tikzpicture}[scale=1,line join=round, line cap=round,font=\footnotesize]
			\def\a{4}
			\path 	(0:0) coordinate (D)
			(20:\a) coordinate (C)
			(135:.55*\a) coordinate (D')
			(75:.32*\a) coordinate (A)
			($(D')-(D)+(C)$) coordinate (C')
			($(C)-(D)+(A)$) coordinate (At)
			($(A)!.42!(At)$) coordinate (B);
			\draw	(C)--(D)--(D')--(C')--cycle	(A)--(B)	(D)--(A)--(D')	(C)--(B)--(C');
			\foreach \x/\g in {C/0,D/-90,D'/180,C'/90,A/90,B/-90}
			\fill[black] 	(\x) circle (1pt)
			($(\g:3mm)+(\x)$) node {$\x$};	
		\end{tikzpicture}
	}
	\par  \shortans{121}
	\loigiai{
		\begin{center}
			\begin{tikzpicture}[scale=1.2,line join=round, line cap=round,font=\footnotesize]
				\def\a{4}
				\path 	(0:0) coordinate (D)
				(20:\a) coordinate (C)
				(135:.55*\a) coordinate (D')
				(80:.32*\a) coordinate (A)
				($(D')-(D)+(C)$) coordinate (C')
				($(C)-(D)+(A)$) coordinate (At)
				($(A)!.42!(At)$) coordinate (B)
				($(D)!.5!(C)$) coordinate (M)
				($(D)!.5!(A)$) coordinate (N)
				($(D')!.75!(D)$) coordinate (P);
				\draw	(C)--(D)--(D')--(C')--cycle	(A)--(B)	(D)--(A)--(D')	(C)--(B)--(C')	(P)--(N)--(M)	(A)--(C);
				\draw[dashed](M)--(P);
				\foreach \x/\g in {C/0,D/-90,D'/180,C'/90,A/90,B/120,M/-45,N/150,P/-150}
				\fill[black] 	(\x) circle (1pt)
				($(\g:3mm)+(\x)$) node {$\x$};	
				\draw pic[draw,angle radius=2mm]{right angle=P--N--A}
				pic[draw,angle radius=2mm]{right angle=M--N--A};
			\end{tikzpicture}
		\end{center}
		Goi $M$, $N$ lần lượt là trung điểm của $DC$, $DA$. \\
		Kẻ $NP \perp AD$ tại $N$, $P \in DD'$. Suy ra $\widehat{MNP}$ chính là góc phẳng nhị diện $[D', AD, C]$. \\
		Từ giả thiết suy ra $AM = AD = AD' = DM = 3$.\\
		Tam giác $ADM$ đều nên $MN = \dfrac{3\sqrt{3}}{2}$. \\
		Xét $\triangle ADD'$ có $\cos\widehat{ADD'} = \dfrac{DA^2 + DD'^2-AD'^2}{2 \cdot DA \cdot DD'} = \dfrac{3^2 + 4^2-3^2}{2 \cdot 3 \cdot 4} = \dfrac{2}{3}$. \\
		Xét $\triangle DNP$ vuông tại $N$ có 
		\begin{eqnarray*}
			\cos\widehat{NDP} = \dfrac{DN}{DP} &\Rightarrow& DP = \dfrac{DN}{\cos\widehat{NDP}} = \dfrac{9}{4} \\
			&\Rightarrow& PN = \sqrt{DP^2-DN^2} = \dfrac{3\sqrt{5}}{4}.
		\end{eqnarray*}
		Xét $\triangle DMP$ vuông tại $D$ có $MP = \sqrt{DM^2 + DP^2} = \dfrac{15}{4}$. \\
		Xét $\triangle MNP$ có 
		\begin{eqnarray*}
			\cos\widehat{MNP} &=& \dfrac{NM^2 + NP^2-PM^2}{2 \cdot NM \cdot NP} \\
			&=& \dfrac{\left(\dfrac{3\sqrt{3}}{2}\right)^2 + \left(\dfrac{3\sqrt{5}}{4}\right)^2-\left(\dfrac{15}{4}\right)^2}{2 \cdot \dfrac{3\sqrt{3}}{2} \cdot \dfrac{3\sqrt{5}}{4}}\\
			&=&-\dfrac{2\sqrt{15}}{15}.
		\end{eqnarray*}
		Suy ra $\widehat{MNP}\approx 121^\circ$. \\
		Vậy Số đo góc phẳng nhị diện $[D', AD, C]$ xấp xỉ $121^\circ$. 
	}
\end{ex}

\begin{ex}%[1H8V6-2]%[Dự án C đề thi thử TN Môn Toán NH25-26 - Dot 4 - Phạm Đăng Đăng]%[THPT Lê Quý Đôn-Đống Đa-Hà Nội]%[GVPB - Phan Quốc Trí]
	Cho hình chóp tứ giác đều $S.ABCD$ có cạnh đáy và cạnh bên đều bằng $a$. Côsin của số đo góc phẳng nhị diện $[D,SA,B]$ là $-\dfrac{a}{b}$ (với $\dfrac{a}{b}$ là phân số tối giản và $a$, $b\in\mathbb{N}^*$). Tính $P=a-b$.
	\shortans{$-2$}
	\loigiai{
		\begin{center}
			\begin{tikzpicture}[scale=1,font=\footnotesize,line join=round
				,line cap=round,>=stealth]
				\def\a{4}
				\def\h{4.5}
				\path 	(0:0) coordinate (A)
				++(0:\a) coordinate (D)
				++(-130:\a/2) coordinate (C)
				($(A)+(C)-(D)$) coordinate (B)
				(intersection of A--C and B--D) coordinate (O)
				($(O)+(90:\h)$) coordinate (S)
				
				($(A)!1/2!(S)$) coordinate (M);
				\draw[dashed,thick] 	(B)--(A)--(D)	(A)--(S)
				(A)--(C)	(B)--(D)	(S)--(O) (D)--(M)--(B)	;
				\draw[thick] 			(B)--(C)--(D)
				(B)--(S)	(C)--(S)	(D)--(S);
				\foreach \x/\g in {A/270,B/-135,C/-45,D/45,S/90,O/-90,M/-70}
				\fill[black] 	(\x) circle (1.5pt)
				($(\g:3mm)+(\x)$) node {$\x$};	
				\draw pic[draw,angle radius=3mm]{right angle=D--O--S};%Theo chiều dương	
			\end{tikzpicture}
		\end{center}
		Gọi $M$ là trung điểm của $SA$. Do các mặt bên của hình chóp là các tam giác đều cạnh $a$ nên $BM = DM = \dfrac{a\sqrt{3}}{2}$.\\
		Lại có $BM \perp SA$, $DM \perp SA$ nên số đo của nhị diện $[D,SA,B]$ bằng số đo góc $\widehat{BMD}$.\\
		Trong tam giác $BMD$ có
		$$
		\cos \widehat{BMD} = \dfrac{BM^2 + DM^2 - BD^2}{2\cdot BM \cdot DM}
		= \dfrac{\dfrac{3a^2}{4} + \dfrac{3a^2}{4} - 2a^2}{2\cdot \dfrac{a\sqrt{3}}{2} \cdot \dfrac{a\sqrt{3}}{2}}
		= \dfrac{-\dfrac{a^2}{2}}{\dfrac{3a^2}{2}}
		= -\dfrac{1}{3}.
		$$
		Vậy $a=1$, $b=3 \Rightarrow P=a-b=1-3=-2$.}
\end{ex}

\begin{ex}%[1H8V6-7]%[KNTT - Lớp 12 - Dự án C - Đợt 2]%[Loc do]%[PB: Phan Anh]
	\immini[thm]{Anh Nghĩa muốn xây một nhà kho dạng hình hộp chữ nhật có chiều rộng $AB=5$ m và làm thêm phần mái nhà $ABCDEF$ như hình vẽ, biết rằng $ADFE$ là một nửa của hình lục giác đều cạnh bằng $4$ m, $BCFE$ là hình thang cân có $EB=3$ m. Để đảm bảo tính chính xác cho việc thi công, Anh Nghĩa muốn xác định góc nhị diện $[B,AE,D]$. Hãy giúp anh Nghĩa tính góc nhị diện đó? (\textit{đơn vị: Độ, làm tròn kết quả đến hàng đơn vị}). }{\begin{tikzpicture}[declare function={a=1;b=4;h=2;},line join=round]
			\path (0,0) coordinate (B')
			(135:a) coordinate (A')
			(b,0) coordinate (C')
			($(C')-(B')+(A')$) coordinate (D')
			($(A')+(90:h)$) coordinate (A)
			($(B')-(A')+(A)$) coordinate (B)
			($(C')-(A')+(A)$) coordinate (C)
			($(D')-(A')+(A)$) coordinate (D);
			\path (0.3,3.5) coordinate (E);
			\path (2.8,3.5) coordinate (F);
			\draw ( B)--(B')--(C') (A)--(B)--(C) (A)--(E)--(B)  (C')--(C)--(F) (A)--(A')--(B') (E)--(F);
			\draw[dashed]  (A')--(D') (C')--(D')--(D)--(A) (F)--(D)--(C);
			\foreach \t/\g in {A/90,B/180,C/90,D/60,E/90,F/90}{
				\draw[fill=black] (\t) circle (1pt) node[shift={(\g:7pt)},font=\scriptsize]{$ \t $};
			}
	\end{tikzpicture}}
	\par\shortans[0]{140}
	\loigiai{
		\begin{center}
			\begin{tikzpicture}[declare function={a=1;b=4;h=2;},line join=round]
				\path (0,0) coordinate (B')
				(135:a) coordinate (A')
				(b,0) coordinate (C')
				($(C')-(B')+(A')$) coordinate (D')
				($(A')+(90:h)$) coordinate (A)
				($(B')-(A')+(A)$) coordinate (B)
				($(C')-(A')+(A)$) coordinate (C)
				($(D')-(A')+(A)$) coordinate (D);
				\path (0.3,3.5) coordinate (E);
				\path (2.8,3.5) coordinate (F);
				\path ($(A)!0.5!(D)$) coordinate (G)
				($(A)!0.5!(E)$) coordinate (H)
				($(A)!0.5!(B)$) coordinate (I);				
				\draw ( B)--(B')--(C') (A)--(B)--(C) (A)--(E)--(B)  (C')--(C)--(F) (A)--(A')--(B') (E)--(F) (H)--(I);
				\draw[dashed]  (A')--(D') (C')--(D')--(D)--(A) (F)--(D)--(C) (I)--(G)--(H);
				\foreach \t/\g in {A/90,B/180,C/90,D/60,E/90,F/90,G/-90,H/120,I/-120}{
					\draw[fill=black] (\t) circle (1pt) node[shift={(\g:7pt)},font=\scriptsize]{$ \t $};
				}
			\end{tikzpicture}
		\end{center}
		Gọi $G$, $H$, $I$ lần lượt là trung điểm của $AD$, $AE$, $AB$.\\
		Tam giác $AGE$ đều nên $GH\perp AE$.\\
		Ta có $AE^2+BE^2=4^2+3^2=5^2=AB^2$ suy ra tam giác $ABE$ vuông tại $E$.\\
		Suy ra $BE\perp AE$, mà $HI \parallel BE \Rightarrow HI\perp AE$.\\
		Mà $GH\perp AE$ nên $[B,AE,D]=\widehat{GHI}$.\\
		Ta có
		\begin{itemize}
			\item $GH=4\cdot \dfrac{\sqrt{3}}{2}=2\sqrt{3}$ m;
			\item $HI=\dfrac{1}{2} BE=\dfrac{3}{2}$ m;
			\item $AD=8$ m.
		\end{itemize}
		Hình chưc nhật $ABCD$ có $BD=\sqrt{BA^2+AD^2}=\sqrt{89} \Rightarrow GI=\dfrac{1}{2}DB=\dfrac{\sqrt{89}}{2}$ m.\\
		Từ đó suy ra $\cos GHI=\dfrac{HG^2+HI^2-GI^2}{2\cdot HG\cdot HI}=-\dfrac{4}{3\sqrt{3}}$.\\
		Vậy $\widehat{GHI}\approx 140^\circ$.
	}
\end{ex}

\begin{ex}%[TT-THPT-HamRong-NH25-26 ]%[GV soạn: Nguyễn Quang Hiệp GVPB2 Đỗ Minh Vũ ]%[1H8V7-2]
	Tính thể tích của khối chóp tứ giác đều $S.ABCD$ có chiều cao bằng $\sqrt{3}$ và độ dài cạnh đáy bằng $2\sqrt{3}$. (Kết quả làm tròn đến hàng phần trăm).
	\par\shortans[oly]{6,93}
	\loigiai{
		\begin{center}
			\begin{tikzpicture}[>=stealth,line join=round,line cap=round,font=\footnotesize,scale=1]
				\tikzset{
					pics/hinhChopTuGiacDeu/.style  n args={6}{
						code={
							\tikzset{
								declare function={a=3;b=1.5;h=3;goc=-140;}
							}	
							\path 
							(0,0)coordinate (#1)+(0:a)coordinate (#2)+(goc:b)coordinate (#4)			
							($(#2)+(#4)-(#1)$)coordinate (#3)
							(intersection of #1--#3 and #2--#4)coordinate (#5) ($(#5)+(90:h)$)coordinate (#6)
							;
						}
				}}
				\path 
				(0,0)pic {hinhChopTuGiacDeu={A}{B}{C}{D}{H}{S}}	;	
				
				\foreach \pointo/\pointt in {S/B,S/C,S/D,C/D,B/C}{
					\draw[fill=black](\pointo)--(\pointt);
				}	
				\foreach \pointo/\pointt in {S/A,A/B,A/D,A/C,B/D,S/H}{
					\draw[fill=black,dashed](\pointo)--(\pointt);
				}
				\foreach \point/\goc in {A/120,S/90,B/11,C/-45,D/200,H/-90}{
					\draw[fill=black](\point)circle(.8pt)+(\goc:2mm)node[scale=.8]{$\point$};
				}
			\end{tikzpicture}
		\end{center}
		Gọi $H$ là tâm của hình vuông đáy $ABCD$. Khi đó $SH$ là chiều cao của hình chóp.\\
		Theo giả thiết, ta có chiều cao $SH = \sqrt{3}$ và cạnh đáy $AB = 2\sqrt{3}$.\\
		Diện tích đáy của khối chóp tứ giác đều là $S_{ABCD} = \left(2\sqrt{3}\right)^2 = 12$.\\
		Thể tích của khối chóp là
		\[V = \dfrac{1}{3} \cdot S_{ABCD} \cdot SH = \dfrac{1}{3} \cdot 12 \cdot \sqrt{3} = 4\sqrt{3} \approx 6{,}93.\] 
	}
\end{ex}

\begin{ex}%[1H8V7-6]%[Dự án C đề thi thử TN Môn Toán NH25-26 - Dot 4 - Phạm Đăng Đăng]%[THPT Lê Quý Đôn-Đống Đa-Hà Nội]%[GVPB - Phan Quốc Trí]
	Để xây dựng các cột điện đường dây $220$ KV, người ta phải đổ bê tông những cột chân đế cột điện. Biết rằng các chân đế cột điện là các hình chóp cụt tứ giác đều có cạnh trên $4$ mét, cạnh dưới $6$ mét và các cạnh bên $4$ mét. Giá bê tông tươi là $1\,400\,000$ đồng trên $1$ mét khối. hỏi cần bao nhiêu tiền (đơn vị triệu đồng, kết quả làm tròn đến hàng đơn vị) để mua bê tông tươi làm một chân đế cột điện? 
	\shortans{$133$}
	\loigiai{\begin{center}
			\begin{tikzpicture}[scale=1,font=\footnotesize,line join=round
				,line cap=round,>=stealth]
				\def\a{4}
				\def\b{2}
				\def\h{3}
				\def\gocB{40}
				\def\goc{70}
				\path
				(0,0) coordinate (B')
				(\gocB:\b) coordinate (A')
				(0:\a) coordinate (C')
				(\goc:\h) coordinate (B)
				($(C')+(A')-(B')$) coordinate (D')
				(B)--++(0:\a*.7) coordinate (C)
				(B)--++(\gocB:\b*.5) coordinate (A)
				($(A)+(C)-(B)$) coordinate (D)
				;
				
				\coordinate (O) at (intersection of A--C and B--D);
				
				\coordinate (O') at (intersection of A'--C' and B'--D');
				\path ($(O')-(O)+(A)$) coordinate (E')
				;
				\coordinate (E) at (intersection of A--E' and A'--C');
				\draw (A)--(C)(B)--(D)(B)--(A)--(D)--(C)--cycle (B)--(B')--(C')--(C) (C')--(D')--(D);
				\draw[dashed] (O)--(O')(A')--(C') (A)--(E) (B')--(D')(B')--(A')--(D') (A')--(A);
				\foreach \x/\g in {A'/135,B'/180,C'/-45,D'/45,A/135,B/180,C/-45,D/45,O'/270,O/40,E/270}
				\fill (\x) circle (1pt) +(\g:0.3) node {$\x$};
			\end{tikzpicture}
		\end{center}
		Kẻ $AE \perp A'C'$, ta có $A'E = \dfrac{1}{2}(A'C' - AC) = \dfrac{1}{2}(6\sqrt{2} - 4\sqrt{2}) = \sqrt{2}$.\\
		Suy ra $AE = \sqrt{AA'^2 - A'E^2} = \sqrt{16 - 2} = \sqrt{14}$.\\Diện tích hai đáy khối chóp cụt là
		$S_1 = 4^2 = 16, S_2 = 6^2 = 36$.\\
		Vậy thể tích khối chóp cụt là $$V = \dfrac{1}{3}AE\left(S_1 + S_2 + \sqrt{S_1S_2}\right) = \dfrac{1}{3}\sqrt{14}\left(16 + 36 + \sqrt{16 \cdot 36}\right) = \dfrac{76\sqrt{14}}{3}.$$
		Số tiền cần tính là $T = V\cdot 1{,}4 \approx 133$ triệu đồng.
	}
\end{ex}

\begin{ex}%[1H8V7-9]%[Dự án đề thi thử NH25-26]%[THPT Chuyên Phan Bội Châu - Nghệ An]%[BC Tuan]
	\immini[thm]{Hai con thằn lằn A và B đang bám ở hai bức tường đối diện của một căn phòng dạng hình hộp chữ nhật với chiều rộng, chiều dài, chiều cao lần lượt là $8$ (m), $12$ (m), $5$ (m). Ban đầu thằn lằn A ở vị trí cách bức tường phía trước và trần nhà lần lượt là $7$ (m) và $3$ (m), còn thằn lằn B ở vị trí cách bức tường phía trước và trần nhà lần lượt là $9$ (m) và $4$ (m) (tham khảo hình vẽ bên). Sau đó chúng nhìn thấy nhau và chạy lại gặp nhau. Biết rằng hai con thằn lằn chỉ chạy trên các bức tường và trần nhà, hỏi tổng quãng đường ngắn nhất hai con thằn lằn di chuyển là bao nhiêu mét (làm tròn kết quả đến hàng đơn vị)?}{
		\begin{tikzpicture}[line join=round, line cap=round, declare function={w=4.5; h=3; d=6.2; goc=75; xd=1.6; wd=1.3; hd=1.8;},font=\footnotesize,scale=0.7]
			\path 
			(0,0) coordinate (A)
			(w,0) coordinate (B)
			(w,h) coordinate (C)
			(0,h) coordinate (D)
			(goc:d) coordinate (V)
			($(A)+(V)$) coordinate (A')
			($(B)+(V)$) coordinate (B')
			($(C)+(V)$) coordinate (C')
			($(D)+(V)$) coordinate (D')
			(xd,0) coordinate (D1)
			(xd,hd) coordinate (D2)
			(xd+wd,hd) coordinate (D3)
			(xd+wd,0) coordinate (D4)
			($(A)!0.55!(D')$) coordinate (F)
			($(B)!0.62!(C)$) coordinate (S_f)
			($(B')!0.62!(C')$) coordinate (S_b)
			($(S_f)!0.3!(S_b)$) coordinate (S);			
			\draw[dashed] (A) -- (A') (A') -- (B') (A') -- (D');
			\draw (A) -- (D1) -- (D2) -- (D3) -- (D4) -- (B) -- (C) -- (D) -- cycle;
			\draw (D) -- (D') -- (C') -- (C) (B) -- (B') -- (C');
			\node[below, font=\footnotesize] at ($(D1)!0.5!(D4)$) {Cửa};
			\draw (F) node[above]{$A$} (S) node[above right]{$B$};
		\end{tikzpicture}
	}
	\shortans{$15$}
	\loigiai{
		Ta dùng phương pháp trải phẳng hình hộp.\\
		Nếu hai con thằn lằn chạy gặp nhau bằng cách đi qua trần nhà, ta trải hình như sau:
		\begin{center}
			\begin{tikzpicture}[line join=round, line cap=round,font=\footnotesize,scale=0.4]
				\draw (0,0) rectangle (5,12) (5,0) rectangle (13,12) (13,0) rectangle (18,12);
				\draw (2,7) circle (1pt) node[above]{$A$} (17,8) circle (1.5pt) node[above]{$B$};
				\draw[dashed](2,7)--(2,0) node[midway,right]{$7$ (m)} (2,7)--(5,7) node[midway,above]{$3$ (m)} (17,8)--(17,0) node[midway,left]{$9$ (m)} (17,8)--(13,8) node[midway,above]{$4$ (m)};
				\draw[red] (2,7)--(17,8);
			\end{tikzpicture}
		\end{center}
		Khi đó khoảng cách ngắn nhất là đoạn thẳng nối hai điểm trên mặt phẳng trải \[AB=\sqrt{(3+8+4)^2+(9-7)^2} = \sqrt{15^2+2^2} = \sqrt{229} \approx 15{,}13 \,(\text{m}).\]
		Nếu hai con thằn lằn chạy gặp nhau bằng cách đi qua bức tường phía sau, ta trải hình như sau:
		\begin{center}
			\begin{tikzpicture}[line join=round, line cap=round,font=\footnotesize,scale=0.4]
				\draw (0,0) rectangle (12,5) (12,0) rectangle (20,5) (20,0) rectangle (32,5);
				\draw (7,2) circle (1.5pt) node[below]{$A$} (24,1) circle (1.5pt) node[below]{$B$};
				\draw[dashed](7,2)--(12,2) node[midway,above]{$5$ (m)} (24,1)--(20,1) node[midway,below]{$3$ (m)} (7,2)--(7,5) node[midway,left]{$3$ (m)} (24,1)--(24,5) node[midway,right]{$4$ (m)};
				\draw[red] (7,2)--(24,1);
			\end{tikzpicture}
		\end{center}
		Khi đó quãng đường là $AB=\sqrt{(5+8+3)^2+(4-3)^2} = \sqrt{16^2+1^2} = \sqrt{257} \approx 16{,}03$ (m).\\
		So sánh hai trường hợp, tổng quãng đường ngắn nhất mà hai con thằn lằn phải di chuyển là khoảng $15$ (m).
	}
\end{ex}

\begin{ex}%[1H8V7-9]%[Dự án C đề thi thử TN Môn Toán NH25-26 - Dot 4 - Phạm Hoàng Đăng]%[THPT Nguyễn Trãi-Hà Nội]%[GVPB - Khắc Thiên]
	\immini[thm]{Một chi tiết máy mẫu bằng nhựa đúc đặc có dạng hình chóp tứ giác đều, cạnh đáy bằng $20$ (cm). Do bị lỗi kĩ thuật, bên trong chi tiết này có một khoang rỗng. Để kiểm tra, kỹ sư thả chìm hoàn toàn chi tiết máy vào một bể chứa dung dịch chống gỉ có dạng hình lăng trụ đứng có đáy là hình vuông cạnh bằng $30$ (cm). Khi đó, mực dung dịch trong bể dâng lên thêm $2$ (cm) và dung dịch đã tràn vào lấp đầy khoang rỗng bên trong. Sau khi vớt chi tiết ra và lau khô bề mặt, khối lượng của nó tăng thêm $160$ (gam) so với ban đầu. Biết khối lượng riêng của dung dịch là $0{,}8$ (gam/cm$^3$). Hãy tính độ dài cạnh bên của hình chóp nói trên (\textit{làm tròn kết quả đến hàng phần mười}).}{
		\begin{tikzpicture}[scale=1,font=\footnotesize,line join=round,line cap=round,>=stealth]
			\draw[dashed,black,scale=2.5] (0,0) coordinate (A)--(-130:0.5) coordinate (B)
			(1,0) coordinate (D)--(A)--($(A)+(90:1.2)$) coordinate (A');
			\draw[black,scale=2.5] (B)--($(D)-(A)+(B)$) coordinate (C)--(D)--($(D)+(90:1.2)$) coordinate (D')--(A')--($(B)+(90:1.2)$) coordinate (B')--($(D')-(A')+(B')$) coordinate (C')--(D')
			(C')--(C)
			(B')--(B);
			\fill[cyan!55] ($(B)+(90:1.5)$)--($(C)+(90:1.5)$)--($(D)+(90:1.5)$)--($(A)+(90:1.5)$)--($(B)+(90:1.5)$);
			\fill[cyan!75] (B)--(C)--($(C)+(90:1.5)$)--($(B)+(90:1.5)$)--(B);
			\fill[cyan!90] (D)--(C)--($(C)+(90:1.5)$)--($(D)+(90:1.5)$)--(D);
			\draw[dashed,black,scale=2.5] (0,0) coordinate (A)--(-130:0.5) coordinate (B)
			(1,0) coordinate (D)--(A)--($(A)+(90:1.2)$) coordinate (A');
			\draw[black,scale=2.5] (B)--($(D)-(A)+(B)$) coordinate (C)--(D)--($(D)+(90:1.2)$) coordinate (D')--(A')--($(B)+(90:1.2)$) coordinate (B')--($(D')-(A')+(B')$) coordinate (C')--(D')
			(C')--(C)
			(B')--(B);
			\path ($(A)!.5!(C)$) coordinate (O)
			($(O)+(0,1.5)$) coordinate (S)
			($(A)!1/3!(O)$) coordinate (Aa)
			($(C)!1/3!(O)$) coordinate (Cc)
			($(B)!1/3!(O)$) coordinate (Bb)
			($(D)!1/3!(O)$) coordinate (Dd)
			;
			\draw[line width=0.25]
			(Bb)--(Cc)--(Dd) (S)--(Bb) (S)--(Dd) (S)--(Cc)
			;
			\draw[dashed, line width=0.25]
			(S)--(Aa) (Bb)--(Aa)--(Dd)
			;
		\end{tikzpicture}
	}
	\par
	\shortans[oly]{20{,}6}
	\loigiai{
		\immini[thm]{
			Chi tiết máy có dạng hình chóp tứ giác đều $S.ABCD$ như hình vẽ.\\
			Ta có thể tích khoang rỗng là $V_1=\dfrac{160}{0{,}8}=200$ (cm$^3$).\\
			Thể tích dung dịch dâng lên $V_2=30^2\cdot 2=1\,800$ (cm$^3$).\\
			Suy ra thể tích khối chóp là $V=200+1\,800=2\,000$ (cm$^3$).\\
			Vì $V=\dfrac{1}{3} \cdot 20^2 \cdot SO$ nên suy ra $SO=15$.\\
			Lại có $ABCD$ là hình vuông nên $BD=20\sqrt{2}$.\\
			Suy ra $BO=10\sqrt{2}$.\\
			Do đó $SB=\sqrt{SO^2+BO^2}=5\sqrt{17}\approx 20{,}6$ (cm).
		}{
			\begin{tikzpicture}[>=stealth,line join=round,line cap=round,font=\footnotesize,scale=1]
				\def\a{3}
				\def\b{1}
				\path (0,0) coordinate (O) +(90:\a) coordinate (S)
				+(140:\b) coordinate (A)
				(A)+(0:\a) coordinate (D)
				($(A)!2!(O)$) coordinate (C)
				($(D)!2!(O)$) coordinate (B)
				($(D)!.5!(C)$) coordinate (M);
				\draw (S)--(B)--(C)--(D)--cycle (S)--(C);
				\draw[dashed](S)--(A)--(D)
				(A)--(B)
				(S)--(O)
				(A)--(C)
				(B)--(D);
				\foreach \x/\g in {S/90,A/150,B/210,C/-30,D/0,O/-90} \fill (\x) circle (0.03)+(\g:3mm) node[scale=.9] {$\x$};
			\end{tikzpicture}
		}
	}
\end{ex}

\begin{ex}%[1H8V7-9]%[TN-THPT-NguyenTrungThien-HaTinh-NH25-26]%[Loc Do]%[PB: Phan Anh]
	Có một tấm bìa ngũ giác đều $MNPQR$ tâm $O$ cạnh bằng $30$ cm, cắt bỏ $5$ tam giác cân bằng nhau có cạnh đáy là cạnh của ngũ giác đều ban đầu, đỉnh là đỉnh của ngũ giác đều bên trong như hình vẽ rồi gập lên thành một khối chóp ngũ giác đều có cạnh đáy bằng $x$ (cm). Tìm $x$ để thể tích tạo thành là lớn nhất (\textit{kết quả làm tròn đến hàng phần chục}).
	\begin{center}
		\begin{tikzpicture}[scale=0.7, font=\footnotesize, line join=round, line cap=round, >=stealth]	
			\coordinate (A) at (240:1.5cm);
			\coordinate (B) at (312:1.5cm);
			\coordinate (C) at (24:1.5cm);
			\coordinate (D) at (96:1.5cm);
			\coordinate (E) at (168:1.5cm);
			\coordinate (M) at (60:4cm);
			\coordinate (R) at (132:4cm);
			\coordinate (Q) at (204:4cm);
			\coordinate (P) at (276:4cm);
			\coordinate (N) at (348:4cm);
			\path (intersection of A--M and B--R) coordinate (O);
			\path (intersection of E--D and O--R) coordinate (I);
			\fill [orange!10] (D)--(R)--(E)--(Q)--(A)--(P)--(B)--(N)--(C)--(M)--cycle;
			% Ngũ giác lớn (bắt đầu từ 60°)
			\draw[thick] (60:4cm)
			\foreach \i in {132,204,276,348} { -- (\i:4cm) } -- cycle;
			
			% Ngũ giác nhỏ (xoay 180° so với ngũ giác lớn)
			\draw[thick] (240:1.5cm)
			\foreach \i in {312,24,96,168} { -- (\i:1.5cm) } -- cycle;
			% Nối các đỉnh
			\foreach \i in {60,132,204,276,348} {
				\draw[thick] (\i:4cm) -- (\i+36:1.5cm); 
				\draw[thick] (\i:4cm) -- (\i-36:1.5cm);
			}
			\fill (O) circle (1.5pt) node[below right] {$O$};
%			\fill (I) circle (1.5pt) node[left] {$I$};
%			\draw[thick] (R)--(O);
			% Vẽ điểm
			\foreach \i/\name/\g in {60/M/30,132/R/120,204/Q/200,276/P/-90,348/N/0} {
				\fill[red] (\i:4cm) circle (1.5pt)+(\g:3mm) node[scale=0.8] {$\name$};
			}
			\foreach \i/\name/\t in {240/A/-110,312/B/-45,24/C/10,96/D/90,168/E/170} {
				\fill[blue] (\i:1.5cm) circle (1.5pt)+(\t:3mm) node[scale=0.8] {$\name$};
			}		
		\end{tikzpicture}
	\end{center}
	\shortans{14,8}
	\loigiai{\begin{center}
			\begin{tikzpicture}[scale=0.7, font=\footnotesize, line join=round, line cap=round, >=stealth]	
				\coordinate (A) at (240:1.5cm);
				\coordinate (B) at (312:1.5cm);
				\coordinate (C) at (24:1.5cm);
				\coordinate (D) at (96:1.5cm);
				\coordinate (E) at (168:1.5cm);
				\coordinate (M) at (60:4cm);
				\coordinate (R) at (132:4cm);
				\coordinate (Q) at (204:4cm);
				\coordinate (P) at (276:4cm);
				\coordinate (N) at (348:4cm);
				\path (intersection of A--M and B--R) coordinate (O);
				\path (intersection of E--D and O--R) coordinate (I);
				\coordinate (K) at ($(Q)!0.5!(R)$);
				\fill [orange!10] (D)--(R)--(E)--(Q)--(A)--(P)--(B)--(N)--(C)--(M)--cycle;
				% Ngũ giác lớn (bắt đầu từ 60°)
				\draw[thick] (60:4cm)
				\foreach \i in {132,204,276,348} { -- (\i:4cm) } -- cycle;
				
				% Ngũ giác nhỏ (xoay 180° so với ngũ giác lớn)
				\draw[thick] (240:1.5cm)
				\foreach \i in {312,24,96,168} { -- (\i:1.5cm) } -- cycle;
				% Nối các đỉnh
				\foreach \i in {60,132,204,276,348} {
					\draw[thick] (\i:4cm) -- (\i+36:1.5cm); 
					\draw[thick] (\i:4cm) -- (\i-36:1.5cm);
				}
				\fill (O) circle (1.5pt) node[below right] {$O$};
				\fill (I) circle (1.5pt) node[left] {$I$};
				\draw[thick] (R)--(O)--(M) (N)--(O)--(P) (K)--(O)--(Q);
				% Vẽ điểm
				\foreach \i/\name/\g in {60/M/30,132/R/120,204/Q/200,276/P/-90,348/N/0} {
					\fill[red] (\i:4cm) circle (1.5pt)+(\g:3mm) node[scale=0.8] {$\name$};
				}
				\foreach \i/\name/\t in {240/A/-110,312/B/-45,24/C/10,96/D/90,168/E/-100} {
					\fill[blue] (\i:1.5cm) circle (1.5pt)+(\t:3mm) node[scale=0.8] {$\name$};
				\fill[red] (K) circle (1.5pt)+(135:3mm) node[scale=0.8] {$K$};
				}		
			\end{tikzpicture}
		\qquad\qquad
		\begin{tikzpicture}[declare function={gocx=90;goc=-34;a=4;b=a*0.3;d=0.5*a;c=0.4*a;h=4;}]
			\path 
			(0,0) coordinate (A)
			(-0.5*goc:c) coordinate (E)
			(3.6,0) coordinate (D)
			(goc:b) coordinate (B)
			($(B)+(d,0)$) coordinate (C)
			($(D)!0.5!(C)$) coordinate (O1)	
			($(D)!0.5!(E)$) coordinate (I)	
			($(A)!0.55!(O1)$) coordinate (O)	
			($(gocx:h)+(O)$) coordinate (S);
			
			\draw[dashed] (A)--(E)--(D) (I)--(O)--(S)--(E) (S)--(I);
			\draw (A)--(B)--(C)--(D)--(S)--(A) (S)--(C) (S)--(B);
			\path pic[angle radius=3mm,draw=blue] {right angle = S--O--I};
			
			\foreach \t/\g in {A/-180,D/0,C/-60,S/90,B/-100,E/-90,O/-90, I/-90}{
				\draw[fill=black] (\t) circle (1pt)
				node[shift={(\g:7pt)},font=\scriptsize]{$\t$};
			}
		\end{tikzpicture}
		\end{center}
		$MNPQR$ là ngũ giác đều nên ta có \[\widehat{ROM}=\widehat{MON}=\widehat{NOP}=\widehat{POQ}=\widehat{QOR}=\dfrac{360^{^\circ}}{5} =72^{^\circ}.\]
		Xét tam giác $OQR$ cân tại $O$, đường cao $OK$ \[\widehat{ROK}=\widehat{KOQ}=\dfrac{\widehat{ROQ}}{2} =\dfrac{72^{^\circ}}{2} =36^{^\circ}.\]
		Xét tam giác $ROK$ vuông tại $K$
		\[\sin \widehat{ROK}=\dfrac{RK}{OR} \Rightarrow OR=\dfrac{RK}{\sin \widehat{ROK}} =\dfrac{RQ}{2\cdot\sin 36^{^\circ}} =\dfrac{30}{2\cdot\sin 36^{^\circ}} =\dfrac{15}{\sin 36^{^\circ}}\, (\rm{cm}).\]
		Ta có $EI=ID=\dfrac{ED}{2} =\dfrac{x}{2}$ (cm).\\
		Xét tam giác $OIE$ vuông tại $I$
		\[\tan \widehat{EOI}=\dfrac{EI}{OI} \Leftrightarrow OI=\dfrac{EI}{\tan \widehat{EOI}} =\dfrac{x}{2\cdot\tan 36^{^\circ}}\, (\rm{cm}).\]
		Suy ra $IR=OR-OI=\dfrac{15}{\sin 36^{^\circ}} -\dfrac{x}{2\cdot\tan 36^{^\circ}}$ (cm).\\
		Gọi đỉnh chóp là $S$, sau khi được gấp lên thì các điểm $M$, $N$, $P$, $Q$, $R$, $S$ trùng nhau.\\
		Chiều cao khối chóp là
		\[h=SO=\sqrt{IR^{2} -OI^{2}} =\sqrt{\left(\dfrac{15}{\sin 36^{^\circ}} -\dfrac{x}{2\tan 36^{^\circ}} \right)^{2} -\left(\dfrac{x}{2\tan 36^{^\circ}} \right)^{2}} \quad (\rm{cm}).\]
		Diện tích đáy khối chóp là
		\[S_{MNPQR} =5S_{OED} =5\cdot \dfrac{1}{2} \cdot OI\cdot DE =5\cdot \dfrac{1}{2} \cdot \dfrac{x}{2\tan 36^{^\circ}} \cdot x=\dfrac{5x^{2}}{4\tan 36^{^\circ}} \quad (\rm{cm}^2).\]
		Thể tích khối chóp là
		\begin{align*}
			V=\dfrac{1}{3} S_{MNPQR} \cdot h=  \dfrac{1}{3} \sqrt{\left(\dfrac{15}{\sin 36^{^\circ}} -\dfrac{x}{2\tan 36^{^\circ}} \right)^{2} -\left(\dfrac{x}{2\tan 36^{^\circ}} \right)^{2}} \cdot \dfrac{5x^{2}}{4\tan 36^{^\circ}}.		
		\end{align*}
		Ta có $V'=0 \Rightarrow x\approx 14{,}8$ (cm).\\
		\begin{center}
			\begin{tikzpicture}[>=stealth]
				\tkzTabInit[nocadre=false,lgt=1,espcl=2,deltacl=0.5]{$x$/1.2 ,$V'$/.7,$V$/2}
				{$ $ , $14{,}8$ , $ $}
				\tkzTabLine{ , + , $0$ , - , }
				\tkzTabVar{-/$ $ , +/$V_{max}$ , -/$ $}
			\end{tikzpicture}
			
		\end{center}
		Vậy $V_{max}$  khi  $x\approx 14{,}8$ (cm).}
\end{ex}

\begin{ex}%[2D1B4-3]%[Dự án C-TT THPT NH25-26-Đợt 4- ĐỖ CHÍ TÂM]%[THPT Phan Bội Châu-Khánh Hòa]%[GVPB: Phạm Văn Long]
	Cho hàm số $y=\dfrac{ax+b}{cx+d}$ có đồ thị như hình vẽ. Tính $a+d$.
	\shortans{$0$}
	\begin{center}
		\begin{tikzpicture}[x=1cm,y=1cm, scale=0.7]
			% Vẽ lưới (tùy chọn, để khớp với ảnh gốc)
			%\draw[step=0.5, lightgray!30, very thin] (-3,-2) grid (5,5);
			
			% Vẽ trục tọa độ
			\draw[->] (-3.2,0) -- (5.5,0) node[below] {$x$};
			\draw[->] (0,-2.2) -- (0,4.5) node[left] {$y$};
			%\node at (-0.2,-0.2) {$O$};
			
			% Vẽ đường tiệm cận
			\draw[dashed, thick] (1,-2.2) -- (1,4.5); % Tiệm cận đứng x=1
			\draw[dashed, thick] (-3,1) -- (5.5,1); % Tiệm cận ngang y=1
			
			% Vẽ đồ thị hàm số y = (x+2)/(x-1)
			% Nhánh 1: x < 1
			\draw[thick, red, domain=-3:0.7, samples=100] plot (\x, {(\x-2)/(\x-1)});
			% Nhánh 2: x > 1
			\draw[thick, red, domain=1.33:5.5, samples=100] plot (\x, {(\x-2)/(\x-1)});
			
			% Đánh dấu các số 1, 2 trên trục
			\fill (1,0) circle(1.5pt) node[shift={(45:3mm)}]{$1$};
			\fill (0,1) circle(1.5pt) node[shift={(45:3mm)}]{$1$};
			\fill (2,0) circle(1.5pt) node[below]{$2$};
			\fill (0,2) circle(1.5pt) node[left]{$2$};
			\fill (0,0) circle(1.5pt) node[shift={(-120:3mm)}]{$0$};
		\end{tikzpicture}
	\end{center}	
	\loigiai{
		Từ đồ thị hàm số, ta có$\colon$\\
		Tiệm cận đứng $x=-\dfrac{d}{c}=1\Leftrightarrow c=-d\,\,\,\, (1)$.\\
		Tiệm cận ngang $y=\dfrac{a}{c}=1\Leftrightarrow a=c\,\,\,\, (2)$.\\
		Từ $(1)(2)$, suy ra $a=-d\Leftrightarrow a+d=0$.
	}
\end{ex}

\begin{ex}%[2D1C3-6]%[Dự án C đợt 3 NH25-26- Phạm Đăng Đăng]%[TT-THPT-TranNhanTong-HaNoi-N25-26]%[GVPB - Phan Quốc Trí]
	Một chiếc thuyền ở vị trí $A$ cách bờ một khoảng $AB$ bằng $5$ km. Trên bờ biển có một cái kho ở vị trí $C$ cách $B$ một khoảng là $7$ km. Người lái thuyền dự định chèo thuyền từ $A$ đến điểm $M$ trên bờ biển với vận tốc $1{,}2$ m/s rồi đi bộ đến $C$ với vận tốc $1{,}8$ m/s (xem hình vẽ dưới đây). Điểm $M$ cách $B$ bao nhiêu km để người đó đến kho nhanh nhất?
	\begin{center}
		\begin{tikzpicture}[scale=1, line cap=round, line join=round, font=\small,>=stealth]
			
			% Các điểm
			\coordinate (B) at (0,0);
			\coordinate (C) at (7,0);
			\coordinate (M) at (3.5,0);
			\coordinate (A) at (0,5);
			
			% Vẽ đoạn BC
			\draw (-2,0)--(B) -- (C)--(9,0);
			
			% Vẽ AB
			\draw (B) -- (A);
			
			% Vẽ AM nét đứt
			\draw[dashed] (A) -- (M);
			
			% Ghi tên điểm
			\fill (A) circle (1pt) node[left] at (A) {A};
			\fill (B) circle (1pt)	node[below] at (B) {B};
			\fill (C) circle (1pt)	node[below] at (C) {C};
			\fill (M) circle (1pt)	node[below] at (M) {M};
			
			% Ghi độ dài AB = 5km
			
			\node[left] at (0,2.5) {5 km};
			
			% Ghi độ dài BC = 7km
			\draw[<->] (0,-0.7) -- (7,-0.7);
			\node[below] at (3.5,-0.7) {$7$ km};
			
		\end{tikzpicture}
	\end{center}
	\shortans{$4{,}47$}
	\loigiai{Vận tốc chèo thuyền $1{,}2$ m/s $=0{,}0012$ km/s.\\
		Vận tốc đi bộ $1{,}8$ m/s $= 0{,}0018$ km/s.\\
		Gọi $MB = x$ (km) với $0 \le x \le 7$.\\
		Khi đó quãng đường chèo thuyền là
		$S_1 = AM = \sqrt{AB^2 + BM^2} = \sqrt{25 + x^2}
		$.\\
		Thời gian chèo thuyền là
		$
		t_1 = \dfrac{AM}{0{,}0012} = \dfrac{\sqrt{25 + x^2}}{0{,}0012} = \dfrac{2500}{3}\sqrt{25 + x^2}\,(\text{s})
		$.\\
		Quãng đường đi bộ là
		$MC = 7 - x$.\\
		Thời gian đi bộ là $t_2 = \dfrac{MC}{0{,}0018} = \dfrac{7 - x}{0{,}0018} = \dfrac{5000}{9}(7 - x)\,(\text{s})$.\\
		Để người đó đến kho nhanh nhất thì
		$
		f(x) = t_1 + t_2 = \dfrac{2500}{3}\sqrt{25 + x^2} + \dfrac{5000}{9}(7 - x)
		$
		đạt giá trị nhỏ nhất trên đoạn $[0;7]$.\\
		Ta có 
		$
		f'(x) = \dfrac{2500}{3} \cdot \dfrac{x}{\sqrt{25 + x^2}} - \dfrac{5000}{9}
		$.\\Khi đó $
		f'(x) = 0\Leftrightarrow \dfrac{2500}{3} \cdot \dfrac{x}{\sqrt{25 + x^2}} - \dfrac{5000}{9}=0 \Leftrightarrow 2\sqrt{25+x^2}=3x\Leftrightarrow  x = 2\sqrt{5}$.\\
		Bảng biến thiên
		
		\begin{center}
			\begin{tikzpicture}
				\tkzTabInit[lgt=1.2,espcl=2.5,deltacl=0.6]
				{$x$/0.6,$f'(x)$/0.6,$f(x)$/2}{$0$,$2\sqrt{5}$,$7$}
				\tkzTabLine{,-,0,+,}
				\tkzTabVar{+/$f(0)$,-/$f(2\sqrt{5})$,+/$f(7)$}
			\end{tikzpicture}
		\end{center}
		Ta thấy $f(x)$ đạt giá trị nhỏ nhất tại $x=2\sqrt{5}\approx 4{,}47$ km.
	}
\end{ex}

\begin{ex}%[2D1H3-6]%[Dự án đề TT Toán NH25-26-Dot 3- Hải Phụng]%[Cụm 13-Hải Phòng]%[GVPB - Viet Hoang Pham]
Giả sử doanh số (tính bằng số sản phẩm) của một sản phẩm mới (trong vòng một số năm nhất định) tuân theo quy luật logistic được mô hình hóa bằng hàm số $f(t) = \dfrac{3\,000}{1 + 3\mathrm{e}^{-t}}$, $t \geq 0$, trong đó thời gian $t$ được tính bằng năm, kể từ khi phát hành sản phẩm mới. Khi đó đạo hàm $f'(t)$ sẽ biểu thị tốc độ bán hàng. Hỏi sau khi phát hành bao nhiêu năm thì tốc độ bán hàng là lớn nhất? (kết quả làm tròn đến hàng phần mười).
\shortans[oly]{$1{,}1$}
\loigiai{
Ta có $f(t) = \dfrac{3\,000}{1 + 3\mathrm{e}^{-t}} = \dfrac{3\,000\mathrm{e}^t}{\mathrm{e}^t + 3} \Rightarrow f'(t) = \dfrac{9\,000\mathrm{e}^t}{\left( \mathrm{e}^t + 3 \right)^2}$.\\
Xét $h(t) = f'(t) = \dfrac{9\,000\mathrm{e}^t}{\left( \mathrm{e}^t + 3 \right)^2}; (t \geq 0)$.\\
Ta có $h'(t) = \dfrac{ \left( 9\,000\mathrm{e}^t \right)' \cdot \left( \mathrm{e}^t + 3 \right)^2 - 2 \left(\mathrm{e}^t + 3\right) \cdot \left(\mathrm{e}^t + 3\right)' \cdot 9\,000\mathrm{e}^t}{\left( \mathrm{e}^t + 3 \right)^4} = \dfrac{9\,000\mathrm{e}^t(3 - \mathrm{e}^t)}{(\mathrm{e}^t + 3)^3}$.\\
Cho $h'(t) = 0 \Leftrightarrow 3 - \mathrm{e}^t = 0 \Leftrightarrow \mathrm{e}^t = 3 \Leftrightarrow t = \ln 3$.\\
Bảng biến thiên của hàm số $h(t)$:
\begin{center}
\begin{tikzpicture}
\tkzTabInit[lgt=1.2,espcl=3]
{$t$/0.8,$h'(t)$/0.8,$h(t)$/1.5}
{0,$\ln 3$,$+\infty$}
\tkzTabLine{,+,0,-,}
\tkzTabVar{-/ , +/ , -/ }
\end{tikzpicture}
\end{center}
Từ bảng biến thiên, suy ra tốc độ bán hàng $h(t)$ lớn nhất khi $t = \ln 3 \approx 1{,}1$.
}
\end{ex}

\begin{ex}%[TT-LienTruong-HaNoi-N25-26]%[Phạm Quốc Toàn - Lê Phúc]%[2D1H3-6]
	Một xí nghiệp chế biến cà phê bán trong nước và xuất khẩu ra nước ngoài. Nếu giá sản xuất mỗi kg cà phê là $x$\,(USD) thì sản lượng xí nghiệp sản suất được là $(2\,000x-150)$\,(kg) và lượng tiêu thụ trong nước là $(4\,000-500x)$\,(kg). Phần cà phê dư sẽ được xuất khẩu với giá cố định là $10$\,(USD) mỗi kg. Biết rằng với mỗi kg cà phê xuất khẩu thì xí nghiệp phải chịu thuế là $0{,}5$\,(USD). Hỏi giá sản xuất mỗi kg cà phê của xí nghiệp là bao nhiêu để lợi nhuận thu được từ xuất khẩu là lớn nhất?
	\shortans{$5{,}58$}
	\loigiai{
		Sản lượng cà phê dư thừa dành cho xuất khẩu là
		\[
		(2\,000x-150)-(4\,000-500x)=2\,500x-4\,150 \text{ (kg)}.
		\]
		Chi phí để sản xuất lượng cà phê dành cho xuất khẩu là $C(x)=(2\,500x-4\,150)x$. \\ 
		Doanh thu từ việc xuất khẩu là $D(x)=(2\,500x-4\,150)\cdot 10$. \\ 
		Chi phí chịu thuế $(2\,500x-4\,150)\cdot 0{,}5$. \\ 
		Lợi nhuận thu được từ xuất khẩu là 
		\begin{align*}
			L(x) & =(2\,500x-4\,150)\cdot 10-(2\,500x-4\,150)\cdot 0{,}5-C(x) \\ 
			& =(2\,500x-4\,150)\cdot 10-(2\,500x-4\,150)\cdot 0{,}5-(2\,500x-4\,150)\cdot x \\ 
			& =-2\,500 x^2+27\,900x-39\,425. 
		\end{align*}
		Hàm số $L(x)$ là một tam thức bậc hai có hệ số $a = -2\,500<0$ nên hàm số đạt giá trị lớn nhất tại $x = - \dfrac{27\,900}{2 \cdot (-2\,500)} = \dfrac{279}{50} \approx 5{,}58$. 
	}
\end{ex}

\begin{ex}%[TT-LienTruong-HaNoi-N25-26]%[Phạm Quốc Toàn - Lê Phúc]%[2D1H3-6]
	Để làm công tác cứu nạn ở một ngôi nhà có một hàng rào cao $2{,}2$ (m) song song và cách bức tường ngôi nhà một khoảng bằng $1{,}6$ (m). Người ta dựng một cái thang thẳng có một đầu chạm đất, đầu kia chạm vào bức tường ngôi nhà, thân của thang chạm vào mép trên hàng rào (tham khảo hình vẽ). Chiều dài ngắn nhất của cái thang là bao nhiêu? (\textit{đơn vị là mét, không làm tròn kết quả các phép tính trung gian, chỉ làm tròn kết quả cuối cùng đến hàng phần trăm}). 
	\begin{center}
		\begin{tikzpicture}[x={(1cm,0cm)}, y={(0cm,1cm)}, z={(0.4cm,0.3cm)}, scale=0.5, line join=round, line cap=round]
			\def\W{7}         % Chiều rộng ngôi nhà
			\def\H{5.5}       % Chiều cao tường nhà
			\def\D{4}         % Chiều sâu ngôi nhà
			\def\Pc{3.5}      % Vị trí tâm cổng
			\def\ladX{11.5}   % Vị trí chân thang (đặt ngoài bãi cỏ)
			\def\ladZ{-1}     % Vị trí thang theo chiều sâu z
			\def\xL{0}        
			\def\xR{\W}       
			\def\Pd{0.5}      
			% Thảm cỏ
			\draw[fill=green!80!black!20, draw=none] (-6, 0, -6) -- (14, 0, -6) -- (14, 0, \D+5) -- (-6, 0, \D+5) -- cycle;
			% Tường nhà bên phải và Cửa sổ cứu hộ
			\draw[fill=yellow!40!white, draw=gray!60!black, thick] (\W, 0, 0) -- (\W, 0, \D) -- (\W, \H, \D) -- (\W, \H, 0) -- cycle;
			% Khung viền ngoài cửa sổ bên phải
			\draw[fill=gray!30, draw=gray!60!black, thick] (\W, 4.15, 1.3) -- (\W, 4.15, 3.7) -- (\W, 5.3, 3.7) -- (\W, 5.3, 1.3) -- cycle;
			% Lớp kính cửa sổ
			\draw[fill=cyan!20, draw=gray!60!black, thick] (\W, 4.25, 1.5) -- (\W, 4.25, 3.5) -- (\W, 5.2, 3.5) -- (\W, 5.2, 1.5) -- cycle;
			% Thanh chia kính
			\draw[gray!60!black, thick] (\W, 4.725, 1.5) -- (\W, 4.725, 3.5);
			\draw[gray!60!black, thick] (\W, 4.25, 2.5) -- (\W, 5.2, 2.5);
			% Tường nhà mặt trước và Cửa chính
			\draw[fill=yellow!20!white, draw=gray!60!black, thick] (0, 0, 0) -- (\W, 0, 0) -- (\W, \H, 0) -- (0, \H, 0) -- cycle;
			\draw[fill=gray!30, draw=gray!60!black, thick] (\Pc-0.8, 0, 0) rectangle (\Pc+0.8, 2.5, 0);
			% --- CODE 2 CỬA SỔ MẶT TRƯỚC DUY THÊM VÀO ---
			% Cửa sổ bên trái
			% Khung viền ngoài
			\draw[fill=gray!30, draw=gray!60!black, thick] (0.6, 0.7, 0) rectangle (2.1, 2.3, 0);
			% Lớp kính
			\draw[fill=cyan!20, draw=gray!60!black, thick] (0.7, 0.8, 0) rectangle (2.0, 2.2, 0);
			% Thanh chia kính
			\draw[gray!60!black, thick] (1.35, 0.8, 0) -- (1.35, 2.2, 0);
			\draw[gray!60!black, thick] (0.7, 1.5, 0) -- (2.0, 1.5, 0);
			% Cửa sổ bên phải (mặt trước)
			% Khung viền ngoài
			\draw[fill=gray!30, draw=gray!60!black, thick] (4.9, 0.7, 0) rectangle (6.4, 2.3, 0);
			% Lớp kính
			\draw[fill=cyan!20, draw=gray!60!black, thick] (5.0, 0.8, 0) rectangle (6.3, 2.2, 0);
			% Thanh chia kính
			\draw[gray!60!black, thick] (5.65, 0.8, 0) -- (5.65, 2.2, 0);
			\draw[gray!60!black, thick] (5.0, 1.5, 0) -- (6.3, 1.5, 0);
			% ---------------------------------------------
			% Mái nhà
			\draw[fill=gray!60!black, draw=gray!20!black, thick] (0, \H, 0) -- (\W, \H, 0) -- (\W, \H+1.5, \D/2) -- (0, \H+1.5, \D/2) -- cycle;
			\draw[fill=gray!70!black, draw=gray!20!black, thick] (\W, \H, 0) -- (\W, \H, \D) -- (\W, \H+1.5, \D/2) -- cycle;
			% Cây cối (Bụi cây)
			\shade[ball color=green!50!black] (\W*0.9, 0.5, -0.5) circle (1.2cm);
			\shade[ball color=green!40!black] (\W*0.1, 0.5, -0.5) circle (1.0cm);
			\shade[ball color=green!45!black] (\xL+0.5, 0.5, -\Pd) circle (0.8cm);
			\shade[ball color=green!45!black] (\xR-0.5, 0.5, -\Pd) circle (0.8cm);
			\shade[ball color=green!35!black] (-1, 0.5, \D*0.2) circle (1.0cm);
			% Hàng rào bên trái
			\draw[fill=gray!20, draw=gray] (-3, 0, -3) -- (-3, 0, \D+3) -- (-3, 0.3, \D+3) -- (-3, 0.3, -3) -- cycle;
			\foreach \i in {0, 1, ..., 20} {
				\pgfmathsetmacro{\pz}{-3 + \i*(\D+6)/20}
				\draw[fill=white, draw=gray!50] (-3, 0.3, \pz-0.05) -- (-3, 0.3, \pz+0.05) -- (-3, 1.2, \pz+0.05) -- (-3, 1.2, \pz-0.05) -- cycle;
			}
			\draw[fill=gray!20, draw=gray] (-3, 1.0, -3) -- (-3, 1.0, \D+3) -- (-3, 1.1, \D+3) -- (-3, 1.1, -3) -- cycle;
			% Hàng rào bên phải
			\draw[fill=gray!20, draw=gray] (\W+2, 0, -3) -- (\W+2, 0, \D+3) -- (\W+2, 0.3, \D+3) -- (\W+2, 0.3, -3) -- cycle;
			\foreach \i in {0, 1, ..., 20} {
				\pgfmathsetmacro{\pz}{-3 + \i*(\D+6)/20}
				\draw[fill=white, draw=gray!50] (\W+2, 0.3, \pz-0.05) -- (\W+2, 0.3, \pz+0.05) -- (\W+2, 1.2, \pz+0.05) -- (\W+2, 1.2, \pz-0.05) -- cycle;
			}
			\draw[fill=gray!20, draw=gray] (\W+2, 1.0, -3) -- (\W+2, 1.0, \D+3) -- (\W+2, 1.1, \D+3) -- (\W+2, 1.1, -3) -- cycle;
			% Hàng rào mặt trước (hai bên cổng)
			\draw[fill=gray!10, draw=gray] (\W+2, 0, -3) -- (\Pc+1.5, 0, -3) -- (\Pc+1.5, 0.3, -3) -- (\W+2, 0.3, -3) -- cycle;
			\draw[fill=gray!10, draw=gray] (\Pc-1.5, 0, -3) -- (-3, 0, -3) -- (-3, 0.3, -3) -- (\Pc-1.5, 0.3, -3) -- cycle;
			\foreach \i in {0, 1, ..., 15} {
				\pgfmathsetmacro{\px}{\W+2 - \i*(\W+2 - (\Pc+1.5))/15}
				\draw[fill=white, draw=gray!50] (\px+0.05, 0.3, -3) -- (\px-0.05, 0.3, -3) -- (\px-0.05, 1.2, -3) -- (\px+0.05, 1.2, -3) -- cycle;
			}
			\foreach \i in {0, 1, ..., 15} {
				\pgfmathsetmacro{\px}{\Pc-1.5 - \i*(\Pc-1.5 - (-3))/15}
				\draw[fill=white, draw=gray!50] (\px+0.05, 0.3, -3) -- (\px-0.05, 0.3, -3) -- (\px-0.05, 1.2, -3) -- (\px+0.05, 1.2, -3) -- cycle;
			}
			\draw[fill=gray!20, draw=gray] (\W+2, 1.0, -3) -- (\Pc+1.5, 1.0, -3) -- (\Pc+1.5, 1.1, -3) -- (\W+2, 1.1, -3) -- cycle;
			\draw[fill=gray!20, draw=gray] (\Pc-1.5, 1.0, -3) -- (-3, 1.0, -3) -- (-3, 1.1, -3) -- (\Pc-1.5, 1.1, -3) -- cycle;
			% Cổng rào
			\draw[fill=gray!90!black, draw=black, thick] (\Pc+1.4, 0.1, -3) -- (\Pc-1.4, 0.1, -3) -- (\Pc-1.4, 1.4, -3) .. controls (\Pc, 1.7, -3) .. (\Pc+1.4, 1.4, -3) -- cycle;
			\foreach \bar in {1.2, 0.9, 0.6, 0.3, 0, -0.3, -0.6, -0.9, -1.2} {
				\pgfmathsetmacro{\bpos}{\Pc+\bar}
				\draw[white!80!gray, thick] (\bpos, 0.1, -3) -- (\bpos, 1.4, -3);
			}
			% Các cột trụ hàng rào
			\foreach \gx in {\Pc+1.5, \Pc-1.5, -3, \W+2} {
				\draw[fill=gray!20, draw=gray] (\gx+0.3, 0, -3-0.3) -- (\gx-0.3, 0, -3-0.3) -- (\gx-0.3, 1.6, -3-0.3) -- (\gx+0.3, 1.6, -3-0.3) -- cycle;
				\draw[fill=gray!30, draw=gray] (\gx-0.3, 0, -3-0.3) -- (\gx-0.3, 0, -3+0.3) -- (\gx-0.3, 1.6, -3+0.3) -- (\gx-0.3, 1.6, -3-0.3) -- cycle;
				\draw[fill=gray!40, draw=gray] (\gx+0.4, 1.6, -3-0.4) -- (\gx-0.4, 1.6, -3-0.4) -- (\gx-0.4, 1.7, -3+0.4) -- (\gx+0.4, 1.7, -3+0.4) -- cycle;
				\fill (\gx, 1.75, -3) circle (0.12cm);
				\fill[yellow!60!white] (\gx, 1.8, -3) circle (0.08cm);
			}
			% Chiếc thang cứu hộ
			\path
			(9,0) coordinate (A)
			+(-150:0.4) coordinate (A')
			(8,5) coordinate (B)
			($(A')+(B)-(A)$) coordinate (B')
			;
			\foreach \x in {1,...,10}{
				\pgfmathsetmacro{\k}{0.01+0.09*\x}
				\draw[gray!30!black, line width=1pt]($(A)!\k!(B)$)--+(-150:0.4);
			}
			\draw[gray!30!black, line width=1.5pt] (A')--(B') (A)--(B);
			% \node at (0,0){\href{2VFHU2}{\tiny \phantom{2VFHU2}}};
		\end{tikzpicture}
	\end{center}
	\par\shortans[0]{5{,}35}
	\loigiai{
		\immini[thm]{Gọi góc hợp bởi thang và mặt đất là $\widehat{CBA}=\alpha$ $\left(0 < \alpha < \dfrac{\pi}{2}\right)$.\\
		Xét $\triangle MHB$ vuông tại $H$, ta có $BM = \dfrac{2{,}2}{\sin \alpha}$.\\
		Xét $\triangle MNC$ vuông tại $N$, ta có $MC = \dfrac{1{,}6}{\cos \alpha}$.\\
		Khi đó, chiều dài chiếc thang là
		\[BC = BM + MC = \dfrac{2{,}2}{\sin \alpha} + \dfrac{1{,}6}{\cos \alpha}.\]
		Xét hàm số $f(x) = \dfrac{2{,}2}{\sin x} + \dfrac{1{,}6}{\cos x}$, với $x \in \left(0; \dfrac{\pi}{2}\right)$.\\
		Ta có $f'(x) = \dfrac{-2{,}2 \cdot \cos x}{\sin^2 x} + \dfrac{1{,}6 \cdot \sin x}{\cos^2 x} = \dfrac{-2{,}2 \cdot \cos^3 x + 1{,}6 \cdot \sin^3 x}{\sin^2 x \cdot \cos^2 x}$.}
		{\begin{tikzpicture}[scale=1, font=\footnotesize, line join=round, line cap=round, >=stealth]
				\path 
				(0,0) coordinate (A)
				(0,4) coordinate (C)
				(3,0) coordinate (B)
				($(B)!3/5!(C)$) coordinate (M)
				($(B)!(M)!(A)$) coordinate (H)
				($(A)!(M)!(C)$) coordinate (N)
				;
				\draw (A)--(B)--(C)--cycle (N)--(M)--(H);
				\path (M)--(H) node[pos=0.5,above,sloped]{$2{,}2$\,m};
				\path (A)--(H) node[pos=0.5,below]{$1{,}6$\,m};
				\foreach \x/\pos in {A/-135,H/-80,B/-45,M/45,C/90,N/180} \fill (\x) circle(1pt) node[{shift=(\pos:0.25)}]{$\x$}; 
				\pic[draw,angle radius=3mm] {angle = C--B--A}; 
				\node[shift=(150:0.5)] at (B) {$\alpha$};
				% \node at (0,0){\href{2VFHU2}{\tiny \phantom{2VFHU2}}};
		\end{tikzpicture}}
	\noindent
		Phương trình 
		\allowdisplaybreaks
		\begin{eqnarray*} 
			f'(x) = 0 &\Leftrightarrow& 2{,}2 \cdot \cos^3 x = 1{,}6 \cdot \sin^3 x  \\ 
			&\Leftrightarrow&  \tan^3 x = \dfrac{2{,}2}{1{,}6} \\ 
			&\Leftrightarrow&  \tan x = \sqrt[3]{\dfrac{11}{8}} .
		\end{eqnarray*}
		Chọn $x_0 \in \left(0; \dfrac{\pi}{2}\right)$ thỏa mãn $\tan x = \sqrt[3]{\dfrac{11}{8}}$.\\
		Ta có bảng biến thiên sau
		\begin{center}
			\begin{tikzpicture}
				\tkzTabInit[nocadre=false, lgt=1.2, espcl=4, deltacl=0.5]
				{$x$/0.7,$f'(x)$/0.7,$f(x)$/2}
				{$0$,$x_0$,$\tfrac{\pi}{2}$}
				\tkzTabLine{,-,0,+,}
				\tkzTabVar{+/$+\infty$,-/$f\left(x_0\right)$,+/$+\infty$}
				% \node at (T11){\href{2VFHU2}{\tiny \phantom{2VFHU2}}};
			\end{tikzpicture}
		\end{center}
		Dựa vào bảng biến thiên, ta thấy chiều dài thang bé nhất là $f\left(x_0\right)\approx 5{,}35$ mét.
	}
\end{ex}

\begin{ex}%[2D1H3-6]%[TN-THPT-NguyenTrungThien-HaTinh-NH25-26]%[Loc Do]%[PB: Phan Anh]
	Một đại lý nhập khẩu trái cây tươi để phân phối cho các cửa hàng. Mỗi lần nhập khẩu trái cây, khoán chi phí vận chuyển (không đổi) là $25$ triệu đồng. Chi phí bảo quản mỗi tạ trái cây dự trữ trong kho là $80$ nghìn đồng/ngày. Thời gian bảo quản trái cây trong kho tối đa $10$ ngày. Biết rằng, kể từ ngày đầu tiên nhập hàng, đại lý sẽ phân phối tới các cửa hàng $25$ tạ trái cây mỗi ngày. Mỗi lần nhập hàng, đại lý phải nhập đủ trái cây cho bao nhiêu ngày phân phối để chi phí trung bình cho mỗi ngày là thấp nhất (bao gồm chi phí vận chuyển và chi phí bảo quản trong kho)?
	\shortans{5}
	\loigiai{
		Giả sử mỗi lần nhập hàng, đại lý nhập đủ trái cây phân phối cho $n$ ngày $(1\le n\le 10)$.\\
		Chi phí vận chuyển của mỗi lần nhập hàng là $25$ triệu.\\
		Chi phí lưu kho là $25\cdot80\cdot(1+2+\ldots+n)=1\,000(n^{2} +n)$ nghìn đồng bằng $n^{2} +n$ triệu đồng.\\
		Vậy chi phí trung bình của mỗi ngày là $f(n)=\dfrac{n^{2} +n+25}{n} =n+1+\dfrac{25}{n}$.
		\[f'(n)=0\Leftrightarrow 1-\dfrac{25}{n^{2}} =0\Leftrightarrow n=5.\]
		Ta có $f(1)=27$; $f(5)=11$; $f(10)=13{,}5$.\\
		Vậy chi phí hàng ngày nhỏ nhất khi $n=5$.}
\end{ex}

\begin{ex}%[2D1K3-6]%[Dự án C-TT THPT NH25-26-Đợt 4- ĐỖ CHÍ TÂM]%[THPT Phan Bội Châu-Khánh Hòa]%[GVPB: Phạm Văn Long]
	Một tàu chở hàng di chuyển trên sông với vận tốc $v$ (km/giờ), $v>0$. Biết chi phí nhiên liệu cho mỗi kilômét là $0{,}3\sqrt{v}+\dfrac{1{,}6}{\sqrt{v}}$ (triệu đồng/km); chi phí bảo trì tính theo thời gian là $0{,}04v\sqrt{v}$ (triệu đồng/giờ). Ngoài ra, chi phí vận hành cố định khác (nhân công, bảo hiểm\ldots) là $17{,}25$ (triệu đồng/giờ). Gọi $v_0$ là vận tốc để tổng chi phí trên mỗi kilômét là nhỏ nhất và $C_{min}$ là giá trị chi phí thấp nhất đó. Tính giá trị biểu thức $P=v_0+100\cdot C_{min}$.
	\shortans{296}
	\loigiai{
		Chi phí nhiên liệu cho mỗi km là $0{,}3\cdot \sqrt{v}+\dfrac{1{,}6}{\sqrt{v}}$ (triệu đồng/km).\\
		Tàu chở hàng di chuyển với vận tốc $v$ (km/giờ) nên tàu đi được $1$ (km) trong $\dfrac{1}{v}$ (giờ).\\
		Vậy chi phí bảo trì trên $1$ (km) hay trong $\dfrac{1}{v}$ (giờ) là $0{,}04v\sqrt{v}\cdot\dfrac{1}{v}=0{,}04\sqrt{v}$ (triệu đồng).\\
		Chi phí vận hành cố định khác trên $1$ (km) hay trong $\dfrac{1}{v}$ (giờ) là $\dfrac{17{,}25}{v}$ (triệu đồng).\\
		Vậy tổng chi phí trên mỗi km của tàu chở hàng là\\
		$C\left(v \right)=0,3\sqrt{v}+\dfrac{1,6}{\sqrt{v}}+0{,}04\sqrt{v}+\dfrac{69}{4}\cdot \dfrac{1}{v}=0{,}34\sqrt{v}+\dfrac{1,6}{\sqrt{v}}+\dfrac{69}{4}\cdot \dfrac{1}{v}$\\
		$C'\left(v \right)=\dfrac{17}{100\sqrt{v}}-\dfrac{4}{5\sqrt{{v^3}}}-\dfrac{69}{4{v^2}}$.\\
		Xét $C'\left(v \right)=0$, ta đặt $\dfrac{1}{\sqrt{v}}=t\left(t{>}0\right)$ $\Rightarrow-\dfrac{69}{4}{t^4}-\dfrac{4}{5}{t^3}+\dfrac{17}{100}t=0\Leftrightarrow \hoac{&t=0\,\text{(loại)}\\&t=\dfrac{1}{5}\Rightarrow v=25.\\
		}$\\
		Ta lập được bảng biến thiên
		\begin{center}
			\begin{tikzpicture}
				\tkzTabInit[nocadre=true,lgt=1.2,espcl=2.5,deltacl=0.6]
				{$v$ /0.6,$C'(v)$ /0.6,$C(v)$ /2}
				{$0$,$25$,$+\infty$}
				\tkzTabLine{,-,$0$,+,}
				\tkzTabVar{+/, -/$2{,}71$,+/}
			\end{tikzpicture}
		\end{center}
		Vậy ${C_{\min}}=2{,}71\Leftrightarrow {v_0}=25\Rightarrow P=25+100\cdot 2{,}71=296$.
	}
\end{ex}

\begin{ex}%[2D1V2-1]%[Dự án C đề thi thử TN Môn Toán NH25-26 - Dot 4 - Vũ Hồng Toàn]%[THPT Mai Anh Tuấn - Thanh Hóa]%[GVPB - Lê Hữu Kiệt]
	Cho đồ thị hàm số $y = x^4 -3x^2 + ax+b$ có $A(2;-2)$ là một điểm cực tiểu. Giá trị của biểu thức $S = 59b-a$ bằng bao nhiêu?
	\par\shortans{2026}
	\loigiai{
		Ta có $y' = 4x^3-6x+a$, $y''=12x^2 -6$. \\
		Đồ thị hàm số $y = x^4 -3x^2 +ax+b$ có $A(2;-2)$ là một điểm cực tiểu suy ra
		\allowdisplaybreaks
		\begin{eqnarray*}
			\heva{&y(2) = -2\\ &y'(2) = 0\\ &y''(2) = 42 > 0}\Leftrightarrow \heva{&2a+b=-6\\ &a+20=0}\Leftrightarrow \heva{&a=-20\\ &b=34.}
		\end{eqnarray*}
		Vậy $S = 59b- a = 59 \cdot 34 + 20 = 2\,026$.
	}
\end{ex}

\begin{ex}%[2D1V2-7]%[TN-THPT-NguyenTrungThien-HaTinh-NH25-26]%[Loc Do]%[PB: Phan Anh]
	Một phần đường chạy của tàu lượn siêu tốc  khi gắn hệ trục tọa độ $Oxy$ được mô phỏng ở hình bên dưới, đơn vị trên mỗi trục là mét. Biết đường chạy của nó là một phần đồ thị hàm số bậc ba $y=ax^{3} +bx^{2} +cx+d$ $(0\le x<90)$; tàu lượn siêu tốc xuất phát từ điểm $A$, đi qua các điểm $C,\, D$ (ba điểm $A$, $C$, $D$ nằm trên đường thẳng song song với trục $Ox$) đồng thời đạt độ cao nhỏ nhất so với mặt đất là $4$ m. Độ cao lớn nhất mà tàu lượn siêu tốc đạt được là bao nhiêu mét so với mặt đất (\textit{kết quả làm tròn đến hàng phần chục})?
	\begin{center}
		\begin{tikzpicture}[x=1cm,y=1cm,scale=1,>=stealth]
			\def\aa{-0.08} \def\bb{1.04} \def\cc{-3.2} \def\dd{3}
			\draw[->] (-1.5,0)--(8.84,0)node[below right]{$x$};
			\draw[->] (0,-1)--(0,4.22)node[left]{$y$};
			\fill (0,0)node[below left]{$O$};
			\draw[black,samples=150,smooth,domain=-0.36:8.7] plot(\x,{\aa*(\x)^3+\bb*(\x)^2+\cc*\x+\dd});
			\draw [dashed] (0,3)--(8,3)--(8,0) (5,3)--(5,0) (0,0.1)--(2,0.1)--(2,0)--(2,0.2);
			\foreach \x/\g in {2/-90,5/-90,8/-90 }{
				\fill (\x,0) circle (1pt) +(\g:3mm) node[scale=0.8] {$\x0$};
			}
			
			\foreach \y/\g in {3/180 }{
				\fill (0,\y) circle (1pt) +(\g:3mm) node[scale=0.8] {$\y0$};
			}
			\foreach \x/\y/\g in {5/3/-30,8/3/0}{
				\fill (\x,\y) circle (1.5pt);
			}
			\fill (0,0.1) circle (1pt) node[left,scale=0.8] {$4$};
			\node[above right] at (0,3) {$A$};
			\node[above left] at (5,3) {$C$};
			\node[above right] at (8,3) {$D$};
		\end{tikzpicture}
	\end{center}
	\shortans{40,7}
	\loigiai{
		Nhận xét:  ta thấy đồ thị của hàm số $y=f(x)=ax^{3} +bx^{2} +cx+d\,(0\le x<90)$ cắt đường thẳng $d\colon y=30$ tại $3$ điểm phân biệt: $A(0;30)$, $C(50;30)$, $D(80;30)$.\\
		Do đó, phương trình $f(x)=30$ có $3$ nghiệm phân biệt $x_{1} =0$; $x_{2} =50$; $x_{3} =80$ hay ta có
		\begin{eqnarray*}
			&& f(x)-30=ax(x-50)(x-80) \\
			&\Leftrightarrow&  f(x)=ax(x-50)(x-80)+30\\
			&\Leftrightarrow&  f(x)=ax(x^{2} -130x+4\,000)+30\\
			&\Leftrightarrow & f(x)=a(x^{3} -130x^{2} +4\,000x)+30\\
			&\Rightarrow &  f'(x)=a(3x^{2} -260x+4\,000).
		\end{eqnarray*}
		Cho $f'(x)=0\Leftrightarrow a\left(3x^{2} -260x+4\,000\right)=0\Leftrightarrow \hoac{& x=20\\ & x=\dfrac{200}{3}.}$\\
		Ta có
		\begin{itemize}
			\item Hàm số $y=f(x)$ đạt cực tiểu tại $x=20$, giá trị cực tiểu là $f(20)=4$.
			\item Hàm số $y=f(x)$ đạt cực đại tại $x=\dfrac{200}{3}$, giá trị cực đại là $f\left(\dfrac{200}{3} \right)$.
		\end{itemize}
		Do $f(20)=4$, suy ra $a\cdot 20\cdot(20-50)\cdot(20-80)+30=4\Rightarrow a=-\dfrac{13}{18\,000}$.\\
		Giá trị cực đại của hàm số $y=f(x)$ là 
		\[f\left(\dfrac{200}{3} \right)=-\dfrac{13}{18\,000} \cdot\dfrac{200}{3} \cdot\left(\dfrac{200}{3} -50\right)\cdot\left(\dfrac{200}{3} -80\right)+30=\dfrac{9\,890}{243} \approx 40{,}7.\]
		Vậy độ cao lớn nhất mà tàu lượn siêu tốc đạt được là khoảng $40{,}7$ mét so với mặt đất.}
\end{ex}

\begin{ex}%[2D1V3-6]%[TT Cụm 5 - Ninh Bình-NH25-26]%[GVSB:Dương Công Tạo]%[GVPB1: Nguyễn Tấn Tài]
	Một nhà máy A chuyên sản xuất một loại sản phẩm cho nhà máy B, nhà máy A chỉ bán sản phẩm cho nhà máy B và nhà máy B cam kết thu mua hết số sản phẩm mà nhà máy A sản xuất được. Nhà máy A có khả năng sản xuất tối đa là $200$ tấn sản phẩm trong một tháng. Nếu bán ra $x$ tấn sản phẩm cho nhà máy B thì giá bán mỗi tấn sản phẩm là $50-0{,}0002x^2$ triệu đồng. Trong một tháng nhà máy A phải chi phí cho nhân công và chi phí khấu hao máy móc một lượng cố định là $150$ triệu đồng, ngoài ra khi sản xuất mỗi tấn sản phẩm thì nhà máy phải chi thêm cho mua nguyên vật liệu là $35$ triệu đồng. Biết rằng nhà máy A phải nộp $5\%$ doanh thu cho cơ quan thuế. Tính lợi nhuận sau thuế (lợi nhuận khi đã trừ thuế) lớn nhất thu được trong một tháng của nhà máy A (đơn vị tính là tỉ đồng và kết quả làm tròn đến hàng phần trăm).\\
	\par  \shortans{1,08}
	\loigiai{
		Doanh thu của nhà máy A trong một tháng là \[D(x)=\left(50-0{,}0002x^2\right)x = 50x-0{,}0002x^3\text{ (triệu đồng)}.\] 
		Tiền thuế nhà máy A phải nộp là $T(x)= D(x)\cdot 5\%$ (triệu đồng).\\
		Chi phí sản xuất $C(x)=150+35x$ (triệu đồng).\\
		Lợi nhuận sau thuế trong một tháng của nhà máy A là 
		\begin{eqnarray*}
			L(x)&=& D(x)- D(x)\cdot 5\%-C(x) \\
			&=& 0{,}95\cdot (50x-0{,}0002x^3)-(150+35x)\\
			&=& -0{,}00019x^3 +12{,}5x-150\text{ (triệu đồng)}.
		\end{eqnarray*}
		Ta có $L'(x) = -0{,}00057x^2 +12{,}5$; $L'(x) = 0 \Leftrightarrow x = \pm \sqrt{\dfrac{12{,}5}{0{,}00057}}\approx 148$.\\
		Bảng biến thiên
		\begin{center}
			\begin{tikzpicture}[>=stealth]
				\tkzTabInit[nocadre=false,lgt=1.5,espcl=3,deltacl=0.6]{$x$/.7 ,$L'(x)$/.7,$L(x)$/2}
				{$0$ , $148$ , $200$}
				\tkzTabLine{ , + , $0$ , - , }
				\tkzTabVar{-/$-150$ , +/$1\,084$ , -/$830$}
			\end{tikzpicture}
		\end{center}
		Lợi nhuận sau thuế lớn nhất trong một tháng của nhà máy A là \[\max L(x) = L(148)\approx 1\,084\text{ (triệu đồng)}\approx 1{,}08\text{ (tỉ đồng)}.\] 
	}
\end{ex}

\begin{ex}%[2D1V3-6] %[Dự án C đề thi thử TN Môn Toán NH25-26 - Dot 4 - Hector Tran]%[Liên trường THPT cụm 09 - Đợt 2 - Hà Nội]%[GVPB - Thanh Phát]
	Một doanh nghiệp sản xuất độc quyền một loại sản phẩm. Giả sử khi sản xuất và bán hết $x$ sản phẩm ($0 < x \le 2\,500$), tổng số tiền doanh nghiệp thu được là $f(x) = 2\,026x - x^2$ và tổng chi phí là $g(x) = x^2 + 1\,438x - 1\,209$ (đơn vị: nghìn đồng). Giả sử mức thuế phụ thu trên một đơn vị sản phẩm bán được là $t$ (nghìn đồng) ($0 < t < 320$). Giá trị của $t$ bằng bao nhiêu nghìn đồng để nhà nước nhận được số tiền thuế phụ thu lớn nhất và doanh nghiệp cũng nhận được lợi nhuận lớn nhất theo mức thuế phụ thu đó?
	\par\shortans{294}
	\loigiai{
		Ta có hàm lợi nhuận khi sản xuất và bán hết $x$ sản phẩm là
		\begin{eqnarray*}
			P(x) &=& f(x) - g(x) - xt \\
			&=& 2\,026x - x^2 - (x^2 + 1\,438x - 1\,209) - xt \\
			&=& -2x^2 + (588 - t)x + 1\,209\text { (nghìn đồng).}
		\end{eqnarray*}
		Khi lợi nhuận lớn nhất thì $x = \dfrac{-b}{2a} = \dfrac{588 - t}{4}$.\\
		Số tiền thuế thu được là $T = xt = \dfrac{588t - t^2}{4}=-\dfrac{1}{4}t^2+147t$ (nghìn đồng).\\
		Số tiền thuế lớn nhất khi $t = \dfrac{-147}{2 \cdot \left( -\dfrac{1}{4}\right) } = 294\text { (nghìn đồng) } \in (0; 320)$ (thỏa mãn).\\
		Vậy $t=294$ (nghìn đồng).
	}
\end{ex}

\begin{ex}%[2D1V3-6]%[Dự án đề TT Toán NH25-26-Dot 3- Hải Phụng]%[Cụm 13-Hải Phòng]%[GVPB - Viet Hoang Pham]
Để hạn chế vi phạm thời gian làm việc đối với công nhân, giám đốc công ty quyết định xử lý bằng cách phạt tiền. Nhờ sự giám sát chặt chẽ của quản đốc, giám đốc công ty biết được trong một tháng, giữa tỉ lệ công nhân vi phạm đúng $k$ lần ($1 \leq k \leq 2$) là $t_k = \dfrac{N_k}{N}$ (trong đó $N_k$ là số công nhân vi phạm đúng $k$ lần, $N$ là tổng số công nhân) và mức phạt mỗi lần vi phạm có mối liên hệ như sau: Nếu mỗi công nhân nộp phạt $x$ nghìn đồng ($60 \leq x \leq 300$) khi vi phạm lần thứ nhất và nộp phạt $x - 20$ nghìn đồng khi vi phạm lần thứ hai thì $t_1 = \dfrac{36}{x + 10}$ và $t_2 = \dfrac{4}{x - 30}$ (không có công nhân nào vi phạm quá hai lần). Biết rằng $N$ không đổi và bằng $2\,400$. Tổng số tiền nộp phạt trong một tháng ít nhất là bao nhiêu triệu đồng? (kết quả làm tròn đến hàng đơn vị)
\shortans[oly]{$103$}
\loigiai{
Tổng số công nhân $N = 2\,400$.\\
Số công nhân vi phạm đúng $1$ lần $N_1 = 2\,400 \cdot \dfrac{36}{x+10}$, tỉ lệ $t_1 = \dfrac{36}{x+10}$.\\
Số công nhân vi phạm đúng $2$ lần $N_2 = 2\,400 \cdot \dfrac{4}{x-30}$, tỉ lệ $t_2 = \dfrac{4}{x-30}$.\\
Không có công nhân nào vi phạm quá $2$ lần.\\
Tiền phạt mỗi lần vi phạm
\begin{itemize}
\item Lần 1: $x$ (nghìn đồng), với $60 \le x \le 300$.
\item Lần 2: $x - 20$ (nghìn đồng).
\end{itemize}
Tiền phạt công nhân vi phạm $1$ lần $x$ (nghìn đồng).\\
Tiền phạt công nhân vi phạm $2$ lần $x + (x - 20) = 2x - 20$ (nghìn đồng).\\
Tổng tiền phạt hàm số theo $x$ là
\[f(x) = N_1 \cdot x + N_2 \cdot (2x - 20) = \dfrac{86\,400x}{x+10} + \dfrac{9600(2x-20)}{x-30}.\]
Tính đạo hàm ta được $f'(x) = \dfrac{864\,000}{(x+10)^2} - \dfrac{384\,000}{(x-30)^2}$. Cho $f'(x) = 0 \Rightarrow \left(\dfrac{x+10}{x-30}\right)^2 = \dfrac{9}{4}$.\\
Từ đó tìm được nghiệm $x = 110$ (thỏa mãn điều kiện $x \in [60; 300]$).\\
Ta có $f(110) = 103\,200$, $f(60)\approx 106\,057$, $f(300)\approx 104\,235$.\\
Giá trị nhỏ nhất đạt được tại $x = 110$, khi đó $f(110) = 103\,200$ (nghìn đồng).\\
Vậy tổng số tiền phạt ít nhất làm tròn đến hàng đơn vị là $103$ triệu đồng.
}
\end{ex}

\begin{ex}%[2D1V3-6]%[Dự án đề thi thử NH25-26]%[THPT Chuyên Phan Bội Châu - Nghệ An]%[BC Tuan]
	Một cửa hàng kinh doanh máy lọc nước Karofi. Hàm cầu biểu thị liên hệ giữa giá bán của một chiếc máy với số lượng máy bán được trong một tháng là hàm bậc nhất. Khi giá bán là $5$ triệu đồng một chiếc thì một tháng bán được $100$ chiếc, khi giá bán là $4{,}5$ triệu đồng một chiếc thì một tháng bán được $120$ chiếc. Biết chi phí trung bình cho một chiếc máy khi bán được $x$ chiếc là $\dfrac{3x+50}{x}$. Hỏi cửa hàng cần bán với giá bao nhiêu triệu đồng một chiếc máy để lợi nhuận thu được trong một tháng lớn nhất?
	\shortans{$5{,}25$}
	\loigiai{
		Gọi hàm biểu thị mối liên hệ giữa giá bán của một chiếc máy với số lượng máy bán được là $p(x)=ax+b$.\\
		Tại $x=100$ thì $p(x)=5$ và tại $x=120$ thì $p(x)=4{,}5$.\\ Từ đó suy ra
		\[\heva{& 5=100a+b \\ & 4{,}5=120a+b}\Leftrightarrow \heva{& a=-0{,}025 \\ & b=7{,}5.}\]
		Vậy hàm số biểu thị giá bán là $p(x)=-0{,}025x+7{,}5$.\\
		Hàm số lợi nhuận là \[L(x) = x \cdot p(x) - \left( \dfrac{3x+50}{x} \right)x = x(-0{,}025x+7{,}5) - (3x+50) = -0{,}025x^2+4{,}5x-50.\]
		Ta có $L'(x)=-0{,}05x+4{,}5$.\\
		Khi đó $L'(x)=0\Leftrightarrow x=90$.\\
		\begin{center}
			\begin{tikzpicture}
				\tkzTabInit[nocadre=true, lgt=1.2, espcl=4, deltacl=0.5]
				{$x$ /0.7, $L'(x)$ /0.7, $L(x)$ /2}
				{$0$,$90$,$+\infty$}
				\tkzTabLine{ ,+,0,-, }
				\tkzTabVar{-/ $-50$, +/ $152{,}5$, -/ $-\infty$}
			\end{tikzpicture}
		\end{center}
		Bảng biến thiên cho thấy $L(x)$ đạt giá trị lớn nhất tại $x=90$.\\
		Vậy để lợi nhuận thu được lớn nhất thì cửa hàng cần bán một cái máy lọc nước với giá $p(90) = -0{,}025\cdot 90 + 7{,}5 = 5{,}25$ (triệu đồng).
	}
\end{ex}

\begin{ex}%[2D1V3-6]%[Dự án đề thi thử NH25-26]%[THPT Chuyên Phan Bội Châu - Nghệ An]%[Phạm Hoàng Đăng]
	\immini[thm]{Một hòn đảo $A$ cách bờ biển $3$ (km) và cách trạm điện $B$ là $5$ (km). Công ty điện lực dự định làm một tuyến đường điện từ trạm điện $B$ dọc theo bờ biển đến một điểm $C$, sau đó từ $C$ nối thẳng ra đảo $A$ (tham khảo hình vẽ). Biết chi phí cho $1$ (km) đường điện trên bờ là $30$ triệu đồng và trên biển là $50$ triệu đồng. Hỏi cần ít nhất bao nhiêu triệu đồng để hoàn thành kế hoạch trên?}{
		\begin{tikzpicture}[scale=0.8, line join=round, line cap=round, >=stealth, font=\footnotesize]
			\coordinate (H) at (0,0);
			\coordinate (A) at (0,3);
			\coordinate (B) at (4,0);
			\coordinate (C) at (2.25,0);
			\draw  (-1,0) -- (5,0) node[below] {Bờ biển};
			\draw[dashed] (A) -- (H) node[midway, left] {$3$ (km)};
			\draw (A) -- (C) node[midway,sloped, above right] {Biển};
			\draw[red] (B) -- (C);
			\draw[blue] (C) -- (A);
			\fill (A) circle (1pt) node[above] {$A$};
			\fill (B) circle (1pt) node[below] {$B$};
			\fill (C) circle (1pt) node[below] {$C$};
			\fill (H) circle (1pt) node[below] {$H$};
			\draw (H) rectangle +(0.2,0.2);
			\draw[dashed] (A) -- (B) node[midway, sloped, above] {$5$ (km)};
		\end{tikzpicture}
	}
	\shortans{$240$}
	\loigiai{
		Gọi $H$ là hình chiếu vuông góc của $A$ lên bờ biển.\\ Ta có $AH = 3$ (km).\\
		Tam giác $AHB$ vuông tại $H$ nên $HB = \sqrt{AB^2 - AH^2} = \sqrt{5^2 - 3^2} = 4$ (km).\\
		Đặt $CH = x$ (km) với $0 \le x \le 4$.\\ Suy ra quãng đường trên bờ là $BC = 4 - x$ (km).\\
		Khoảng cách từ điểm $C$ trên bờ nối thẳng ra đảo $A$ là $AC = \sqrt{CH^2 + AH^2} = \sqrt{x^2 + 9}$ (km).\\
		Hàm tổng chi phí làm đường điện là $f(x) = 30(4-x) + 50\sqrt{x^2+9}$ (triệu đồng).\\
		Ta có $f'(x) = -30 + \dfrac{50x}{\sqrt{x^2+9}}$.\\
		Khi đó 
		\begin{eqnarray*}
			&&f'(x) = 0 \\
			&\Leftrightarrow& 50x = 30\sqrt{x^2+9}\\ 
			&\Leftrightarrow& 25x^2 = 9(x^2+9)\\ 
			&\Leftrightarrow& 16x^2 = 81\\
			&\Leftrightarrow& x = \dfrac{9}{4} = 2{,}25.
		\end{eqnarray*}
		Bảng biến thiên
		\begin{center}
			\begin{tikzpicture}
				\tkzTabInit[nocadre=true, lgt=1.2, espcl=3, deltacl=0.5]
				{$x$ /0.8, $f'(x)$ /0.8, $f(x)$ /2.5}
				{$0$,$2{,}25$,$4$}
				\tkzTabLine{ ,-,0,+, }
				\tkzTabVar{+/ $270$, -/ $240$, +/ $250$}
			\end{tikzpicture}
		\end{center}
		Bảng biến thiên cho thấy hàm số $f(x)$ đạt giá trị nhỏ nhất tại $x = 2{,}25$.\\
		Chi phí ít nhất là $f(2{,}25) = 30(4 - 2{,}25) + 50\sqrt{2{,}25^2 + 9} = 52{,}5 + 187{,}5 = 240$ (triệu đồng).
	}
\end{ex}

\begin{ex}%[2D1V3-6]%[Dự án đề thi thử đợt 3-Nguyễn Đức Lợi]%[SGD tỉnh Lạng Sơn]%[GVPB - ThanhTa]
	Một nhà máy sản xuất $x$ sản phẩm trong mỗi tháng. Chi phí sản xuất $x$ sản phẩm được cho bởi hàm chi phí $C(x) = 16\,000 + 500x - 1{,}6x^2 + 0{,}004x^3$ (nghìn đồng). Biết giá bán của mỗi sản phẩm là một hàm số phụ thuộc vào số lượng sản phẩm $x$ và được cho bởi công thức \break$p(x) = 1\,700 - 7x$ (nghìn đồng). Hỏi mỗi tháng nhà máy nên sản xuất bao nhiêu sản phẩm để lợi nhuận thu được là lớn nhất? Biết rằng kết quả khảo sát thị trường cho thấy sản phẩm sản xuất ra sẽ được tiêu thụ hết.
	\par\shortans{100}
	\loigiai{
		Hàm doanh thu của nhà máy là
		$R(x)=x\cdot p(x)=x(1\,700-7x)=1\,700x-7x^2$ (nghìn đồng).\\
		Hàm lợi nhuận thu được là 
		\begin{align*}
			L(x)&=R(x)-C(x)\\
			&=\left(1\,700x-7x^2\right)-\left(16\,000+500x-1{,}6x^2+0{,}004x^3\right)\\
			&=-0{,}004x^3-5{,}4x^2+1\,200x-16\,000.
		\end{align*}
		Ta có $L'(x)=-0{,}012x^2-10{,}8x+1\,200$.\\
		Cho $L'(x)=0\Leftrightarrow-0{,}012x^2-10{,}8x+1\,200=0\Leftrightarrow\hoac{&x= 100\text{ (nhận)}\\&x=-1\,000\text{ (loại)}.}$\\
		Bảng biến thiên của hàm số $L(x)$ trên khoảng $(0;+\infty)$.
		\begin{center}
			\begin{tikzpicture}    
				\tkzTabInit[nocadre=false, lgt=1.2,deltacl=0.5 ,espcl=4]
				{$x$/0.7, $L'(x)$/0.7, $L(x)$/2}
				{$0$,$100$,$+\infty$}
				\tkzTabLine{,+,0,-,}
				\tkzTabVar{-/,+/$L(100)$,-/}
			\end{tikzpicture}
		\end{center}
		Dựa vào bảng biến thiên, hàm số $L(x)$ đạt giá trị lớn nhất tại $x=100$.\\
		Vậy mỗi tháng nhà máy nên sản xuất $100$ sản phẩm để lợi nhuận thu được là lớn nhất.
	}
\end{ex}

\begin{ex}%[2D1V3-6]%[Dự án C NH25-26 - Đợt 3 - Hector Tran]%[SGD-Tuyên Quang - Lần 1]%[GVPB: Thanh Phát]
	Một khu chung cư có $120$ căn hộ cho thuê. Người quản lí của khu chung cư nhận thấy rằng nếu giá thuê một căn hộ là $7$ triệu đồng một tháng thì tất cả các căn hộ đều sẽ có người thuê. Một cuộc khảo sát thị trường cho thấy, trung bình cứ mỗi lần tăng giá thuê một căn hộ mỗi tháng thêm $250$ nghìn đồng thì sẽ có thêm ba căn hộ bị bỏ trống. Người quản lí nên đặt giá thuê mỗi căn hộ là bao nhiêu triệu đồng một tháng để doanh thu một tháng là lớn nhất?
	\par\shortans{8,5}
	\loigiai{
		Gọi $x$ là số lần tăng giá thuê $250$\ nghìn đồng ($x \ge 0$).\\
		Khi đó, giá thuê mỗi căn hộ một tháng là $7 + 0{,}25x$ (triệu đồng).\\
		Số căn hộ có người thuê là $120 - 3x$ (căn hộ).\\
		Doanh thu một tháng của khu chung cư là
		\begin{eqnarray*}
			R(x) &=& (7 + 0{,}25x)(120 - 3x)\\
			&=& 840 - 21x + 30x - 0{,}75x^2\\
			&=& -0{,}75x^2 + 9x + 840.
		\end{eqnarray*}
		Đây là một hàm số bậc hai với hệ số $a = -0{,}75 < 0$, nên hàm số đạt giá trị lớn nhất tại hoành độ đỉnh parabol
		\[x = -\dfrac{b}{2a} = -\dfrac{9}{2 \cdot (-0{,}75)} = 6.\]
		Khi đó, giá thuê mỗi căn hộ để doanh thu lớn nhất là $7 + 0{,}25 \cdot 6 = 8{,}5$ (triệu đồng).
	}
\end{ex}

\begin{ex}%[TT-THPT-HamRong-NH25-26 ]%[GV soạn: Nguyễn Quang Hiệp GVPB2 Đỗ Minh Vũ ]%[2D1V3-6]
	Trong một thí nghiệm y học, người ta cấy $1\,000$ vi khuẩn vào môi trường dinh dưỡng. Bằng thực nghiệm, người ta xác định được số lượng vi khuẩn thay đổi theo thời gian bởi công thức $N(t)=1\,000+\dfrac{100t}{100+t^{2}}$ (con), trong đó $t$ là thời gian tính bằng giây ($t \ge 0$). Tính số lượng vi khuẩn lớn nhất kể từ khi thực hiện cấy vi khuẩn vào môi trường dinh dưỡng.
	\par\shortans[oly]{1005}
	\loigiai{
		Ta có đạo hàm:
		\[N'(t) = \dfrac{100\left(100+t^2\right) - 100t \cdot 2t}{\left(100+t^2\right)^2} = \dfrac{10\,000 - 100t^2}{\left(100+t^2\right)^2}.\]
		Xét $N'(t) = 0 \Leftrightarrow 10\,000 - 100t^2 = 0 \Leftrightarrow t = 10$ (vì $t \ge 0$).\\
		Ta có bảng biến thiên:
		\begin{center}
			\begin{tikzpicture}
				\tkzTabInit[nocadre=true, lgt=1.2, espcl=3.5, deltacl=0.5]
				{$t$/0.7, $N'(t)$/0.7, $N(t)$/2}
				{$0$, $10$, $+\infty$}
				
				\tkzTabLine{,+,0,-,}
				\tkzTabVar{-/$1\,000$, +/$1\,005$, -/$1\,000$}
			\end{tikzpicture}
		\end{center}
		Từ bảng biến thiên suy ra số lượng vi khuẩn lớn nhất kể từ khi thực hiện cấy vi khuẩn vào môi trường dinh dưỡng là $1\,005$ con.
	}
\end{ex}

\begin{ex}%[TT-THPT-HamRong-NH25-26 ]%[GV soạn: Nguyễn Quang Hiệp GVPB2 Đỗ Minh Vũ ]%[2D1V3-6]
	Một công ty công nghệ sản xuất linh kiện điện tử. Chi phí sản xuất $x$ linh kiện được tính theo công thức $C(x) = 0{,}5x^{2} + 200x + 5\,000$ (USD). Khách hàng ký hợp đồng bao tiêu sản phẩm với mức giá như sau:
	\begin{itemize}
		\item Với $500$ sản phẩm đầu tiên được mua với giá $1\,000$ USD/sản phẩm.
		\item Từ sản phẩm thứ $501$ trở đi, giá thu mua là $800$ USD/sản phẩm.
	\end{itemize}
	Để điều tiết sản xuất, Nhà nước áp đặt một mức thuế phụ thu là $t$ (USD) trên mỗi đơn vị sản phẩm, nhưng chỉ áp dụng cho những sản phẩm từ thứ $501$ trở đi và khi thu thuế thì mong muốn số tiền phụ thu lớn nhất. Với cách tính thuế và yêu cầu như vậy, doanh nghiệp nên sản xuất bao nhiêu linh kiện để lợi nhuận thu được lớn nhất?
	\par\shortans[oly]{550}
	\loigiai{
		Trường hợp 1: Công ty sản xuất không quá $500$ linh kiện, tức là $1 \le x \le 500$.\\
		Lợi nhuận là $L(x) = \text{Doanh thu} - \text{Chi phí sản xuất}$.\\
		Khi đó, $L(x) = 1\,000x - C(x) = 1\,000x - \left(0{,}5x^{2} + 200x + 5\,000\right) = -0{,}5x^{2} + 800x - 5\,000$.\\
		Ta có $L'(x) = -x + 800 = 0 \Leftrightarrow x = 800$ (không thỏa mãn điều kiện $1 \le x \le 500$).\\
		Mặt khác, $L(x)$ đồng biến trên khoảng $(1; 500)$ nên khi $x = 500$ thì $L(x)$ đạt giá trị lớn nhất là $L(500) = 270\,000$ (USD).\\
		Trường hợp 2: Công ty sản xuất nhiều hơn $500$ linh kiện, tức là $x > 500$.\\
		Khi đó, số sản phẩm từ thứ $501$ trở đi là $x - 500$ (sản phẩm).\\
		Doanh thu là $1\,000 \cdot 500 + 800(x - 500)$ (USD).\\
		Chi phí sản xuất là $C(x) = 0{,}5x^{2} + 200x + 5\,000$ (USD).\\
		Số tiền thuế Nhà nước phụ thu là $T(x) = t(x - 500)$ (USD).\\
		Lợi nhuận là:
		\begin{eqnarray*}
			L(x) &=& 1\,000 \cdot 500 + 800(x - 500) - \left(0{,}5x^{2} + 200x + 5\,000\right) - t(x - 500) \\
			&=& -0{,}5x^{2} + (600 - t)x + (95\,000 + 500t) \text{ (USD)}.
		\end{eqnarray*}
		Vì $L(x)$ là hàm số bậc hai có hệ số $a = -0{,}5 < 0$ nên đồ thị là parabol có bề lõm quay xuống dưới và $L(x)$ đạt giá trị lớn nhất tại $x = \dfrac{-b}{2a} = \dfrac{-(600 - t)}{2 \cdot (-0{,}5)} = 600 - t$.\\
		Như vậy, ứng với mỗi mức thuế $t$ (USD/sản phẩm), doanh nghiệp sẽ sản xuất tại điểm mà lợi nhuận thu được lớn nhất ứng với mức thuế $t$ đó.\\
		Khi đó, số linh kiện sản xuất là $x = 600 - t$ (sản phẩm).\\
		Số tiền thuế thu được là $T(t) = t(600 - t - 500) = -t^{2} + 100t$ (USD).\\
		Để số tiền thuế thu được là lớn nhất thì hàm số $T(t) = -t^{2} + 100t$ đạt giá trị lớn nhất tại $t = \dfrac{-100}{2 \cdot (-1)} = 50$ (USD).\\
		Số linh kiện doanh nghiệp sản xuất là $x = 600 - 50 = 550$ (sản phẩm).\\
		Lợi nhuận khi đó là $L = -0{,}5 \cdot 550^{2} + (600 - 50) \cdot 550 + (95\,000 + 500 \cdot 50) = 271\,250$ (USD).\\
		So sánh với trường hợp 1, ta thấy $271\,250 > 270\,000$.\\
		Vậy doanh nghiệp nên sản xuất $550$ linh kiện để lợi nhuận thu được là lớn nhất.
	}
\end{ex}

\begin{ex}%[2D1V3-6]%[KNTT - Lớp 12 - Dự án C - Đợt 2]%[Loc do]%[PB: Phan Anh]
	Một công ty công nghệ cung cấp gói lưu trữ dữ liệu doanh nghiệp với giá niêm yết $200\,000$ đồng/tháng và đang có $400$ khách hàng. Phòng kinh doanh xác định được quy luật: Cứ giảm giá $10\,000$ đồng thì sẽ có thêm $50$ khách hàng mới. Tuy nhiên, do hiện tượng quá tải băng thông, chi phí vận hành $C(n)$ (nghìn đồng) không cố định mà biến thiên theo hàm bậc hai của số lượng khách hàng $n$, được xác định bởi công thức $C(n)=28n+0,01n^2 +15\, 000$. Công ty cần chốt giá bán bao nhiêu nghìn đồng để lợi nhuận đạt mức tối đa?
	\par\shortans[0]{$160$}
	\loigiai{
		Gọi $x$ là số lần giảm giá $10\,000$ đồng $(0<x<20)$.\\
		Giá bán mới $P=200-10x$ (nghìn đồng).\\
		Số lượng khách hàng khi đó là $n=400+50x\Rightarrow x=\dfrac{n-400}{50} $.\\
		Khi đó $P=200-10x=200-10\cdot\dfrac{n-400}{50} =-\dfrac{n}{5} +280$.\\
		Ta có hàm lợi nhuận là
		\begin{align*}
			L=& n\cdot P(n)-C(n)\\
			=&\, n\left(-\dfrac{n}{5} +280\right)-(28n+0{,}01n^2 +15\, 000)\\
			=&\,-\dfrac{21}{100} n^2 +252n-15\,000.
		\end{align*}		
		Suy ra  $L'=-\dfrac{21}{50} n+252$; $L'=0\Rightarrow -\dfrac{21}{50} n+252=0\Rightarrow n=600$.\\
		Khi đó giá bán ra là $P=-\dfrac{600}{5} +280=160$ (nghìn đồng).\\
		Vậy công ty cần chốt giá bán $160$ nghìn đồng để lợi nhuận đạt mức tối đa.}
\end{ex}

\begin{ex}%[2D1V3-6] %[Dự án C đề thi thử TN Môn Toán NH25-26 - Dot 4 - Vũ Hồng Toàn]%[THPT Mai Anh Tuấn - Thanh Hóa]%[GVPB - Lê Hữu Kiệt]
	Nhà thầy Quang cách bờ biển $1$\,km. Mỗi buổi sáng thầy chạy bộ từ nhà ra bờ biển sau đó chạy dọc bờ biển $500$\,m, rồi thầy chạy qua chợ hải sản để lấy thức ăn trong ngày, cuối cùng thầy chạy về nhà. Biết chợ hải sản cách bờ biển $400$\,m và cách nhà thầy Quang $1$\,km. Quãng đường ngắn nhất mà thầy Quang đã chạy trong mỗi buổi sáng là bao nhiêu mét {\it (làm tròn đến hàng đơn vị)}?
	\begin{center}
		\begin{tikzpicture}[scale=1, font=\footnotesize, line join=round, line cap=round, >=stealth]
			\tikzset{declare function={a=4;b1=0.4*a;b2=b1+1;goc=40;}}
			\path (0,0)coordinate (H) (0,b1+b2)coordinate (A)node[shift={(90:7pt)}]{Nhà}
			(b1/2,0)coordinate (B)++(b2,0)coordinate (C)++(b1,0)coordinate (K)++(0,b2)coordinate (D)node[shift={(90:7pt)}]{Chợ};
			\fill[cyan!20] (-1,0) rectangle (6, -1);
			\draw (-1,0) -- (6,0); % Đường mặt đất
			\draw [dashed] (A)--(H) (D)--(K);
			\foreach \x/\y in {A/B,B/C, C/D,D/A}{\draw[red, thick, ->] (\x)--(\y);}
			\path (B)--(C)node[midway,shift={(-90:7pt)}]{Biển};
		\end{tikzpicture}
	\end{center}
	\par\shortans{2932}
	\loigiai{
		\begin{center}
			\begin{tikzpicture}[scale=1, font=\footnotesize, line join=round, line cap=round, >=stealth]
				\tikzset{declare function={a=4;b1=0.4*a;b2=b1+1;goc=40;}}
				\path (0,0)coordinate (H) (0,b1+b2)coordinate (A)node[shift={(90:7pt)}]{Nhà}
				(0,b2)coordinate(I)
				(b1/2,0)coordinate (B)++(b2,0)coordinate (C)++(b1,0)coordinate (K)++(0,b2)coordinate (D)node[shift={(90:7pt)}]{Chợ};
				\fill[cyan!20] (-1,0) rectangle (6, -1);
				\draw (-1,0) -- (6,0); % Đường mặt đất
				\draw [dashed] (A)--(H) (K)--(D)--(I);
				\foreach \x/\y/\z in {A/H/B,C/K/D,D/I/H} \draw pic[draw,angle radius=2mm]{right angle= \x--\y--\z};
				\foreach \x/\y in {A/B,B/C, C/D,D/A}{\draw[red, thick, ->] (\x)--(\y);}
				\path (B)--(C)node[midway,shift={(-90:7pt)}]{Biển};
				\foreach \x /\gN in {A/200,B/-90,C/-90,K/-90,H/-90,D/-20,I/180}
				\fill (\x) circle (1pt)($(\x)+(\gN:2.5mm)$) node {$\x$};
			\end{tikzpicture}
		\end{center}
		Đổi $500m = 0{,}5$\,km và $400m = 0{,}4$\,km.\\
		Đặt tên các điểm như hình vẽ.\\
		Ta có $AI = AH - DK = 1-0{,}4 = 0{,}6$.\\
		$HK = ID = \sqrt{AD^2 - AI^2} = \sqrt{1^2 -0{,}6^2} = 0{,}8$ và $BC = 0{,}5$.\\
		Đặt $HB = x$ (km) suy ra $CK = HK - HC = 0{,}3-x$ điều kiện $0 \le x \le 0{,}3$.\\
		Quãng đường mà thầy Quang chạy mỗi sáng là
		\allowdisplaybreaks
		\begin{eqnarray*}
			AB + BC + CD + DA &=& \sqrt{x^2 +1} +1{,}5 + \sqrt{(0{,}3-x)^2 +0{,}4^2} \\&=& \sqrt{x^2 +1} +1{,}5+ \sqrt{x^2-0{,}6x+0{,}25}.
		\end{eqnarray*}
		Xét hàm $f (x) = \sqrt{x^2 +1} + \sqrt{x^2 - 0{,}6x+0{,}25} +1{,}5$ với $x \in [0;0{,}3]$.\\
		Ta có $f'(x) = \dfrac{x}{\sqrt{x^2+1}} + \dfrac{x-0{,}3}{\sqrt{x^2-0{,}6x+0{,}25}}$.    \\
		Cho $f'(x)=0 \Leftrightarrow x=\dfrac{3}{14}=x_0$.\\
		Bảng biến thiên
		\begin{center}
			\begin{tikzpicture}
				\tkzTabInit[lgt=1.2,espcl=5,deltacl=0.5]
				{$x$/0.7,$f'(x)$/0.7,$f(x)$/2}
				{$-\infty$,$x_0$,$+\infty$}
				\tkzTabLine{, -,0,+,}
				\tkzTabVar{+/,-/$f(x_0)$,+/}
				\tkzTabVal[draw]{1}{2}{0.5}{$0$}{$f(0)$}
				\tkzTabVal[draw]{2}{3}{0.5}{$0{,}3$}{$f(0{,}3)$}%
				\fill[pattern = north west lines,opacity=0.3]
				(T11)rectangle (M13) (M21)rectangle (T23);
			\end{tikzpicture}
		\end{center}
		Từ bảng biến thiên ta có $\min\limits_{[0;0{,}3]} f(x) = f\left(\dfrac{3}{14}\right) \approx 2{,}932$.\\
		Vậy, quãng đường ngắn nhất mà thầy Quang chạy mỗi sáng là $2{,}932$\,km hay $2\,932$\,(mét).
	}
\end{ex}

\begin{ex}%[2D1V3-6]%[Dự án C đề thi thử TN Môn Toán NH25-26 - Dot 4 - Phạm Hoàng Đăng]%[THPT Nguyễn Trãi-Hà Nội]%[GVPB - Khắc Thiên]
	Một nông trại dâu tây sạch tại Đà Lạt (Bên A) ký hợp đồng cung cấp độc quyền cho một chuỗi cửa hàng tại TP.HCM (Bên B). Dựa trên dữ liệu thị trường, Bên B dự tính rằng nếu mỗi ngày nhập và bán hết $x$ (kg) dâu tây ($0<x< 150$), thì tổng doanh thu bán lẻ là $R(x)=300x-x^2$ (nghìn đồng). Tuy nhiên, để vận hành việc bán hàng (bảo quản, nhân sự, mặt bằng,\ldots), bên B phải chịu chi phí vận hành (không bao gồm tiền nhập hàng) là $C(x)=x^2+20x+500$ (nghìn đồng). Bên A sẽ quyết định giá bán sỉ là $p$ cho Bên B. Giả sử Bên B là đơn vị kinh doanh tối ưu, họ sẽ luôn chọn phương án nhập hàng là $x$ sao cho lợi nhuận thu được là cao nhất ứng với mức giá $p$. Hãy xác định mức giá sỉ $p$ (nghìn đồng) mà nông trại nên thiết lập để tổng doanh thu của nông trại từ việc bán hàng cho Bên B là lớn nhất (\textit{làm tròn kết quả đến hàng đơn vị}).
	\par
	\shortans[oly]{140}
	\loigiai{
		Với Bên B (Cửa hàng) ta có\\
		\begin{itemize}
			\item Doanh thu $R(x)=300x-x^2$ (nghìn đồng).
			\item Chi phí vận hành $C(x)=x^2+20x+500$ (nghìn đồng).
			\item Chi phí nhập hàng $p\cdot x$ (nghìn đồng).
			\item Lợi nhuận Bên B thu được là
			\[L_B(x)=R(x)-C(x)-px=-2x^2+(280-p)x-500.\]
			Ta có $L_B'(x)=-4x+280-p\ (0<x<150)$.\\
			Khi đó $L_B'(x)=0 \Leftrightarrow x=x_0=\dfrac{280-p}{4}$.
			\begin{center}
				\begin{tikzpicture}
					\tkzTabInit[lgt=1.5,espcl=3]
					{$x$ /0.8,$L'_B$ / 0.8,$L_B$ / 2}
					{$0$,$x_0$,$150$}
					\tkzTabLine{,+,0,-}
					\tkzTabVar{-/ $ $,+/ $ $,-/ $ $}
				\end{tikzpicture}
			\end{center}
			Do đó $L_B(x)$ đạt giá trị lớn nhất trên $(0;150)$ tại $x=\dfrac{280-p}{4} \Rightarrow p=280-4x$.
		\end{itemize}
		Với Bên A (Nông trại) ta có doanh thu của Bên A là
		\[L_A(x)=px=(280-4x) \cdot x=-4x^2+280x.\]
		Suy ra $L_A'(x)=-8x+280$.\\ 
		Khi đó $L_A'(x)=0 \Leftrightarrow x=35$.\\
		Bảng biến thiên
		\begin{center}
			\begin{tikzpicture}
				\tkzTabInit[lgt=1.5,espcl=3]
				{$x$ /0.8,$L'_A$ / 0.8,$L_A$ / 2}
				{$0$,$35$,$150$}
				\tkzTabLine{,+,0,-}
				\tkzTabVar{-/ $ $,+/ $ $,-/ $ $}
			\end{tikzpicture}
		\end{center}
		Do đó $L_A(x)$ đạt giá trị lớn nhất trên $(0; 150)$ tại $x=35$.\\
		Suy ra $p=280-4\cdot 35=140$ (nghìn đồng).
	}
\end{ex}

\begin{ex}%[2D1V4-1]%[Dự án C đề thi thử TN Môn Toán NH25-26 - Dot 4 - Phạm Hoàng Đăng]%[THPT Nguyễn Trãi-Hà Nội]%[GVPB - Khắc Thiên]
	Cho hàm số $y=g(x)=\ln\dfrac{(x-2)^2(x+5)}{(x-4)(x+1)}$. Hỏi đồ thị hàm số có bao nhiêu đường tiệm cận (chỉ xét tiệm cận đứng và tiệm cận ngang)?
	\par
	\shortans[oly]{3}
	\loigiai{
		Điều kiện xác định $\dfrac{(x-2)^2(x+5)}{(x-4)(x+1)}>0\Leftrightarrow\hoac{&-5<x<-1\\&x>4}$.\\
		Vậy tập xác định $\mathscr{D}=(-5;-1)\cup(4;+\infty)$.
		\begin{itemize}
			\item Trước hết ta tìm tiệm cận đứng.\\
			Ta biết với $a>1$ thì $\lim\limits_{x\rightarrow 0^+}\log_a x=-\infty$; $\lim\limits_{x \rightarrow+\infty} \log_a x=+\infty$.\\
			Đặt $f(x)=\dfrac{(x-2)^2(x+5)}{(x-4)(x+1)}$.\\ Ta xét các giá trị của $x$ để $f(x) \rightarrow 0^{+}$ và $f(x) \rightarrow+\infty$.
			\begin{itemize}
				\item Xét $x \rightarrow-5^{+} \Rightarrow f(x) \rightarrow 0^{+}$.\\ Do đó $\lim\limits_{x \rightarrow-5^{+}} \ln f(x)=-\infty$.\\
				Vậy $x=-5$ là một đường tiệm cận đứng của đồ thị $y=g(x)=\ln \dfrac{(x-2)^2(x+5)}{(x-4)(x+1)}$.
				\item Xét $x \rightarrow-1^{-} \Rightarrow f(x) \rightarrow+\infty$.\\ Do đó $\lim\limits_{x \rightarrow-1^{-}} \ln f(x)=+\infty$.\\
				Vậy $x=-1$ là một đường tiệm cận đứng của đồ thị $y=g(x)=\ln \dfrac{(x-2)^2(x+5)}{(x-4)(x+1)}$.
				\item Xét $x \rightarrow 4^{+} \Rightarrow f(x) \rightarrow+\infty$.\\ Do đó $\lim\limits_{x \rightarrow 4^{+}} \ln f(x)=+\infty$.\\
				Vậy $x=4$ là một đường tiệm cận đứng của đồ thị $y=g(x)=\ln \dfrac{(x-2)^2(x+5)}{(x-4)(x+1)}$.
			\end{itemize}
			Vậy đồ thị đã cho có $3$ đường tiệm cận đứng.
			\item Ta có $\lim\limits_{x \rightarrow+\infty }f(x)=+\infty$ nên $\lim\limits_{x \rightarrow+\infty} \ln f(x)=+\infty$.\\
			Vậy đồ thị hàm số đã cho không có tiệm cận ngang.
		\end{itemize}
		Vậy đồ thị hàm số có $3$ đường tiệm cận.
	}
\end{ex}

\begin{ex}%[2D1V5-8]%[TT Cụm 5 - Ninh Bình-NH25-26]%[GVSB:Dương Công Tạo]%[GVPB1: Nguyễn Tấn Tài]
	\immini[thm]{Trong hệ trục tọa độ $Oxy$, đơn vị mỗi trục là mét, một đường trượt mới sẽ được xây dựng theo bản thiết kế đã trình bày như hình vẽ. Thanh trượt bắt đầu từ $A$ và kết thúc tại $C$, đường cong của thanh trượt là một phần của đồ thị hàm số $f(x) = \dfrac{ax^2 + bx + c}{x+d}$, biết đồ thị hàm số $f(x)$ tiếp xúc với trục $Ox$ tại điểm $B$. Bạn Nam bắt đầu trượt từ điểm $A$, hỏi khi Nam cách vị trí ban đầu theo phương ngang một khoảng $5$ mét thì Nam cách mặt đất bao nhiêu mét, biết trục $Ox$ nằm trên mặt đất (\textit{kết quả làm tròn đến hàng phần trăm}).}
	{\begin{tikzpicture}[scale=.5,>=stealth, line join=round, line cap=round, x=0.7cm, y=0.7cm, every node/.style={font=\footnotesize},
			declare function={
				hA = 10;
				xB = 10;
				dxC = 5;
				xC = xB + dxC;
				yC = 2.5; 
				a = hA/(xB*xB);
				f(\x) = a*(\x-xB)*(\x-xB);
				wWall = 2.5;
				hTop = 1.2;
				xSlantTop = -wWall - 0.8;
				xSlantBot = -wWall - 1.5;
				hBase = 1.5;
				xBaseLeft = xSlantBot - 0.2;
			}]
			
			\fill[black!85] (xBaseLeft, -hBase) rectangle (0, 0);
			
			\fill[pattern=north west lines, pattern color=gray!70] (xSlantBot, 0) -- (xSlantTop, hA) -- (-wWall, hA) -- (-wWall, 0) -- cycle;
			\draw (xSlantBot, 0) -- (xSlantTop, hA) -- (-wWall, hA) -- (-wWall, 0) -- cycle;
			
			\fill[gray!30] (-wWall, 0) rectangle (0, hA - hTop);
			
			\fill[pattern=north west lines, pattern color=gray!70] (-wWall, hA - hTop) rectangle (0, hA);
			
			\draw (-wWall, 0) rectangle (0, hA);
			\draw (-wWall, hA - hTop) -- (0, hA - hTop);
			
			\draw[->] (0, 0) -- (xC + 2, 0) node[below] {$x$};
			\draw[->] (0, 0) -- (0, hA + 2) node[left] {$y$};
			
			\draw[smooth, samples=200, domain=0:xC] plot (\x, {f(\x)});
			
			\fill (0, 0) circle (1pt) node[shift={(45:10pt)}] {$O$};
			\fill (0, hA) circle (1pt) node[right=4pt] {$A$};
			\fill (xB, 0) circle (1pt) node[above=6pt] {$B$};
			\fill (xC, yC) circle (1pt) node[above=4pt] {$C$};
			
			\draw[<->] (-wWall/2, 0) -- (-wWall/2, hA) node[midway, fill=gray!30, inner sep=2pt] {$10$ m};
			
			\draw (xB, 0) -- (xB, -1.3);
			\draw (xC, 0) -- (xC, -1.3);
			
			\draw[<->] (0, -1) -- (xB, -1) node[midway, fill=white, inner sep=2pt] {$10$ m};
			\draw[<->] (xB, -1) -- (xC, -1) node[midway, fill=white, inner sep=2pt] {$5$ m};
			
			\draw[<->] (xC, 0) -- (xC, yC) node[midway, right] {$1$ m};
			
	\end{tikzpicture}}
	\par  \shortans{1,67}
	\loigiai{
		Ta có đồ thị hàm số $f(x)$ tiếp xúc với trục $Ox$ tại điểm $B(10;0)$ nên $f(x) = \dfrac{a(x-10)^2}{x+d}$. \\
		Đồ thị đi qua $A(0;10) \Rightarrow 10 = \dfrac{a \cdot 100}{d} \Rightarrow d = 10a$. \\
		Đồ thị đi qua $C(15;1) \Rightarrow 1 = \dfrac{a \cdot 25}{15+d} \Rightarrow 15+d = 25a \Rightarrow 15+10a = 25a \Rightarrow a = 1$; $d =10$. \\
		Do đó $f(x) = \dfrac{(x-10)^2}{x+10}$. \\
		Có $x = 5 \Rightarrow y = \dfrac{25}{15} \approx 1{,}67$. \\
		Vậy Nam cách vị trí ban đầu theo phương ngang một khoảng $5$ mét thì Nam cách mặt đất $1{,}67$ (m).
	}
\end{ex}

\begin{ex}%[2D1V5-8]%[Dự án C - Đợt 3 - Sở Giáo Dục Sơn La]%[Lê Hùng Thắng - Nguyễn Tiến]%%[Thy Nguyen Vo Diem - Nguyễn Tiến]
	Một nhà địa chất học đang ở địa điểm $A$ trên sa mạc. Anh ta muốn đến điểm $B$ (cũng ở trên sa mạc) và cách $A$ một khoảng bằng $70$ km. Trong sa mạc, xe anh ta chỉ có thể di chuyển với vận tốc $30$ km/h. Nhà địa chất phải đến được điểm $B$ sau $2$ giờ, vì vậy nếu anh ta đi thẳng từ $A$ đến $B$ sẽ không thể đến đúng giờ được. Rất may, có một con đường nhựa song song với đường nối $A$ và $B$ và cách $AB$ một đoạn $10$ km (tham khảo hình vẽ). 
	\begin{center}
		\begin{tikzpicture}[scale=1, font=\footnotesize, line join=round, line cap=round, >=stealth]
			\coordinate (A1) at (0,0);
			\coordinate (A) at (0,2);
			\coordinate (B) at (7,2);
			\coordinate (B1) at (7,0);
			\draw[thick] (-1,0) -- (9,0) node[below] {Đường nhựa};
			\draw[dashed] (A)--(A1) node[midway, left] {$10$ km};
			\draw[dashed] (B)--(B1) node[midway, right] {$10$ km};
			\draw[<->] (0,2.5) -- (7,2.5) node[midway, above] {$70$ km};
			\draw[fill=black] (A) circle (1pt) node[above] {$A$};
			\draw[fill=black] (B) circle (1pt) node[above] {$B$};
			\draw[blue, thick] (A) -- (1,0) -- (6,0) -- (B);
		\end{tikzpicture}
	\end{center}
	Trên đường nhựa đó, xe nhà địa chất có thể di chuyển với vận tốc $50$ km/h. Thời gian ngắn nhất để nhà địa chất di chuyển từ $A$ đến $B$ là bao nhiêu phút? 
	\shortans{116}
	\loigiai{
		Giả sử lộ trình của nhà địa chất là từ $A$ đi đến một điểm $M$ trên đường nhựa, sau đó di chuyển trên đường nhựa đến điểm $N$, rồi từ $N$ đi đến $B$.\\
		Gọi $A'$ và $B'$ lần lượt là hình chiếu của $A$ và $B$ trên đường nhựa. 
		\begin{center}
			\begin{tikzpicture}[scale=1, font=\footnotesize, line join=round, line cap=round, >=stealth]
				\coordinate (A1) at (0,0);
				\coordinate (A) at (0,2);
				\coordinate (B) at (7,2);
				\coordinate (B1) at (7,0);
				\draw[thick] (-1,0) -- (9,0) node[below] {Đường nhựa};
				\draw[dashed] (A)--(A1) node[midway, left] {$10$ km};
				\draw[dashed] (B)--(B1) node[midway, right] {$10$ km};
				\draw[<->] (0,2.5) -- (7,2.5) node[midway, above] {$70$ km};
				\draw[fill=black] (A) circle (1pt) node[above] {$A$};
				\draw[fill=black] (B) circle (1pt) node[above] {$B$};
				\draw[blue, thick] (A) -- (1,0) -- (6,0) -- (B);
				\node[below] at (1,0) {$M$};
				\node[below] at (6,0) {$N$};
				\node[below] at (0,0) {$A'$};
				\node[below] at (7,0) {$B'$};
			\end{tikzpicture}
		\end{center}
		Khi đó $A'B' = AB = 70$ km và $AA' = BB' = 10$ km. \\
		Đặt $A'M = x$ (km) với $0 \le x \le 35$.\\
		Do tính đối xứng của bài toán, để thời gian là ngắn nhất thì $B'N = x$.
		\\Khi đó, quãng đường di chuyển trong sa mạc là 
		$$AM = NB = \sqrt{x^2 + 10^2} = \sqrt{x^2 + 100} \text{(km).}$$
		Quãng đường di chuyển trên đường nhựa là $MN = 70 - 2x$ (km).	\\
		Tổng thời gian di chuyển là:
		$$T(x) = \dfrac{2\sqrt{x^2 + 100}}{30} + \dfrac{70 - 2x}{50} = \dfrac{\sqrt{x^2 + 100}}{15} + \dfrac{35 - x}{25} \text{(giờ)}.$$ 
		Xét hàm số $T(x)$ liên tục trên $[0; 35]$.\\
		$T'(x) = \dfrac{x}{15\sqrt{x^2 + 100}} - \dfrac{1}{25}$. \\
		\begin{eqnarray*}
			& & T'(x) = 0 \Leftrightarrow 25x = 15\sqrt{x^2 + 100} \\
			& \Leftrightarrow  & 5x = 3\sqrt{x^2 + 100} \\
			& \Leftrightarrow  & x^2 = \dfrac{225}{4} \Leftrightarrow x = 7{,}5.
		\end{eqnarray*}
		Ta có $T(7{,}5) = \dfrac{\sqrt{7{,}5^2 + 100}}{15} + \dfrac{35 - 7{,}5}{25} = \dfrac{29}{15}$ (giờ).\\
		Vậy thời gian ngắn nhất để nhà địa chất di chuyển từ $A$ đến $B$ là $\dfrac{29}{15}$ giờ = $116$ phút.
		\textbf{Cách khác:}\\
		Gọi $H$, $K$, $C$, $D$ là các điểm như hình vẽ.
		\begin{center}
			\begin{tikzpicture}[scale=1, font=\footnotesize, line join=round, line cap=round, >=stealth]
				\path (-1,0) coordinate (C) 
				(3,0) coordinate (D)
				(-3,2) coordinate (A) 
				(4.5,2) coordinate (B)
				(-3,0) coordinate (H)
				(4.5,0) coordinate (K)
				;
				\draw (-5,0)--(5,0) (-5,2)--(5,2) (A)--(C) (B)--(D) (A)--(H) node[below]{$H$} (B)--(K) node[below]{$K$};
				\fill (C) circle (1pt) (D) circle(1pt) (A) circle (1pt) (B) circle (1pt) (H) circle (1pt) (K) circle (1pt);
				\draw (C) node[below]{$C$} --(D)node[below]{$D$} node[midway,below]{Đường nhựa};
				\draw (A) node[above]{$A$}--(B)node[above]{$B$} node[midway,above]{Sa mạc} node[midway,below]{$70$ km}
				;
				\draw[<->] (-4.5,0)--(-4.5,2) node[midway, right]{$10$ km};
				\draw pic[draw,angle radius=2mm]{right angle=B--K--D};
				\draw pic[draw,angle radius=2mm]{right angle=C--H--A};
			\end{tikzpicture}
		\end{center}
		Ta chia quãng đường đi thành ba giai đoạn 
		\begin{itemize}
			\item\textbf{ Giai đoạn 1:} đi từ $A$ đến $C$ (từ sa mạc đến đường nhựa song song).
			\item \textbf{Giai đoạn 2:} đi từ $C$ đến $D$ (một quãng đường nào đó trên đường nhựa).
			\item \textbf{Giai đoạn 3:} đi từ $D$ đến $B$ (từ điểm kết thúc $D$ trên đường nhựa đi tiếp đến $B$ băng qua sa mạc).
		\end{itemize}
		Đặt $HC=x$ ($0<x<70$) và $DK=y$ ($0<y<70$).\\
		Ta có
		\begin{itemize}
			\item $AC=\sqrt{10^2+x^2} \Rightarrow t_{AC}=\dfrac{AC}{v_{AC}}=\dfrac{\sqrt{100+x^2}}{30}$;
			\item $DB=\sqrt{10^2+y^2} \Rightarrow t_{DB}=\dfrac{DB}{v_{DB}}=\dfrac{\sqrt{100+y^2}}{30}$;
			\item $CD=70-(x+y) \Rightarrow t_{CD}=\dfrac{CD}{v_{CD}}=\dfrac{70-x-y}{50}$.
		\end{itemize}
		Tổng thời gian nhà địa chất học đi từ $A$ đến $B$ là \[T(x;y)=\dfrac{\sqrt{100+x^2}}{30}+\dfrac{\sqrt{100+y^2}}{30}+\dfrac{70-x-y}{50}.\]
		Đây là một biểu thức có dạng đối xứng $2$ biến $x$, $y$ và ta cần tìm $\min T(x;y)$.\\
		Ta có $$T(x;y)=\dfrac{100+x^2}{30}+\dfrac{35-x}{50}+\dfrac{100+y^2}{30}+\dfrac{35-y}{50}=f(x)+f(y).$$
		Xét $f(u)=\dfrac{\sqrt{100+u^2}}{30}+\dfrac{35-u}{50}$, $0<u<70$.\\
		Suy ra $f'(u)=\dfrac{u}{30\sqrt{100+u^2}}-\dfrac{1}{50}$.\\
		Cho $f'(u)=0 \Leftrightarrow \sqrt{100+u^2}=\dfrac{5u}{3}>0 \Leftrightarrow u=\dfrac{15}{2}$.\\
		Bảng biến thiên
		\begin{center}
			\begin{tikzpicture}
				\tkzTabInit[nocadre=true, lgt=1.2, espcl=4, deltacl=0.5]
				{$u$/0.7,$f’(u)$/0.7,$f(u)$/2}
				{$0$,$7{,}5$,$70$}
				\tkzTabLine{ ,-,0,+, }
				\tkzTabVar{+/$ $,-/$ $,+/$ $}
			\end{tikzpicture}
		\end{center}
		Dựa vào bảng biến thiên, ta có $\min\limits_{u \in (0;70)} f(u)=f\left(\dfrac{15}{2}\right)=\dfrac{29}{30}$.\\
		Do đó ta có $T(x;y)=f(x)+f(y) \ge \dfrac{29}{30}+\dfrac{29}{30}=\dfrac{29}{15}$ giờ, tức là $116$ phút.
	}
\end{ex}

\begin{ex}%[Dự án đề TT Toán NH25-26-Dot 4- Hải Phụng]%[THCS-THPT LÊ THÁNH TÔNG - TP.HCM]%[GVPB - Viet Hoang Pham]%[2D4C3-2]
	\immini[thm]{
		Anh Nghĩa có một mảnh đất dạng hình thang cong $OABC$  ($B$ là điểm cực đại của đồ thị hàm số $y=f(x)$ được mô hình hóa trong mặt phẳng $Oxy$  (đơn vị mỗi trục là $10$ m). Anh Nghĩa chia mảnh đất hình thang cong $OABC$  thành hai phần để làm hồ bơi và làm vườn trồng cỏ, có đường ngăn cách bởi một phần của đồ thị hàm bậc ba  $y=f(x)$ như hình vẽ. Biết đơn giá làm hồ bơi là $400\,000$ đồng/m$^2$, giá trồng cỏ là $200\,000$ đồng/m$^2$. Tổng chi phí anh Nghĩa phải trả là $295$ triệu đồng. Bên cạnh đó có một con đường nhựa được mô hình hóa bởi hàm số $g(x) = \dfrac{x+1}{x-2}$ $(x>2)$. Anh Nghĩa muốn làm một đoạn đường $MN$ từ vườn anh đến con đường nhựa đó. Hãy tính độ dài ngắn nhất của đoạn đường $MN$. (Đơn vị: mét; làm tròn đến hàng phần trăm)
	}{
		\begin{tikzpicture}[>=stealth,scale=1, line join = round, line cap = round]
			\tikzset{label style/.style={font=\footnotesize}}
			\draw[->] (-1,0)--(0,0) node[below left]{$O$} -- (5,0) node[above]{$x$};
			\draw[->] (0.,-1) -- (0.,5) node[left]{$y$};
			\foreach \x in {3}
			\draw[shift={(\x,0)}] (0pt,2pt) -- (0pt,-2pt) node[below] {\footnotesize $\x$};
			\draw [smooth, samples=100, domain=-0.5:3] plot (\x,{-(\x)^3+3*(\x)^2}) node[above right]{$f(x)$};
			
			\draw[fill=blue, opacity=0.1] (0,0) plot[domain=0:2] (\x,{-(\x)^3+3*(\x)^2})--(0,4)--(0,0)--cycle;
			\draw[fill=green!80!black, opacity=0.2] (0,0) plot[domain=0:3] (\x,{-(\x)^3+3*(\x)^2})--(3,0)--(0,0)--cycle;
			\draw[fill=black] (0,4) circle(1.3pt) node[left]{$A$} -- (2,4) circle(1.3pt) node[above]{$B$};
			\draw [smooth, samples=100, domain=2.8:4.8] plot (\x, {(\x+1)/(\x-2)}) node[above right]{$g(x)$};
			
			\draw[fill=black] (3.3,3.3) circle(1.3pt) node[above right]{$N$} -- (2.6,2.704) circle(1.3pt) node[below left]{$M$};
			\draw (0.6,3) node{\footnotesize Hồ bơi} (2,1) node{\footnotesize Vườn anh Nghĩa};
		\end{tikzpicture}
	}
	\shortans{$7{,}49$}
	\loigiai{
		Dựa vào đồ thị ta thấy $f(x) = ax^2(x-3) = ax^3-3ax^2$.\\
		Suy ra $f'(x) = 3ax^2-6ax = 3ax(x-2)$ có nghiệm  $x=0$ hoặc $x =2$.\\
		Ta có $f(2) = -4a$. \\
		Diện tích hồ bơi là $S_1 = 100 \displaystyle\int\limits_{0}^{2} [-4a-f(x)] \mathrm{\,d}x = 100  \displaystyle\int\limits_{0}^{2} [-4a-a(x^3-3x^2)] \mathrm{\,d}x = -400a$. \\
		Diện tích vườn là $S_2 = 100  \displaystyle\int\limits_{0}^{3} f(x) \mathrm{\,d}x = 100a  \displaystyle\int\limits_{0}^{3} (x^3-3x^2) \mathrm{\,d}x = -675a$. \\
		Theo đề bài ta có 
		\begin{eqnarray*}
			400\,000S_1 + 200\,000S_2 = 295\,000\,000 &\Leftrightarrow& 400\,000(-400a) + 200\,000(-675a) = 295\,000\,000 \\
			&\Leftrightarrow& a=-1
		\end{eqnarray*}
		Vậy $f(x) = -x^3 +3x^2$.\\
		Gọi $M(m;-m^3+3m^2)$, $N\left(n;1+\dfrac{3}{n-2}\right)$, với $n>2$. \\
		Để $MN$ nhỏ nhất thì tiếp tuyến của hai đồ thị tại $M,N$ phải song song với nhau nên 
		$$f'(m)=g'(n) \Leftrightarrow -3m^2+6m = \dfrac{-3}{(n-2)^2} \Leftrightarrow n = 2 + \dfrac{1}{\sqrt{m^2-2m}} \text{ (vì } n>2).$$\\
		Khi đó $N(2 + \dfrac{1}{\sqrt{m^2-2m}};1+3\sqrt{m^2-2m})$. \\
		Từ đó $MN = \sqrt{\left(2+\dfrac{1}{\sqrt{m^2-2m}}-m\right)^2 + \left(1+3\sqrt{m^2-2m} + m^3 -3m^2\right)^2}$. \\
		Sử dụng MTCT tìm được $MN_{\min} \approx 0{,}7485$ khi $ m\approx 2{,}37$. \\
		Vậy độ dài ngắn nhất của đoạn $MN$ là $d_{\min} = 0{,}7485 \cdot 10 \approx 7{,}49$.
	}
\end{ex}

\begin{ex}%[2D4C3-2]%[Dự án C đợt 3 NH25-26- Phạm Đăng Đăng]%[TT-THPT-TranNhanTong-HaNoi-N25-26]%[GVPB - Phan Quốc Trí]
	Hoạ sĩ thiết kế logo hình con cá cho doanh nghiệp kinh doanh hải sản. Logo là hình phẳng giới hạn bởi hai parabol với kích thước được cho trong hình sau (đơn vị trên mỗi trục toạ độ là cm).
	\begin{center}
		\begin{tikzpicture}[scale=0.5, line cap=round, line join=round, font=\small,>=stealth]
			\fill[gray,smooth]
			plot[domain=-7:6] (\x,{3 - \x*\x/12})
			-- plot[domain=6:-7] (\x,{-5 + 5*\x*\x/36})
			-- cycle;
			% Khung chữ nhật
			\draw (-7,-5) rectangle (7,3);
			
			% Trục tọa độ
			\draw[->] (-7,0) -- (9,0) node[right] {$x$};
			\draw[->] (0,-5) -- (0,3.5) node[above] {$y$};
			\node at (0,0) [below left] {$O$};
			
			% Parabol trên (kéo dài tới biên trái)
			\draw[domain=-7:6,smooth,variable=\x]
			plot (\x,{3 - \x*\x/12});
			
			% Parabol dưới (kéo dài tới biên trái)
			\draw[domain=-7:6,smooth,variable=\x]
			plot (\x,{-5 + 5*\x*\x/36});
			
			% Hai điểm
			\coordinate (A) at (3,2.25);
			\coordinate (B) at (3,-3.75);
			
			% Đường cong tròn hơn
			\draw
			(A)
			.. controls (1.8,1) and (1.8,-2.5) ..
			(B);
			
			% Đánh dấu điểm
			\fill (A) circle (1.5pt);
			\fill (B) circle (1.5pt);
			
			
			
			% Điểm trong hình
			\fill (4,0.8) circle (0.15);
			
			% Kích thước ngang
			\draw[<->] (-7,-5.8) -- (-6,-5.8);
			\node at (-6.5,-6.3) {$1$ cm};
			
			\draw[<->] (-6,-5.8) -- (0,-5.8);
			\node at (-3,-6.3) {$6$ cm};
			
			\draw[<->] (0,-5.8) -- (6,-5.8);
			\node at (3,-6.3) {$6$ cm};
			
			% Kích thước dọc
			\draw[<->] (7.5,0) -- (7.5,3);
			\node[right] at (7.5,1.5) {$3$ cm};
			
			\draw[<->] (7.5,0) -- (7.5,-5);
			\node[right] at (7.5,-2.5) {$5$ cm};
			
		\end{tikzpicture}
	\end{center}
	Diện tích logo bằng bao nhiêu cm$^2$ (không làm tròn kết quả các phép tính trung gian, chỉ làm tròn kết quả cuối cùng)?
	\shortans{$65{,}4$}
	\loigiai{Do hình con cá giới hạn bởi hai parabol nên ta gọi $(P_1)$ là parabol có hệ số $a_1>0$ có công thức dạng 
		$
		y_1=a_1x^2+b_1x+c_1$.\\
		Khi đó $(P_1)$ đi qua các điểm có tọa độ là $(6;0)$, $(0;-5)$, $(-6;0)$.\\
		Từ đó ta có hệ phương trình
		$\heva{
			&36a_1+6b_1+c_1=0\\
			&36a_1-6b_1+c_1=0\\
			&c_1=-5
		}
		\Leftrightarrow
		\heva{
			&a_1=\dfrac{5}{36}\\
			&b_1=0\\
			&c_1=-5.
		}
		$\\
		Vậy $(P_1)$ có công thức là
		$
		y_1=\dfrac{5}{36}x^2-5$.\\
		Gọi $(P_2)$ là parabol có hệ số $a_2<0$ có công thức dạng $y_2=a_2x^2+b_2x+c_2$.\\
		Khi đó $(P_2)$ đi qua các điểm có tọa độ là $(6;0)$, $(0;3)$, $(-6;0)$.\\
		Từ đó ta có hệ phương trình
		$\heva{
			&36a_2+6b_2+c_2=0\\
			&36a_2-6b_2+c_2=0\\
			&c_2=3}
		\Leftrightarrow
		\heva{
			&a_2=-\dfrac{1}{12}\\
			&b_2=0\\
			&c_2=3.}$\\
		Vậy $(P_2)$ có công thức là $
		y_2=-\dfrac{1}{12}x^2+3.
		$\\
		Khi đó diện tích con cá được tính bởi công thức
		$$
		S=\displaystyle\int \limits_{-7}^{6}\left|y_1-y_2\right|\mathrm{\,d}x
		=\displaystyle\int\limits_{-7}^{6}\left|\dfrac{2}{9}x^2-8\right|\mathrm{\,d}x
		\approx 65{,}4.
		$$
		Vậy diện tích logo là $65{,}4$ cm$^2$.}
\end{ex}

\begin{ex}%[2D4C3-4]%[Dự án C đề thi thử TN Môn Toán NH25-26 - Dot 4 - Phạm Đăng Đăng]%[THPT Lê Quý Đôn-Đống Đa-Hà Nội]%[GVPB - Phan Quốc Trí]
	Sân vận động Sport Hub (Singapore) là sân có mái vòm kỳ vĩ nhất thế giới. Đây là nơi diễn ra lễ khai mạc Đại hội thể thao Đông Nam Á năm $2015$. Nền sân phẳng là một elip $(E)$ có trục lớn dài $150$ m, trục bé dài $90$ m. Một mặt phẳng vuông góc với trục lớn của $(E)$ và cắt elip ở $A$, $B$ và cắt phần dưới mái vòm theo một cung tròn của đường tròn tâm $I$ thoả mãn góc $\widehat{AIB}=90^\circ$ (tham khảo hình bên dưới). Các kĩ sư cần tính toán để lắp máy điều hòa cho sân vận động, mỗi máy điều hoà có công suất $100\,000$  BTU đáp ứng được không gian $500$ m$^3$. Hỏi các kĩ sư cần dùng tối thiểu bao nhiêu máy để đáp ứng đủ theo yêu cầu cho sân vận động?
	\begin{center}
		\begin{tikzpicture}[scale=1,font=\footnotesize,line join=round
			,line cap=round,>=stealth]
			\draw[thick] (0,0) ellipse (3cm and 1.5cm);
			\coordinate (A) at ({3*cos(50)}, {1.5*sin(50)});
			\coordinate (B) at ({3*cos(-50)}, {1.5*sin(-50)});
			\draw[thick] (A) -- (B);
			\node[above] at (A) {$A$};
			\node[below] at (B) {$B$};
			\fill[black] (A) circle (1.5pt);
			\fill[black] (B) circle (1.5pt);
			
			\begin{scope}[xshift=8.5cm]
		\coordinate (I) at (0,0);
				\draw[thick] (I) circle (2cm);
				\coordinate (A) at ({2*cos(45)}, {2*sin(45)});
				\coordinate (B) at ({2*cos(-45)}, {2*sin(-45)});
				\draw[thick] (I) -- (A);
				\draw[thick] (I) -- (B);
				\node[left] at (I) {$I$};
				% Label A nằm phía trên bên phải điểm
				\node[above right] at (A) {$A$};
				% Label B nằm phía dưới bên phải điểm
				\node[below right] at (B) {$B$};
				
				% Vẽ các dấu chấm tròn tại vị trí các điểm để đánh dấu rõ hơn
				\fill[black] (I) circle (1.5pt);
				\fill[black] (A) circle (1.5pt);
				\fill[black] (B) circle (1.5pt);
			\end{scope}
		\end{tikzpicture}
	\end{center}
	\shortans{$232$}
	\loigiai{\begin{center}
			\begin{tikzpicture}[scale=1,font=\footnotesize,line join=round
				,line cap=round,>=stealth]
				% 1. Định nghĩa tâm O
				\coordinate (O) at (0,0);
				\draw[thick] (O) ellipse (3cm and 1.2cm);
				\draw[->, thick] (O) -- (0,3.5) node[left] {$z$};
				\draw[->, thick] (O) -- (3.8, 1) node[above] {$x$};
				\draw[->, thick] (O) -- (-2.8, 2) node[below] {$y$};
				\fill (O) circle (1.5pt);
				\begin{scope}[xshift=8cm]
					\coordinate (O) at (0,0);
					\draw[->, thick] (-3.5, 0) -- (3.5, 0) node[above right] {$x$};
					\draw[thick] (O) ellipse (3cm and 1.8cm);
					\node[above] at (O) {$O$};
					\fill[black] (O) circle (1.5pt);
					\coordinate (A) at ({3*cos(50)}, {1.8*sin(50)});
					\coordinate (B) at ({3*cos(-50)}, {1.8*sin(-50)});
					\draw[thick] (A) -- (B);
					\node[above right] at (A) {$A$};
					\node[below right] at (B) {$B$};
					\fill[black] (A) circle (1.5pt);
					\fill[black] (B) circle (1.5pt);
				\end{scope}
			\end{tikzpicture}
		\end{center}
		Chọn hệ trục như hình vẽ.\\ Ta cần tìm diện tích của $S(x)$ thiết diện.
		Gọi $\mathrm{d}(O, AB) = x$.\\
		$(E): \dfrac{x^2}{75^2} + \dfrac{y^2}{45^2} = 1$.
		Khi đó $AB = 2y = 2\sqrt{45^2 \left( 1 - \dfrac{x^2}{75^2} \right)} = 90\sqrt{1 - \dfrac{x^2}{75^2}}$.\\
		Suy ra $R = \dfrac{AB}{\sqrt{2}} = \dfrac{90}{\sqrt{2}}\sqrt{1 - \dfrac{x^2}{75^2}} \Rightarrow R^2 = \dfrac{90^2}{2} \left( 1 - \dfrac{x^2}{75^2} \right)$.\\
		$S(x) = \dfrac{1}{4}\pi R^2 - \dfrac{1}{2}R^2 = \left( \dfrac{1}{4}\pi - \dfrac{1}{2} \right)R^2 = (\pi - 2)\dfrac{2025}{2}\left( 1 - \dfrac{x^2}{75^2} \right)$.\\
		Vậy thể tích khoảng không bên dưới mái che và bên trên mặt sân là
		$$V = \displaystyle\int\limits_{-75}^{75} (\pi - 2)\dfrac{2025}{2}\left( 1 - \dfrac{x^2}{75^2} \right)\, \mathrm{\,d}x \approx 115\,586 \, \text{m}^3.$$
		Số chiếc điều hòa cần lắp là
		$115586 : 500 = 231{,}172$.\\
		Vậy cần $232$ chiếc điều hòa.}
\end{ex}

\begin{ex}%[2D4V1-6]%[Dự án đề thi thử đợt 3-Nguyễn Đức Lợi]%[SGD tỉnh Lạng Sơn]%[GVPB - ThanhTa]
	Ở giai đoạn thải trừ, giai đoạn cuối sau khi một người uống một liều thuốc, nồng độ thuốc trong máu, ký hiệu là $C(t)$ (đơn vị: mg/l), giảm dần sau $t$ giờ kể từ khi giai đoạn này bắt đầu. Khi đó, tốc độ giảm nồng độ $C'(t)$ tỉ lệ với chính nồng độ hiện có, tức là $\dfrac{C'(t)}{C(t)}=-k$ ($k$ là một hằng số dương). Biết rằng khi bắt đầu giai đoạn thải trừ, nồng độ thuốc còn lại là $12$ mg/l và sau $6$ giờ kể từ lúc bắt đầu thải trừ, nồng độ đo được là $3$ mg/l. Sau khoảng bao nhiêu giờ thì nồng độ còn lại bằng $2$ mg/l ({\it kết quả làm tròn đến hàng phần mười})?
	\par
	\shortans{7,8}
	\loigiai{
		Ta có $\dfrac{C'(t)}{C(t)}=-k \Rightarrow \displaystyle \int\dfrac{C'(t)}{C(t)}\mathrm{\,d}t=\displaystyle \int-k\mathrm{\,d}t \Rightarrow \ln \big|C(t)\big|=-kt+m$ ($m$ là hằng số).\\
		Vì $C(t)>0$ nên ta có $C(t)=\mathrm{e}^{-kt+m}=\mathrm{e}^m \cdot \mathrm{e}^{-kt}=A\mathrm{e}^{-kt}$ (đặt hằng số $\mathrm{e}^m=A$).\\
		$C(0)=12 \Leftrightarrow A\mathrm{e}^{-k.0}=12 \Leftrightarrow A=12$.\\
		$C(6)=3 \Leftrightarrow 12\mathrm{e}^{-k\cdot6}=3 \Leftrightarrow \mathrm{e}^{-6k}=\dfrac{1}{4} \Leftrightarrow k=-\dfrac{1}{6} \ln \dfrac{1}{4}$.\\
		Suy ra $C(t)=12\mathrm{e}^{\tfrac{1}{6}\left(\ln \tfrac{1}{4}\right)t}$.\\
		Mà ta lại có \[C(t)=2\Leftrightarrow 12\mathrm{e}^{\tfrac{1}{6}\left(\ln \tfrac{1}{4}\right)t}=2\Leftrightarrow \mathrm{e}^{\tfrac{1}{6}\left(\ln \tfrac{1}{4}\right)t}=\dfrac{1}{6}\Leftrightarrow \dfrac{1}{6}\left(\ln \dfrac{1}{4}\right)t=\ln \dfrac{1}{6}\Leftrightarrow t=\dfrac{6 \ln \dfrac{1}{6}}{\ln \dfrac{1}{4}} \approx 7{,}8.\]
		Vậy sau khoảng $7{,}8$ giờ thì nồng độ còn lại bằng $2$ mg/l.
	}
\end{ex}

\begin{ex}%[2D4V2-6]%[Dự án C đề thi thử TN Môn Toán NH25-26 - Dot 4 - Vũ Hồng Toàn]%[THPT Mai Anh Tuấn - Thanh Hóa]%[GVPB - Lê Hữu Kiệt]
	Một vật chuyển động trong $4$ giờ với vận tốc $v$\,(km/h) phụ thuộc thời gian $t$\,(h) có đồ thị của vận tốc như hình bên. Trong khoảng thời gian $2$ giờ kể từ khi bắt đầu chuyển động, đồ thị đó là một phần của đường parabol có đỉnh $I(2;7)$ và trục đối xứng của parabol song song với trục tung, khoảng thời gian còn lại đồ thị là đoạn thẳng $IA$. Quãng đường $s$ mà vật di chuyển được trong $4$ giờ đó là bao nhiêu km {\it (kết quả được làm tròn đến hàng phần chục)}?
	\begin{center}
		\begin{tikzpicture}[xscale=1,yscale=0.6, font=\footnotesize, line join=round, line cap=round,>=stealth]
			\tikzset{declare function={f(\x)=-(\x)^2+4*(\x)+3;
			xmin=0;xmax=4;ymin=0;ymax=7;},
			smooth,samples=100
		}
		\path (4,3)coordinate (A) (2,7)coordinate (I);
		\draw[->] (xmin-0.5,0)--(0,0) node[below left]{$O$} --(xmax+1,0) node[shift={(-100:7pt)}] {$x$};
		\draw[->] (0,ymin-1)--(0,ymax+1.5) node[shift={(200:7pt)}] {$y$};
		\fill (0,0) circle (1.2pt);
		\draw[red,domain=xmin:2] plot (\x,{f(\x)});
		\draw[dashed,red,domain=2:xmax] plot (\x,{f(\x)});
		\draw (I)--(A);
		\foreach \y/\goc in {3/180}{
			\fill (0,\y) circle (1pt) node[shift={(\goc:7pt)}]{$\y$};
		}
		\foreach \x/\y/\px/\py in {2/7/-90/180,4/3/-90/180}{
			\fill (\x,0) circle (1.2pt) node[shift={(\px:7pt)}]{$\x$};
			\fill (0,\y) circle (1.2pt) node[shift={(\py:7pt)}]{$\y$};
			\fill (\x,\y) circle (1.2pt);
			\draw[dashed, thin] (\x,0) |-(0,\y);
		}
		\foreach \x /\gN in {A/20,I/90}{
			\fill (\x) circle (1pt)($(\x)+(\gN:9pt)$) node {$\x$};
		}
	\end{tikzpicture}
\end{center}
\par\shortans{21,3}
\loigiai{
	Phương trình parabol là $v = -t^2 + 4t + 3$\,(km/h).\\
	Quãng đường vật chuyển động được trong khoảng thời gian $2$ giờ đầu là
	$$\displaystyle\int\limits_0^2 v(t)\mathrm{\,d}t = \displaystyle\int\limits_0^2 \left(-t^2 + 4t + 3\right)\mathrm{d}t = \dfrac{34}{3}\,\text{(km)}.$$
	Quãng đường vật đi được trong $2$ giờ sau là $\dfrac{(7+3)\cdot 2}{2} = 10$\,km.\\
	Vậy tổng quãng đường vật đi được trong $4$ giờ là $\dfrac{34}{3}+10\approx 21{,}3$\,(km).
}
\end{ex}

\begin{ex}%[2D4V3-2]%[TT-SGD-TuyenQuang-L1-NH25-26]%[Võ Thị Thùy Trang]%[GVPB: Thanh Phát]
	\immini[thm]{
		Người ta dự định trồng hoa để trang trí vào phần tô đậm trong hình vẽ bên. Biết rằng phần tô đậm là diện tích hình phẳng giới hạn bởi đồ thị hai hàm số $y=f(x)=ax^3+bx^2+cx+6$ và $y=g(x)=-bx^2+mx+n$ trong đó $a$, $b$, $c$, $m$, $n \in \mathbb{R}$. Biết đồ thị hai hàm số đã cho cắt nhau tại các điểm có hoành độ lần lượt bằng $-2$; $1$; $3$. Chi phí trồng hoa là $150\,000$ đồng/$1$ m$^2$ và đơn vị trên mỗi trục tọa độ là mét. Tổng chi phí để trồng hoa theo dự định là bao nhiêu nghìn đồng (\textit{làm tròn kết quả đến hàng đơn vị})?
		}
	{\begin{tikzpicture}[>=stealth,line join=round,line cap=round,font=\footnotesize,scale=.7]
			\def\aa{1} \def\bb{0} \def\cc{0}
			\draw[->] (-5,0)--(5,0)node[below right]{$x$};
			\draw[->] (0,-1)--(0,10)node[left]{$y$};
			\fill (0,0)node[below left]{$O$};
			\draw[black,samples=150,smooth,domain=-3.2:3.2] plot(\x,{\aa*(\x)^2+\bb*\x+\cc});
			\draw[black,samples=150,smooth,domain=-2.37:3.1] plot(\x,{1*(\x)^3-1*(\x)^2-5*\x+6});
			\fill (3,9)circle(1pt)
			(1,1)circle(1pt)
			(-2,4)  circle(1pt)
			;
			\draw[dashed](3,0)--(3,9) (-2,0)--(-2,4) (1,0)--(1,1);
			\node at (-2,0)[below]{$-2$};
			\node at (1,0)[below]{$1$};
			\node at (3,0)[below]{$3$};
			\path[pattern = north east lines, pattern color = black][domain=-2.0:1.0] plot(\x,{\aa*(\x)^2+\bb*\x+\cc})-- (1.,1.) {[smooth,samples=50,domain=1.0:-2.0] -- plot(\x,{(\x)^(3)-(\x)^(2)-5*(\x)+6})} -- (-2.,4.) -- cycle;
			\path[pattern = north east lines, pattern color = black]{[smooth,samples=50,domain=1.0:3.0]plot(\x,{\aa*(\x)^2+\bb*\x+\cc})} -- (3.,9.) {[smooth,samples=50,domain=3.0:1.0] -- plot(\x,{(\x)^(3)-(\x)^(2)-5*(\x)+6})} -- (1.,1.) -- cycle;
	\end{tikzpicture}}
	\par\shortans{3163}
	\loigiai{Phương trình hoành độ giao điểm của hai đồ thị hàm số $y=f(x)$ và $y=g(x)$ là 
		\[f(x) - g(x) = 0 \Leftrightarrow ax^3 + 2bx^2 + (c-m)x + 6 - n = 0 \quad (*)\]
		Theo giả thiết, đồ thị hai hàm số cắt nhau tại các điểm có hoành độ $x = -2$, $x = 1$, $x = 3$ nên phương trình $(*)$ có $3$ nghiệm phân biệt là $-2$; $1$; $3$. Do đó ta có 
		\begin{align*}
			f(x) - g(x) &= a(x+2)(x-1)(x-3) \\
			&= a\left( x^2 + x - 2\right) (x - 3) \\
			&= a\left(x^3 - 2x^2 - 5x + 6\right) \\
			&= ax^3 - 2ax^2 - 5ax + 6a.
		\end{align*}
		Đồng nhất hệ số với biểu thức ban đầu của $f(x) - g(x)$, ta được $\heva{
			&2b = -2a \\
			&c - m = -5a \\
			&6 - n = 6a.}$\\
		Quan sát đồ thị, ta thấy đường Parabol $y = g(x)$ đi qua gốc tọa độ $O(0;0)$.\\
		Suy ra $g(0) = 0 \Rightarrow n = 0$.\\
		Thay $n = 0$ vào phương trình $6 - n = 6a$, ta được $6 = 6a \Leftrightarrow a = 1$.\\
		Khi đó $f(x) - g(x) = x^3 - 2x^2 - 5x + 6$.\\
		Diện tích phần tô đậm là
		\[S =\displaystyle\int\limits_{-2}^{3} \left| f(x) - g(x)\right|  \mathrm{\,d}x =\displaystyle\int\limits_{-2}^{1} \left| x^3 - 2x^2 - 5x + 6\right|  \mathrm{\,d}x + \int_{1}^{3} \left| x^3 - 2x^2 - 5x + 6\right|  \mathrm{\,d}x.\]
		Dựa vào đồ thị, trên khoảng $(-2; 1)$ thì $f(x)$ nằm trên $g(x)$ và trên khoảng $(1; 3)$ thì $f(x)$ nằm dưới $g(x)$. Ta có
		\begin{align*}
			S &=\displaystyle\int\limits_{-2}^{1} \left( x^3 - 2x^2 - 5x + 6\right) \mathrm{\,d}x - \displaystyle\int\limits_{1}^{3} \left( x^3 - 2x^2 - 5x + 6\right)  \mathrm{\,d}x \\
			&= \left( \dfrac{x^4}{4} - \dfrac{2x^3}{3} - \dfrac{5x^2}{2} + 6x \right) \Bigg|_{-2}^{1} - \left( \dfrac{x^4}{4} - \dfrac{2x^3}{3} - \dfrac{5x^2}{2} + 6x \right) \Bigg|_{1}^{3} \\
			&= \left( \dfrac{37}{12} - \dfrac{-152}{12} \right) - \left( \dfrac{-27}{12} - \dfrac{37}{12} \right) \\
			&= \dfrac{189}{12} - \left( -\dfrac{64}{12} \right) = \dfrac{63}{4} + \dfrac{16}{3} = \dfrac{253}{12}\,\text{m}^2.
		\end{align*}
		Tổng chi phí dự định để trồng hoa là
		\[T = \dfrac{253}{12} \cdot 150\,000 = 3\,162\,500\, \text{(đồng)}\approx 3\,163\, \text{(nghìn đồng).}\]
	}
\end{ex}

\begin{ex}%[2D4V3-2]%[KNTT - Lớp 12 - Dự án C - Đợt 2]%[Loc do]%[PB: Phan Anh]
	\immini[thm]{Một bể chứa nước có diện tích mặt cắt ngang được tô màu như hình bên, ở đó đơn vị trên các trục đo bằng mét. Trên mặt cắt ngang, phần đáy của bể chứa có phương trình $y=k(x-8)^2$. Tính diện tích của mặt cắt ngang theo mét vuông.}{\begin{tikzpicture}[x=0.5cm,y=0.5cm,scale=0.8,>=stealth]
			\def\aa{0.094} \def\bb{-1.5} \def\cc{6}
			\draw[->] (-1.5,0)--(13,0)node[below right]{$x$};
			\draw[->] (0,-1.25)--(0,7)node[left]{$y$};
			\fill (0,0)node[below left]{$O$};
			\begin{scope}
				\clip (-1,-1) rectangle (14,8);
				\fill[pattern=crosshatch,pattern color=green]
				plot[samples=200,domain=0:12,variable=\x] (\x,{6})-- plot[samples=200,domain=12:0,variable=\x] (\x,{0.094*(\x)^2-1.5*(\x)+6})-- cycle;
				\draw plot[samples=200,domain=0:12,variable=\x] (\x,{0*(\x)+6});
				\draw plot[samples=200,domain=12:0,variable=\x] (\x,{0.094*(\x)^2-1.5*(\x)+6});
			\end{scope}
			\draw (12,0)--(12,6) node[above] {$N(12,6)$};
			\fill (0,6) circle (1pt) node[left] {$M$};
			\fill (8,0) circle (1pt) node[below] {$(8,0)$};
			\fill (12,0) circle (1pt) node[below] {$P$};
	\end{tikzpicture}}
	\par\shortans[0]{54}
	\loigiai{
		Đáy của bể nước là một phần của parabol có phương trình $y=k(x-8)^2$.\\
		Parabol đi qua điểm $M(0;\, 6)\Rightarrow k\cdot(-8)^2 =6\Rightarrow k=\dfrac{3}{32}$.\\ 
		Khi đó ta có Parabol $y=\dfrac{3}{32} (x-8)^2$.\\
		Vậy diện tích của mặt cắt ngang là
		$S=\displaystyle\int\limits _0^{12}\left[6-\left(\dfrac{3}{32} \left(x-8\right)^2 \right)\right] \mathrm{\,d}x=54$ m$^2$.
	}
\end{ex}

\begin{ex}%[2D4V3-4]%[Dự án đề thi thử NH25-26]%[THPT Chuyên Phan Bội Châu - Nghệ An]%[Sỹ Trường]
	\immini[thm]
	{
		Khi thi công tuyến cao tốc Nghi Sơn-Bãi Vọt, đơn vị thi công cần san một ngọn núi nhỏ. Biết đường viền chân núi là một elip có trục lớn $AB=400$ (m) và trục bé $CD=200$ (m); mặt cắt bởi mặt phẳng chứa $AB$ và vuông góc với mặt đất là nửa hình tròn đường kính $AB$; mặt cắt bởi các mặt phẳng vuông góc với $AB$ có dạng parabol với đỉnh thuộc nửa đường tròn đường kính $AB$ (tham khảo hình vẽ). Tính thể tích ngọn núi đó theo đơn vị triệu $\text{m}^3$ (làm tròn đến hàng phần trăm).
	}
	{
		\begin{tikzpicture}[line join=round, line cap=round, >=stealth, font=\footnotesize, scale=1]
			\begin{scope}[rotate=15]
				\coordinate (O) at (0,0);
				\coordinate (A) at (-3,0);
				\coordinate (B) at (3,0);
				\coordinate (C) at (0,-1);
				\coordinate (D) at (0,1);
				\def\xh{-2.2}
				\def\yh{0.68} 
				\coordinate (H) at (\xh, 0);
				\coordinate (E) at (\xh, -\yh);
				\coordinate (F) at (\xh, \yh);
				\draw (B) arc (0:-180:3 and 1);
				\draw[dashed] (B) arc (0:180:3 and 1);
				\draw[dashed] (A) -- (B);
				\draw[dashed] (C) -- (D);
				\draw[dashed] (E) -- (F);
				\draw ($(H)+(-0.15,0)$) -- ($(H)+(-0.15,-0.15)$) -- ($(H)+(0,-0.15)$);
			\end{scope}
			\draw (B) [x={(B)}, y={(0,3)}] arc (0:180:1 and 1);
			\def\zh{2.04}
			\coordinate (G) at ($(H) + (0, \zh)$);
			\begin{scope}[shift={(H)}]
				\coordinate (vecX) at ($(F)-(H)$);
				\coordinate (vecY) at (0, \zh);
				\draw [x={(vecX)}, y={(vecY)}] plot[domain=-1:0, samples=50, variable=\t] (\t, {1-\t*\t});
				\draw[dashed][x={(vecX)}, y={(vecY)}] plot[domain=0.45:1, samples=50, variable=\t] (\t, {1-\t*\t});
			\end{scope}
			\draw[dashed] (H) -- (G);
			\draw[dashed] ($(H)!0.25!(E)$) -- ++(0, 0.2) -- ($(H)+(0,0.2)$);
			\foreach \p/\ang in {A/150, B/30, C/-45, D/135} {
				\fill (\p) circle (1pt);
				\node at ($(\p)+(\ang:3.5mm)$) {$\p$};
			}
		\end{tikzpicture}	
	}
	\shortans{$7{,}11$}
	\loigiai{
		Chọn hệ trục tọa độ $Oxyz$ sao cho gốc $O$ trùng với tâm của đường elip chân núi, trục $Ox$ chứa trục lớn $AB$, trục $Oy$ chứa trục bé $CD$ và trục $Oz$ hướng thẳng đứng lên trên.
		\begin{itemize}
			\item Vì trục lớn $AB = 400$ nên bán trục lớn $a = 200$.
			\item Vì trục bé $CD = 200$ nên bán trục bé $b = 100$.
		\end{itemize}
		Phương trình đường elip viền chân núi trong mặt phẳng $Oxy$ là
		\[\frac{x^2}{200^2} + \frac{y^2}{100^2} = 1 \Rightarrow y = \pm 100\sqrt{1 - \frac{x^2}{200^2}} = \pm \frac{1}{2}\sqrt{200^2 - x^2}.\]
		Xét mặt phẳng vuông góc với trục $Ox$ tại điểm có hoành độ $x$ ($x \in [-200; 200]$), mặt phẳng này cắt ngọn núi theo một thiết diện là hình parabol.
		\begin{itemize}
			\item Đáy của parabol này là đoạn thẳng nối hai điểm trên elip, có độ dài là
			\[d = 2y = 2 \cdot \frac{1}{2}\sqrt{200^2 - x^2} = \sqrt{200^2 - x^2}.\]
			\item Theo giả thiết, mặt cắt chứa $AB$ và vuông góc với mặt đất là nửa đường tròn đường kính $AB$. Do đó, phương trình của nửa đường tròn này trong mặt phẳng $Oxz$ là $x^2 + z^2 = 200^2$ (với $z \ge 0$).
			\item Chiều cao của thiết diện parabol (từ đáy lên đỉnh thuộc nửa đường tròn) chính là cao độ $z$
			\[h = \sqrt{200^2 - x^2}.\]
		\end{itemize}
		Diện tích của thiết diện parabol tại hoành độ $x$ được tính bằng $S = \dfrac{2}{3} \cdot \text{đáy} \cdot \text{chiều cao}$.\\ Suy ra 
		\[S(x) = \dfrac{2}{3} \cdot d \cdot h = \dfrac{2}{3} \cdot \sqrt{200^2 - x^2} \cdot \sqrt{200^2 - x^2} = \dfrac{2}{3}\left(200^2 - x^2\right).\]
		Thể tích ngọn núi cần san lấp là
		\allowdisplaybreaks
		\[\begin{aligned}[t]
			V &=\displaystyle\int\limits_{-200}^{200} S(x) \mathrm{\,d}x = \displaystyle\int\limits_{-200}^{200} \dfrac{2}{3}(200^2 - x^2) \mathrm{\,d}x \\
			&= \dfrac{4}{3}\displaystyle\int\limits_{0}^{200} (200^2 - x^2) \mathrm{\,d}x = \dfrac{4}{3} \left( 200^2x - \dfrac{x^3}{3} \right) \Bigg|_{0}^{200} \\
			&= \dfrac{4}{3} \left( 200^3 - \dfrac{200^3}{3} \right) = \frac{4}{3} \cdot \dfrac{2}{3} \cdot 200^3 = \dfrac{8}{9} \cdot 8\,000\,000 = \dfrac{64\,000\,000}{9} \text{ (m}^3).
		\end{aligned}\]
		Tính xấp xỉ ta được $V \approx 7\,111\,111{,}11 \text{ m}^3 \approx 7{,}11$ triệu $\text{m}^3$.\\
		Vậy thể tích ngọn núi cần tính xấp xỉ là $7{,}11$ triệu $\text{m}^3$.
	}
\end{ex}

\begin{ex}%[2D4V3-5]%[Dự án C - Đợt 3 - Sở Giáo Dục Sơn La]%[Nguyễn Chiến - Nguyễn Tiến]
	Từ một quả cầu bằng đá trắng bán kính bằng $1$ dm người ta khoan rút lõi ngay \lq\lq chính giữa\rq\rq\ quả cầu (trục đối xứng của lõi và quả cầu trùng nhau) như hình minh họa, đường kính lõi là $1$ dm. Thể tích còn lại của quả cầu bằng bao nhiêu dm$^3$? (không làm tròn kết quả ở các bước trung gian, làm tròn kết quả ở bước cuối cùng đến hàng phần trăm).
	\begin{center}
		\begin{tikzpicture}[scale=1.5]
			\def\rx{.5}\def\ry{.1}\def\a{433/500}\def\b{.5}\def\c{.15}
			\draw (0,0) circle(1)
			(-.5,-\a)--(-.5,\a) (.5,-\a)--(.5,\a);
			\foreach \x/\s in {.5/, .5/dashed}
			\foreach \y in {\a,-\a}
			\draw[\s] (\x,\y) arc(0:\ifx\s\empty-180\else180\fi:{\rx} and {\ry});
			\node at(0,-1.3){Quả cầu đá};
			\begin{scope}[shift={(2,-1)}]
				\tikzset{declare function={d = 3; r = 0.5;}}
				\draw[] (0,0)--++(d,0) to [out=60,in=-60] ++(0,2*r)
				--++(-d,0) to [out=-120,in=120] cycle;
				\draw (0,0) arc(270:90:{0.1} and {r})
				(d,0) arc(270:90:{0.1} and {r});
				\draw[dashed] (0,0) arc(-90:90:{0.1} and {r})
				(d,0) arc(-90:90:{0.1} and {r});
				
				\node at(1.5,-0.3){Lõi};
			\end{scope}
		\end{tikzpicture}
	\end{center}
	\shortans[oly]{$2{,}72$}
	\loigiai{
		\immini[thm]{
			Bán kính quả cầu là $R = 1$ dm.\\
			Lõi bị khoan rút có phần hình trụ với bán kính đáy $r = \dfrac{1}{2} = 0{,}5$ dm.\\
			Chiều cao của phần khối trụ là $$h = 2\sqrt{R^2 - r^2} = 2\sqrt{1^2 - 0{,}5^2} = \sqrt{3} \, \text{dm.}$$
			Chọn hệ trục tọa độ $Oxy$ với gốc $O$ là tâm quả cầu, trục $Ox$ dọc theo trục của khối trụ.\\
			Khi đó, phần thể tích còn lại của quả cầu bằng thể tích khối tròn xoay tạo thành khi cho hình phẳng giới hạn bởi đường cong $y = \sqrt{1 - x^2}$, đường thẳng $y = 0{,}5$ và hai đường thẳng $x = -\dfrac{\sqrt{3}}{2}$, $x = \dfrac{\sqrt{3}}{2}$ quay quanh trục $Ox$.
		}{
			\begin{tikzpicture}[scale=1.2, font=\footnotesize, line join=round, line cap=round, >=stealth]
				\def\R{2}
				\def\r{1}
				\def\h{1.73205}
				\draw[->] (-2.3,0) -- (2.3,0) node[below] {$x$};
				\draw[->] (0,-2.3) -- (0,2.3) node[left] {$y$};
				\draw[thick] (0,0) circle (\R);
				\draw[dashed] (-2.05,\r) -- (2.05,\r);
				\draw[dashed] (-2.05,-\r) -- (2.05,-\r);
				\coordinate [label=below left:$O$] (O) at (0,0);
				\coordinate [label=above right:$A$] (A) at (\h,\r);
				\coordinate [label=above left:$H$] (H) at (\h,0);
				\draw[thick] (O) -- (A) node[midway, above left] {$R$};
				\draw[dashed] (-\h,0) -- (-\h,\r);
				\draw[dashed] (H) -- (A) node[midway, right] {$r$};
				\foreach \i in {A,H,O} \draw[fill=black] (\i) circle(1pt);
				\draw pic[draw,angle radius=5]{right angle=A--H--O};
				\node[below] at (\h,0) {$\dfrac{\sqrt{3}}{2}$};
				\node[below] at (-\h,0) {$-\dfrac{\sqrt{3}}{2}$};
				\node[above left] at (0,\r) {$0{,}5$};
				\foreach \p in {(\h,0), (-\h,0), (0,\r)} \fill \p circle (1pt);
			\end{tikzpicture}
		}
		\noindent Thể tích cần tìm là
		{\allowdisplaybreaks
			\begin{eqnarray*}
				V &=& \pi \int\limits_{-\frac{\sqrt{3}}{2}}^{\frac{\sqrt{3}}{2}} \left[ \left(\sqrt{1 - x^2}\right)^2 - (0{,}5)^2 \right] \mathrm{\,d}x \\
				&=& \pi \int\limits_{-\frac{\sqrt{3}}{2}}^{\frac{\sqrt{3}}{2}} \left( 0{,}75 - x^2 \right) \mathrm{\,d}x \\
				&=& \pi \left( 0{,}75x - \dfrac{x^3}{3} \right) \Bigg|_{-\frac{\sqrt{3}}{2}}^{\frac{\sqrt{3}}{2}} = \dfrac{\pi\sqrt{3}}{2} \approx 2{,}72.
		\end{eqnarray*}}
	}
\end{ex}

\begin{ex}%[2H2H2-3]%[TT Cụm 5 - Ninh Bình-NH25-26]%[GVSB:Dương Công Tạo]%[GVPB1: Nguyễn Tấn Tài]
	Trong không gian với hệ trục toạ độ $Oxyz$, cho hình hộp $ABCD.A'B'C'D'$ có $A(-3;2;1)$, $C(4;2;0)$, $B'(-2;1;1)$, $D'(3;5;4)$. Gọi toạ độ điểm $A'(a;b;c)$. Tính giá trị của biểu thức $\dfrac{a-b-c}{2}$.
	\par  \shortans{-4,5}
	\loigiai{
		\immini[thm]{
			Gọi $I$, $I'$ lần lượt là trung điểm của $AC$ và $B'D'$ thì $I\left(\dfrac{1}{2};2;\dfrac{1}{2}\right)$. $I'\left(\dfrac{1}{2};3;\dfrac{5}{2}\right)$. \\
			Mặt khác $AA'I'I$ là hình bình hành nên \[\overrightarrow{AA'}=\overrightarrow{II'}\Leftrightarrow\heva{&a+3=\dfrac{1}{2}-\dfrac{1}{2}\\ &b-2=3-2\\ &c-1=\dfrac{5}{2}-\dfrac{1}{2}}\Leftrightarrow\heva{&a=-3\\ &b=3\\ &c=3}\Rightarrow A'(-3;3;3).\] 
		}
		{\begin{tikzpicture}[scale=.65,line join=round, line cap=round,font=\footnotesize]
				\def\a{3.5}
				\def\h{4}
				\path 	(0:0) coordinate (A)
				++(0:\a) coordinate (D)
				++(-130:\a/2) coordinate (C)
				($(A)+(C)-(D)$) coordinate (B)	
				($(A)!0.25!(C)$) coordinate (H)
				($(H)+(90:\h)$) coordinate (A')
				($(A')+(C)-(A)$) coordinate (C')
				($(C')+(B)-(C)$) coordinate (B')
				($(A')+(D)-(A)$) coordinate (D')
				($(C)!.5!(A)$) coordinate (I)
				($(C')!.5!(A')$) coordinate (I');
				\draw[dashed] 	(B)--(A)--(D)	(A)--(A') (A)--(C)	(I)--(I');
				\draw	(C)--(C') 	(D)--(D') 	(B)--(B')  (B)--(C)--(D) (A')--(B')--(C')--(D')--cycle	(A')--(C');
				\foreach \x/\g in {A/180,B/180,C/0,D/0,A'/180,B'/180,C'/0,D'/0,I/-135,I'/45}
				\fill[black] 	(\x) circle (1pt)
				($(\g:4mm)+(\x)$) node {$\x$};
			\end{tikzpicture}
		}
		\noindent Vậy $\dfrac{a-b-c}{2} = -4{,}5$.
	}
\end{ex}

\begin{ex}%[2H2K1-5]%[Dự án C-TT THPT NH25-26-Đợt 4- ĐỖ CHÍ TÂM]%[THPT Phan Bội Châu-Khánh Hòa]%[GVPB: Phạm Văn Long]
	Một nhóm học sinh thiết kế một chiếc lều dã ngoại có dạng hình nón theo cách sau: Từ một tấm bạt hình tròn có diện tích $4{,}84\pi$ m$^2$, nhóm cắt lấy một hình quạt có góc ở tâm là $\alpha \left(0{<}\alpha {<}2\pi \right)$ (như hình 1) để làm thành mặt xung quanh của lều (như hình 2). Sau đó, nhóm dùng một tấm bạt chống thấm cắt thành hình tròn để trải kín đáy lều. Biết chi phí loại bạt làm mặt xung quanh lều là $100\,000$ đồng/m$^2$ và chi phí loại bạt làm đáy lều là $150\,000$ đồng/m$^2$. Giả sử phần bạt thừa sau khi cắt không tính vào chi phí và kích thước các mép nối không đáng kể. Để thiết kế được chiếc lều có diện tích đáy lớn nhất với tổng chi phí không vượt quá $2\,300\,000$ đồng, chiều cao $h$ của lều bằng bao nhiêu mét (kết quả làm tròn đến hàng phần trăm).
	\begin{center}
		\begin{tikzpicture}
			% --- Hình 1: Hình quạt tròn ---
			\begin{scope}[shift={(0,0)}]
				\draw (0,0) coordinate (O) circle (1.5cm);
				\draw (0,0) -- (30:1.5cm) arc (30:330:1.5cm) -- (0,0);
				\filldraw[fill=green!20, draw=green!50!black] (0,0) -- (30:0.5cm) arc (30:330:0.5cm) -- cycle;
				\node at (-0.3, 0) {$\alpha$};
				\node at (0, -2) {\textbf{Hình 1}};
			\end{scope}	
		\end{tikzpicture} \hspace*{1cm}
		\begin{tikzpicture}[scale=1.2]
			% --- Hình 2: Hình nón ---
			\begin{scope}[shift={(5,0)}]
				% Đáy hình nón (elip)
				\draw[dashed] (1,0) arc (0:180:1cm and 0.3cm);
				\draw (1,0) arc (0:-180:1cm and 0.3cm);
				\draw (-1,0) -- (0,2.5) -- (1,0);
				
				% Các đường nét bên trong
				\draw[dashed] (0,0) -- (0,2.5) node[midway, left] {$h$};
				\draw[dashed] (0,0) -- (1,0) node[midway, below] {$r$};
				\draw (0,2.5) -- (1,0) node[midway, right] {$\ell$};
				
				\node at (0, -0.7) {\textbf{Hình 2}};
			\end{scope}	
		\end{tikzpicture}
	\end{center}
	\shortans{$1{,}52$}
	\loigiai{
		Tấm bạt ban đầu có diện tích $S=\pi {R^2}\Leftrightarrow {R^2}=\dfrac{S}{\pi}=4,84\Rightarrow R=2{,}2$ (m).\\
		Vì mặt xung quanh của lều được cắt từ tấm bạt hình tròn này, nên bán kính của tấm bạt là độ dài đường sinh của hình nón. Do đó $\ell=2,2$ (m).\\
		Diện tích xung quanh của hình nón là: ${S_{xq}}=\pi r\ell=2{,}2\pi r$.\\
		Diện tích đáy của hình nón là: ${S_d}=\pi{r^2}$.\\
		Tổng chi phí làm lều là $100\,000\cdot {S_{xq}}+150\,000\cdot {S_d}=100\,000\cdot 2{,}2\pi r+150\,000\pi r^2$.\\
		Theo đề bài $100\,000\cdot 2{,}2\pi r+150\,000\cdot \pi {r^2}\le 2\,300\,000\Leftrightarrow 22\pi r+15\pi r^2\le 230$\\
		\hspace*{9.8cm}$\Leftrightarrow-3{,}061\le r\le 1{,}594 $.\\
		Do đó, diện tích đáy lớn nhất thỏa mãn $22\pi r+15\pi {r^2}\le 230$ khi $r\approx 1{,}594$.\\
		Lại có $h=\sqrt{{\ell^2}-{r^2}}\approx \sqrt{{{\left(2{,}2\right)}^2}-{{\left(1{,}594\right)}^2}}\approx 1,52$.\\
		Vậy chiều cao của lều là $h\approx 1{,}52$ ($m$).
	}
\end{ex}

\begin{ex}%[TT-THPT-HamRong-NH25-26 ]%[GV soạn: Nguyễn Quang Hiệp GVPB2 Đỗ Minh Vũ ]%[2H2V1-4]
	Một giá treo vật được lắp đặt vào góc tường tại ba điểm $C$, $O$ và $B$ bằng các bản lề cố định. Cấu tạo giá treo gồm ba thanh $AB$, $AC$ và $AO$ hội tụ tại điểm $A$ (khối lượng các thanh không đáng kể). Biết các thanh $AB$, $AC$ có chiều dài bằng nhau và vuông góc nhau tại $A$. Thanh $OA$ đóng vai trò thanh chống, chịu được lực nén tối đa là $150$ N. Khi treo vật có trọng lượng $P$ vào điểm $A$, lực căng trên hai thanh $AC$ và $AB$ có độ lớn luôn bằng một phần ba lực nén trên thanh $OA$. Trọng lượng lớn nhất có thể treo vào đầu $A$ của giá treo bằng bao nhiêu để hệ thống vẫn an toàn? (Kết quả làm tròn đến hàng đơn vị).
	\begin{center}
		\begin{tikzpicture}[>=stealth, line join=round, line cap=round, font=\footnotesize, scale=1.2]
			
			% Định nghĩa màu sắc cho thanh cứng
			\definecolor{rodcolor}{RGB}{205, 145, 55}
			
			% Định nghĩa các tọa độ chính
			\coordinate (O) at (0,0);
			\coordinate (A) at (2.8, 2.0);
			\coordinate (C) at (-2, 2.0);
			\coordinate (B) at (4.2, 3.5);
			\coordinate (D) at (0, 3.5); % Điểm giao của các nét đứt trên trục Z
			
			% 1. Vẽ các trục tọa độ
			\draw[thick,->] (O) -- (5,0) node[above left] {\textbf{y}}; % Trục Y
			\draw[thick,->] (O) -- (0,4.5)node[above left] {\textbf{z}}; % Trục Z (thẳng đứng).
			\draw[thick,->] (O) -- (-2.5,-1.5)node[above left] {\textbf{x}}; % Trục X (xiên)
			
			% 2. Vẽ các đường nét đứt (mặt phẳng chiếu)
			\draw[dashed, thick] (C) -- (D) -- (B);
			\draw[dashed, thick] (A) -- (D);
			
			% 3. Vẽ các thanh cứng (dùng nét đôi để tạo cảm giác hình khối 3D)
			\tikzstyle{rod}=[line width=4pt, rodcolor]
			\tikzstyle{rodborder}=[line width=5pt, black]
			
			% Thanh OA
			\draw[rodborder] (O) -- (A);
			\draw[rod] (O) -- (A);
			% Thanh CA
			\draw[rodborder] (C) -- (A);
			\draw[rod] (C) -- (A);
			% Thanh AB
			\draw[rodborder] (A) -- (B);
			\draw[rod] (A) -- (B);
			
			% 4. Vẽ dây treo và vật nặng (màu xanh lá)
			\draw[thick] (A) -- ([yshift=-1.2cm]A);
			\filldraw[fill=green!40!black, draw=black, thick] 
			([xshift=-0.4cm, yshift=-1.2cm]A) rectangle ([xshift=0.4cm, yshift=-2.5cm]A);
			
			\coordinate (W1) at ([xshift=-3pt, yshift=0.7cm]C);
			\coordinate (W2) at ([xshift=-3pt, yshift=-0.7cm]C);
			\draw[thick] (W1) -- (W2);
			\foreach \i in {0, 0.2, 0.4, 0.6, 0.8, 1.0, 1.2, 1.4} {
				\draw[thick] ([yshift=-\i cm]W1) -- ++(-0.25,-0.2);
			}
			
			% Ngàm tại B (Trần ngang)
			\coordinate (C1) at ([yshift=3pt, xshift=-0.7cm]B);
			\coordinate (C2) at ([yshift=3pt, xshift=0.7cm]B);
			\draw[thick] (C1) -- (C2);
			\foreach \i in {0, 0.2, 0.4, 0.6, 0.8, 1.0, 1.2, 1.4} {
				\draw[thick] ([xshift=\i cm]C1) -- ++(0.2, 0.25);
			}
			\foreach \point/\goc in {A/-60,C/75,B/-60,O/130}{
				\draw[fill=white](\point)circle(2.8pt)+(\goc:2mm)node[scale=1]{$\point$};
			}
			
		\end{tikzpicture}
	\end{center}
	\par\shortans[oly]{132}
	\loigiai{
		\begin{center}
			\begin{tikzpicture}[>=stealth, line join=round, line cap=round, font=\footnotesize, scale=1]
				
				% =========================================================
				% 1. ĐỊNH NGHĨA PIC KHỐI HỘP ĐỘNG
				% =========================================================
				\tikzset{
					pics/hinhHopLuc/.style n args={8}{
						code={
							% Thiết lập độ dài và góc cho 3 chiều không gian (3 véc-tơ)
							\tikzset{
								declare function={
									lenFtwo=3.5; gocFtwo=180;   % Trục véc-tơ F2
									lenFone=2.12; gocFone=55;   % Trục véc-tơ F1
									lenP=3.5; gocP=-90;         % Trục véc-tơ P
								}
							}
							% Tính toán 4 đỉnh của mặt trên (gốc A)
							\path
							(0,0) coordinate (#1) % Đỉnh gốc A
							+(gocFtwo:lenFtwo) coordinate (#2) % C'
							+(gocFone:lenFone) coordinate (#4) % B'
							+(gocP:lenP) coordinate (#5)       % P
							($(#2)+(#4)-(#1)$) coordinate (#3) % TopBackLeft
							;
							% Tịnh tiến mặt trên xuống mặt đáy theo véc-tơ P
							\foreach \pone/\pname in {#2/#6, #3/#7, #4/#8}{
								\path ($(\pone)+(#5)-(#1)$) coordinate (\pname);
							}
							
							% Vẽ các cạnh khuất (hội tụ tại đỉnh đối diện A là O')
							\foreach \pointo/\pointt in {#7/#3, #7/#6, #7/#8}{
								\draw[dashed] (\pointo)--(\pointt);
							}
							
							% Vẽ các cạnh thấy (không vẽ các cạnh từ A để vẽ véc-tơ đè lên sau)
							\foreach \pointo/\pointt in {#2/#3, #3/#4, #6/#5, #5/#8, #2/#6, #4/#8}{
								\draw (\pointo)--(\pointt);
							}
						}
					}
				}
				
				% =========================================================
				% 2. GỌI PIC VÀ TRUYỀN TÊN CÁC ĐIỂM
				% =========================================================
				% Truyền vào 8 node tương ứng: A, C', TopBackLeft, B', P, BottomFrontLeft, O', BottomBackRight
				\path (2,1) pic{hinhHopLuc={A}{C'}{TopBackLeft}{B'}{P_vec}{BottomFrontLeft}{O'}{BottomBackRight}};
				
				% =========================================================
				% 3. XÁC ĐỊNH CÁC TIA KÉO DÀI
				% =========================================================
				\path 
				($(A)!1.3!(C')$) coordinate (C)
				($(A)!1.3!(B')$) coordinate (B)
				($(A)!1.6!(O')$) coordinate (O);
				
				\draw (C') -- (C);
				\draw (B') -- (B);
				\draw[dashed] (O') -- (O);
				
				% =========================================================
				% 4. VẼ CÁC VÉC-TƠ LỰC
				% =========================================================
				\draw[->, thick] (A) -- (C') node[midway, above] {$\overrightarrow{F}_2$};
				\draw[->, thick] (A) -- (B') node[midway, right] {$\overrightarrow{F}_1$};
				\draw[->, thick] (A) -- (P_vec) node[midway, above right] {$\overrightarrow{P}$};
				\draw[->, thick, dashed] (A) -- (O') node[pos=0.6, below right] {$\overrightarrow{F}_3$};
				
				% =========================================================
				% 5. ĐÁNH DẤU ĐIỂM VÀ TÊN DÙNG VÒNG LẶP (\foreach)
				% =========================================================
				% Cấu trúc: Tọa_độ / Tên_hiển_thị / Góc_đặt_nhãn
				\foreach \point/\name/\goc in {
					A/A/135, B'/B'/0, C'/C'/90, 
					P_vec/P/0, O'/O'/135, 
					B/B/90, C/C/90, O/O/135
				}{
					\fill (\point) circle (1pt);
					\path (\point) ++(\goc:3mm) node {$\name$};
				}
				
			\end{tikzpicture}
		\end{center}
		Gọi lực căng lên hai thanh $AB$, $AC$ và lực nén lên thanh $OA$ lần lượt là $\overrightarrow{F}_1$, $\overrightarrow{F}_2$, $\overrightarrow{F}_3$.\\
		Đặt $\overrightarrow{F}_1 = \overrightarrow{AB'}$, $\overrightarrow{F}_2 = \overrightarrow{AC'}$, $\overrightarrow{P} = \overrightarrow{AP}$.\\
		Để hệ thống vẫn an toàn thì
		\[\overrightarrow{F}_1 + \overrightarrow{F}_2 + \overrightarrow{F}_3 + \overrightarrow{P} = \overrightarrow{0} \Leftrightarrow -\overrightarrow{F}_3 = \overrightarrow{F}_1 + \overrightarrow{F}_2 + \overrightarrow{P} \Leftrightarrow -\overrightarrow{F}_3 = \overrightarrow{AB'} + \overrightarrow{AC'} + \overrightarrow{AP} \quad (*).\]
		Từ giả thiết ta thấy $AB'$, $AC'$, $AP$ lần lượt là ba cạnh của hình hộp chữ nhật, từ $(*)$ và quy tắc hình hộp suy ra $-\overrightarrow{F}_3 = \overrightarrow{AO'}$ với $O' \in OA$ và $AO'$ là đường chéo của hình hộp chữ nhật trên.\\
		Khi đó ta có
		\[\left| -\overrightarrow{F}_3 \right| = \sqrt{AB'^2 + AC'^2 + AP^2} \Rightarrow \left| \overrightarrow{F}_3 \right|^2 = AB'^2 + AC'^2 + AP^2 \Rightarrow AP^2 = \left| \overrightarrow{F}_3 \right|^2 - AB'^2 - AC'^2.\]
		Theo giả thiết, ta có $\left| \overrightarrow{F}_1 \right| = \left| \overrightarrow{F}_2 \right| = \dfrac{1}{3}\left| \overrightarrow{F}_3 \right| \Rightarrow AB' = AC' = \dfrac{1}{3}\left| \overrightarrow{F}_3 \right| \Rightarrow AB'^2 = AC'^2 = \dfrac{1}{9}\left| \overrightarrow{F}_3 \right|^2$.\\
		Suy ra
		\[AP^2 = \left| \overrightarrow{F}_3 \right|^2 - \dfrac{1}{9}\left| \overrightarrow{F}_3 \right|^2 - \dfrac{1}{9}\left| \overrightarrow{F}_3 \right|^2 = \dfrac{7}{9}\left| \overrightarrow{F}_3 \right|^2.\]
		Do vậy $\left| \overrightarrow{P} \right|^2 = \dfrac{7}{9}\left| \overrightarrow{F}_3 \right|^2 \Rightarrow \left| \overrightarrow{P} \right| = \dfrac{\sqrt{7}}{3}\left| \overrightarrow{F}_3 \right| \le \dfrac{\sqrt{7}}{3} \cdot 150 \approx 132$.\\
		Vậy trọng lượng lớn nhất có thể treo vào đầu $A$ của giá treo là $132$ N.
	}
\end{ex}

\begin{ex}%[2H2V2-5]%[Dự án đề thi thử đợt 3-Nguyễn Đức Lợi]%[SGD tỉnh Lạng Sơn]%[GVPB - ThanhTa]
	Cho hình lăng trụ đứng $ABC.A'B'C'$ có đáy $ABC$ là tam giác vuông tại $A$, biết $AB = 2$, $AC = 3$, $AA' = 4$. Gọi $M$ là trung điểm của $AB$. Khoảng cách giữa hai đường thẳng $A'B$ và $CM$ bằng bao nhiêu (\textit{kết quả làm tròn đến hàng phần trăm})?
	\par
	\shortans{0,86}
	\loigiai{
		\begin{center}
			\begin{tikzpicture}[scale=0.8, font=\footnotesize, line join=round, line cap=round, >=stealth]
				% Định nghĩa các điểm đáy dưới
				\coordinate (A) at (0,0);
				\coordinate (B) at (4,0);
				\coordinate (C) at (1,-2);
				% Định nghĩa các điểm đáy trên
				\coordinate (A') at ($(A)+(0,4)$);
				\coordinate (B') at ($(B)+(0,4)$);
				\coordinate (C') at ($(C)+(0,4)$);
				\coordinate (z) at ($(A)!1.2!(A')$);
				\coordinate (x) at ($(A)!1.2!(B)$);
				\coordinate (y) at ($(A)!1.2!(C)$);
				% Trung điểm M của AB
				\coordinate (M) at ($(A)!0.5!(B)$);
				\draw[->] (A') -- (z) node[left] {$z$};
				\draw[->] (B) -- (x) node[above] {$x$};
				\draw[->] (C) -- (y) node[right] {$y$};
				% Vẽ các cạnh khuất
				\draw[dashed] (A)--(B)  (C) -- (M);
				% Vẽ các cạnh thấy
				\draw (C) -- (C') -- (A') -- (B') -- (B) -- (C) (C') -- (B') (A) -- (A') (C) -- (A);
				% Vẽ hai đường thẳng cần tính khoảng cách
				\draw[thick, blue,dashed] (A') -- (B);
				\draw[thick, blue, dashed] (C) -- (M);
				
				% Tên các điểm
				\foreach \p/\pos in {A/left, B/below, C/left, A'/left, B'/right, C'/left, M/above}
				\fill (\p) circle (1.5pt) node[\pos] {$\p$};
				
				% Ký hiệu vuông góc tại A
				\draw (0,0.2) -- (0.2,0.2) -- (0.2,0);
			\end{tikzpicture}
		\end{center}
		Chọn hệ trục tọa độ $Oxyz$ sao cho $A \equiv O(0;0;0)$, $B(2;0;0)$, $C(0;3;0)$ và $A'(0;0;4)$.\\
		Vì $M$ là trung điểm của $AB$ nên $M(1;0;0)$.\\
		Ta có $\vec{A'B} = (2;0;-4)$ và $\vec{CM} = (1;-3;0)$.\\
		Vectơ chỉ phương của $A'B$ là $\vec{u} = (1;0;-2)$, của $CM$ là $\vec{v} = (1;-3;0)$.\\
		Suy ra $\vec{n} =[\vec{v}, \vec{u}] = (6; 2; 3)$.\\
		Lấy điểm $M(1;0;0)$ thuộc $CM$ và $B(2;0;0)$ thuộc $A'B$, ta có $\vec{MB} = (1;0;0)$.\\
		Khoảng cách giữa hai đường thẳng là
		\[ d(A'B, CM) = \dfrac{|\vec{n} \cdot \vec{MB}|}{|\vec{n}|} = \dfrac{|6 \cdot 1 + 2 \cdot 0 + 3 \cdot 0|}{\sqrt{6^2 + 2^2 + 3^2}} = \dfrac{6}{7} \approx 0{,}86. \]
	}
\end{ex}

\begin{ex}%[2H2V2-6]%[Dự án đề TT Toán NH25-26-Dot 3- Hải Phụng]%[Cụm 13-Hải Phòng]%[GVPB - Viet Hoang Pham]
Một ngôi nhà gồm hai phần. Phần thân nhà dạng hình hộp chứ nhật $ABCD.OMNK$ có chiều dài $1\,200$ cm, chiều rộng $900$ cm, chiều cao $450$ cm. Phần mái nhà $S.ABCD$ có dạng một hình chóp tứ giác có các cạnh bên bằng nhau. Biết $ABCD$ là hình chữ nhật và có chiều cao của ngôi nhà $\left(\text{bằng khoảng cách từ }S\text{ đến }(OMNK)\right)$ bằng $600$ cm. Chọn hệ trục tọa độ $Oxyz$ sao cho $M$ thuộc tia $Ox$, $K$ thuộc tia $Oy$, $A$ thuộc tia $Oz$ (như hình vẽ), (mỗi đơn vị trên hệ trục ứng với $1$ m).
\begin{center}
\begin{tikzpicture}[line join=round,line cap=round, font=\footnotesize, scale=1,>=stealth]
\coordinate[label=left:$M$] (M) at (0,0);
\coordinate[label=above left:$O$] (O) at (1.2,.8);
\coordinate[label=below right:$N$] (N) at (4,0);
\coordinate[label=above left:$K$] (K) at ($(N)-(M)+(O)$);
\coordinate[label=below right:$A$] (A) at ($(O)+(90:3)$);
\coordinate[label=left:$B$] (B) at ($(M)-(O)+(A)$); 
\coordinate[label=below right:$C$] (C) at ($(N)-(O)+(A)$);
\coordinate[label=above right:$D$] (D) at ($(K)-(O)+(A)$); 
\coordinate[label=left:$x$] (x) at ($(O)!1.7!(M)$);
\coordinate[label=below:$y$] (y) at ($(O)!1.3!(K)$);
\coordinate[label=left:$z$] (z) at ($(O)!1.8!(A)$);
\coordinate (s) at ($(A)!0.5!(C)$);
\coordinate[label=above:$S$] (S) at ($(s)+(0,2)$);

\draw (B)--(M)--(N)--(K)--(D) (B)--(C)--(D) (N)--(C) (S)--(B) (S)--(C) (S)--(D);
\draw[dashed] (O)--(K) (A)--(O)--(M) (S)--(A)--(B) (A)--(D);
\draw[->] (K)--(y);
\draw[->] (M)--(x);
\draw[->] (A)--(z);

\draw [<->] ($(M)+(0,-0.3)$)--($(N)+(0,-0.3)$) node[below,pos=0.5]{$1200$};
\draw [<->] ($(N)+(0.3,0)$)--($(K)+(0.3,0)$) node[below right,pos=0.5]{$900$}; 
\draw [<->] ($(K)+(0.3,0)$)--($(D)+(0.3,0)$) node[below right,pos=0.5]{$450$}; 

\foreach \x in {A,B,C,D,O,M,N,K,S} 
 \draw[fill=black] (\x) circle (1.5pt); 
\end{tikzpicture}
\end{center}
Gọi tọa độ điểm $S$ là $S(a;b;c)$. Hãy tính $a-b+c$.
\shortans[oly]{4,5}
\loigiai{
\immini[thm]{
Đổi $1\,200 \text{ cm} = 12\text{ m}$; $900\text{ cm} = 9\text{ m}$; $600\text{ cm} = 6\text{ m}$.\\  
Gọi $I$ là hình chiếu vuông góc của điểm $S$ lên mặt phẳng $(OMNK)$. Khi đó $I$ là tâm hình chữ nhật $OMNK$. \\
Do đó tọa độ của điểm $I$ là $(4{,}5;6;6)$. \\
Suy ra tọa độ của điểm $S$ là $S(4{,}5;6;6)$.\\
Do đó $a-b+c=4{,}5 - 6 + 6 = 4{,}5$.	
}{
\begin{tikzpicture}[line join=round,line cap=round, font=\footnotesize, scale=1,>=stealth]
\coordinate[label=left:$M$] (M) at (0,0);
\coordinate[label=above left:$O$] (O) at (1.2,.8);
\coordinate[label=below right:$N$] (N) at (4,0);
\coordinate[label=above left:$K$] (K) at ($(N)-(M)+(O)$);
\coordinate[label=below right:$A$] (A) at ($(O)+(90:3)$);
\coordinate[label=left:$B$] (B) at ($(M)-(O)+(A)$); 
\coordinate[label=below right:$C$] (C) at ($(N)-(O)+(A)$);
\coordinate[label=above right:$D$] (D) at ($(K)-(O)+(A)$); 
\coordinate[label=left:$x$] (x) at ($(O)!1.7!(M)$);
\coordinate[label=below:$y$] (y) at ($(O)!1.3!(K)$);
\coordinate[label=left:$z$] (z) at ($(O)!1.8!(A)$);
\coordinate (s) at ($(A)!0.5!(C)$);
\coordinate[label=above:$S$] (S) at ($(s)+(0,2)$);
\coordinate[label=right:$I$] (I) at ($(O)!0.5!(N)$);

\draw (B)--(M)--(N)--(K)--(D) (B)--(C)--(D) (N)--(C) (S)--(B) (S)--(C) (S)--(D);
\draw[dashed] (O)--(K) (A)--(O)--(M) (S)--(A)--(B) (A)--(D) (S)--(I);
\draw[->] (K)--(y);
\draw[->] (M)--(x);
\draw[->] (A)--(z);

\draw [<->] ($(M)+(0,-0.3)$)--($(N)+(0,-0.3)$) node[below,pos=0.5]{$1200$};
\draw [<->] ($(N)+(0.3,0)$)--($(K)+(0.3,0)$) node[below right,pos=0.5]{$900$}; 
\draw [<->] ($(K)+(0.3,0)$)--($(D)+(0.3,0)$) node[below right,pos=0.5]{$450$}; 

\foreach \x in {A,B,C,D,O,M,N,K,S,I} 
 \draw[fill=black] (\x) circle (1.5pt); 
\end{tikzpicture}
}
}
\end{ex}

\begin{ex}%[2H2V2-6]%[Dự án đề TT Toán NH25-26-Dot 3- Hải Phụng]%[Cụm 13-Hải Phòng]%[GVPB - Viet Hoang Pham]
Hai chiếc khinh khí cầu bay lên từ cùng một địa điểm trong không gian. Để theo dõi hành trình của hai khinh khí cầu, người ta chọn hệ trục toạ độ $Oxyz$  với gốc $O$  đặt tại điểm xuất phát của hai khinh khí cầu, mặt phẳng $(Oxy)$ trùng với mặt đất với trục $Ox$  hướng về phía Nam, trục $Oy$  hướng về phía Đông và trục $Oz$  hướng thẳng lên trời (đơn vị đo lấy theo kilômét). Vào lúc $9$ giờ sáng, chiếc thứ nhất nằm cách điểm xuất phát $3$ km về phía Đông và $2$ km về phía Nam, đồng thời cách mặt đất $0{,}5$ km; chiếc thứ hai nằm cách điểm xuất phát $1$ km về phía Bắc và $1$ km về phía Tây, đồng thời cách mặt đất $0{,}3$ km. Cùng thời điểm đó, một người đứng
trên mặt đất và nhìn thấy hai khinh khí cầu nói trên. Biết rằng, so với các vị trí quan sát khác trên mặt đất, vị trí người đó đứng có tổng khoảng cách đến hai khinh khí cầu là nhỏ nhất. Vị trí người quan sát đứng lúc đó là $M(a;b;c)$. Tính tổng $32a+60b+5c$.

\shortans[oly]{$34$}
\loigiai{ \begin{center}
		\begin{tikzpicture}[line join=round, line cap=round, >=stealth, font=\footnotesize, scale=1]
			\def\x{3}
			\def\y{5}
			\def\z{4}
			\path (0:0) coordinate (O)
			++(0:\y) coordinate (y)
			(O)++(-140:\x) coordinate (x)
			(O)++(90:\z) coordinate (z)
			($(O)!0.6!(x)$) coordinate (X)
			($(O)!0.8!(y)$) coordinate (Y)
			($(O)!0.6!(z)$) coordinate (Z)
			($(X)+(Y)-(O)$) coordinate (XY)
			($(XY)+(Z)-(O)$) coordinate (A)
			($(O)!-0.4!(x)$) coordinate (X')
			($(O)!-0.25!(y)$) coordinate (Y')
			($(O)!0.5!(z)$) coordinate (Z')
			($(X')+(Y')-(O)$) coordinate (X'Y')
			($(X'Y')+(Z')-(O)$) coordinate (B)
			;
			%Vẽ các trục
			\draw[->,thick] (O)--(x) node[right]{$x$ (Nam)};
			\draw[->,thick] (O)--(y) node[above]{$y$ (Đông)};
			\draw[->,thick] (O)--(z) node[left]{$z$};
			%Các hướng
			\draw[dashed,thick] (O)--($(O)!-0.9!(y)$);
			\draw[dashed,thick] (X)--(XY)--(Y) (XY)--(A);
			\draw[dashed,thick] (X') --(X'Y')--(Y') (X'Y')--(B);
			\draw[dashed,thick] (O)--($(O)!-1!(x)$);
			
			\draw ($(X)!0.5!(XY)$) node[below]{$3$ km} ($(XY)!0.5!(Y)$) node[right]{$2$ km} ($(X')!0.5!(O)$) node[right]{$1$ km} ($(Y')!0.3!(O)$) node[below]{$1$ km} ($(XY)!0.7!(A)$) node[right]{0,5 km} ($(X'Y')!0.7!(B)$) node[left]{0,3 km} ; 
			\draw ($(Y')+(-3,0)$) node[below]{(Tây)} ($(O)!-1.7!(X)$) node[right]{(Bắc)};
			% Máy bay tại A:
			\node[font=\Huge, rotate=0] at (A) {\faBullseye};
			%% Máy bay tại B:
			\node[font=\Huge, rotate=120] at (B) {\faBullseye};
		\end{tikzpicture}
	\end{center}
Vào lúc 9 giờ sáng:
\begin{itemize}
\item Chiếc khinh khí cầu thứ nhất ở vị trí $A(2;3;0{,}5)$.
\item Chiếc khinh khí cầu thứ hai ở vị trí $B(-1;-1;0{,}3)$.
\end{itemize}
Mặt đất là mặt phẳng $(Oxy)$, hai điểm $A$, $B$ nằm về cùng phía so với mặt phẳng $(Oxy)$. \\
Gọi  $A'$ là điểm đối xứng với $A$ qua mặt phẳng $(Oxy)$. 
Khi đó $A'(2;3;-0,5)$. \\
Tổng khoảng cách từ vị trí người đó đứng trên mặt đất đến hai khinh khí cầu là
$$T=MA+MB=MA'+MB \geq A'B.$$
Vậy $T$ đạt giá trị nhỏ nhất bằng $A'B$ khi và chỉ khi $A'$, $M$, $B$ thẳng hàng. \\
Do $M(a;b;c) \in (Oxy) $ nên $c=0$. \\
Ta có $\vec{A'M} = (a-2;b-3;0,5)$ và $\vec{A'B} = (-3;-4;0,8)$. \\
Ba điểm $A',M,B$ thẳng hàng khi và chỉ khi hai vectơ $\vec{A'M}$ và $\vec{A'B}$ cùng phương. Khi đó 
$$\dfrac{a-2}{-3} = \dfrac{b-3}{-4} = \dfrac{0,5}{0,8} \Rightarrow a=\dfrac{1}{8}; b=\dfrac{1}{2}.$$
Vậy $32a+60b+5c = 34$
}
\end{ex}

\begin{ex}%[2H2V2-6]%[TN-THPT-NguyenTrungThien-HaTinh-NH25-26]%[Loc Do]%[PB: Phan Anh]
	Một cuộc diễn tập phòng không được đặt trong không gian với hệ tọa độ $Oxyz$ với $1$ đơn vị độ dài trên các trục tính bằng $1$ km trên thực tế. Một bệ phóng tên lửa phòng không được đặt tại vị trí $O(0;0;0)$ và tên lửa bay theo một đường thẳng với vận tốc không đổi là $600$ m/s. Một máy bay không người lái bay theo một đường thẳng với vận tốc không đổi là $1\,080$ km/h và bay theo hướng của vectơ $\overrightarrow{u}=(-3;4;0)$. Khi phát hiện máy bay không người lái tại vị trí $A(6;-20;1)$ thì tên lửa khai hỏa, rời bệ phóng và sau đó bắn hạ được mục tiêu. Hỏi khoảng cách từ bệ phóng tên lửa đến vị trí máy bay không người bị bắn hạ bằng bao nhiêu km (\textit{kết quả làm tròn đến hàng phần chục, giả sử cả máy bay không người lái và tên lửa đều không chịu tác động của trọng lực hay lực cản không khí})?
	\shortans{14,4}
	\loigiai{
		Đổi $600$ m/s $=2160$ km/h.\\
		Phương trình đường bay của máy bay không người lái là $d\colon \heva{&x=6-3 t \\ &y=-20+4t \\ &z=1}$, ($t \in \mathbb{R}$, $t>0$).\\
		Gọi điểm máy bay bị bắn trúng là $M$. Vì $M$ thuộc $d$ nên $M(6-3t ;-20+4t ; 1)$.\\
		Khi máy bay ở vị trí $A$ thì tên lửa khai hỏa và bắn trúng máy bay, do đó thời gian máy bay di chuyển từ $A$ đến $M$ và thời gian tên lửa di chuyển từ $O$ đến $M$ là như nhau. \\	
		Khi đó	
		$$
		\dfrac{AM}{1080}=\dfrac{O M}{2160} \Leftrightarrow \dfrac{\sqrt{(-3 t)^2+(4t)^2+0^2}}{1\,080}=\dfrac{\sqrt{(6-3t)^2+(-20+4t)^2+1^2}}{2\,160}.
		$$	
		Giải phương trình trên nhận được nghiệm thỏa mãn là $t=\dfrac{36}{25}$.\\
		Từ đó suy ra $OM \approx 14{,}4$ (km).
	}
\end{ex}

\begin{ex}%[2H3K1-1]%[Dự án C-TT THPT NH25-26-Đợt 4- ĐỖ CHÍ TÂM]%[THPT Phan Bội Châu-Khánh Hòa]%[GVPB: Phạm Văn Long]
	\immini[thm]{Sân hiên hình chữ nhật của một ngôi nhà là khoảng đất $ABCD$ được lợp mái bằng kính màu để hạn chế ánh sáng đi qua với mái dốc. Các bề mặt bên $ADHE$ và $CGHD$ nằm ở bức tường bên ngoài ngôi nhà. Đặt vào mô hình hệ trục tọa độ như hình vẽ thì ta có $B\left(5;\dfrac{7}{2};0\right)$; $E\left(5;0;2\right)$ và $H\left(0;0;3\right)$. Trên tường nhà có một ngọn đèn đặt tại điểm $L$ cách điểm $D$ một khoảng $6$ m theo phương thẳng đứng. Phần có mái của sân hiên in bóng lên khu vườn bằng phẳng phía trước ngôi nhà dưới ánh đèn tạo thành khoảng đất hạn chế ánh sáng. Tính diện tích khoảng đất đó (\textit{Kết quả làm tròn kết quả đến hàng phần chục}).
		\shortans{$45{,}9$}
	}{\begin{tikzpicture}[>=stealth,line join=round,line cap=round,font=\footnotesize,scale=1, black]
			\path (0,0) coordinate (D) (0,0.7) coordinate (H) (0,2) coordinate (D')
			(2,0) coordinate (C) (3,0) coordinate (C') (0,0.5) coordinate (H')
			(-120:3) coordinate (A) (-120:4) coordinate (A') (0,1.5) coordinate (L);
			\path ($(C)-(D)+(A)$) coordinate (B)
			($(H')-(D)+(A)$) coordinate (E)
			($(H')-(D)+(B)$) coordinate (F)
			($(H)-(D)+(C)$) coordinate (G);
			\fill[gray!30] (E)--(H)--(G)--(F)--cycle;
			\draw (A)--(B)--(C);
			\draw[dashed] (A)--(D)--(C) (D)--(H);
			\draw [->,thick] (A)--(A') node[left] {$x$};
			\draw [->,thick] (C)--(C') node[above] {$y$};
			\draw [->,thick] (H)--(D') node[right] {$z$};
			\draw (E)--(H)--(G)--(F)--cycle (A)--(E) (B)--(F) (C)--(G);
			\foreach \t/\g in {A/180,B/-90,C/-40,D/-80,E/180,F/120,G/0,H/30,L/180}{
				\fill (\t) circle (1pt) node[shift={(\g:7pt)},font=\scriptsize]{$ \t $};
			}
	\end{tikzpicture}}
	\loigiai{\begin{center}
			\begin{tikzpicture}[>=stealth,line join=round,line cap=round,font=\footnotesize,scale=1]
				\path (0,0) coordinate (D) (0,0.7) coordinate (H) (0,2) coordinate (D')
				(2,0) coordinate (C) (5,0) coordinate (C') (0,0.5) coordinate (H')
				(-120:3) coordinate (A) (-120:5) coordinate (A') (0,1.5) coordinate (L);
				\path ($(C)-(D)+(A)$) coordinate (B)
				($(H')-(D)+(A)$) coordinate (E)
				($(H')-(D)+(B)$) coordinate (F)
				($(H)-(D)+(C)$) coordinate (G);
				\path (intersection of L--E and A--A') coordinate (E');
				\path (intersection of L--G and C--C') coordinate (G');
				\path ($(C)-(D)+(E')$) coordinate (M);
				\path (intersection of L--F and E'--M) coordinate (F');
				\fill[gray!40, opacity=0.6] (E)--(H)--(G)--(F)--cycle;
				\fill[pattern=dots, pattern color=red!30] (E')--(A)--(B)--(C)--(G')--(F')--cycle;
				\draw (A)--(B)--(C);
				\draw[dashed] (A)--(D)--(C) (D)--(H) (E)--(H');
				\draw [->,thick] (A)--(A') node[left] {$x$};
				\draw [->,thick] (C)--(4.2,0) node[above] {$y$};
				\draw [->,thick] (H)--(D') node[right] {$z$};
				\draw (E)--(H)--(G)--(F)--cycle (A)--(E) (B)--(F) (C)--(G);
				\draw (L)--(E')--(F')--(G')--cycle (L)--(F');
				\foreach \t/\g in {A/-90,C/-40,F/120,G/0}{
					\fill (\t) circle (1pt) node[shift={(\g:7pt)},font=\scriptsize]{$ \t $};
				}
				\foreach \t/\g in {B/-90,E/180,H/30, L/180}{
					\fill[black] (\t) circle (1.5pt) node[shift={(\g:7pt)},font=\scriptsize]{$ \t $};
				}
				\foreach \t/\g in {D/-80,E'/180,G'/90, F'/-90}{
					\fill[black] (\t) circle (1.5pt) node[shift={(\g:7pt)},font=\scriptsize]{$ \t $};
				}
			\end{tikzpicture}
		\end{center}
		$B\left(5;3{,}5;0\right)$, $E\left(5;0;2\right)$, $H\left(0;0;3\right)$, $L(6;0;0)$. Suy ra $C(0;3{,5};0)$, $F(5;3{,}5; 2)$, $A(5; 0;0)$, $G(0; 3{,}5; 3)$.\\
		$GC$ là đường trung bình của tam giác $LDG'$. Suy ra $G'\left(0;7;0\right)$.\\
		Ta có $\dfrac{E'A}{E'D}=\dfrac{EA}{LD}=\dfrac{2}{6}=\dfrac{1}{3}\Rightarrow E'A=\dfrac{1}{3}E'D$\\
		Suy ra $\vec{E'A}=\dfrac{1}{3}\vec{E'D}\Rightarrow E'\left(7{,}5;0;0\right)$.\\
		Mặt khác $\dfrac{EF}{E'F'}=\dfrac{LE}{LE'}=\dfrac{2}{3}\Rightarrow \vec{E'F'}=1{,}5\vec{EF}$\\
		Suy ra $F'(7{,}5; 5{,}25; 0)$\\
		Diện tích khoảng đất thiếu ánh sáng là diện tích hình thang vuông $DG'F'E'$.\\
		Ta có $DG'=7$, $E'F'=5{,}25$, $DE'=7{,}5$.\\
		\hspace*{2cm}$S_{DG'F'E'}=\dfrac{\left(DG'+E'F'\right)\cdot DE'}{2}=\dfrac{\left(7+5{,}25\right)\cdot 7{,}5}{2}=45{,}9$.
	}
\end{ex}

\begin{ex}%[Dự án đề TT Toán NH25-26-Dot 4- Hải Phụng]%[THCS-THPT LÊ THÁNH TÔNG - TP.HCM]%[GVPB - Viet Hoang Pham]%[2H5C1-7]
		Cận kề ngày Tết cổ truyền, anh Nghĩa dự định trang trí hệ thống đèn LED cho khoảng sân nhà hình chữ nhật $MNPQ$ có chiều dài $MQ=12$ m, $MN=8$ m. Để cung cấp điện cho hệ thống đèn LED, anh Nghĩa đi dây điện theo trình tự: $A \to B \to C \to O \to D \to E \to F \to H$ (như hình vẽ bên). Trong đó:
		\begin{itemize}
			\item Điểm $A$ nằm trên cột cổng cao $3$ m, vị trí chân cột cách cây Nêu $3$ m và chân cột nằm trên đoạn $MQ$.
			\item Đoạn $BC$ nằm trên bức tường hình chữ nhật $PQIJ$ có $PJ=4$ m (bức tường vuông góc với mặt đất), độ dài $BC=3$ m và $BC \parallel IJ$.
			\item Điểm $O$ nằm trên mái che hình chữ nhật $IJUV$ có $JU=10$ m (mái che song song với mặt đất).
			\item Cột đèn đặt tại góc $N$ và điểm $D$ nằm trên cột đèn cao $3$ m.
			\item Điểm $E, F$ nằm trên cây Nêu sao cho $EF=2$ m. Cây Nêu $MG$ vuông góc với mặt đất tại $M$, chiều cao $MG=10$ m.
			\item Điểm $H$ cách đỉnh $G$ một khoảng $HG=0{,}5$ m ($HG$ song song với mặt đất).
		\end{itemize}
		Tính tổng độ dài dây điện ngắn nhất mà anh Nghĩa cần sử dụng để hoàn thành hệ thống đèn trang trí này (Đơn vị: mét, làm tròn kết quả đến hàng phần chục).
\begin{center}
	\begin{tikzpicture}[scale=0.8, line join=round, line cap=round, >=stealth]
		% --- ĐỊNH NGHĨA CÁC ĐIỂM CƠ SỞ (PHỐI CẢNH) ---
		\coordinate (M) at (0,0);
		\coordinate (Q) at (6,0);
		\coordinate (N) at (-2.5,-1.8);
		\coordinate (P) at ($(Q)+(N)-(M)$);
		% --- VẼ SÂN MNPQ ---
		\draw[thick] (M) -- (Q) -- (P) -- (N) -- cycle;
		\node at ($(M)!0.5!(N)$)[sloped, above, scale=0.8]{$8$ m};
		\node at ($(N)!0.5!(P)$)[below, scale=0.8]{$12$ m};
		% --- VẼ BỨC TƯỜNG PQIJ ---
		\coordinate (I) at ($(Q)+(0,4)$);
		\coordinate (J) at ($(P)+(0,4)$);
		\draw[fill=gray!10] (P) -- (Q) -- (I) -- (J) -- cycle;
		% Vẽ vân gạch cho tường
		\begin{scope}
			\clip (P) -- (Q) -- (I) -- (J) -- cycle;
			\foreach \i in {0,0.4,...,4} \draw[gray!40] ($(P)+(0,\i)$) -- ($(Q)+(0,\i)$);
			\foreach \i in {0,0.5,...,6} \draw[gray!40] ($(P)+(\i,0)$) -- ($(J)+(\i,0)$);
		\end{scope}
		\node at ($(P)!0.5!(J)$)[left, scale=0.8]{$4$ m};
		% --- VẼ CÂY NÊU MG ---
		\coordinate (G) at (0,7);
		\draw[line width=2pt, green!50!black] (M) -- (G);
		\foreach \i in {1,2,3,4,5,6} \draw[green!20!black] (-0.1,\i) -- (0.1,\i+0.1); % Đốt tre
		% --- VẼ MÁI CHE IJUV ---
		\coordinate (U) at ($(J)+(-4.5,0)$);
		\coordinate (V) at ($(I)+(-4.5,0)$);
		\draw[fill=red!10, opacity=0.8] (U) -- (J) -- (I) -- (V) -- cycle;
		\node[above] at (I) {$I$};
		\node[above] at (J) {$J$};
		\node[left] at (U) {$U$};
		\node[above] at (V) {$V$};
		% --- VẼ CỘT ĐÈN TẠI N ---
		\begin{scope}[shift={(N)}]
			\draw[fill=black] (-0.2,0) rectangle (0.2,0.2); % Chân đế
			\draw[fill=black] (-0.05,0.2) rectangle (0.05,2.5); % Thân cột
			\draw[fill=black] (-0.3,2.5) -- (0.3,2.5) -- (0.4,3.2) -- (-0.4,3.2) -- cycle; % Đầu đèn
			\draw[fill=yellow!60] (-0.25,2.6) -- (0.25,2.6) -- (0.3,3) -- (-0.3,3) -- cycle; % Kính đèn
			\coordinate (D) at (0,3); % Điểm D trên đỉnh cột đèn
		\end{scope}
		% --- VẼ CÁC ĐIỂM CHỐT A, B, C, O, E, F, H ---
		\coordinate (A) at (1.5,1.4); % Ước lượng vị trí cột cổng
		\draw[fill=gray!30] (1.3,0) rectangle (1.7,1.4); % Cột cổng A
		\coordinate (B) at ($(I)+(-0.6,-2.3)$);
		\coordinate (C) at ($(J)+(0.5,-1.5)$);
		\coordinate (O) at ($(U)!0.5!(I)$); % Điểm O trên mái
		\coordinate (E) at (0,4.5);
		\coordinate (F) at (0,5.8);
		\coordinate (H) at (-1.2,6.5);
		% --- VẼ LỒNG ĐÈN TẠI H ---
		\draw[thick] (G) -- (H);
		\draw[fill=red, rounded corners=2pt] ($(H)+(-0.3,-1)$)-|($(H)+(0.3,0)$)-|cycle;
		\draw[yellow, line width=1pt] ($(H)+(-0.3,-0.2)$) -- ($(H)+(0.3,-0.2)$);
		\draw[yellow, line width=1pt] ($(H)+(-0.3,-0.8)$) -- ($(H)+(0.3,-0.8)$);
		\node at (H) [above] {H};
		% --- VẼ HỆ THỐNG DÂY ĐIỆN (NÉT ĐỨT VÀ LIỀN) ---
		\draw[red, thick, dashed] (A) -- (B);
		\draw[red, thick, dashed] (B) -- (C) node[midway, right, black, scale=0.7]{$3$ m};
		\draw[red, thick, dashed] (C) -- (O);
		\draw[red, thick, dashed] (O) -- (D);
		\draw[red, thick] (D) -- (E);
		\draw[red, thick] (E) -- (F) node[midway, left, black, scale=0.7]{$2$ m};
		\draw[red, thick] (F) -- (H);
		
		% --- ĐÁNH DẤU CÁC ĐIỂM ---
		\foreach \p in {A,B,C,O,D,E,F,G,M,N,P,Q}
		\filldraw[black] (\p) circle (1pt);
		\node[above] at (A) {$A$};
		\node[right] at (B) {$B$};
		\node[below] at (C) {$C$};
		\node[above] at (O) {$O$};
		\node[left] at (D) {$D$};
		\node[above right] at (E) {$E$};
		\node[right] at (F) {$F$};
		\node[right] at (G) {$G$};
		\node[below] at (N) {$N$};
		\node[below] at (P) {$P$};
		\node[right] at (Q) {$Q$};
		\node[left] at (M) {$M$};
	\end{tikzpicture}
\end{center}
	\shortans{$36{,}5$}
	\loigiai{
		Chọn hệ trục $Oxyz$ sao cho $M \equiv O$, $ M(0;0;0)$, tia $Ox$ đi qua $Q$, tia $Oy$ đi qua $N$.\\			
		Khi đó:	$Q(12;0;0)$, $N(0;8;0)$, trục $O z \perp (MNPQ)$, $G(0;0;10)$,  $P(12;8;0)$.\\			
		Để độ dài dây điện ngắn nhất thì các đoạn dây phải đi theo đoạn thẳng.\\
		Theo giả thiết ta có
		$$T = AB + BC + CO + OD + DE + EF + FH
		= AB + CO + OD + DE + FH + 5.$$
		Ta có:
		$A(3;0;3)$, $D(0;8;3)$, $H\left(0;\dfrac{1}{2};10\right)$,	mặt phẳng $(PQIJ)\colon x = 12$, mặt phẳng $(IJUV)\colon z = 4 $.\\		
		Gọi $A'$ là điểm đối xứng với $A$ qua mặt phẳng $(IJUV) \Rightarrow A'(3;0;5)$. \\			
		Gọi $B'$, $C'$ lần lượt là điểm đối xứng của $B$, $C$ qua $(IJUV)$.\\
		Khi đó: $AB = A'B'$. \\
		Gọi \( A'' \) là điểm đối xứng của \( A' \) qua mặt phẳng $(PQIJ) \Rightarrow A''(21;0;5)$, suy ra $A'B' = A''B'.$\\			
		Gọi $ S $ là điểm thỏa mãn $ \overrightarrow{A''S}  =\overrightarrow{B'C'} \Rightarrow A''B'C'S $ là hình bình hành
		$\Rightarrow A''B' = SC'.$\\
		Vì $ BC \parallel IJ \Rightarrow A''S \parallel IJ \Rightarrow S(21;3;5).$\\		
		Do đó, $AB + CO + OD$ nhỏ nhất bằng $SD$ khi và chỉ khi $S, O, C', D$ thẳng hàng, $SD = \sqrt{470}$.\\				
		Gọi $K$ là điểm thỏa mãn $\overrightarrow{HK} = \overrightarrow{FE} = (0; 0; -2) \Rightarrow K\left( 0; \dfrac{1}{2}; 8 \right)$ và $KHFE$ là hình bình hành.\\				
		Gọi $D'$ là điểm đối xứng với $D$ qua trục $Mz \Rightarrow D'(0; -8; 3)$. \\
		Lại có $ DE + FH = DE + EK = D'E + EK. $
		Để $DE + FH$ nhỏ nhất thì $D'$, $E$, $K$ thẳng hàng và $D'K = \dfrac{\sqrt{389}}{2}$.\\			
		Vậy $\min T = \sqrt{470} + \dfrac{\sqrt{389}}{2} + 5 \approx 36{,}5$.		
	}
	
\end{ex}

\begin{ex}%[Đề thi thử môn Toán TRƯỜNG THPT Cổ Loa - Hà Nội năm học 2025-2026]%[Phạm Quốc Toàn - Phạm Hoàng Việt]%[2H5H1-7]
	Tại một khu quảng trường, người ta muốn trang trí dải đèn led để chào mừng một sự kiện. Từ đỉnh $M$ của một cái cây cao (minh họa bằng đoạn $MK$), người ta kéo dây đèn led thẳng xuống một vị trí $H$ nào đó trên mặt đất, rồi tiếp tục kéo dây đèn led thẳng từ $H$ đến một vị trí $A$ cố định trên mặt đất sao cho $MH$ luôn vuông góc với $HA$. Sau đó, từ vị trí $H$, người ta sẽ đi dây điện thẳng đến công tắc nằm trên một bức tường vuông góc với mặt đất. Chọn hệ trục toạ độ $Oxyz$ với mặt phẳng $(Oxy)$ trùng với mặt đất (đơn vị trên mỗi trục tính theo mét), người ta xác định được $M(4;-2;3)$, $A(-2;6;0)$ và bức tường chứa công tắc nằm trên mặt phẳng $(\alpha) \colon 3x+4y+89=0$. Để độ dài đường dây điện phải đi là ngắn nhất, người ta đã tính toán khoảng cách nhỏ nhất từ vị trí $H$ đến bức tường bằng $a$ (m). Giá trị của $a$ bằng bao nhiêu?
	\begin{center}
		\begin{tikzpicture}[font=\footnotesize, line join=round, line cap=round, >=stealth,scale=.7]
			\coordinate (A) at (3,2);
			\coordinate (H) at (5,1);
			\coordinate (M) at (8,7);
			\coordinate (K) at (8,2);
			%%%%%%%%%%%%%%%%%%%%%%%%%%%%%%%%
			\draw (0,0)--(0,5) -- (2,8) -- (2, 3) -- (0,0) (0,0) -- (8,0) -- (10,3) -- (2, 3);
			\draw (A) -- (H) -- (M) -- (K);
			\foreach \i/\g in {A/90,H/-90,M/90,K/-90}{\draw(\i) circle (1pt) ($(\i)+(\g:3mm)$) node[scale=1]{$\i$};}
			\pic[draw,thin,angle radius=2mm] {right angle=A--H--M};
			%%%%%%%%%%%%%%%%%%%%%%%%%%%%%%%%
			\fill (0.2,0) node[above right]{$(Oxy)$}; 
			\fill (2,3) node[left]{$(\alpha)$}; 
		\end{tikzpicture}
	\end{center}
	\par\shortans[0]{15}
	\loigiai{
		Ta có $\widehat{AHM}=90^\circ$.\\
		Suy ra $H$ thuộc mặt cầu $(S)$ có đường kính $AM$.\\
		Mặt cầu $(S)$ có tâm $I\left(1;2;\dfrac{3}{2}\right)$, bán kính $R=\dfrac{\sqrt{109}}{2}$.\\
		Ta có $H=(S)\cap (Oxy)$.\\
		Suy ra $H$ thuộc đường tròn $(C)$ có tâm $J(1;2;0)$ là hình chiếu của $I$ lên mặt phẳng $(Oxy)$. \\
		Bán kính của đường tròn là $r=\sqrt{R^2-IJ^2}=5$.\\
		Ta có $\mathrm{d} \big(J, (\alpha)\big)=20$ và $(\alpha) \perp (Oxy)$ nên $\mathrm{d} \big(H, (\alpha)\big)_{\min}=\mathrm{d} \big(J, (\alpha)\big)-r=20-5=15$.
	}
\end{ex}

\begin{ex}%[2H5V1-5]%[TT-LienTruong-HaNoi-N25-26]%[Phạm Quốc Toàn - Lê Phúc]
	Cho hình chóp $S. ABCD$ có đáy là hình thoi cạnh $3$, góc $\widehat{ABC}=60^\circ$ và cạnh bên $SA$ vuông góc với mặt phẳng đáy. Gọi $M$, $N$ lần lượt là trung điểm của $SB$, $SD$. Biết góc nhị diện $[M, AC, N]$ bằng $120^\circ$. Tính khoảng cách từ điểm $C$ đến mặt phẳng $(SBD)$. (\textit{Không làm tròn kết quả các phép tính trung gian, chỉ làm tròn kết quả cuối cùng đến hàng phần mười})
		\par\shortans[0]{1{,}1}
	\loigiai{
		\immini[thm]{Đặt hệ trục tọa độ $Oxyz$ với gốc $O$ là giao điểm của hai đường chéo $AC$ và $BD$.\\
		Vì $ABCD$ là hình thoi cạnh $3$ và $\widehat{ABC}=60^\circ$, ta có tam giác $ABC$ đều, suy ra $AC=3$ và $OB=OD=\dfrac{3\sqrt{3}}{2}$.\\
		Tọa độ các đỉnh là $O(0;0;0)$, $A\left(-\dfrac{3}{2};0;0\right)$, $C\left(\dfrac{3}{2};0;0\right)$, $B\left(0;\dfrac{3\sqrt{3}}{2};0\right)$, $D\left(0;-\dfrac{3\sqrt{3}}{2};0\right)$. \\ 
		Vì $SA\perp (ABCD)$, gọi $SA=h>0$ là chiều cao của hình chóp $S.ABCD$ nên $S\left(-\dfrac{3}{2};0;h\right)$.\\ 
		Tọa độ trung điểm của $SB$ là $M\left(-\dfrac{3}{4};\dfrac{3\sqrt{3}}{4};\dfrac{h}{2}\right)$. \\ 
		Tọa độ trung điểm của $SD$ là $N\left(-\dfrac{3}{4};-\dfrac{3\sqrt{3}}{4};\dfrac{h}{2}\right)$. }
		{	\begin{tikzpicture}[font=\footnotesize, line join=round, line cap=round, >=stealth,scale=0.6]
				\coordinate (S) at (3,7);
				\coordinate (A) at (3,2);
				\coordinate (B) at (0,0);
				\coordinate (C) at (5,0);
				\coordinate (D) at (8,2);
				\coordinate (O) at (4,1);
				\coordinate (G) at (4,4.5);
				\coordinate (M) at (1.5,3.5);
				\coordinate (N) at (5.5,4.5);
				%%%%%%%%%%%%%%%%%%%%%%%%%%%
				\draw[->] (5,0)--(6,-1) node[below]{$x$};
				\draw[->] (B)--($(O)!1.3!(B)$) node[below]{$y$};
				\draw[->] (4,4.5)--(4,7.6) node[right]{$z$};
				%%%%%%%%%%%%%%%%%%%%%%%%%%%
				\draw (S) -- (B)  -- (C) -- (D) -- (S) -- (C);
				\draw[dashed] (S)--(A) -- (B)  (A) -- (D) -- (B) (A) -- (C) (O) -- (G);
				%%%%%%%%%%%%%%%%%%%%%%%%%%%
				\foreach \i/\g in {S/90,A/180,B/-90,C/-90,D/0,O/-90,M/180,N/90}{\fill(\i) circle (1.5pt) ($(\i)+(\g:3mm)$) node[scale=1]{$\i$};}
				\pic[draw,thin,angle radius=1.5mm] {right angle=G--O--B} 
				pic[draw,thin,angle radius=1.5mm] {right angle=G--O--C};
		\end{tikzpicture}}
	\noindent
		Góc nhị diện $\left[M, AC, N\right]$ bằng $120^\circ$ nên góc giữa hai mặt phẳng $(MAC)$ và $(NAC)$ bằng $60^\circ$. \\ 
		Vectơ pháp tuyến của mặt phẳng $(MAC)$ là $\overrightarrow{n}_1=\left[\overrightarrow{AC}, \overrightarrow{AM}\right]=\left(0;-h;\dfrac{3\sqrt{3}}{2}\right)$.\\
		Vectơ pháp tuyến của mặt phẳng $(NAC)$ là $\overrightarrow{n}_2=\left[\overrightarrow{AC}, \overrightarrow{AN}\right]=\left(0;h;\dfrac{3\sqrt{3}}{2}\right)$.\\
		Vậy 
		\begin{align*}
			& \cos 60^\circ =\dfrac{\left|\overrightarrow{n}_1\cdot{\overrightarrow{n}_2}\right|}{\left|\overrightarrow{n}_1\right|\cdot \left|\overrightarrow{n}_2\right|}=\dfrac{\left| \dfrac{27}{4}-h^2\right|}{\sqrt{h^2+\dfrac{27}{4}}\cdot \sqrt{h^2+\dfrac{27}{4}}}  \\
			& \Leftrightarrow \dfrac{1}{2}=\dfrac{\left| \dfrac{27}{4}-h^2\right|}{\sqrt{h^2+\dfrac{27}{4}}\cdot \sqrt{h^2+\dfrac{27}{4}}} \\ 
			& \Leftrightarrow \sqrt{h^2+\dfrac{27}{4}}\cdot \sqrt{h^2+\dfrac{27}{4}}=2\left| \dfrac{27}{4}-h^2\right| \\ 
			& \Leftrightarrow h^2+\dfrac{27}{4}=2\left| \dfrac{27}{4}-h^2\right| \\ 
			& \Leftrightarrow h=\dfrac{3}{2}. 
		\end{align*}
		Phương trình mặt phẳng  $(SBD)$ đi qua ba điểm $S\left(-\dfrac{3}{2};0;\dfrac{3}{2}\right)$, $B\left(0;\dfrac{3\sqrt{3}}{2};0\right)$, $D\left(0;-\dfrac{3\sqrt{3}}{2};0\right)$ là $x+z=0$.\\
		Khoảng cách từ điểm $C\left(\dfrac{3}{2};0;0\right)$ đến mặt phẳng $(SBD)$ là
		\[
		\mathrm{d}\big(C;(SBD)\big)=\dfrac{3}{2\sqrt{2}}=\dfrac{3\sqrt{2}}{4}\approx 1{,}1. 
		\]
	}
\end{ex}

\begin{ex}%[2H5V1-5]%[Dự án đề thi thử đợt 3-Nguyễn Đức Lợi]%[SGD tỉnh Lạng Sơn]%[GVPB - ThanhTa]
	Trong không gian $Oxyz$, cho điểm $M(1;-2;-5)$. Gọi $(P)$ là mặt phẳng đi qua $M$ và cách gốc tọa độ một khoảng lớn nhất. Biết mặt phẳng $(P)$ cắt $Ox$, $Oy$, $Oz$ lần lượt tại ba điểm $A$, $B$, $C$. Thể tích tứ diện $OABC$ bằng bao nhiêu?
	\par
	\shortans{450}
	\loigiai{
		\begin{center}
			\begin{tikzpicture}[scale=0.8, font=\footnotesize, line join=round, line cap=round, >=stealth]
				% Trục tọa độ
				\draw[->] (3,0,0) -- (4,0,0) node[below] {$x$};
				\draw[->] (0,4,0) -- (0,5,0) node[left] {$z$};
				\draw[->] (0,0,5) -- (0,0,6) node[right] {$y$};
				
				% Các điểm giao tuyến (ước lệ minh họa)
				\coordinate (A) at (3,0,0);
				\coordinate (B) at (0,0,5);
				\coordinate (C) at (0,4,0);
				\coordinate (O) at (0,0,0);
				\coordinate (M) at (1,1,1.5); % Minh họa điểm M
				
				% Vẽ mặt phẳng (P)
				\draw[thick, fill=blue!10, opacity=0.6] (A) -- (B) -- (C) -- cycle;
				
				% Vẽ đoạn OM
				\draw[dashed, thick, red] (O) -- (M) (O)--(A) (O)--(B) (O)--(C);
				\fill[red] (M) circle (1.5pt) node[right] {$M$};
				
				% Tên các điểm giao
				\fill (A) circle (1pt) node[below] {$A$};
				\fill (O) circle (1pt) node[below] {$O$};
				\fill (B) circle (1pt) node[left] {$B$};
				\fill (C) circle (1pt) node[left] {$C$};
			\end{tikzpicture}
		\end{center}
		Gọi $h$ là khoảng cách từ $O$ đến mặt phẳng $(P)$.\\\
		Ta luôn có $h = \mathrm{d}(O, (P)) \le OM$.\\
		Do đó, khoảng cách $\mathrm{d}(O, (P))$ lớn nhất khi và chỉ khi $(P)$ vuông góc với $OM$ tại $M$.\\
		Khi đó, phương trình mặt phẳng $(P)$ đi qua $M(1; -2; -5)$ và có vectơ pháp tuyến $\overrightarrow{OM} = (1; -2; -5)$ là
		\[1(x - 1) - 2(y + 2) - 5(z + 5) = 0 \Leftrightarrow x - 2y - 5z - 30 = 0.\]
		Tọa độ giao điểm của $(P)$ với các trục tọa độ
		\begin{itemize}
			\item Cho $y=0$, $z=0$ suy ra $x = 30 \Rightarrow A(30; 0; 0) \Rightarrow OA = 30$.
			\item Cho $x=0$, $z=0$ suy ra $-2y = 30 \Rightarrow y = -15 \Rightarrow B(0; -15; 0) \Rightarrow OB = 15$.
			\item Cho $x=0$, $y=0$ suy ra $-5z = 30 \Rightarrow z = -6 \Rightarrow C(0; 0; -6) \Rightarrow OC = 6$.
		\end{itemize}
		Tứ diện $OABC$ có $OA$, $OB$, $OC$ đôi một vuông góc nên thể tích là
		\[V = \dfrac{1}{6} \cdot OA \cdot OB \cdot OC = \dfrac{1}{6} \cdot 30 \cdot 15 \cdot 6 = 450.\]
	}
\end{ex}

\begin{ex}%[2H5V1-7]%[TT Cụm 5 - Ninh Bình-NH25-26]%[GVSB:Dương Công Tạo]%[GVPB1: Nguyễn Tấn Tài]
	\immini[thm]{Một hộ gia đình muốn xây dựng một kho chứa mini dạng khối tứ diện $OABC$ tận dụng góc tường nhà (với $O$ là góc tường, ba cạnh $OA$, $OB$, $OC$ nằm trên ba mép tường vuông góc). Để lắp đặt hệ thống thông gió, mặt phía trước của kho (mặt phẳng $ABC$) bắt buộc phải đi qua một van điều tiết đặt tại vị trí $M$ có tọa độ $M(1;2;1)$ (đơn vị tính là mét). Gia đình muốn thiết kế kho sao cho thể tích của kho là nhỏ nhất để không làm ảnh hưởng quá nhiều đến diện tích sinh hoạt của sân chung. Khi kho được thiết kế với thể tích nhỏ nhất, hãy tính khoảng cách từ góc tường $O$ đến mặt phẳng $(ABC)$.}
	{\begin{tikzpicture}[scale=.75,
			x={(220:1.3cm)}, 
			y={(10:1.1cm)}, 
			z={(90:1cm)},  
			>=stealth    
			]
			\coordinate (O) at (0,0,0);
			\coordinate (A) at (3.5,0,0); 
			\coordinate (B) at (0,4.0,0); 
			\coordinate (C) at (0,0,3.5);
			\fill[lightgray!40] (0,0,3.7) -- (0,0,0) -- (0,4.2,0) -- (0,4.2,3.7) -- cycle;
			
			\fill[lightgray!40] (3.7,0,0) -- (0,0,0) -- (0,0,3.7) -- (3.7,0,3.7) -- cycle;
			
			\draw[->] (O) -- (4,0,0) node[anchor=north east, pos=1] {$x$};
			\draw[->] (O) -- (0,4.5,0) node[anchor=west, pos=1] {$y$};
			\draw[->] (O) -- (0,0,4) node[anchor=south, pos=1] {$z$};
			
			\draw[thick] (A) -- (B) -- (C) -- cycle;
			
			\node[above left] at (A) {$A$};
			\node[below right] at (B) {$B$};
			\node[above left] at (C) {$C$};
			\node[below] at (O) {$O$};
			
			\coordinate (M) at (1.5, 2.5,1.2); 
			
			\fill[black] (M) circle (1pt); 
			\node[above=3pt] at (M) {$M(1; 2;1)$}; 
			
			\draw[dashed, black!60] (M) -- (0, 1, -0.3); 
			
	\end{tikzpicture}}
	\par  \shortans{2}
	\loigiai{
		Đặt hệ trục tọa độ $Oxyz$ có gốc $O(0;0;0)$, các tia $Ox$, $Oy$, $Oz$ lần lượt chứa các cạnh $OA$, $OB$, $OC$ của tứ diện. \\
		Gọi tọa độ các điểm là $A(a;0;0)$, $B(0;b;0)$, $C(0;0;c)$ với $a$, $b$, $c > 0$. \\
		Phương trình mặt phẳng $(ABC)$ có dạng mặt phẳng theo đoạn chắn $\dfrac{x}{a} + \dfrac{y}{b} + \dfrac{z}{c} = 1$. \\
		Vì mặt phẳng $(ABC)$ đi qua điểm $M(1;2;1)$ nên ta có hệ thức $\dfrac{1}{a} + \dfrac{2}{b} + \dfrac{1}{c} = 1$. \\
		Thể tích khối tứ diện $OABC$ là $V = \dfrac{1}{6}abc$. \\
		Áp dụng bất đẳng thức Cauchy (AM-GM) cho ba số dương $\dfrac{1}{a}$, $\dfrac{2}{b}$, $\dfrac{1}{c}$, ta có \[1 = \dfrac{1}{a} + \dfrac{2}{b} + \dfrac{1}{c} \ge 3\sqrt[3]{\dfrac{1}{a} \cdot \dfrac{2}{b} \cdot \dfrac{1}{c}} = 3\sqrt[3]{\dfrac{2}{abc}} \Rightarrow 1 \ge 27 \cdot \dfrac{2}{abc} \Rightarrow abc \ge 54.\]
		Thể tích nhỏ nhất $V_{\min} = \dfrac{1}{6} \cdot 54 = 9$ m$^3$ khi và chỉ khi $\dfrac{1}{a} = \dfrac{2}{b} = \dfrac{1}{c} = \dfrac{1}{3} \Rightarrow a = 3$, $b = 6$, $c = 3$. \\
		Khi đó, phương trình mặt phẳng $(ABC)$ là \[\dfrac{x}{3} + \dfrac{y}{6} + \dfrac{z}{3} = 1 \Leftrightarrow 2x+y+2z-6=0.\]
		Khoảng cách từ gốc tọa độ $O$ đến mặt phẳng $(ABC)$ là \[\mathrm{d}\left(O,(ABC)\right)= \dfrac{\left|2 \cdot 0+0+2 \cdot 0-6\right|}{\sqrt{2^2 + 1^2 + 2^2}} = \dfrac{|-6|}{\sqrt{9}} = \dfrac{6}{3} = 2\text{ (m)}.\]
	}
\end{ex}

\begin{ex}%[2H5V1-7]%[Dự án C đề thi thử TN Môn Toán NH25-26 - Dot 4 - Hector Tran]%[Liên trường THPT cụm 09 - Đợt 2 - Hà Nội]%[GVPB - Thanh Phát]
	\immini[thm]{
		Trong giờ thể dục học về kỹ thuật chuyền bóng hơi, Bình và An tập chuyền bóng cho nhau. Ở một động tác Bình chuyền bóng cho An, quả bóng bay lên cao nhưng lại lệch sang bên trái của An và rơi xuống vị trí cách chỗ An đứng $0{,}5$ m và cách chỗ Bình $4{,}5$ m. Chọn hệ trục tọa độ $Oxyz$ sao cho gốc tọa độ $O$ tại vị trí của Bình, vị trí của An nằm trên tia $Oy$ và mặt phẳng $(Oxy)$ là mặt đất. Biết rằng quỹ đạo của quả bóng nằm trong mặt phẳng $(\alpha) \colon x + by + cz + d = 0$ và $(\alpha)$ vuông góc với mặt đất. Khi đó, giá trị của $200b^2 - 100c^2 + 200d^2$ bằng bao nhiêu?
	}{
		\begin{tikzpicture}[line join = round, line cap=round,>=stealth,font=\footnotesize,scale=1]
			\draw[->] (2,2)coordinate (O)--(7.5,2)coordinate (y);
			\draw[->] (2,2)--(0.6,0.8)coordinate (x);
			\draw[->] (2,2)--(2,5.2)coordinate (z);
			\fill (2.3,3.1) circle(2pt) (6.55,1.35) circle(2pt);
			\draw[dashed] (2.3,3.1)..controls +(28:3.0) and +(90:1) .. (6.55,1.35)coordinate (H)node[below]{$A$};
			\fill (1.9,2.85) circle(3pt);
			\draw[line width=3pt] (1.9, 2.85) -- (2.0, 2.6) -- (2.0, 2.3) -- (1.8, 2.1)
			(2.0, 2.3) -- (2.05, 2.2) -- (2.0, 2.0)
			(2.0, 2.75) -- (2.2, 2.95)
			;
			\fill (6.65,2.7) circle(3pt);
			\draw[line width=3pt]
			(6.65, 2.7) -- (6.85, 2.4) -- (6.75, 2.2) -- (6.95, 2.0)
			(6.85, 2.4) -- (6.65, 2.3) -- (6.65, 2.0)
			(6.7, 2.6) -- (6.5, 2.5)
			(6.7, 2.6) -- (6.5, 2.4)
			;
			\draw (O)--(H)--($(O)!(H)!(y)+(0.13,0)$)coordinate (t);
			\draw pic[draw,angle radius=2mm] {right angle = H--t--y};
			\path
			($(t)+(0,-0.4)$) node[right]{$0{,}5$\,m}
			($(O)+(2,-0.75)$) node[]{$4{,}5$\,m}
			;
			\foreach \p/\r in {O/-90,y/90,x/180,z/180}
			\fill (\p) node[shift={(\r:3mm)}]{$\p$};
		\end{tikzpicture}
	}
	\par\shortans{2,5}
	\loigiai{
		Bình ở $O(0; 0; 0)$. Quả bóng rơi xuống tại điểm $A(x; y; 0)$ thỏa mãn $OA = 4{,}5$ và cách An một đoạn $0{,}5$.\\
		Suy ra $A\left( 0{,}5; \sqrt{20};0\right) $.\\
		Mặt phẳng $(\alpha) \colon x + by + cz + d = 0$ đi qua $O$ nên $d = 0$.\\
		Điểm $A\left( 0{,}5; \sqrt{20};0\right) $ thuộc $(\alpha)$ nên $0{,}5+ \sqrt{20} b = 0 \Leftrightarrow b = -\dfrac{\sqrt{5}}{20}$.\\
		Mặt khác $(\alpha)$ vuông góc với mặt đất $(Oxy)$ nên $\overrightarrow{n}_{\alpha} \perp \overrightarrow{k} \Rightarrow c = 0$.\\
		Vậy giá trị $200b^2 - 100c^2 + 200d^2 =200\left( -\dfrac{\sqrt{5}}{20}\right) ^2 =2{,}5$.
	}
\end{ex}

\begin{ex}%[Đề thi thử môn Toán Sở giáo dục tỉnh Tuyên Quang năm học 2025-2026]%[Đoàn Minh Tâm]%[2H5V1-7]%[GVPB: Thanh Phát]
	Ông An đang xây một ngôi nhà, trong quá trình xây phải đổ bê-tông cho một mái vát để lợp ngói. Ông tính toán việc ghép cốt pha đi qua điểm $B$ trên một chân tường và điểm $C$ trên cột góc nhà, đồng thời mặt ghép cốt pha phải đi qua điểm $A$ trên chân tường còn lại cách điểm $O$ ở góc giao hai chân tường một khoảng $5$ m, ông cũng tận dụng một chiếc cột có sẵn để chống mặt ghép (xem hình dưới). Biết rằng hai bức tường được xây vuông góc với nhau, mỗi bức tường đều vuông góc với sàn mái nhà, cột có chiều cao $1$ m và cách hai bức tường với cùng khoảng cách $1$ m (đỉnh cột là điểm $M$). Diện tích nhỏ nhất của khung ghép cốt pha $ABC$ là bao nhiêu mét vuông (\textit{làm tròn kết quả đến hàng phần trăm})?
	\begin{center}
		\begin{tikzpicture}[
			% Tùy chỉnh góc nhìn 3D
			x={({0.866cm,-0.15cm})},
			y={({-0.6cm,-0.5cm})},
			z={({0cm,1cm})},
			scale=1.2,
			font=\footnotesize,
			line join=round,line cap=round,line width=.6pt,>=stealth
			]
			\coordinate (O) at (0,0,0);
			\coordinate (A_ext) at (7.5,0,0);
			\coordinate (A) at (7,0,0);
			\coordinate (B_ext) at (0,4.5,0);
			\coordinate (B) at (0,3.5,0);
			\coordinate (C_ext) at (0,0,4.5);
			\coordinate (C) at (0,0,3.8);
			\coordinate (Mx) at (2, 0, 0);
			\coordinate (My) at (0, 2, 0);
			\coordinate (Mz) at (0, 0, 2);
			\coordinate (Mxy) at (2, 2, 0);
			\coordinate (M) at (2, 2, 2);
			\begin{scope}
				\clip (A) -- (B) -- (C) -- cycle;
				\fill[gray!10] (A) -- (B) -- (C) -- cycle;
				\foreach \x in {0.2, 0.4, ..., 7.4} {
					\pgfmathsetmacro{\factor}{1-\x/7.5}
					\draw[very thin, gray!60] (\x, {\factor*3.5}, 0) -- (\x, 0, {\factor*3.8});
				}
				\foreach \y in {0.2, 0.4, ..., 3.4} {
					\pgfmathsetmacro{\factor}{1-\y/3.5}
					\draw[very thin, gray!60] ({\factor*7.5}, \y, 0) -- (0, \y, {\factor*3.8});
				}
			\end{scope}
			
			\draw[dashed, thick] (O) -- (A);
			\draw[dashed, thick] (O) -- (B);
			\draw[dashed, thick] (O) -- (C);
			\draw[thick] (A) -- (A_ext);
			\draw[thick] (B) -- (B_ext);
			\draw[thick] (C) -- (C_ext);
			\draw[thick] (A) -- (B) -- (C) -- cycle;
			\draw[dashed, thick] (Mx) -- (Mxy) -- (My);
			\draw[dashed, thick] (O) -- (Mxy);
			\draw[dashed, thick] (Mxy) -- (M)--(Mz);
			\fill (M) circle (1pt);
			\fill (O) circle (1pt);
			\fill (A) circle (1pt);
			\fill (B) circle (1pt);
			\fill (C) circle (1pt);
			\node[above right ] at (A) {$ A $};
			\node[left ] at (B) {$ B $};
			\node[above right ] at (C) {$ C $};
			\node[above right ] at (O) {$ O $};
			\node[above right=-1pt ] at (M) {$ M $};
			\tikzset{lbl/.style={fill=white, inner sep=1.5pt}}
			\path (O) -- (A) node[midway, sloped, above= 1pt, lbl] {Chân tường $1$};
			\path (O) -- (B) node[pos=0.38, sloped, above= 1pt, lbl] {Chân tường $2$};
			\path (O) -- (C) node[pos=0.45, sloped, above= 1pt, lbl] {Cột góc nhà};
			\path (Mxy) -- (Mx) node[midway, sloped, below=1pt, lbl] {$ 1 $ m};
			\path (Mxy) -- (My) node[midway, sloped, below=1pt, lbl] {$ 1 $ m};
			\path (Mxy) -- (M) node[pos=0.7, right=1pt, lbl] {$ 1 $ m};
		\end{tikzpicture}
	\end{center}
	\par\shortans[oly]{9,38}
	\loigiai{
		\immini[thm]{
			Chọn hệ trục tọa độ $Oxyz$ với $O(0;0;0)$ là góc giao hai chân tường, tia $Oy$ nằm trên chân tường chứa điểm $A$, tia $Ox$ nằm trên chân tường chứa điểm $B$, tia $Oz$ trùng với cột góc nhà chứa điểm $C$.\\
			Theo giả thiết bài toán, ta có tọa độ các điểm $A(0;5;0)$.\\
			Giả sử $B(b;0;0)$ và $C(0;0;c)$ với $b>0$, $c>0$.\\
			Mặt phẳng $(ABC)$ chứa mặt ghép cốt pha có phương trình theo đoạn chắn là
			\[\dfrac{x}{b} + \dfrac{y}{5} + \dfrac{z}{c} = 1.\]
		}{
			\begin{tikzpicture}[
				x={({0.866cm,-0.15cm})},
				y={({-0.6cm,-0.5cm})},
				z={({0cm,1cm})},
				scale=.7,
				font=\footnotesize,
				line join=round,line cap=round,line width=.6pt,>=stealth
				]
				\coordinate (O) at (0,0,0);
				\coordinate (A_ext) at (7.5,0,0);
				\coordinate (A) at (7,0,0);
				\coordinate (B_ext) at (0,4.5,0);
				\coordinate (B) at (0,3.5,0);
				\coordinate (C_ext) at (0,0,4.5);
				\coordinate (C) at (0,0,3.8);
				\coordinate (Mx) at (2, 0, 0);
				\coordinate (My) at (0, 2, 0);
				\coordinate (Mz) at (0, 0, 2);
				\coordinate (Mxy) at (2, 2, 0);
				\coordinate (M) at (2, 2, 2);
				\draw[dashed, thick] (O) -- (A);
				\draw[dashed, thick] (O) -- (B);
				\draw[dashed, thick] (O) -- (C);
				\draw[thick,->] (A) -- (A_ext) node[right]{$y$};
				\draw[thick,->] (B) -- (B_ext) node[below]{$x$};
				\draw[thick,->] (C) -- (C_ext) node[above]{$z$};
				\draw[thick] (A) -- (B) -- (C) -- cycle;
				\draw[dashed, thick] (Mx) -- (Mxy) -- (My);
				\draw[dashed, thick] (O) -- (Mxy);
				\draw[dashed, thick] (Mxy) -- (M)--(Mz);
				\fill (M) circle (1.5pt);
				\fill (O) circle (1.5pt);
				\fill (A) circle (1.5pt);
				\fill (B) circle (1.5pt);
				\fill (C) circle (1.5pt);
				\node[above right ] at (A) {$A$};
				\node[left ] at (B) {$B$};
				\node[above right ] at (C) {$C$};
				\node[above left ] at (O) {$O$};
				\node[above right=-1pt ] at (M) {$M$};
			\end{tikzpicture}
		}
		\noindent Cột có chiều cao $1$ m, cách hai bức tường $1$ m nên đỉnh cột là $M(1;1;1)$.\\
		Vì mặt ghép cốt pha đi qua $M$ nên điểm $M \in (ABC)$.\\
		Thay tọa độ $M$ vào phương trình mặt phẳng ta được
		\[\dfrac{1}{b} + \dfrac{1}{5} + \dfrac{1}{c} = 1 \Leftrightarrow \dfrac{1}{b} + \dfrac{1}{c} = \dfrac{4}{5} \Leftrightarrow \dfrac{b+c}{bc} = \dfrac{4}{5} \Leftrightarrow b+c = \dfrac{4}{5}bc.\]
		Áp dụng bất đẳng thức AM-GM, ta có
		\[b+c \ge 2\sqrt{bc} \Leftrightarrow \dfrac{4}{5}bc \ge 2\sqrt{bc} \Leftrightarrow \sqrt{bc} \ge \dfrac{5}{2} \Leftrightarrow bc \ge \dfrac{25}{4}.\]
		Diện tích khung ghép cốt pha là diện tích tam giác $ABC$ là
		\[S = \dfrac{1}{2} \left| \left[ \overrightarrow{AB}, \overrightarrow{AC} \right] \right| = \dfrac{1}{2} \sqrt{b^2c^2 + 25b^2 + 25c^2}.\]
		Ta cần tìm giá trị nhỏ nhất của $S$, ta xét $S^2$ như sau
		\[S^2 = \dfrac{1}{4} \left[ b^2c^2 + 25(b^2+c^2) \right] = \dfrac{1}{4} \left[ b^2c^2 + 25 \left( (b+c)^2 - 2bc \right) \right].\]
		Thay $b+c = \dfrac{4}{5}bc$ vào biểu thức trên, ta có
		\[S^2 = \dfrac{1}{4} \left[ b^2c^2 + 25 \left( \dfrac{16}{25}b^2c^2 - 2bc \right) \right] = \dfrac{1}{4} \left( 17b^2c^2 - 50bc \right).\]
		Đặt $t = bc$ với $t \ge \dfrac{25}{4} = 6{,}25$.\\
		Xét hàm số $f(t) = 17t^2 - 50t$ trên $[6{,}25; +\infty)$.\\
		Ta có $f'(t) = 34t - 50$.\\
		Với $t \ge 6{,}25$ thì $f'(t) \ge 34\cdot 6{,}25 - 50 = 162{,}5 > 0$.\\
		Do đó hàm số $f(t)$ đồng biến trên $[6{,}25; +\infty)$.\\
		Suy ra $\min f(t) = f(6{,}25) = 17\cdot (6{,}25)^2 - 50\cdot 6{,}25 = 351{,}5625$.\\
		Khi đó $S^2 \ge \dfrac{1}{4} \cdot 351{,}5625 = 87{,}890625 \Rightarrow S \ge \sqrt{87{,}890625} = 9{,}375 \approx 9{,}38 \text{ m}^2$.\\
		Dấu \lq\lq $=$\rq\rq\ xảy ra khi $b=c=2{,}5$.
	}
\end{ex}

\begin{ex}%[2H5V2-8]%[Dự án C đề thi thử TN Môn Toán NH25-26 - Dot 4 - Vũ Hồng Toàn]%[THPT Mai Anh Tuấn - Thanh Hóa]%[GVPB - Lê Hữu Kiệt]
	Trong không gian một bức tường hình chữ nhật $ABCD$ có chiều cao $3$\,m được dựng vuông góc với mặt phẳng $(Oxy)$ với hai điểm $A(3;1;0)$ và $B(1;4;0)$. Giả sử rằng mặt phẳng $(Oxy)$ là mặt đất, trục $Oz$ hướng lên trên và mỗi đơn vị trên hệ trục tương ứng là $1$\,m. Một trạm phát sóng được đặt tại điểm $M(5;6;5)$. Trạm phát sóng phát ra các sóng thẳng và truyền về mọi hướng nhưng khi gặp bức tường thì bị chắn lại vì vậy có một vùng sau bức tường không có sóng. Một chất điểm bắt đầu chuyển động thẳng đều trong mặt phẳng mặt phẳng $(Oxy)$ từ điểm $N(-6;-4;0)$ qua gốc tọa độ $O$ đến chân tường với vận tốc $0{,}8$\,m/s. Hỏi sau bao nhiêu giây từ lúc chất điểm bắt đầu đi vào vùng mất sóng đến lúc chất điểm chạm vào chân tường và quay về vị trí ban đầu, giả thiết độ dày của bức tường là không đáng kể {\it(kết quả làm tròn đến hàng phần chục)}?
	\begin{center}
		\begin{tikzpicture}[scale=1, font=\footnotesize, line join=round, line cap=round, >=stealth]
			\tikzset{declare function={a=6;b=0.4*a;goc=-130;h=0.65*a;}}
			\path (0:0) coordinate (O)
			(goc+20:b-1) coordinate (A)++(0,b-0.5)coordinate (D)
			(goc+125:b) coordinate (B)++(0,b-0.5)coordinate (C)
			(a-1,-b+1)coordinate (Hx)++(0,h)coordinate (M)
			(-goc+40:h)coordinate (N)
			(intersection of M--C and Hx--B) coordinate (E)
			(intersection of M--D and Hx--A) coordinate (F)
			($(N)+(Hx)-(F)$)coordinate (Ix)
			(intersection of N--O and A--D) coordinate (No)
			(intersection of N--O and A--B) coordinate (H)
			(intersection of N--O and E--F) coordinate (I)
			(intersection of {O--$(goc:1)$} and D--A) coordinate (Ox)
			(intersection of {O--$(1,0)$} and B--C) coordinate (Oy)
			(intersection of {O--$(0,1)$} and D--C) coordinate (Oz) ;
			\fill[gray!10] ($(F) + (-1,-0.5)$) -- ($(N) + (-1,0.5)$) -- ($(Ix) + (0,0.7)$) -- ($(Hx) + (1,-0.5)$) -- cycle;
			\draw[gray, line width=2pt] (M)--(Hx);
			\draw[->] (Oy)--(a+0.5,0) node[shift={(-100:7pt)}]{$y$};
			\draw[->] (Ox)--(goc:b)node[shift={(0:7pt)}]{$x$};
			\draw[->] (Oz)--(0,h-1)node[shift={(200:7pt)}]{$z$};
			\fill[gray!60, opacity=0.6, pattern=north west lines, pattern color=red!30] (A) -- (B) -- (E) -- (F) -- cycle;
			\node at (-2,-0.7,0) [teal, font=\footnotesize, align=center] {Vùng mất sóng};
			\shade[shading=ball](M)circle(4pt);
			\draw[dashed] (Oy)--(O)--(Ox) (Oz)--(O);
			\fill[draw, pattern=bricks, pattern color=red!50] (A)--(B)--(C)--(D)--cycle;
			\draw[red] (N)--(No);
			\draw[red, dashed, thin] (H)--(No);
			\foreach \x /\gN in {A/-90,B/-10,C/90,D/-40,O/-80,N/90,M/90}
			\fill (\x) circle (1pt)($(\x)+(\gN:2.5mm)$) node {$\x$};
		\end{tikzpicture}
	\end{center}
	\par\shortans{21,1}
	\loigiai{
		\begin{center}
			\begin{tikzpicture}[scale=1, font=\footnotesize, line join=round, line cap=round, >=stealth]
				\tikzset{declare function={a=6;b=0.4*a;goc=-130;h=0.65*a;}}
				\path (0:0) coordinate (O)
				(goc+20:b-1) coordinate (A)++(0,b-0.5)coordinate (D)
				(goc+125:b) coordinate (B)++(0,b-0.5)coordinate (C)
				(a-1,-b+1)coordinate (Hx)++(0,h)coordinate (M)
				(-goc+40:h)coordinate (N)
				(intersection of M--C and Hx--B) coordinate (E)
				(intersection of M--D and Hx--A) coordinate (F)
				($(N)+(Hx)-(F)$)coordinate (Ix)
				(intersection of N--O and A--D) coordinate (No)
				(intersection of N--O and A--B) coordinate (H)
				(intersection of N--O and E--F) coordinate (I)
				(intersection of {O--$(goc:1)$} and D--A) coordinate (Ox)
				(intersection of {O--$(1,0)$} and B--C) coordinate (Oy)
				(intersection of {O--$(0,1)$} and D--C) coordinate (Oz) ;
				\fill[gray!10] ($(F) + (-1,-0.5)$) -- ($(N) + (-1,0.5)$) -- ($(Ix) + (0,0.7)$) -- ($(Hx) + (1,-0.5)$) -- cycle;
				\draw[gray, line width=2pt] (M)--(Hx);
				\draw[->] (Oy)--(a+0.5,0) node[shift={(-100:7pt)}]{$y$};
				\draw[->] (Ox)--(goc:b)node[shift={(0:7pt)}]{$x$};
				\draw[->] (Oz)--(0,h-1)node[shift={(200:7pt)}]{$z$};
				\fill[gray!60, opacity=0.6, pattern=north west lines, pattern color=red!30] (A) -- (B) -- (E) -- (F) -- cycle;
				\node at (-2,-0.7,0) [teal, align=center] {Vùng mất sóng};
				\shade[shading=ball](M)circle(4pt);
				\draw[dashed] (Oy)--(O)--(Ox) (Oz)--(O);
				\fill[draw, pattern=bricks, pattern color=red!50] (A)--(B)--(C)--(D)--cycle;
				\draw[red] (N)--(No);
				\draw[red, dashed, thin] (H)--(No);
				\draw[thin, gray] (E)--(M) --(F)--cycle;
				\foreach \x /\gN in {A/-90,B/-10,C/90,D/-40,O/-80,H/-30,M/90,E/130,F/180,N/90,I/90}
				\fill (\x) circle (1pt)($(\x)+(\gN:2.5mm)$) node {$\x$};
			\end{tikzpicture}
		\end{center}
		Ta có $C(1;4;3)$, $D(3;1;3)$.\\
		Gọi $E(x;y;0)$, $F$ lần lượt là giao điểm của $MC$, $MD$ với mặt phẳng $(Oxy)$.\\
		Khi đó
		\allowdisplaybreaks
		\begin{eqnarray*}
			\overrightarrow{MC} = k\overrightarrow{ME} \Leftrightarrow \heva{&1-5 = k(x-5)\\& 4-6 = k(y-6)\\& 3-5 = k(0-5)} \Leftrightarrow \heva{&x = -5\\& y = 1\\&k = \dfrac{2}{5}} \Rightarrow E(-5;1;0).
		\end{eqnarray*}
		Tương tự, ta cũng có $F(0;-\dfrac{13}{2};0)$.\\
		Gọi $I$, $H$ lần lượt là giao điểm của đường thẳng $NO$ với $EF$ và $AB$.\\
		Vì $I$, $H$, $O$, $N \in (Oxy)$ nên chỉ xét trên mặt phẳng tọa độ $(Oxy)$.\\
		Ta có $N(-6;-4)$; $O(0;0)$ nên phương trình đường thẳng $NO$ là $2x - 3y = 0$.\\
		$E(-5;1)$; $F\left(0;-\dfrac{13}{2}\right)$ nên phương trình đường thẳng $EF$ là $3x + 2y + 13 = 0$.\\
		Với $A(3;1)$; $B(1;4)$ nên phương trình đường thẳng $AB$ là $3x + 2y - 11 = 0$.\\
		Suy ra $I(-3;-2)$, $H\left(\dfrac{33}{13};\dfrac{22}{13}\right) \Rightarrow IH = \dfrac{24\sqrt{13}}{13}$, $IN = \sqrt{13}$.\\
		Tổng quãng đường chất điểm di chuyển cần tính thời gian là
		$$IH + HN = 2IH + IN = \dfrac{61\sqrt{13}}{13}.$$
		Vậy thời gian cần tìm là $t = \dfrac{\dfrac{61\sqrt{13}}{13}}{0{,}8}  \approx 21{,}1$\,(giây).
	}
\end{ex}

\begin{ex}%[2H5V3-4]%[KNTT - Lớp 12 - Dự án C - Đợt 2]%[Loc do]%[PB: Phan Anh]
	Trong không gian $Oxyz$ (đơn vị km), Trái đất được mô phỏng là mặt cầu tâm $O$ bán kính $R=6\,400$ và đường xích đạo được xem là đường tròn giao tuyến giữa mặt cầu mô phỏng trái đất và mặt phẳng $(Oxy)$. Tại Bắc Cực $I(0;0;6\,400)$, Bộ chỉ huy thiết lập một \lq\lq Vòm An Ninh\rq\rq.Vùng an toàn này là một mặt cầu tâm I có bán kính quét là $R_q =1\,280\sqrt{10} $. Bất kỳ tàu nào đi vào biên giới của vùng này sẽ được bảo vệ tuyệt đối. Siêu điệp viên Alpha đang ẩn náu tại tọa độ $M(0;5\,120;3\,840)$. Anh nhận được mật lệnh:  Phải di chuyển từ vị trí hiện tại xuống đường Xích đạo để kích hoạt thiết bị định vị, sau đó lập tức rút lui về ranh giới Vòm An Ninh tại điểm nhập $K$ theo lộ trình ngắn nhất để tàu thoát. Biết vận tốc di chuyển của Alpha là $v_A =60$ km/h. Cùng lúc đó, tàu truy kích Beta đang phục kích tại vị trí $N(0;-6\,400;0)$. Ngay khi Alpha xuất phát, Beta tính toán được chính xác điểm $K$ và lao đến để chặn đầu. Tàu Beta phải duy trì vận tốc tối thiểu là bao nhiêu km/h để có thể chặn được Alpha trước lúc điệp viên này vào ranh giới vùng an toàn? (xem vận tốc của Alpha và Beta là chuyển động đều và hai tàu di chuyển trên mặt cầu. Làm tròn kết quả đến hàng phần chục)
	\par\shortans[0]{$84,6$}
	\loigiai{
		\begin{center}
			\begin{tikzpicture}[>=stealth,line join=round,line cap=round,font=\footnotesize,scale=1]
				\draw[->] (-4,0)--(4.5,0) node[below]{$x$};
				\draw [->] (0,-3.5)--(0,5) node [right]{$y$};
				\path (0,0) coordinate (O);
				\draw[black] (O) circle (3cm);
				\coordinate (H) at ($(O)+(0:3)$);
				\coordinate (M) at ($(O)+(25:3)$);
				\fill (2.7,0) circle (1pt) node[below] {$5\,120$};
				\fill (0,2.7) circle (1pt) node[left] {$5\,120$};
				\fill (1.3,0) circle (1pt) node[below] {$3\,840$};
				
				\coordinate (K) at ($(O)+(65:3)$);
				\coordinate (I) at ($(O)+(90:3)$);
				\coordinate (N) at ($(O)+(180:3)$);
				\draw[black] (I) circle (1.3cm);
				\draw [dashed] (0,1.3)--(M)--(2.7,0) (0,2.7)--(K)--(1.3,0) (K)--(N) ;
				\fill (0,1.3) circle (1pt) node[left] {$3\,840$};
				\foreach \t/\g in {H/30,M/20,K/30,I/60,N/-120}{
					\fill (\t) circle (1pt) node[shift={(\g:7pt)},font=\scriptsize]{$ \t $};
				}
				
			\end{tikzpicture}
		\end{center}
		Ta có $\overrightarrow{OI}=(0;0;6\,400)$, $\overrightarrow{OM}=(0;5\,120;3\,840)\Rightarrow \left[\overrightarrow{OI},\overrightarrow{OM}\right]=(-2\,048\,000;0;0)$.\\
		Suy ra $(OIM)\colon x=0$ hay $O$, $I$, $M$ thuộc mặt phẳng $(Oyz)$, ta có độ dài cung tròn $l=R\cdot\alpha $.\\
		Alpha đi theo hai cung $\wideparen{MH}$ và $\wideparen{HK}$. \\
		Ta có phương trình đường tròn 
		\begin{itemize}
			\item $(O)\colon y^2 +z^2 =6\,400^2$; \hfill$(1)$
			\item $(I)\colon y^2 +(z-6\,400)^2 =\left(1\,280\sqrt{10}\right)^2$. \hfill$(2)$
		\end{itemize}
		Lấy phương trình ($1$) trừ phương trình ($2$) ta được $z=5\,120\Rightarrow y=3\,840$. Suy ra $K(0;3\,840;5\,120)$.\\
		Độ dài cung Alpha đi là 
		\[l_{\wideparen{MH}} +l_{\wideparen{HK}} =\tan ^{-1} \left( \dfrac{3\,840}{5\,120} \right)\cdot 6\,400+\tan ^{-1} \left( \dfrac{5\,120}{3\,840} \right)\cdot 6\,400=3\,200\pi.\]
		Do đó thời gian Alpha đi là $t_{Alpha} =\dfrac{3200\pi}{60} $.\\
		Beta đi theo cung $\wideparen{NK}$, nên có $l_{\wideparen{NK}} =6\,400\left(\pi -\tan ^{-1} \dfrac{5\,120}{3\,840}\right) =14\,171{,}50359$.\\
		Vận tốc của tàu Beta là $v_{Beta} =\dfrac{14\,171{,}50\,359}{\dfrac{3\,200\pi}{60}} \approx 84{,}57993178$.\\
		Vậy $v_{Beta} \approx 84{,}6$ (km/h).
	}
=======
\setcounter{deso}{0}
\begin{name}
	{\tenchude}
	{TỔNG HỢP TRẢ LỜI NGẮN}
	{LỚP TOÁN THẦY PHÁT}
	{\lq\lq  It's not how much time you have, it's how you use it.\rq\rq}
\end{name}
\setcounter{ex}{0}
\begin{ex}%[0D0C2-3]%[KNTT - Lớp 12 - Dự án C - Đợt 2]%[Loc do]%[PB: Phan Anh]
	Bác Nghĩa đang giúp sắp xếp $16$ cuốn sách ôn thi vào một chiếc kệ có $5$ ngăn phân biệt, $16$ cuốn sách này thuộc $8$ môn học khác nhau: Toán, Lý, Hóa, Sinh, Sử, Địa, Văn, Anh. Mỗi môn học gồm đúng hai cuốn: một cuốn sách giáo khoa và một cuốn sách bài tập.  Để việc ôn tập đạt hiệu quả cao nhất theo từng khối thi, bác Nghĩa đặt ra các quy tắc khắc khe như sau: 
	\begin{itemize}
		\item Do ngăn kệ nhỏ, mỗi ngăn chỉ chứa được tối đa $5$ cuốn sách và không được để ngăn nào trống.
		\item Hai cuốn sách của cùng một môn học phải luôn nằm chung một ngăn với nhau.
		\item Các môn học trong cùng một tổ hợp môn thi phải nằm ở ba ngăn liên tiếp để thuận tiện cho việc tra cứu. Các tổ hợp bao gồm: (Văn, Sử, Địa); (Toán, Lý, Hóa); (Toán, Hóa, Sinh) và (Toán, Lý, Anh).
		\item Các cuốn sách trong mỗi ngăn được xếp theo hàng ngang với gáy sách quay ra ngoài ở mỗi ngăn, thứ tự từ trái qua phải.
	\end{itemize}
	Tổng số cách sắp xếp $16$ cuốn sách này vào $5$ kệ thỏa mãn điều kiện trên là $T$. Tính $\dfrac{T}{864} $.
	\par\shortans[0]{$6912$}
	\loigiai{
		Để xếp $16$ cuốn sách vào $5$ ngăn sao cho không có ngăn nào được để trống và các sách cùng môn sẽ được xếp chung ngăn thì sẽ xếp được $3$ ngăn chứa $4$ cuốn ($2$ môn) và $2$ ngăn chứa $2$ cuốn ($1$ môn).\\
		Số cách hoán vị sách trong các ngăn này (các sách cùng môn ở chung ngăn nhưng không nhất thiết phải ở cạnh nhau) là $(4!)^3 \cdot (2!)^2 $.\\
		Theo đề yêu cầu, ta đánh giá được như sau\\
		Ta thấy Toán tham gia vào $3$ tổ hợp liên tiếp nên cố định Toán ở ngăn $3$ để thỏa tính liền kề của ba tổ hợp môn.	Nhóm xã hội ( Văn, Sử, Địa) phải nằm ở ba ngăn liên tiếp nên hoán vị nội bộ của $3$ môn này là $3!$.\\
		Xét các trường hợp sau:\\
		\textbf{Trường hợp $1$}: Các môn tự nhiên xếp đều trên $5$ ngăn ( mỗi ngăn một môn tự nhiên). Khi đó có hai khả năng xảy ra là
		\begin{itemize}
			\item Sinh -- Hóa -- Toán -- Lý -- Anh.
			\item Anh -- Lý -- Toán -- Hóa -- Sinh.
		\end{itemize}
		Lúc này mỗi ngăn chỉ chứa $2$ cuốn của tự nhiên nên có thể đặt thêm ba môn Xã Hội vào bất kỳ $3$ ngăn liên tiếp trong $5$ ngăn đó là ($1,2,3$); ($2,3,4$) và ($3,4,5$).\\
		Vậy trường hợp này có: $2\cdot  3\cdot  3!=36$ (cách).\\
		\textbf{Trường hợp $2$}: Có $2$ môn tự nhiên nằm ở ngăn $2$. Khi đó có $4$ khả năng xảy ra là
		\begin{itemize}
			\item Anh ( ngăn $1$) -- Lý, Sinh ( ngăn $2$) -- Toán (ngăn $3$) -- Hóa (ngăn $4$);
			\item Hóa ( ngăn $1$) -- Lý, Sinh ( ngăn $2$) -- Toán (ngăn $3$) -- Anh (ngăn $4$);
			\item Sinh ( ngăn $1$) -- Hóa, Anh (ngăn $2$) -- Toán (ngăn $3$) -- Lý ( ngăn $4$)
			\item Lý (ngăn $1$) -- Hóa, Anh (ngăn $2$) -- Toán (ngăn $3$) -- Sinh (ngăn $4$)
		\end{itemize}
		Khi đó tổ hợp Xã Hội chỉ được xếp vào $3$ ngăn liên tiếp là ($3,4,5$).\\
		Vậy trường hợp này có: $4\cdot  1\cdot  3!=24$ (cách).\\
		\textbf{Trường hợp $3$}: Có một cặp tự nhiên nằm ở ngăn $4$. Trường hợp này là cách làm đối xứng của trường hợp $2$ nên có $24$ (cách).\\
		\textbf{Trường hợp $4$}: Có hai cặp môn tự nhiên cùng nằm ở ngăn $1$ và ngăn $2$. Khi đó có $2$ khả năng xảy ra là
		\begin{itemize}
			\item Lý, Sinh ( ngăn $1$) -- Hóa, Anh (ngăn $2$) -- Toán ( ngăn $3$);
			\item Hóa, Anh (ngăn $1$) -- Lý, Sinh ( ngăn $2$) -- Toán ( ngăn $3$)
		\end{itemize}
		Khi đó tổ hợp Xã Hội chỉ được xếp vào $3$ ngăn liên tiếp là ($3,4,5$).\\
		Vậy trường hợp này có $2\cdot  1\cdot  3!=12$ (cách).\\
		\textbf{Trường hợp $5$}: Có hai cặp môn tự nhiên cùng nằm ở ngăn $4$ và ngăn $5$. Đây là cách làm đối xứng với trường hợp $4$ nên có $12$ (cách).\\
		Suy ra tổng số xếp sách ở các trường hợp là $36+24+24+12+12=108$ ( cách).\\
		Vậy để xếp $18$ cuốn sách theo yêu cầu sẽ có $T=(4!)^3 \cdot  (2!)^2 \cdot  108=5\,971\,968$ (cách).\\
		Khi đó $\dfrac{T}{864} =6\,912$.}
\end{ex}

\begin{ex}%[0D0C2-5]%[Dự án C đề thi thử TN Môn Toán NH25-26 - Dot 4 - Phạm Đăng Đăng]%[THPT Lê Quý Đôn-Đống Đa-Hà Nội]%[GVPB - Phan Quốc Trí]
	Trong một trò chơi có thưởng bạn An được chơi một lần bằng cách lấy ngẫu nhiên ra $3$ quả cầu từ một hộp $150$ quả đánh số từ $1$ đến $150$. Bạn An sẽ nhận được phần thưởng nếu $3$ quả cầu bạn lấy ra thoả mãn đồng thời hai yêu cầu sau:\begin{itemize}
		\item Tổng các số được đánh số trên ba quả cầu  không vượt quả $300$.
		\item Các số trên $3$ quả cầu lấy được lập thành một cấp số cộng.
	\end{itemize}
	Xác suất để bạn An được phần thưởng là $p$, tính giá trị của $1\,000p$ (Kết quả là tròn đến hàng phần chục).
	\shortans{$7{,}8$}
	\loigiai{Chọn ngẫu nhiên $3$ quả cầu trong $150$ quả cầu nên $n(\Omega)=\mathrm{C}_{150}^3$.\\
		Giả sử 3 số chọn được là $a-d$; $a$; $a+d$ lập thành cấp số cộng với $a\ge 1$,  $d\ge 1$.\\
		Theo đề bài ta có
		$a-d+a+a+d\le 300 \Leftrightarrow 3a\le 300 \Leftrightarrow a\le 100
		$.\\
		Suy ra $\heva{&1\le a\le 100\\&a+d\le 150\\&a-d\ge 1}
		\Leftrightarrow\heva{&1\le a\le 100\\&d\le 150-a\\&d\le a-1.}$
		\begin{itemize}
			\item[\textbf{TH1:}] 
			$a-1\le 150-a \Leftrightarrow 2a\le 151 \Leftrightarrow a\le \dfrac{151}{2} \Leftrightarrow 1\le a\le 75.$\\
			Vậy có $75$ trường hợp xảy ra đối với $a$.\\
			Khi đó $d\le a-1 \Leftrightarrow 0< d\le 74$ nên có $74$ trường hợp xảy ra đối với $d$.\\
			Số trường hợp xảy ra là
			$\dfrac{75\cdot 74}{2}=2\,775$.
			\item[\textbf{TH2:}] $a-1\ge 150-a \Leftrightarrow 2a\ge 151 \Leftrightarrow a\ge \dfrac{151}{2} \Leftrightarrow 76\le a\le 100.$\\
			Khi đó $d\le 150-a \Leftrightarrow 150-100 \le d \le 150-76 \Leftrightarrow 50\le d\le 74 \ (\text{có }25\text{ số})$.
			Vậy số trường hợp xảy ra là $\dfrac{(50+74)\cdot 25}{2}=1\,550$.\\
			Tổng các trường hợp xảy ra là $2\,775+1\,550=4\,325$.
		\end{itemize}
		Vậy	xác suất cần tìm là $p=\dfrac{4\,325}{\mathrm{C}_{150}^3}\approx 0{,}0078469 \Rightarrow 1\,000\cdot p \approx 7{,}8.
		$}
\end{ex}

\begin{ex}%[Dự án đề TT Toán NH25-26-Dot 4- Hải Phụng]%[THCS-THPT LÊ THÁNH TÔNG - TP.HCM]%[GVPB - Viet Hoang Pham]%[0D0C2-6] 
	\immini[thm]{
		Trên một bàn cờ vua, Thông và Nghĩa cùng tham gia một trò chơi xác suất thú vị. Họ đặt một quân Mã tại ô D4 và một quân Vua tại ô F6. Quy ước rằng trong trò chơi này, hai quân cờ không ăn nhau. Trò chơi được tiến hành như sau:
		\begin{itemize}
			\item Thông: Thực hiện hành trình $4$ bước đi sao cho các ô đi qua trong $3$ bước đầu không trùng nhau và sau bước thứ $4$ Mã phải quay về lại đúng vị trí D4 ban đầu.
			\item Nghĩa: Thực hiện hành trình $3$ bước đi sao cho các ô đi qua trong $2$ bước đầu không trùng nhau và sau bước thứ $3$ Vua phải quay về lại đúng vị trí F6 ban đầu.
		\end{itemize}
	}
	{
		\begin{tikzpicture}[x=0.7cm, y=0.7cm]%Khai báo thêm gói chessfss để tạo quân cờ \usepackage{chessfss}
			%Thêm gói lệnh   \usepackage{chessfss}
			% 1. Định nghĩa màu sắc
			\definecolor{lightsq}{HTML}{F0D9B5} 
			\definecolor{darksq}{HTML}{B58863}  
			% 2. Vẽ bàn cờ
			\foreach \x in {0,...,7}
			\foreach \y in {0,...,7}
			{
				\pgfmathparse{mod(\x+\y,2) ? "darksq" : "lightsq"}
				\fill[\pgfmathresult] (\x,\y) rectangle (\x+1,\y+1);
			}
			% 3. Khung viền kép (Double border)
			\draw[thick] (0,0) rectangle (8,8);
			\draw[black, thin] (-0.8,-0.8) rectangle (8.8,8.8);
			\draw[black, thick] (-0.9,-0.9) rectangle (8.9,8.9);
			% 4. Tọa độ A-H
			\foreach \x [count=\i from 0] in {A,B,C,D,E,F,G,H} {
				\node at (\i+0.5, -0.4) {\small\textbf{\x}};
				\node at (\i+0.5, 8.4) {\small\textbf{\x}};
			}
			% 5. Tọa độ 1-8
			\foreach \y [count=\i from 0] in {1,2,3,4,5,6,7,8} {
				\node at (-0.4, \i+0.5) {\small\textbf{\y}};
				\node at (8.4, \i+0.5) {\small\textbf{\y}};
			}
			% 6. Đặt quân cờ bằng lệnh của gói skak
			% f6 tương ứng tọa độ (5.5, 5.5)
			% d4 tương ứng tọa độ (3.5, 3.5)
			% Quân Vua Đen (Black King)
			%\node at (5.5, 5.5) {\fontsize{35}{35}\selectfont \blackking}; 
			\node at (5.5, 5.5) {\fontsize{35}{35}\selectfont \BlackKingOnWhite};
			% Quân Mã Đen (Black Knight)
			\node at (3.5, 3.5) {\fontsize{35}{35}\selectfont \BlackKnightOnWhite};
			
		\end{tikzpicture}
	}
	\noindent Trình tự: Hai quân di chuyển xen kẽ nhau, bắt đầu bằng Mã.  Cụ thể: Mã đi bước $1$, Vua đi bước $1$, Mã đi bước $2$, Vua đi bước $2$,... cho đến khi cả hai hoàn tất hành trình. Xác suất để trong suốt quá trình di chuyển trên, có ít nhất một thời điểm Mã và Vua cùng đứng trên một ô cờ có dạng phân số tối giản  $\dfrac{a}{b}$ với $a, b \in \mathbb{N}^*$. Hãy tính giá trị của  $b-5a$.\\
	Cách di chuyển của quân Mã: Mã di chuyển theo đường chéo hình chữ nhật $2\times3$ ô vuông (hoặc $3\time2$ ô vuông).\\
	Cách di chuyển của quân Vua: Vua có thể di chuyển sang bất kỳ ô liền kề chung cạnh hoặc chung đỉnh xung quanh nó theo hướng ngang, dọc hoặc chéo.
	\shortans{$71$}
	\loigiai{
		\immini[thm]{
			Để dễ lập luận về sau, ta sẽ gọi hình vuông được tạo thành bởi vị trí ban đầu của Vua và các ô vuông xung quanh vị trí ban đầu này là hình vuông chứa Vua.\\
			Các bước đi của Mã và Vua sẽ được gọi là M1, M2, M3, M4, V1, V2, V3. Vị trí ban đầu của Mã và Vua gọi là M và V.\\
			Bây giờ ta sẽ tính số cách mà bộ Vua Mã có thể đi, đây chính là không gian mẫu.\\
			M1 có $8$ cách đi và tạo thành $2$ nhóm. Nhóm $1$ gồm $4$ vị trí mà cách lề bàn cờ nhiều hơn $1$ ô vuông, nhóm $2$ gồm $4$ vị trí cách ít nhất $1$ lề bàn cờ chỉ $1$ ô vuông.\\
			Để M4 về lại vị trí ban đầu thì bốn vị trí của Mã phải tạo thành hình thoi, khi đó M, M1, M2 phải không thẳng hàng.
		}
		{
			\begin{tikzpicture}[x=0.8cm, y=0.8cm]
				% 1. Định nghĩa màu sắc
				\definecolor{lightsq}{HTML}{F0D9B5} 
				\definecolor{darksq}{HTML}{B58863}  
				% 2. Vẽ bàn cờ
				\foreach \x in {0,...,7}
				\foreach \y in {0,...,7}
				{
					\pgfmathparse{mod(\x+\y,2) ? "darksq" : "lightsq"}
					\fill[\pgfmathresult] (\x,\y) rectangle (\x+1,\y+1);
				}
				% 3. Khung viền kép (Double border)
				\draw[thick] (0,0) rectangle (8,8);
				\draw[black, thin] (-0.8,-0.8) rectangle (8.8,8.8);
				\draw[black, thick] (-0.9,-0.9) rectangle (8.9,8.9);
				% 4. Tọa độ A-H
				\foreach \x [count=\i from 0] in {A,B,C,D,E,F,G,H} {
					\node at (\i+0.5, -0.4) {\small\textbf{\x}};
					\node at (\i+0.5, 8.4) {\small\textbf{\x}};
				}
				% 5. Tọa độ 1-8
				\foreach \y [count=\i from 0] in {1,2,3,4,5,6,7,8} {
					\node at (-0.4, \i+0.5) {\small\textbf{\y}};
					\node at (8.4, \i+0.5) {\small\textbf{\y}};
				}
				% Quân Vua Đen (Black King)
				\node at (5.5, 5.5) {\fontsize{35}{35}\selectfont \BlackKingOnWhite};
				% Khung bao từ góc dưới trái E5 (4,4) đến góc trên phải G7 (7,7)
				\draw[red, ultra thick] (4,4) rectangle (7,7);
				%Vẽ các quân cờ
				\node at (3.5, 3.5) {\fontsize{35}{35}\selectfont \BlackKnightOnWhite};
				\draw[solid, green!50!black, ultra thick] (3.5, 3.5) circle (0.42);
				% List tọa độ: B3(1,2), B5(1,4), C2(2,1), C6(2,5), E2(4,1), E6(4,5), F3(5,2), F5(5,4).
				% Cộng thêm 0.5 để vào tâm ô.
				\foreach \cx/\cy in {1.5/2.5, 1.5/4.5, 2.5/1.5, 2.5/5.5, 4.5/1.5, 4.5/5.5, 5.5/2.5, 5.5/4.5} {
					% 1. Vòng tròn nét đứt màu xám bao quanh ô đích
					\draw[dashed, black!50] (\cx, \cy) circle (0.35);
					% 2. Quân mã đen nhỏ trên ô tối
					\node at (\cx, \cy) {\fontsize{18}{18}\selectfont \BlackKnightOnBlack};
				}
			\end{tikzpicture}
		}
		\begin{itemize}
			\item  \textbf{Trường hợp 1:} M1 thuộc nhóm $1$. Khi đó M2 có $6$ cách đi do phải khác M và M, M1, M2 không thẳng hàng.
			\item  \textbf{Trường hợp 2:} M1 thuộc nhóm $2$. Khi đó M2 có $5$ cách đi do có $1$ vị trí nằm bên ngoài bàn cờ.
		\end{itemize}
		Vậy Mã có $4 \cdot 6 + 4 \cdot 5 = 44$ cách di chuyển.\\
		Để Vua có thể về vị trí đầu sau $3$ lần đi thì các bước đi của Vua phải nằm trong hình vuông chứa Vua đã nói ban đầu. Do đó, V1 có $8$ cách đi và chia thành $2$ nhóm.\\
		Nhóm $1$ gồm các ô vuông nhỏ chung đỉnh với hình vuông ban đầu và nhóm $2$ gồm các ô vuông nhỏ chung cạnh với hình vuông ban đầu.
		\begin{itemize}
			\item  \textbf{Trường hợp 1:} V1 thuộc nhóm $1$. Khi đó V2 có $2$ cách đi.
			\item  \textbf{Trường hợp 2: } V1 thuộc nhóm $2$. Khi đó V2 có $4$ cách đi.
		\end{itemize}
		Vậy Vua có $4 \cdot 2 + 4 \cdot 4 = 24$ cách.\\
		Do đó $n(\Omega) = 24 \cdot 44 = 1\,056$.\\
		Bây giờ ta sẽ tính số cách đi mà Mã và Vua có thể gặp nhau ít nhất $1$ lần.
		\begin{itemize}
			\item \textbf{Trường hợp 1:} M1 gặp V1.\\
			Khi đó M1 sẽ ở vị trí E6 hoặc F5. Với mỗi vị trí của M1 thì có $6$ cách chọn vị trí M2 nên Mã có $12$ cách di chuyển.\\
			Vua đi sau mã nên V1 sẽ phải đi vào ô M1. Với mỗi vị trí của V1 thì có $4$ cách chọn ví trí cho V2.\\
			Vậy số cách đi thoả mãn trường hợp này là $12 \cdot 4 =48$ cách.
			\item \textbf{Trường hợp 2:} M2 gặp V2.\\
			Để điều này xảy ra thì M2 phải nằm trong hình vuông chứa Vua nên M1 có $4$ cách chọn vị trí. Khi đó M2 có $2$ cách chọn vị trí tương ứng và điều nằm ở $4$ đỉnh hình vuông chứa Vua.\\
			Để V2 rơi vào ô M2 thì V1 phải nằm kề với ô mà M2 sẽ đi tới nên V1 có $2$ cách chọn vị trí.\\
			Vậy số cách đi thoả mãn trường hợp này là $4 \cdot 2 \cdot 2 = 16$ cách.
			\item  \textbf{Trường hợp 3:} M2 gặp V1.\\
			Cách di chuyển của M2 trong trường hợp này giống hệt \textbf{trường hợp 2}.\\
			Để V1 rơi vào ô M2 thì V1 sẽ đi vào ô mà M2 sẽ đi, khi đó V1 nằm ở  đỉnh hình vuông chứa Vua nên V2 sẽ có $2$ cách đi. \\
			Vậy có $4 \cdot 2 \cdot 2 = 16$ cách.
			\item  \textbf{Trường hợp 4:} M3 gặp V2.\\
			Do tính đối xứng nên trường hợp này thật ra là đi theo chiều ngược lại của trường hợp 1 nên cũng có $48$ cách.\\
			Tuy nhiên, có $2$ cách đi mà M1 gặp V1 thuộc trường hợp này.
		\end{itemize}
		Vậy số cách đi thoả mãn là $48+16+16+48-2=126$.\\
		Vậy xác suất để Mã và Vua gặp nhau ít nhất 1 lần là $\dfrac{126}{1\,056} = \dfrac{21}{176}$.\\
		Do đó $b-5a = 71$.
	}
\end{ex}

\begin{ex}%[0D0C2-7]%[Dự án C NH25-26 - Đợt 3 - Hector Tran]%[SGD-Tuyên Quang - Lần 1]%[GVPB: Thanh Phát]
	Gọi $S$ là tập hợp tất cả các số tự nhiên có $6$ chữ số khác nhau được lập từ các chữ số $1$, $2$, $3$, $4$, $5$, $6$. Lấy ngẫu nhiên một số thuộc $S$, gọi $T$ là xác suất số lấy được là số lẻ đồng thời tổng của ba chữ số đầu lớn hơn tổng của ba chữ số cuối một đơn vị. Giá trị của $230T$ bằng bao nhiêu?
	\par\shortans{23}
	\loigiai{
		Số phần tử của tập hợp $S$ tất cả các số tự nhiên có $6$ chữ số khác nhau được lập từ các chữ số $1$, $2$, $3$, $4$, $5$, $6$ là $n(\Omega) = 6! = 720$.\\
		Gọi $A$ là biến cố chọn từ $S$ được một số lẻ đồng thời tổng của ba chữ số đầu lớn hơn tổng của ba chữ số cuối một đơn vị.\\
		Giả sử số được chọn có dạng $\overline{a_1a_2a_3a_4a_5a_6}$. Theo giả thiết
		\begin{itemize}
			\item Số là số lẻ nên $a_6 \in \{1; 3; 5\}$.
			\item Tổng tất cả các chữ số là $1 + 2 + 3 + 4 + 5 + 6 = 21$.
			\item Gọi $S_1 = a_1 + a_2 + a_3$ là tổng ba chữ số đầu và $S_2 = a_4 + a_5 + a_6$ là tổng ba chữ số cuối.
		\end{itemize}
		Ta có hệ phương trình
		\[\heva{&S_1 + S_2 = 21 \\ &S_1 = S_2 + 1} \Leftrightarrow \heva{&S_1 = 11 \\ &S_2 = 10.}\]
		Các bộ gồm $3$ chữ số có tổng bằng $10$ từ tập $\{1; 2; 3; 4; 5; 6\}$ là $\{1; 3; 6\}$; $\{1; 4; 5\}$; $\{2; 3; 5\}$.
		\begin{itemize}
			\item Trường hợp $1$. $\{a_4; a_5; a_6\} = \{1; 3; 6\}$. Khi đó bộ còn lại là $\{a_1; a_2; a_3\} = \{2; 4; 5\}$.\\
			Vì số là số lẻ nên $a_6 \in \{1; 3\}$ ($2$\ cách chọn).\\
			Hai vị trí $a_4$; $a_5$ có $2!$ cách hoán vị.\\
			Ba vị trí $a_1$; $a_2$; $a_3$ có $3!$ cách hoán vị.\\
			Số các số thỏa mãn là $2 \cdot 2! \cdot 3! = 24$.
			\item Trường hợp $2$. $\{a_4; a_5; a_6\} = \{1; 4; 5\}$. Khi đó bộ còn lại là $\{a_1; a_2; a_3\} = \{2; 3; 6\}$.\\
			Vì số là số lẻ nên $a_6 \in \{1; 5\}$ ($2$\ cách chọn).\\ 
			Tương tự trường hợp $1$; có $2 \cdot 2! \cdot 3! = 24$ số.
			\item Trường hợp $3$. $\{a_4; a_5; a_6\} = \{2; 3; 5\}$. Khi đó bộ còn lại là $\{a_1; a_2; a_3\} = \{1; 4; 6\}$.\\
			Vì số là số lẻ nên $a_6 \in \{3; 5\}$ ($2$\ cách chọn).\\ 
			Tương tự trường hợp $1$; có $2 \cdot 2! \cdot 3! = 24$ số.
		\end{itemize}
		Tổng số kết quả thuận lợi cho biến cố $A$ là $n(A)=24 + 24 + 24 = 72$.\\
		Xác suất $T =\dfrac{n(A)}{n(\Omega)}= \dfrac{72}{720} = \dfrac{1}{10} = 0{,}1$.\\
		Vậy $230T = 230 \cdot 0{,}1 = 23$.
	}
\end{ex}

\begin{ex}%[0D0C2-9]%[Dự án C đề thi thử TN Môn Toán NH25-26 - Dot 4 - Phạm Hoàng Đăng]%[THPT Nguyễn Trãi-Hà Nội]%[GVPB - Khắc Thiên]
	\immini[thm]{Chọn ngẫu nhiên $5$ số phân biệt từ tập hợp $X=\{2; 4; 6; 8; 10; 12; 13; 14; 16; 18; 20\}$ để xếp vào $5$ ô $A$, $B$, $C$, $D$, $O$ như hình vẽ (mỗi ô xếp một số). Gọi $Y$ là biến cố \lq\lq Số $13$ xếp ở ô $O$ và tổng các số xếp ở các ô $A$, $O$, $C$ bằng tổng các số xếp ở các ô $B$, $O$, $D$\rq\rq. Biết xác suất của biến cố $Y$ là $\mathrm{P}(Y)=\dfrac{a}{b}$ (với $a$, $b \in \mathbb{N}$, $\dfrac{a}{b}$ là phân số tối giản). Khi đó giá trị $a+b$ bằng bao nhiêu?}{
		\begin{tikzpicture}[scale=1,>=stealth, line join=round, line cap=round, smooth]
			\def\R{1.5}
			\path (0,0) coordinate (O)
			(90:\R) coordinate (A)
			(180:\R) coordinate (B)
			(270:\R) coordinate (C)
			(0:\R) coordinate (D);
			\draw[blue!60!black] (O) -- (A) (O) -- (B) (O) -- (C) (O) -- (D);
			\foreach \p in {O,A,B,C,D} {
				\node[draw=green!60!black, circle, fill=white, inner sep=4pt, font=\large\itshape] at (\p) {\p};
			}
		\end{tikzpicture}
	}
	\par
	\shortans[oly]{698}
	\loigiai{
		Tập $X$ có $11$ phần tử.\\
		Chọn ngẫu nhiên $5$ phần tử từ tập $X$ và xếp vào $5$ vị trí.\\ Khi đó $n(\Omega)=\mathrm{A}_{11}^5=55\,440$.\\
		Xét biến cố $Y$: \lq\lq Số $13$ xếp ở ô $O$ và tổng các số xếp ở các ô $A$, $O$, $C$ bằng tổng các số xếp ở các ô $B$, $O$, $D$\rq\rq.
		Ta có:\\
		Số $13$ được xếp ở ô $O$ có $1$ (cách).\\
		Các số xếp ở các ô $A$, $B$, $C$, $D$ thỏa mãn $A+C=B+D=S$.\\ Ta chọn các bộ thỏa mãn:
		\begin{itemize}
			\item $S=10 \Rightarrow \{2;8\}$, $\{4;6\}$ có $1$ (bộ).
			\item $S=12 \Rightarrow \{2;10\}$, $\{4;8\}$ có $1$ (bộ).
			\item $S=14 \Rightarrow \{2;12\}$ và $\{4;10\}$, $\{2;12\}$ và $\{6;8\}$, $\{4;10\}$ và $\{6;8\}$ có $3$ (bộ).
			\item $S=16 \Rightarrow \{2;14\}$ và $\{6;10\}$, $\{2;14\}$ và $\{4;12\}$, $\{4;12\}$ và $\{6;10\}$ có $3$ (bộ).
			\item $S=18 \Rightarrow \{2;16\}$, $\{4;14\}$, $\{6;12\}$, $\{8;10\}$ suy ra chọn $2$ trong $4$ cặp có $\mathrm{C}_4^2=6$ (bộ).
			\item $S=20 \Rightarrow \{2;18\}$, $\{4;16\}$, $\{6;14\}$, $\{8;12\}$ suy ra chọn $2$ trong $4$ cặp có $\mathrm{C}_4^2=6$ (bộ).
			\item $S=22 \Rightarrow \{2;20\}$, $\{4;18\}$, $\{6;16\}$, $\{8;14\}$, $\{10;12\}$ suy ra chọn $2$ trong $5$ cặp có $\mathrm{C}_5^2=10$ (bộ).
			\item $S=24 \Rightarrow \{4;20\}$, $\{6;18\}$, $\{8;16\}$, $\{10;14\}$ suy ra chọn $2$ trong $4$ cặp có $\mathrm{C}_4^2=6$ (bộ).
			\item $S=26 \Rightarrow \{6;20\}$, $\{8;18\}$, $\{10;16\}$, $\{12;14\}$ suy ra chọn $2$ trong $4$ cặp có $\mathrm{C}_4^2=6$ (bộ).
			\item $S=28 \Rightarrow \{8;20\}$, $\{10;18\}$, $\{12;16\}$ suy ra chọn $2$ trong $3$ cặp có $\mathrm{C}_3^2=3$ (bộ).
			\item $S=30 \Rightarrow \{10;20\}$, $\{12;18\}$, $\{14;16\}$ suy ra chọn $2$ trong $3$ cặp có $\mathrm{C}_3^2=3$ (bộ).
			\item $S=32 \Rightarrow \{12;20\}$, $\{14;18\}$ có $1$ (bộ).
			\item $S=34 \Rightarrow \{14;20\}$, $\{16;18\}$ có $1$ (bộ).
		\end{itemize}
		Tổng số các bộ thỏa mãn là $1+1+3+3+6+6+10+6+6+3+3+1+1=50$ (bộ).\\
		Ứng với mỗi bộ ta chọn $1$ cặp xếp vào ô $A-C$, cặp còn lại cho ô $B-D$ có $2$ (cách).\\
		Hoán vị $2$ số trong cặp $A$, $C$ có $2$ (cách).\\
		Hoán vị $2$ số trong cặp $B$, $D$ có $2$ (cách).\\
		Suy ra mỗi bộ có $2 \cdot 2 \cdot 2=8$ (cách).\\
		Số phần tử thuận lợi của biến cố $Y$ là $n(Y)=50 \cdot 8=400$.\\
		Do đó $\mathrm{P}(Y)=\dfrac{n(Y)}{n(\Omega)}=\dfrac{400}{55\,440}=\dfrac{5}{693} \Rightarrow a+b=698$.
	}
\end{ex}

\begin{ex}%[0D0C2-9]%[Dự án C đợt 3 NH25-26- Phạm Đăng Đăng]%[TT-THPT-TranNhanTong-HaNoi-N25-26]%[GVPB - Phan Quốc Trí]
	Kết thúc học kì $1$, cô giáo chuẩn bị $8$ hộp bút bi, $7$ tập vở viết và $5$ bộ dụng cụ vẽ hình (thước kẻ, compa, êke, thước đo độ) để làm phần thưởng cho $10$ học sinh có kết quả xuất sắc nhất lớp. Mỗi học sinh nhận thưởng sẽ được $2$ phần thưởng khác loại. Trong số $10$ học sinh trên có $2$ học sinh tên Dũng và Nam. Tìm xác suất để Dũng và Nam có phần thưởng giống nhau.   
	\shortans{$0{,}31$}
	\loigiai{
		Giả sử số phần thưởng gồm $1$ hộp bút và $1$ tập vở là $x$ ($x \in \mathbb{N}$).\\
		Số phần thưởng gồm $1$ tập vở và $1$ bộ dụng cụ là $y$ ($y \in \mathbb{N}$).\\
		Số phần thưởng gồm $1$ hộp bút và $1$ bộ dụng cụ là $z$ ($z \in \mathbb{N}$).\\
		Vì cô có $8$ hộp bút bi, $7$ tập vở viết và $5$ bộ dụng cụ nên
		$
		\heva{&x + z = 8 \\ &x + y = 7 \\ &y + z = 5} \Leftrightarrow \heva{&x = 5 \\ &y = 2 \\ &z = 3.}
		$\\
		Tổng số cách tặng thưởng là $\mathrm{C}_{10}^5 \cdot \mathrm{C}_5^2 \cdot \mathrm{C}_3^3 = 2\,520$.\\
		Để các bạn nhận thưởng sao cho Dũng và Nam nhận thưởng như nhau, ta chia thành các trường hợp:
		\begin{itemize}
			\item \textbf{TH1:} Dũng và Nam nhận $1$ bút và $1$ vở, còn lại $3$ phần gồm $1$ bút và $1$ vở, $2$ phần $1$ vở và $1$ dụng cụ, $3$ phần $1$ bút và $1$ dụng cụ. Ta chọn $3$ phần gồm $1$ bút và $1$ vở cho $8$ bạn, chọn $2$ phần gồm $1$ tập vở và $1$ bộ dụng cụ cho $5$ bạn và chọn $3$ phần gồm $1$ hộp bút và $1$ bộ dụng cụ cho $3$  bạn còn lại.\\ Khi đó
			số cách chọn ở trường hợp này là $\mathrm{C}_8^3 \cdot \mathrm{C}_5^2 \cdot \mathrm{C}_3^3 = 560$.
			\item \textbf{TH2:} Dũng và Nam nhận $1$ vở và $1$ dụng cụ, còn lại $5$ phần gồm $1$ bút và $1$ vở, $3$ phần $1$ bút và $1$ dụng cụ nên có số cách phát $\mathrm{C}_8^5 \cdot \mathrm{C}_3^3 = 56$.
			\item 	\textbf{TH3:} Dũng và Nam nhận $1$ bút và $1$ dụng cụ, còn lại $5$ phần gồm $1$ bút và $1$ vở, $2$ phần $1$ vở và $1$ dụng cụ, 1 phần bút và dụng cụ nên có số cách phát $\mathrm{C}_8^5 \cdot \mathrm{C}_3^2 = 168$.
		\end{itemize}
		Suy ra tổng số cách phát phần thưởng là $560 + 56+168 = 784$.\\
		Vậy	xác suất cần tìm  $\dfrac{784}{2\,520} \approx 0{,}31$.}
\end{ex}

\begin{ex} %[0D0V2-5]%[TN-THPT-NguyenTrungThien-HaTinh-NH25-26]%[Loc Do]%[PB: Phan Anh]
	Một người môi giới bất động sản có $8$ chìa khoá để mở $8$ ngôi nhà mới. Mỗi chìa chỉ mở được đúng $1$ ngôi nhà. Biết có $3$ ngôi nhà thường không khoá cửa, người môi giới chọn ngẫu nhiên $3$ chìa khoá mang theo. Hỏi nếu người môi giới chọn ngẫu nhiên $1$ ngôi nhà để vào, thì xác suất để người môi giới này có thể vào được là bao nhiêu (\textit{kết quả làm tròn đến hàng phần trăm})?
	\shortans{0,61}
	\loigiai{
		Chọn $1$ nhà từ $8$ nhà có $8$ cách, sau đó chọn $3$ chìa từ $8$ chìa có $\mathrm{C}_{8}^{3}$ cách.\\
		Suy ra số phần tử của không gian mẫu là
		\[n(\Omega)=8\cdot \mathrm{C}_{8}^{3} =448.\]
		Gọi $A$ là biến cố: \lq\lq Để người môi giới này có thể vào được\rq\rq.
		\begin{itemize}
			\item Trường hợp $1$: Chọn đúng một trong ba nhà không khoá cửa, chọn $3$ chìa từ $8$ chìa bất kỳ (khi đó không cần quan tâm chìa chọn như thế nào). Nên có $3\cdot\mathrm{C}_{8}^{3} =168$ cách
			\item Trường hợp $2$: Chọn đúng nhà có khoá cửa.
			\begin{itemize}
				\item Chọn $1$ nhà từ $5$ nhà có $5$ cách;
				\item Chọn $1$ chìa đúng, và $2$ chìa còn lại từ $7$ chìa có: $1\cdot\mathrm{C}_{7}^{2} =21$ cách.
			\end{itemize}
		\end{itemize}
		Suy ra có $5\cdot 21=105$ cách.\\
		Ta có số phần tử của biến cố $A$ là $n(A)=168+105=273$.
		\[\Rightarrow \mathrm{P}(A)=\dfrac{n(A)}{n(\Omega)} =\dfrac{39}{64} \approx 0{,}61.\]}
\end{ex}

\begin{ex}%[0D0V2-7]%[Dự án đề thi thử NH25-26]%[THPT Chuyên Phan Bội Châu - Nghệ An]%[Sỹ Trường]
	Trong một hộp kín đựng $20$ viên bi được đánh số từ $1$ đến $20$. Bạn An chọn ngẫu nhiên $4$ viên bi trong hộp, xác suất để $4$ số trên $4$ viên bi được chọn lập thành một cấp số cộng là $\dfrac{1}{a}$. Tính $a$.
	\shortans{$85$}
	\loigiai{
		Số cách chọn $4$ viên bi từ hộp là $\mathrm{C}_{20}^4 = 4\,845$.\\
		Gọi cấp số cộng tạo thành là $a$, $b$, $c$, $d$. \\
		Ta có công sai $x = \dfrac{d-a}{3}$, do đó $a$, $d$ có cùng số dư khi chia cho $3$.\\
		Từ $1$ đến $20$ ta có $6$ số có dạng $3k$, có $7$ số có dạng $3k+1$ và $7$ số có dạng $3k+2$ với $k \in \mathbb{N}$.\\
		Mỗi cặp $a$, $d$ (với $a<d$) thỏa mãn tính chất trên sẽ xác định duy nhất hai số $b$, $c$ tạo thành cấp số cộng gồm $4$ phần tử trong tập hợp đã cho.\\
		Suy ra số cấp số cộng bằng số cách chọn cặp $a$, $d$ và bằng $\mathrm{C}_6^2+\mathrm{C}_7^2+\mathrm{C}_7^2 = 57$.\\
		Vậy xác suất là $\mathrm{P}(A) = \dfrac{57}{4\,845} = \dfrac{1}{85} \Rightarrow a=85$.
	}
\end{ex}

\begin{ex}%[0D0V2-8]%[Dự án C đề thi thử TN Môn Toán NH25-26 - Dot 4 - Vũ Hồng Toàn]%[THPT Mai Anh Tuấn - Thanh Hóa]%[GVPB - Lê Hữu Kiệt]
	Có $8$ bạn cùng ngồi xung quanh một cái bàn tròn, mỗi bạn cầm một đồng xu như sau. Tất cả $8$ bạn cùng tung đồng xu của mình, bạn có đồng xu ngửa thì đứng, bạn có đồng xu sấp thì ngồi. Xác suất để không có hai bạn liền kề cùng đứng là bao nhiêu {\it (kết quả được làm tròn đến hàng phần trăm)}?
	\par\shortans{0,18}
	\loigiai{
		Gọi $A$ là biến cố không có hai người liền kề cùng đứng.\\
		Số phần tử của không gian mẫu là $n(\Omega) = 2^8 = 256$.\\
		Rõ ràng nếu nhiều hơn $4$ đồng xu ngửa thì biến cố $A$ không xảy ra.        \\
		Để biến cố $A$ xảy ra có các trường hợp sau
		\begin{enumerate}[{\bf TH 1.}]
			\item Có nhiều nhất $1$ đồng xu ngửa.
			\begin{itemize}
				\item $0$ ngửa có $1$ cách.
				\item $1$ ngửa: chọn $1$ trong $8$ vị trí có $8$ cách.
			\end{itemize}
			\item Có $2$ đồng xu ngửa.\\
			Hai đồng xu ngửa kề nhau: có $8$ khả năng.\\
			Suy ra số kết quả của trường hợp này là $\mathrm{C}_8^2 - 8 = 20$.
			\item Có $3$ đồng xu ngửa.\\
			Cả $3$ đồng xu ngửa kề nhau: có $8$ kết quả.\\
			Trong $3$ đồng xu ngửa, có đúng một cặp kề nhau có $8\cdot 4 = 32$ cách.\\
			Suy ra số kết quả của trường hợp này là $\mathrm{C}_8^3 - 8 - 32 = 16$ cách.
			\item Có $4$ đồng xu ngửa.\\
			Trường hợp này có $2$ kết quả thỏa mãn biến cố $ A$ xảy ra.
		\end{enumerate}
		Suy ra $n(A) = 9 + 20 + 16 + 2 = 47$.\\
		Xác suất để không có hai bạn liền kề cùng đứng là $\mathrm{P}(A) = \dfrac{n(A)}{n(\Omega)} = \dfrac{47}{256} \approx 0{,}18$.
	}
\end{ex}

\begin{ex}%[0D0V2-9]%[TT-LienTruong-HaNoi-N25-26]%[Phạm Quốc Toàn - Lê Phúc]
	\immini[thm]{Cho một bảng ô vuông $3\times 3$ như hình vẽ. Điền ngẫu nhiên $9$ số thuộc tập hợp $X=\left\{0;1;2;3;4;5;6;7;8;9\right\}$ vào $9$ ô vuông trong bảng (mỗi ô điền một số khác nhau). Gọi $Y$ là biến cố \lq\lq Mỗi hàng, mỗi cột bất kì trong bảng đều có ít nhất một số lẻ\rq\rq. Biết xác suất $P(Y)=\dfrac{a}{b}$ (với $a$, $b\in \mathbb{N}$ và $\dfrac{a}{b}$ là phân số tối giản). Khi đó $a+b$ bằng bao nhiêu?}
	{	\begin{tikzpicture}[font=\footnotesize, line join=round, line cap=round, >=stealth,scale=0.7]
			%%%%%%%%%%%%%%%%%%%%%%%%%%%
			\draw[thick] (0,0) -- (3,0) -- (3,3) -- (0,3) -- (0,0) (1,0) -- (1,3) (2, 0) -- (2,3) (0, 1) -- (3, 1) (0,2) -- (3,2);
			%%%%%%%%%%%%%%%%%%%%%%%%%%%
	\end{tikzpicture}}
	\par\shortans[0]{43}
	\loigiai{
		Số phần tử của không gian mẫu là $n(\Omega)=\mathrm{A}_{10}^9$.\\
		Ta có $\overline{Y}$ là biến cố \lq\lq Có ít nhất một hàng hoặc một cột toàn chữ số chẵn\rq\rq.\\
		Tập $X$ gồm $5$ chữ số chẵn và $5$ chữ số lẻ nên ta có hai trường hợp sau.
		\begin{itemize}
			\item \textbf{Trường hợp $1$.} Bảng được xếp $5$ chữ số lẻ và $4$ chữ số chẵn\\
				Do chỉ có $4$ chữ số chẵn nên chỉ có tối đa đúng $1$ cột hoặc $1$ hàng toàn số chẵn.
			\begin{itemize}
				\item Chọn bất kì một hàng hoặc một cột để xếp toàn số chẵn có $6$ cách.
				\item Chọn $5$ trong $6$ ô còn lại và xếp $5$ số lẻ vào có $\mathrm{A}_6^5$ cách.
				\item Chọn $4$ trong $5$ số chẵn và xếp vào $4$ ô còn lại có $\mathrm{A}_5^4$ cách.
			\end{itemize}
			Trường hợp này có $6 \cdot \mathrm{A}_6^5 \cdot \mathrm{A}_5^4=518\,400$ cách xếp thỏa mãn yêu cầu đề bài.
			\item \textbf{Trường hợp $2$.} Bảng được xếp $4$ chữ số lẻ và $5$ chữ số chẵn
			\begin{itemize}
				\item Chọn bất kì một hàng hoặc một cột để xếp toàn số chẵn có $6$ cách.
				\item Chọn $4$ trong $6$ ô còn lại để xếp số lẻ vào có $\mathrm{C}_6^4$ cách.
				\item Chọn $4$ trong $5$ số lẻ và xếp vào $4$ ô đã chọn có $\mathrm{A}_5^4$ cách.
				\item Xếp nốt $5$ số chẵn vào $5$ ô còn lại có $5!$ cách.
			\end{itemize}
			Do $5$ số chẵn có thể xếp đầy một hàng và một cột (chung nhau một ô) nên ta đã đếm lặp trường hợp có cả một cột và một hàng toàn số chẵn, trường hợp này có $3 \cdot 3 \cdot \mathrm{A}_5^4 \cdot 5!$ cách.\\
			Suy ra trường hợp này có $6\cdot \mathrm{C}_6^4 \cdot \mathrm{A}_5^4\cdot 5!-3 \cdot 3\cdot \mathrm{A}_5^4 \cdot 5!=1\,166\,400$ cách.
		\end{itemize} 
		Vậy xác suất cần tính là $\mathrm{P}(Y)=1-\mathrm{P} \left(\overline{Y} \right)=1-\dfrac{518\,400+1\,166\,400}{\mathrm{A}_{10}^9}=\dfrac{15}{28}$.\\
		Khi đó $a+b=43$.
	}
\end{ex}

\begin{ex}%[0D1V3-5]%[Dự án C đề thi thử TN Môn Toán NH25-26 - Dot 4 - Hector Tran]%[Liên trường THPT cụm 09 - Đợt 2 - Hà Nội]%[GVPB - Thanh Phát]
	Lớp 10C có $45$ học sinh trong đó có $25$ em thích môn Văn, $20$ em thích môn Toán, $18$ em thích môn Sử, $6$ em không thích môn nào, $5$ em thích cả ba môn. Có bao nhiêu học sinh thích chỉ một trong ba môn học trên?
	\par\shortans{20}
	\loigiai{
		\immini {
			Gọi $a$, $b$, $c$ theo thứ tự là số học sinh chỉ thích môn Văn, Sử, Toán.\\
			Gọi $x$ là số học sinh chỉ thích hai môn là Văn và Toán.\\
			Gọi $y$ là số học sinh chỉ thích hai môn là Sử và Toán.\\
			Gọi $z$ là số học sinh chỉ thích hai môn là Văn và Sử.\\
			Ta có số em thích ít nhất một môn là $45 - 6 = 39$.\\
			Dựa vào biểu đồ Ven ta có hệ phương trình
			\[\heva{&a + x + z + 5 = 25&\quad (1) \\ &b + y + z + 5 = 18&\quad (2) \\ &c + x + y + 5 = 20&\quad (3) \\ &x + y + z + a + b + c + 5 = 39.& \quad (4)}\]
		}{
			\begin{tikzpicture}[scale=0.6, node/.style={scale=0.8}, declare function={
					xV=0; yV=0; rV=2.8;
					xT=2.6; yT=0.8; rT=1.9;
					xS=2.2; yS=-2.1; rS=2.3;
					nx25=-0.6; ny25=1.4;
					nxa=-0.6; nya=-0.2;
					nx20=3.2; ny20=1.5;
					nxc=3.6; nyc=0.6;
					nx18=2.2; ny18=-3.8;
					nxb=2.2; nyb=-2.1;
					nxx=1.4; nyx=1.1;
					nxy=3.1; nyy=-0.45;
					nxz=1; nyz=-1.3;
					nx5=1.9; ny5=-0.4;
				}]
				\tikzset{every node/.style={scale=0.8}}  
				\draw[line width=0.8pt] (xV, yV) circle (rV);
				\draw[line width=0.8pt] (xT, yT) circle (rT);
				\draw[line width=0.8pt] (xS, yS) circle (rS);				
				\node at (nx25, ny25) {$25(V)$};
				\node at (nxa, nya) {$a$};				
				\node at (nx20, ny20) {$20(T)$};
				\node at (nxc, nyc) {$c$};				
				\node at (nx18, ny18) {$18(S)$};
				\node at (nxb, nyb) {$b$};				
				\node at (nxx, nyx) {$x$};
				\node at (nxy, nyy) {$y$};
				\node at (nxz, nyz) {$z$};				
				\node at (nx5, ny5) {$5$};				
			\end{tikzpicture}
		}	\noindent
		Cộng vế với vế của phương trình $(1)$, $(2)$, $(3)$ ta có
		\begin{eqnarray*}
			&&(a + b + c) + 2(x + y + z) + 15 = 25 + 18 + 20 \\
			&\Leftrightarrow& (a + b + c) + 2(x + y + z) = 48. \quad (5)
		\end{eqnarray*}
		Từ $(4)$ ta có $x + y + z = 39 - 5 - (a + b + c) = 34 - (a + b + c)$.\\
		Thay $x + y + z = 34 - (a + b + c)$ vào $(5)$ ta được
		\begin{eqnarray*}
			&&(a + b + c) + 2[34 - (a + b + c)] = 48 \\
			&\Leftrightarrow& -(a + b + c) + 68 = 48 \\
			&\Leftrightarrow& a + b + c = 20.
		\end{eqnarray*}
		Vậy có $20$ em thích chỉ một trong ba môn trên.
	}
\end{ex}

\begin{ex}%[Đề thi thử môn Toán TRƯỜNG THPT Cổ Loa - Hà Nội năm học 2025-2026]%[Phạm Quốc Toàn - Phạm Hoàng Việt]%[0D2H1-3]
	Mỗi tuần, một cửa hàng bán điện thoại di động trung bình bán được $1\,000$ chiếc điện thoại $A$ với giá $14$ triệu đồng một chiếc. Biết rằng, nếu cứ giảm giá bán $500$ nghìn đồng một chiếc thì số lượng điện thoại $A$ bán ra sẽ tăng thêm $100$ chiếc mỗi tuần. Mặt khác, để bán được $n$ chiếc điện thoại $A$ mỗi tuần thì chi phí hàng tuần của cửa hàng là $C(n)=12\,000-3n$ (triệu đồng). Để lợi nhuận thu được lớn nhất, cửa hàng nên bán mỗi chiếc điện thoại $A$ với giá bao nhiêu triệu đồng?
	\par\shortans[0]{8}
	\loigiai{
		Gọi $x$ ($x\in \left[0;14\right]$; triệu đồng) là giá bán một chiếc điện thoại $A$.\\
		Đổi $500$ nghìn đồng bằng $0{,}5$ triệu đồng.\\
		Số lần giảm giá là $\dfrac{14-x}{0{,}5}$ .\\
		Số điện thoại bán ra là $1\,000+\dfrac{14-x}{0{,}5} \cdot 100=-200 x+3800=n$.\\
		Doanh thu của cửa hàng là $(-200x+3\,800) x$ (triệu đồng).\\
		Chi phí của cửa hàng là $12\,000-3(-200x+3\,800)=600x+600$ (triệu đồng).\\
		Lợi nhuận của cửa hàng là $f(x)=(-200x+3\,800)x-(600x+600)=-200x^2+3\,200x-600$. \\ 
		Ta có $f(x)$ là hàm số bậc hai nên đạt giá trị lớn nhất tại $x=-\dfrac{b}{2a}=8$.\\
		Để lợi nhuận thu được lớn nhất, cửa hàng nên bán mỗi chiếc điện thoại $A$ với giá là $8$ triệu đồng.
	}
\end{ex}

\begin{ex}%[0D2H2-3]%[Dự án C - Đợt 3 - Sở Giáo Dục Sơn La]%[Thy Nguyen Vo Diem - Nguyễn Tiến]
	Trong một cuộc thi làm bánh, mỗi đội chỉ được sử dụng tối đa $6$ kg bột mì và $4$ kg đậu để làm ra
	bánh loại I và bánh loại II. Để làm ra một chiếc bánh loại I cần $0{,}06$ kg bột mì và $0{,}08$ kg đậu, để làm ra một chiếc bánh loại II cần $0{,}08$ kg bột mì và $0{,}04$ kg đậu. Mỗi chiếc bánh loại I được $8$ điểm, mỗi chiếc bánh loại II được $10$ điểm. Số điểm tối đa có thể đạt được của mỗi đội bằng bao nhiêu?
	\shortans{760}
	\loigiai{
		Gọi $x$, $y$ lần lượt là số bánh loại I và loại II đội làm được ($x$, $y \ge 0$; $x$, $y \in \mathbb{Z}$).\\
		Theo giả thiết, ta có hệ bất phương trình $\heva{&x \ge 0\\&y \ge 0\\&0{,}06x+0{,}08y \le 6\\&0{,}08x+0{,}04y \le 4} \Leftrightarrow \heva{&x \ge 0\\&y \ge 0\\&3x+4y \le 300\\&2x+y \le 100.}$\\
		Số điểm đội đạt được là $P=8x+10y$.\\
		Miền nghiệm hệ bất phương trình là miền tứ giác $OABC$ với $O(0;0)$, $A(50;0)$, $B(20;60)$, $C(0;75)$.\\
		\begin{center}
			\begin{tikzpicture}[scale=0.7, font=\footnotesize, line join=round, line cap=round, >=stealth]
				\fill[blue!15] (0,0) -- (5,0) -- (2,6) -- (0,7.5) -- cycle;
				
				\draw[->] (-1,0) -- (11,0) node[right] {$x$};
				\draw[->] (0,-1) -- (0,11) node[above] {$y$};
				
				\draw[blue] (-1, 8.25) -- (11, -0.75) node[right] {$3x + 4y = 300$};
				
				\draw[red] (-0.5, 11) -- (5.5, -1) node[right] {$2x + y = 100$};
				
				\fill (0,0) circle (1pt) node[below left] {$O$};
				\fill (5,0) circle (1pt) node[below left] {$50$} node[above] {$A$};
				\fill (2,6) circle (1pt) (2,0) circle (1pt) (0,6) circle (1pt);
				\fill (0,7.5) circle (1pt) node[left] {$75$} node[right] {$C$};
				\fill (10,0) circle (1pt) node[below] {$100$};
				\fill (0,10) circle (1pt) node[left] {$100$};
				\draw[dashed] (2,0)node[below]{$20$}--(2,6) node[above right]{$B$}--(0,6) node[left]{$60$};
			\end{tikzpicture}
		\end{center}
		Giá trị của $P$ tại các đỉnh $P(0;0)=0$, $P(50;0)=400$, $P(20;60)=760$, $P(0;75)=750$.\\
		Vậy số điểm tối đa đạt được là $760$.
	}
\end{ex}

\begin{ex}%[0D2V2-3]%[TT-SGD-TuyenQuang-L1-NH25-26]%[Võ Thị Thùy Trang]%[GVPB: Thanh Phát]
	Để gây quỹ từ thiện, câu lạc bộ thiện nguyện của một trường THPT tổ chức hoạt động bán hàng với hai mặt hàng là trà sữa và bánh ngọt. Câu lạc bộ thiết kế hai thực đơn. Thực đơn $1$ có giá $40$ nghìn đồng, bao gồm hai ly trà sữa và một chiếc bánh ngọt. Thực đơn $2$ có giá $65$ nghìn đồng, bao gồm ba ly trà sữa và hai chiếc bánh ngọt. Biết rằng câu lạc bộ chỉ làm được không quá $180$ ly trà sữa và $110$ chiếc bánh ngọt. Số tiền lớn nhất mà câu lạc bộ có thể nhận được sau khi bán hết hàng bằng bao nhiêu nghìn đồng?
	\par\shortans{3800}
	\loigiai{
		Gọi $x$, $y$ lần lượt là số lượng thực đơn $1$ và thực đơn $2$ mà câu lạc bộ dự định bán ($x \ge 0$, $y \ge 0$).\\
		Theo đề bài, ta có
		\begin{itemize}
			\item Tổng số ly trà sữa cần dùng là $2x + 3y$. Do câu lạc bộ chỉ làm được không quá $180$ ly nên ta có $2x + 3y \le 180$.
			\item Tổng số chiếc bánh ngọt cần dùng là $x + 2y$. Do câu lạc bộ chỉ làm được không quá $110$ chiếc nên ta có $x + 2y \le 110$.
		\end{itemize}
		Từ đó ta có hệ bất phương trình sau $ \heva{&x \ge 0, y \ge 0\\&2x + 3y \le 180\\&x + 2y \le 110.}$\\
		Gọi $F(x,y)$ (nghìn đồng) là số tiền mà câu lạc bộ nhận được sau khi bán hết hàng.\\
		Tìm $x$, $y$ thỏa mãn hệ trên sao cho biểu thức $F(x, y) = 40x + 65y$ (nghìn đồng) đạt giá trị lớn nhất.
		\begin{center}
			\begin{tikzpicture}[>=stealth,line join=round,line cap=round,font=\footnotesize,scale=0.7]
				\draw[->] (-1,0) -- (12,0) node[above]{$x$};
				\draw (0,0) node[below left]{$O$};
				\fill[gray!20] (0,0) -- (9,0) -- (3,4) -- (0,5.5) -- cycle;
				\draw[domain=-0.5:10, samples=2] plot (\x, {6-2/3*\x}) node[right]{$2x+3y=180$};
				\draw[domain=-0.5:11.5, samples=2] plot (\x, {5.5-0.5*\x});
				\node at (9,1.1) [right]{$x+2y=110$};
				\fill (9,0) circle (1.5pt) node[below left] {$A(90;0)$};
				\fill (3,4) circle (1.5pt) node[above right] {$B(30;40)$};
				\fill (0,5.5) circle (1.5pt) node[below left] {$C(0;55)$};
				\draw[dashed] (3,0) node[below]{$30$} -- (3,4) -- (0,4) node[left]{$40$};
				\draw[->] (0,-1) -- (0,7) node[left]{$y$};
			\end{tikzpicture}
		\end{center}
		Miền nghiệm của hệ bất phương trình là miền tứ giác $OABC$ với các đỉnh $O(0;0)$, $A(90;0)$, $B(30;40)$ và $C(0;55)$.
		\begin{itemize}
			\item $F(0; 0) = 0$.
			\item $F(90; 0) = 40 \cdot 90 + 65 \cdot 0 = 3\,600$.
			\item $F(30; 40) = 40 \cdot 30 + 65 \cdot 40 = 3\,800$.
			\item $F(0; 55) = 40 \cdot 0 + 65 \cdot 55 = 3\,575$.
		\end{itemize}
		Ta thấy $F(x, y)$ đạt giá trị lớn nhất bằng $3\,800$ tại $x = 30$ và $y = 40$.\\
		Vậy số tiền lớn nhất mà câu lạc bộ có thể nhận được là $3\,800$ nghìn đồng.
	}
\end{ex}

\begin{ex}%[Đề thi thử môn Toán TRƯỜNG THPT Cổ Loa - Hà Nội năm học 2025-2026]%[Phạm Quốc Toàn - Phạm Hoàng Việt]%[0D2V2-3]
	Một cơ sở đóng thuyền thủ công cần $10$ giờ lao động để đóng một chiếc thuyền loại $A$ và $15$ giờ lao động để đóng một chiếc thuyền loại $B$. Mỗi tuần cơ sở bố trí được tối đa $120$ giờ lao động cho việc đóng hai loại thuyền này. Qua thực tế, người ta thấy mỗi tuần cơ sở bán được tối đa $6$ chiếc thuyền loại $A$ và tối thiểu $2$ chiếc thuyền loại $B$. Mỗi chiếc thuyền loại $A$ cho thuận lợi là $0{,}5$ triệu đồng và mỗi chiếc thuyền loại $B$ cho lợi nhuận là $0{,}7$ triệu đồng. Qua thời gian tìm hiểu, chủ cơ sở đã tìm được số lượng nên đóng thuyền mỗi loại để có thể thu được lợi nhuận cao nhất là $T$ (triệu đồng). Giá trị $T$ bằng bao nhiêu? 
	\par\shortans[0]{5{,}8}
	\loigiai{
		Gọi $x$, $y$ lần lượt là số thuyền loại $A$ và số thuyền loại $B$ được cơ sở đóng trong một tuần $\left(x, y\in \mathbb{N}\right)$. Theo đề bài ta có hệ bất phương trình sau 
		\[
		\heva{&10x+15y \le 120 \\&0 \le x \le 6 \\& y \ge 2} \Leftrightarrow \heva{&2x+3y\le 24 \\&0\le x\le 6 \\&y\ge 2.}.
		\]
		Hàm lợi nhuận là $T=0{,}5x+0{,}7y$.\\
		Ta có miền nghiệm của hệ bất phương trình trên là tứ giác $ABCD$ như hình vẽ 
		\begin{center}
			\begin{tikzpicture}[scale=0.8, font=\footnotesize, line join=round, line cap=round, >=stealth]
				\tikzset{every node/.style={scale=0.9}}
				\begin{scope}
					\clip (-1,-1) rectangle (13,9);
					\fill[black!35] (-2,9.33)--(14,9.33)--(14,-1.33)--cycle;
					\fill[black!35] (0,-1)--(-1,-1)--(-1,9)--(0,9)--cycle;
					\fill[black!35] (6,-1)--(13,-1)--(13,9)--(6,9)--cycle;
					\fill[black!35] (-1,2)--(-1,-1)--(13,-1)--(13,2)--cycle;
					\draw (-1.5,9)--(13.5,-1) node [pos=0.6, above, sloped] {$2x+3y-24=0$};
					\draw (6,-1)--(6,9) (-1,2)--(13,2) ;
				\end{scope}
				\draw[->] (-1,0)--(13,0) node[below]{$x$};
				\draw[->] (0,-1)--(0,9) node[left]{$y$};
				\draw (0,0) node[below left]{$O$};
				\foreach \x in {6,12}
				\draw[thin] (\x,1pt)--(\x,-1pt) node [below left] {$\x$};
				\foreach \y in {2,4,8}
				\draw[thin] (1pt,\y)--(-1pt,\y) node [below left] {$\y$};
				\draw(0,2)circle(1pt)node[above right]{$A$};
				\draw(6,2)circle(1pt)node[above left]{$B$};
				\draw(6,4)circle(1pt)node[below left]{$C$};
				\draw(0,8)circle(1pt)node[right]{$D$};
				\draw[dashed](0,4)--(6,4);
			\end{tikzpicture}
		\end{center}
		\begin{itemize}
			\item Tại $A(0;2)$, ta có $T=0{,}5 \cdot 0 +0{,}7 \cdot 2=1{,}4$.
			\item Tại $B(6;2)$, ta có $T=0{,}5\cdot 6+0{,}7 \cdot 2=4{,}4$. 
			\item Tại $C(6;4)$, ta có $T=0{,}5 \cdot 6+0{,}7 \cdot 4=5{,}8$.
			\item Tại $D(0;8)$, ta có $T=0{,}5 \cdot 0 +0{,}7 \cdot 8=5{,}6$.
		\end{itemize}
		Vậy lợi nhuận nhiều nhất là $5{,}8$ triệu đồng khi đóng $6$ thuyền loại $A$ và $4$ thuyền loại $B$ mỗi tuần.
	}
\end{ex}

\begin{ex}%[TT-THPT-HamRong-NH25-26 ]%[GV soạn: Nguyễn Quang Hiệp GVPB2 Đỗ Minh Vũ ]%[0D2V2-3]
	Một cửa hàng gạo tháng này cần nhập hai loại gạo là gạo loại I và gạo loại II. Số tiền chi cho việc nhập hai loại gạo này không quá $100$ triệu đồng. Gạo loại I mua vào với giá $30$ triệu đồng/tấn, lợi nhuận dự kiến là $2$ triệu đồng/tấn. Gạo loại II nhập vào với giá $40$ triệu đồng/tấn và lợi nhuận dự kiến là $2{,}5$ triệu đồng/tấn. Cửa hàng ước tính doanh số bán ra một tháng của hai loại gạo này không quá $3$ tấn. Hỏi cửa hàng có thể lãi nhiều nhất là bao nhiêu (triệu đồng) khi kinh doanh hai loại gạo này?
	\par\shortans[oly]{6,5}
	\loigiai{
		Gọi $x$, $y$ lần lượt là số tấn gạo loại I và loại II cần nhập trong tháng ($x \ge 0$, $y \ge 0$).\\
		Tiền lãi cửa hàng nhận được là $T = 2x + 2{,}5y$ (triệu đồng).\\
		Số tiền cửa hàng cần chi là $30x + 40y$ (triệu đồng).\\
		Tổng doanh số bán ra không quá $3$ tấn trong tháng nên ta có $x + y \le 3$.\\
		Do cửa hàng dự định chi không quá $100$ triệu đồng nên $30x + 40y \le 100 \Leftrightarrow 3x + 4y \le 10$.\\
		Bài toán đưa về tìm $x$, $y$ là nghiệm của hệ bất phương trình
		$ \heva{& x \ge 0 \\ & y \ge 0 \\ & x + y \le 3 \\ & 3x + 4y \le 10} $
		sao cho $T = 2x + 2{,}5y$ có giá trị lớn nhất.\\
		Miền nghiệm của hệ bất phương trình là miền tứ giác $OABC$ với $O(0;0)$, $A\left(0;\dfrac{5}{2}\right)$, $B(2;1)$, $C(3;0)$.
		\begin{center}
			\begin{tikzpicture}[line cap=round, line join=round, font=\footnotesize, >=stealth, scale=1.5]
				% Yêu cầu khai báo thư viện trong preamble: \usetikzlibrary{patterns}
				
				% Giới hạn của hệ trục (Bounding Box)
				\def\xmin{-1.0}
				\def\xmax{4.5}
				\def\ymin{-1.0}
				\def\ymax{4.0}
				
				% Vẽ lưới tọa độ (mỗi ô vuông 0.5 đơn vị)
				\draw[step=0.5cm, help lines, gray!40, thin] (\xmin,\ymin) grid (\xmax,\ymax);
				
				\begin{scope}
					\clip (\xmin,\ymin) rectangle (\xmax,\ymax);
					
					% === TÔ MIỀN KHÔNG NGHIỆM ===
					
					% 1. Miền KHÔNG thỏa: x >= 0 (Gạch bỏ x < 0)
					\fill[pattern=north west lines, pattern color=green!60!black]
					(\xmin,\ymin) rectangle (0,\ymax);
					\fill[opacity=0.05, green!60!black]
					(\xmin,\ymin) rectangle (0,\ymax);
					
					% 2. Miền KHÔNG thỏa: y >= 0 (Gạch bỏ y < 0)
					\fill[pattern=north east lines, pattern color=brown!80!red]
					(\xmin,\ymin) rectangle (\xmax,0);
					\fill[opacity=0.05, brown!80!red]
					(\xmin,\ymin) rectangle (\xmax,0);
					
					% 3. Miền KHÔNG thỏa: x + 2y <= 5 (Gạch bỏ x + 2y > 5 => y > 2.5 - 0.5x)
					\fill[pattern=north west lines, pattern color=blue!80]
					(\xmin, {2.5 - 0.5*\xmin}) -- (\xmax, {2.5 - 0.5*\xmax}) -- (\xmax, \ymax) -- (\xmin, \ymax) -- cycle;
					\fill[opacity=0.05, blue!80]
					(\xmin, {2.5 - 0.5*\xmin}) -- (\xmax, {2.5 - 0.5*\xmax}) -- (\xmax, \ymax) -- (\xmin, \ymax) -- cycle;
					
					% 4. Miền KHÔNG thỏa: x + y <= 3 (Gạch bỏ x + y > 3 => y > 3 - x)
					\fill[pattern=north east lines, pattern color=red!80]
					(\xmin, {3 - \xmin}) -- (\xmax, {3 - \xmax}) -- (\xmax, \ymax) -- (\xmin, \ymax) -- cycle;
					\fill[opacity=0.05, red!80]
					(\xmin, {3 - \xmin}) -- (\xmax, {3 - \xmax}) -- (\xmax, \ymax) -- (\xmin, \ymax) -- cycle;
					
					% === VẼ CÁC ĐƯỜNG THẲNG ===
					\draw[blue, thick] (\xmin, {2.5 - 0.5*\xmin}) -- (\xmax, {2.5 - 0.5*\xmax});
					\draw[red, thick] (\xmin, {3 - \xmin}) -- (\xmax, {3 - \xmax});
				\end{scope}
				
				% === TRỤC TỌA ĐỘ ===
				\draw[->, thick] (\xmin,0)--(\xmax,0) node[above left] {$x$};
				\draw[->, thick] (0,\ymin)--(0,\ymax) node[below right] {$y$};
				\node[below left] at (0,0) {$O$};
				
				% === CÁC ĐIỂM ĐÁNH DẤU TỌA ĐỘ ===
				% Trên trục Ox
				\draw[thick] (2, 0.05) -- (2, -0.05) node[below] {$2$};
				\draw[thick] (3, 0.05) -- (3, -0.05) node[below] {$3$};
				
				% Trên trục Oy
				\draw[thick] (0.05, 1) -- (-0.05, 1) node[left] {$1$};
				\draw[thick] (0.05, 2.5) -- (-0.05, 2.5) node[left] {$\dfrac{5}{2}$};
				
				% === VẼ ĐIỂM GIAO & DÁN NHÃN ===
				\filldraw[fill=yellow, draw=orange] (0, 2.5) circle (1.2pt) node[blue, above right] {$A$};
				\filldraw[fill=yellow, draw=orange] (1, 2) circle (1.2pt) node[blue, below left] {$B$};
				\filldraw[fill=yellow, draw=orange] (3, 0) circle (1.2pt) node[blue, above right] {$C$};
				
				% Khung bao viền (Bounding Box giống file xuất GeoGebra)
				\draw[thick, gray!80] (\xmin,\ymin) rectangle (\xmax,\ymax);
				
			\end{tikzpicture}
		\end{center}
		Giá trị biểu thức $T = 2x + 2{,}5y$ đạt giá trị lớn nhất tại một trong các đỉnh của tứ giác $OABC$.\\
		Ta có $T(O) = 0$, $T(A) = 2 \cdot 0 + 2{,}5 \cdot \dfrac{5}{2} = 6{,}25$, $T(B) = 2 \cdot 2 + 2{,}5 \cdot 1 = 6{,}5$, $T(C) = 2 \cdot 3 + 2{,}5 \cdot 0 = 6$.\\
		Suy ra $T$ đạt giá trị lớn nhất bằng $6{,}5$ khi $x = 2$, $y = 1$ ứng với tọa độ đỉnh $B$.\\
		Vậy cửa hàng lãi nhiều nhất là $6{,}5$ triệu đồng trong tháng khi kinh doanh hai loại gạo này.
	}
\end{ex}

\begin{ex}%[0D2V2-3]%[Dự án C đề thi thử TN Môn Toán NH25-26 - Dot 4 - Phạm Đăng Đăng]%[THPT Lê Quý Đôn-Đống Đa-Hà Nội]%[GVPB - Phan Quốc Trí]
	Một cơ sở rang xay cà phê để sản xuất và đóng gói cà phê  thành phẩm. Mỗi tháng cơ sở nhập về $300$ kg cà phê vối và $240$ kg cà phê chè để làm nguyên liệu. Hiện tại cơ sở sản xuất hai loại cà phê là loại $I$ và loại $II$.
	\begin{itemize}
		\item Loại $I$: Đóng gói $100$ g với tỉ lệ trộn $40$ g cà phê vối và $60$ g cà phê chè.
		\item Loại $II$: Đóng gói $100$ g với tỉ lệ $60$ g cà phê vối và $40$ g cà phê chè.
	\end{itemize}
	Biết rằng mỗi gói cà phê loại $I$ cơ sở lãi $10\,000$ đồng, mỗi gói cà phê loại $II$ cơ sở lãi $8\,000$ đồng. Nếu không phải lo khâu tiêu thụ thì cơ sở sản xuất thu được lợi nhuận mỗi tháng tối đa là bao nhiêu triệu đồng ( làm tròn đến hàng phần mười)?
	\shortans{$45{,}6$}
	\loigiai{Ta có $300$ kg = $300\,000$ g cà phê vối; $240$ kg = $240\,000$ g cà phê chè.\\
		Gọi $x$ là số gói loại $I$, $y$ là số gói loại $II$ ($x, y \geq 0$).\\
		Tổng cà phê vối: $40x + 60y \leq 300\,000 \Leftrightarrow 2x + 3y \leq 15\,000$.\\
		Tổng cà phê chè: $60x + 40y \leq 240\,000 \Leftrightarrow 3x + 2y \leq 12\,000$.\\
		Tổng tiền lãi của cơ sở $P = 10\,000x + 8\,000y$ (đơn vị: đồng).\\
		Bài toán trở thành tìm nghiệm của hệ bất phương trình 
		$\heva{ 
			&x \geq 0 \\ 
			&y \geq 0 \\ 
			&2x + 3y \leq 15\,000 \\ 
			&3x + 2y \leq 12\,000 
		}$ để biểu thức
		$P = 10\,000x + 8\,000y$ đạt giá trị lớn nhất.
		\begin{center}
			\begin{tikzpicture}[scale=0.7,font=\footnotesize,line join=round
				,line cap=round,>=stealth]
				\coordinate (A) at (0,5);
				\coordinate (B) at (1.2,4.2);
				\coordinate (C) at (4,0);
				\coordinate (O) at (0,0);
				\draw[fill=black] (A) circle (1pt) node[below,xshift=1.5mm,yshift=-4mm] {$A$};
				\draw[fill=black] (B) circle (1pt) node[below] {$B$};
				\draw[fill=black] (C) circle (1pt) node[above,xshift=-4mm]  {$C$};
				\draw[fill=black] (O) circle (1pt) node[ above right]  {$O$};
				\fill[pattern=crosshatch dots
				] 
				(-3,0)--(-3,-2)--(10,-2)--(10,0)--cycle
				(-3,0)--(-3,10)--(0,10)--(0,0)--cycle
				(0,10)--(A)--(B)--(C)--(10,0)--(10,10)--cycle;
					\draw[->] (-3,0)--(10,0) node[above right]{$x$};
				\draw[->] (0,-2)--(0,10) node[above right]{$y$};
				
				\draw[thick,blue,domain=9:-3,samples=2]
				plot (\x,{5-2/3*\x})
				node[below right,rotate=-33] {$2x+3y=15\,000$};

				\draw[thick,red,domain=5:-2.5,samples=2]
				plot (\x,{6-1.5*\x})
				node[above right,rotate=-55,] {$3x+2y=12\,000$};
				
			\end{tikzpicture}
		\end{center}
		Miền nghiệm của hệ bất phương trình là miền tứ giác $OABC$ với 
		$O(0; 0)$, $A(0; 5\,000)$, $C(4\,000; 0)$, $B(1\,200; 4\,200)$ (Miền không bị tô đậm).\\
		Tại $(0; 0): P = 0$\\
		Tại $(1\,200; 4\,200): P = 10\,000(1\,200) + 8\,000(4\,200)  = 45\,600\,000$ đồng.\\
		Tại $(0;5\,000): P = 8\,000(5\,000) = 40\,000\,000$ đồng.\\
		Tại $(4\,000; 0): P = 10\,000(4\,000) = 40\,000\,000$ đồng.\\
		Vậy lợi nhuận tối đa là $45{,}6$ triệu đồng/tháng.}
\end{ex}

\begin{ex}%[0D2V2-3]%[TN-THPT-NguyenTrungThien-HaTinh-NH25-26]%[Loc Do]%[PB: Phan Anh]
	Một cơ quan X chuẩn bị tổ chức cho $500$ người đi tham quan trải nghiệm. Để chuẩn bị cho chuyến đi, cơ quan cần vận chuyển tổng cộng $29$ tấn hàng hoá (bao gồm vật dụng, thực phẩm,...). Công ty vận tải báo giá cho thuê xe như sau: 
	\begin{itemize}
		\item Xe lớn: Có thể chở tối đa $50$ người và $2$ tấn hàng. Chi phí thuê là $10$ triệu đồng/xe. Công ty có $13$ xe loại này.
		\item Xe nhỏ: Có thể chở tối đa $30$ người và $3$ tấn hàng. Chi phí thuê là $7$ triệu đồng/xe. Công ty có $15$ xe loại này.
	\end{itemize}
	Sau khi tính toán, cơ quan X chọn phương án để chi phí thuê xe là thấp nhất. Số tiền thuê xe thấp nhất mà cơ quan X phải trả là bao nhiêu triệu đồng?
	\shortans{137}
	\loigiai{
		Giả sử, cơ quan X thuê số xe lớn và xe nhỏ lần lượt là $x$, $y$ (xe), ($x,y\in \mathbb{N}$).\\
		Khi đó, số tiền mà cơ quan phải trả là $F(x,y)=10x+7y$ (triệu đồng).\\
		Bất phương trình biểu thị số người là $50x+30y\ge 500\Leftrightarrow 5x+3y\ge 50$.\\
		Bất phương trình biểu thị khối lượng hàng hoá là  $2x+3y\ge 29$.\\
		Do đó, từ đề bài ta có hệ bất phương trình sau $\heva{& 0\le x\le 13\\ & 0\le y\le 15\\ & 5x+3y\ge 50\\ & 2x+3y\ge 29.}$ (*)\\
		Bài toán trở thành tìm giá trị nhỏ nhất của biểu thức $F(x,y)=10x+7y$ với $x,y$ thoả mãn (*).\\
		Biểu diễn miền nghiệm của hệ (*) bằng phần không bị tô đậm (kể cả bờ) được minh hoạ ở hình vẽ dưới đây:
		\begin{center}
			\begin{tikzpicture}[line join=round, line cap=round,>=stealth,thick,x=0.2cm,y=0.2cm]
				\tikzset{every node/.style={scale=0.9}}
				\begin{scope}
					\clip (-5,-5) rectangle (20,20);
					\fill[pattern=north east lines, pattern color=blue] (-5,-5) -- (0,-5) -- (0,20) -- (-5,20) -- cycle;
					\draw (0,-5)--(0,20)  ;
					\fill[pattern=north west lines, pattern color=green] (13,-5) -- (20,-5) -- (20,20) -- (13,20) -- cycle;
					\draw (13,-5)--(13,20)  ;
					\fill[pattern=vertical lines, pattern color=cyan] (-5,0) -- (-5,-5) -- (20,-5) -- (20,0) -- cycle;
					\draw (-5,0)--(20,0)  ;
					\fill[pattern=grid, pattern color=orange] (-5,15) -- (20,15) -- (20,20) -- (-5,20) -- cycle;
					\draw (-5,15)--(20,15)  ;
					\fill[pattern=dots, pattern color=magenta] (-5,-5) -- (13,-5) -- (-2,20) -- (-5,20) -- cycle;
					\draw (13,-5)--(-2,20)  ;
					\fill[pattern=crosshatch, pattern color=brown] (-5,13) -- (-5,-5) -- (20,-5) -- (20,-3.67) -- cycle;
					\draw (-5,13)--(20,-3.67)  ;
				\end{scope}
				\draw[->] (-5,0)--(22.2,0) node[below]{$x$};
				\draw[->] (0,-5)--(0,22.2) node[left]{$y$};
				\draw (0,0) node[below left]{$O$};
				\fill (1,15) circle (1.5pt) node[below right,scale=0.8] {$(1,15)$};
				\fill (13,15) circle (1.5pt) node[below left,scale=0.8] {$(13,15)$};
				\fill (13,1) circle (1.5pt) node[ right,scale=0.8] {$(13,1)$};
				\fill (7,5) circle (1.5pt) node[right,scale=0.8] {$(7,5)$};
				\foreach \x/\g in {5/-90,10/-90,15/-90}{
					\fill (\x,0) circle (1pt) +(\g:3mm) node[scale=0.8] {$\x$};
				}
				
				\foreach \y/\g in {5/180,10/180,15/160}{
					\fill (0,\y) circle (1pt) +(\g:3mm) node[scale=0.8] {$\y$};
				}
			\end{tikzpicture}
		\end{center}
		Trong đó, toạ độ các đỉnh được xác định như sau
		\begin{itemize}
			\item $A(1;15)$ là giao điểm của $2$ đường thẳng $y=15$ và $5x+3y=50$;
			\item $B(13;15)$ là giao điểm của $2$ đường thẳng $x=13$ và $y=15$;
			\item $C(7;5)$ là giao điểm của $2$ đường thẳng $5x+3y=50$ và $2x+3y=29$;
			\item $D(13; 1)$ là giao điểm của $2$ đường thẳng $x=13$ và $2x+3y=29$.
		\end{itemize}
		Kiểm tra giá trị của $F$ tại các đỉnh của tứ giác ta được
		\[F(1;15)=115, \,F(13;15)=235,\, F(7;5)=105,\, F(13;1)=137.\]
		Vậy số tiền ít nhất mà cơ quan X phải trả là $105$ triệu đồng (khi thuê $7$ xe lớn và $5$ xe nhỏ).
	}
\end{ex}

\begin{ex}%[0D2V2-3]%[Dự án C đợt 3 NH25-26- Phạm Đăng Đăng]%[TT-THPT-TranNhanTong-HaNoi-N25-26]
	Một xưởng mộc sản xuất hai loại sản phẩm: bàn gỗ và ghế gỗ. Xưởng có tối đa $240$ đơn vị gỗ và tổng thời gian nhân công là $100$ giờ. Để làm $1$ cái bàn cần $5$ đơn vị gỗ và $2$ giờ công. Để làm $1$ cái ghế cần $2$ đơn vị gỗ và $1$ giờ công. Mỗi cái bàn lãi $600\,000$ đồng, mỗi cái ghế lãi $300\,000$ đồng. Lợi nhuận của xưởng mộc lớn nhất là bao nhiêu triệu đồng?
	\shortans{$30$}
	\loigiai{
	\immini[thm]{	Gọi $x$, $y$ lần lượt là số bàn, ghế mà xưởng mộc sản xuất $(x,y>0)$.\\
		Theo đề bài, ta có hệ bất phương trình $$\heva{&5x+2y\le240\\
			&2x+y\le 100\\
			&x,y\ge0.}$$
	Miền nghiệm của hệ là miền tứ giác $OABC$ với $O(0;0)$, $A(0;100)$, $B(40;20)$ và $C(48;0)$.\\
			Lợi nhuận của xưởng thu được là \break $L(x;y)=600x+300y$ (đơn vị: nghìn đồng).\\
			Ta có $L(0;0)=0$, $L(0;100)=30\,000$,\\ $L(40;20)=30\,000$, $L(48;0)=28\,800$.\\Vậy lợi nhuận lớn nhất thu được là 30 triệu đồng.}{	\begin{tikzpicture}[scale=0.05, line cap=round, line join=round, font=\small,>=stealth]
				\clip (-50,-20) rectangle(100,200);
				% Trục tọa độ
				\draw[->] (-40,0)--(80,0) node[above right]{$x$};
				\draw[->] (0,-20)--(0,180) node[above right]{$y$};
				
				
				\draw[thick,blue,domain=110:-30,samples=2]
				plot (\x,{100-2*\x})
				node[below right,rotate=-65] {$2x+y=100$};
				
				% x + 4y = 200 <=> y = 50 - x/4
				\draw[thick,red,domain=220:-20,samples=2]
				plot (\x,{120-2.5*\x})
				node[above right,rotate=-65,] {$5x+2y=240$};
				
				
				
				% Các điểm A, B, C
				\coordinate (A) at (0,100);
				\coordinate (B) at (40,20);
				\coordinate (C) at (48,0);
				
				\draw[fill=black] (A) circle (20pt) node[below,xshift=1.5mm,yshift=-4mm] {$A$};
				\draw[fill=black] (B) circle (20pt) node[left] {$B$};
				\draw[fill=black] (C) circle (20pt) node[above,xshift=-4mm]  {$C$};
				\node[below left] at (0,0) {$O$};
				
				\node[below left] at (0,100) {$100$};
				\fill[pattern=horizontal lines] 
				(80,0)--(C)--(B)--(A)--(0,180)--(80,180)
				(80,0)--(80,-20)--(-40,-20)--(-40,0)--cycle
				(0,0)--(-40,0)--(-40,180)--(0,180)--cycle;
				%\fill[pattern=horizontal lines] plot[domain=-20:200/3]
				%(\x,{100-\x}) --
				%(220,-5) --(220,120)--cycle;
				%\fill[pattern=horizontal lines] plot[domain=90:20]
				%	(\x,{155-2*\x}) --
				%	(-20,120) --(-20,-20)--cycle;
		\end{tikzpicture}}
	}
\end{ex}

\begin{ex}%[Dự án đề TT Toán NH25-26-Dot 4- Hải Phụng]%[THCS-THPT LÊ THÁNH TÔNG - TP.HCM]%[GVPB - Viet Hoang Pham]%[0D3H2-6]
	Anh Trọng quyết định mua $2$ cây đào và $1$ cây mai tại một nhà vườn để trang trí nhà dịp Tết Nguyên Đán. Giá mỗi cây đào là $1\,000\,000$ đồng, giá mỗi cây mai là $2\,000\,000$ đồng. Để kích cầu, nhà vườn đưa ra chính sách: nếu khách hàng mua thêm $x$  cây cảnh nhỏ (với $x$  là số nguyên không âm), mỗi cây nhỏ có giá $50\,000$ đồng thì chi phí vận chuyển toàn bộ đơn hàng được tính theo hàm số $f(x)=x^2-100x+3\,000$  (đơn vị: nghìn đồng). Biết rằng số lượng cây mua thêm không vượt quá 40 cây để đảm bảo tải trọng xe $(0 \leq x \leq 40)$. Tổng chi phí (tiền mua cây và tiền vận chuyển) thấp nhất mà Anh Trọng phải bỏ ra để trang trí nhà dịp Tết là bao nhiêu nghìn đồng?
	\shortans{$6375$}
	\loigiai{
		Anh Trọng mua cố định:
		\begin{itemize}
			\item $2$ cây đào: \( 2 \cdot 1\,000\,000 = 2\,000\,000 \) đồng
			\item $1$ cây mai: \( 1 \cdot 2\,000\,000 = 2\,000\,000 \) đồng
		\end{itemize}
		Vậy tiền mua cây cố định là \( 4\,000\,000 \) đồng, tức là \( 4\,000 \) nghìn đồng.\\
		Tổng chi phí \( C(x) \) (đơn vị nghìn đồng) là
		\[
		C(x) = 4\,000 + 50x + (x^2 - 100x + 3\,000) = x^2 - 50x + 7\,000.
		\]
		Hàm \( C(x) = x^2 - 50x + 7\,000 \) là hàm bậc hai. Giá trị nhỏ nhất đạt tại
		\[
		x = -\frac{b}{2a} = \frac{50}{2 \cdot 1} = 25.
		\]
		Chi phí tại \( x = 25 \) là
		\[
		C(25) = 25^2 - 50 \cdot 25 + 7\,000 = 625 - 1250 + 7\,000 = 6\,375.
		\]
		Vậy chi phí nhỏ nhất là \( 6\,375 \) nghìn đồng khi mua thêm \( x = 25 \) cây cảnh nhỏ.
	}
\end{ex}

\begin{ex}%[0D8C1-3]%[Dự án đề thi thử đợt 3-Nguyễn Đức Lợi]%[SGD tỉnh Lạng Sơn]%[GVPB - ThanhTa]
	\immini[thm]{Một màn chơi của một trò chơi điện tử được thiết kế như sau: có năm vị trí $A$, $B$, $C$, $D$, $E$, được đặt ở năm đỉnh của một hình chóp tứ giác. Nhân vật sẽ được đặt ở một vị trí bất kỳ và có thể di chuyển tự do giữa các đỉnh với nhau, mỗi lần di chuyển đều phải từ đỉnh này di chuyển đến đỉnh khác. Giả sử khi bắt đầu, nhân vật được đặt ở vị trí $A$. Số cách di chuyển để sau sáu bước nhảy, nhân vật quay lại vị trí $A$ là bao nhiêu?
	}
	{
		\begin{tikzpicture}[font=\footnotesize, line join=round, line cap=round,>=stealth,scale=.8]
			
			\path (0,0) coordinate (B)
			(1,-1.5) coordinate (C)
			(3,-1.5) coordinate (D)
			(5,0) coordinate (E)
			(2,3) coordinate (A);
			\draw (A)--(B)--(C)--cycle (A)--(D)--(E)--cycle (C)--(D);
			\draw [dashed] (D)--(B)--(E)--(C);
			\foreach \t/\g in {A/90,B/180,C/-135,D/-45,E/0}{
				\fill (\t) circle (1pt) node[shift={(\g:10pt)}]{$ \t $};
			}
			\foreach \t/\g in {A,B,C,D,E}{
				\draw (\t) circle (5pt);}
		\end{tikzpicture}
	}
	\par
	\shortans{820}
	\loigiai{
		Từ một đỉnh bất kỳ, ta luôn có $4$ cách di chuyển đến $4$ đỉnh còn lại (các đỉnh $B$, $C$, $D$, $E$ có vai trò như nhau). Quá trình di chuyển gồm $6$ bước, bắt đầu từ $A$ và kết thúc tại $A$.\\
		Ta chia thành các trường hợp dựa vào số lần nhân vật quay trở lại $A$ ở các bước giữa (từ bước $1$ đến bước $5$).
		\begin{enumerate}[\it TH 1:]
			\item Trong $5$ bước ở giữa, không có lần nào quay lại $A$.\\
			Đường đi có dạng $A \to X_1 \to X_2 \to X_3 \to X_4 \to X_5 \to A$ (với $X_i \ne A$).
			\begin{itemize}
				\item Từ $A \to X_1$ có $4$ cách.
				\item Từ $X_1 \to X_2$:có $3$ cách (do $X_2 \ne A$ và $X_2 \ne X_1$).
				\item Tương tự, các bước $X_2 \to X_3$, $X_3 \to X_4$, $X_4 \to X_5$ mỗi bước đều có $3$ cách.
				\item Từ $X_5 \to A$ có $1$ cách.
			\end{itemize}
			Số cách di chuyển ở TH$_1$ là $4 \cdot 3 \cdot 3 \cdot 3 \cdot 3 \cdot 1=324$ (cách).
			\item Trong $5$ bước ở giữa, có đúng $1$ lần quay lại $A$.\\
			Do không thể đứng yên tại $A$ trong $2$ bước liên tiếp, lần quay lại $A$ không thể là bước $1$ và bước $5$, nên chỉ có thể rơi vào bước $2$, $3$ hoặc $4$.
			\begin{itemize}
				\item Về $A$ ở bước $2$ ($A \to X_1 \to A \to X_3 \to X_4 \to X_5 \to A$).\\
				Số cách là $4 \cdot 1 \cdot 4 \cdot 3 \cdot 3 \cdot 1=144$ (cách).
				\item Về $A$ ở bước $3$ ($A \to X_1 \to X_2 \to A \to X_4 \to X_5 \to A$).\\
				Số cách là $4 \cdot 3 \cdot 1 \cdot 4 \cdot 3 \cdot 1=144$ (cách).
				\item Về $A$ ở bước $4$ ($A \to X_1 \to X_2 \to X_3 \to A \to X_5 \to A$).\\
				Số cách là $4 \cdot 3 \cdot 3 \cdot 1 \cdot 4 \cdot 1=144$ (cách).
			\end{itemize}
			Số cách di chuyển ở TH$_2$ là $144 \cdot 3=432$ (cách).
			\item Trong $5$ bước ở giữa, có đúng $2$ lần quay lại $A$.\\
			Để $A$ không lặp lại kề nhau, $2$ lần quay lại $A$ bắt buộc phải ở bước $2$ và bước $4$.\\
			Đường đi có dạng $A \to X_1 \to A \to X_3 \to A \to X_5 \to A$.
			\begin{itemize}
				\item Bước $A \to X_1 \to A$ có $4 \cdot 1=4$ cách.
				\item Bước $A \to X_3 \to A$ có $4 \cdot 1=4$ cách.
				\item Bước $A \to X_5 \to A$ có $4 \cdot 1=4$ cách.
			\end{itemize}
			Số cách di chuyển ở TH$_3$ là $4 \cdot 4 \cdot 4=64$ (cách).
		\end{enumerate}
		(Không thể có $3$ lần quay lại $A$ vì sẽ dẫn đến $2$ bước liên tiếp đứng tại $A$).\\
		Vậy tổng số cách di chuyển thỏa mãn là $324 + 432 + 64 = 820$ (cách).
	}
\end{ex}

\begin{ex}%[0D8C2-5]%[Dự án C - Đợt 3 - Sở Giáo Dục Sơn La]%[Nguyễn Chiến - Nguyễn Tiến]
	\immini[thm]{
		Cho hình vẽ bên gồm $4$ tam giác. Người ta chọn $3$ số phân biệt từ tập hợp $S = \{1; 2; \ldots; 26\}$ để xếp vào $3$ tam giác ở $3$ góc. Sau đó, tính tổng bình phương của $3$ số đó rồi ghi kết quả vào tam giác còn lại ở giữa. Hỏi có bao nhiêu cách xếp sao cho số ghi ở tam giác giữa là một số chia hết cho $5$?
	}{
		\begin{tikzpicture}[scale=0.8, line join=round, line cap=round]
			\coordinate (A) at (0,0);
			\coordinate (B) at (3,0);
			\coordinate (C) at (1.5,{1.5*sqrt(3)});
			\coordinate (M) at ($(A)!0.5!(B)$);
			\coordinate (N) at ($(B)!0.5!(C)$);
			\coordinate (P) at ($(C)!0.5!(A)$);
			\draw[blue] (A) -- (B) -- (C) -- cycle;
			\draw[blue] (M) -- (N) -- (P) -- cycle;
		\end{tikzpicture}
	}
	\shortans[oly]{$3360$}
	\loigiai{
		Gọi $a$, $b$, $c$ là $3$ số phân biệt được chọn từ tập hợp $S$ để xếp vào $3$ góc.\\
		Theo bài ra, số ghi ở tam giác giữa là $a^2 + b^2 + c^2$.\\
		Ta cần tìm số cách chọn và xếp có thứ tự $3$ số $a$, $b$, $c$ sao cho tổng $a^2 + b^2 + c^2$ chia hết cho $5$.\\
		Với mọi số nguyên $x$, bình phương của nó chia cho $5$ chỉ có thể dư $0$, $1$ hoặc $4$.\\
		Gọi $R_a$, $R_b$, $R_c \in \{0; 1; 4\}$ lần lượt là số dư của $a^2$, $b^2$, $c^2$ khi chia cho $5$.\\
		Để $a^2 + b^2 + c^2$ chia hết cho $5$ thì tổng $R_a + R_b + R_c$ phải chia hết cho $5$.\\
		Ta chia tập $S = \{1; 2; \ldots; 26\}$ thành $3$ tập con theo số dư của bình phương khi chia cho $5$:
		\begin{itemize} 
			\item Tập $A_0$ gồm các số $x \in S$ thỏa mãn $x$ chia hết cho $5$.\\
			Suy ra $A_0 = \{5; 10; 15; 20; 25\}$ gồm $5$ phần tử.
			\item Tập $A_1$ gồm các số $x \in S$ thỏa mãn $x$ chia cho $5$ dư $1$ hoặc $x$ chia cho $5$ dư $4$.\\
			Suy ra $A_1 = \{1; 4; 6; 9; 11; 14; 16; 19; 21; 24; 26\}$ gồm $11$ phần tử.
			\item Tập $A_2$ gồm các số $x \in S$ thỏa mãn $x$ chia cho $5$ dư $2$ hoặc $x$ chia cho $5$ dư $3$.\\
			Suy ra $A_2 = \{2; 3; 7; 8; 12; 13; 17; 18; 22; 23\}$ gồm $10$ phần tử.
		\end{itemize}
		Trường hợp 1: Cả $3$ số $a$, $b$, $c$ đều có bình phương chia hết cho $5$.\\
		Ta chọn $3$ số phân biệt từ tập $A_0$ và xếp vào $3$ vị trí, số cách xếp là $\mathrm{A}_5^3 = 5 \cdot 4 \cdot 3 = 60$ (cách).\\
		Trường hợp 2: $3$ số được chọn có bình phương chia $5$ dư lần lượt là $0$, $1$ và $4$.\\
		Ta chọn $1$ số từ $A_0$, $1$ số từ $A_1$ và $1$ số từ $A_2$, sau đó hoán vị vào $3$ vị trí góc.\\
		Số cách thực hiện là $\mathrm{C}_5^1 \cdot \mathrm{C}_{11}^1 \cdot \mathrm{C}_{10}^1 \cdot 3! = 5 \cdot 11 \cdot 10 \cdot 6 = 3300$ (cách).\\
		Tổng số cách xếp thỏa mãn yêu cầu bài toán là $60 + 3300 = 3360$ (cách).
	}
\end{ex}

\begin{ex}%[0D8C2-7]%[Dự án C đề thi thử TN Môn Toán NH25-26 - Dot 4 - Phạm Hoàng Đăng]%[THPT Nguyễn Trãi-Hà Nội]%[GVPB - Khắc Thiên]
	Một kỹ sư tiến hành lắp ráp một rotor của động cơ phản lực. Rotor có $9$ khe cắm cánh quạt được đánh số cố định từ $1$ đến $9$ theo vòng tròn như hình vẽ. Khoảng cách giữa các khe đều nhau, tạo thành các đỉnh của một đa giác đều có $9$ cạnh. Do sai số chế tạo, $9$ cánh quạt có khối lượng thực tế là các số nguyên phân biệt từ $1$ gam đến $9$ gam. Để đảm bảo rotor cân bằng động học khi quay, kỹ sư lựa chọn phương án lắp đặt thỏa mãn đồng thời các điều kiện sau:
	\begin{itemize}
		\item Chia $9$ cánh quạt thành $3$ nhóm (mỗi nhóm $3$ cánh).
		\item Mỗi nhóm được lắp vào $3$ khe cắm tạo thành một tam giác đều (ba khe cắm tạo thành một tam giác đều khi và chỉ khi chúng cách nhau đúng $3$ khe theo vòng tròn).
		\item Tổng khối lượng của $3$ cánh quạt trong mỗi nhóm phải bằng nhau.
	\end{itemize}
	\begin{center}
		\begin{minipage}{0.45\textwidth}
			\begin{tikzpicture}[scale=1,>=stealth, line join=round, line cap=round, declare function={R=3; h=0.4; ecc=0.6;}]
				\path (0,0) coordinate (O);
				\fill[gray!30] (-R,0) arc (180:360:R and R*ecc) -- (R,-h) arc (0:-180:R and R*ecc) -- cycle;
				\filldraw[fill=white] (O) ellipse (R and R*ecc);
				\draw (-R,0) -- (-R,-h) arc (180:360:R and R*ecc) -- (R,0);
				\draw[fill=gray!10] (O) ellipse (0.5 and 0.5*ecc);
				\foreach \g in {0,40,...,320}{
					\filldraw[fill=gray!5] (\g:2.2 and 2.2*ecc) ellipse (0.35 and 0.35*ecc);
				}
			\end{tikzpicture}
		\end{minipage}
		\begin{minipage}{0.45\textwidth}
			\begin{tikzpicture}[scale=1,>=stealth, line join=round, line cap=round, smooth, samples=450]
				\def\R{3}
				\def\n{9}
				\def\ang{40}
				\def\xmin{-3.5} \def\xmax{3.5} \def\ymin{-3.5} \def\ymax{4}
				\path (0,0) coordinate (O)
				\foreach \i in {1,...,\n}{
					({90-(\i-1)*\ang}:\R) coordinate (P\i)
				};
				\draw (P1) \foreach \i in {2,...,\n}{-- (P\i)} -- cycle;
				\draw(O) circle (\R)
				(P1) \foreach \i in {2,...,\n}{-- (P\i)} -- cycle;
				\foreach \i/\g in {1/90,2/10,3/0,4/-10,5/-30,6/-120,7/-150,8/-180,9/140}
				\draw[fill=black] (P\i) circle (1pt) +(\g:.4) node{$\i$};
			\end{tikzpicture}
		\end{minipage}
	\end{center}
	Hỏi có bao nhiêu cách sắp xếp $9$ cánh quạt vào $9$ khe cắm thỏa mãn các điều kiện kỹ thuật trên (hai cách sắp xếp được coi là khác nhau nếu có ít nhất một cánh quạt ở một vị trí khe cắm khác nhau, không đồng nhất các cách lắp khác nhau bởi phép quay hay phép đối xứng của rotor)?
	\par
	\shortans[oly]{2\,592}
	\loigiai{
		Tổng khối lượng của các cánh quạt là $1+2+\ldots+9=45$ (gam).\\
		Ta chia $3$ nhóm cánh quạt sao cho tổng khối lượng của $3$ cánh quạt trong mỗi nhóm phải bằng nhau.\\ Suy ra, khối lượng của mỗi nhóm bằng $\dfrac{45}{3}=15$ (gam).\\
		Các bộ gồm $3$ số lấy từ tập $\{1;2;\ldots;9\}$ mà tổng của chúng bằng $15$ là\\
		\[\{1;5;9\},\{1;6;8\},\{2;4;9\},\{2;5;8\},\{2;6;7\},\{3;4;8\},\{3;5;7\},\{4;5;6\}.\]
		Chọn $3$ bộ trong số $8$ bộ trên sao cho các chữ số không trùng nhau để có được $3$ nhóm cánh quạt có tổng khối lượng thỏa mãn.
		\begin{itemize}
			\item \textbf{Trường hợp 1:} Ba bộ đó là $\{1;5;9\}$, $\{2;6;7\}$, $\{3;4;8\}$.\\
			Chọn ra $3$ tam giác đều thỏa đề bài có $3!$ (cách).\\
			Với mỗi tam giác đều, để lắp $3$ số vào $3$ đỉnh ta có $3!$ (cách).\\
			Vậy có tất cả $3$ tam giác đều.\\ 
			Do đó số cách xếp trong trường hợp này là $3!\cdot(3!)^3=1\,296$ (cách).
			\item \textbf{Trường hợp 2:} Ba bộ là $\{1;6;8\}$, $\{2;4;9\}$, $\{3;5;7\}$.\\
			Tương tự trường hợp $1$, có $1\,296$ (cách).
		\end{itemize}
		Vậy có tất cả $1\,296+1\,296=2\,592$ (cách).
	}
\end{ex}

\begin{ex}%[Đề thi thử môn Toán TRƯỜNG THPT Cổ Loa - Hà Nội năm học 2025-2026]%[Phạm Quốc Toàn - Phạm Hoàng Việt]%[0D8V2-5]
	Cho hình vuông $ABCD$ có tâm $O$. Gọi $M$, $N$, $P$, $Q$ lần lượt là trung điểm của các cạnh $AB$, $BC$, $CD$, $DA$. Kí hiệu $H_1$, $H_2$, $H_3$, $H_4$ lần lượt là các hình vuông $AMOQ$, $BMON$, $CPON$, $DPOQ$. Bạn An muốn tô màu mỗi hình vuông $H_1$, $H_2$, $H_3$, $H_4$ bởi một màu trong $5$ màu (đỏ, xanh, vàng, tím, nâu) sao cho hai hình vuông có chung cạnh sẽ được tô bởi hai màu khác nhau. Hỏi bạn An có bao nhiêu cách tô màu?
	\begin{center}
		\begin{tikzpicture}[scale=0.8, font=\footnotesize, line join=round, line cap=round, >=stealth]
			% Tọa độ các điểm
			\coordinate (A) at (0,4);
			\coordinate (B) at (4,4);
			\coordinate (C) at (4,0);
			\coordinate (D) at (0,0);
			%%%%%%%%%%%%
			\coordinate (M) at (2,4);
			\coordinate (N) at (4,2);
			\coordinate (P) at (2,0);
			\coordinate (Q) at (0,2);
			\coordinate (O) at (2,2);
			% Tô nền
			\fill[gray!20] (A) rectangle (C);
			% Khung ngoài
			\draw[thick] (A) -- (B) -- (C) -- (D) -- cycle;
			% Các đường chia
			\draw[thick] (M) -- (P);
			\draw[thick] (Q) -- (N);
			% Vẽ các điểm
			\foreach \p in {A,B,C,D,M,N,P,Q,O} \fill (\p) circle (2pt);
			% Nhãn các điểm
			\node[above left] at (A) {$A$};
			\node[above right] at (B) {$B$};
			\node[below right] at (C) {$C$};
			\node[below left] at (D) {$D$};
			%%%%%%%%%%%%%%%%%
			\node[above] at (M) {$M$};
			\node[right] at (N) {$N$};
			\node[below] at (P) {$P$};
			\node[left] at (Q) {$Q$};
			\node[below right] at (O) {$O$};
			% Nhãn các miền
			\node at (1,3) {$H_1$};
			\node at (3,3) {$H_2$};
			\node at (3,1) {$H_3$};
			\node at (1,1) {$H_4$};
		\end{tikzpicture}
	\end{center}
	\par\shortans[0]{260}
	\loigiai{
		\begin{itemize}
			\item 	\textbf{Trường hợp $1$:} Hình $H_1$, $H_3$ được tô cùng màu. 
			\begin{itemize}
				\item Có $5$ cách chọn màu cho hình $H_1$, $H_3$.
				\item Có $4$ cách chọn màu cho hình $H_2$.
				\item Có $4$ cách chọn màu cho hình $H_4$. \\ 
				Suy ra trường hợp $1$ có $5 \times 4 \times 4=80$ cách.
			\end{itemize}
		\item \textbf{Trường hợp $2$:} Hình $H_1$, $H_3$ được tô khác màu. 
		\begin{itemize}
			\item Có $5$ cách chọn màu cho hình $H_1$.
			\item Có $4$ cách chọn màu cho hình $H_3$.
			\item Có $3$ cách chọn màu cho hình $H_2$.
			\item Có $3$ cách chọn màu cho hình $H_4$.\\
			Suy ra trường hợp $2$ có $5 \times 4 \times 3 \times 3=180$ cách.
		\end{itemize}
	\end{itemize}
		Vậy tất cả có $80+180=260$ cách.
	}
\end{ex}

\begin{ex}%[Đề thi thử môn Toán TRƯỜNG THPT Cổ Loa - Hà Nội năm học 2025-2026]%[Phạm Quốc Toàn - Phạm Hoàng Việt]%[0D8V2-7]
	Sắp xếp ngẫu nhiên ba bạn nam $A$, $B$, $C$ và $5$ bạn nữ vào $8$ cái ghế được xếp theo hàng ngang (mỗi bạn ngồi một cái ghế). Gọi xác suất để $3$ bạn $A$, $B$, $C$ ngồi theo thứ tự đó từ trái qua phải, đồng thời giữa $A$ và $B$ có ít nhất một bạn nữ, giữa $B$ và $C$ có nhiều nhất một bạn nữ là $p$. Giá trị của $\dfrac{5}{p}$ bằng bao nhiêu?
	\par\shortans[0]{67{,}2}
	\loigiai{
		Số phần tử của không gian mẫu là $n(\Omega)=8!=40\,320$.
		\begin{itemize}
			\item \textbf{Trường hợp $1$:} Giữa $B$, $C$ không có bạn nữ nào.\\
			Ta gọi $x_1$, $x_2$, $x_3$ lần lượt là số bạn nữ xếp như sau 
			\[
			x_1 \quad A \quad x_2 \quad B\quad C \quad x_3\,\, \text{với}\,\, x_1\ge 0,\, x_2\ge 1,\, x_3\ge 0.
			\]
			Ta có $x_1+x_2+x_3=5 \Leftrightarrow x'_1+x_2+x'_3=7$ với $x'_1\ge 1$, $x_2\ge 1$, $x'_3 \ge 1$. \\
			Do vậy, ta có $\mathrm{C}_6^2 \cdot 5!$ cách.
			\item \textbf{Trường hợp $2$:} Giữa $B$, $C$ có một bạn nữ.\\
			Ta gọi $x_1$, $x_2$, $x_3$, $x_4$ lần lượt là số bạn nữ xếp như sau 
			\[
			x_1 \quad A \quad x_2 \quad B \quad x_3 \quad C \quad x_4\,\, \text{với}\,\, x_1\ge 0,\, x_2\ge 1,\, x_3=1,\, x_4\ge 0. 
			\]
			Ta có $x_1+x_2+1+x_4=5 \Leftrightarrow x_1+x_2+x_4=4\Leftrightarrow x'_1+x_2+x'_4=6$, với $x'_1\ge 1$, $x_2\ge 1$, $x'_4\ge 1$. \\
			Do vậy sẽ có $\mathrm{C}_5^2\cdot 5!$ cách.
		\end{itemize}
		Vậy có $\mathrm{C}_6^2 \cdot 5!+\mathrm{C}_5^2\cdot 5!=3\, 000$ (cách). \\ 
		Vậy xác suất là $p=\dfrac{3\,000}{8!}$. \\
		Do vậy $\dfrac{5}{p}=\dfrac{336}{5}=67{,}2$. 
	}
\end{ex}

\begin{ex}%[0H5V3-5]%[Dự án C-TT THPT NH25-26-Đợt 4- ĐỖ CHÍ TÂM]%[THPT Phan Bội Châu-Khánh Hòa]%[GVPB: Phạm Văn Long]
	Cho hình hộp $ABCD.A'B'C'D'$. Gọi $N$ là điểm trên đoạn $C'D$ sao cho $DN=\dfrac{1}{3}D{C}'$. Khi phân tích $\overrightarrow{BN}=x\cdot \overrightarrow{BA}+y\cdot\overrightarrow{BC}+z\cdot \overrightarrow{B{B}'}$ thì giá trị $3x+2y+3z$ bằng bao nhiêu?
	\shortans{$5$}
	\loigiai{
		\begin{center}
			\begin{tikzpicture}[scale=0.7]
				%%\draw[gray!20] (-3,-2) grid (6,5);
				\path (0,0) coordinate (A)  
				(-2,-2) coordinate (B)
				(5,0) coordinate (D)
				(3,-2) coordinate (C)
				(0,4) coordinate (A')
				(-2,2) coordinate (B')
				(3,2) coordinate (C')
				(5, 4) coordinate (D')
				($(D)!0.33!(C')$) coordinate(N);
				\draw[dashed] (A')--(A)--(B) (A)--(D) (B)--(N);
				\draw (B)--(C)--(D)--(D')--(C')--(B')--(A') (B')--(B) (C')--(C) (A')--(D') (C')--(D) (B)--(C');
				\foreach \p/\q in {A/180, B/-90, C/-90, D/0, A'/90, B'/90, C'/90, D'/90, N/45}
				\fill[blue] (\p) circle(2pt) node[shift={(\q:3mm)}]{\bfseries $\p$};
			\end{tikzpicture}
		\end{center}
		Ta có $\overrightarrow{BN}=\overrightarrow{B{C}'}+\overrightarrow{C'N}=\left(\overrightarrow{B{B}'}+\overrightarrow{BC} \right)+\dfrac{2}{3}\overrightarrow{C'D}$ $=\left(\overrightarrow{B{B}'}+\overrightarrow{BC} \right)+\dfrac{2}{3}\left(\overrightarrow{C'D'}+\overrightarrow{C'C} \right)$\\
		\hspace*{4.5cm}$=\left(\overrightarrow{B{B}'}+\overrightarrow{BC} \right)+\dfrac{2}{3}\left(\overrightarrow{BA}+\left(-\overrightarrow{B{B}'} \right) \right)$ $=\dfrac{2}{3}\overrightarrow{BA}+\overrightarrow{BC}+\dfrac{1}{3}\overrightarrow{B{B}'}$.\\
		Suy ra $x=\dfrac{2}{3}$; $y=1$; $z=\dfrac{1}{3}$.\\
		Vậy $3x+2y+3z=3\cdot \dfrac{2}{3}+2\cdot 1+3\cdot \dfrac{1}{3}=5$.
	}
\end{ex}

\begin{ex}%[0H9V4-7]%[KNTT - Lớp 12 - Dự án C - Đợt 2]%[Loc do]%[PB: Phan Anh]
	\immini[thm]{Miền Trung Việt Nam vốn luôn phải oằn mình chống chọi với thiên tai. Giả sử trong một đợt áp thấp nhiệt đới mạnh lên thành bão, Trung tâm Dự báo Khí tượng Thủy văn Quốc gia phát đi thông báo khẩn về cơn bão số $4$ có tên quốc tế là Sao La. Trên bản đồ quy hoạch phòng chống thiên tai được gắn hệ trục tọa độ Oxy với đơn vị đo là $10$ km (hướng Đông là trục $Ox$, hướng Bắc là trục $Oy$), vị trí ba thành phố trọng điểm là Hà Tĩnh, Vinh (Nghệ An) và Thanh Hóa được xác định lần lượt tại các điểm $T(0;0)$, $N(-2;3)$ và $H(-1;5)$. Tại thời điểm bản tin phát đi, tâm bão Sao La đang ở ngoài khơi Biển Đông, cách thành phố Hà Tĩnh $200$ km. Vị trí tâm bão nằm ở hướng Đông Nam so với Hà Tĩnh, tạo với hướng Đông một góc $\alpha $ sao cho $\cos \alpha =0{,}8$. }{\begin{tikzpicture}[>=stealth,line join=round,line cap=round,font=\footnotesize,scale=0.8]
			\draw[->] (-1,0)--(6,0) node[below]{$x$};
			\draw [->] (0,-4)--(0,4) node [right]{$y$};
			\path (0,0) coordinate (T)
			(5,-3) coordinate (S) (3.5,-1.2) coordinate (S1) 
			(3,0) coordinate (B)  (2,-3) coordinate (C) ;
			\fill (3.5,-1.2) circle (1pt) node[right] {$S$};
			\fill (-0.2,3) circle (1pt) node[left] {$H$};
			\fill (-0.5,2) circle (1pt) node[left] {$N$};
			\draw[dotted,green,thick] (S) circle (0.8cm);
			\draw[thick,red] (S1) circle (1cm);
			\coordinate (A) at ($(S1)+(180:1)$);
			\path pic[angle radius=5mm,fill=green,draw=blue,"$\alpha$",angle eccentricity=1.5] {angle = S--T--B};
			\path pic[angle radius=5mm,fill=green,draw=blue,"$60^\circ$",angle eccentricity=1.5] {angle = S1--S--C};
			\draw (T)--(S)--(S1)--(A) (S)--(C) ;
			\node at (3,-1) {$r(t)$};
			\node at (5.6,0.2) {Đông};
			\node at (-0.5,4) {Bắc};
			\foreach \t/\g in {T/-120,S/0}{
				\fill (\t) circle (1pt) node[shift={(\g:7pt)},font=\scriptsize]{$ \t $};
			}
			
	\end{tikzpicture}}
	Cơn bão di chuyển phức tạp theo hướng Tây Bắc lệch $60^\circ$ so với hướng Tây với vận tốc $20$ km/h. Sức tàn phá của bão rất lớn với vùng nguy hiểm là một hình tròn có bán kính ban đầu $20$ km và liên tục mở rộng thêm $10$ km mỗi giờ. Để lên phương án sơ tán dân cư đồng bộ, Ban chỉ đạo cần biết khoảng thời gian mà cả $3$ thành phố này cùng nằm trong vùng nguy hiểm của bão Sao La kéo dài bao nhiêu giờ? (làm tròn kết quả đến hàng phần chục).
	\par\shortans[0]{$10,8$}
	\loigiai{
		Ta có $T(0;0)$, $N(-2;3)$ và $H(-1;5)$.\\
		Tại thời điểm bản tin phát đi, tâm bão Sao La có tọa độ là $(x;y)$.\\
		Suy ra $\heva{&x=20\cdot \cos\alpha\\& y=-20\cdot \sin\alpha}\Rightarrow S(16;-12)$.\\
		Bão di chuyển theo hướng Tây Bắc lệch $60^\circ$ so với hướng Tây nên vectơ chỉ phương có tọa độ 
		\[\left(\cos120^\circ; \sin120^\circ\right)=\left(-\dfrac{1}{2}; -\dfrac{\sqrt{3}}{2}\right).\]
		Vì vận tốc $20$ km/h nên vectơ vận tốc \[\overrightarrow{v}=2\left(-\dfrac{1}{2}; -\dfrac{\sqrt{3}}{2}\right)=\left(-1;\sqrt{3}\right).\]
		Tọa độ tâm bão sau $t$ (giờ) là $S\left(16-t;\; -12+\sqrt{3} t\right)$.\\
		Bán kính vùng nguy hiểm sau $t$ (giờ) là $r=2+t$. Suy ra 
		\begin{itemize}
			\item $SH^2 =\left(17-t\right)^2 +\left(\sqrt{3} t-17\right)^2 $
			\item $ST^2 =\left(16-t\right)^2 +\left(\sqrt{3} t-12\right)^2 $
			\item $SN^2 =\left(18-t\right)^2 +\left(\sqrt{3} t-15\right)^2$.
		\end{itemize} 
		Thành phố Thanh Hóa, Hà Tĩnh, Vinh nằm trong vùng nguy hiểm khi khoảng cách từ tâm bão đến thành phố nhỏ hơn hoặc bằng bán kính vùng nguy hiểm do vậy ta được
		\begin{eqnarray*}
			\heva{& SH^2 \le r^2\\ & ST^2 \le r^2\\ & SN^2 \le r^2} &\Leftrightarrow& \heva{& 3t^2 -\left(38+34\sqrt{3}\right)t+574\le 0\\ & t^2 -\left(12+8\sqrt{3}\right)t+132\le 0\\ & 3t^2 -\left(40+30\sqrt{3}\right)t+545\le 0}\\
			&\Leftrightarrow& t\in \left[\dfrac{40+30\sqrt{3} -\sqrt{2400\sqrt{3} -2240}}{6} ;\; \dfrac{12+8\sqrt{3} +\sqrt{192\sqrt{3} -192}}{2} \right]
		\end{eqnarray*}
		Khoảng thời gian kéo dài là 
		\[\dfrac{12+8\sqrt{3} +\sqrt{192\sqrt{3} -192}}{2} -\dfrac{40+30\sqrt{3} -\sqrt{2400\sqrt{3} -2240}}{6} \approx 10{,}8\,\text{giờ}.\]
	}
\end{ex}

\begin{ex}%[1C2V3-1]%[Dự án C đề thi thử TN Môn Toán NH25-26 - Dot 4 - Hector Tran]%[Liên trường THPT cụm 09 - Đợt 2 - Hà Nội]%[GVPB - Thanh Phát]
	\immini [thm]{
		Hằng năm, trường THPT M tổ chức các đoàn đi thăm hỏi và chúc Tết tứ thân phụ mẫu cao tuổi của cán bộ, giáo viên và nhân viên trong trường. Đoàn số $1$ có nhiệm vụ đến thăm $4$ gia đình tại các địa điểm $A$, $B$, $C$, $D$. Đoàn xuất phát từ trường THPT M, đến thăm đủ $4$ gia đình (mỗi gia đình đúng một lần). Sau khi thăm xong gia đình cuối cùng, đoàn kết thúc hành trình tại đó (không quay về trường). Do điều kiện giao thông dịp cuối năm, đoạn đường $AB$ chỉ có thể đi theo chiều từ $A$ đến $B$, không đi được theo chiều ngược lại (\textit{hình vẽ}). Đoàn số $1$ đã chọn được lộ trình di chuyển có tổng quãng đường di chuyển ngắn nhất. Tổng quãng đường ngắn nhất đó dài bao nhiêu km?
	}{
		\begin{tikzpicture}[scale=0.5, line join=round, line cap=round, every node/.style={scale=0.8}, declare function={xM=-4; yM=1; xC=0; yC=4.5; xB=4.5; yB=0; xA=2.5; yA=-3.5; xD=-2; yD=-3; R=0.5;}]		
			\path
			(xM, yM) coordinate (M)
			(xC, yC) coordinate (C)
			(xB, yB) coordinate (B)
			(xA, yA) coordinate (A)
			(xD, yD) coordinate (D);		
			\draw[-stealth, shorten >=0.2cm, shorten <=0.2cm] (M) -- (C) node[pos=0.5, sloped, above] {$4$ km};
			\draw[stealth-stealth, shorten >=0.2cm, shorten <=0.2cm] (C) -- (B) node[pos=0.5, sloped, above] {$3$ km};
			\draw[-stealth, shorten >=0.2cm, shorten <=0.2cm] (A) -- (B) node[pos=0.5, sloped, below] {$3$ km};
			\draw[stealth-stealth, shorten >=0.2cm, shorten <=0.2cm] (D) -- (A) node[pos=0.5, sloped, below] {$5$ km};
			\draw[-stealth, shorten >=0.2cm, shorten <=0.2cm] (M) -- (D) node[pos=0.5, sloped, below] {$7$ km};		
			\draw[-stealth, shorten >=0.2cm, shorten <=0.2cm] (M) -- (B) node[pos=0.5, sloped, above] {$6$ km};
			\draw[-stealth, shorten >=0.2cm, shorten <=0.2cm] (M) -- (A) node[pos=0.3, sloped, above] {$8$ km};
			\draw[stealth-stealth, shorten >=0.2cm, shorten <=0.2cm] (D) -- (C) node[pos=0.7, sloped, above] {$5$ km};
			\draw[stealth-stealth, shorten >=0.2cm, shorten <=0.2cm] (D) -- (B) node[pos=0.7, sloped, above] {$4$ km};
			\draw[stealth-stealth, shorten >=0.2cm, shorten <=0.2cm] (C) -- (A) node[pos=0.25, sloped, above] {$6$ km};		
			\foreach \p/\pos in {M/135, C/90, B/0, A/-45, D/-135} {
				\draw (\p) circle (0.6pt);
				\draw (\p) circle (0.5cm) node[shift={(\pos:20pt)}] {$\p$};
			}
		\end{tikzpicture}
	}
	\par\shortans{16}
	\loigiai{
		- Gán nhãn cho điểm xuất phát $\ell(M)=0$ và coi đây là nhãn vĩnh viễn.\\ 		
		Các đỉnh $A$, $B$, $C$, $D$ lúc này có nhãn tạm thời dựa trên khoảng cách từ $M$ là $\ell(A)=8$, $\ell(B)=6$, $\ell(C)=4$, $\ell(D)=7$.\\
		- So sánh các nhãn tạm thời, ta thấy $\ell(C)=4$ là nhỏ nhất.\\ 
		Ta chọn $C$ làm điểm tiếp theo của hành trình.\\
		- Từ $C$, ta xem xét các gia đình chưa thăm ($A$, $B$, $D$).
		\begin{itemize}
			\item Khoảng cách $M$ đến $A$ là $\ell(C)+AC=4+6=10$ suy ra $\ell(A)=10$.
			\item Khoảng cách $M$ đến $B$ là $\ell(C)+CB=4+3=7$ suy ra $\ell(B)=7$.
			\item Khoảng cách $M$ đến $D$ là $\ell(C)+CD=4+5=9$ suy ra $\ell(D)=9$.
		\end{itemize}
		Hiện tại nhãn nhỏ nhất từ các lựa chọn là $\ell(B)=7$.\\ 		
		Ta chọn $B$ là điểm tiếp theo.\\
		- Từ $B$, ta còn $A$ và $D$.\\ 		
		Do tuyến đường từ $B$ đến $A$ bị chặn, nên đi đến $D$ trước, rồi tới $A$.\\ 
		Vậy lộ trình di chuyển có tổng quãng đường di chuyển ngắn nhất là 		
		\[MC+CB+BD+DA=\ell(B)+BD+DA=7+4+5=16\text { (km).}\]
	}
\end{ex}

\begin{ex}%[1D2C2-5]%[Dự án C đề thi thử TN Môn Toán NH25-26 - Dot 4 - Hector Tran]%[Liên trường THPT cụm 09 - Đợt 2 - Hà Nội]%[GVPB - Thanh Phát]
	Một hộp chứa $100$ cái thẻ được đánh số thứ tự liên tiếp từ $1$ đến $100$. Hai thẻ khác nhau thì đánh số thứ tự khác nhau. Chọn ngẫu nhiên $3$ thẻ trong hộp. Có bao nhiêu cách chọn $3$ thẻ có số thứ tự lập thành cấp số cộng, đồng thời có tổng không vượt quá $150$?
	\par\shortans{1225}
	\loigiai{
		Giả sử các số trên $3$ thẻ lập thành cấp số cộng với số hạng đầu là $a$, công sai $d$, với $a$, $d \in \mathbb{N}^*$.\\
		Tổng các số trên $3$ thẻ không vượt quá $150$ nên ta có
		\begin{eqnarray*}
			&&a + (a + d) + (a + 2d) \le 150 \\
			&\Leftrightarrow& 3a + 3d \le 150 \\
			&\Leftrightarrow& a + d \le 50.
		\end{eqnarray*}
		Vì $a$, $d \in \mathbb{N}^*$ nên
		\begin{itemize}
			\item Với $d = 1$ thì $a \le 49$. Với mỗi số tự nhiên $a \in \mathbb{N}^*$ ta có $1$ cấp số cộng. Do đó trường hợp này ta thu được $49$ cấp số cộng thỏa mãn.
			\item Với $d = 2$ thì $a \le 48$. Với mỗi số tự nhiên $a \in \mathbb{N}^*$ ta có $1$ cấp số cộng. Do đó trường hợp này ta thu được $48$ cấp số cộng thỏa mãn.
			\item \ldots
			\item Cứ tiếp tục như vậy với $d = 49$ thì $a \le 1$, trường hợp này có $1$ cấp số cộng thỏa mãn.
		\end{itemize}
		Do đó số cách chọn thỏa yêu cầu đề là $49 + 48 + \cdots + 1 = \dfrac{49 \cdot 50}{2} = 1\,225$.
	}
\end{ex}

\begin{ex}%[Dự án đề TT Toán NH25-26-Dot 4- Hải Phụng]%[THCS-THPT LÊ THÁNH TÔNG - TP.HCM]%[GVPB - Viet Hoang Pham]%[1D2V2-7]
	Trong dịp Tết Nguyên Đán, bạn Kiên nhận được tiền lì xì ngày mùng $1$ là $3\,000\,000$ đồng. Kể từ mùng $2$, mỗi ngày tiền lì xì giảm $300\,000$ đồng so với ngày trước đó. Nhưng đến ngày mùng $3$, mẹ bạn Kiên lại nói: \lq\lq Để mẹ giữ tiền cho sau này cưới vợ nhé\rq\rq. Nên bạn Kiên quyết định từ ngày hôm đó trở đi (từ mùng $3$ đến mùng $9$) mỗi ngày sẽ đưa cho mẹ $X$ triệu đồng. Biết rằng đến hết ngày mùng $9$, Kiên còn lại đúng $5\,000\,000$ đồng. Tính giá trị của $X$.
	\shortans{$1,6$}
	\loigiai{
		Gọi $u_n$ là số tiền lì xì Kiên nhận được vào ngày mùng $n$. \\
		Theo giả thiết, dãy số $(u_n)$ là một cấp số cộng có số hạng đầu $u_1 = 3\,000\,000$ và công sai $d = -300\,000$.\\
		Tổng số tiền lì xì Kiên nhận được từ mùng $1$ đến hết mùng $9$ là tổng của $9$ số hạng đầu tiên của cấp số cộng:
		$$S_9 = \dfrac{9}{2} \left[ 2u_1 + (9-1)d \right] = \dfrac{9}{2} \left[ 2 \cdot 3\,000\,000 + 8 \cdot (-300\,000) \right] = 16~200\,000 \text{ (đồng).}$$
		Từ ngày mùng $3$ đến ngày mùng $9$ có tất cả $9 - 3 + 1 = 7$ ngày.\\
		Mỗi ngày Kiên đưa cho mẹ $X$ triệu đồng (tức là $1\,000\,000X$ đồng).\\
		Tổng số tiền Kiên đưa cho mẹ là: $7 \cdot 1\,000\,000X = 7\,000\,000X$ (đồng).\\
		Số tiền còn lại của Kiên sau ngày mùng $9$ là
		{\allowdisplaybreaks
			\begin{eqnarray*}
				& & 16~200\,000 - 7\,000\,000X = 5\,000\,000  \\
				&\Leftrightarrow& 7\,000\,000X = 11~200\,000\\
				&\Leftrightarrow& X = \dfrac{11~200\,000}{7\,000\,000} = 1{,}6.
			\end{eqnarray*}
		}
	}
\end{ex}

\begin{ex}%[1D2V3-8]%[Dự án C đợt 3 NH25-26- Phạm Đăng Đăng]%[TT-THPT-TranNhanTong-HaNoi-N25-26]%[GVPB - Phan Quốc Trí]
	Bác An dự định gửi vào ngân hàng một số tiền với lãi suất không đổi là $7{,}5\%$ một năm. Biết rằng cứ sau mỗi năm số tiền lãi sẽ được nhập vào vốn ban đầu. Bác An cần gửi ít nhất vào ngân hàng bao nhiêu triệu đồng (kết quả làm tròn đến hàng đơn vị) để sau $5$ năm Bác rút được số tiền đủ mua $1$ chiếc xe máy điện trị giá $57$ triệu đồng.
	\shortans{$40$}
	\loigiai{Gọi $u_n$ là số tiền Bác An nhận được sau $n-1$ năm.\\ Khi đó $(u_n)$ là một cấp số nhân với công bội $q=1,075$.\\
		Ta có $u_6=u_1\cdot q^5\Leftrightarrow u_1=\dfrac{57}{1,075^5}\approx 40$ triệu đồng.
	}
\end{ex}

\begin{ex}%[1D5H1-3]%[Dự án C-TT THPT NH25-26-Đợt 4- ĐỖ CHÍ TÂM]%[THPT Phan Bội Châu-Khánh Hòa]%[GVPB: Phạm Văn Long]
	Thâm niên công tác của các công nhân hai nhà máy A và B được cho như trong bảng dưới đây.
	\begin{table}[h]
		\centering
		\begin{tabular}{|l|c|c|c|c|c|}
			\hline
			Thâm niên công tác (năm) & \textbf{$[0; 5)$} & \textbf{$[5; 10)$}\, & \textbf{$[10; 15)$} & \textbf{$[15; 20)$} & \textbf{$[20; 25)$} \\ \hline
			Số công nhân nhà máy $A$ & $30$ & $18$ & $20$ & $22$ & $10$ \\ \hline
			Số công nhân nhà máy $B$ & $27$ & $15$ & $30$ & $16$ & $12$ \\ \hline
		\end{tabular}
	\end{table}\\
	Tính tổng hai phương sai về thâm niên công tác của công nhân làm việc tại hai nhà máy trên (kết quả được qui tròn đến hàng phần chục).\shortans{$91{,}4$}
	\loigiai{
		\begin{table}[h]
			\centering
			\begin{tabular}{|l|c|c|c|c|c|}
				\hline
				{Thâm niên công tác (năm)} & \textbf{$[0; 5)$} & \textbf{$[5; 10)$}\, & \textbf{$[10; 15)$} & \textbf{$[15; 20)$} & \textbf{$[20; 25)$} \\ \hline
				Giá trị đại diện &$2{,}5$& $7{,}5$ & $12{,}5$ & $17{,}5$ & $22{,}5$ \\ \hline
			\end{tabular}
		\end{table}\\
		{\bf Nhà máy A:}\\
		Tổng số công nhân: $N_A=30+18+20+22+10=100$.\\
		Số trung bình:
		$\overline{x}_A=\dfrac{30\cdot2{,}5+18\cdot7{,}5+20\cdot12{,}5+22\cdot17{,}5+10\cdot22{,}5}{100}=10{,}7$.\\  
		Phương sai:\\ $s_A^2=\dfrac{30(2{,}5-10{,}7)^2+18(7{,}5-10{,}7)^2+20(12{,}5-10{,}7)^2+22(17{,}5-10{,}7)^2+10(22{,}5-10{,}7)^2}{100}$\\
		\hspace*{0.6cm}$=46{,}76$.\\
		{\bf Nhà máy B:}\\
		Tổng số công nhân: $N_B=27+15+30+16+12=100$.\\
		Số trung bình:
		$\overline{x}_B=\dfrac{27\cdot2{,}5+15\cdot7{,}5+30\cdot12{,}5+16\cdot17{,}5+12\cdot22{,}5}{100}=11{,}05$.\\ 
		Phương sai:\\
		$s_B^2=\dfrac{27(2{,}5-11{,}05)^2+15(7{,}5-11{,}05)^2+30(12{,}5-11{,}05)^2+16(17{,}5-11{,}05)^2+12(22{,}5-11{,}05)^2}{100}$\\
		\hspace*{0.6cm}$=44{,}6475$.\\  
		Tổng hai phương sai: $s_A^2+s_B^2=46{,}76+44{,}6475=91,4075\approx 91{,}4$.
	}
\end{ex}

\begin{ex}%[1D6H2-5]%[TT-LienTruong-HaNoi-N25-26]%[Phạm Quốc Toàn - Lê Phúc]
	Cường độ trận động đất (Richter) được cho bởi công thức $M = \log A - \log A_0$, với $A$ là biên độ rung chấn tối đa và $A_0$ là biên độ chuẩn (hằng số). Đầu thế kỷ $20$, một trận động đất ở San Francisco có cường độ $6{,}7$ độ Richter. Trong cùng năm đó, một trận động đất khác ở Nam Mỹ có biên độ mạnh hơn gấp $4$ lần. Cường độ của trận động đất ở Nam Mỹ là bao nhiêu? (\textit{không làm tròn kết quả các phép tính trung gian, chỉ làm tròn kết quả cuối cùng đến hàng phần mười}).
	\par\shortans[0]{7{,}3}
	\loigiai{
		Gọi $M_1$, $A_1$ lần lượt là cường độ và biên độ rung chấn tối đa trận động đất tại San Francisco.\\
		Ta có 
		\[
		M_1 = 6{,}7 = \log A_1 - \log A_0 \Leftrightarrow 6{,}7 = \log \dfrac{A_1}{A_0} \Leftrightarrow A_1=10^{6{,}7} \cdot A_0. 
		\]
		Gọi $M_2$, $A_2$ lần lượt là cường độ và biên độ rung chấn tối đa trận động đất tại Nam Mỹ.\\
		Ta có $A_2=4A_1=4\cdot 10^{6{,}7}\cdot A_0$. \\ 
		Vậy cường độ trận động đất tại Nam Mỹ là 
		\[
		M_2 = \log \dfrac{A_2}{A_0} =\log \dfrac{4\cdot 10^{6{,}7} \cdot A_0}{A_0} = \log \left(4 \cdot 10^{6{,}7}\right) \approx 7{,}3. 
		\].
	}
\end{ex}

\begin{ex}%[1D6H3-5]%[Dự án đề TT Toán NH25-26-Dot 3- Hải Phụng]%[Cụm 13-Hải Phòng]%[GVPB - Viet Hoang Pham]
Một người gửi tiết kiệm vào ngân hàng $200$ triệu đồng theo hình thức lãi kép (tức là tiền lãi được cộng vào vốn của kỳ kế tiếp). Ban đầu người đó gửi với kỳ hạn $3$ tháng, lãi suất $2{,}1 \%$/kỳ hạn, sau $2$ năm người đó thay đổi phương thức gửi, chuyển thành kỳ hạn $1$ tháng với lãi suất $0{,}65 \%$/tháng. Tính tổng số tiền lãi nhận được sau $5$ năm (đơn vị là triệu đồng và kết quả làm tròn đến hàng phần mười).
\shortans[oly]{$98{,}2$}
\loigiai{Xét $2$ năm đầu tiên, số tiền nhận được là
$$T_1 = 200 (1 + 2{,}1\%)^{2 \cdot \frac{ 12}{3}} \approx 236{,}176 \text{ (triệu đồng).}$$
Số tiền nhận được sau $5$ năm là
$$T_2 = T_1 (1 + 0{,}65\%)^{3 \cdot 12} = 298{,}217 \text{ (triệu đồng).}$$ 
Tổng số tiền lãi nhận được sau $5$ năm là $298{,}217 - 200 = 98{,}2$  (triệu đồng). 			
}
\end{ex}

\begin{ex}%[1D7V2-8]%[TT Cụm 5 - Ninh Bình-NH25-26]%[GVSB:Dương Công Tạo]%[GVPB1: Nguyễn Tấn Tài]
	Giả sử dân số Việt Nam được dự báo theo mô hình logistic, giai đoạn từ năm 2\,023 đến năm 2\,035 là hàm số $P(t) = \dfrac{120}{1+0{,}2\mathrm{e}^{-0{,}06t}}$ (triệu người), trong đó $t$ là số năm tính từ 2\,023. Chi phí an sinh xã hội bình quân theo đầu người được mô hình hóa bằng hàm số $C(t) = 25-20\mathrm{e}^{-0{,}05t}$ (triệu đồng/đầu người/năm). Tính tốc độ thay đổi của tổng chi phí an sinh xã hội toàn quốc (nghìn tỷ đồng/năm) vào đầu năm 2\,030 (\textit{làm tròn kết quả đến hàng đơn vị}).
	\par  \shortans{83}
	\loigiai{
		Gọi $T(t)$ là tổng chi phí an sinh xã hội toàn quốc tại thời điểm $t$ (tính từ đầu năm 2\,023). \\
		Khi đó $T(t) = P(t)\cdot C(t)=\dfrac{120}{1+0{,}2\mathrm{e}^{-0{,}06t}}\cdot \left(25-20\mathrm{e}^{-0{,}05t}\right)$. \\
		Ta có $T'(t) = \dfrac{1{,}44\mathrm{e}^{-0{,}06t}}{(1+0{,}2\mathrm{e}^{-0{,}06t})^2}(25-20\mathrm{e}^{-0{,}05t}) + \dfrac{120}{1+0{,}2\mathrm{e}^{-0{,}06t}}\cdot\mathrm{e}^{-0{,}05t}$ \\
		Tốc độ thay đổi của tổng chi phí an sinh xã hội toàn quốc (nghìn tỷ đồng/năm) vào đầu năm 2\,030 là $T'(7) \approx 83$ (nghìn tỷ đồng/năm).
	}
\end{ex}

\begin{ex}%[1D9V2-5]%[Dự án C - Đợt 3 - Sở Giáo Dục Sơn La]%[Lê Hùng Thắng - Nguyễn Tiến]% 
	Trong một cuộc thi đấu Robotics, sân đấu được thiết kế dạng lưới ô vuông như hình vẽ.
	\begin{center}
		\begin{tikzpicture}[scale=1,>=stealth, font=\footnotesize, line join=round, line cap=round]
			\draw (0,0) grid (6,4);
			\foreach \x in {0,1,...,6}{
				\foreach \y in {0,1,...,4}
				\fill (\x,\y) circle (2.5pt);}
			\draw (2,2) node [above right] {$M$};
			\draw (3,1) node [above right] {$N$};
			\draw (0,4) node [above left] {$A$};
			\draw (6,0) node [below right] {$B$};
		\end{tikzpicture}
	\end{center}
	Các robot xuất phát từ vị trí điểm ${A}$, di chuyển ngẫu nhiên theo cạnh của các ô vuông theo hướng xuống dưới hoặc sang phải đến vị trí điểm ${B}$. Tính xác suất robot đi từ ${A}$ đến ${B}$ mà không đi qua cả ${M}$ và ${N}$ (kết quả làm tròn đến chữ số hàng phần trăm).
	\shortans{$0,31$}
	\loigiai{
		\begin{center}
			\begin{tikzpicture}[scale=1,>=stealth, font=\footnotesize, line join=round, line cap=round]
				\draw (0,0) grid (6,4);
				\foreach \x in {0,1,...,6}{
					\foreach \y in {0,1,...,4}
					\fill (\x,\y) circle (2.5pt);}
				\draw (2,2) node [above right] {$M$};
				\draw (3,1) node [above right] {$N$};
				\draw (0,4) node [above left] {$A$};
				\draw (6,0) node [below right] {$B$};
			\end{tikzpicture}
		\end{center}
		Số cách di chuyển từ $A$ đến $B$ là $n(\Omega)=\mathrm{C}_{10}^4=210$. \\
		Gọi $E_M\colon$\lq\lq Robot đi qua điểm $M$\rq\rq.\\
		Ta có $n(E_M)=\mathrm{C}_5^2\cdot \mathrm{C}_5^2=10\cdot 10=100$.\\
		Gọi $E_N\colon$\lq\lq Robot đi qua điểm $N$\rq\rq.\\
		Ta có $n(E_N)=\mathrm{C}_7^3\cdot \mathrm{C}_3^1=35\cdot 3=105$.\\
		Số cách robot đi qua cả $M$ và $N$ là\\
		$n(E_M\cap E_N)=\mathrm{C}_5^2\cdot \mathrm{C}_2^1\cdot \mathrm{C}_3^1=10\cdot2\cdot3=60$.\\
		Số cách robot đi qua ít nhất một trong hai điểm $M$ hoặc $N$ là
		$$n(E_M\cup E_N)=n(E_M)+n(E_N)-n(E_M\cap E_N)=100+105-60=145>$$
		Số cách robot đi từ $A$ đến $B$ mà không đi qua $M$ và $N$ là
		$$n(X)=210-145=65.$$
		Xác suất cần tìm là $\mathrm{P}(X)=\dfrac{65}{210}\approx 0{,}31$.
	}
\end{ex}

\begin{ex}%[GV soạn: Nguyễn Quang Hiệp GVPB2 Đỗ Minh Vũ ]%[1D9V2-5]
	Trong thư viện có $3$ quyển sách toán, $3$ quyển sách lý, $3$ quyển sách hóa, $3$ quyển sách sinh. Biết các quyển sách cùng môn giống nhau, xếp $12$ quyển sách trên lên giá thành một hàng. Tính xác suất sao cho khi xếp không có $3$ quyển nào cùng môn đứng cạnh nhau (Kết quả làm tròn đến hàng phần trăm).
	\par\shortans[oly]{0{,}84}
	\loigiai{
		Có tất cả $12$ quyển sách, trong đó có $4$ nhóm (gồm toán, lý, hóa và sinh) các quyển sách cùng môn giống nhau. Số cách xếp là $n(\Omega) = \dfrac{12!}{3! \cdot 3! \cdot 3! \cdot 3!} = 369\,600$ (cách).\\
		Gọi $A_1, A_2, A_3, A_4$ lần lượt là tập hợp các cách xếp mà môn toán, lý, hóa, sinh các quyển sách cùng môn đứng cạnh nhau.\\
		\textbf{Trường hợp 1} Một nhóm $3$ quyển sách cạnh nhau (xem $3$ quyển cùng môn là một khối và $9$ quyển còn lại)
		\[n(A_i) = \dfrac{10!}{3! \cdot 3! \cdot 3! \cdot 1!} \cdot \mathrm{C}_4^1 = 67\,200.\]
		\textbf{Trường hợp 2} Hai nhóm mà mỗi nhóm gồm $3$ quyển sách cạnh nhau (xem $2$ khối và $6$ quyển còn lại)
		\[n(A_i \cap A_j) = \dfrac{8!}{3! \cdot 3! \cdot 1! \cdot 1!} \cdot \mathrm{C}_4^2 = 6\,720.\]
		\textbf{Trường hợp 3} Ba nhóm mà mỗi nhóm gồm $3$ quyển sách cạnh nhau (xem $3$ khối và $3$ quyển còn lại)
		\[n(A_i \cap A_j \cap A_k) = \dfrac{6!}{3! \cdot 1! \cdot 1! \cdot 1!} \cdot \mathrm{C}_4^3 = 480.\]
		\textbf{Trường hợp 4} Cả bốn nhóm mà mỗi nhóm gồm $3$ quyển sách cạnh nhau (xem $4$ khối)
		\[n(A_1 \cap A_2 \cap A_3 \cap A_4) = \dfrac{4!}{1! \cdot 1! \cdot 1! \cdot 1!} \cdot \mathrm{C}_4^4 = 24.\]
		Khi đó \allowdisplaybreaks
		\begin{eqnarray*}
			n(A_1 \cup A_2 \cup A_3 \cup A_4) &= &n(A_i) - n(A_i \cap A_j) + n(A_i \cap A_j \cap A_k) - n(A_1 \cap A_2 \cap A_3 \cap A_4)\\
			&=& 67\,200 - 6\,720 + 480 - 24 \\
			&=& 60\,936.\\
		\end{eqnarray*}
		Gọi $E$ là biến cố \lq\lq xếp không có $3$ quyển nào cùng môn đứng cạnh nhau\rq\rq\\
		$n(E) = 369\,600 - 60\,936 = 308\,664$.\\
		Xác suất sao cho khi xếp không có $3$ quyển nào cùng môn đứng cạnh nhau là
		\[\mathrm{P}(E) = \dfrac{n(E)}{n(\Omega)} = \dfrac{308\,664}{369\,600} \approx 0{,}84.\]
	}
\end{ex}

\begin{ex}%[1H8C1-4]%[Dự án C đề thi thử TN Môn Toán NH25-26 - Dot 4 - Phạm Đăng Đăng]%[THPT Lê Quý Đôn-Đống Đa-Hà Nội]%[GVPB - Phan Quốc Trí]		
	Để chuẩn bị chào đón năm mới xuân Bính Ngọ 2026, người ta cần trang trí dây đèn Led quanh một tòa nhà hình kim tự tháp là hình chóp tứ giác đều $S.ABCD$. Biết rằng cạnh bên của tòa nhà dài $40$ m, góc $\widehat{BSC}=15^\circ$. Dây đèn Led được trang trí theo đường $AMNPQRTUVS$ (tham khảo hình vẽ), trong đó $VS=3$ m. Hỏi cần dùng ít nhất bao nhiêu mét dây đèn Led để trang trí (kết quả làm tròn đến hàng phần chục)?
	
	\begin{center}
		\begin{tikzpicture}[scale=1.2,font=\footnotesize,line join=round
			,line cap=round,>=stealth]
			% 1. Định nghĩa các đỉnh chính của hình chóp S.ABCD
			\coordinate (S) at (2.2,4.5);
			\coordinate (A) at (0,0);
			\coordinate (B) at (1.6,-1.3);
			\coordinate (C) at (6.5,-0.7);
			\coordinate (D) at (4.5,0.8);
			
			% 2. Xác định các điểm trên các cạnh
			\coordinate (O) at ($(A)!0.45!(S)$);
			\coordinate (V) at ($(A)!0.8!(S)$);
			
			\coordinate (M) at ($(B)!0.15!(S)$); % M trên SB
			\coordinate (N) at ($(C)!0.4!(S)$);  % N trên SC
			
			\coordinate (R_SB) at ($(B)!0.52!(S)$); % R duy nhất trên SB
			\coordinate (T) at ($(C)!0.7!(S)$);      % T trên SC
			
			% P nằm trên SD
			\coordinate (P) at ($(D)!0.42!(S)$);    
			
			% SỬA ĐIỂM U: U là giao của VT và SD (cạnh khuất)
			\coordinate (U) at ($(D)!0.7!(S)$);
			
			% 3. Vẽ các nét khuất (Dashed)
			\draw[dashed, gray!80] (A) -- (D) -- (C);
			\draw[dashed, gray!80] (S) -- (D);
			\draw[dashed, gray!80] (A) -- (D);
			\draw[dashed, gray!80] (O) -- (P); 
			\draw[dashed, gray!80] (P) -- (N);
			\draw[dashed, gray!80] (V) -- (U); % Đoạn VU khuất vì U nằm trên SD
			\draw[dashed, gray!80] (U) -- (T); % Đoạn UT khuất
			% 4. Vẽ các nét thấy (Solid)
			\draw(R_SB) -- (T); 
			\draw (S) -- (A) -- (B) -- (C) -- (S);
			\draw (A)--(M);
			\draw (S) -- (B);
			\draw (M) -- (N);
			\draw (O) -- (R_SB);
			% 5. Vẽ ký hiệu góc 15 độ tại đỉnh S
			\begin{scope}
				\path[clip] (S) -- (D) -- (R_SB) -- cycle;
				\draw (S) circle (0.7);
				\node at ($(S)+(-72:0.85)$) {\footnotesize $15^\circ$};
			\end{scope}
			
			% 6. Gán nhãn và vẽ các điểm
			\foreach \p/\pos/\txt in {
				S/above/S, A/left/A, B/below left/B, C/right/C, D/below/D,
				M/left/M, N/right/N, P/right/P, O/left/O, V/above left/V,
				T/right/T, R_SB/right/R%
			} {
				\fill (\p) circle (1pt);
				\node[\pos] at (\p) {\small $\txt$};
			}
			
			% Vẽ điểm U đặc biệt (màu xám, nằm trên SD)
			\fill (U) circle (1pt);
			\draw (U) circle (1pt);
			\node[below left] at (U) {\small $U$};
			
		\end{tikzpicture}
	\end{center}
	\shortans{$44{,}6$}
	\loigiai{\begin{center}\begin{tikzpicture}[scale=1,font=\footnotesize,line join=round
				,line cap=round,>=stealth]
				% 1. Định nghĩa đỉnh S ở trên cùng làm gốc
				\coordinate (S) at (0,3); % Đặt S cao lên
				
				% 2. Tinh chỉnh các góc và bán kính để tạo độ lệch ABCD, xòe rộng và hướng xuống dưới
				% Các góc được tính toán trong khoảng từ 190 đến 350 độ (nửa mặt phẳng dưới S)
				\foreach \i/\ang/\r in {
					1/195/4.8,  % A trái cùng
					2/215/4.5,  % B
					3/235/4.2,  % C
					4/255/4.0,  % D
					5/275/3.8,  % A
					6/295/3.6,  % B
					7/315/3.7,  % C
					8/335/4.1,  % D
					9/355/4.6    % A phải cùng
				} {
					\coordinate (P\i) at ($(S) + (\ang:\r)$);
				}
				
				% 3. Vẽ các tia từ S xuống
				\foreach \i in {1,...,9} {
					\draw (S) -- (P\i);
				}
				
				% 4. Vẽ đường biên ngoài (đáy quạt, không đều)
				\draw (P1) -- (P2) -- (P3) -- (P4) -- (P5) -- (P6) -- (P7) -- (P8) -- (P9);
				
				% 5. Xác định các điểm nút bên trong (M, N, P, Q, R, T, U, V)
				% Giữ nguyên tỉ lệ lệch để các điểm này trông tự nhiên
				\coordinate (M) at ($(S)!0.9!(P2)$);
				\coordinate (N) at ($(S)!0.85!(P3)$);
				\coordinate (P) at ($(S)!0.8!(P4)$);
				\coordinate (Q) at ($(S)!0.7!(P5)$);
				\coordinate (R) at ($(S)!0.9!(P6)$);
				\coordinate (T) at ($(S)!0.8!(P7)$);
				\coordinate (U) at ($(S)!0.85!(P8)$);
				\coordinate (V) at ($(S)!0.3!(P9)$); % V nằm trên tia A (thứ 5)
				
				% Vẽ đường gấp khúc thiết diện bên trong
				\draw[thick] (P1)--(M) -- (N) -- (P) -- (Q) -- (R) -- (T) -- (U);
				\draw[thick] (U) -- (V)--(P1); 
				
				% 6. Gán nhãn các điểm ngoài (A, B, C, D...)
				\node[left] at (P1) {$A$};
				\node[below left] at (P2) {$B$};
				\node[below] at (P3) {$C$};
				\node[below] at (P4) {$D$};
				\node[below] at (P5) {$A$};
				\node[below right] at (P6) {$B$};
				\node[right] at (P7) {$C$};
				\node[right] at (P8) {$D$};
				\node[right] at (P9) {$A$};
				\node[above] at (S) {$S$}; % S ở trên cùng
				
				% 7. Gán nhãn các điểm nút bên trong
				\node[below right=-2pt] at (M) {$M$};
				\node[below right=-2pt] at (N) {$N$};
				\node[below right=-1pt] at (P) {$P$};
				\node[right] at (Q) {$Q$};
				\node[right] at (R) {$R$};
				\node[left] at (T) {$T$};
				\node[above] at (U) {$U$};
				\node[above] at (V) {$V$};
				
				% 8. Vẽ chấm điểm
				\foreach \p in {S, P1, P2, P3, P4, P5, P6, P7, P8, P9, M, N, P, Q, R, T, U, V} {
					\fill (\p) circle (1.2pt);
				}
			\end{tikzpicture}
		\end{center}
		Vì dây đèn led quấn đủ $2$ vòng quanh tòa nhà trước khi kết thúc tại $V$ và đi về $S$. Khi đó, chúng ta phải trải phẳng $4$ mặt tam giác của hình chóp đều ($8$ mặt là $8$ tam giác cân tại $S$).\\
		Vì mỗi mặt bên là một tam giác cân có góc ở đỉnh $S$ là $15^\circ$ nên tổng góc ở đỉnh $S$ trên hình trải phẳng sẽ là $15^\circ \cdot 8 = 120^\circ$.\\
		Xét tam giác $SAV$ trên mặt phẳng trải hình (với góc $\widehat{ASV} = 120^\circ$, $SA = 40$ m, $SV = 3$ m)
		$$AV = \sqrt{SA^2 + SV^2 - 2 \cdot SA \cdot SV \cdot \cos 120^\circ} = \sqrt{40^2 + 3^2 - 2 \cdot 40 \cdot 3 \cdot (-0,5)} = \sqrt{1729} \approx 41{,}6 \text{ m}.$$
		Để độ dài đường gấp khúc $AMNPQRTUVS$ ngắn nhất thì $AV$ ngắn nhất. Khi đó tất cả các điểm $A$, $M$, $N$, $P$, $Q$, $R$, $T$, $U$, $V$ thẳng hàng.\\
		Do đó tổng chiều dài dây đèn ngắn nhất
		$L = 41{,}6 + 3 = 44{,}6$ m.
	}
\end{ex}

\begin{ex}%[1H8C6-1]%[Dự án C đợt 3 NH25-26- Phạm Đăng Đăng]%[TT-THPT-TranNhanTong-HaNoi-N25-26]%[GVPB - Phan Quốc Trí]
	Cho hình chóp $S.ABC$ có đáy là tam giác cân tại $A$. Cạnh bên $SA$ vuông góc với đáy. Gọi $H$ là trực tâm tam giác $ABC$. Tính $\cos\alpha$ với $\alpha$ là góc giữa đường thẳng $CH$ và mặt phẳng $(SBC)$. Biết $AB=AC=3$, $SA=BC=4$.
	\shortans{$0{,}81$}
	\loigiai{
		\begin{center}
			\begin{tikzpicture}[scale=1, line cap=round, line join=round, font=\small,>=stealth]
				\def\a{4} %Khai báo cạnh
				\def\h{4}
				\path 	(0:0) coordinate (A)
				++(0:\a) coordinate (C)
				++(-150:4*\a/5) coordinate (B)
				($(A)+(90:\h)$) coordinate (S)
				($(B)!0.5!(C)$) coordinate (M)
				($(A)!2/3!(B)$) coordinate (F)
				($(S)!5/7!(M)$) coordinate (K)
				;
				
				\coordinate (H) at (intersection of A--M and C--F);
				\draw[thick] 	(A)--(B)--(C)
				(A)--(S)	(C)--(S)	(B)--(S)--(M) (K)--(C);
				\draw[dashed,thick] (C)--(F)	(A)--(C) (A)--(M) (H)--(K);
				\foreach \x /\goc in {A/180,B/270,C/0,S/90,M/270,F/180,H/250,K/90}
				\fill[black] (\x) circle (1.5pt)
				($(\x)+(\goc:3mm)$) node {$\x$};
				;%Theo chiều dương
			\end{tikzpicture}
		\end{center}
		Gọi $M$ là trung điểm $BC$, $F$ là giao điểm của $CH$ và $AB$. Kẻ $HK$ vuông góc với $SM$ tại $K$.\\
		Ta có $\heva{&BC\perp AM\\&BC\perp SA}\Rightarrow BC \perp (SAM)\Rightarrow BC\perp HK$.\\Ta lại có $HK\perp SM$ nên ta suy ra $HK\perp (SBC)$. Do đó $\left(CH,(SBC)\right)=(CH,CK)=\widehat{HCK}=\alpha$.\\
		Ta có $AM=\sqrt{AB^2-BM^2}=\sqrt{5}$, $\tan\widehat{ BAM}=\dfrac{BM}{AM}=\dfrac{2}{\sqrt{5}}$.\\
		Ta có $\widehat {BAM}=\widehat{BCF}$ (cùng phụ $\widehat{ABC}$) nên $\tan \widehat{BCF}=\dfrac{FB}{FC}=\dfrac{HM}{MC}$. Suy ra $HM=\dfrac{4}{\sqrt{5}}$.\\
		$SM=\sqrt{SA^2+AM^2}=\sqrt{21}$.\\
		Lại có $\Delta SAM \sim \Delta HKM$, suy ra $\dfrac{SA}{HK}=\dfrac{MS}{MH}$, do đó $HK=\dfrac{SA\cdot MH}{MS}=\dfrac{16}{\sqrt{105}}$.\\
		$HC=\sqrt{HM^2+MC^2}=\dfrac{6}{\sqrt{5}}$ và $KC=\sqrt{HC^2-HK^2}=\dfrac{10}{\sqrt{21}}$. \\Do đó $\cos\alpha=\cos\widehat{ HCK}=\dfrac{KC}{HC}\approx 0{,}81$.}
\end{ex}

\begin{ex}%[1H8C6-6]%[Dự án đề TT Toán NH25-26-Dot 3- Hải Phụng]%[Cụm 13-Hải Phòng]%[GVPB - Viet Hoang Pham]
Cho hình chóp $S.ABC$ có đáy $ABC$ là tam giác vuông tại $A$, $AB=1$ cm, $AC=\sqrt{2}$ cm; $\widehat{SBA}=\widehat{SCA}=90^\circ$, góc giữa $BC$ và mặt phẳng $(SAB)$ bằng $45^\circ$. Tính khoảng cách giữa hai đường thẳng $SA$ và $BC$, với đơn vị là cm và kết quả làm tròn đến hàng phần trăm.
\shortans[oly]{$0{,}68$}
\loigiai{
\begin{center}
	\begin{tikzpicture}[line join=round, line cap=round, scale=0.9]
		%\draw[dashed,cyan] (-5,-5) grid (5,5);
		\def\d{5} \def\r{1.8} \def\h{4} \def\l{2.2}
		
		\coordinate[label={below left}:$B$] (B) at (-3,-3);
		\coordinate[label={below right}:$A$] (A) at ($(B)+(\d,0)$);
		\coordinate[label={above left}:$H$] (H) at ($(B)+(\l,\r)$);
		\coordinate[label={above right}:$C$] (C) at ($(H)+(\d,0)$);
		\coordinate[label={above}:$S$] (S) at ($(H)+(0,\h)$);
		\coordinate[label={above left}:{$K$}] (K) at ($(S)!0.6!(B)$);
		\coordinate[label={right}:{$O$}] (O) at ($(H)!0.5!(A)$);
		\coordinate[label={above right}:{$I$}] (I) at ($(S)!0.5!(H)$);
		\coordinate[label={below left}:{$M$}] (M) at ($(B)!0.8!(O)$);
		\coordinate[label={right}:{$N$}] (N) at ($(I)!0.5!(M)$);
		\draw (S)--(B)--(A)--(C)--(S)--(A);
		\draw[dashed] (S)--(H)--(B) (H)--(C)--(B) (H)--(A) (H)--(K) (O)--(I)--(C) (B)--(I) (H)--(M)--(I) (H)--(N);
		\draw 
		pic[draw, angle radius=2mm]{right angle=S--B--A}
		pic[draw, angle radius=2mm]{right angle=S--C--A}
		pic[draw, angle radius=2mm]{right angle=S--H--B}
		pic[draw, angle radius=2mm]{right angle=B--A--C}
		pic[draw, angle radius=2mm]{right angle=H--K--S}
		pic[draw, angle radius=2mm]{right angle=H--M--B}
		pic[draw, angle radius=2mm]{right angle=H--N--I}
		; 
		
		\foreach \x in {A,B,C,H,K,S,O,I,M,N} 
		\draw[fill=black] (\x) circle (1.5pt); 
		
	\end{tikzpicture}
\end{center}
Kẻ $SH\perp (ABC)$. Khi đó $AB\perp SH$ và $AB\perp SB$ nên $AB\perp (SBH)$.\\
Tương tự $AC\perp (SCH)$.\\
Do đó $AB\perp BH$ và $AC\perp CH$. Nghĩa là $ACHB$ là hình chữ nhật.\\
Vì góc giữa $BC$ và mặt phẳng $(SAB)$ bằng $45^\circ$ nên
$$\sin 45^\circ = \dfrac{\mathrm{d}\left[ C,(SAB)\right]}{BC}\Rightarrow \dfrac{\sqrt{2}}{2} = \dfrac{\mathrm{d}\left[ C,(SAB)\right]}{\sqrt{3}} \Rightarrow \mathrm{d}\left[ C,(SAB)\right] = \dfrac{\sqrt{6}}{2}. $$
Vì $CH\parallel (SAB)$ nên $\mathrm{d}\left[ H,(SAB)\right] =\mathrm{d}\left[ C,(SAB)\right] = \dfrac{\sqrt{6}}{2}$.\\
Kẻ $HK\perp SB$ tại $K$. Khi đó $HK\perp SB$ và $HK\perp AB$ nên $HK\perp (SAB)$. Suy ra $\mathrm{d}\left[ H,(SAB)\right] = HK =\dfrac{\sqrt{6}}{2}$.\\
Mà $HK=\dfrac{HS\cdot HB}{\sqrt{SH^2 + HB^2}}\Rightarrow \dfrac{\sqrt{6}}{2} = \dfrac{HS\cdot \sqrt{2}}{\sqrt{SH^2+2}}\Rightarrow SH=\sqrt{6}$.
\\
Gọi $I$ là trung điểm $SH$ và $O=HA\cap BC$. Khi đó $(BIC)\parallel SA$.\\
Suy ra $\mathrm{d}(SA,BC) = \mathrm{d}(SA,(BIC)) = \mathrm{d}(A,(BIC)) = \mathrm{d}(H,(BIC))$. \hfill $(1)$\\
Kẻ $HM\perp BC$ tại $M$, $HN\perp IM$ tại $N$. Khi đó $\mathrm{d}(H,(BIC)) = HN = \dfrac{HM\cdot HI}{\sqrt{HM^2+HI^2}}$.\\
Mà $HM=\dfrac{HC\cdot HB}{\sqrt{HC^2+HB^2}} = \dfrac{1\cdot \sqrt{2}}{\sqrt{1^2+\sqrt{2}^2}} =\dfrac{\sqrt{6}}{3} $ và $HI=\dfrac{SH}{2} = \dfrac{\sqrt{6}}{2}$.\\
Suy ra $HN= \dfrac{\sqrt{78}}{13}$. \hfill $(2)$\\
Từ $(1)$ và $(2)$ suy ra $\mathrm{d}(SA,BC) = \dfrac{\sqrt{78}}{13}\approx 0{,}68$.
}
\end{ex}

\begin{ex}%[Dự án đề TT Toán NH25-26-Dot 4- Hải Phụng]%[THCS-THPT LÊ THÁNH TÔNG - TP.HCM]%[GVPB - Viet Hoang Pham]%[1H8C7-9]
	\immini[thm]{
		Một tấm pin năng lượng mặt trời hình chữ nhật $ABCD$ có kích thước $AB = 2$ m và $BC = 3$ m. Tấm pin được đặt nghiêng sao cho cạnh $AB$ nằm sát trên mặt đất phẳng. Một bóng đèn (xem như một điểm) được đặt tại vị trí $S$ cao $4$ m có hình chiếu vuông góc lên mặt đất trùng với trung điểm $I$ của cạnh $AB$. Vào buổi tối khi bật đèn lên bóng của tấm pin trên mặt đất tạo thành một hình thang cân $ABC'D'$ (với $C', D'$ lần lượt là bóng của $CD$). Biết rằng hình thang cân $ABC'D'$ có chiều cao bằng $4{,}5$ m. Hãy tính diện tích lớn nhất của hình thang cân $ABC'D'$ là bóng của tấm pin trên mặt đất? (Đơn vị: m$^2$, làm tròn kết quả đến hàng phần chục).
	}{
		\begin{tikzpicture}[line join=round, line cap=round, scale=1]
			\def\d{5} \def\r{1.8} \def\h{4} \def\l{1}
			
			\coordinate[label={below left}:$B$] (B) at (-3,-3);
			\coordinate[label={below right}:$C'$] (C') at ($(B)+(\d,-0.3)$);
			\coordinate[label={below left}:$A$] (A) at ($(B)+(-\l,\r)$);
			\coordinate[label={above right}:$D'$] (D') at ($(A)+(\d+0.5,0.7)$);
			\coordinate[label={below left}:$I$] (I) at ($(A)!0.5!(B)$);
			\coordinate[label={right,xshift=1mm, yshift=2mm}:$S$] (S) at ($(I)+(0,\h)$);
			\coordinate[label={above right}:{$C$}] (C) at ($(S)!0.7!(C')$);
			\coordinate[label={above right}:{$D$}] (D) at ($(S)!0.7!(D')$);
			
			\draw (S)--(B)--(C')--(D')--(S)--(C') (S)--(A)--(B) (S)--(I) (B)--(C)--(D);
			\draw[dashed] (A)--(D') (D)--(A);
			\draw 
			pic[draw, angle radius=2mm]{right angle=S--I--B}
			; 
			\draw (S)  node[above,yshift=-1mm]{\Large \faLightbulbO};
			\draw[draw=none, fill=gray, opacity=0.3] (A)--(B)--(C)--(D)--cycle;
			
		\end{tikzpicture}
	}
	\shortans{$15{,}8$}
	\loigiai{
		\immini[thm]{
			Gọi $H$, $K$ lần lượt là trung điểm của $CD$ và $C'D'$.\\
			Vì $ABC'D'$ là hình thang cân, mà $C'D'>CD=AB$ nên đáy hình thang này bắt buộc phải là $AB$ và $C'D'$ (vì hai cạnh bên hình thang cân bằng nhau). Từ đó suy ra $IK$ là đường cao hình thang $ABC'D'$ nên  $IK=4{,}5$.\\
			Diện tích hình thang $ABC'D'$ là
			$$S=\dfrac{1}{2}\cdot (AB+C'D')\cdot IK = \dfrac{1}{2}\cdot (2+C'D')\cdot 4{,}5.$$
			Như vậy $S_{\max}$ khi và chỉ khi $C'D'$ lớn nhất. \\
			Dễ thấy $AB$ vuông góc với các đường thẳng $SI$, $IH$ và $IK$. \\
			Ta có $SK=\sqrt{SI^2+IK^2} = \dfrac{\sqrt{145}}{2}$. Khi đó
		}{%Bổ sung gói vẽ bóng đèn \usepackage{fontawesome5}
			\begin{tikzpicture}[line join=round, line cap=round, scale=1]
				\def\d{5} \def\r{1.8} \def\h{4} \def\l{1}
				
				\coordinate[label={below left}:$B$] (B) at (-3,-3);
				\coordinate[label={below right}:$C'$] (C') at ($(B)+(\d,-0.3)$);
				\coordinate[label={below left}:$A$] (A) at ($(B)+(-\l,\r)$);
				\coordinate[label={above right}:$D'$] (D') at ($(A)+(\d+0.5,0.7)$);
				\coordinate[label={below left}:$I$] (I) at ($(A)!0.5!(B)$);
				\coordinate[label={right,xshift=1mm, yshift=2mm}:$S$] (S) at ($(I)+(0,\h)$);
				\coordinate[label={above right}:{$C$}] (C) at ($(S)!0.7!(C')$);
				\coordinate[label={above right}:{$D$}] (D) at ($(S)!0.7!(D')$);
				\coordinate[label={right}:$H$] (H) at ($(C)!0.5!(D)$);
				\coordinate[label={right}:$K$] (K) at ($(C')!0.5!(D')$);
				
				\draw (S)--(B)--(C')--(D')--(S)--(C') (S)--(A)--(B) (S)--(I) (B)--(C)--(D) (S)--(K);
				\draw[dashed] (A)--(D') (D)--(A) (K)--(I)--(H);
				\draw 
				pic[draw, angle radius=2mm]{right angle=S--I--B}
				; 
				\draw 
				pic[draw, angle radius=2mm]{right angle=B--I--K}
				; 
				\draw 
				pic[draw, angle radius=2mm]{right angle=C'--K--S}
				; 
				\draw 
				pic[draw, angle radius=2mm]{right angle=C'--K--I}
				; 
				\draw (S)  node[above,yshift=-1mm]{\Large \faLightbulbO};
				\draw[draw=none, fill=gray, opacity=0.3] (A)--(B)--(C)--(D)--cycle;
			\end{tikzpicture}
		}
		\begin{eqnarray*}
			\dfrac{CD}{C'D'}=\dfrac{SH}{SK} & \Rightarrow& \dfrac{2}{C'D'}=\dfrac{2\cdot SH}{\sqrt{145}}\\
			& \Rightarrow& C'D'=\dfrac{\sqrt{145}}{SH}.
		\end{eqnarray*}
		Từ đó, để $C'D'$ lớn nhất thì $SH$ nhỏ nhất.\\
		Ta có
		{\allowdisplaybreaks
			\begin{eqnarray*}
				IH^2=SI^2+SH^2-2SI\cdot SH\cdot \cos \widehat{ISH} &\Leftrightarrow& 3^2=4^2+SH^2 -2\cdot 4\cdot SH\cdot \dfrac{8}{\sqrt{145}}\\
				&\Leftrightarrow& SH=2{,}4.
			\end{eqnarray*}
		}
		\noindent Suy ra $C'D'=5$ (chọn $SH$ nhỏ nhất để $C'D'$ lớn nhất).\\
		Vậy $S_{\max} = \dfrac{1}{2}\cdot (2+5)\cdot 4{,}5 = 15{,}75\approx 15{,}8$.
	}
\end{ex}

\begin{ex}%[1H8C7-9]%[Dự án C đề thi thử TN Môn Toán NH25-26 - Dot 4 - Phạm Hoàng Đăng]%[THPT Nguyễn Trãi-Hà Nội]%[GVPB - Khắc Thiên]
	\immini[thm]{Cho một tòa nhà đồ chơi có dạng hình hộp chữ nhật $ABCD.A'B'C'D'$ với đáy là hình vuông cạnh $20$ (cm) và chiều cao $80$ (cm), $A'B'C'D'$ là mặt dưới, $ABCD$ là mặt trên. Một con kiến xuất phát từ điểm $A'$ và bò lên bề mặt xung quanh của tòa nhà để đến đích là điểm $A$ (hình vẽ minh họa). Đường đi của con kiến là một đường gấp khúc liên tục, nằm hoàn toàn trên bốn mặt bên (không bò trên mặt trên và mặt dưới), không trùng với các cạnh của hình hộp và lần lượt gặp các cạnh bên theo đúng thứ tự $BB'$, $CC'$, $BB'$, $CC'$, $DD'$. Biết rằng con kiến chạm vào cạnh $BB'$ ở các điểm cách $B'$ một khoảng không nhỏ hơn $16$ (cm) và chạm vào cạnh $CC'$ ở các điểm cách $C$ một khoảng không lớn hơn $16$ (cm). Tính độ dài ngắn nhất của quãng đường mà con kiến đi (\textit{làm tròn kết quả đến hàng đơn vị, đơn vị độ dài là cm}).}{
		\begin{tikzpicture}[line cap=round,line join=round,scale=.7,font=\footnotesize]
			\path
			(0,0) coordinate (A)
			(4,0) coordinate (B)
			(2,2) coordinate (D)
			($(B)+(D)-(A)$) coordinate (C)
			($(A)!-6cm!90:(B)$) coordinate (A')
			($(A')+(D)-(A)$) coordinate (D')
			($(A')+(B)-(A)$) coordinate (B')
			($(B')+(D')-(A')$) coordinate (C');
			\draw (A)--(B)--(C)--(D)--cycle (A')--(B')--(C')
			(B)--(B') (C)--(C')(A)--(A')
			(A')--($(B')!1cm!(B)$)--($(C')!2cm!(C)$)--($(B')!3cm!(B)$)--($(C')!4cm!(C)$);
			\draw[dashed] (C')--(D')--(A') (D)--(D') ($(C')!4cm!(C)$)--($(D')!5cm!(D)$)--(A)
			;
			\foreach \i/\g in {A'/-150,B'/-40,C'/0,D'/140,A/180,B/-140,C/45,D/90}
			\draw[fill=black] (\i) circle (1pt) +(\g:.4) node{$\i$};
	\end{tikzpicture}}
	\shortans[oly]{157}
	\loigiai{
		Trải phẳng các mặt bên của hình hộp chữ nhật theo đường đi của con kiến ta được hình phẳng như hình vẽ.
		\begin{center}
			\begin{tikzpicture}[scale=1,>=stealth, line join=round, line cap=round, smooth]
				\path (0,0) coordinate (Ap1)
				(1,0) coordinate (Bp1)
				(2,0) coordinate (Cp1)
				(3,0) coordinate (Bp2)
				(4,0) coordinate (Cp2)
				(5,0) coordinate (Dp1)
				(6,0) coordinate (Ap2)
				(0,4) coordinate (A1)
				(1,4) coordinate (B1)
				(2,4) coordinate (C1)
				(3,4) coordinate (B2)
				(4,4) coordinate (C2)
				(5,4) coordinate (D1)
				(6,4) coordinate (A2)
				(2, 3.25) coordinate (H)
				(1, 1.625) coordinate (K)
				(1, 0.666) coordinate (E)
				(0,0) coordinate (O);
				\foreach \x in {1,2,3,4,5} \draw (\x,0) -- (\x,4);
				\draw (Ap1) -- (H) -- (A2)
				(Ap1) -- (A2)
				(A1) -- (A2) (Ap1) -- (Ap2) (0,0) -- (0,4) (6,0) -- (6,4);
				\foreach \x/\y/\z in {A1/135/A, B1/90/B, C1/90/C, B2/90/B, C2/90/C, D1/90/D, A2/45/A, Ap1/-140/A', Bp1/-90/B', Cp1/-90/C', Bp2/-90/B', Cp2/-90/C', Dp1/-90/D', Ap2/-40/A', H/160/H, K/135/K, E/135/E}
				\path[fill=black, draw] (\x) circle (.035)+(\y:.3) node{$\z$};
			\end{tikzpicture}
		\end{center}
		Xét trường hợp con kiến đi thẳng từ $A'$ đến $A$.\\ Gọi $E$ là điểm con kiến chạm vào cạnh $BB'$.\\
		Khi đó
		\[\dfrac{EB'}{AA'}=\dfrac{20}{120} \Rightarrow EB'=\dfrac{80\cdot 20}{120} \approx 13{,}3 \text{ (cm)}<16 \text{ (cm)} \text{ (không thỏa mãn).}\]
		Xét trường hợp con kiến chạm vào cạnh $CC'$ tại $H$ (như hình vẽ), với $0<CH \leq 16$.\\
		Khi đó $A'H$ đạt giá trị nhỏ nhất khi $CH=16$.\\
		Suy ra $C'H=64$ và $A'H=\sqrt{40^2+64^2}=8 \sqrt{89}$.\\
		Gọi $K$ là điểm con kiếm chạm cạnh $BB'$ trong trường hợp này.\\
		Ta có $\dfrac{KB'}{HC'}=\dfrac{20}{40} \Rightarrow KB'=\dfrac{64\cdot 20}{40}=32 \text{ (cm)}>16 \text{ (cm)}$ (thỏa mãn yêu cầu).\\
		Con kiến tiếp tục đi thẳng từ $H$ đến $A$, $AH=\sqrt{16^2+80^2}=16\sqrt{26}$.\\
		Vậy độ dài ngắn nhất của quãng đường mà con kiến đi là
		\[A'H+HA=8\sqrt{89}+16\sqrt{26} \approx 157{,}056 \text{ (cm)}.\]
	}
\end{ex}

\begin{ex}%[Đề thi thử môn Toán TRƯỜNG THPT Cổ Loa - Hà Nội năm học 2025-2026]%[Phạm Quốc Toàn - Phạm Hoàng Việt]%[1H8H5-4]
	Cho hình chóp $S\cdot ABCD$ có đáy $ABCD$ là hình vuông tâm $O$, cạnh bằng $1$. Cạnh bên $SA$ vuông góc với đáy và góc $\widehat{SBD}=60^{\circ}$. Khoảng cách giữa hai đường thẳng $AB$ và $SO$ bằng bao nhiêu? (không làm tròn kết quả các phép tính trung gian, chỉ làm tròn kết quả cuối cùng đến hàng phần trăm)
	\par\shortans[0]{0{,}45}
	\loigiai{
		\begin{center}
			\begin{tikzpicture}[font=\footnotesize, line join=round, line cap=round, >=stealth,scale=.7]
				\coordinate (S) at (3,8);
				\coordinate (A) at (3,3);
				\coordinate (B) at (0,0);
				\coordinate (C) at (6,0);
				\coordinate (D) at (9,3);
				\coordinate (O) at (4.5,1.5);
				\coordinate (E) at (6,3);
				\coordinate (H) at (5.1,4.5);
				%%%%%%%%%%%%%%%%%%%%%%%%%%%%%%%%
				\draw (S) -- (B) -- (C) -- (D) -- (S) -- (C);
				\draw[dashed] (S) -- (A) -- (B) -- (D) -- (A) (A) -- (C) (O) -- (E) -- (S) (A) -- (H);
				\foreach \i/\g in {S/90,A/180,B/180,C/0, D/0, O/-90, E/90, H/0}{\draw(\i) circle (1pt) ($(\i)+(\g:3mm)$) node[scale=1]{$\i$};}
				\pic[draw,thin,angle radius=2mm] {right angle=S--A--D}
				pic[draw,thin,angle radius=2mm] {right angle=A--H--E};
				%%%%%%%%%%%%%%%%%%%%%%%%%%%%%%%%
			\end{tikzpicture}
		\end{center}
		Vì $ABCD$ là hình vuông có cạnh bằng $1$ nên $BD=\sqrt{2}$. \\ 
		$SA \perp (ABCD) \Rightarrow \heva{& SA \perp AB \\ & SA \perp AD} \Rightarrow \Delta SAB=\Delta SAD \Rightarrow SB=SD$. \\ 
		Tam giác $SBD$ có $SB=SD$ và $\widehat{SBD}=60^{\circ} \Rightarrow \Delta SBD$ đều $\Rightarrow SB=BD=\sqrt{2}$. \\ 
		Tam giác $SAB$ có $\widehat{A}=90^{\circ} \Rightarrow SA=\sqrt{SB^2-AB^2}=\sqrt{\left(\sqrt{2}\right)^2-1^2}=1$. \\ 
		Gọi $E$ là trung điểm của $AD \Rightarrow OE$ song song $AB$, $OE \perp AD$, $AE=\dfrac{1}{2}$. \\
		Lại có $SA\perp OE \Rightarrow OE \perp (SAD)$. \\ 
		Kẻ $AH \perp SE \Rightarrow OE\perp AH\Rightarrow AH\perp (SEO)$. \\
		Do $AB$ song song $OE \Rightarrow AB$ song song $(SEO)$. \\
		Từ đó, ta có $\mathrm{d} \left(AB, SO\right)=\mathrm{d} \big(AB, (SEO)\big)=\mathrm{d}\big(A, (SEO)\big)=AH$. \\ 
		Ta có $AH=\dfrac{SA\cdot AE}{\sqrt{SA^2+AE^2}}=\dfrac{1\cdot \dfrac{1}{2}}{\sqrt{1^2+\left(\dfrac{1}{2}\right)^2}}=\dfrac{1}{\sqrt{5}}\approx 0{,}45$.
	}
\end{ex}

\begin{ex}%[1H8V5-3]%[Dự án C - Đợt 3 - Sở Giáo Dục Sơn La]%[Lê Hùng Thắng - Nguyễn Tiến]% 
	Cho hình chóp $S.ABCD$ có đáy $ABCD$ là hình thoi tâm $O$, cạnh bằng $1$, góc $\widehat{ABC} = 60^{\circ}$. Biết rằng $SO \perp (ABCD)$, $SO = \dfrac{3}{2}$. Tính khoảng cách từ $A$ đến mặt phẳng $(SCD)$ (không làm tròn kết quả ở các bước trung gian, làm tròn kết quả ở bước cuối cùng đến chữ số thập phân hàng phần trăm).
	\shortans{0,83}
	\loigiai{
		\immini[thm]{
			Ta có $AO \cap(SCD)=C$ nên
			\begin{eqnarray*}
				&& \dfrac{\mathrm{d}(A, (SCD))}{\mathrm{d}(O, (SCD))} = \dfrac{AC}{OC} = 2  \\
				&\Rightarrow & \mathrm{d}(A, (SCD)) = 2\mathrm{d}(O, (SCD)).
			\end{eqnarray*}
			Tam giác $ACD$ cân tại $D$ có $\widehat{ADC}=60^{\circ}$ nên $ACD$ là tam giác đều cạnh bằng $1$.\\
			Gọi $M$ là trung điểm $CD$ và $H$ là trung điểm của $CM$, khi đó\\
			$AM \perp CD \Rightarrow OH \perp CD$ tại $H$  (vì $OH \parallel AM$).\\
			Kẻ $OK \perp SH$ tại $K$, suy ra $OK \perp (SCD)$.\\
			Khi đó $\mathrm{d}(O, (SCD)) = OK$.
			
		}
		{\begin{tikzpicture}[scale=1, font=\footnotesize, line join=round, line cap=round, >=stealth]
				\path (0,0) coordinate (A)
				(-1.5,-1.5) coordinate (B)
				(2.5,-1.5) coordinate (C) 
				($(A)+(C)-(B)$) coordinate (D)
				($(A)!0.5!(C)$) coordinate (O)
				($(O)+(0,4)$) coordinate (S)
				($(C)!0.5!(D)$) coordinate (M)
				($(C)!0.25!(D)$) coordinate (H)
				($(S)!0.65!(H)$) coordinate (K);
				\draw (H)--(S)--(B)--(C)--(D)--(S)--(C);
				\draw[dashed] (S)--(A)--(B) (A)--(D) (A)--(C) (B)--(D) (S)--(O) (O)--(H) (O)--(K) (A)--(M);
				\foreach \p/\r in {S/90,A/135,B/-135,M/0,C/-45,D/45,O/-90,H/-10,K/30}
				\fill (\p) circle (1.5pt) node[shift={(\r:3mm)}]{$\p$};
			\end{tikzpicture}
		}
		\noindent
		$AM$ là đường cao của tam giác đều cạnh bằng $1$ nên đường cao $AM = \dfrac{\sqrt{3}}{2}$.	\\
		Ta có $OH=\dfrac{1}{2}AM = \dfrac{\sqrt{3}}{4}$.\\
		Xét tam giác $SOH$ vuông tại $O$ có $OK$ là đường cao nên ta có
		\begin{eqnarray*}
			& & \dfrac{1}{OK^2} = \dfrac{1}{SO^2} + \dfrac{1}{OH^2} =\dfrac{1}{\left(\dfrac{3}{2}\right)^2} + \dfrac{1}{\left(\dfrac{\sqrt{3}}{4}\right)^2}= \dfrac{4}{9} + \dfrac{16}{3} = \dfrac{52}{9} \\
			&\Rightarrow & OK = \sqrt{\dfrac{9}{52}} = \dfrac{3}{2\sqrt{13}}.
		\end{eqnarray*}
		Vậy $\mathrm{d}(A, (SCD)) = 2 \cdot \dfrac{3}{2\sqrt{13}} = \dfrac{3}{\sqrt{13}} \approx 0,83$.
	}
\end{ex}

\begin{ex}%[Đề thi thử môn Toán Sở giáo dục tỉnh Tuyên Quang năm học 2025-2026]%[Đoàn Minh Tâm]%[1H8V5-3]%[GVPB: Thanh Phát]
	Cho hình lăng trụ đứng $ABC.A'B'C'$ có đáy $ABC$ là tam giác vuông tại $A$, $BC=2\sqrt{3}$, $AB=3$. Khoảng cách giữa đường thẳng $AA'$ và mặt phẳng $(BCC'B')$ bằng bao nhiêu?
	\shortans[oly]{1,5}
	\loigiai{
		\immini[thm]{
			Vì $ABC.A'B'C'$ là hình lăng trụ đứng nên cạnh bên $AA'$ song song với cạnh bên $BB'$.\\
			Suy ra $AA' \parallel (BCC'B')$.\\
			Do đó $\mathrm{d}\left(AA',(BCC'B')\right)=\mathrm{d}\left(A,(BCC'B')\right)$.\\
			Trong mặt phẳng $(ABC)$, kẻ $AH \perp BC$ tại $H$.\\
			Ta có lăng trụ đứng nên $BB' \perp (ABC) \Rightarrow BB' \perp AH$.\\
			Từ $\heva{&AH \perp BC \\ &AH \perp BB'}\Rightarrow AH \perp (BCC'B')$.\\
			Vậy khoảng cách cần tìm chính là độ dài đoạn thẳng $AH$.\\
			Xét tam giác $ABC$ vuông tại $A$, áp dụng định lí Pythagore, ta có
			\[ AC = \sqrt{BC^2 - AB^2} = \sqrt{\left(2\sqrt{3}\right)^2 - 3^2} = \sqrt{12 - 9} = \sqrt{3}. \]}{
			\begin{tikzpicture}[line join=round,line cap=round,line width=.6pt,font=\footnotesize,scale=1,>=stealth]
				\coordinate[label=left:$A$] (A) at (0,0);
				\coordinate[label=below left:$B$] (B) at (1,-1);
				\coordinate[label=right:$C$] (C) at (4,0);
				\coordinate[label=left:$A'$] (A1) at ($(A)+(90:4)$);
				\coordinate[label=below left:$B'$] (B1) at ($(B)-(A)+(A1)$);
				\coordinate[label=right:$C'$] (C1) at ($(C)-(A)+(A1)$);
				\coordinate[label=below right:$H$] (H) at ($(B)!1/3!(C)$);
				\draw (A1)--(A)--(B)--(C)--(C1)--(A1)--(B1)--(C1) (B)--(B1);
				\draw[dashed] (A)--(C) (A)--(H);
				\draw ($(H)!5pt!(A)$)--($(H)!5pt!(A)+(H)!5pt!(B)-(H)$)--($(H)!5pt!(B)$)--(H)--cycle;
				\fill (A)circle(1pt) (B)circle(1pt) (C)circle(1pt) (A1)circle(1pt) (B1)circle(1pt) (C1)circle(1pt) (H)circle(1pt);
			\end{tikzpicture}
		}
		\noindent Áp dụng hệ thức lượng trong tam giác vuông $ABC$ với đường cao $AH$ là
		\[ AH \cdot BC = AB \cdot AC \Rightarrow AH = \dfrac{AB \cdot AC}{BC} = \dfrac{3 \cdot \sqrt{3}}{2\sqrt{3}} = \dfrac{3}{2} = 1{,}5. \]
		Vậy $\mathrm{d}\left(AA',(BCC'B')\right)=1{,}5$.
	}
\end{ex}

\begin{ex}%[1H8V5-4]%[Dự án đề thi thử NH25-26]%[THPT Chuyên Phan Bội Châu - Nghệ An]%[Phạm Hoàng Đăng]
	\immini[thm]{Cho hình lăng trụ đều $ABC.A'B'C'$. Biết $AA'=6$, góc giữa đường thẳng $A'B$ và mặt phẳng $(ABC)$ bằng $60^\circ$. Tính khoảng cách giữa hai đường thẳng $AA'$ và $BC$.}{
		\begin{tikzpicture}[declare function={a=3; b=0.55*a; h=0.9*a;}]
			\path (0:0) coordinate (A)
			(-43:b) coordinate (B)
			(a,0) coordinate (C)
			\foreach \x in {A,B,C}{(\x)++(90:h) coordinate (\x1)};
			\draw[dashed] (A)--(C);
			\draw (B1)--(A1)--(C1)--cycle
			(A1)--(A)--(B)--(B1)
			(C1)--(C)--(B);
			\foreach \x/\goc/\t in {A/180/A,B/-90/B,C/0/C,A1/100/A',B1/85/B',C1/80/C'}
			\fill (\x) circle (1pt) node[shift={(\goc:9pt)},font=\scriptsize]{$\t$};
		\end{tikzpicture}
	}
	\shortans{$3$}
	\loigiai{
		Vì $ABC.A'B'C'$ là hình lăng trụ đều nên cạnh bên $AA' \perp (ABC)$.\\
		Hình chiếu vuông góc của đường thẳng $A'B$ lên mặt phẳng $(ABC)$ là $AB$.\\
		Do đó, góc giữa $A'B$ và mặt phẳng $(ABC)$ là $\left(A'B, AB\right) = \widehat{A'BA} = 60^\circ$.\\
		\begin{center}
			\begin{tikzpicture}[declare function={a=3; b=0.55*a; h=0.9*a;}]
				\path (0:0) coordinate (A)
				(-43:b) coordinate (B)
				(a,0) coordinate (C)
				\foreach \x in {A,B,C}{(\x)++(90:h) coordinate (\x1)};
				\path (B) -- (C) coordinate[midway] (H);
				\draw[dashed] (A)--(C)
				(A)--(H);
				\draw (B1)--(A1)--(C1)--cycle
				(A1)--(A)--(B)--(B1)
				(C1)--(C)--(B)
				(A1)--(B);
				\foreach \x/\y/\z in {A/H/C}{
					\path (\y) pic[draw,angle radius=5pt]{right angle = \x--\y--\z};
				}
				\foreach \x/\goc/\t in {A/180/A,B/-90/B,C/0/C,A1/100/A',B1/85/B',C1/80/C',H/-45/H}
				\fill (\x) circle (1pt) node[shift={(\goc:9pt)},font=\scriptsize]{$\t$};
			\end{tikzpicture}
		\end{center}
		Xét tam giác vuông $A'AB$ tại $A$, ta có $AB = \dfrac{AA'}{\tan 60^\circ} = \dfrac{6}{\sqrt{3}} = 2\sqrt{3}$.\\
		Vì $AA' \parallel BB'$ nên $AA' \parallel (BCC'B')$.\\ 
		Khoảng cách cần tìm là
		\[\mathrm{d}(AA', BC) = \mathrm{d}\big(AA', (BCC'B')\big) = \mathrm{d}\big(A, (BCC'B')\big).\]
		Gọi $H$ là trung điểm của $BC$, vì tam giác đáy $ABC$ là tam giác đều cạnh $2\sqrt{3}$ nên $AH \perp BC$.\\
		Mặt khác, $AH \perp BB'$ (do $BB' \perp (ABC)$).\\ Suy ra $AH \perp (BCC'B')$.\\
		Vậy khoảng cách $\mathrm{d}\big(A, (BCC'B')\big) = AH = AB \cdot \dfrac{\sqrt{3}}{2} = 2\sqrt{3} \cdot \dfrac{\sqrt{3}}{2} = 3$.
	}
\end{ex}

\begin{ex}%[1H8V5-4]%[Dự án C đề thi thử TN Môn Toán NH25-26 - Dot 4 - Vũ Hồng Toàn]%[THPT Mai Anh Tuấn - Thanh Hóa]%[GVPB - Lê Hữu Kiệt]
	Cho hình lăng trụ đứng $ABC.A'B'C'$ có $AB = 5$, $BC = 5$, $CA = 7$. Khoảng cách giữa hai đường thẳng $AA'$ và $BC$ bằng bao nhiêu {\it (làm tròn kết quả đến hàng phần chục)}?
	\par\shortans{5,0}
	\loigiai{
		\begin{center}
			\begin{tikzpicture}[scale=1, font=\footnotesize, line join=round, line cap=round, >=stealth]
				\tikzset{declare function={a=3.5;b=0.65*a;goc=-40;h=0.75*a;}}
				\path (0:0) coordinate (A)
				(a,0) coordinate (C)
				(goc:b) coordinate (B)
				\foreach \p in {A,B,C} {($(\p)+(90:h)$) coordinate (\p')}
				($(C)!0.5!(B)$)coordinate (H);
				\draw[dashed]  (H)--(A)--(C);
				\draw (B')--(B)--(A)--(A')--cycle--(C')--(C)--(B) (A')--(C');
				\foreach \x /\gN in {A/180,B/-90,C/-20,A'/90,B'/80,C'/90,H/0}{
					\fill (\x) circle (1pt)($(\x)+(\gN:2.5mm)$) node {$\x$};
				}
			\end{tikzpicture}
		\end{center}
		Trong tam giác $ABC$, dựng $AH \perp BC$ với $H \in BC$.\\
		Khi đó $\heva{AH \perp BC\\ AH \perp BB'} \Rightarrow AH \perp (BCC'B')$.\\
		Ta có $p = \dfrac{a+b+c}{2} = \dfrac{5+5+7}{2} = \dfrac{17}{2}$.\\
		$S_{ABC} = \sqrt{\dfrac{17}{2}\left(\dfrac{17}{2}-5\right)\left(\dfrac{17}{2}-5\right)\left(\dfrac{17}{2}-7\right)} = \dfrac{7\sqrt{51}}{4}$.\\
		Mặt khác $S_{ABC} = \dfrac{1}{2}\cdot AH\cdot BC \Rightarrow AH = \dfrac{2S_{ABC}}{BC} = \dfrac{2 \cdot \dfrac{7\sqrt{51}}{4}}{5} \approx 5,0$.\\
		Vậy $\mathrm{d}(AA', BC) = \mathrm{d}\left(A,(BCC'B')\right) = AH \approx 5{,}0$.
	}
\end{ex}

\begin{ex}%[1H8V6-2]%[TT Cụm 5 - Ninh Bình-NH25-26]%[GVSB:Dương Công Tạo]%[GVPB1: Nguyễn Tấn Tài]
	\immini[thm]{Một mái nhà được tạo bởi hai nửa lục giác đều $ABCD$, $ABC'D'$ và hai tam giác bằng nhau $ADD'$, $BCC'$. Biết $CDD'C'$ là hình chữ nhật và $AB\parallel CD\parallel C'D'$, $CD=C'D'=2AB=6$ m, $DD'=4$ m. Tìm số đo góc nhị diện $[D', AD, C]$ (\textit{làm tròn kết quả đến hàng đơn vị của độ}).}
	{\begin{tikzpicture}[scale=1,line join=round, line cap=round,font=\footnotesize]
			\def\a{4}
			\path 	(0:0) coordinate (D)
			(20:\a) coordinate (C)
			(135:.55*\a) coordinate (D')
			(75:.32*\a) coordinate (A)
			($(D')-(D)+(C)$) coordinate (C')
			($(C)-(D)+(A)$) coordinate (At)
			($(A)!.42!(At)$) coordinate (B);
			\draw	(C)--(D)--(D')--(C')--cycle	(A)--(B)	(D)--(A)--(D')	(C)--(B)--(C');
			\foreach \x/\g in {C/0,D/-90,D'/180,C'/90,A/90,B/-90}
			\fill[black] 	(\x) circle (1pt)
			($(\g:3mm)+(\x)$) node {$\x$};	
		\end{tikzpicture}
	}
	\par  \shortans{121}
	\loigiai{
		\begin{center}
			\begin{tikzpicture}[scale=1.2,line join=round, line cap=round,font=\footnotesize]
				\def\a{4}
				\path 	(0:0) coordinate (D)
				(20:\a) coordinate (C)
				(135:.55*\a) coordinate (D')
				(80:.32*\a) coordinate (A)
				($(D')-(D)+(C)$) coordinate (C')
				($(C)-(D)+(A)$) coordinate (At)
				($(A)!.42!(At)$) coordinate (B)
				($(D)!.5!(C)$) coordinate (M)
				($(D)!.5!(A)$) coordinate (N)
				($(D')!.75!(D)$) coordinate (P);
				\draw	(C)--(D)--(D')--(C')--cycle	(A)--(B)	(D)--(A)--(D')	(C)--(B)--(C')	(P)--(N)--(M)	(A)--(C);
				\draw[dashed](M)--(P);
				\foreach \x/\g in {C/0,D/-90,D'/180,C'/90,A/90,B/120,M/-45,N/150,P/-150}
				\fill[black] 	(\x) circle (1pt)
				($(\g:3mm)+(\x)$) node {$\x$};	
				\draw pic[draw,angle radius=2mm]{right angle=P--N--A}
				pic[draw,angle radius=2mm]{right angle=M--N--A};
			\end{tikzpicture}
		\end{center}
		Goi $M$, $N$ lần lượt là trung điểm của $DC$, $DA$. \\
		Kẻ $NP \perp AD$ tại $N$, $P \in DD'$. Suy ra $\widehat{MNP}$ chính là góc phẳng nhị diện $[D', AD, C]$. \\
		Từ giả thiết suy ra $AM = AD = AD' = DM = 3$.\\
		Tam giác $ADM$ đều nên $MN = \dfrac{3\sqrt{3}}{2}$. \\
		Xét $\triangle ADD'$ có $\cos\widehat{ADD'} = \dfrac{DA^2 + DD'^2-AD'^2}{2 \cdot DA \cdot DD'} = \dfrac{3^2 + 4^2-3^2}{2 \cdot 3 \cdot 4} = \dfrac{2}{3}$. \\
		Xét $\triangle DNP$ vuông tại $N$ có 
		\begin{eqnarray*}
			\cos\widehat{NDP} = \dfrac{DN}{DP} &\Rightarrow& DP = \dfrac{DN}{\cos\widehat{NDP}} = \dfrac{9}{4} \\
			&\Rightarrow& PN = \sqrt{DP^2-DN^2} = \dfrac{3\sqrt{5}}{4}.
		\end{eqnarray*}
		Xét $\triangle DMP$ vuông tại $D$ có $MP = \sqrt{DM^2 + DP^2} = \dfrac{15}{4}$. \\
		Xét $\triangle MNP$ có 
		\begin{eqnarray*}
			\cos\widehat{MNP} &=& \dfrac{NM^2 + NP^2-PM^2}{2 \cdot NM \cdot NP} \\
			&=& \dfrac{\left(\dfrac{3\sqrt{3}}{2}\right)^2 + \left(\dfrac{3\sqrt{5}}{4}\right)^2-\left(\dfrac{15}{4}\right)^2}{2 \cdot \dfrac{3\sqrt{3}}{2} \cdot \dfrac{3\sqrt{5}}{4}}\\
			&=&-\dfrac{2\sqrt{15}}{15}.
		\end{eqnarray*}
		Suy ra $\widehat{MNP}\approx 121^\circ$. \\
		Vậy Số đo góc phẳng nhị diện $[D', AD, C]$ xấp xỉ $121^\circ$. 
	}
\end{ex}

\begin{ex}%[1H8V6-2]%[Dự án C đề thi thử TN Môn Toán NH25-26 - Dot 4 - Phạm Đăng Đăng]%[THPT Lê Quý Đôn-Đống Đa-Hà Nội]%[GVPB - Phan Quốc Trí]
	Cho hình chóp tứ giác đều $S.ABCD$ có cạnh đáy và cạnh bên đều bằng $a$. Côsin của số đo góc phẳng nhị diện $[D,SA,B]$ là $-\dfrac{a}{b}$ (với $\dfrac{a}{b}$ là phân số tối giản và $a$, $b\in\mathbb{N}^*$). Tính $P=a-b$.
	\shortans{$-2$}
	\loigiai{
		\begin{center}
			\begin{tikzpicture}[scale=1,font=\footnotesize,line join=round
				,line cap=round,>=stealth]
				\def\a{4}
				\def\h{4.5}
				\path 	(0:0) coordinate (A)
				++(0:\a) coordinate (D)
				++(-130:\a/2) coordinate (C)
				($(A)+(C)-(D)$) coordinate (B)
				(intersection of A--C and B--D) coordinate (O)
				($(O)+(90:\h)$) coordinate (S)
				
				($(A)!1/2!(S)$) coordinate (M);
				\draw[dashed,thick] 	(B)--(A)--(D)	(A)--(S)
				(A)--(C)	(B)--(D)	(S)--(O) (D)--(M)--(B)	;
				\draw[thick] 			(B)--(C)--(D)
				(B)--(S)	(C)--(S)	(D)--(S);
				\foreach \x/\g in {A/270,B/-135,C/-45,D/45,S/90,O/-90,M/-70}
				\fill[black] 	(\x) circle (1.5pt)
				($(\g:3mm)+(\x)$) node {$\x$};	
				\draw pic[draw,angle radius=3mm]{right angle=D--O--S};%Theo chiều dương	
			\end{tikzpicture}
		\end{center}
		Gọi $M$ là trung điểm của $SA$. Do các mặt bên của hình chóp là các tam giác đều cạnh $a$ nên $BM = DM = \dfrac{a\sqrt{3}}{2}$.\\
		Lại có $BM \perp SA$, $DM \perp SA$ nên số đo của nhị diện $[D,SA,B]$ bằng số đo góc $\widehat{BMD}$.\\
		Trong tam giác $BMD$ có
		$$
		\cos \widehat{BMD} = \dfrac{BM^2 + DM^2 - BD^2}{2\cdot BM \cdot DM}
		= \dfrac{\dfrac{3a^2}{4} + \dfrac{3a^2}{4} - 2a^2}{2\cdot \dfrac{a\sqrt{3}}{2} \cdot \dfrac{a\sqrt{3}}{2}}
		= \dfrac{-\dfrac{a^2}{2}}{\dfrac{3a^2}{2}}
		= -\dfrac{1}{3}.
		$$
		Vậy $a=1$, $b=3 \Rightarrow P=a-b=1-3=-2$.}
\end{ex}

\begin{ex}%[1H8V6-7]%[KNTT - Lớp 12 - Dự án C - Đợt 2]%[Loc do]%[PB: Phan Anh]
	\immini[thm]{Anh Nghĩa muốn xây một nhà kho dạng hình hộp chữ nhật có chiều rộng $AB=5$ m và làm thêm phần mái nhà $ABCDEF$ như hình vẽ, biết rằng $ADFE$ là một nửa của hình lục giác đều cạnh bằng $4$ m, $BCFE$ là hình thang cân có $EB=3$ m. Để đảm bảo tính chính xác cho việc thi công, Anh Nghĩa muốn xác định góc nhị diện $[B,AE,D]$. Hãy giúp anh Nghĩa tính góc nhị diện đó? (\textit{đơn vị: Độ, làm tròn kết quả đến hàng đơn vị}). }{\begin{tikzpicture}[declare function={a=1;b=4;h=2;},line join=round]
			\path (0,0) coordinate (B')
			(135:a) coordinate (A')
			(b,0) coordinate (C')
			($(C')-(B')+(A')$) coordinate (D')
			($(A')+(90:h)$) coordinate (A)
			($(B')-(A')+(A)$) coordinate (B)
			($(C')-(A')+(A)$) coordinate (C)
			($(D')-(A')+(A)$) coordinate (D);
			\path (0.3,3.5) coordinate (E);
			\path (2.8,3.5) coordinate (F);
			\draw ( B)--(B')--(C') (A)--(B)--(C) (A)--(E)--(B)  (C')--(C)--(F) (A)--(A')--(B') (E)--(F);
			\draw[dashed]  (A')--(D') (C')--(D')--(D)--(A) (F)--(D)--(C);
			\foreach \t/\g in {A/90,B/180,C/90,D/60,E/90,F/90}{
				\draw[fill=black] (\t) circle (1pt) node[shift={(\g:7pt)},font=\scriptsize]{$ \t $};
			}
	\end{tikzpicture}}
	\par\shortans[0]{140}
	\loigiai{
		\begin{center}
			\begin{tikzpicture}[declare function={a=1;b=4;h=2;},line join=round]
				\path (0,0) coordinate (B')
				(135:a) coordinate (A')
				(b,0) coordinate (C')
				($(C')-(B')+(A')$) coordinate (D')
				($(A')+(90:h)$) coordinate (A)
				($(B')-(A')+(A)$) coordinate (B)
				($(C')-(A')+(A)$) coordinate (C)
				($(D')-(A')+(A)$) coordinate (D);
				\path (0.3,3.5) coordinate (E);
				\path (2.8,3.5) coordinate (F);
				\path ($(A)!0.5!(D)$) coordinate (G)
				($(A)!0.5!(E)$) coordinate (H)
				($(A)!0.5!(B)$) coordinate (I);				
				\draw ( B)--(B')--(C') (A)--(B)--(C) (A)--(E)--(B)  (C')--(C)--(F) (A)--(A')--(B') (E)--(F) (H)--(I);
				\draw[dashed]  (A')--(D') (C')--(D')--(D)--(A) (F)--(D)--(C) (I)--(G)--(H);
				\foreach \t/\g in {A/90,B/180,C/90,D/60,E/90,F/90,G/-90,H/120,I/-120}{
					\draw[fill=black] (\t) circle (1pt) node[shift={(\g:7pt)},font=\scriptsize]{$ \t $};
				}
			\end{tikzpicture}
		\end{center}
		Gọi $G$, $H$, $I$ lần lượt là trung điểm của $AD$, $AE$, $AB$.\\
		Tam giác $AGE$ đều nên $GH\perp AE$.\\
		Ta có $AE^2+BE^2=4^2+3^2=5^2=AB^2$ suy ra tam giác $ABE$ vuông tại $E$.\\
		Suy ra $BE\perp AE$, mà $HI \parallel BE \Rightarrow HI\perp AE$.\\
		Mà $GH\perp AE$ nên $[B,AE,D]=\widehat{GHI}$.\\
		Ta có
		\begin{itemize}
			\item $GH=4\cdot \dfrac{\sqrt{3}}{2}=2\sqrt{3}$ m;
			\item $HI=\dfrac{1}{2} BE=\dfrac{3}{2}$ m;
			\item $AD=8$ m.
		\end{itemize}
		Hình chưc nhật $ABCD$ có $BD=\sqrt{BA^2+AD^2}=\sqrt{89} \Rightarrow GI=\dfrac{1}{2}DB=\dfrac{\sqrt{89}}{2}$ m.\\
		Từ đó suy ra $\cos GHI=\dfrac{HG^2+HI^2-GI^2}{2\cdot HG\cdot HI}=-\dfrac{4}{3\sqrt{3}}$.\\
		Vậy $\widehat{GHI}\approx 140^\circ$.
	}
\end{ex}

\begin{ex}%[TT-THPT-HamRong-NH25-26 ]%[GV soạn: Nguyễn Quang Hiệp GVPB2 Đỗ Minh Vũ ]%[1H8V7-2]
	Tính thể tích của khối chóp tứ giác đều $S.ABCD$ có chiều cao bằng $\sqrt{3}$ và độ dài cạnh đáy bằng $2\sqrt{3}$. (Kết quả làm tròn đến hàng phần trăm).
	\par\shortans[oly]{6,93}
	\loigiai{
		\begin{center}
			\begin{tikzpicture}[>=stealth,line join=round,line cap=round,font=\footnotesize,scale=1]
				\tikzset{
					pics/hinhChopTuGiacDeu/.style  n args={6}{
						code={
							\tikzset{
								declare function={a=3;b=1.5;h=3;goc=-140;}
							}	
							\path 
							(0,0)coordinate (#1)+(0:a)coordinate (#2)+(goc:b)coordinate (#4)			
							($(#2)+(#4)-(#1)$)coordinate (#3)
							(intersection of #1--#3 and #2--#4)coordinate (#5) ($(#5)+(90:h)$)coordinate (#6)
							;
						}
				}}
				\path 
				(0,0)pic {hinhChopTuGiacDeu={A}{B}{C}{D}{H}{S}}	;	
				
				\foreach \pointo/\pointt in {S/B,S/C,S/D,C/D,B/C}{
					\draw[fill=black](\pointo)--(\pointt);
				}	
				\foreach \pointo/\pointt in {S/A,A/B,A/D,A/C,B/D,S/H}{
					\draw[fill=black,dashed](\pointo)--(\pointt);
				}
				\foreach \point/\goc in {A/120,S/90,B/11,C/-45,D/200,H/-90}{
					\draw[fill=black](\point)circle(.8pt)+(\goc:2mm)node[scale=.8]{$\point$};
				}
			\end{tikzpicture}
		\end{center}
		Gọi $H$ là tâm của hình vuông đáy $ABCD$. Khi đó $SH$ là chiều cao của hình chóp.\\
		Theo giả thiết, ta có chiều cao $SH = \sqrt{3}$ và cạnh đáy $AB = 2\sqrt{3}$.\\
		Diện tích đáy của khối chóp tứ giác đều là $S_{ABCD} = \left(2\sqrt{3}\right)^2 = 12$.\\
		Thể tích của khối chóp là
		\[V = \dfrac{1}{3} \cdot S_{ABCD} \cdot SH = \dfrac{1}{3} \cdot 12 \cdot \sqrt{3} = 4\sqrt{3} \approx 6{,}93.\] 
	}
\end{ex}

\begin{ex}%[1H8V7-6]%[Dự án C đề thi thử TN Môn Toán NH25-26 - Dot 4 - Phạm Đăng Đăng]%[THPT Lê Quý Đôn-Đống Đa-Hà Nội]%[GVPB - Phan Quốc Trí]
	Để xây dựng các cột điện đường dây $220$ KV, người ta phải đổ bê tông những cột chân đế cột điện. Biết rằng các chân đế cột điện là các hình chóp cụt tứ giác đều có cạnh trên $4$ mét, cạnh dưới $6$ mét và các cạnh bên $4$ mét. Giá bê tông tươi là $1\,400\,000$ đồng trên $1$ mét khối. hỏi cần bao nhiêu tiền (đơn vị triệu đồng, kết quả làm tròn đến hàng đơn vị) để mua bê tông tươi làm một chân đế cột điện? 
	\shortans{$133$}
	\loigiai{\begin{center}
			\begin{tikzpicture}[scale=1,font=\footnotesize,line join=round
				,line cap=round,>=stealth]
				\def\a{4}
				\def\b{2}
				\def\h{3}
				\def\gocB{40}
				\def\goc{70}
				\path
				(0,0) coordinate (B')
				(\gocB:\b) coordinate (A')
				(0:\a) coordinate (C')
				(\goc:\h) coordinate (B)
				($(C')+(A')-(B')$) coordinate (D')
				(B)--++(0:\a*.7) coordinate (C)
				(B)--++(\gocB:\b*.5) coordinate (A)
				($(A)+(C)-(B)$) coordinate (D)
				;
				
				\coordinate (O) at (intersection of A--C and B--D);
				
				\coordinate (O') at (intersection of A'--C' and B'--D');
				\path ($(O')-(O)+(A)$) coordinate (E')
				;
				\coordinate (E) at (intersection of A--E' and A'--C');
				\draw (A)--(C)(B)--(D)(B)--(A)--(D)--(C)--cycle (B)--(B')--(C')--(C) (C')--(D')--(D);
				\draw[dashed] (O)--(O')(A')--(C') (A)--(E) (B')--(D')(B')--(A')--(D') (A')--(A);
				\foreach \x/\g in {A'/135,B'/180,C'/-45,D'/45,A/135,B/180,C/-45,D/45,O'/270,O/40,E/270}
				\fill (\x) circle (1pt) +(\g:0.3) node {$\x$};
			\end{tikzpicture}
		\end{center}
		Kẻ $AE \perp A'C'$, ta có $A'E = \dfrac{1}{2}(A'C' - AC) = \dfrac{1}{2}(6\sqrt{2} - 4\sqrt{2}) = \sqrt{2}$.\\
		Suy ra $AE = \sqrt{AA'^2 - A'E^2} = \sqrt{16 - 2} = \sqrt{14}$.\\Diện tích hai đáy khối chóp cụt là
		$S_1 = 4^2 = 16, S_2 = 6^2 = 36$.\\
		Vậy thể tích khối chóp cụt là $$V = \dfrac{1}{3}AE\left(S_1 + S_2 + \sqrt{S_1S_2}\right) = \dfrac{1}{3}\sqrt{14}\left(16 + 36 + \sqrt{16 \cdot 36}\right) = \dfrac{76\sqrt{14}}{3}.$$
		Số tiền cần tính là $T = V\cdot 1{,}4 \approx 133$ triệu đồng.
	}
\end{ex}

\begin{ex}%[1H8V7-9]%[Dự án đề thi thử NH25-26]%[THPT Chuyên Phan Bội Châu - Nghệ An]%[BC Tuan]
	\immini[thm]{Hai con thằn lằn A và B đang bám ở hai bức tường đối diện của một căn phòng dạng hình hộp chữ nhật với chiều rộng, chiều dài, chiều cao lần lượt là $8$ (m), $12$ (m), $5$ (m). Ban đầu thằn lằn A ở vị trí cách bức tường phía trước và trần nhà lần lượt là $7$ (m) và $3$ (m), còn thằn lằn B ở vị trí cách bức tường phía trước và trần nhà lần lượt là $9$ (m) và $4$ (m) (tham khảo hình vẽ bên). Sau đó chúng nhìn thấy nhau và chạy lại gặp nhau. Biết rằng hai con thằn lằn chỉ chạy trên các bức tường và trần nhà, hỏi tổng quãng đường ngắn nhất hai con thằn lằn di chuyển là bao nhiêu mét (làm tròn kết quả đến hàng đơn vị)?}{
		\begin{tikzpicture}[line join=round, line cap=round, declare function={w=4.5; h=3; d=6.2; goc=75; xd=1.6; wd=1.3; hd=1.8;},font=\footnotesize,scale=0.7]
			\path 
			(0,0) coordinate (A)
			(w,0) coordinate (B)
			(w,h) coordinate (C)
			(0,h) coordinate (D)
			(goc:d) coordinate (V)
			($(A)+(V)$) coordinate (A')
			($(B)+(V)$) coordinate (B')
			($(C)+(V)$) coordinate (C')
			($(D)+(V)$) coordinate (D')
			(xd,0) coordinate (D1)
			(xd,hd) coordinate (D2)
			(xd+wd,hd) coordinate (D3)
			(xd+wd,0) coordinate (D4)
			($(A)!0.55!(D')$) coordinate (F)
			($(B)!0.62!(C)$) coordinate (S_f)
			($(B')!0.62!(C')$) coordinate (S_b)
			($(S_f)!0.3!(S_b)$) coordinate (S);			
			\draw[dashed] (A) -- (A') (A') -- (B') (A') -- (D');
			\draw (A) -- (D1) -- (D2) -- (D3) -- (D4) -- (B) -- (C) -- (D) -- cycle;
			\draw (D) -- (D') -- (C') -- (C) (B) -- (B') -- (C');
			\node[below, font=\footnotesize] at ($(D1)!0.5!(D4)$) {Cửa};
			\draw (F) node[above]{$A$} (S) node[above right]{$B$};
		\end{tikzpicture}
	}
	\shortans{$15$}
	\loigiai{
		Ta dùng phương pháp trải phẳng hình hộp.\\
		Nếu hai con thằn lằn chạy gặp nhau bằng cách đi qua trần nhà, ta trải hình như sau:
		\begin{center}
			\begin{tikzpicture}[line join=round, line cap=round,font=\footnotesize,scale=0.4]
				\draw (0,0) rectangle (5,12) (5,0) rectangle (13,12) (13,0) rectangle (18,12);
				\draw (2,7) circle (1pt) node[above]{$A$} (17,8) circle (1.5pt) node[above]{$B$};
				\draw[dashed](2,7)--(2,0) node[midway,right]{$7$ (m)} (2,7)--(5,7) node[midway,above]{$3$ (m)} (17,8)--(17,0) node[midway,left]{$9$ (m)} (17,8)--(13,8) node[midway,above]{$4$ (m)};
				\draw[red] (2,7)--(17,8);
			\end{tikzpicture}
		\end{center}
		Khi đó khoảng cách ngắn nhất là đoạn thẳng nối hai điểm trên mặt phẳng trải \[AB=\sqrt{(3+8+4)^2+(9-7)^2} = \sqrt{15^2+2^2} = \sqrt{229} \approx 15{,}13 \,(\text{m}).\]
		Nếu hai con thằn lằn chạy gặp nhau bằng cách đi qua bức tường phía sau, ta trải hình như sau:
		\begin{center}
			\begin{tikzpicture}[line join=round, line cap=round,font=\footnotesize,scale=0.4]
				\draw (0,0) rectangle (12,5) (12,0) rectangle (20,5) (20,0) rectangle (32,5);
				\draw (7,2) circle (1.5pt) node[below]{$A$} (24,1) circle (1.5pt) node[below]{$B$};
				\draw[dashed](7,2)--(12,2) node[midway,above]{$5$ (m)} (24,1)--(20,1) node[midway,below]{$3$ (m)} (7,2)--(7,5) node[midway,left]{$3$ (m)} (24,1)--(24,5) node[midway,right]{$4$ (m)};
				\draw[red] (7,2)--(24,1);
			\end{tikzpicture}
		\end{center}
		Khi đó quãng đường là $AB=\sqrt{(5+8+3)^2+(4-3)^2} = \sqrt{16^2+1^2} = \sqrt{257} \approx 16{,}03$ (m).\\
		So sánh hai trường hợp, tổng quãng đường ngắn nhất mà hai con thằn lằn phải di chuyển là khoảng $15$ (m).
	}
\end{ex}

\begin{ex}%[1H8V7-9]%[Dự án C đề thi thử TN Môn Toán NH25-26 - Dot 4 - Phạm Hoàng Đăng]%[THPT Nguyễn Trãi-Hà Nội]%[GVPB - Khắc Thiên]
	\immini[thm]{Một chi tiết máy mẫu bằng nhựa đúc đặc có dạng hình chóp tứ giác đều, cạnh đáy bằng $20$ (cm). Do bị lỗi kĩ thuật, bên trong chi tiết này có một khoang rỗng. Để kiểm tra, kỹ sư thả chìm hoàn toàn chi tiết máy vào một bể chứa dung dịch chống gỉ có dạng hình lăng trụ đứng có đáy là hình vuông cạnh bằng $30$ (cm). Khi đó, mực dung dịch trong bể dâng lên thêm $2$ (cm) và dung dịch đã tràn vào lấp đầy khoang rỗng bên trong. Sau khi vớt chi tiết ra và lau khô bề mặt, khối lượng của nó tăng thêm $160$ (gam) so với ban đầu. Biết khối lượng riêng của dung dịch là $0{,}8$ (gam/cm$^3$). Hãy tính độ dài cạnh bên của hình chóp nói trên (\textit{làm tròn kết quả đến hàng phần mười}).}{
		\begin{tikzpicture}[scale=1,font=\footnotesize,line join=round,line cap=round,>=stealth]
			\draw[dashed,black,scale=2.5] (0,0) coordinate (A)--(-130:0.5) coordinate (B)
			(1,0) coordinate (D)--(A)--($(A)+(90:1.2)$) coordinate (A');
			\draw[black,scale=2.5] (B)--($(D)-(A)+(B)$) coordinate (C)--(D)--($(D)+(90:1.2)$) coordinate (D')--(A')--($(B)+(90:1.2)$) coordinate (B')--($(D')-(A')+(B')$) coordinate (C')--(D')
			(C')--(C)
			(B')--(B);
			\fill[cyan!55] ($(B)+(90:1.5)$)--($(C)+(90:1.5)$)--($(D)+(90:1.5)$)--($(A)+(90:1.5)$)--($(B)+(90:1.5)$);
			\fill[cyan!75] (B)--(C)--($(C)+(90:1.5)$)--($(B)+(90:1.5)$)--(B);
			\fill[cyan!90] (D)--(C)--($(C)+(90:1.5)$)--($(D)+(90:1.5)$)--(D);
			\draw[dashed,black,scale=2.5] (0,0) coordinate (A)--(-130:0.5) coordinate (B)
			(1,0) coordinate (D)--(A)--($(A)+(90:1.2)$) coordinate (A');
			\draw[black,scale=2.5] (B)--($(D)-(A)+(B)$) coordinate (C)--(D)--($(D)+(90:1.2)$) coordinate (D')--(A')--($(B)+(90:1.2)$) coordinate (B')--($(D')-(A')+(B')$) coordinate (C')--(D')
			(C')--(C)
			(B')--(B);
			\path ($(A)!.5!(C)$) coordinate (O)
			($(O)+(0,1.5)$) coordinate (S)
			($(A)!1/3!(O)$) coordinate (Aa)
			($(C)!1/3!(O)$) coordinate (Cc)
			($(B)!1/3!(O)$) coordinate (Bb)
			($(D)!1/3!(O)$) coordinate (Dd)
			;
			\draw[line width=0.25]
			(Bb)--(Cc)--(Dd) (S)--(Bb) (S)--(Dd) (S)--(Cc)
			;
			\draw[dashed, line width=0.25]
			(S)--(Aa) (Bb)--(Aa)--(Dd)
			;
		\end{tikzpicture}
	}
	\par
	\shortans[oly]{20{,}6}
	\loigiai{
		\immini[thm]{
			Chi tiết máy có dạng hình chóp tứ giác đều $S.ABCD$ như hình vẽ.\\
			Ta có thể tích khoang rỗng là $V_1=\dfrac{160}{0{,}8}=200$ (cm$^3$).\\
			Thể tích dung dịch dâng lên $V_2=30^2\cdot 2=1\,800$ (cm$^3$).\\
			Suy ra thể tích khối chóp là $V=200+1\,800=2\,000$ (cm$^3$).\\
			Vì $V=\dfrac{1}{3} \cdot 20^2 \cdot SO$ nên suy ra $SO=15$.\\
			Lại có $ABCD$ là hình vuông nên $BD=20\sqrt{2}$.\\
			Suy ra $BO=10\sqrt{2}$.\\
			Do đó $SB=\sqrt{SO^2+BO^2}=5\sqrt{17}\approx 20{,}6$ (cm).
		}{
			\begin{tikzpicture}[>=stealth,line join=round,line cap=round,font=\footnotesize,scale=1]
				\def\a{3}
				\def\b{1}
				\path (0,0) coordinate (O) +(90:\a) coordinate (S)
				+(140:\b) coordinate (A)
				(A)+(0:\a) coordinate (D)
				($(A)!2!(O)$) coordinate (C)
				($(D)!2!(O)$) coordinate (B)
				($(D)!.5!(C)$) coordinate (M);
				\draw (S)--(B)--(C)--(D)--cycle (S)--(C);
				\draw[dashed](S)--(A)--(D)
				(A)--(B)
				(S)--(O)
				(A)--(C)
				(B)--(D);
				\foreach \x/\g in {S/90,A/150,B/210,C/-30,D/0,O/-90} \fill (\x) circle (0.03)+(\g:3mm) node[scale=.9] {$\x$};
			\end{tikzpicture}
		}
	}
\end{ex}

\begin{ex}%[1H8V7-9]%[TN-THPT-NguyenTrungThien-HaTinh-NH25-26]%[Loc Do]%[PB: Phan Anh]
	Có một tấm bìa ngũ giác đều $MNPQR$ tâm $O$ cạnh bằng $30$ cm, cắt bỏ $5$ tam giác cân bằng nhau có cạnh đáy là cạnh của ngũ giác đều ban đầu, đỉnh là đỉnh của ngũ giác đều bên trong như hình vẽ rồi gập lên thành một khối chóp ngũ giác đều có cạnh đáy bằng $x$ (cm). Tìm $x$ để thể tích tạo thành là lớn nhất (\textit{kết quả làm tròn đến hàng phần chục}).
	\begin{center}
		\begin{tikzpicture}[scale=0.7, font=\footnotesize, line join=round, line cap=round, >=stealth]	
			\coordinate (A) at (240:1.5cm);
			\coordinate (B) at (312:1.5cm);
			\coordinate (C) at (24:1.5cm);
			\coordinate (D) at (96:1.5cm);
			\coordinate (E) at (168:1.5cm);
			\coordinate (M) at (60:4cm);
			\coordinate (R) at (132:4cm);
			\coordinate (Q) at (204:4cm);
			\coordinate (P) at (276:4cm);
			\coordinate (N) at (348:4cm);
			\path (intersection of A--M and B--R) coordinate (O);
			\path (intersection of E--D and O--R) coordinate (I);
			\fill [orange!10] (D)--(R)--(E)--(Q)--(A)--(P)--(B)--(N)--(C)--(M)--cycle;
			% Ngũ giác lớn (bắt đầu từ 60°)
			\draw[thick] (60:4cm)
			\foreach \i in {132,204,276,348} { -- (\i:4cm) } -- cycle;
			
			% Ngũ giác nhỏ (xoay 180° so với ngũ giác lớn)
			\draw[thick] (240:1.5cm)
			\foreach \i in {312,24,96,168} { -- (\i:1.5cm) } -- cycle;
			% Nối các đỉnh
			\foreach \i in {60,132,204,276,348} {
				\draw[thick] (\i:4cm) -- (\i+36:1.5cm); 
				\draw[thick] (\i:4cm) -- (\i-36:1.5cm);
			}
			\fill (O) circle (1.5pt) node[below right] {$O$};
%			\fill (I) circle (1.5pt) node[left] {$I$};
%			\draw[thick] (R)--(O);
			% Vẽ điểm
			\foreach \i/\name/\g in {60/M/30,132/R/120,204/Q/200,276/P/-90,348/N/0} {
				\fill[red] (\i:4cm) circle (1.5pt)+(\g:3mm) node[scale=0.8] {$\name$};
			}
			\foreach \i/\name/\t in {240/A/-110,312/B/-45,24/C/10,96/D/90,168/E/170} {
				\fill[blue] (\i:1.5cm) circle (1.5pt)+(\t:3mm) node[scale=0.8] {$\name$};
			}		
		\end{tikzpicture}
	\end{center}
	\shortans{14,8}
	\loigiai{\begin{center}
			\begin{tikzpicture}[scale=0.7, font=\footnotesize, line join=round, line cap=round, >=stealth]	
				\coordinate (A) at (240:1.5cm);
				\coordinate (B) at (312:1.5cm);
				\coordinate (C) at (24:1.5cm);
				\coordinate (D) at (96:1.5cm);
				\coordinate (E) at (168:1.5cm);
				\coordinate (M) at (60:4cm);
				\coordinate (R) at (132:4cm);
				\coordinate (Q) at (204:4cm);
				\coordinate (P) at (276:4cm);
				\coordinate (N) at (348:4cm);
				\path (intersection of A--M and B--R) coordinate (O);
				\path (intersection of E--D and O--R) coordinate (I);
				\coordinate (K) at ($(Q)!0.5!(R)$);
				\fill [orange!10] (D)--(R)--(E)--(Q)--(A)--(P)--(B)--(N)--(C)--(M)--cycle;
				% Ngũ giác lớn (bắt đầu từ 60°)
				\draw[thick] (60:4cm)
				\foreach \i in {132,204,276,348} { -- (\i:4cm) } -- cycle;
				
				% Ngũ giác nhỏ (xoay 180° so với ngũ giác lớn)
				\draw[thick] (240:1.5cm)
				\foreach \i in {312,24,96,168} { -- (\i:1.5cm) } -- cycle;
				% Nối các đỉnh
				\foreach \i in {60,132,204,276,348} {
					\draw[thick] (\i:4cm) -- (\i+36:1.5cm); 
					\draw[thick] (\i:4cm) -- (\i-36:1.5cm);
				}
				\fill (O) circle (1.5pt) node[below right] {$O$};
				\fill (I) circle (1.5pt) node[left] {$I$};
				\draw[thick] (R)--(O)--(M) (N)--(O)--(P) (K)--(O)--(Q);
				% Vẽ điểm
				\foreach \i/\name/\g in {60/M/30,132/R/120,204/Q/200,276/P/-90,348/N/0} {
					\fill[red] (\i:4cm) circle (1.5pt)+(\g:3mm) node[scale=0.8] {$\name$};
				}
				\foreach \i/\name/\t in {240/A/-110,312/B/-45,24/C/10,96/D/90,168/E/-100} {
					\fill[blue] (\i:1.5cm) circle (1.5pt)+(\t:3mm) node[scale=0.8] {$\name$};
				\fill[red] (K) circle (1.5pt)+(135:3mm) node[scale=0.8] {$K$};
				}		
			\end{tikzpicture}
		\qquad\qquad
		\begin{tikzpicture}[declare function={gocx=90;goc=-34;a=4;b=a*0.3;d=0.5*a;c=0.4*a;h=4;}]
			\path 
			(0,0) coordinate (A)
			(-0.5*goc:c) coordinate (E)
			(3.6,0) coordinate (D)
			(goc:b) coordinate (B)
			($(B)+(d,0)$) coordinate (C)
			($(D)!0.5!(C)$) coordinate (O1)	
			($(D)!0.5!(E)$) coordinate (I)	
			($(A)!0.55!(O1)$) coordinate (O)	
			($(gocx:h)+(O)$) coordinate (S);
			
			\draw[dashed] (A)--(E)--(D) (I)--(O)--(S)--(E) (S)--(I);
			\draw (A)--(B)--(C)--(D)--(S)--(A) (S)--(C) (S)--(B);
			\path pic[angle radius=3mm,draw=blue] {right angle = S--O--I};
			
			\foreach \t/\g in {A/-180,D/0,C/-60,S/90,B/-100,E/-90,O/-90, I/-90}{
				\draw[fill=black] (\t) circle (1pt)
				node[shift={(\g:7pt)},font=\scriptsize]{$\t$};
			}
		\end{tikzpicture}
		\end{center}
		$MNPQR$ là ngũ giác đều nên ta có \[\widehat{ROM}=\widehat{MON}=\widehat{NOP}=\widehat{POQ}=\widehat{QOR}=\dfrac{360^{^\circ}}{5} =72^{^\circ}.\]
		Xét tam giác $OQR$ cân tại $O$, đường cao $OK$ \[\widehat{ROK}=\widehat{KOQ}=\dfrac{\widehat{ROQ}}{2} =\dfrac{72^{^\circ}}{2} =36^{^\circ}.\]
		Xét tam giác $ROK$ vuông tại $K$
		\[\sin \widehat{ROK}=\dfrac{RK}{OR} \Rightarrow OR=\dfrac{RK}{\sin \widehat{ROK}} =\dfrac{RQ}{2\cdot\sin 36^{^\circ}} =\dfrac{30}{2\cdot\sin 36^{^\circ}} =\dfrac{15}{\sin 36^{^\circ}}\, (\rm{cm}).\]
		Ta có $EI=ID=\dfrac{ED}{2} =\dfrac{x}{2}$ (cm).\\
		Xét tam giác $OIE$ vuông tại $I$
		\[\tan \widehat{EOI}=\dfrac{EI}{OI} \Leftrightarrow OI=\dfrac{EI}{\tan \widehat{EOI}} =\dfrac{x}{2\cdot\tan 36^{^\circ}}\, (\rm{cm}).\]
		Suy ra $IR=OR-OI=\dfrac{15}{\sin 36^{^\circ}} -\dfrac{x}{2\cdot\tan 36^{^\circ}}$ (cm).\\
		Gọi đỉnh chóp là $S$, sau khi được gấp lên thì các điểm $M$, $N$, $P$, $Q$, $R$, $S$ trùng nhau.\\
		Chiều cao khối chóp là
		\[h=SO=\sqrt{IR^{2} -OI^{2}} =\sqrt{\left(\dfrac{15}{\sin 36^{^\circ}} -\dfrac{x}{2\tan 36^{^\circ}} \right)^{2} -\left(\dfrac{x}{2\tan 36^{^\circ}} \right)^{2}} \quad (\rm{cm}).\]
		Diện tích đáy khối chóp là
		\[S_{MNPQR} =5S_{OED} =5\cdot \dfrac{1}{2} \cdot OI\cdot DE =5\cdot \dfrac{1}{2} \cdot \dfrac{x}{2\tan 36^{^\circ}} \cdot x=\dfrac{5x^{2}}{4\tan 36^{^\circ}} \quad (\rm{cm}^2).\]
		Thể tích khối chóp là
		\begin{align*}
			V=\dfrac{1}{3} S_{MNPQR} \cdot h=  \dfrac{1}{3} \sqrt{\left(\dfrac{15}{\sin 36^{^\circ}} -\dfrac{x}{2\tan 36^{^\circ}} \right)^{2} -\left(\dfrac{x}{2\tan 36^{^\circ}} \right)^{2}} \cdot \dfrac{5x^{2}}{4\tan 36^{^\circ}}.		
		\end{align*}
		Ta có $V'=0 \Rightarrow x\approx 14{,}8$ (cm).\\
		\begin{center}
			\begin{tikzpicture}[>=stealth]
				\tkzTabInit[nocadre=false,lgt=1,espcl=2,deltacl=0.5]{$x$/1.2 ,$V'$/.7,$V$/2}
				{$ $ , $14{,}8$ , $ $}
				\tkzTabLine{ , + , $0$ , - , }
				\tkzTabVar{-/$ $ , +/$V_{max}$ , -/$ $}
			\end{tikzpicture}
			
		\end{center}
		Vậy $V_{max}$  khi  $x\approx 14{,}8$ (cm).}
\end{ex}

\begin{ex}%[2D1B4-3]%[Dự án C-TT THPT NH25-26-Đợt 4- ĐỖ CHÍ TÂM]%[THPT Phan Bội Châu-Khánh Hòa]%[GVPB: Phạm Văn Long]
	Cho hàm số $y=\dfrac{ax+b}{cx+d}$ có đồ thị như hình vẽ. Tính $a+d$.
	\shortans{$0$}
	\begin{center}
		\begin{tikzpicture}[x=1cm,y=1cm, scale=0.7]
			% Vẽ lưới (tùy chọn, để khớp với ảnh gốc)
			%\draw[step=0.5, lightgray!30, very thin] (-3,-2) grid (5,5);
			
			% Vẽ trục tọa độ
			\draw[->] (-3.2,0) -- (5.5,0) node[below] {$x$};
			\draw[->] (0,-2.2) -- (0,4.5) node[left] {$y$};
			%\node at (-0.2,-0.2) {$O$};
			
			% Vẽ đường tiệm cận
			\draw[dashed, thick] (1,-2.2) -- (1,4.5); % Tiệm cận đứng x=1
			\draw[dashed, thick] (-3,1) -- (5.5,1); % Tiệm cận ngang y=1
			
			% Vẽ đồ thị hàm số y = (x+2)/(x-1)
			% Nhánh 1: x < 1
			\draw[thick, red, domain=-3:0.7, samples=100] plot (\x, {(\x-2)/(\x-1)});
			% Nhánh 2: x > 1
			\draw[thick, red, domain=1.33:5.5, samples=100] plot (\x, {(\x-2)/(\x-1)});
			
			% Đánh dấu các số 1, 2 trên trục
			\fill (1,0) circle(1.5pt) node[shift={(45:3mm)}]{$1$};
			\fill (0,1) circle(1.5pt) node[shift={(45:3mm)}]{$1$};
			\fill (2,0) circle(1.5pt) node[below]{$2$};
			\fill (0,2) circle(1.5pt) node[left]{$2$};
			\fill (0,0) circle(1.5pt) node[shift={(-120:3mm)}]{$0$};
		\end{tikzpicture}
	\end{center}	
	\loigiai{
		Từ đồ thị hàm số, ta có$\colon$\\
		Tiệm cận đứng $x=-\dfrac{d}{c}=1\Leftrightarrow c=-d\,\,\,\, (1)$.\\
		Tiệm cận ngang $y=\dfrac{a}{c}=1\Leftrightarrow a=c\,\,\,\, (2)$.\\
		Từ $(1)(2)$, suy ra $a=-d\Leftrightarrow a+d=0$.
	}
\end{ex}

\begin{ex}%[2D1C3-6]%[Dự án C đợt 3 NH25-26- Phạm Đăng Đăng]%[TT-THPT-TranNhanTong-HaNoi-N25-26]%[GVPB - Phan Quốc Trí]
	Một chiếc thuyền ở vị trí $A$ cách bờ một khoảng $AB$ bằng $5$ km. Trên bờ biển có một cái kho ở vị trí $C$ cách $B$ một khoảng là $7$ km. Người lái thuyền dự định chèo thuyền từ $A$ đến điểm $M$ trên bờ biển với vận tốc $1{,}2$ m/s rồi đi bộ đến $C$ với vận tốc $1{,}8$ m/s (xem hình vẽ dưới đây). Điểm $M$ cách $B$ bao nhiêu km để người đó đến kho nhanh nhất?
	\begin{center}
		\begin{tikzpicture}[scale=1, line cap=round, line join=round, font=\small,>=stealth]
			
			% Các điểm
			\coordinate (B) at (0,0);
			\coordinate (C) at (7,0);
			\coordinate (M) at (3.5,0);
			\coordinate (A) at (0,5);
			
			% Vẽ đoạn BC
			\draw (-2,0)--(B) -- (C)--(9,0);
			
			% Vẽ AB
			\draw (B) -- (A);
			
			% Vẽ AM nét đứt
			\draw[dashed] (A) -- (M);
			
			% Ghi tên điểm
			\fill (A) circle (1pt) node[left] at (A) {A};
			\fill (B) circle (1pt)	node[below] at (B) {B};
			\fill (C) circle (1pt)	node[below] at (C) {C};
			\fill (M) circle (1pt)	node[below] at (M) {M};
			
			% Ghi độ dài AB = 5km
			
			\node[left] at (0,2.5) {5 km};
			
			% Ghi độ dài BC = 7km
			\draw[<->] (0,-0.7) -- (7,-0.7);
			\node[below] at (3.5,-0.7) {$7$ km};
			
		\end{tikzpicture}
	\end{center}
	\shortans{$4{,}47$}
	\loigiai{Vận tốc chèo thuyền $1{,}2$ m/s $=0{,}0012$ km/s.\\
		Vận tốc đi bộ $1{,}8$ m/s $= 0{,}0018$ km/s.\\
		Gọi $MB = x$ (km) với $0 \le x \le 7$.\\
		Khi đó quãng đường chèo thuyền là
		$S_1 = AM = \sqrt{AB^2 + BM^2} = \sqrt{25 + x^2}
		$.\\
		Thời gian chèo thuyền là
		$
		t_1 = \dfrac{AM}{0{,}0012} = \dfrac{\sqrt{25 + x^2}}{0{,}0012} = \dfrac{2500}{3}\sqrt{25 + x^2}\,(\text{s})
		$.\\
		Quãng đường đi bộ là
		$MC = 7 - x$.\\
		Thời gian đi bộ là $t_2 = \dfrac{MC}{0{,}0018} = \dfrac{7 - x}{0{,}0018} = \dfrac{5000}{9}(7 - x)\,(\text{s})$.\\
		Để người đó đến kho nhanh nhất thì
		$
		f(x) = t_1 + t_2 = \dfrac{2500}{3}\sqrt{25 + x^2} + \dfrac{5000}{9}(7 - x)
		$
		đạt giá trị nhỏ nhất trên đoạn $[0;7]$.\\
		Ta có 
		$
		f'(x) = \dfrac{2500}{3} \cdot \dfrac{x}{\sqrt{25 + x^2}} - \dfrac{5000}{9}
		$.\\Khi đó $
		f'(x) = 0\Leftrightarrow \dfrac{2500}{3} \cdot \dfrac{x}{\sqrt{25 + x^2}} - \dfrac{5000}{9}=0 \Leftrightarrow 2\sqrt{25+x^2}=3x\Leftrightarrow  x = 2\sqrt{5}$.\\
		Bảng biến thiên
		
		\begin{center}
			\begin{tikzpicture}
				\tkzTabInit[lgt=1.2,espcl=2.5,deltacl=0.6]
				{$x$/0.6,$f'(x)$/0.6,$f(x)$/2}{$0$,$2\sqrt{5}$,$7$}
				\tkzTabLine{,-,0,+,}
				\tkzTabVar{+/$f(0)$,-/$f(2\sqrt{5})$,+/$f(7)$}
			\end{tikzpicture}
		\end{center}
		Ta thấy $f(x)$ đạt giá trị nhỏ nhất tại $x=2\sqrt{5}\approx 4{,}47$ km.
	}
\end{ex}

\begin{ex}%[2D1H3-6]%[Dự án đề TT Toán NH25-26-Dot 3- Hải Phụng]%[Cụm 13-Hải Phòng]%[GVPB - Viet Hoang Pham]
Giả sử doanh số (tính bằng số sản phẩm) của một sản phẩm mới (trong vòng một số năm nhất định) tuân theo quy luật logistic được mô hình hóa bằng hàm số $f(t) = \dfrac{3\,000}{1 + 3\mathrm{e}^{-t}}$, $t \geq 0$, trong đó thời gian $t$ được tính bằng năm, kể từ khi phát hành sản phẩm mới. Khi đó đạo hàm $f'(t)$ sẽ biểu thị tốc độ bán hàng. Hỏi sau khi phát hành bao nhiêu năm thì tốc độ bán hàng là lớn nhất? (kết quả làm tròn đến hàng phần mười).
\shortans[oly]{$1{,}1$}
\loigiai{
Ta có $f(t) = \dfrac{3\,000}{1 + 3\mathrm{e}^{-t}} = \dfrac{3\,000\mathrm{e}^t}{\mathrm{e}^t + 3} \Rightarrow f'(t) = \dfrac{9\,000\mathrm{e}^t}{\left( \mathrm{e}^t + 3 \right)^2}$.\\
Xét $h(t) = f'(t) = \dfrac{9\,000\mathrm{e}^t}{\left( \mathrm{e}^t + 3 \right)^2}; (t \geq 0)$.\\
Ta có $h'(t) = \dfrac{ \left( 9\,000\mathrm{e}^t \right)' \cdot \left( \mathrm{e}^t + 3 \right)^2 - 2 \left(\mathrm{e}^t + 3\right) \cdot \left(\mathrm{e}^t + 3\right)' \cdot 9\,000\mathrm{e}^t}{\left( \mathrm{e}^t + 3 \right)^4} = \dfrac{9\,000\mathrm{e}^t(3 - \mathrm{e}^t)}{(\mathrm{e}^t + 3)^3}$.\\
Cho $h'(t) = 0 \Leftrightarrow 3 - \mathrm{e}^t = 0 \Leftrightarrow \mathrm{e}^t = 3 \Leftrightarrow t = \ln 3$.\\
Bảng biến thiên của hàm số $h(t)$:
\begin{center}
\begin{tikzpicture}
\tkzTabInit[lgt=1.2,espcl=3]
{$t$/0.8,$h'(t)$/0.8,$h(t)$/1.5}
{0,$\ln 3$,$+\infty$}
\tkzTabLine{,+,0,-,}
\tkzTabVar{-/ , +/ , -/ }
\end{tikzpicture}
\end{center}
Từ bảng biến thiên, suy ra tốc độ bán hàng $h(t)$ lớn nhất khi $t = \ln 3 \approx 1{,}1$.
}
\end{ex}

\begin{ex}%[TT-LienTruong-HaNoi-N25-26]%[Phạm Quốc Toàn - Lê Phúc]%[2D1H3-6]
	Một xí nghiệp chế biến cà phê bán trong nước và xuất khẩu ra nước ngoài. Nếu giá sản xuất mỗi kg cà phê là $x$\,(USD) thì sản lượng xí nghiệp sản suất được là $(2\,000x-150)$\,(kg) và lượng tiêu thụ trong nước là $(4\,000-500x)$\,(kg). Phần cà phê dư sẽ được xuất khẩu với giá cố định là $10$\,(USD) mỗi kg. Biết rằng với mỗi kg cà phê xuất khẩu thì xí nghiệp phải chịu thuế là $0{,}5$\,(USD). Hỏi giá sản xuất mỗi kg cà phê của xí nghiệp là bao nhiêu để lợi nhuận thu được từ xuất khẩu là lớn nhất?
	\shortans{$5{,}58$}
	\loigiai{
		Sản lượng cà phê dư thừa dành cho xuất khẩu là
		\[
		(2\,000x-150)-(4\,000-500x)=2\,500x-4\,150 \text{ (kg)}.
		\]
		Chi phí để sản xuất lượng cà phê dành cho xuất khẩu là $C(x)=(2\,500x-4\,150)x$. \\ 
		Doanh thu từ việc xuất khẩu là $D(x)=(2\,500x-4\,150)\cdot 10$. \\ 
		Chi phí chịu thuế $(2\,500x-4\,150)\cdot 0{,}5$. \\ 
		Lợi nhuận thu được từ xuất khẩu là 
		\begin{align*}
			L(x) & =(2\,500x-4\,150)\cdot 10-(2\,500x-4\,150)\cdot 0{,}5-C(x) \\ 
			& =(2\,500x-4\,150)\cdot 10-(2\,500x-4\,150)\cdot 0{,}5-(2\,500x-4\,150)\cdot x \\ 
			& =-2\,500 x^2+27\,900x-39\,425. 
		\end{align*}
		Hàm số $L(x)$ là một tam thức bậc hai có hệ số $a = -2\,500<0$ nên hàm số đạt giá trị lớn nhất tại $x = - \dfrac{27\,900}{2 \cdot (-2\,500)} = \dfrac{279}{50} \approx 5{,}58$. 
	}
\end{ex}

\begin{ex}%[TT-LienTruong-HaNoi-N25-26]%[Phạm Quốc Toàn - Lê Phúc]%[2D1H3-6]
	Để làm công tác cứu nạn ở một ngôi nhà có một hàng rào cao $2{,}2$ (m) song song và cách bức tường ngôi nhà một khoảng bằng $1{,}6$ (m). Người ta dựng một cái thang thẳng có một đầu chạm đất, đầu kia chạm vào bức tường ngôi nhà, thân của thang chạm vào mép trên hàng rào (tham khảo hình vẽ). Chiều dài ngắn nhất của cái thang là bao nhiêu? (\textit{đơn vị là mét, không làm tròn kết quả các phép tính trung gian, chỉ làm tròn kết quả cuối cùng đến hàng phần trăm}). 
	\begin{center}
		\begin{tikzpicture}[x={(1cm,0cm)}, y={(0cm,1cm)}, z={(0.4cm,0.3cm)}, scale=0.5, line join=round, line cap=round]
			\def\W{7}         % Chiều rộng ngôi nhà
			\def\H{5.5}       % Chiều cao tường nhà
			\def\D{4}         % Chiều sâu ngôi nhà
			\def\Pc{3.5}      % Vị trí tâm cổng
			\def\ladX{11.5}   % Vị trí chân thang (đặt ngoài bãi cỏ)
			\def\ladZ{-1}     % Vị trí thang theo chiều sâu z
			\def\xL{0}        
			\def\xR{\W}       
			\def\Pd{0.5}      
			% Thảm cỏ
			\draw[fill=green!80!black!20, draw=none] (-6, 0, -6) -- (14, 0, -6) -- (14, 0, \D+5) -- (-6, 0, \D+5) -- cycle;
			% Tường nhà bên phải và Cửa sổ cứu hộ
			\draw[fill=yellow!40!white, draw=gray!60!black, thick] (\W, 0, 0) -- (\W, 0, \D) -- (\W, \H, \D) -- (\W, \H, 0) -- cycle;
			% Khung viền ngoài cửa sổ bên phải
			\draw[fill=gray!30, draw=gray!60!black, thick] (\W, 4.15, 1.3) -- (\W, 4.15, 3.7) -- (\W, 5.3, 3.7) -- (\W, 5.3, 1.3) -- cycle;
			% Lớp kính cửa sổ
			\draw[fill=cyan!20, draw=gray!60!black, thick] (\W, 4.25, 1.5) -- (\W, 4.25, 3.5) -- (\W, 5.2, 3.5) -- (\W, 5.2, 1.5) -- cycle;
			% Thanh chia kính
			\draw[gray!60!black, thick] (\W, 4.725, 1.5) -- (\W, 4.725, 3.5);
			\draw[gray!60!black, thick] (\W, 4.25, 2.5) -- (\W, 5.2, 2.5);
			% Tường nhà mặt trước và Cửa chính
			\draw[fill=yellow!20!white, draw=gray!60!black, thick] (0, 0, 0) -- (\W, 0, 0) -- (\W, \H, 0) -- (0, \H, 0) -- cycle;
			\draw[fill=gray!30, draw=gray!60!black, thick] (\Pc-0.8, 0, 0) rectangle (\Pc+0.8, 2.5, 0);
			% --- CODE 2 CỬA SỔ MẶT TRƯỚC DUY THÊM VÀO ---
			% Cửa sổ bên trái
			% Khung viền ngoài
			\draw[fill=gray!30, draw=gray!60!black, thick] (0.6, 0.7, 0) rectangle (2.1, 2.3, 0);
			% Lớp kính
			\draw[fill=cyan!20, draw=gray!60!black, thick] (0.7, 0.8, 0) rectangle (2.0, 2.2, 0);
			% Thanh chia kính
			\draw[gray!60!black, thick] (1.35, 0.8, 0) -- (1.35, 2.2, 0);
			\draw[gray!60!black, thick] (0.7, 1.5, 0) -- (2.0, 1.5, 0);
			% Cửa sổ bên phải (mặt trước)
			% Khung viền ngoài
			\draw[fill=gray!30, draw=gray!60!black, thick] (4.9, 0.7, 0) rectangle (6.4, 2.3, 0);
			% Lớp kính
			\draw[fill=cyan!20, draw=gray!60!black, thick] (5.0, 0.8, 0) rectangle (6.3, 2.2, 0);
			% Thanh chia kính
			\draw[gray!60!black, thick] (5.65, 0.8, 0) -- (5.65, 2.2, 0);
			\draw[gray!60!black, thick] (5.0, 1.5, 0) -- (6.3, 1.5, 0);
			% ---------------------------------------------
			% Mái nhà
			\draw[fill=gray!60!black, draw=gray!20!black, thick] (0, \H, 0) -- (\W, \H, 0) -- (\W, \H+1.5, \D/2) -- (0, \H+1.5, \D/2) -- cycle;
			\draw[fill=gray!70!black, draw=gray!20!black, thick] (\W, \H, 0) -- (\W, \H, \D) -- (\W, \H+1.5, \D/2) -- cycle;
			% Cây cối (Bụi cây)
			\shade[ball color=green!50!black] (\W*0.9, 0.5, -0.5) circle (1.2cm);
			\shade[ball color=green!40!black] (\W*0.1, 0.5, -0.5) circle (1.0cm);
			\shade[ball color=green!45!black] (\xL+0.5, 0.5, -\Pd) circle (0.8cm);
			\shade[ball color=green!45!black] (\xR-0.5, 0.5, -\Pd) circle (0.8cm);
			\shade[ball color=green!35!black] (-1, 0.5, \D*0.2) circle (1.0cm);
			% Hàng rào bên trái
			\draw[fill=gray!20, draw=gray] (-3, 0, -3) -- (-3, 0, \D+3) -- (-3, 0.3, \D+3) -- (-3, 0.3, -3) -- cycle;
			\foreach \i in {0, 1, ..., 20} {
				\pgfmathsetmacro{\pz}{-3 + \i*(\D+6)/20}
				\draw[fill=white, draw=gray!50] (-3, 0.3, \pz-0.05) -- (-3, 0.3, \pz+0.05) -- (-3, 1.2, \pz+0.05) -- (-3, 1.2, \pz-0.05) -- cycle;
			}
			\draw[fill=gray!20, draw=gray] (-3, 1.0, -3) -- (-3, 1.0, \D+3) -- (-3, 1.1, \D+3) -- (-3, 1.1, -3) -- cycle;
			% Hàng rào bên phải
			\draw[fill=gray!20, draw=gray] (\W+2, 0, -3) -- (\W+2, 0, \D+3) -- (\W+2, 0.3, \D+3) -- (\W+2, 0.3, -3) -- cycle;
			\foreach \i in {0, 1, ..., 20} {
				\pgfmathsetmacro{\pz}{-3 + \i*(\D+6)/20}
				\draw[fill=white, draw=gray!50] (\W+2, 0.3, \pz-0.05) -- (\W+2, 0.3, \pz+0.05) -- (\W+2, 1.2, \pz+0.05) -- (\W+2, 1.2, \pz-0.05) -- cycle;
			}
			\draw[fill=gray!20, draw=gray] (\W+2, 1.0, -3) -- (\W+2, 1.0, \D+3) -- (\W+2, 1.1, \D+3) -- (\W+2, 1.1, -3) -- cycle;
			% Hàng rào mặt trước (hai bên cổng)
			\draw[fill=gray!10, draw=gray] (\W+2, 0, -3) -- (\Pc+1.5, 0, -3) -- (\Pc+1.5, 0.3, -3) -- (\W+2, 0.3, -3) -- cycle;
			\draw[fill=gray!10, draw=gray] (\Pc-1.5, 0, -3) -- (-3, 0, -3) -- (-3, 0.3, -3) -- (\Pc-1.5, 0.3, -3) -- cycle;
			\foreach \i in {0, 1, ..., 15} {
				\pgfmathsetmacro{\px}{\W+2 - \i*(\W+2 - (\Pc+1.5))/15}
				\draw[fill=white, draw=gray!50] (\px+0.05, 0.3, -3) -- (\px-0.05, 0.3, -3) -- (\px-0.05, 1.2, -3) -- (\px+0.05, 1.2, -3) -- cycle;
			}
			\foreach \i in {0, 1, ..., 15} {
				\pgfmathsetmacro{\px}{\Pc-1.5 - \i*(\Pc-1.5 - (-3))/15}
				\draw[fill=white, draw=gray!50] (\px+0.05, 0.3, -3) -- (\px-0.05, 0.3, -3) -- (\px-0.05, 1.2, -3) -- (\px+0.05, 1.2, -3) -- cycle;
			}
			\draw[fill=gray!20, draw=gray] (\W+2, 1.0, -3) -- (\Pc+1.5, 1.0, -3) -- (\Pc+1.5, 1.1, -3) -- (\W+2, 1.1, -3) -- cycle;
			\draw[fill=gray!20, draw=gray] (\Pc-1.5, 1.0, -3) -- (-3, 1.0, -3) -- (-3, 1.1, -3) -- (\Pc-1.5, 1.1, -3) -- cycle;
			% Cổng rào
			\draw[fill=gray!90!black, draw=black, thick] (\Pc+1.4, 0.1, -3) -- (\Pc-1.4, 0.1, -3) -- (\Pc-1.4, 1.4, -3) .. controls (\Pc, 1.7, -3) .. (\Pc+1.4, 1.4, -3) -- cycle;
			\foreach \bar in {1.2, 0.9, 0.6, 0.3, 0, -0.3, -0.6, -0.9, -1.2} {
				\pgfmathsetmacro{\bpos}{\Pc+\bar}
				\draw[white!80!gray, thick] (\bpos, 0.1, -3) -- (\bpos, 1.4, -3);
			}
			% Các cột trụ hàng rào
			\foreach \gx in {\Pc+1.5, \Pc-1.5, -3, \W+2} {
				\draw[fill=gray!20, draw=gray] (\gx+0.3, 0, -3-0.3) -- (\gx-0.3, 0, -3-0.3) -- (\gx-0.3, 1.6, -3-0.3) -- (\gx+0.3, 1.6, -3-0.3) -- cycle;
				\draw[fill=gray!30, draw=gray] (\gx-0.3, 0, -3-0.3) -- (\gx-0.3, 0, -3+0.3) -- (\gx-0.3, 1.6, -3+0.3) -- (\gx-0.3, 1.6, -3-0.3) -- cycle;
				\draw[fill=gray!40, draw=gray] (\gx+0.4, 1.6, -3-0.4) -- (\gx-0.4, 1.6, -3-0.4) -- (\gx-0.4, 1.7, -3+0.4) -- (\gx+0.4, 1.7, -3+0.4) -- cycle;
				\fill (\gx, 1.75, -3) circle (0.12cm);
				\fill[yellow!60!white] (\gx, 1.8, -3) circle (0.08cm);
			}
			% Chiếc thang cứu hộ
			\path
			(9,0) coordinate (A)
			+(-150:0.4) coordinate (A')
			(8,5) coordinate (B)
			($(A')+(B)-(A)$) coordinate (B')
			;
			\foreach \x in {1,...,10}{
				\pgfmathsetmacro{\k}{0.01+0.09*\x}
				\draw[gray!30!black, line width=1pt]($(A)!\k!(B)$)--+(-150:0.4);
			}
			\draw[gray!30!black, line width=1.5pt] (A')--(B') (A)--(B);
			% \node at (0,0){\href{2VFHU2}{\tiny \phantom{2VFHU2}}};
		\end{tikzpicture}
	\end{center}
	\par\shortans[0]{5{,}35}
	\loigiai{
		\immini[thm]{Gọi góc hợp bởi thang và mặt đất là $\widehat{CBA}=\alpha$ $\left(0 < \alpha < \dfrac{\pi}{2}\right)$.\\
		Xét $\triangle MHB$ vuông tại $H$, ta có $BM = \dfrac{2{,}2}{\sin \alpha}$.\\
		Xét $\triangle MNC$ vuông tại $N$, ta có $MC = \dfrac{1{,}6}{\cos \alpha}$.\\
		Khi đó, chiều dài chiếc thang là
		\[BC = BM + MC = \dfrac{2{,}2}{\sin \alpha} + \dfrac{1{,}6}{\cos \alpha}.\]
		Xét hàm số $f(x) = \dfrac{2{,}2}{\sin x} + \dfrac{1{,}6}{\cos x}$, với $x \in \left(0; \dfrac{\pi}{2}\right)$.\\
		Ta có $f'(x) = \dfrac{-2{,}2 \cdot \cos x}{\sin^2 x} + \dfrac{1{,}6 \cdot \sin x}{\cos^2 x} = \dfrac{-2{,}2 \cdot \cos^3 x + 1{,}6 \cdot \sin^3 x}{\sin^2 x \cdot \cos^2 x}$.}
		{\begin{tikzpicture}[scale=1, font=\footnotesize, line join=round, line cap=round, >=stealth]
				\path 
				(0,0) coordinate (A)
				(0,4) coordinate (C)
				(3,0) coordinate (B)
				($(B)!3/5!(C)$) coordinate (M)
				($(B)!(M)!(A)$) coordinate (H)
				($(A)!(M)!(C)$) coordinate (N)
				;
				\draw (A)--(B)--(C)--cycle (N)--(M)--(H);
				\path (M)--(H) node[pos=0.5,above,sloped]{$2{,}2$\,m};
				\path (A)--(H) node[pos=0.5,below]{$1{,}6$\,m};
				\foreach \x/\pos in {A/-135,H/-80,B/-45,M/45,C/90,N/180} \fill (\x) circle(1pt) node[{shift=(\pos:0.25)}]{$\x$}; 
				\pic[draw,angle radius=3mm] {angle = C--B--A}; 
				\node[shift=(150:0.5)] at (B) {$\alpha$};
				% \node at (0,0){\href{2VFHU2}{\tiny \phantom{2VFHU2}}};
		\end{tikzpicture}}
	\noindent
		Phương trình 
		\allowdisplaybreaks
		\begin{eqnarray*} 
			f'(x) = 0 &\Leftrightarrow& 2{,}2 \cdot \cos^3 x = 1{,}6 \cdot \sin^3 x  \\ 
			&\Leftrightarrow&  \tan^3 x = \dfrac{2{,}2}{1{,}6} \\ 
			&\Leftrightarrow&  \tan x = \sqrt[3]{\dfrac{11}{8}} .
		\end{eqnarray*}
		Chọn $x_0 \in \left(0; \dfrac{\pi}{2}\right)$ thỏa mãn $\tan x = \sqrt[3]{\dfrac{11}{8}}$.\\
		Ta có bảng biến thiên sau
		\begin{center}
			\begin{tikzpicture}
				\tkzTabInit[nocadre=false, lgt=1.2, espcl=4, deltacl=0.5]
				{$x$/0.7,$f'(x)$/0.7,$f(x)$/2}
				{$0$,$x_0$,$\tfrac{\pi}{2}$}
				\tkzTabLine{,-,0,+,}
				\tkzTabVar{+/$+\infty$,-/$f\left(x_0\right)$,+/$+\infty$}
				% \node at (T11){\href{2VFHU2}{\tiny \phantom{2VFHU2}}};
			\end{tikzpicture}
		\end{center}
		Dựa vào bảng biến thiên, ta thấy chiều dài thang bé nhất là $f\left(x_0\right)\approx 5{,}35$ mét.
	}
\end{ex}

\begin{ex}%[2D1H3-6]%[TN-THPT-NguyenTrungThien-HaTinh-NH25-26]%[Loc Do]%[PB: Phan Anh]
	Một đại lý nhập khẩu trái cây tươi để phân phối cho các cửa hàng. Mỗi lần nhập khẩu trái cây, khoán chi phí vận chuyển (không đổi) là $25$ triệu đồng. Chi phí bảo quản mỗi tạ trái cây dự trữ trong kho là $80$ nghìn đồng/ngày. Thời gian bảo quản trái cây trong kho tối đa $10$ ngày. Biết rằng, kể từ ngày đầu tiên nhập hàng, đại lý sẽ phân phối tới các cửa hàng $25$ tạ trái cây mỗi ngày. Mỗi lần nhập hàng, đại lý phải nhập đủ trái cây cho bao nhiêu ngày phân phối để chi phí trung bình cho mỗi ngày là thấp nhất (bao gồm chi phí vận chuyển và chi phí bảo quản trong kho)?
	\shortans{5}
	\loigiai{
		Giả sử mỗi lần nhập hàng, đại lý nhập đủ trái cây phân phối cho $n$ ngày $(1\le n\le 10)$.\\
		Chi phí vận chuyển của mỗi lần nhập hàng là $25$ triệu.\\
		Chi phí lưu kho là $25\cdot80\cdot(1+2+\ldots+n)=1\,000(n^{2} +n)$ nghìn đồng bằng $n^{2} +n$ triệu đồng.\\
		Vậy chi phí trung bình của mỗi ngày là $f(n)=\dfrac{n^{2} +n+25}{n} =n+1+\dfrac{25}{n}$.
		\[f'(n)=0\Leftrightarrow 1-\dfrac{25}{n^{2}} =0\Leftrightarrow n=5.\]
		Ta có $f(1)=27$; $f(5)=11$; $f(10)=13{,}5$.\\
		Vậy chi phí hàng ngày nhỏ nhất khi $n=5$.}
\end{ex}

\begin{ex}%[2D1K3-6]%[Dự án C-TT THPT NH25-26-Đợt 4- ĐỖ CHÍ TÂM]%[THPT Phan Bội Châu-Khánh Hòa]%[GVPB: Phạm Văn Long]
	Một tàu chở hàng di chuyển trên sông với vận tốc $v$ (km/giờ), $v>0$. Biết chi phí nhiên liệu cho mỗi kilômét là $0{,}3\sqrt{v}+\dfrac{1{,}6}{\sqrt{v}}$ (triệu đồng/km); chi phí bảo trì tính theo thời gian là $0{,}04v\sqrt{v}$ (triệu đồng/giờ). Ngoài ra, chi phí vận hành cố định khác (nhân công, bảo hiểm\ldots) là $17{,}25$ (triệu đồng/giờ). Gọi $v_0$ là vận tốc để tổng chi phí trên mỗi kilômét là nhỏ nhất và $C_{min}$ là giá trị chi phí thấp nhất đó. Tính giá trị biểu thức $P=v_0+100\cdot C_{min}$.
	\shortans{296}
	\loigiai{
		Chi phí nhiên liệu cho mỗi km là $0{,}3\cdot \sqrt{v}+\dfrac{1{,}6}{\sqrt{v}}$ (triệu đồng/km).\\
		Tàu chở hàng di chuyển với vận tốc $v$ (km/giờ) nên tàu đi được $1$ (km) trong $\dfrac{1}{v}$ (giờ).\\
		Vậy chi phí bảo trì trên $1$ (km) hay trong $\dfrac{1}{v}$ (giờ) là $0{,}04v\sqrt{v}\cdot\dfrac{1}{v}=0{,}04\sqrt{v}$ (triệu đồng).\\
		Chi phí vận hành cố định khác trên $1$ (km) hay trong $\dfrac{1}{v}$ (giờ) là $\dfrac{17{,}25}{v}$ (triệu đồng).\\
		Vậy tổng chi phí trên mỗi km của tàu chở hàng là\\
		$C\left(v \right)=0,3\sqrt{v}+\dfrac{1,6}{\sqrt{v}}+0{,}04\sqrt{v}+\dfrac{69}{4}\cdot \dfrac{1}{v}=0{,}34\sqrt{v}+\dfrac{1,6}{\sqrt{v}}+\dfrac{69}{4}\cdot \dfrac{1}{v}$\\
		$C'\left(v \right)=\dfrac{17}{100\sqrt{v}}-\dfrac{4}{5\sqrt{{v^3}}}-\dfrac{69}{4{v^2}}$.\\
		Xét $C'\left(v \right)=0$, ta đặt $\dfrac{1}{\sqrt{v}}=t\left(t{>}0\right)$ $\Rightarrow-\dfrac{69}{4}{t^4}-\dfrac{4}{5}{t^3}+\dfrac{17}{100}t=0\Leftrightarrow \hoac{&t=0\,\text{(loại)}\\&t=\dfrac{1}{5}\Rightarrow v=25.\\
		}$\\
		Ta lập được bảng biến thiên
		\begin{center}
			\begin{tikzpicture}
				\tkzTabInit[nocadre=true,lgt=1.2,espcl=2.5,deltacl=0.6]
				{$v$ /0.6,$C'(v)$ /0.6,$C(v)$ /2}
				{$0$,$25$,$+\infty$}
				\tkzTabLine{,-,$0$,+,}
				\tkzTabVar{+/, -/$2{,}71$,+/}
			\end{tikzpicture}
		\end{center}
		Vậy ${C_{\min}}=2{,}71\Leftrightarrow {v_0}=25\Rightarrow P=25+100\cdot 2{,}71=296$.
	}
\end{ex}

\begin{ex}%[2D1V2-1]%[Dự án C đề thi thử TN Môn Toán NH25-26 - Dot 4 - Vũ Hồng Toàn]%[THPT Mai Anh Tuấn - Thanh Hóa]%[GVPB - Lê Hữu Kiệt]
	Cho đồ thị hàm số $y = x^4 -3x^2 + ax+b$ có $A(2;-2)$ là một điểm cực tiểu. Giá trị của biểu thức $S = 59b-a$ bằng bao nhiêu?
	\par\shortans{2026}
	\loigiai{
		Ta có $y' = 4x^3-6x+a$, $y''=12x^2 -6$. \\
		Đồ thị hàm số $y = x^4 -3x^2 +ax+b$ có $A(2;-2)$ là một điểm cực tiểu suy ra
		\allowdisplaybreaks
		\begin{eqnarray*}
			\heva{&y(2) = -2\\ &y'(2) = 0\\ &y''(2) = 42 > 0}\Leftrightarrow \heva{&2a+b=-6\\ &a+20=0}\Leftrightarrow \heva{&a=-20\\ &b=34.}
		\end{eqnarray*}
		Vậy $S = 59b- a = 59 \cdot 34 + 20 = 2\,026$.
	}
\end{ex}

\begin{ex}%[2D1V2-7]%[TN-THPT-NguyenTrungThien-HaTinh-NH25-26]%[Loc Do]%[PB: Phan Anh]
	Một phần đường chạy của tàu lượn siêu tốc  khi gắn hệ trục tọa độ $Oxy$ được mô phỏng ở hình bên dưới, đơn vị trên mỗi trục là mét. Biết đường chạy của nó là một phần đồ thị hàm số bậc ba $y=ax^{3} +bx^{2} +cx+d$ $(0\le x<90)$; tàu lượn siêu tốc xuất phát từ điểm $A$, đi qua các điểm $C,\, D$ (ba điểm $A$, $C$, $D$ nằm trên đường thẳng song song với trục $Ox$) đồng thời đạt độ cao nhỏ nhất so với mặt đất là $4$ m. Độ cao lớn nhất mà tàu lượn siêu tốc đạt được là bao nhiêu mét so với mặt đất (\textit{kết quả làm tròn đến hàng phần chục})?
	\begin{center}
		\begin{tikzpicture}[x=1cm,y=1cm,scale=1,>=stealth]
			\def\aa{-0.08} \def\bb{1.04} \def\cc{-3.2} \def\dd{3}
			\draw[->] (-1.5,0)--(8.84,0)node[below right]{$x$};
			\draw[->] (0,-1)--(0,4.22)node[left]{$y$};
			\fill (0,0)node[below left]{$O$};
			\draw[black,samples=150,smooth,domain=-0.36:8.7] plot(\x,{\aa*(\x)^3+\bb*(\x)^2+\cc*\x+\dd});
			\draw [dashed] (0,3)--(8,3)--(8,0) (5,3)--(5,0) (0,0.1)--(2,0.1)--(2,0)--(2,0.2);
			\foreach \x/\g in {2/-90,5/-90,8/-90 }{
				\fill (\x,0) circle (1pt) +(\g:3mm) node[scale=0.8] {$\x0$};
			}
			
			\foreach \y/\g in {3/180 }{
				\fill (0,\y) circle (1pt) +(\g:3mm) node[scale=0.8] {$\y0$};
			}
			\foreach \x/\y/\g in {5/3/-30,8/3/0}{
				\fill (\x,\y) circle (1.5pt);
			}
			\fill (0,0.1) circle (1pt) node[left,scale=0.8] {$4$};
			\node[above right] at (0,3) {$A$};
			\node[above left] at (5,3) {$C$};
			\node[above right] at (8,3) {$D$};
		\end{tikzpicture}
	\end{center}
	\shortans{40,7}
	\loigiai{
		Nhận xét:  ta thấy đồ thị của hàm số $y=f(x)=ax^{3} +bx^{2} +cx+d\,(0\le x<90)$ cắt đường thẳng $d\colon y=30$ tại $3$ điểm phân biệt: $A(0;30)$, $C(50;30)$, $D(80;30)$.\\
		Do đó, phương trình $f(x)=30$ có $3$ nghiệm phân biệt $x_{1} =0$; $x_{2} =50$; $x_{3} =80$ hay ta có
		\begin{eqnarray*}
			&& f(x)-30=ax(x-50)(x-80) \\
			&\Leftrightarrow&  f(x)=ax(x-50)(x-80)+30\\
			&\Leftrightarrow&  f(x)=ax(x^{2} -130x+4\,000)+30\\
			&\Leftrightarrow & f(x)=a(x^{3} -130x^{2} +4\,000x)+30\\
			&\Rightarrow &  f'(x)=a(3x^{2} -260x+4\,000).
		\end{eqnarray*}
		Cho $f'(x)=0\Leftrightarrow a\left(3x^{2} -260x+4\,000\right)=0\Leftrightarrow \hoac{& x=20\\ & x=\dfrac{200}{3}.}$\\
		Ta có
		\begin{itemize}
			\item Hàm số $y=f(x)$ đạt cực tiểu tại $x=20$, giá trị cực tiểu là $f(20)=4$.
			\item Hàm số $y=f(x)$ đạt cực đại tại $x=\dfrac{200}{3}$, giá trị cực đại là $f\left(\dfrac{200}{3} \right)$.
		\end{itemize}
		Do $f(20)=4$, suy ra $a\cdot 20\cdot(20-50)\cdot(20-80)+30=4\Rightarrow a=-\dfrac{13}{18\,000}$.\\
		Giá trị cực đại của hàm số $y=f(x)$ là 
		\[f\left(\dfrac{200}{3} \right)=-\dfrac{13}{18\,000} \cdot\dfrac{200}{3} \cdot\left(\dfrac{200}{3} -50\right)\cdot\left(\dfrac{200}{3} -80\right)+30=\dfrac{9\,890}{243} \approx 40{,}7.\]
		Vậy độ cao lớn nhất mà tàu lượn siêu tốc đạt được là khoảng $40{,}7$ mét so với mặt đất.}
\end{ex}

\begin{ex}%[2D1V3-6]%[TT Cụm 5 - Ninh Bình-NH25-26]%[GVSB:Dương Công Tạo]%[GVPB1: Nguyễn Tấn Tài]
	Một nhà máy A chuyên sản xuất một loại sản phẩm cho nhà máy B, nhà máy A chỉ bán sản phẩm cho nhà máy B và nhà máy B cam kết thu mua hết số sản phẩm mà nhà máy A sản xuất được. Nhà máy A có khả năng sản xuất tối đa là $200$ tấn sản phẩm trong một tháng. Nếu bán ra $x$ tấn sản phẩm cho nhà máy B thì giá bán mỗi tấn sản phẩm là $50-0{,}0002x^2$ triệu đồng. Trong một tháng nhà máy A phải chi phí cho nhân công và chi phí khấu hao máy móc một lượng cố định là $150$ triệu đồng, ngoài ra khi sản xuất mỗi tấn sản phẩm thì nhà máy phải chi thêm cho mua nguyên vật liệu là $35$ triệu đồng. Biết rằng nhà máy A phải nộp $5\%$ doanh thu cho cơ quan thuế. Tính lợi nhuận sau thuế (lợi nhuận khi đã trừ thuế) lớn nhất thu được trong một tháng của nhà máy A (đơn vị tính là tỉ đồng và kết quả làm tròn đến hàng phần trăm).\\
	\par  \shortans{1,08}
	\loigiai{
		Doanh thu của nhà máy A trong một tháng là \[D(x)=\left(50-0{,}0002x^2\right)x = 50x-0{,}0002x^3\text{ (triệu đồng)}.\] 
		Tiền thuế nhà máy A phải nộp là $T(x)= D(x)\cdot 5\%$ (triệu đồng).\\
		Chi phí sản xuất $C(x)=150+35x$ (triệu đồng).\\
		Lợi nhuận sau thuế trong một tháng của nhà máy A là 
		\begin{eqnarray*}
			L(x)&=& D(x)- D(x)\cdot 5\%-C(x) \\
			&=& 0{,}95\cdot (50x-0{,}0002x^3)-(150+35x)\\
			&=& -0{,}00019x^3 +12{,}5x-150\text{ (triệu đồng)}.
		\end{eqnarray*}
		Ta có $L'(x) = -0{,}00057x^2 +12{,}5$; $L'(x) = 0 \Leftrightarrow x = \pm \sqrt{\dfrac{12{,}5}{0{,}00057}}\approx 148$.\\
		Bảng biến thiên
		\begin{center}
			\begin{tikzpicture}[>=stealth]
				\tkzTabInit[nocadre=false,lgt=1.5,espcl=3,deltacl=0.6]{$x$/.7 ,$L'(x)$/.7,$L(x)$/2}
				{$0$ , $148$ , $200$}
				\tkzTabLine{ , + , $0$ , - , }
				\tkzTabVar{-/$-150$ , +/$1\,084$ , -/$830$}
			\end{tikzpicture}
		\end{center}
		Lợi nhuận sau thuế lớn nhất trong một tháng của nhà máy A là \[\max L(x) = L(148)\approx 1\,084\text{ (triệu đồng)}\approx 1{,}08\text{ (tỉ đồng)}.\] 
	}
\end{ex}

\begin{ex}%[2D1V3-6] %[Dự án C đề thi thử TN Môn Toán NH25-26 - Dot 4 - Hector Tran]%[Liên trường THPT cụm 09 - Đợt 2 - Hà Nội]%[GVPB - Thanh Phát]
	Một doanh nghiệp sản xuất độc quyền một loại sản phẩm. Giả sử khi sản xuất và bán hết $x$ sản phẩm ($0 < x \le 2\,500$), tổng số tiền doanh nghiệp thu được là $f(x) = 2\,026x - x^2$ và tổng chi phí là $g(x) = x^2 + 1\,438x - 1\,209$ (đơn vị: nghìn đồng). Giả sử mức thuế phụ thu trên một đơn vị sản phẩm bán được là $t$ (nghìn đồng) ($0 < t < 320$). Giá trị của $t$ bằng bao nhiêu nghìn đồng để nhà nước nhận được số tiền thuế phụ thu lớn nhất và doanh nghiệp cũng nhận được lợi nhuận lớn nhất theo mức thuế phụ thu đó?
	\par\shortans{294}
	\loigiai{
		Ta có hàm lợi nhuận khi sản xuất và bán hết $x$ sản phẩm là
		\begin{eqnarray*}
			P(x) &=& f(x) - g(x) - xt \\
			&=& 2\,026x - x^2 - (x^2 + 1\,438x - 1\,209) - xt \\
			&=& -2x^2 + (588 - t)x + 1\,209\text { (nghìn đồng).}
		\end{eqnarray*}
		Khi lợi nhuận lớn nhất thì $x = \dfrac{-b}{2a} = \dfrac{588 - t}{4}$.\\
		Số tiền thuế thu được là $T = xt = \dfrac{588t - t^2}{4}=-\dfrac{1}{4}t^2+147t$ (nghìn đồng).\\
		Số tiền thuế lớn nhất khi $t = \dfrac{-147}{2 \cdot \left( -\dfrac{1}{4}\right) } = 294\text { (nghìn đồng) } \in (0; 320)$ (thỏa mãn).\\
		Vậy $t=294$ (nghìn đồng).
	}
\end{ex}

\begin{ex}%[2D1V3-6]%[Dự án đề TT Toán NH25-26-Dot 3- Hải Phụng]%[Cụm 13-Hải Phòng]%[GVPB - Viet Hoang Pham]
Để hạn chế vi phạm thời gian làm việc đối với công nhân, giám đốc công ty quyết định xử lý bằng cách phạt tiền. Nhờ sự giám sát chặt chẽ của quản đốc, giám đốc công ty biết được trong một tháng, giữa tỉ lệ công nhân vi phạm đúng $k$ lần ($1 \leq k \leq 2$) là $t_k = \dfrac{N_k}{N}$ (trong đó $N_k$ là số công nhân vi phạm đúng $k$ lần, $N$ là tổng số công nhân) và mức phạt mỗi lần vi phạm có mối liên hệ như sau: Nếu mỗi công nhân nộp phạt $x$ nghìn đồng ($60 \leq x \leq 300$) khi vi phạm lần thứ nhất và nộp phạt $x - 20$ nghìn đồng khi vi phạm lần thứ hai thì $t_1 = \dfrac{36}{x + 10}$ và $t_2 = \dfrac{4}{x - 30}$ (không có công nhân nào vi phạm quá hai lần). Biết rằng $N$ không đổi và bằng $2\,400$. Tổng số tiền nộp phạt trong một tháng ít nhất là bao nhiêu triệu đồng? (kết quả làm tròn đến hàng đơn vị)
\shortans[oly]{$103$}
\loigiai{
Tổng số công nhân $N = 2\,400$.\\
Số công nhân vi phạm đúng $1$ lần $N_1 = 2\,400 \cdot \dfrac{36}{x+10}$, tỉ lệ $t_1 = \dfrac{36}{x+10}$.\\
Số công nhân vi phạm đúng $2$ lần $N_2 = 2\,400 \cdot \dfrac{4}{x-30}$, tỉ lệ $t_2 = \dfrac{4}{x-30}$.\\
Không có công nhân nào vi phạm quá $2$ lần.\\
Tiền phạt mỗi lần vi phạm
\begin{itemize}
\item Lần 1: $x$ (nghìn đồng), với $60 \le x \le 300$.
\item Lần 2: $x - 20$ (nghìn đồng).
\end{itemize}
Tiền phạt công nhân vi phạm $1$ lần $x$ (nghìn đồng).\\
Tiền phạt công nhân vi phạm $2$ lần $x + (x - 20) = 2x - 20$ (nghìn đồng).\\
Tổng tiền phạt hàm số theo $x$ là
\[f(x) = N_1 \cdot x + N_2 \cdot (2x - 20) = \dfrac{86\,400x}{x+10} + \dfrac{9600(2x-20)}{x-30}.\]
Tính đạo hàm ta được $f'(x) = \dfrac{864\,000}{(x+10)^2} - \dfrac{384\,000}{(x-30)^2}$. Cho $f'(x) = 0 \Rightarrow \left(\dfrac{x+10}{x-30}\right)^2 = \dfrac{9}{4}$.\\
Từ đó tìm được nghiệm $x = 110$ (thỏa mãn điều kiện $x \in [60; 300]$).\\
Ta có $f(110) = 103\,200$, $f(60)\approx 106\,057$, $f(300)\approx 104\,235$.\\
Giá trị nhỏ nhất đạt được tại $x = 110$, khi đó $f(110) = 103\,200$ (nghìn đồng).\\
Vậy tổng số tiền phạt ít nhất làm tròn đến hàng đơn vị là $103$ triệu đồng.
}
\end{ex}

\begin{ex}%[2D1V3-6]%[Dự án đề thi thử NH25-26]%[THPT Chuyên Phan Bội Châu - Nghệ An]%[BC Tuan]
	Một cửa hàng kinh doanh máy lọc nước Karofi. Hàm cầu biểu thị liên hệ giữa giá bán của một chiếc máy với số lượng máy bán được trong một tháng là hàm bậc nhất. Khi giá bán là $5$ triệu đồng một chiếc thì một tháng bán được $100$ chiếc, khi giá bán là $4{,}5$ triệu đồng một chiếc thì một tháng bán được $120$ chiếc. Biết chi phí trung bình cho một chiếc máy khi bán được $x$ chiếc là $\dfrac{3x+50}{x}$. Hỏi cửa hàng cần bán với giá bao nhiêu triệu đồng một chiếc máy để lợi nhuận thu được trong một tháng lớn nhất?
	\shortans{$5{,}25$}
	\loigiai{
		Gọi hàm biểu thị mối liên hệ giữa giá bán của một chiếc máy với số lượng máy bán được là $p(x)=ax+b$.\\
		Tại $x=100$ thì $p(x)=5$ và tại $x=120$ thì $p(x)=4{,}5$.\\ Từ đó suy ra
		\[\heva{& 5=100a+b \\ & 4{,}5=120a+b}\Leftrightarrow \heva{& a=-0{,}025 \\ & b=7{,}5.}\]
		Vậy hàm số biểu thị giá bán là $p(x)=-0{,}025x+7{,}5$.\\
		Hàm số lợi nhuận là \[L(x) = x \cdot p(x) - \left( \dfrac{3x+50}{x} \right)x = x(-0{,}025x+7{,}5) - (3x+50) = -0{,}025x^2+4{,}5x-50.\]
		Ta có $L'(x)=-0{,}05x+4{,}5$.\\
		Khi đó $L'(x)=0\Leftrightarrow x=90$.\\
		\begin{center}
			\begin{tikzpicture}
				\tkzTabInit[nocadre=true, lgt=1.2, espcl=4, deltacl=0.5]
				{$x$ /0.7, $L'(x)$ /0.7, $L(x)$ /2}
				{$0$,$90$,$+\infty$}
				\tkzTabLine{ ,+,0,-, }
				\tkzTabVar{-/ $-50$, +/ $152{,}5$, -/ $-\infty$}
			\end{tikzpicture}
		\end{center}
		Bảng biến thiên cho thấy $L(x)$ đạt giá trị lớn nhất tại $x=90$.\\
		Vậy để lợi nhuận thu được lớn nhất thì cửa hàng cần bán một cái máy lọc nước với giá $p(90) = -0{,}025\cdot 90 + 7{,}5 = 5{,}25$ (triệu đồng).
	}
\end{ex}

\begin{ex}%[2D1V3-6]%[Dự án đề thi thử NH25-26]%[THPT Chuyên Phan Bội Châu - Nghệ An]%[Phạm Hoàng Đăng]
	\immini[thm]{Một hòn đảo $A$ cách bờ biển $3$ (km) và cách trạm điện $B$ là $5$ (km). Công ty điện lực dự định làm một tuyến đường điện từ trạm điện $B$ dọc theo bờ biển đến một điểm $C$, sau đó từ $C$ nối thẳng ra đảo $A$ (tham khảo hình vẽ). Biết chi phí cho $1$ (km) đường điện trên bờ là $30$ triệu đồng và trên biển là $50$ triệu đồng. Hỏi cần ít nhất bao nhiêu triệu đồng để hoàn thành kế hoạch trên?}{
		\begin{tikzpicture}[scale=0.8, line join=round, line cap=round, >=stealth, font=\footnotesize]
			\coordinate (H) at (0,0);
			\coordinate (A) at (0,3);
			\coordinate (B) at (4,0);
			\coordinate (C) at (2.25,0);
			\draw  (-1,0) -- (5,0) node[below] {Bờ biển};
			\draw[dashed] (A) -- (H) node[midway, left] {$3$ (km)};
			\draw (A) -- (C) node[midway,sloped, above right] {Biển};
			\draw[red] (B) -- (C);
			\draw[blue] (C) -- (A);
			\fill (A) circle (1pt) node[above] {$A$};
			\fill (B) circle (1pt) node[below] {$B$};
			\fill (C) circle (1pt) node[below] {$C$};
			\fill (H) circle (1pt) node[below] {$H$};
			\draw (H) rectangle +(0.2,0.2);
			\draw[dashed] (A) -- (B) node[midway, sloped, above] {$5$ (km)};
		\end{tikzpicture}
	}
	\shortans{$240$}
	\loigiai{
		Gọi $H$ là hình chiếu vuông góc của $A$ lên bờ biển.\\ Ta có $AH = 3$ (km).\\
		Tam giác $AHB$ vuông tại $H$ nên $HB = \sqrt{AB^2 - AH^2} = \sqrt{5^2 - 3^2} = 4$ (km).\\
		Đặt $CH = x$ (km) với $0 \le x \le 4$.\\ Suy ra quãng đường trên bờ là $BC = 4 - x$ (km).\\
		Khoảng cách từ điểm $C$ trên bờ nối thẳng ra đảo $A$ là $AC = \sqrt{CH^2 + AH^2} = \sqrt{x^2 + 9}$ (km).\\
		Hàm tổng chi phí làm đường điện là $f(x) = 30(4-x) + 50\sqrt{x^2+9}$ (triệu đồng).\\
		Ta có $f'(x) = -30 + \dfrac{50x}{\sqrt{x^2+9}}$.\\
		Khi đó 
		\begin{eqnarray*}
			&&f'(x) = 0 \\
			&\Leftrightarrow& 50x = 30\sqrt{x^2+9}\\ 
			&\Leftrightarrow& 25x^2 = 9(x^2+9)\\ 
			&\Leftrightarrow& 16x^2 = 81\\
			&\Leftrightarrow& x = \dfrac{9}{4} = 2{,}25.
		\end{eqnarray*}
		Bảng biến thiên
		\begin{center}
			\begin{tikzpicture}
				\tkzTabInit[nocadre=true, lgt=1.2, espcl=3, deltacl=0.5]
				{$x$ /0.8, $f'(x)$ /0.8, $f(x)$ /2.5}
				{$0$,$2{,}25$,$4$}
				\tkzTabLine{ ,-,0,+, }
				\tkzTabVar{+/ $270$, -/ $240$, +/ $250$}
			\end{tikzpicture}
		\end{center}
		Bảng biến thiên cho thấy hàm số $f(x)$ đạt giá trị nhỏ nhất tại $x = 2{,}25$.\\
		Chi phí ít nhất là $f(2{,}25) = 30(4 - 2{,}25) + 50\sqrt{2{,}25^2 + 9} = 52{,}5 + 187{,}5 = 240$ (triệu đồng).
	}
\end{ex}

\begin{ex}%[2D1V3-6]%[Dự án đề thi thử đợt 3-Nguyễn Đức Lợi]%[SGD tỉnh Lạng Sơn]%[GVPB - ThanhTa]
	Một nhà máy sản xuất $x$ sản phẩm trong mỗi tháng. Chi phí sản xuất $x$ sản phẩm được cho bởi hàm chi phí $C(x) = 16\,000 + 500x - 1{,}6x^2 + 0{,}004x^3$ (nghìn đồng). Biết giá bán của mỗi sản phẩm là một hàm số phụ thuộc vào số lượng sản phẩm $x$ và được cho bởi công thức \break$p(x) = 1\,700 - 7x$ (nghìn đồng). Hỏi mỗi tháng nhà máy nên sản xuất bao nhiêu sản phẩm để lợi nhuận thu được là lớn nhất? Biết rằng kết quả khảo sát thị trường cho thấy sản phẩm sản xuất ra sẽ được tiêu thụ hết.
	\par\shortans{100}
	\loigiai{
		Hàm doanh thu của nhà máy là
		$R(x)=x\cdot p(x)=x(1\,700-7x)=1\,700x-7x^2$ (nghìn đồng).\\
		Hàm lợi nhuận thu được là 
		\begin{align*}
			L(x)&=R(x)-C(x)\\
			&=\left(1\,700x-7x^2\right)-\left(16\,000+500x-1{,}6x^2+0{,}004x^3\right)\\
			&=-0{,}004x^3-5{,}4x^2+1\,200x-16\,000.
		\end{align*}
		Ta có $L'(x)=-0{,}012x^2-10{,}8x+1\,200$.\\
		Cho $L'(x)=0\Leftrightarrow-0{,}012x^2-10{,}8x+1\,200=0\Leftrightarrow\hoac{&x= 100\text{ (nhận)}\\&x=-1\,000\text{ (loại)}.}$\\
		Bảng biến thiên của hàm số $L(x)$ trên khoảng $(0;+\infty)$.
		\begin{center}
			\begin{tikzpicture}    
				\tkzTabInit[nocadre=false, lgt=1.2,deltacl=0.5 ,espcl=4]
				{$x$/0.7, $L'(x)$/0.7, $L(x)$/2}
				{$0$,$100$,$+\infty$}
				\tkzTabLine{,+,0,-,}
				\tkzTabVar{-/,+/$L(100)$,-/}
			\end{tikzpicture}
		\end{center}
		Dựa vào bảng biến thiên, hàm số $L(x)$ đạt giá trị lớn nhất tại $x=100$.\\
		Vậy mỗi tháng nhà máy nên sản xuất $100$ sản phẩm để lợi nhuận thu được là lớn nhất.
	}
\end{ex}

\begin{ex}%[2D1V3-6]%[Dự án C NH25-26 - Đợt 3 - Hector Tran]%[SGD-Tuyên Quang - Lần 1]%[GVPB: Thanh Phát]
	Một khu chung cư có $120$ căn hộ cho thuê. Người quản lí của khu chung cư nhận thấy rằng nếu giá thuê một căn hộ là $7$ triệu đồng một tháng thì tất cả các căn hộ đều sẽ có người thuê. Một cuộc khảo sát thị trường cho thấy, trung bình cứ mỗi lần tăng giá thuê một căn hộ mỗi tháng thêm $250$ nghìn đồng thì sẽ có thêm ba căn hộ bị bỏ trống. Người quản lí nên đặt giá thuê mỗi căn hộ là bao nhiêu triệu đồng một tháng để doanh thu một tháng là lớn nhất?
	\par\shortans{8,5}
	\loigiai{
		Gọi $x$ là số lần tăng giá thuê $250$\ nghìn đồng ($x \ge 0$).\\
		Khi đó, giá thuê mỗi căn hộ một tháng là $7 + 0{,}25x$ (triệu đồng).\\
		Số căn hộ có người thuê là $120 - 3x$ (căn hộ).\\
		Doanh thu một tháng của khu chung cư là
		\begin{eqnarray*}
			R(x) &=& (7 + 0{,}25x)(120 - 3x)\\
			&=& 840 - 21x + 30x - 0{,}75x^2\\
			&=& -0{,}75x^2 + 9x + 840.
		\end{eqnarray*}
		Đây là một hàm số bậc hai với hệ số $a = -0{,}75 < 0$, nên hàm số đạt giá trị lớn nhất tại hoành độ đỉnh parabol
		\[x = -\dfrac{b}{2a} = -\dfrac{9}{2 \cdot (-0{,}75)} = 6.\]
		Khi đó, giá thuê mỗi căn hộ để doanh thu lớn nhất là $7 + 0{,}25 \cdot 6 = 8{,}5$ (triệu đồng).
	}
\end{ex}

\begin{ex}%[TT-THPT-HamRong-NH25-26 ]%[GV soạn: Nguyễn Quang Hiệp GVPB2 Đỗ Minh Vũ ]%[2D1V3-6]
	Trong một thí nghiệm y học, người ta cấy $1\,000$ vi khuẩn vào môi trường dinh dưỡng. Bằng thực nghiệm, người ta xác định được số lượng vi khuẩn thay đổi theo thời gian bởi công thức $N(t)=1\,000+\dfrac{100t}{100+t^{2}}$ (con), trong đó $t$ là thời gian tính bằng giây ($t \ge 0$). Tính số lượng vi khuẩn lớn nhất kể từ khi thực hiện cấy vi khuẩn vào môi trường dinh dưỡng.
	\par\shortans[oly]{1005}
	\loigiai{
		Ta có đạo hàm:
		\[N'(t) = \dfrac{100\left(100+t^2\right) - 100t \cdot 2t}{\left(100+t^2\right)^2} = \dfrac{10\,000 - 100t^2}{\left(100+t^2\right)^2}.\]
		Xét $N'(t) = 0 \Leftrightarrow 10\,000 - 100t^2 = 0 \Leftrightarrow t = 10$ (vì $t \ge 0$).\\
		Ta có bảng biến thiên:
		\begin{center}
			\begin{tikzpicture}
				\tkzTabInit[nocadre=true, lgt=1.2, espcl=3.5, deltacl=0.5]
				{$t$/0.7, $N'(t)$/0.7, $N(t)$/2}
				{$0$, $10$, $+\infty$}
				
				\tkzTabLine{,+,0,-,}
				\tkzTabVar{-/$1\,000$, +/$1\,005$, -/$1\,000$}
			\end{tikzpicture}
		\end{center}
		Từ bảng biến thiên suy ra số lượng vi khuẩn lớn nhất kể từ khi thực hiện cấy vi khuẩn vào môi trường dinh dưỡng là $1\,005$ con.
	}
\end{ex}

\begin{ex}%[TT-THPT-HamRong-NH25-26 ]%[GV soạn: Nguyễn Quang Hiệp GVPB2 Đỗ Minh Vũ ]%[2D1V3-6]
	Một công ty công nghệ sản xuất linh kiện điện tử. Chi phí sản xuất $x$ linh kiện được tính theo công thức $C(x) = 0{,}5x^{2} + 200x + 5\,000$ (USD). Khách hàng ký hợp đồng bao tiêu sản phẩm với mức giá như sau:
	\begin{itemize}
		\item Với $500$ sản phẩm đầu tiên được mua với giá $1\,000$ USD/sản phẩm.
		\item Từ sản phẩm thứ $501$ trở đi, giá thu mua là $800$ USD/sản phẩm.
	\end{itemize}
	Để điều tiết sản xuất, Nhà nước áp đặt một mức thuế phụ thu là $t$ (USD) trên mỗi đơn vị sản phẩm, nhưng chỉ áp dụng cho những sản phẩm từ thứ $501$ trở đi và khi thu thuế thì mong muốn số tiền phụ thu lớn nhất. Với cách tính thuế và yêu cầu như vậy, doanh nghiệp nên sản xuất bao nhiêu linh kiện để lợi nhuận thu được lớn nhất?
	\par\shortans[oly]{550}
	\loigiai{
		Trường hợp 1: Công ty sản xuất không quá $500$ linh kiện, tức là $1 \le x \le 500$.\\
		Lợi nhuận là $L(x) = \text{Doanh thu} - \text{Chi phí sản xuất}$.\\
		Khi đó, $L(x) = 1\,000x - C(x) = 1\,000x - \left(0{,}5x^{2} + 200x + 5\,000\right) = -0{,}5x^{2} + 800x - 5\,000$.\\
		Ta có $L'(x) = -x + 800 = 0 \Leftrightarrow x = 800$ (không thỏa mãn điều kiện $1 \le x \le 500$).\\
		Mặt khác, $L(x)$ đồng biến trên khoảng $(1; 500)$ nên khi $x = 500$ thì $L(x)$ đạt giá trị lớn nhất là $L(500) = 270\,000$ (USD).\\
		Trường hợp 2: Công ty sản xuất nhiều hơn $500$ linh kiện, tức là $x > 500$.\\
		Khi đó, số sản phẩm từ thứ $501$ trở đi là $x - 500$ (sản phẩm).\\
		Doanh thu là $1\,000 \cdot 500 + 800(x - 500)$ (USD).\\
		Chi phí sản xuất là $C(x) = 0{,}5x^{2} + 200x + 5\,000$ (USD).\\
		Số tiền thuế Nhà nước phụ thu là $T(x) = t(x - 500)$ (USD).\\
		Lợi nhuận là:
		\begin{eqnarray*}
			L(x) &=& 1\,000 \cdot 500 + 800(x - 500) - \left(0{,}5x^{2} + 200x + 5\,000\right) - t(x - 500) \\
			&=& -0{,}5x^{2} + (600 - t)x + (95\,000 + 500t) \text{ (USD)}.
		\end{eqnarray*}
		Vì $L(x)$ là hàm số bậc hai có hệ số $a = -0{,}5 < 0$ nên đồ thị là parabol có bề lõm quay xuống dưới và $L(x)$ đạt giá trị lớn nhất tại $x = \dfrac{-b}{2a} = \dfrac{-(600 - t)}{2 \cdot (-0{,}5)} = 600 - t$.\\
		Như vậy, ứng với mỗi mức thuế $t$ (USD/sản phẩm), doanh nghiệp sẽ sản xuất tại điểm mà lợi nhuận thu được lớn nhất ứng với mức thuế $t$ đó.\\
		Khi đó, số linh kiện sản xuất là $x = 600 - t$ (sản phẩm).\\
		Số tiền thuế thu được là $T(t) = t(600 - t - 500) = -t^{2} + 100t$ (USD).\\
		Để số tiền thuế thu được là lớn nhất thì hàm số $T(t) = -t^{2} + 100t$ đạt giá trị lớn nhất tại $t = \dfrac{-100}{2 \cdot (-1)} = 50$ (USD).\\
		Số linh kiện doanh nghiệp sản xuất là $x = 600 - 50 = 550$ (sản phẩm).\\
		Lợi nhuận khi đó là $L = -0{,}5 \cdot 550^{2} + (600 - 50) \cdot 550 + (95\,000 + 500 \cdot 50) = 271\,250$ (USD).\\
		So sánh với trường hợp 1, ta thấy $271\,250 > 270\,000$.\\
		Vậy doanh nghiệp nên sản xuất $550$ linh kiện để lợi nhuận thu được là lớn nhất.
	}
\end{ex}

\begin{ex}%[2D1V3-6]%[KNTT - Lớp 12 - Dự án C - Đợt 2]%[Loc do]%[PB: Phan Anh]
	Một công ty công nghệ cung cấp gói lưu trữ dữ liệu doanh nghiệp với giá niêm yết $200\,000$ đồng/tháng và đang có $400$ khách hàng. Phòng kinh doanh xác định được quy luật: Cứ giảm giá $10\,000$ đồng thì sẽ có thêm $50$ khách hàng mới. Tuy nhiên, do hiện tượng quá tải băng thông, chi phí vận hành $C(n)$ (nghìn đồng) không cố định mà biến thiên theo hàm bậc hai của số lượng khách hàng $n$, được xác định bởi công thức $C(n)=28n+0,01n^2 +15\, 000$. Công ty cần chốt giá bán bao nhiêu nghìn đồng để lợi nhuận đạt mức tối đa?
	\par\shortans[0]{$160$}
	\loigiai{
		Gọi $x$ là số lần giảm giá $10\,000$ đồng $(0<x<20)$.\\
		Giá bán mới $P=200-10x$ (nghìn đồng).\\
		Số lượng khách hàng khi đó là $n=400+50x\Rightarrow x=\dfrac{n-400}{50} $.\\
		Khi đó $P=200-10x=200-10\cdot\dfrac{n-400}{50} =-\dfrac{n}{5} +280$.\\
		Ta có hàm lợi nhuận là
		\begin{align*}
			L=& n\cdot P(n)-C(n)\\
			=&\, n\left(-\dfrac{n}{5} +280\right)-(28n+0{,}01n^2 +15\, 000)\\
			=&\,-\dfrac{21}{100} n^2 +252n-15\,000.
		\end{align*}		
		Suy ra  $L'=-\dfrac{21}{50} n+252$; $L'=0\Rightarrow -\dfrac{21}{50} n+252=0\Rightarrow n=600$.\\
		Khi đó giá bán ra là $P=-\dfrac{600}{5} +280=160$ (nghìn đồng).\\
		Vậy công ty cần chốt giá bán $160$ nghìn đồng để lợi nhuận đạt mức tối đa.}
\end{ex}

\begin{ex}%[2D1V3-6] %[Dự án C đề thi thử TN Môn Toán NH25-26 - Dot 4 - Vũ Hồng Toàn]%[THPT Mai Anh Tuấn - Thanh Hóa]%[GVPB - Lê Hữu Kiệt]
	Nhà thầy Quang cách bờ biển $1$\,km. Mỗi buổi sáng thầy chạy bộ từ nhà ra bờ biển sau đó chạy dọc bờ biển $500$\,m, rồi thầy chạy qua chợ hải sản để lấy thức ăn trong ngày, cuối cùng thầy chạy về nhà. Biết chợ hải sản cách bờ biển $400$\,m và cách nhà thầy Quang $1$\,km. Quãng đường ngắn nhất mà thầy Quang đã chạy trong mỗi buổi sáng là bao nhiêu mét {\it (làm tròn đến hàng đơn vị)}?
	\begin{center}
		\begin{tikzpicture}[scale=1, font=\footnotesize, line join=round, line cap=round, >=stealth]
			\tikzset{declare function={a=4;b1=0.4*a;b2=b1+1;goc=40;}}
			\path (0,0)coordinate (H) (0,b1+b2)coordinate (A)node[shift={(90:7pt)}]{Nhà}
			(b1/2,0)coordinate (B)++(b2,0)coordinate (C)++(b1,0)coordinate (K)++(0,b2)coordinate (D)node[shift={(90:7pt)}]{Chợ};
			\fill[cyan!20] (-1,0) rectangle (6, -1);
			\draw (-1,0) -- (6,0); % Đường mặt đất
			\draw [dashed] (A)--(H) (D)--(K);
			\foreach \x/\y in {A/B,B/C, C/D,D/A}{\draw[red, thick, ->] (\x)--(\y);}
			\path (B)--(C)node[midway,shift={(-90:7pt)}]{Biển};
		\end{tikzpicture}
	\end{center}
	\par\shortans{2932}
	\loigiai{
		\begin{center}
			\begin{tikzpicture}[scale=1, font=\footnotesize, line join=round, line cap=round, >=stealth]
				\tikzset{declare function={a=4;b1=0.4*a;b2=b1+1;goc=40;}}
				\path (0,0)coordinate (H) (0,b1+b2)coordinate (A)node[shift={(90:7pt)}]{Nhà}
				(0,b2)coordinate(I)
				(b1/2,0)coordinate (B)++(b2,0)coordinate (C)++(b1,0)coordinate (K)++(0,b2)coordinate (D)node[shift={(90:7pt)}]{Chợ};
				\fill[cyan!20] (-1,0) rectangle (6, -1);
				\draw (-1,0) -- (6,0); % Đường mặt đất
				\draw [dashed] (A)--(H) (K)--(D)--(I);
				\foreach \x/\y/\z in {A/H/B,C/K/D,D/I/H} \draw pic[draw,angle radius=2mm]{right angle= \x--\y--\z};
				\foreach \x/\y in {A/B,B/C, C/D,D/A}{\draw[red, thick, ->] (\x)--(\y);}
				\path (B)--(C)node[midway,shift={(-90:7pt)}]{Biển};
				\foreach \x /\gN in {A/200,B/-90,C/-90,K/-90,H/-90,D/-20,I/180}
				\fill (\x) circle (1pt)($(\x)+(\gN:2.5mm)$) node {$\x$};
			\end{tikzpicture}
		\end{center}
		Đổi $500m = 0{,}5$\,km và $400m = 0{,}4$\,km.\\
		Đặt tên các điểm như hình vẽ.\\
		Ta có $AI = AH - DK = 1-0{,}4 = 0{,}6$.\\
		$HK = ID = \sqrt{AD^2 - AI^2} = \sqrt{1^2 -0{,}6^2} = 0{,}8$ và $BC = 0{,}5$.\\
		Đặt $HB = x$ (km) suy ra $CK = HK - HC = 0{,}3-x$ điều kiện $0 \le x \le 0{,}3$.\\
		Quãng đường mà thầy Quang chạy mỗi sáng là
		\allowdisplaybreaks
		\begin{eqnarray*}
			AB + BC + CD + DA &=& \sqrt{x^2 +1} +1{,}5 + \sqrt{(0{,}3-x)^2 +0{,}4^2} \\&=& \sqrt{x^2 +1} +1{,}5+ \sqrt{x^2-0{,}6x+0{,}25}.
		\end{eqnarray*}
		Xét hàm $f (x) = \sqrt{x^2 +1} + \sqrt{x^2 - 0{,}6x+0{,}25} +1{,}5$ với $x \in [0;0{,}3]$.\\
		Ta có $f'(x) = \dfrac{x}{\sqrt{x^2+1}} + \dfrac{x-0{,}3}{\sqrt{x^2-0{,}6x+0{,}25}}$.    \\
		Cho $f'(x)=0 \Leftrightarrow x=\dfrac{3}{14}=x_0$.\\
		Bảng biến thiên
		\begin{center}
			\begin{tikzpicture}
				\tkzTabInit[lgt=1.2,espcl=5,deltacl=0.5]
				{$x$/0.7,$f'(x)$/0.7,$f(x)$/2}
				{$-\infty$,$x_0$,$+\infty$}
				\tkzTabLine{, -,0,+,}
				\tkzTabVar{+/,-/$f(x_0)$,+/}
				\tkzTabVal[draw]{1}{2}{0.5}{$0$}{$f(0)$}
				\tkzTabVal[draw]{2}{3}{0.5}{$0{,}3$}{$f(0{,}3)$}%
				\fill[pattern = north west lines,opacity=0.3]
				(T11)rectangle (M13) (M21)rectangle (T23);
			\end{tikzpicture}
		\end{center}
		Từ bảng biến thiên ta có $\min\limits_{[0;0{,}3]} f(x) = f\left(\dfrac{3}{14}\right) \approx 2{,}932$.\\
		Vậy, quãng đường ngắn nhất mà thầy Quang chạy mỗi sáng là $2{,}932$\,km hay $2\,932$\,(mét).
	}
\end{ex}

\begin{ex}%[2D1V3-6]%[Dự án C đề thi thử TN Môn Toán NH25-26 - Dot 4 - Phạm Hoàng Đăng]%[THPT Nguyễn Trãi-Hà Nội]%[GVPB - Khắc Thiên]
	Một nông trại dâu tây sạch tại Đà Lạt (Bên A) ký hợp đồng cung cấp độc quyền cho một chuỗi cửa hàng tại TP.HCM (Bên B). Dựa trên dữ liệu thị trường, Bên B dự tính rằng nếu mỗi ngày nhập và bán hết $x$ (kg) dâu tây ($0<x< 150$), thì tổng doanh thu bán lẻ là $R(x)=300x-x^2$ (nghìn đồng). Tuy nhiên, để vận hành việc bán hàng (bảo quản, nhân sự, mặt bằng,\ldots), bên B phải chịu chi phí vận hành (không bao gồm tiền nhập hàng) là $C(x)=x^2+20x+500$ (nghìn đồng). Bên A sẽ quyết định giá bán sỉ là $p$ cho Bên B. Giả sử Bên B là đơn vị kinh doanh tối ưu, họ sẽ luôn chọn phương án nhập hàng là $x$ sao cho lợi nhuận thu được là cao nhất ứng với mức giá $p$. Hãy xác định mức giá sỉ $p$ (nghìn đồng) mà nông trại nên thiết lập để tổng doanh thu của nông trại từ việc bán hàng cho Bên B là lớn nhất (\textit{làm tròn kết quả đến hàng đơn vị}).
	\par
	\shortans[oly]{140}
	\loigiai{
		Với Bên B (Cửa hàng) ta có\\
		\begin{itemize}
			\item Doanh thu $R(x)=300x-x^2$ (nghìn đồng).
			\item Chi phí vận hành $C(x)=x^2+20x+500$ (nghìn đồng).
			\item Chi phí nhập hàng $p\cdot x$ (nghìn đồng).
			\item Lợi nhuận Bên B thu được là
			\[L_B(x)=R(x)-C(x)-px=-2x^2+(280-p)x-500.\]
			Ta có $L_B'(x)=-4x+280-p\ (0<x<150)$.\\
			Khi đó $L_B'(x)=0 \Leftrightarrow x=x_0=\dfrac{280-p}{4}$.
			\begin{center}
				\begin{tikzpicture}
					\tkzTabInit[lgt=1.5,espcl=3]
					{$x$ /0.8,$L'_B$ / 0.8,$L_B$ / 2}
					{$0$,$x_0$,$150$}
					\tkzTabLine{,+,0,-}
					\tkzTabVar{-/ $ $,+/ $ $,-/ $ $}
				\end{tikzpicture}
			\end{center}
			Do đó $L_B(x)$ đạt giá trị lớn nhất trên $(0;150)$ tại $x=\dfrac{280-p}{4} \Rightarrow p=280-4x$.
		\end{itemize}
		Với Bên A (Nông trại) ta có doanh thu của Bên A là
		\[L_A(x)=px=(280-4x) \cdot x=-4x^2+280x.\]
		Suy ra $L_A'(x)=-8x+280$.\\ 
		Khi đó $L_A'(x)=0 \Leftrightarrow x=35$.\\
		Bảng biến thiên
		\begin{center}
			\begin{tikzpicture}
				\tkzTabInit[lgt=1.5,espcl=3]
				{$x$ /0.8,$L'_A$ / 0.8,$L_A$ / 2}
				{$0$,$35$,$150$}
				\tkzTabLine{,+,0,-}
				\tkzTabVar{-/ $ $,+/ $ $,-/ $ $}
			\end{tikzpicture}
		\end{center}
		Do đó $L_A(x)$ đạt giá trị lớn nhất trên $(0; 150)$ tại $x=35$.\\
		Suy ra $p=280-4\cdot 35=140$ (nghìn đồng).
	}
\end{ex}

\begin{ex}%[2D1V4-1]%[Dự án C đề thi thử TN Môn Toán NH25-26 - Dot 4 - Phạm Hoàng Đăng]%[THPT Nguyễn Trãi-Hà Nội]%[GVPB - Khắc Thiên]
	Cho hàm số $y=g(x)=\ln\dfrac{(x-2)^2(x+5)}{(x-4)(x+1)}$. Hỏi đồ thị hàm số có bao nhiêu đường tiệm cận (chỉ xét tiệm cận đứng và tiệm cận ngang)?
	\par
	\shortans[oly]{3}
	\loigiai{
		Điều kiện xác định $\dfrac{(x-2)^2(x+5)}{(x-4)(x+1)}>0\Leftrightarrow\hoac{&-5<x<-1\\&x>4}$.\\
		Vậy tập xác định $\mathscr{D}=(-5;-1)\cup(4;+\infty)$.
		\begin{itemize}
			\item Trước hết ta tìm tiệm cận đứng.\\
			Ta biết với $a>1$ thì $\lim\limits_{x\rightarrow 0^+}\log_a x=-\infty$; $\lim\limits_{x \rightarrow+\infty} \log_a x=+\infty$.\\
			Đặt $f(x)=\dfrac{(x-2)^2(x+5)}{(x-4)(x+1)}$.\\ Ta xét các giá trị của $x$ để $f(x) \rightarrow 0^{+}$ và $f(x) \rightarrow+\infty$.
			\begin{itemize}
				\item Xét $x \rightarrow-5^{+} \Rightarrow f(x) \rightarrow 0^{+}$.\\ Do đó $\lim\limits_{x \rightarrow-5^{+}} \ln f(x)=-\infty$.\\
				Vậy $x=-5$ là một đường tiệm cận đứng của đồ thị $y=g(x)=\ln \dfrac{(x-2)^2(x+5)}{(x-4)(x+1)}$.
				\item Xét $x \rightarrow-1^{-} \Rightarrow f(x) \rightarrow+\infty$.\\ Do đó $\lim\limits_{x \rightarrow-1^{-}} \ln f(x)=+\infty$.\\
				Vậy $x=-1$ là một đường tiệm cận đứng của đồ thị $y=g(x)=\ln \dfrac{(x-2)^2(x+5)}{(x-4)(x+1)}$.
				\item Xét $x \rightarrow 4^{+} \Rightarrow f(x) \rightarrow+\infty$.\\ Do đó $\lim\limits_{x \rightarrow 4^{+}} \ln f(x)=+\infty$.\\
				Vậy $x=4$ là một đường tiệm cận đứng của đồ thị $y=g(x)=\ln \dfrac{(x-2)^2(x+5)}{(x-4)(x+1)}$.
			\end{itemize}
			Vậy đồ thị đã cho có $3$ đường tiệm cận đứng.
			\item Ta có $\lim\limits_{x \rightarrow+\infty }f(x)=+\infty$ nên $\lim\limits_{x \rightarrow+\infty} \ln f(x)=+\infty$.\\
			Vậy đồ thị hàm số đã cho không có tiệm cận ngang.
		\end{itemize}
		Vậy đồ thị hàm số có $3$ đường tiệm cận.
	}
\end{ex}

\begin{ex}%[2D1V5-8]%[TT Cụm 5 - Ninh Bình-NH25-26]%[GVSB:Dương Công Tạo]%[GVPB1: Nguyễn Tấn Tài]
	\immini[thm]{Trong hệ trục tọa độ $Oxy$, đơn vị mỗi trục là mét, một đường trượt mới sẽ được xây dựng theo bản thiết kế đã trình bày như hình vẽ. Thanh trượt bắt đầu từ $A$ và kết thúc tại $C$, đường cong của thanh trượt là một phần của đồ thị hàm số $f(x) = \dfrac{ax^2 + bx + c}{x+d}$, biết đồ thị hàm số $f(x)$ tiếp xúc với trục $Ox$ tại điểm $B$. Bạn Nam bắt đầu trượt từ điểm $A$, hỏi khi Nam cách vị trí ban đầu theo phương ngang một khoảng $5$ mét thì Nam cách mặt đất bao nhiêu mét, biết trục $Ox$ nằm trên mặt đất (\textit{kết quả làm tròn đến hàng phần trăm}).}
	{\begin{tikzpicture}[scale=.5,>=stealth, line join=round, line cap=round, x=0.7cm, y=0.7cm, every node/.style={font=\footnotesize},
			declare function={
				hA = 10;
				xB = 10;
				dxC = 5;
				xC = xB + dxC;
				yC = 2.5; 
				a = hA/(xB*xB);
				f(\x) = a*(\x-xB)*(\x-xB);
				wWall = 2.5;
				hTop = 1.2;
				xSlantTop = -wWall - 0.8;
				xSlantBot = -wWall - 1.5;
				hBase = 1.5;
				xBaseLeft = xSlantBot - 0.2;
			}]
			
			\fill[black!85] (xBaseLeft, -hBase) rectangle (0, 0);
			
			\fill[pattern=north west lines, pattern color=gray!70] (xSlantBot, 0) -- (xSlantTop, hA) -- (-wWall, hA) -- (-wWall, 0) -- cycle;
			\draw (xSlantBot, 0) -- (xSlantTop, hA) -- (-wWall, hA) -- (-wWall, 0) -- cycle;
			
			\fill[gray!30] (-wWall, 0) rectangle (0, hA - hTop);
			
			\fill[pattern=north west lines, pattern color=gray!70] (-wWall, hA - hTop) rectangle (0, hA);
			
			\draw (-wWall, 0) rectangle (0, hA);
			\draw (-wWall, hA - hTop) -- (0, hA - hTop);
			
			\draw[->] (0, 0) -- (xC + 2, 0) node[below] {$x$};
			\draw[->] (0, 0) -- (0, hA + 2) node[left] {$y$};
			
			\draw[smooth, samples=200, domain=0:xC] plot (\x, {f(\x)});
			
			\fill (0, 0) circle (1pt) node[shift={(45:10pt)}] {$O$};
			\fill (0, hA) circle (1pt) node[right=4pt] {$A$};
			\fill (xB, 0) circle (1pt) node[above=6pt] {$B$};
			\fill (xC, yC) circle (1pt) node[above=4pt] {$C$};
			
			\draw[<->] (-wWall/2, 0) -- (-wWall/2, hA) node[midway, fill=gray!30, inner sep=2pt] {$10$ m};
			
			\draw (xB, 0) -- (xB, -1.3);
			\draw (xC, 0) -- (xC, -1.3);
			
			\draw[<->] (0, -1) -- (xB, -1) node[midway, fill=white, inner sep=2pt] {$10$ m};
			\draw[<->] (xB, -1) -- (xC, -1) node[midway, fill=white, inner sep=2pt] {$5$ m};
			
			\draw[<->] (xC, 0) -- (xC, yC) node[midway, right] {$1$ m};
			
	\end{tikzpicture}}
	\par  \shortans{1,67}
	\loigiai{
		Ta có đồ thị hàm số $f(x)$ tiếp xúc với trục $Ox$ tại điểm $B(10;0)$ nên $f(x) = \dfrac{a(x-10)^2}{x+d}$. \\
		Đồ thị đi qua $A(0;10) \Rightarrow 10 = \dfrac{a \cdot 100}{d} \Rightarrow d = 10a$. \\
		Đồ thị đi qua $C(15;1) \Rightarrow 1 = \dfrac{a \cdot 25}{15+d} \Rightarrow 15+d = 25a \Rightarrow 15+10a = 25a \Rightarrow a = 1$; $d =10$. \\
		Do đó $f(x) = \dfrac{(x-10)^2}{x+10}$. \\
		Có $x = 5 \Rightarrow y = \dfrac{25}{15} \approx 1{,}67$. \\
		Vậy Nam cách vị trí ban đầu theo phương ngang một khoảng $5$ mét thì Nam cách mặt đất $1{,}67$ (m).
	}
\end{ex}

\begin{ex}%[2D1V5-8]%[Dự án C - Đợt 3 - Sở Giáo Dục Sơn La]%[Lê Hùng Thắng - Nguyễn Tiến]%%[Thy Nguyen Vo Diem - Nguyễn Tiến]
	Một nhà địa chất học đang ở địa điểm $A$ trên sa mạc. Anh ta muốn đến điểm $B$ (cũng ở trên sa mạc) và cách $A$ một khoảng bằng $70$ km. Trong sa mạc, xe anh ta chỉ có thể di chuyển với vận tốc $30$ km/h. Nhà địa chất phải đến được điểm $B$ sau $2$ giờ, vì vậy nếu anh ta đi thẳng từ $A$ đến $B$ sẽ không thể đến đúng giờ được. Rất may, có một con đường nhựa song song với đường nối $A$ và $B$ và cách $AB$ một đoạn $10$ km (tham khảo hình vẽ). 
	\begin{center}
		\begin{tikzpicture}[scale=1, font=\footnotesize, line join=round, line cap=round, >=stealth]
			\coordinate (A1) at (0,0);
			\coordinate (A) at (0,2);
			\coordinate (B) at (7,2);
			\coordinate (B1) at (7,0);
			\draw[thick] (-1,0) -- (9,0) node[below] {Đường nhựa};
			\draw[dashed] (A)--(A1) node[midway, left] {$10$ km};
			\draw[dashed] (B)--(B1) node[midway, right] {$10$ km};
			\draw[<->] (0,2.5) -- (7,2.5) node[midway, above] {$70$ km};
			\draw[fill=black] (A) circle (1pt) node[above] {$A$};
			\draw[fill=black] (B) circle (1pt) node[above] {$B$};
			\draw[blue, thick] (A) -- (1,0) -- (6,0) -- (B);
		\end{tikzpicture}
	\end{center}
	Trên đường nhựa đó, xe nhà địa chất có thể di chuyển với vận tốc $50$ km/h. Thời gian ngắn nhất để nhà địa chất di chuyển từ $A$ đến $B$ là bao nhiêu phút? 
	\shortans{116}
	\loigiai{
		Giả sử lộ trình của nhà địa chất là từ $A$ đi đến một điểm $M$ trên đường nhựa, sau đó di chuyển trên đường nhựa đến điểm $N$, rồi từ $N$ đi đến $B$.\\
		Gọi $A'$ và $B'$ lần lượt là hình chiếu của $A$ và $B$ trên đường nhựa. 
		\begin{center}
			\begin{tikzpicture}[scale=1, font=\footnotesize, line join=round, line cap=round, >=stealth]
				\coordinate (A1) at (0,0);
				\coordinate (A) at (0,2);
				\coordinate (B) at (7,2);
				\coordinate (B1) at (7,0);
				\draw[thick] (-1,0) -- (9,0) node[below] {Đường nhựa};
				\draw[dashed] (A)--(A1) node[midway, left] {$10$ km};
				\draw[dashed] (B)--(B1) node[midway, right] {$10$ km};
				\draw[<->] (0,2.5) -- (7,2.5) node[midway, above] {$70$ km};
				\draw[fill=black] (A) circle (1pt) node[above] {$A$};
				\draw[fill=black] (B) circle (1pt) node[above] {$B$};
				\draw[blue, thick] (A) -- (1,0) -- (6,0) -- (B);
				\node[below] at (1,0) {$M$};
				\node[below] at (6,0) {$N$};
				\node[below] at (0,0) {$A'$};
				\node[below] at (7,0) {$B'$};
			\end{tikzpicture}
		\end{center}
		Khi đó $A'B' = AB = 70$ km và $AA' = BB' = 10$ km. \\
		Đặt $A'M = x$ (km) với $0 \le x \le 35$.\\
		Do tính đối xứng của bài toán, để thời gian là ngắn nhất thì $B'N = x$.
		\\Khi đó, quãng đường di chuyển trong sa mạc là 
		$$AM = NB = \sqrt{x^2 + 10^2} = \sqrt{x^2 + 100} \text{(km).}$$
		Quãng đường di chuyển trên đường nhựa là $MN = 70 - 2x$ (km).	\\
		Tổng thời gian di chuyển là:
		$$T(x) = \dfrac{2\sqrt{x^2 + 100}}{30} + \dfrac{70 - 2x}{50} = \dfrac{\sqrt{x^2 + 100}}{15} + \dfrac{35 - x}{25} \text{(giờ)}.$$ 
		Xét hàm số $T(x)$ liên tục trên $[0; 35]$.\\
		$T'(x) = \dfrac{x}{15\sqrt{x^2 + 100}} - \dfrac{1}{25}$. \\
		\begin{eqnarray*}
			& & T'(x) = 0 \Leftrightarrow 25x = 15\sqrt{x^2 + 100} \\
			& \Leftrightarrow  & 5x = 3\sqrt{x^2 + 100} \\
			& \Leftrightarrow  & x^2 = \dfrac{225}{4} \Leftrightarrow x = 7{,}5.
		\end{eqnarray*}
		Ta có $T(7{,}5) = \dfrac{\sqrt{7{,}5^2 + 100}}{15} + \dfrac{35 - 7{,}5}{25} = \dfrac{29}{15}$ (giờ).\\
		Vậy thời gian ngắn nhất để nhà địa chất di chuyển từ $A$ đến $B$ là $\dfrac{29}{15}$ giờ = $116$ phút.
		\textbf{Cách khác:}\\
		Gọi $H$, $K$, $C$, $D$ là các điểm như hình vẽ.
		\begin{center}
			\begin{tikzpicture}[scale=1, font=\footnotesize, line join=round, line cap=round, >=stealth]
				\path (-1,0) coordinate (C) 
				(3,0) coordinate (D)
				(-3,2) coordinate (A) 
				(4.5,2) coordinate (B)
				(-3,0) coordinate (H)
				(4.5,0) coordinate (K)
				;
				\draw (-5,0)--(5,0) (-5,2)--(5,2) (A)--(C) (B)--(D) (A)--(H) node[below]{$H$} (B)--(K) node[below]{$K$};
				\fill (C) circle (1pt) (D) circle(1pt) (A) circle (1pt) (B) circle (1pt) (H) circle (1pt) (K) circle (1pt);
				\draw (C) node[below]{$C$} --(D)node[below]{$D$} node[midway,below]{Đường nhựa};
				\draw (A) node[above]{$A$}--(B)node[above]{$B$} node[midway,above]{Sa mạc} node[midway,below]{$70$ km}
				;
				\draw[<->] (-4.5,0)--(-4.5,2) node[midway, right]{$10$ km};
				\draw pic[draw,angle radius=2mm]{right angle=B--K--D};
				\draw pic[draw,angle radius=2mm]{right angle=C--H--A};
			\end{tikzpicture}
		\end{center}
		Ta chia quãng đường đi thành ba giai đoạn 
		\begin{itemize}
			\item\textbf{ Giai đoạn 1:} đi từ $A$ đến $C$ (từ sa mạc đến đường nhựa song song).
			\item \textbf{Giai đoạn 2:} đi từ $C$ đến $D$ (một quãng đường nào đó trên đường nhựa).
			\item \textbf{Giai đoạn 3:} đi từ $D$ đến $B$ (từ điểm kết thúc $D$ trên đường nhựa đi tiếp đến $B$ băng qua sa mạc).
		\end{itemize}
		Đặt $HC=x$ ($0<x<70$) và $DK=y$ ($0<y<70$).\\
		Ta có
		\begin{itemize}
			\item $AC=\sqrt{10^2+x^2} \Rightarrow t_{AC}=\dfrac{AC}{v_{AC}}=\dfrac{\sqrt{100+x^2}}{30}$;
			\item $DB=\sqrt{10^2+y^2} \Rightarrow t_{DB}=\dfrac{DB}{v_{DB}}=\dfrac{\sqrt{100+y^2}}{30}$;
			\item $CD=70-(x+y) \Rightarrow t_{CD}=\dfrac{CD}{v_{CD}}=\dfrac{70-x-y}{50}$.
		\end{itemize}
		Tổng thời gian nhà địa chất học đi từ $A$ đến $B$ là \[T(x;y)=\dfrac{\sqrt{100+x^2}}{30}+\dfrac{\sqrt{100+y^2}}{30}+\dfrac{70-x-y}{50}.\]
		Đây là một biểu thức có dạng đối xứng $2$ biến $x$, $y$ và ta cần tìm $\min T(x;y)$.\\
		Ta có $$T(x;y)=\dfrac{100+x^2}{30}+\dfrac{35-x}{50}+\dfrac{100+y^2}{30}+\dfrac{35-y}{50}=f(x)+f(y).$$
		Xét $f(u)=\dfrac{\sqrt{100+u^2}}{30}+\dfrac{35-u}{50}$, $0<u<70$.\\
		Suy ra $f'(u)=\dfrac{u}{30\sqrt{100+u^2}}-\dfrac{1}{50}$.\\
		Cho $f'(u)=0 \Leftrightarrow \sqrt{100+u^2}=\dfrac{5u}{3}>0 \Leftrightarrow u=\dfrac{15}{2}$.\\
		Bảng biến thiên
		\begin{center}
			\begin{tikzpicture}
				\tkzTabInit[nocadre=true, lgt=1.2, espcl=4, deltacl=0.5]
				{$u$/0.7,$f’(u)$/0.7,$f(u)$/2}
				{$0$,$7{,}5$,$70$}
				\tkzTabLine{ ,-,0,+, }
				\tkzTabVar{+/$ $,-/$ $,+/$ $}
			\end{tikzpicture}
		\end{center}
		Dựa vào bảng biến thiên, ta có $\min\limits_{u \in (0;70)} f(u)=f\left(\dfrac{15}{2}\right)=\dfrac{29}{30}$.\\
		Do đó ta có $T(x;y)=f(x)+f(y) \ge \dfrac{29}{30}+\dfrac{29}{30}=\dfrac{29}{15}$ giờ, tức là $116$ phút.
	}
\end{ex}

\begin{ex}%[Dự án đề TT Toán NH25-26-Dot 4- Hải Phụng]%[THCS-THPT LÊ THÁNH TÔNG - TP.HCM]%[GVPB - Viet Hoang Pham]%[2D4C3-2]
	\immini[thm]{
		Anh Nghĩa có một mảnh đất dạng hình thang cong $OABC$  ($B$ là điểm cực đại của đồ thị hàm số $y=f(x)$ được mô hình hóa trong mặt phẳng $Oxy$  (đơn vị mỗi trục là $10$ m). Anh Nghĩa chia mảnh đất hình thang cong $OABC$  thành hai phần để làm hồ bơi và làm vườn trồng cỏ, có đường ngăn cách bởi một phần của đồ thị hàm bậc ba  $y=f(x)$ như hình vẽ. Biết đơn giá làm hồ bơi là $400\,000$ đồng/m$^2$, giá trồng cỏ là $200\,000$ đồng/m$^2$. Tổng chi phí anh Nghĩa phải trả là $295$ triệu đồng. Bên cạnh đó có một con đường nhựa được mô hình hóa bởi hàm số $g(x) = \dfrac{x+1}{x-2}$ $(x>2)$. Anh Nghĩa muốn làm một đoạn đường $MN$ từ vườn anh đến con đường nhựa đó. Hãy tính độ dài ngắn nhất của đoạn đường $MN$. (Đơn vị: mét; làm tròn đến hàng phần trăm)
	}{
		\begin{tikzpicture}[>=stealth,scale=1, line join = round, line cap = round]
			\tikzset{label style/.style={font=\footnotesize}}
			\draw[->] (-1,0)--(0,0) node[below left]{$O$} -- (5,0) node[above]{$x$};
			\draw[->] (0.,-1) -- (0.,5) node[left]{$y$};
			\foreach \x in {3}
			\draw[shift={(\x,0)}] (0pt,2pt) -- (0pt,-2pt) node[below] {\footnotesize $\x$};
			\draw [smooth, samples=100, domain=-0.5:3] plot (\x,{-(\x)^3+3*(\x)^2}) node[above right]{$f(x)$};
			
			\draw[fill=blue, opacity=0.1] (0,0) plot[domain=0:2] (\x,{-(\x)^3+3*(\x)^2})--(0,4)--(0,0)--cycle;
			\draw[fill=green!80!black, opacity=0.2] (0,0) plot[domain=0:3] (\x,{-(\x)^3+3*(\x)^2})--(3,0)--(0,0)--cycle;
			\draw[fill=black] (0,4) circle(1.3pt) node[left]{$A$} -- (2,4) circle(1.3pt) node[above]{$B$};
			\draw [smooth, samples=100, domain=2.8:4.8] plot (\x, {(\x+1)/(\x-2)}) node[above right]{$g(x)$};
			
			\draw[fill=black] (3.3,3.3) circle(1.3pt) node[above right]{$N$} -- (2.6,2.704) circle(1.3pt) node[below left]{$M$};
			\draw (0.6,3) node{\footnotesize Hồ bơi} (2,1) node{\footnotesize Vườn anh Nghĩa};
		\end{tikzpicture}
	}
	\shortans{$7{,}49$}
	\loigiai{
		Dựa vào đồ thị ta thấy $f(x) = ax^2(x-3) = ax^3-3ax^2$.\\
		Suy ra $f'(x) = 3ax^2-6ax = 3ax(x-2)$ có nghiệm  $x=0$ hoặc $x =2$.\\
		Ta có $f(2) = -4a$. \\
		Diện tích hồ bơi là $S_1 = 100 \displaystyle\int\limits_{0}^{2} [-4a-f(x)] \mathrm{\,d}x = 100  \displaystyle\int\limits_{0}^{2} [-4a-a(x^3-3x^2)] \mathrm{\,d}x = -400a$. \\
		Diện tích vườn là $S_2 = 100  \displaystyle\int\limits_{0}^{3} f(x) \mathrm{\,d}x = 100a  \displaystyle\int\limits_{0}^{3} (x^3-3x^2) \mathrm{\,d}x = -675a$. \\
		Theo đề bài ta có 
		\begin{eqnarray*}
			400\,000S_1 + 200\,000S_2 = 295\,000\,000 &\Leftrightarrow& 400\,000(-400a) + 200\,000(-675a) = 295\,000\,000 \\
			&\Leftrightarrow& a=-1
		\end{eqnarray*}
		Vậy $f(x) = -x^3 +3x^2$.\\
		Gọi $M(m;-m^3+3m^2)$, $N\left(n;1+\dfrac{3}{n-2}\right)$, với $n>2$. \\
		Để $MN$ nhỏ nhất thì tiếp tuyến của hai đồ thị tại $M,N$ phải song song với nhau nên 
		$$f'(m)=g'(n) \Leftrightarrow -3m^2+6m = \dfrac{-3}{(n-2)^2} \Leftrightarrow n = 2 + \dfrac{1}{\sqrt{m^2-2m}} \text{ (vì } n>2).$$\\
		Khi đó $N(2 + \dfrac{1}{\sqrt{m^2-2m}};1+3\sqrt{m^2-2m})$. \\
		Từ đó $MN = \sqrt{\left(2+\dfrac{1}{\sqrt{m^2-2m}}-m\right)^2 + \left(1+3\sqrt{m^2-2m} + m^3 -3m^2\right)^2}$. \\
		Sử dụng MTCT tìm được $MN_{\min} \approx 0{,}7485$ khi $ m\approx 2{,}37$. \\
		Vậy độ dài ngắn nhất của đoạn $MN$ là $d_{\min} = 0{,}7485 \cdot 10 \approx 7{,}49$.
	}
\end{ex}

\begin{ex}%[2D4C3-2]%[Dự án C đợt 3 NH25-26- Phạm Đăng Đăng]%[TT-THPT-TranNhanTong-HaNoi-N25-26]%[GVPB - Phan Quốc Trí]
	Hoạ sĩ thiết kế logo hình con cá cho doanh nghiệp kinh doanh hải sản. Logo là hình phẳng giới hạn bởi hai parabol với kích thước được cho trong hình sau (đơn vị trên mỗi trục toạ độ là cm).
	\begin{center}
		\begin{tikzpicture}[scale=0.5, line cap=round, line join=round, font=\small,>=stealth]
			\fill[gray,smooth]
			plot[domain=-7:6] (\x,{3 - \x*\x/12})
			-- plot[domain=6:-7] (\x,{-5 + 5*\x*\x/36})
			-- cycle;
			% Khung chữ nhật
			\draw (-7,-5) rectangle (7,3);
			
			% Trục tọa độ
			\draw[->] (-7,0) -- (9,0) node[right] {$x$};
			\draw[->] (0,-5) -- (0,3.5) node[above] {$y$};
			\node at (0,0) [below left] {$O$};
			
			% Parabol trên (kéo dài tới biên trái)
			\draw[domain=-7:6,smooth,variable=\x]
			plot (\x,{3 - \x*\x/12});
			
			% Parabol dưới (kéo dài tới biên trái)
			\draw[domain=-7:6,smooth,variable=\x]
			plot (\x,{-5 + 5*\x*\x/36});
			
			% Hai điểm
			\coordinate (A) at (3,2.25);
			\coordinate (B) at (3,-3.75);
			
			% Đường cong tròn hơn
			\draw
			(A)
			.. controls (1.8,1) and (1.8,-2.5) ..
			(B);
			
			% Đánh dấu điểm
			\fill (A) circle (1.5pt);
			\fill (B) circle (1.5pt);
			
			
			
			% Điểm trong hình
			\fill (4,0.8) circle (0.15);
			
			% Kích thước ngang
			\draw[<->] (-7,-5.8) -- (-6,-5.8);
			\node at (-6.5,-6.3) {$1$ cm};
			
			\draw[<->] (-6,-5.8) -- (0,-5.8);
			\node at (-3,-6.3) {$6$ cm};
			
			\draw[<->] (0,-5.8) -- (6,-5.8);
			\node at (3,-6.3) {$6$ cm};
			
			% Kích thước dọc
			\draw[<->] (7.5,0) -- (7.5,3);
			\node[right] at (7.5,1.5) {$3$ cm};
			
			\draw[<->] (7.5,0) -- (7.5,-5);
			\node[right] at (7.5,-2.5) {$5$ cm};
			
		\end{tikzpicture}
	\end{center}
	Diện tích logo bằng bao nhiêu cm$^2$ (không làm tròn kết quả các phép tính trung gian, chỉ làm tròn kết quả cuối cùng)?
	\shortans{$65{,}4$}
	\loigiai{Do hình con cá giới hạn bởi hai parabol nên ta gọi $(P_1)$ là parabol có hệ số $a_1>0$ có công thức dạng 
		$
		y_1=a_1x^2+b_1x+c_1$.\\
		Khi đó $(P_1)$ đi qua các điểm có tọa độ là $(6;0)$, $(0;-5)$, $(-6;0)$.\\
		Từ đó ta có hệ phương trình
		$\heva{
			&36a_1+6b_1+c_1=0\\
			&36a_1-6b_1+c_1=0\\
			&c_1=-5
		}
		\Leftrightarrow
		\heva{
			&a_1=\dfrac{5}{36}\\
			&b_1=0\\
			&c_1=-5.
		}
		$\\
		Vậy $(P_1)$ có công thức là
		$
		y_1=\dfrac{5}{36}x^2-5$.\\
		Gọi $(P_2)$ là parabol có hệ số $a_2<0$ có công thức dạng $y_2=a_2x^2+b_2x+c_2$.\\
		Khi đó $(P_2)$ đi qua các điểm có tọa độ là $(6;0)$, $(0;3)$, $(-6;0)$.\\
		Từ đó ta có hệ phương trình
		$\heva{
			&36a_2+6b_2+c_2=0\\
			&36a_2-6b_2+c_2=0\\
			&c_2=3}
		\Leftrightarrow
		\heva{
			&a_2=-\dfrac{1}{12}\\
			&b_2=0\\
			&c_2=3.}$\\
		Vậy $(P_2)$ có công thức là $
		y_2=-\dfrac{1}{12}x^2+3.
		$\\
		Khi đó diện tích con cá được tính bởi công thức
		$$
		S=\displaystyle\int \limits_{-7}^{6}\left|y_1-y_2\right|\mathrm{\,d}x
		=\displaystyle\int\limits_{-7}^{6}\left|\dfrac{2}{9}x^2-8\right|\mathrm{\,d}x
		\approx 65{,}4.
		$$
		Vậy diện tích logo là $65{,}4$ cm$^2$.}
\end{ex}

\begin{ex}%[2D4C3-4]%[Dự án C đề thi thử TN Môn Toán NH25-26 - Dot 4 - Phạm Đăng Đăng]%[THPT Lê Quý Đôn-Đống Đa-Hà Nội]%[GVPB - Phan Quốc Trí]
	Sân vận động Sport Hub (Singapore) là sân có mái vòm kỳ vĩ nhất thế giới. Đây là nơi diễn ra lễ khai mạc Đại hội thể thao Đông Nam Á năm $2015$. Nền sân phẳng là một elip $(E)$ có trục lớn dài $150$ m, trục bé dài $90$ m. Một mặt phẳng vuông góc với trục lớn của $(E)$ và cắt elip ở $A$, $B$ và cắt phần dưới mái vòm theo một cung tròn của đường tròn tâm $I$ thoả mãn góc $\widehat{AIB}=90^\circ$ (tham khảo hình bên dưới). Các kĩ sư cần tính toán để lắp máy điều hòa cho sân vận động, mỗi máy điều hoà có công suất $100\,000$  BTU đáp ứng được không gian $500$ m$^3$. Hỏi các kĩ sư cần dùng tối thiểu bao nhiêu máy để đáp ứng đủ theo yêu cầu cho sân vận động?
	\begin{center}
		\begin{tikzpicture}[scale=1,font=\footnotesize,line join=round
			,line cap=round,>=stealth]
			\draw[thick] (0,0) ellipse (3cm and 1.5cm);
			\coordinate (A) at ({3*cos(50)}, {1.5*sin(50)});
			\coordinate (B) at ({3*cos(-50)}, {1.5*sin(-50)});
			\draw[thick] (A) -- (B);
			\node[above] at (A) {$A$};
			\node[below] at (B) {$B$};
			\fill[black] (A) circle (1.5pt);
			\fill[black] (B) circle (1.5pt);
			
			\begin{scope}[xshift=8.5cm]
		\coordinate (I) at (0,0);
				\draw[thick] (I) circle (2cm);
				\coordinate (A) at ({2*cos(45)}, {2*sin(45)});
				\coordinate (B) at ({2*cos(-45)}, {2*sin(-45)});
				\draw[thick] (I) -- (A);
				\draw[thick] (I) -- (B);
				\node[left] at (I) {$I$};
				% Label A nằm phía trên bên phải điểm
				\node[above right] at (A) {$A$};
				% Label B nằm phía dưới bên phải điểm
				\node[below right] at (B) {$B$};
				
				% Vẽ các dấu chấm tròn tại vị trí các điểm để đánh dấu rõ hơn
				\fill[black] (I) circle (1.5pt);
				\fill[black] (A) circle (1.5pt);
				\fill[black] (B) circle (1.5pt);
			\end{scope}
		\end{tikzpicture}
	\end{center}
	\shortans{$232$}
	\loigiai{\begin{center}
			\begin{tikzpicture}[scale=1,font=\footnotesize,line join=round
				,line cap=round,>=stealth]
				% 1. Định nghĩa tâm O
				\coordinate (O) at (0,0);
				\draw[thick] (O) ellipse (3cm and 1.2cm);
				\draw[->, thick] (O) -- (0,3.5) node[left] {$z$};
				\draw[->, thick] (O) -- (3.8, 1) node[above] {$x$};
				\draw[->, thick] (O) -- (-2.8, 2) node[below] {$y$};
				\fill (O) circle (1.5pt);
				\begin{scope}[xshift=8cm]
					\coordinate (O) at (0,0);
					\draw[->, thick] (-3.5, 0) -- (3.5, 0) node[above right] {$x$};
					\draw[thick] (O) ellipse (3cm and 1.8cm);
					\node[above] at (O) {$O$};
					\fill[black] (O) circle (1.5pt);
					\coordinate (A) at ({3*cos(50)}, {1.8*sin(50)});
					\coordinate (B) at ({3*cos(-50)}, {1.8*sin(-50)});
					\draw[thick] (A) -- (B);
					\node[above right] at (A) {$A$};
					\node[below right] at (B) {$B$};
					\fill[black] (A) circle (1.5pt);
					\fill[black] (B) circle (1.5pt);
				\end{scope}
			\end{tikzpicture}
		\end{center}
		Chọn hệ trục như hình vẽ.\\ Ta cần tìm diện tích của $S(x)$ thiết diện.
		Gọi $\mathrm{d}(O, AB) = x$.\\
		$(E): \dfrac{x^2}{75^2} + \dfrac{y^2}{45^2} = 1$.
		Khi đó $AB = 2y = 2\sqrt{45^2 \left( 1 - \dfrac{x^2}{75^2} \right)} = 90\sqrt{1 - \dfrac{x^2}{75^2}}$.\\
		Suy ra $R = \dfrac{AB}{\sqrt{2}} = \dfrac{90}{\sqrt{2}}\sqrt{1 - \dfrac{x^2}{75^2}} \Rightarrow R^2 = \dfrac{90^2}{2} \left( 1 - \dfrac{x^2}{75^2} \right)$.\\
		$S(x) = \dfrac{1}{4}\pi R^2 - \dfrac{1}{2}R^2 = \left( \dfrac{1}{4}\pi - \dfrac{1}{2} \right)R^2 = (\pi - 2)\dfrac{2025}{2}\left( 1 - \dfrac{x^2}{75^2} \right)$.\\
		Vậy thể tích khoảng không bên dưới mái che và bên trên mặt sân là
		$$V = \displaystyle\int\limits_{-75}^{75} (\pi - 2)\dfrac{2025}{2}\left( 1 - \dfrac{x^2}{75^2} \right)\, \mathrm{\,d}x \approx 115\,586 \, \text{m}^3.$$
		Số chiếc điều hòa cần lắp là
		$115586 : 500 = 231{,}172$.\\
		Vậy cần $232$ chiếc điều hòa.}
\end{ex}

\begin{ex}%[2D4V1-6]%[Dự án đề thi thử đợt 3-Nguyễn Đức Lợi]%[SGD tỉnh Lạng Sơn]%[GVPB - ThanhTa]
	Ở giai đoạn thải trừ, giai đoạn cuối sau khi một người uống một liều thuốc, nồng độ thuốc trong máu, ký hiệu là $C(t)$ (đơn vị: mg/l), giảm dần sau $t$ giờ kể từ khi giai đoạn này bắt đầu. Khi đó, tốc độ giảm nồng độ $C'(t)$ tỉ lệ với chính nồng độ hiện có, tức là $\dfrac{C'(t)}{C(t)}=-k$ ($k$ là một hằng số dương). Biết rằng khi bắt đầu giai đoạn thải trừ, nồng độ thuốc còn lại là $12$ mg/l và sau $6$ giờ kể từ lúc bắt đầu thải trừ, nồng độ đo được là $3$ mg/l. Sau khoảng bao nhiêu giờ thì nồng độ còn lại bằng $2$ mg/l ({\it kết quả làm tròn đến hàng phần mười})?
	\par
	\shortans{7,8}
	\loigiai{
		Ta có $\dfrac{C'(t)}{C(t)}=-k \Rightarrow \displaystyle \int\dfrac{C'(t)}{C(t)}\mathrm{\,d}t=\displaystyle \int-k\mathrm{\,d}t \Rightarrow \ln \big|C(t)\big|=-kt+m$ ($m$ là hằng số).\\
		Vì $C(t)>0$ nên ta có $C(t)=\mathrm{e}^{-kt+m}=\mathrm{e}^m \cdot \mathrm{e}^{-kt}=A\mathrm{e}^{-kt}$ (đặt hằng số $\mathrm{e}^m=A$).\\
		$C(0)=12 \Leftrightarrow A\mathrm{e}^{-k.0}=12 \Leftrightarrow A=12$.\\
		$C(6)=3 \Leftrightarrow 12\mathrm{e}^{-k\cdot6}=3 \Leftrightarrow \mathrm{e}^{-6k}=\dfrac{1}{4} \Leftrightarrow k=-\dfrac{1}{6} \ln \dfrac{1}{4}$.\\
		Suy ra $C(t)=12\mathrm{e}^{\tfrac{1}{6}\left(\ln \tfrac{1}{4}\right)t}$.\\
		Mà ta lại có \[C(t)=2\Leftrightarrow 12\mathrm{e}^{\tfrac{1}{6}\left(\ln \tfrac{1}{4}\right)t}=2\Leftrightarrow \mathrm{e}^{\tfrac{1}{6}\left(\ln \tfrac{1}{4}\right)t}=\dfrac{1}{6}\Leftrightarrow \dfrac{1}{6}\left(\ln \dfrac{1}{4}\right)t=\ln \dfrac{1}{6}\Leftrightarrow t=\dfrac{6 \ln \dfrac{1}{6}}{\ln \dfrac{1}{4}} \approx 7{,}8.\]
		Vậy sau khoảng $7{,}8$ giờ thì nồng độ còn lại bằng $2$ mg/l.
	}
\end{ex}

\begin{ex}%[2D4V2-6]%[Dự án C đề thi thử TN Môn Toán NH25-26 - Dot 4 - Vũ Hồng Toàn]%[THPT Mai Anh Tuấn - Thanh Hóa]%[GVPB - Lê Hữu Kiệt]
	Một vật chuyển động trong $4$ giờ với vận tốc $v$\,(km/h) phụ thuộc thời gian $t$\,(h) có đồ thị của vận tốc như hình bên. Trong khoảng thời gian $2$ giờ kể từ khi bắt đầu chuyển động, đồ thị đó là một phần của đường parabol có đỉnh $I(2;7)$ và trục đối xứng của parabol song song với trục tung, khoảng thời gian còn lại đồ thị là đoạn thẳng $IA$. Quãng đường $s$ mà vật di chuyển được trong $4$ giờ đó là bao nhiêu km {\it (kết quả được làm tròn đến hàng phần chục)}?
	\begin{center}
		\begin{tikzpicture}[xscale=1,yscale=0.6, font=\footnotesize, line join=round, line cap=round,>=stealth]
			\tikzset{declare function={f(\x)=-(\x)^2+4*(\x)+3;
			xmin=0;xmax=4;ymin=0;ymax=7;},
			smooth,samples=100
		}
		\path (4,3)coordinate (A) (2,7)coordinate (I);
		\draw[->] (xmin-0.5,0)--(0,0) node[below left]{$O$} --(xmax+1,0) node[shift={(-100:7pt)}] {$x$};
		\draw[->] (0,ymin-1)--(0,ymax+1.5) node[shift={(200:7pt)}] {$y$};
		\fill (0,0) circle (1.2pt);
		\draw[red,domain=xmin:2] plot (\x,{f(\x)});
		\draw[dashed,red,domain=2:xmax] plot (\x,{f(\x)});
		\draw (I)--(A);
		\foreach \y/\goc in {3/180}{
			\fill (0,\y) circle (1pt) node[shift={(\goc:7pt)}]{$\y$};
		}
		\foreach \x/\y/\px/\py in {2/7/-90/180,4/3/-90/180}{
			\fill (\x,0) circle (1.2pt) node[shift={(\px:7pt)}]{$\x$};
			\fill (0,\y) circle (1.2pt) node[shift={(\py:7pt)}]{$\y$};
			\fill (\x,\y) circle (1.2pt);
			\draw[dashed, thin] (\x,0) |-(0,\y);
		}
		\foreach \x /\gN in {A/20,I/90}{
			\fill (\x) circle (1pt)($(\x)+(\gN:9pt)$) node {$\x$};
		}
	\end{tikzpicture}
\end{center}
\par\shortans{21,3}
\loigiai{
	Phương trình parabol là $v = -t^2 + 4t + 3$\,(km/h).\\
	Quãng đường vật chuyển động được trong khoảng thời gian $2$ giờ đầu là
	$$\displaystyle\int\limits_0^2 v(t)\mathrm{\,d}t = \displaystyle\int\limits_0^2 \left(-t^2 + 4t + 3\right)\mathrm{d}t = \dfrac{34}{3}\,\text{(km)}.$$
	Quãng đường vật đi được trong $2$ giờ sau là $\dfrac{(7+3)\cdot 2}{2} = 10$\,km.\\
	Vậy tổng quãng đường vật đi được trong $4$ giờ là $\dfrac{34}{3}+10\approx 21{,}3$\,(km).
}
\end{ex}

\begin{ex}%[2D4V3-2]%[TT-SGD-TuyenQuang-L1-NH25-26]%[Võ Thị Thùy Trang]%[GVPB: Thanh Phát]
	\immini[thm]{
		Người ta dự định trồng hoa để trang trí vào phần tô đậm trong hình vẽ bên. Biết rằng phần tô đậm là diện tích hình phẳng giới hạn bởi đồ thị hai hàm số $y=f(x)=ax^3+bx^2+cx+6$ và $y=g(x)=-bx^2+mx+n$ trong đó $a$, $b$, $c$, $m$, $n \in \mathbb{R}$. Biết đồ thị hai hàm số đã cho cắt nhau tại các điểm có hoành độ lần lượt bằng $-2$; $1$; $3$. Chi phí trồng hoa là $150\,000$ đồng/$1$ m$^2$ và đơn vị trên mỗi trục tọa độ là mét. Tổng chi phí để trồng hoa theo dự định là bao nhiêu nghìn đồng (\textit{làm tròn kết quả đến hàng đơn vị})?
		}
	{\begin{tikzpicture}[>=stealth,line join=round,line cap=round,font=\footnotesize,scale=.7]
			\def\aa{1} \def\bb{0} \def\cc{0}
			\draw[->] (-5,0)--(5,0)node[below right]{$x$};
			\draw[->] (0,-1)--(0,10)node[left]{$y$};
			\fill (0,0)node[below left]{$O$};
			\draw[black,samples=150,smooth,domain=-3.2:3.2] plot(\x,{\aa*(\x)^2+\bb*\x+\cc});
			\draw[black,samples=150,smooth,domain=-2.37:3.1] plot(\x,{1*(\x)^3-1*(\x)^2-5*\x+6});
			\fill (3,9)circle(1pt)
			(1,1)circle(1pt)
			(-2,4)  circle(1pt)
			;
			\draw[dashed](3,0)--(3,9) (-2,0)--(-2,4) (1,0)--(1,1);
			\node at (-2,0)[below]{$-2$};
			\node at (1,0)[below]{$1$};
			\node at (3,0)[below]{$3$};
			\path[pattern = north east lines, pattern color = black][domain=-2.0:1.0] plot(\x,{\aa*(\x)^2+\bb*\x+\cc})-- (1.,1.) {[smooth,samples=50,domain=1.0:-2.0] -- plot(\x,{(\x)^(3)-(\x)^(2)-5*(\x)+6})} -- (-2.,4.) -- cycle;
			\path[pattern = north east lines, pattern color = black]{[smooth,samples=50,domain=1.0:3.0]plot(\x,{\aa*(\x)^2+\bb*\x+\cc})} -- (3.,9.) {[smooth,samples=50,domain=3.0:1.0] -- plot(\x,{(\x)^(3)-(\x)^(2)-5*(\x)+6})} -- (1.,1.) -- cycle;
	\end{tikzpicture}}
	\par\shortans{3163}
	\loigiai{Phương trình hoành độ giao điểm của hai đồ thị hàm số $y=f(x)$ và $y=g(x)$ là 
		\[f(x) - g(x) = 0 \Leftrightarrow ax^3 + 2bx^2 + (c-m)x + 6 - n = 0 \quad (*)\]
		Theo giả thiết, đồ thị hai hàm số cắt nhau tại các điểm có hoành độ $x = -2$, $x = 1$, $x = 3$ nên phương trình $(*)$ có $3$ nghiệm phân biệt là $-2$; $1$; $3$. Do đó ta có 
		\begin{align*}
			f(x) - g(x) &= a(x+2)(x-1)(x-3) \\
			&= a\left( x^2 + x - 2\right) (x - 3) \\
			&= a\left(x^3 - 2x^2 - 5x + 6\right) \\
			&= ax^3 - 2ax^2 - 5ax + 6a.
		\end{align*}
		Đồng nhất hệ số với biểu thức ban đầu của $f(x) - g(x)$, ta được $\heva{
			&2b = -2a \\
			&c - m = -5a \\
			&6 - n = 6a.}$\\
		Quan sát đồ thị, ta thấy đường Parabol $y = g(x)$ đi qua gốc tọa độ $O(0;0)$.\\
		Suy ra $g(0) = 0 \Rightarrow n = 0$.\\
		Thay $n = 0$ vào phương trình $6 - n = 6a$, ta được $6 = 6a \Leftrightarrow a = 1$.\\
		Khi đó $f(x) - g(x) = x^3 - 2x^2 - 5x + 6$.\\
		Diện tích phần tô đậm là
		\[S =\displaystyle\int\limits_{-2}^{3} \left| f(x) - g(x)\right|  \mathrm{\,d}x =\displaystyle\int\limits_{-2}^{1} \left| x^3 - 2x^2 - 5x + 6\right|  \mathrm{\,d}x + \int_{1}^{3} \left| x^3 - 2x^2 - 5x + 6\right|  \mathrm{\,d}x.\]
		Dựa vào đồ thị, trên khoảng $(-2; 1)$ thì $f(x)$ nằm trên $g(x)$ và trên khoảng $(1; 3)$ thì $f(x)$ nằm dưới $g(x)$. Ta có
		\begin{align*}
			S &=\displaystyle\int\limits_{-2}^{1} \left( x^3 - 2x^2 - 5x + 6\right) \mathrm{\,d}x - \displaystyle\int\limits_{1}^{3} \left( x^3 - 2x^2 - 5x + 6\right)  \mathrm{\,d}x \\
			&= \left( \dfrac{x^4}{4} - \dfrac{2x^3}{3} - \dfrac{5x^2}{2} + 6x \right) \Bigg|_{-2}^{1} - \left( \dfrac{x^4}{4} - \dfrac{2x^3}{3} - \dfrac{5x^2}{2} + 6x \right) \Bigg|_{1}^{3} \\
			&= \left( \dfrac{37}{12} - \dfrac{-152}{12} \right) - \left( \dfrac{-27}{12} - \dfrac{37}{12} \right) \\
			&= \dfrac{189}{12} - \left( -\dfrac{64}{12} \right) = \dfrac{63}{4} + \dfrac{16}{3} = \dfrac{253}{12}\,\text{m}^2.
		\end{align*}
		Tổng chi phí dự định để trồng hoa là
		\[T = \dfrac{253}{12} \cdot 150\,000 = 3\,162\,500\, \text{(đồng)}\approx 3\,163\, \text{(nghìn đồng).}\]
	}
\end{ex}

\begin{ex}%[2D4V3-2]%[KNTT - Lớp 12 - Dự án C - Đợt 2]%[Loc do]%[PB: Phan Anh]
	\immini[thm]{Một bể chứa nước có diện tích mặt cắt ngang được tô màu như hình bên, ở đó đơn vị trên các trục đo bằng mét. Trên mặt cắt ngang, phần đáy của bể chứa có phương trình $y=k(x-8)^2$. Tính diện tích của mặt cắt ngang theo mét vuông.}{\begin{tikzpicture}[x=0.5cm,y=0.5cm,scale=0.8,>=stealth]
			\def\aa{0.094} \def\bb{-1.5} \def\cc{6}
			\draw[->] (-1.5,0)--(13,0)node[below right]{$x$};
			\draw[->] (0,-1.25)--(0,7)node[left]{$y$};
			\fill (0,0)node[below left]{$O$};
			\begin{scope}
				\clip (-1,-1) rectangle (14,8);
				\fill[pattern=crosshatch,pattern color=green]
				plot[samples=200,domain=0:12,variable=\x] (\x,{6})-- plot[samples=200,domain=12:0,variable=\x] (\x,{0.094*(\x)^2-1.5*(\x)+6})-- cycle;
				\draw plot[samples=200,domain=0:12,variable=\x] (\x,{0*(\x)+6});
				\draw plot[samples=200,domain=12:0,variable=\x] (\x,{0.094*(\x)^2-1.5*(\x)+6});
			\end{scope}
			\draw (12,0)--(12,6) node[above] {$N(12,6)$};
			\fill (0,6) circle (1pt) node[left] {$M$};
			\fill (8,0) circle (1pt) node[below] {$(8,0)$};
			\fill (12,0) circle (1pt) node[below] {$P$};
	\end{tikzpicture}}
	\par\shortans[0]{54}
	\loigiai{
		Đáy của bể nước là một phần của parabol có phương trình $y=k(x-8)^2$.\\
		Parabol đi qua điểm $M(0;\, 6)\Rightarrow k\cdot(-8)^2 =6\Rightarrow k=\dfrac{3}{32}$.\\ 
		Khi đó ta có Parabol $y=\dfrac{3}{32} (x-8)^2$.\\
		Vậy diện tích của mặt cắt ngang là
		$S=\displaystyle\int\limits _0^{12}\left[6-\left(\dfrac{3}{32} \left(x-8\right)^2 \right)\right] \mathrm{\,d}x=54$ m$^2$.
	}
\end{ex}

\begin{ex}%[2D4V3-4]%[Dự án đề thi thử NH25-26]%[THPT Chuyên Phan Bội Châu - Nghệ An]%[Sỹ Trường]
	\immini[thm]
	{
		Khi thi công tuyến cao tốc Nghi Sơn-Bãi Vọt, đơn vị thi công cần san một ngọn núi nhỏ. Biết đường viền chân núi là một elip có trục lớn $AB=400$ (m) và trục bé $CD=200$ (m); mặt cắt bởi mặt phẳng chứa $AB$ và vuông góc với mặt đất là nửa hình tròn đường kính $AB$; mặt cắt bởi các mặt phẳng vuông góc với $AB$ có dạng parabol với đỉnh thuộc nửa đường tròn đường kính $AB$ (tham khảo hình vẽ). Tính thể tích ngọn núi đó theo đơn vị triệu $\text{m}^3$ (làm tròn đến hàng phần trăm).
	}
	{
		\begin{tikzpicture}[line join=round, line cap=round, >=stealth, font=\footnotesize, scale=1]
			\begin{scope}[rotate=15]
				\coordinate (O) at (0,0);
				\coordinate (A) at (-3,0);
				\coordinate (B) at (3,0);
				\coordinate (C) at (0,-1);
				\coordinate (D) at (0,1);
				\def\xh{-2.2}
				\def\yh{0.68} 
				\coordinate (H) at (\xh, 0);
				\coordinate (E) at (\xh, -\yh);
				\coordinate (F) at (\xh, \yh);
				\draw (B) arc (0:-180:3 and 1);
				\draw[dashed] (B) arc (0:180:3 and 1);
				\draw[dashed] (A) -- (B);
				\draw[dashed] (C) -- (D);
				\draw[dashed] (E) -- (F);
				\draw ($(H)+(-0.15,0)$) -- ($(H)+(-0.15,-0.15)$) -- ($(H)+(0,-0.15)$);
			\end{scope}
			\draw (B) [x={(B)}, y={(0,3)}] arc (0:180:1 and 1);
			\def\zh{2.04}
			\coordinate (G) at ($(H) + (0, \zh)$);
			\begin{scope}[shift={(H)}]
				\coordinate (vecX) at ($(F)-(H)$);
				\coordinate (vecY) at (0, \zh);
				\draw [x={(vecX)}, y={(vecY)}] plot[domain=-1:0, samples=50, variable=\t] (\t, {1-\t*\t});
				\draw[dashed][x={(vecX)}, y={(vecY)}] plot[domain=0.45:1, samples=50, variable=\t] (\t, {1-\t*\t});
			\end{scope}
			\draw[dashed] (H) -- (G);
			\draw[dashed] ($(H)!0.25!(E)$) -- ++(0, 0.2) -- ($(H)+(0,0.2)$);
			\foreach \p/\ang in {A/150, B/30, C/-45, D/135} {
				\fill (\p) circle (1pt);
				\node at ($(\p)+(\ang:3.5mm)$) {$\p$};
			}
		\end{tikzpicture}	
	}
	\shortans{$7{,}11$}
	\loigiai{
		Chọn hệ trục tọa độ $Oxyz$ sao cho gốc $O$ trùng với tâm của đường elip chân núi, trục $Ox$ chứa trục lớn $AB$, trục $Oy$ chứa trục bé $CD$ và trục $Oz$ hướng thẳng đứng lên trên.
		\begin{itemize}
			\item Vì trục lớn $AB = 400$ nên bán trục lớn $a = 200$.
			\item Vì trục bé $CD = 200$ nên bán trục bé $b = 100$.
		\end{itemize}
		Phương trình đường elip viền chân núi trong mặt phẳng $Oxy$ là
		\[\frac{x^2}{200^2} + \frac{y^2}{100^2} = 1 \Rightarrow y = \pm 100\sqrt{1 - \frac{x^2}{200^2}} = \pm \frac{1}{2}\sqrt{200^2 - x^2}.\]
		Xét mặt phẳng vuông góc với trục $Ox$ tại điểm có hoành độ $x$ ($x \in [-200; 200]$), mặt phẳng này cắt ngọn núi theo một thiết diện là hình parabol.
		\begin{itemize}
			\item Đáy của parabol này là đoạn thẳng nối hai điểm trên elip, có độ dài là
			\[d = 2y = 2 \cdot \frac{1}{2}\sqrt{200^2 - x^2} = \sqrt{200^2 - x^2}.\]
			\item Theo giả thiết, mặt cắt chứa $AB$ và vuông góc với mặt đất là nửa đường tròn đường kính $AB$. Do đó, phương trình của nửa đường tròn này trong mặt phẳng $Oxz$ là $x^2 + z^2 = 200^2$ (với $z \ge 0$).
			\item Chiều cao của thiết diện parabol (từ đáy lên đỉnh thuộc nửa đường tròn) chính là cao độ $z$
			\[h = \sqrt{200^2 - x^2}.\]
		\end{itemize}
		Diện tích của thiết diện parabol tại hoành độ $x$ được tính bằng $S = \dfrac{2}{3} \cdot \text{đáy} \cdot \text{chiều cao}$.\\ Suy ra 
		\[S(x) = \dfrac{2}{3} \cdot d \cdot h = \dfrac{2}{3} \cdot \sqrt{200^2 - x^2} \cdot \sqrt{200^2 - x^2} = \dfrac{2}{3}\left(200^2 - x^2\right).\]
		Thể tích ngọn núi cần san lấp là
		\allowdisplaybreaks
		\[\begin{aligned}[t]
			V &=\displaystyle\int\limits_{-200}^{200} S(x) \mathrm{\,d}x = \displaystyle\int\limits_{-200}^{200} \dfrac{2}{3}(200^2 - x^2) \mathrm{\,d}x \\
			&= \dfrac{4}{3}\displaystyle\int\limits_{0}^{200} (200^2 - x^2) \mathrm{\,d}x = \dfrac{4}{3} \left( 200^2x - \dfrac{x^3}{3} \right) \Bigg|_{0}^{200} \\
			&= \dfrac{4}{3} \left( 200^3 - \dfrac{200^3}{3} \right) = \frac{4}{3} \cdot \dfrac{2}{3} \cdot 200^3 = \dfrac{8}{9} \cdot 8\,000\,000 = \dfrac{64\,000\,000}{9} \text{ (m}^3).
		\end{aligned}\]
		Tính xấp xỉ ta được $V \approx 7\,111\,111{,}11 \text{ m}^3 \approx 7{,}11$ triệu $\text{m}^3$.\\
		Vậy thể tích ngọn núi cần tính xấp xỉ là $7{,}11$ triệu $\text{m}^3$.
	}
\end{ex}

\begin{ex}%[2D4V3-5]%[Dự án C - Đợt 3 - Sở Giáo Dục Sơn La]%[Nguyễn Chiến - Nguyễn Tiến]
	Từ một quả cầu bằng đá trắng bán kính bằng $1$ dm người ta khoan rút lõi ngay \lq\lq chính giữa\rq\rq\ quả cầu (trục đối xứng của lõi và quả cầu trùng nhau) như hình minh họa, đường kính lõi là $1$ dm. Thể tích còn lại của quả cầu bằng bao nhiêu dm$^3$? (không làm tròn kết quả ở các bước trung gian, làm tròn kết quả ở bước cuối cùng đến hàng phần trăm).
	\begin{center}
		\begin{tikzpicture}[scale=1.5]
			\def\rx{.5}\def\ry{.1}\def\a{433/500}\def\b{.5}\def\c{.15}
			\draw (0,0) circle(1)
			(-.5,-\a)--(-.5,\a) (.5,-\a)--(.5,\a);
			\foreach \x/\s in {.5/, .5/dashed}
			\foreach \y in {\a,-\a}
			\draw[\s] (\x,\y) arc(0:\ifx\s\empty-180\else180\fi:{\rx} and {\ry});
			\node at(0,-1.3){Quả cầu đá};
			\begin{scope}[shift={(2,-1)}]
				\tikzset{declare function={d = 3; r = 0.5;}}
				\draw[] (0,0)--++(d,0) to [out=60,in=-60] ++(0,2*r)
				--++(-d,0) to [out=-120,in=120] cycle;
				\draw (0,0) arc(270:90:{0.1} and {r})
				(d,0) arc(270:90:{0.1} and {r});
				\draw[dashed] (0,0) arc(-90:90:{0.1} and {r})
				(d,0) arc(-90:90:{0.1} and {r});
				
				\node at(1.5,-0.3){Lõi};
			\end{scope}
		\end{tikzpicture}
	\end{center}
	\shortans[oly]{$2{,}72$}
	\loigiai{
		\immini[thm]{
			Bán kính quả cầu là $R = 1$ dm.\\
			Lõi bị khoan rút có phần hình trụ với bán kính đáy $r = \dfrac{1}{2} = 0{,}5$ dm.\\
			Chiều cao của phần khối trụ là $$h = 2\sqrt{R^2 - r^2} = 2\sqrt{1^2 - 0{,}5^2} = \sqrt{3} \, \text{dm.}$$
			Chọn hệ trục tọa độ $Oxy$ với gốc $O$ là tâm quả cầu, trục $Ox$ dọc theo trục của khối trụ.\\
			Khi đó, phần thể tích còn lại của quả cầu bằng thể tích khối tròn xoay tạo thành khi cho hình phẳng giới hạn bởi đường cong $y = \sqrt{1 - x^2}$, đường thẳng $y = 0{,}5$ và hai đường thẳng $x = -\dfrac{\sqrt{3}}{2}$, $x = \dfrac{\sqrt{3}}{2}$ quay quanh trục $Ox$.
		}{
			\begin{tikzpicture}[scale=1.2, font=\footnotesize, line join=round, line cap=round, >=stealth]
				\def\R{2}
				\def\r{1}
				\def\h{1.73205}
				\draw[->] (-2.3,0) -- (2.3,0) node[below] {$x$};
				\draw[->] (0,-2.3) -- (0,2.3) node[left] {$y$};
				\draw[thick] (0,0) circle (\R);
				\draw[dashed] (-2.05,\r) -- (2.05,\r);
				\draw[dashed] (-2.05,-\r) -- (2.05,-\r);
				\coordinate [label=below left:$O$] (O) at (0,0);
				\coordinate [label=above right:$A$] (A) at (\h,\r);
				\coordinate [label=above left:$H$] (H) at (\h,0);
				\draw[thick] (O) -- (A) node[midway, above left] {$R$};
				\draw[dashed] (-\h,0) -- (-\h,\r);
				\draw[dashed] (H) -- (A) node[midway, right] {$r$};
				\foreach \i in {A,H,O} \draw[fill=black] (\i) circle(1pt);
				\draw pic[draw,angle radius=5]{right angle=A--H--O};
				\node[below] at (\h,0) {$\dfrac{\sqrt{3}}{2}$};
				\node[below] at (-\h,0) {$-\dfrac{\sqrt{3}}{2}$};
				\node[above left] at (0,\r) {$0{,}5$};
				\foreach \p in {(\h,0), (-\h,0), (0,\r)} \fill \p circle (1pt);
			\end{tikzpicture}
		}
		\noindent Thể tích cần tìm là
		{\allowdisplaybreaks
			\begin{eqnarray*}
				V &=& \pi \int\limits_{-\frac{\sqrt{3}}{2}}^{\frac{\sqrt{3}}{2}} \left[ \left(\sqrt{1 - x^2}\right)^2 - (0{,}5)^2 \right] \mathrm{\,d}x \\
				&=& \pi \int\limits_{-\frac{\sqrt{3}}{2}}^{\frac{\sqrt{3}}{2}} \left( 0{,}75 - x^2 \right) \mathrm{\,d}x \\
				&=& \pi \left( 0{,}75x - \dfrac{x^3}{3} \right) \Bigg|_{-\frac{\sqrt{3}}{2}}^{\frac{\sqrt{3}}{2}} = \dfrac{\pi\sqrt{3}}{2} \approx 2{,}72.
		\end{eqnarray*}}
	}
\end{ex}

\begin{ex}%[2H2H2-3]%[TT Cụm 5 - Ninh Bình-NH25-26]%[GVSB:Dương Công Tạo]%[GVPB1: Nguyễn Tấn Tài]
	Trong không gian với hệ trục toạ độ $Oxyz$, cho hình hộp $ABCD.A'B'C'D'$ có $A(-3;2;1)$, $C(4;2;0)$, $B'(-2;1;1)$, $D'(3;5;4)$. Gọi toạ độ điểm $A'(a;b;c)$. Tính giá trị của biểu thức $\dfrac{a-b-c}{2}$.
	\par  \shortans{-4,5}
	\loigiai{
		\immini[thm]{
			Gọi $I$, $I'$ lần lượt là trung điểm của $AC$ và $B'D'$ thì $I\left(\dfrac{1}{2};2;\dfrac{1}{2}\right)$. $I'\left(\dfrac{1}{2};3;\dfrac{5}{2}\right)$. \\
			Mặt khác $AA'I'I$ là hình bình hành nên \[\overrightarrow{AA'}=\overrightarrow{II'}\Leftrightarrow\heva{&a+3=\dfrac{1}{2}-\dfrac{1}{2}\\ &b-2=3-2\\ &c-1=\dfrac{5}{2}-\dfrac{1}{2}}\Leftrightarrow\heva{&a=-3\\ &b=3\\ &c=3}\Rightarrow A'(-3;3;3).\] 
		}
		{\begin{tikzpicture}[scale=.65,line join=round, line cap=round,font=\footnotesize]
				\def\a{3.5}
				\def\h{4}
				\path 	(0:0) coordinate (A)
				++(0:\a) coordinate (D)
				++(-130:\a/2) coordinate (C)
				($(A)+(C)-(D)$) coordinate (B)	
				($(A)!0.25!(C)$) coordinate (H)
				($(H)+(90:\h)$) coordinate (A')
				($(A')+(C)-(A)$) coordinate (C')
				($(C')+(B)-(C)$) coordinate (B')
				($(A')+(D)-(A)$) coordinate (D')
				($(C)!.5!(A)$) coordinate (I)
				($(C')!.5!(A')$) coordinate (I');
				\draw[dashed] 	(B)--(A)--(D)	(A)--(A') (A)--(C)	(I)--(I');
				\draw	(C)--(C') 	(D)--(D') 	(B)--(B')  (B)--(C)--(D) (A')--(B')--(C')--(D')--cycle	(A')--(C');
				\foreach \x/\g in {A/180,B/180,C/0,D/0,A'/180,B'/180,C'/0,D'/0,I/-135,I'/45}
				\fill[black] 	(\x) circle (1pt)
				($(\g:4mm)+(\x)$) node {$\x$};
			\end{tikzpicture}
		}
		\noindent Vậy $\dfrac{a-b-c}{2} = -4{,}5$.
	}
\end{ex}

\begin{ex}%[2H2K1-5]%[Dự án C-TT THPT NH25-26-Đợt 4- ĐỖ CHÍ TÂM]%[THPT Phan Bội Châu-Khánh Hòa]%[GVPB: Phạm Văn Long]
	Một nhóm học sinh thiết kế một chiếc lều dã ngoại có dạng hình nón theo cách sau: Từ một tấm bạt hình tròn có diện tích $4{,}84\pi$ m$^2$, nhóm cắt lấy một hình quạt có góc ở tâm là $\alpha \left(0{<}\alpha {<}2\pi \right)$ (như hình 1) để làm thành mặt xung quanh của lều (như hình 2). Sau đó, nhóm dùng một tấm bạt chống thấm cắt thành hình tròn để trải kín đáy lều. Biết chi phí loại bạt làm mặt xung quanh lều là $100\,000$ đồng/m$^2$ và chi phí loại bạt làm đáy lều là $150\,000$ đồng/m$^2$. Giả sử phần bạt thừa sau khi cắt không tính vào chi phí và kích thước các mép nối không đáng kể. Để thiết kế được chiếc lều có diện tích đáy lớn nhất với tổng chi phí không vượt quá $2\,300\,000$ đồng, chiều cao $h$ của lều bằng bao nhiêu mét (kết quả làm tròn đến hàng phần trăm).
	\begin{center}
		\begin{tikzpicture}
			% --- Hình 1: Hình quạt tròn ---
			\begin{scope}[shift={(0,0)}]
				\draw (0,0) coordinate (O) circle (1.5cm);
				\draw (0,0) -- (30:1.5cm) arc (30:330:1.5cm) -- (0,0);
				\filldraw[fill=green!20, draw=green!50!black] (0,0) -- (30:0.5cm) arc (30:330:0.5cm) -- cycle;
				\node at (-0.3, 0) {$\alpha$};
				\node at (0, -2) {\textbf{Hình 1}};
			\end{scope}	
		\end{tikzpicture} \hspace*{1cm}
		\begin{tikzpicture}[scale=1.2]
			% --- Hình 2: Hình nón ---
			\begin{scope}[shift={(5,0)}]
				% Đáy hình nón (elip)
				\draw[dashed] (1,0) arc (0:180:1cm and 0.3cm);
				\draw (1,0) arc (0:-180:1cm and 0.3cm);
				\draw (-1,0) -- (0,2.5) -- (1,0);
				
				% Các đường nét bên trong
				\draw[dashed] (0,0) -- (0,2.5) node[midway, left] {$h$};
				\draw[dashed] (0,0) -- (1,0) node[midway, below] {$r$};
				\draw (0,2.5) -- (1,0) node[midway, right] {$\ell$};
				
				\node at (0, -0.7) {\textbf{Hình 2}};
			\end{scope}	
		\end{tikzpicture}
	\end{center}
	\shortans{$1{,}52$}
	\loigiai{
		Tấm bạt ban đầu có diện tích $S=\pi {R^2}\Leftrightarrow {R^2}=\dfrac{S}{\pi}=4,84\Rightarrow R=2{,}2$ (m).\\
		Vì mặt xung quanh của lều được cắt từ tấm bạt hình tròn này, nên bán kính của tấm bạt là độ dài đường sinh của hình nón. Do đó $\ell=2,2$ (m).\\
		Diện tích xung quanh của hình nón là: ${S_{xq}}=\pi r\ell=2{,}2\pi r$.\\
		Diện tích đáy của hình nón là: ${S_d}=\pi{r^2}$.\\
		Tổng chi phí làm lều là $100\,000\cdot {S_{xq}}+150\,000\cdot {S_d}=100\,000\cdot 2{,}2\pi r+150\,000\pi r^2$.\\
		Theo đề bài $100\,000\cdot 2{,}2\pi r+150\,000\cdot \pi {r^2}\le 2\,300\,000\Leftrightarrow 22\pi r+15\pi r^2\le 230$\\
		\hspace*{9.8cm}$\Leftrightarrow-3{,}061\le r\le 1{,}594 $.\\
		Do đó, diện tích đáy lớn nhất thỏa mãn $22\pi r+15\pi {r^2}\le 230$ khi $r\approx 1{,}594$.\\
		Lại có $h=\sqrt{{\ell^2}-{r^2}}\approx \sqrt{{{\left(2{,}2\right)}^2}-{{\left(1{,}594\right)}^2}}\approx 1,52$.\\
		Vậy chiều cao của lều là $h\approx 1{,}52$ ($m$).
	}
\end{ex}

\begin{ex}%[TT-THPT-HamRong-NH25-26 ]%[GV soạn: Nguyễn Quang Hiệp GVPB2 Đỗ Minh Vũ ]%[2H2V1-4]
	Một giá treo vật được lắp đặt vào góc tường tại ba điểm $C$, $O$ và $B$ bằng các bản lề cố định. Cấu tạo giá treo gồm ba thanh $AB$, $AC$ và $AO$ hội tụ tại điểm $A$ (khối lượng các thanh không đáng kể). Biết các thanh $AB$, $AC$ có chiều dài bằng nhau và vuông góc nhau tại $A$. Thanh $OA$ đóng vai trò thanh chống, chịu được lực nén tối đa là $150$ N. Khi treo vật có trọng lượng $P$ vào điểm $A$, lực căng trên hai thanh $AC$ và $AB$ có độ lớn luôn bằng một phần ba lực nén trên thanh $OA$. Trọng lượng lớn nhất có thể treo vào đầu $A$ của giá treo bằng bao nhiêu để hệ thống vẫn an toàn? (Kết quả làm tròn đến hàng đơn vị).
	\begin{center}
		\begin{tikzpicture}[>=stealth, line join=round, line cap=round, font=\footnotesize, scale=1.2]
			
			% Định nghĩa màu sắc cho thanh cứng
			\definecolor{rodcolor}{RGB}{205, 145, 55}
			
			% Định nghĩa các tọa độ chính
			\coordinate (O) at (0,0);
			\coordinate (A) at (2.8, 2.0);
			\coordinate (C) at (-2, 2.0);
			\coordinate (B) at (4.2, 3.5);
			\coordinate (D) at (0, 3.5); % Điểm giao của các nét đứt trên trục Z
			
			% 1. Vẽ các trục tọa độ
			\draw[thick,->] (O) -- (5,0) node[above left] {\textbf{y}}; % Trục Y
			\draw[thick,->] (O) -- (0,4.5)node[above left] {\textbf{z}}; % Trục Z (thẳng đứng).
			\draw[thick,->] (O) -- (-2.5,-1.5)node[above left] {\textbf{x}}; % Trục X (xiên)
			
			% 2. Vẽ các đường nét đứt (mặt phẳng chiếu)
			\draw[dashed, thick] (C) -- (D) -- (B);
			\draw[dashed, thick] (A) -- (D);
			
			% 3. Vẽ các thanh cứng (dùng nét đôi để tạo cảm giác hình khối 3D)
			\tikzstyle{rod}=[line width=4pt, rodcolor]
			\tikzstyle{rodborder}=[line width=5pt, black]
			
			% Thanh OA
			\draw[rodborder] (O) -- (A);
			\draw[rod] (O) -- (A);
			% Thanh CA
			\draw[rodborder] (C) -- (A);
			\draw[rod] (C) -- (A);
			% Thanh AB
			\draw[rodborder] (A) -- (B);
			\draw[rod] (A) -- (B);
			
			% 4. Vẽ dây treo và vật nặng (màu xanh lá)
			\draw[thick] (A) -- ([yshift=-1.2cm]A);
			\filldraw[fill=green!40!black, draw=black, thick] 
			([xshift=-0.4cm, yshift=-1.2cm]A) rectangle ([xshift=0.4cm, yshift=-2.5cm]A);
			
			\coordinate (W1) at ([xshift=-3pt, yshift=0.7cm]C);
			\coordinate (W2) at ([xshift=-3pt, yshift=-0.7cm]C);
			\draw[thick] (W1) -- (W2);
			\foreach \i in {0, 0.2, 0.4, 0.6, 0.8, 1.0, 1.2, 1.4} {
				\draw[thick] ([yshift=-\i cm]W1) -- ++(-0.25,-0.2);
			}
			
			% Ngàm tại B (Trần ngang)
			\coordinate (C1) at ([yshift=3pt, xshift=-0.7cm]B);
			\coordinate (C2) at ([yshift=3pt, xshift=0.7cm]B);
			\draw[thick] (C1) -- (C2);
			\foreach \i in {0, 0.2, 0.4, 0.6, 0.8, 1.0, 1.2, 1.4} {
				\draw[thick] ([xshift=\i cm]C1) -- ++(0.2, 0.25);
			}
			\foreach \point/\goc in {A/-60,C/75,B/-60,O/130}{
				\draw[fill=white](\point)circle(2.8pt)+(\goc:2mm)node[scale=1]{$\point$};
			}
			
		\end{tikzpicture}
	\end{center}
	\par\shortans[oly]{132}
	\loigiai{
		\begin{center}
			\begin{tikzpicture}[>=stealth, line join=round, line cap=round, font=\footnotesize, scale=1]
				
				% =========================================================
				% 1. ĐỊNH NGHĨA PIC KHỐI HỘP ĐỘNG
				% =========================================================
				\tikzset{
					pics/hinhHopLuc/.style n args={8}{
						code={
							% Thiết lập độ dài và góc cho 3 chiều không gian (3 véc-tơ)
							\tikzset{
								declare function={
									lenFtwo=3.5; gocFtwo=180;   % Trục véc-tơ F2
									lenFone=2.12; gocFone=55;   % Trục véc-tơ F1
									lenP=3.5; gocP=-90;         % Trục véc-tơ P
								}
							}
							% Tính toán 4 đỉnh của mặt trên (gốc A)
							\path
							(0,0) coordinate (#1) % Đỉnh gốc A
							+(gocFtwo:lenFtwo) coordinate (#2) % C'
							+(gocFone:lenFone) coordinate (#4) % B'
							+(gocP:lenP) coordinate (#5)       % P
							($(#2)+(#4)-(#1)$) coordinate (#3) % TopBackLeft
							;
							% Tịnh tiến mặt trên xuống mặt đáy theo véc-tơ P
							\foreach \pone/\pname in {#2/#6, #3/#7, #4/#8}{
								\path ($(\pone)+(#5)-(#1)$) coordinate (\pname);
							}
							
							% Vẽ các cạnh khuất (hội tụ tại đỉnh đối diện A là O')
							\foreach \pointo/\pointt in {#7/#3, #7/#6, #7/#8}{
								\draw[dashed] (\pointo)--(\pointt);
							}
							
							% Vẽ các cạnh thấy (không vẽ các cạnh từ A để vẽ véc-tơ đè lên sau)
							\foreach \pointo/\pointt in {#2/#3, #3/#4, #6/#5, #5/#8, #2/#6, #4/#8}{
								\draw (\pointo)--(\pointt);
							}
						}
					}
				}
				
				% =========================================================
				% 2. GỌI PIC VÀ TRUYỀN TÊN CÁC ĐIỂM
				% =========================================================
				% Truyền vào 8 node tương ứng: A, C', TopBackLeft, B', P, BottomFrontLeft, O', BottomBackRight
				\path (2,1) pic{hinhHopLuc={A}{C'}{TopBackLeft}{B'}{P_vec}{BottomFrontLeft}{O'}{BottomBackRight}};
				
				% =========================================================
				% 3. XÁC ĐỊNH CÁC TIA KÉO DÀI
				% =========================================================
				\path 
				($(A)!1.3!(C')$) coordinate (C)
				($(A)!1.3!(B')$) coordinate (B)
				($(A)!1.6!(O')$) coordinate (O);
				
				\draw (C') -- (C);
				\draw (B') -- (B);
				\draw[dashed] (O') -- (O);
				
				% =========================================================
				% 4. VẼ CÁC VÉC-TƠ LỰC
				% =========================================================
				\draw[->, thick] (A) -- (C') node[midway, above] {$\overrightarrow{F}_2$};
				\draw[->, thick] (A) -- (B') node[midway, right] {$\overrightarrow{F}_1$};
				\draw[->, thick] (A) -- (P_vec) node[midway, above right] {$\overrightarrow{P}$};
				\draw[->, thick, dashed] (A) -- (O') node[pos=0.6, below right] {$\overrightarrow{F}_3$};
				
				% =========================================================
				% 5. ĐÁNH DẤU ĐIỂM VÀ TÊN DÙNG VÒNG LẶP (\foreach)
				% =========================================================
				% Cấu trúc: Tọa_độ / Tên_hiển_thị / Góc_đặt_nhãn
				\foreach \point/\name/\goc in {
					A/A/135, B'/B'/0, C'/C'/90, 
					P_vec/P/0, O'/O'/135, 
					B/B/90, C/C/90, O/O/135
				}{
					\fill (\point) circle (1pt);
					\path (\point) ++(\goc:3mm) node {$\name$};
				}
				
			\end{tikzpicture}
		\end{center}
		Gọi lực căng lên hai thanh $AB$, $AC$ và lực nén lên thanh $OA$ lần lượt là $\overrightarrow{F}_1$, $\overrightarrow{F}_2$, $\overrightarrow{F}_3$.\\
		Đặt $\overrightarrow{F}_1 = \overrightarrow{AB'}$, $\overrightarrow{F}_2 = \overrightarrow{AC'}$, $\overrightarrow{P} = \overrightarrow{AP}$.\\
		Để hệ thống vẫn an toàn thì
		\[\overrightarrow{F}_1 + \overrightarrow{F}_2 + \overrightarrow{F}_3 + \overrightarrow{P} = \overrightarrow{0} \Leftrightarrow -\overrightarrow{F}_3 = \overrightarrow{F}_1 + \overrightarrow{F}_2 + \overrightarrow{P} \Leftrightarrow -\overrightarrow{F}_3 = \overrightarrow{AB'} + \overrightarrow{AC'} + \overrightarrow{AP} \quad (*).\]
		Từ giả thiết ta thấy $AB'$, $AC'$, $AP$ lần lượt là ba cạnh của hình hộp chữ nhật, từ $(*)$ và quy tắc hình hộp suy ra $-\overrightarrow{F}_3 = \overrightarrow{AO'}$ với $O' \in OA$ và $AO'$ là đường chéo của hình hộp chữ nhật trên.\\
		Khi đó ta có
		\[\left| -\overrightarrow{F}_3 \right| = \sqrt{AB'^2 + AC'^2 + AP^2} \Rightarrow \left| \overrightarrow{F}_3 \right|^2 = AB'^2 + AC'^2 + AP^2 \Rightarrow AP^2 = \left| \overrightarrow{F}_3 \right|^2 - AB'^2 - AC'^2.\]
		Theo giả thiết, ta có $\left| \overrightarrow{F}_1 \right| = \left| \overrightarrow{F}_2 \right| = \dfrac{1}{3}\left| \overrightarrow{F}_3 \right| \Rightarrow AB' = AC' = \dfrac{1}{3}\left| \overrightarrow{F}_3 \right| \Rightarrow AB'^2 = AC'^2 = \dfrac{1}{9}\left| \overrightarrow{F}_3 \right|^2$.\\
		Suy ra
		\[AP^2 = \left| \overrightarrow{F}_3 \right|^2 - \dfrac{1}{9}\left| \overrightarrow{F}_3 \right|^2 - \dfrac{1}{9}\left| \overrightarrow{F}_3 \right|^2 = \dfrac{7}{9}\left| \overrightarrow{F}_3 \right|^2.\]
		Do vậy $\left| \overrightarrow{P} \right|^2 = \dfrac{7}{9}\left| \overrightarrow{F}_3 \right|^2 \Rightarrow \left| \overrightarrow{P} \right| = \dfrac{\sqrt{7}}{3}\left| \overrightarrow{F}_3 \right| \le \dfrac{\sqrt{7}}{3} \cdot 150 \approx 132$.\\
		Vậy trọng lượng lớn nhất có thể treo vào đầu $A$ của giá treo là $132$ N.
	}
\end{ex}

\begin{ex}%[2H2V2-5]%[Dự án đề thi thử đợt 3-Nguyễn Đức Lợi]%[SGD tỉnh Lạng Sơn]%[GVPB - ThanhTa]
	Cho hình lăng trụ đứng $ABC.A'B'C'$ có đáy $ABC$ là tam giác vuông tại $A$, biết $AB = 2$, $AC = 3$, $AA' = 4$. Gọi $M$ là trung điểm của $AB$. Khoảng cách giữa hai đường thẳng $A'B$ và $CM$ bằng bao nhiêu (\textit{kết quả làm tròn đến hàng phần trăm})?
	\par
	\shortans{0,86}
	\loigiai{
		\begin{center}
			\begin{tikzpicture}[scale=0.8, font=\footnotesize, line join=round, line cap=round, >=stealth]
				% Định nghĩa các điểm đáy dưới
				\coordinate (A) at (0,0);
				\coordinate (B) at (4,0);
				\coordinate (C) at (1,-2);
				% Định nghĩa các điểm đáy trên
				\coordinate (A') at ($(A)+(0,4)$);
				\coordinate (B') at ($(B)+(0,4)$);
				\coordinate (C') at ($(C)+(0,4)$);
				\coordinate (z) at ($(A)!1.2!(A')$);
				\coordinate (x) at ($(A)!1.2!(B)$);
				\coordinate (y) at ($(A)!1.2!(C)$);
				% Trung điểm M của AB
				\coordinate (M) at ($(A)!0.5!(B)$);
				\draw[->] (A') -- (z) node[left] {$z$};
				\draw[->] (B) -- (x) node[above] {$x$};
				\draw[->] (C) -- (y) node[right] {$y$};
				% Vẽ các cạnh khuất
				\draw[dashed] (A)--(B)  (C) -- (M);
				% Vẽ các cạnh thấy
				\draw (C) -- (C') -- (A') -- (B') -- (B) -- (C) (C') -- (B') (A) -- (A') (C) -- (A);
				% Vẽ hai đường thẳng cần tính khoảng cách
				\draw[thick, blue,dashed] (A') -- (B);
				\draw[thick, blue, dashed] (C) -- (M);
				
				% Tên các điểm
				\foreach \p/\pos in {A/left, B/below, C/left, A'/left, B'/right, C'/left, M/above}
				\fill (\p) circle (1.5pt) node[\pos] {$\p$};
				
				% Ký hiệu vuông góc tại A
				\draw (0,0.2) -- (0.2,0.2) -- (0.2,0);
			\end{tikzpicture}
		\end{center}
		Chọn hệ trục tọa độ $Oxyz$ sao cho $A \equiv O(0;0;0)$, $B(2;0;0)$, $C(0;3;0)$ và $A'(0;0;4)$.\\
		Vì $M$ là trung điểm của $AB$ nên $M(1;0;0)$.\\
		Ta có $\vec{A'B} = (2;0;-4)$ và $\vec{CM} = (1;-3;0)$.\\
		Vectơ chỉ phương của $A'B$ là $\vec{u} = (1;0;-2)$, của $CM$ là $\vec{v} = (1;-3;0)$.\\
		Suy ra $\vec{n} =[\vec{v}, \vec{u}] = (6; 2; 3)$.\\
		Lấy điểm $M(1;0;0)$ thuộc $CM$ và $B(2;0;0)$ thuộc $A'B$, ta có $\vec{MB} = (1;0;0)$.\\
		Khoảng cách giữa hai đường thẳng là
		\[ d(A'B, CM) = \dfrac{|\vec{n} \cdot \vec{MB}|}{|\vec{n}|} = \dfrac{|6 \cdot 1 + 2 \cdot 0 + 3 \cdot 0|}{\sqrt{6^2 + 2^2 + 3^2}} = \dfrac{6}{7} \approx 0{,}86. \]
	}
\end{ex}

\begin{ex}%[2H2V2-6]%[Dự án đề TT Toán NH25-26-Dot 3- Hải Phụng]%[Cụm 13-Hải Phòng]%[GVPB - Viet Hoang Pham]
Một ngôi nhà gồm hai phần. Phần thân nhà dạng hình hộp chứ nhật $ABCD.OMNK$ có chiều dài $1\,200$ cm, chiều rộng $900$ cm, chiều cao $450$ cm. Phần mái nhà $S.ABCD$ có dạng một hình chóp tứ giác có các cạnh bên bằng nhau. Biết $ABCD$ là hình chữ nhật và có chiều cao của ngôi nhà $\left(\text{bằng khoảng cách từ }S\text{ đến }(OMNK)\right)$ bằng $600$ cm. Chọn hệ trục tọa độ $Oxyz$ sao cho $M$ thuộc tia $Ox$, $K$ thuộc tia $Oy$, $A$ thuộc tia $Oz$ (như hình vẽ), (mỗi đơn vị trên hệ trục ứng với $1$ m).
\begin{center}
\begin{tikzpicture}[line join=round,line cap=round, font=\footnotesize, scale=1,>=stealth]
\coordinate[label=left:$M$] (M) at (0,0);
\coordinate[label=above left:$O$] (O) at (1.2,.8);
\coordinate[label=below right:$N$] (N) at (4,0);
\coordinate[label=above left:$K$] (K) at ($(N)-(M)+(O)$);
\coordinate[label=below right:$A$] (A) at ($(O)+(90:3)$);
\coordinate[label=left:$B$] (B) at ($(M)-(O)+(A)$); 
\coordinate[label=below right:$C$] (C) at ($(N)-(O)+(A)$);
\coordinate[label=above right:$D$] (D) at ($(K)-(O)+(A)$); 
\coordinate[label=left:$x$] (x) at ($(O)!1.7!(M)$);
\coordinate[label=below:$y$] (y) at ($(O)!1.3!(K)$);
\coordinate[label=left:$z$] (z) at ($(O)!1.8!(A)$);
\coordinate (s) at ($(A)!0.5!(C)$);
\coordinate[label=above:$S$] (S) at ($(s)+(0,2)$);

\draw (B)--(M)--(N)--(K)--(D) (B)--(C)--(D) (N)--(C) (S)--(B) (S)--(C) (S)--(D);
\draw[dashed] (O)--(K) (A)--(O)--(M) (S)--(A)--(B) (A)--(D);
\draw[->] (K)--(y);
\draw[->] (M)--(x);
\draw[->] (A)--(z);

\draw [<->] ($(M)+(0,-0.3)$)--($(N)+(0,-0.3)$) node[below,pos=0.5]{$1200$};
\draw [<->] ($(N)+(0.3,0)$)--($(K)+(0.3,0)$) node[below right,pos=0.5]{$900$}; 
\draw [<->] ($(K)+(0.3,0)$)--($(D)+(0.3,0)$) node[below right,pos=0.5]{$450$}; 

\foreach \x in {A,B,C,D,O,M,N,K,S} 
 \draw[fill=black] (\x) circle (1.5pt); 
\end{tikzpicture}
\end{center}
Gọi tọa độ điểm $S$ là $S(a;b;c)$. Hãy tính $a-b+c$.
\shortans[oly]{4,5}
\loigiai{
\immini[thm]{
Đổi $1\,200 \text{ cm} = 12\text{ m}$; $900\text{ cm} = 9\text{ m}$; $600\text{ cm} = 6\text{ m}$.\\  
Gọi $I$ là hình chiếu vuông góc của điểm $S$ lên mặt phẳng $(OMNK)$. Khi đó $I$ là tâm hình chữ nhật $OMNK$. \\
Do đó tọa độ của điểm $I$ là $(4{,}5;6;6)$. \\
Suy ra tọa độ của điểm $S$ là $S(4{,}5;6;6)$.\\
Do đó $a-b+c=4{,}5 - 6 + 6 = 4{,}5$.	
}{
\begin{tikzpicture}[line join=round,line cap=round, font=\footnotesize, scale=1,>=stealth]
\coordinate[label=left:$M$] (M) at (0,0);
\coordinate[label=above left:$O$] (O) at (1.2,.8);
\coordinate[label=below right:$N$] (N) at (4,0);
\coordinate[label=above left:$K$] (K) at ($(N)-(M)+(O)$);
\coordinate[label=below right:$A$] (A) at ($(O)+(90:3)$);
\coordinate[label=left:$B$] (B) at ($(M)-(O)+(A)$); 
\coordinate[label=below right:$C$] (C) at ($(N)-(O)+(A)$);
\coordinate[label=above right:$D$] (D) at ($(K)-(O)+(A)$); 
\coordinate[label=left:$x$] (x) at ($(O)!1.7!(M)$);
\coordinate[label=below:$y$] (y) at ($(O)!1.3!(K)$);
\coordinate[label=left:$z$] (z) at ($(O)!1.8!(A)$);
\coordinate (s) at ($(A)!0.5!(C)$);
\coordinate[label=above:$S$] (S) at ($(s)+(0,2)$);
\coordinate[label=right:$I$] (I) at ($(O)!0.5!(N)$);

\draw (B)--(M)--(N)--(K)--(D) (B)--(C)--(D) (N)--(C) (S)--(B) (S)--(C) (S)--(D);
\draw[dashed] (O)--(K) (A)--(O)--(M) (S)--(A)--(B) (A)--(D) (S)--(I);
\draw[->] (K)--(y);
\draw[->] (M)--(x);
\draw[->] (A)--(z);

\draw [<->] ($(M)+(0,-0.3)$)--($(N)+(0,-0.3)$) node[below,pos=0.5]{$1200$};
\draw [<->] ($(N)+(0.3,0)$)--($(K)+(0.3,0)$) node[below right,pos=0.5]{$900$}; 
\draw [<->] ($(K)+(0.3,0)$)--($(D)+(0.3,0)$) node[below right,pos=0.5]{$450$}; 

\foreach \x in {A,B,C,D,O,M,N,K,S,I} 
 \draw[fill=black] (\x) circle (1.5pt); 
\end{tikzpicture}
}
}
\end{ex}

\begin{ex}%[2H2V2-6]%[Dự án đề TT Toán NH25-26-Dot 3- Hải Phụng]%[Cụm 13-Hải Phòng]%[GVPB - Viet Hoang Pham]
Hai chiếc khinh khí cầu bay lên từ cùng một địa điểm trong không gian. Để theo dõi hành trình của hai khinh khí cầu, người ta chọn hệ trục toạ độ $Oxyz$  với gốc $O$  đặt tại điểm xuất phát của hai khinh khí cầu, mặt phẳng $(Oxy)$ trùng với mặt đất với trục $Ox$  hướng về phía Nam, trục $Oy$  hướng về phía Đông và trục $Oz$  hướng thẳng lên trời (đơn vị đo lấy theo kilômét). Vào lúc $9$ giờ sáng, chiếc thứ nhất nằm cách điểm xuất phát $3$ km về phía Đông và $2$ km về phía Nam, đồng thời cách mặt đất $0{,}5$ km; chiếc thứ hai nằm cách điểm xuất phát $1$ km về phía Bắc và $1$ km về phía Tây, đồng thời cách mặt đất $0{,}3$ km. Cùng thời điểm đó, một người đứng
trên mặt đất và nhìn thấy hai khinh khí cầu nói trên. Biết rằng, so với các vị trí quan sát khác trên mặt đất, vị trí người đó đứng có tổng khoảng cách đến hai khinh khí cầu là nhỏ nhất. Vị trí người quan sát đứng lúc đó là $M(a;b;c)$. Tính tổng $32a+60b+5c$.

\shortans[oly]{$34$}
\loigiai{ \begin{center}
		\begin{tikzpicture}[line join=round, line cap=round, >=stealth, font=\footnotesize, scale=1]
			\def\x{3}
			\def\y{5}
			\def\z{4}
			\path (0:0) coordinate (O)
			++(0:\y) coordinate (y)
			(O)++(-140:\x) coordinate (x)
			(O)++(90:\z) coordinate (z)
			($(O)!0.6!(x)$) coordinate (X)
			($(O)!0.8!(y)$) coordinate (Y)
			($(O)!0.6!(z)$) coordinate (Z)
			($(X)+(Y)-(O)$) coordinate (XY)
			($(XY)+(Z)-(O)$) coordinate (A)
			($(O)!-0.4!(x)$) coordinate (X')
			($(O)!-0.25!(y)$) coordinate (Y')
			($(O)!0.5!(z)$) coordinate (Z')
			($(X')+(Y')-(O)$) coordinate (X'Y')
			($(X'Y')+(Z')-(O)$) coordinate (B)
			;
			%Vẽ các trục
			\draw[->,thick] (O)--(x) node[right]{$x$ (Nam)};
			\draw[->,thick] (O)--(y) node[above]{$y$ (Đông)};
			\draw[->,thick] (O)--(z) node[left]{$z$};
			%Các hướng
			\draw[dashed,thick] (O)--($(O)!-0.9!(y)$);
			\draw[dashed,thick] (X)--(XY)--(Y) (XY)--(A);
			\draw[dashed,thick] (X') --(X'Y')--(Y') (X'Y')--(B);
			\draw[dashed,thick] (O)--($(O)!-1!(x)$);
			
			\draw ($(X)!0.5!(XY)$) node[below]{$3$ km} ($(XY)!0.5!(Y)$) node[right]{$2$ km} ($(X')!0.5!(O)$) node[right]{$1$ km} ($(Y')!0.3!(O)$) node[below]{$1$ km} ($(XY)!0.7!(A)$) node[right]{0,5 km} ($(X'Y')!0.7!(B)$) node[left]{0,3 km} ; 
			\draw ($(Y')+(-3,0)$) node[below]{(Tây)} ($(O)!-1.7!(X)$) node[right]{(Bắc)};
			% Máy bay tại A:
			\node[font=\Huge, rotate=0] at (A) {\faBullseye};
			%% Máy bay tại B:
			\node[font=\Huge, rotate=120] at (B) {\faBullseye};
		\end{tikzpicture}
	\end{center}
Vào lúc 9 giờ sáng:
\begin{itemize}
\item Chiếc khinh khí cầu thứ nhất ở vị trí $A(2;3;0{,}5)$.
\item Chiếc khinh khí cầu thứ hai ở vị trí $B(-1;-1;0{,}3)$.
\end{itemize}
Mặt đất là mặt phẳng $(Oxy)$, hai điểm $A$, $B$ nằm về cùng phía so với mặt phẳng $(Oxy)$. \\
Gọi  $A'$ là điểm đối xứng với $A$ qua mặt phẳng $(Oxy)$. 
Khi đó $A'(2;3;-0,5)$. \\
Tổng khoảng cách từ vị trí người đó đứng trên mặt đất đến hai khinh khí cầu là
$$T=MA+MB=MA'+MB \geq A'B.$$
Vậy $T$ đạt giá trị nhỏ nhất bằng $A'B$ khi và chỉ khi $A'$, $M$, $B$ thẳng hàng. \\
Do $M(a;b;c) \in (Oxy) $ nên $c=0$. \\
Ta có $\vec{A'M} = (a-2;b-3;0,5)$ và $\vec{A'B} = (-3;-4;0,8)$. \\
Ba điểm $A',M,B$ thẳng hàng khi và chỉ khi hai vectơ $\vec{A'M}$ và $\vec{A'B}$ cùng phương. Khi đó 
$$\dfrac{a-2}{-3} = \dfrac{b-3}{-4} = \dfrac{0,5}{0,8} \Rightarrow a=\dfrac{1}{8}; b=\dfrac{1}{2}.$$
Vậy $32a+60b+5c = 34$
}
\end{ex}

\begin{ex}%[2H2V2-6]%[TN-THPT-NguyenTrungThien-HaTinh-NH25-26]%[Loc Do]%[PB: Phan Anh]
	Một cuộc diễn tập phòng không được đặt trong không gian với hệ tọa độ $Oxyz$ với $1$ đơn vị độ dài trên các trục tính bằng $1$ km trên thực tế. Một bệ phóng tên lửa phòng không được đặt tại vị trí $O(0;0;0)$ và tên lửa bay theo một đường thẳng với vận tốc không đổi là $600$ m/s. Một máy bay không người lái bay theo một đường thẳng với vận tốc không đổi là $1\,080$ km/h và bay theo hướng của vectơ $\overrightarrow{u}=(-3;4;0)$. Khi phát hiện máy bay không người lái tại vị trí $A(6;-20;1)$ thì tên lửa khai hỏa, rời bệ phóng và sau đó bắn hạ được mục tiêu. Hỏi khoảng cách từ bệ phóng tên lửa đến vị trí máy bay không người bị bắn hạ bằng bao nhiêu km (\textit{kết quả làm tròn đến hàng phần chục, giả sử cả máy bay không người lái và tên lửa đều không chịu tác động của trọng lực hay lực cản không khí})?
	\shortans{14,4}
	\loigiai{
		Đổi $600$ m/s $=2160$ km/h.\\
		Phương trình đường bay của máy bay không người lái là $d\colon \heva{&x=6-3 t \\ &y=-20+4t \\ &z=1}$, ($t \in \mathbb{R}$, $t>0$).\\
		Gọi điểm máy bay bị bắn trúng là $M$. Vì $M$ thuộc $d$ nên $M(6-3t ;-20+4t ; 1)$.\\
		Khi máy bay ở vị trí $A$ thì tên lửa khai hỏa và bắn trúng máy bay, do đó thời gian máy bay di chuyển từ $A$ đến $M$ và thời gian tên lửa di chuyển từ $O$ đến $M$ là như nhau. \\	
		Khi đó	
		$$
		\dfrac{AM}{1080}=\dfrac{O M}{2160} \Leftrightarrow \dfrac{\sqrt{(-3 t)^2+(4t)^2+0^2}}{1\,080}=\dfrac{\sqrt{(6-3t)^2+(-20+4t)^2+1^2}}{2\,160}.
		$$	
		Giải phương trình trên nhận được nghiệm thỏa mãn là $t=\dfrac{36}{25}$.\\
		Từ đó suy ra $OM \approx 14{,}4$ (km).
	}
\end{ex}

\begin{ex}%[2H3K1-1]%[Dự án C-TT THPT NH25-26-Đợt 4- ĐỖ CHÍ TÂM]%[THPT Phan Bội Châu-Khánh Hòa]%[GVPB: Phạm Văn Long]
	\immini[thm]{Sân hiên hình chữ nhật của một ngôi nhà là khoảng đất $ABCD$ được lợp mái bằng kính màu để hạn chế ánh sáng đi qua với mái dốc. Các bề mặt bên $ADHE$ và $CGHD$ nằm ở bức tường bên ngoài ngôi nhà. Đặt vào mô hình hệ trục tọa độ như hình vẽ thì ta có $B\left(5;\dfrac{7}{2};0\right)$; $E\left(5;0;2\right)$ và $H\left(0;0;3\right)$. Trên tường nhà có một ngọn đèn đặt tại điểm $L$ cách điểm $D$ một khoảng $6$ m theo phương thẳng đứng. Phần có mái của sân hiên in bóng lên khu vườn bằng phẳng phía trước ngôi nhà dưới ánh đèn tạo thành khoảng đất hạn chế ánh sáng. Tính diện tích khoảng đất đó (\textit{Kết quả làm tròn kết quả đến hàng phần chục}).
		\shortans{$45{,}9$}
	}{\begin{tikzpicture}[>=stealth,line join=round,line cap=round,font=\footnotesize,scale=1, black]
			\path (0,0) coordinate (D) (0,0.7) coordinate (H) (0,2) coordinate (D')
			(2,0) coordinate (C) (3,0) coordinate (C') (0,0.5) coordinate (H')
			(-120:3) coordinate (A) (-120:4) coordinate (A') (0,1.5) coordinate (L);
			\path ($(C)-(D)+(A)$) coordinate (B)
			($(H')-(D)+(A)$) coordinate (E)
			($(H')-(D)+(B)$) coordinate (F)
			($(H)-(D)+(C)$) coordinate (G);
			\fill[gray!30] (E)--(H)--(G)--(F)--cycle;
			\draw (A)--(B)--(C);
			\draw[dashed] (A)--(D)--(C) (D)--(H);
			\draw [->,thick] (A)--(A') node[left] {$x$};
			\draw [->,thick] (C)--(C') node[above] {$y$};
			\draw [->,thick] (H)--(D') node[right] {$z$};
			\draw (E)--(H)--(G)--(F)--cycle (A)--(E) (B)--(F) (C)--(G);
			\foreach \t/\g in {A/180,B/-90,C/-40,D/-80,E/180,F/120,G/0,H/30,L/180}{
				\fill (\t) circle (1pt) node[shift={(\g:7pt)},font=\scriptsize]{$ \t $};
			}
	\end{tikzpicture}}
	\loigiai{\begin{center}
			\begin{tikzpicture}[>=stealth,line join=round,line cap=round,font=\footnotesize,scale=1]
				\path (0,0) coordinate (D) (0,0.7) coordinate (H) (0,2) coordinate (D')
				(2,0) coordinate (C) (5,0) coordinate (C') (0,0.5) coordinate (H')
				(-120:3) coordinate (A) (-120:5) coordinate (A') (0,1.5) coordinate (L);
				\path ($(C)-(D)+(A)$) coordinate (B)
				($(H')-(D)+(A)$) coordinate (E)
				($(H')-(D)+(B)$) coordinate (F)
				($(H)-(D)+(C)$) coordinate (G);
				\path (intersection of L--E and A--A') coordinate (E');
				\path (intersection of L--G and C--C') coordinate (G');
				\path ($(C)-(D)+(E')$) coordinate (M);
				\path (intersection of L--F and E'--M) coordinate (F');
				\fill[gray!40, opacity=0.6] (E)--(H)--(G)--(F)--cycle;
				\fill[pattern=dots, pattern color=red!30] (E')--(A)--(B)--(C)--(G')--(F')--cycle;
				\draw (A)--(B)--(C);
				\draw[dashed] (A)--(D)--(C) (D)--(H) (E)--(H');
				\draw [->,thick] (A)--(A') node[left] {$x$};
				\draw [->,thick] (C)--(4.2,0) node[above] {$y$};
				\draw [->,thick] (H)--(D') node[right] {$z$};
				\draw (E)--(H)--(G)--(F)--cycle (A)--(E) (B)--(F) (C)--(G);
				\draw (L)--(E')--(F')--(G')--cycle (L)--(F');
				\foreach \t/\g in {A/-90,C/-40,F/120,G/0}{
					\fill (\t) circle (1pt) node[shift={(\g:7pt)},font=\scriptsize]{$ \t $};
				}
				\foreach \t/\g in {B/-90,E/180,H/30, L/180}{
					\fill[black] (\t) circle (1.5pt) node[shift={(\g:7pt)},font=\scriptsize]{$ \t $};
				}
				\foreach \t/\g in {D/-80,E'/180,G'/90, F'/-90}{
					\fill[black] (\t) circle (1.5pt) node[shift={(\g:7pt)},font=\scriptsize]{$ \t $};
				}
			\end{tikzpicture}
		\end{center}
		$B\left(5;3{,}5;0\right)$, $E\left(5;0;2\right)$, $H\left(0;0;3\right)$, $L(6;0;0)$. Suy ra $C(0;3{,5};0)$, $F(5;3{,}5; 2)$, $A(5; 0;0)$, $G(0; 3{,}5; 3)$.\\
		$GC$ là đường trung bình của tam giác $LDG'$. Suy ra $G'\left(0;7;0\right)$.\\
		Ta có $\dfrac{E'A}{E'D}=\dfrac{EA}{LD}=\dfrac{2}{6}=\dfrac{1}{3}\Rightarrow E'A=\dfrac{1}{3}E'D$\\
		Suy ra $\vec{E'A}=\dfrac{1}{3}\vec{E'D}\Rightarrow E'\left(7{,}5;0;0\right)$.\\
		Mặt khác $\dfrac{EF}{E'F'}=\dfrac{LE}{LE'}=\dfrac{2}{3}\Rightarrow \vec{E'F'}=1{,}5\vec{EF}$\\
		Suy ra $F'(7{,}5; 5{,}25; 0)$\\
		Diện tích khoảng đất thiếu ánh sáng là diện tích hình thang vuông $DG'F'E'$.\\
		Ta có $DG'=7$, $E'F'=5{,}25$, $DE'=7{,}5$.\\
		\hspace*{2cm}$S_{DG'F'E'}=\dfrac{\left(DG'+E'F'\right)\cdot DE'}{2}=\dfrac{\left(7+5{,}25\right)\cdot 7{,}5}{2}=45{,}9$.
	}
\end{ex}

\begin{ex}%[Dự án đề TT Toán NH25-26-Dot 4- Hải Phụng]%[THCS-THPT LÊ THÁNH TÔNG - TP.HCM]%[GVPB - Viet Hoang Pham]%[2H5C1-7]
		Cận kề ngày Tết cổ truyền, anh Nghĩa dự định trang trí hệ thống đèn LED cho khoảng sân nhà hình chữ nhật $MNPQ$ có chiều dài $MQ=12$ m, $MN=8$ m. Để cung cấp điện cho hệ thống đèn LED, anh Nghĩa đi dây điện theo trình tự: $A \to B \to C \to O \to D \to E \to F \to H$ (như hình vẽ bên). Trong đó:
		\begin{itemize}
			\item Điểm $A$ nằm trên cột cổng cao $3$ m, vị trí chân cột cách cây Nêu $3$ m và chân cột nằm trên đoạn $MQ$.
			\item Đoạn $BC$ nằm trên bức tường hình chữ nhật $PQIJ$ có $PJ=4$ m (bức tường vuông góc với mặt đất), độ dài $BC=3$ m và $BC \parallel IJ$.
			\item Điểm $O$ nằm trên mái che hình chữ nhật $IJUV$ có $JU=10$ m (mái che song song với mặt đất).
			\item Cột đèn đặt tại góc $N$ và điểm $D$ nằm trên cột đèn cao $3$ m.
			\item Điểm $E, F$ nằm trên cây Nêu sao cho $EF=2$ m. Cây Nêu $MG$ vuông góc với mặt đất tại $M$, chiều cao $MG=10$ m.
			\item Điểm $H$ cách đỉnh $G$ một khoảng $HG=0{,}5$ m ($HG$ song song với mặt đất).
		\end{itemize}
		Tính tổng độ dài dây điện ngắn nhất mà anh Nghĩa cần sử dụng để hoàn thành hệ thống đèn trang trí này (Đơn vị: mét, làm tròn kết quả đến hàng phần chục).
\begin{center}
	\begin{tikzpicture}[scale=0.8, line join=round, line cap=round, >=stealth]
		% --- ĐỊNH NGHĨA CÁC ĐIỂM CƠ SỞ (PHỐI CẢNH) ---
		\coordinate (M) at (0,0);
		\coordinate (Q) at (6,0);
		\coordinate (N) at (-2.5,-1.8);
		\coordinate (P) at ($(Q)+(N)-(M)$);
		% --- VẼ SÂN MNPQ ---
		\draw[thick] (M) -- (Q) -- (P) -- (N) -- cycle;
		\node at ($(M)!0.5!(N)$)[sloped, above, scale=0.8]{$8$ m};
		\node at ($(N)!0.5!(P)$)[below, scale=0.8]{$12$ m};
		% --- VẼ BỨC TƯỜNG PQIJ ---
		\coordinate (I) at ($(Q)+(0,4)$);
		\coordinate (J) at ($(P)+(0,4)$);
		\draw[fill=gray!10] (P) -- (Q) -- (I) -- (J) -- cycle;
		% Vẽ vân gạch cho tường
		\begin{scope}
			\clip (P) -- (Q) -- (I) -- (J) -- cycle;
			\foreach \i in {0,0.4,...,4} \draw[gray!40] ($(P)+(0,\i)$) -- ($(Q)+(0,\i)$);
			\foreach \i in {0,0.5,...,6} \draw[gray!40] ($(P)+(\i,0)$) -- ($(J)+(\i,0)$);
		\end{scope}
		\node at ($(P)!0.5!(J)$)[left, scale=0.8]{$4$ m};
		% --- VẼ CÂY NÊU MG ---
		\coordinate (G) at (0,7);
		\draw[line width=2pt, green!50!black] (M) -- (G);
		\foreach \i in {1,2,3,4,5,6} \draw[green!20!black] (-0.1,\i) -- (0.1,\i+0.1); % Đốt tre
		% --- VẼ MÁI CHE IJUV ---
		\coordinate (U) at ($(J)+(-4.5,0)$);
		\coordinate (V) at ($(I)+(-4.5,0)$);
		\draw[fill=red!10, opacity=0.8] (U) -- (J) -- (I) -- (V) -- cycle;
		\node[above] at (I) {$I$};
		\node[above] at (J) {$J$};
		\node[left] at (U) {$U$};
		\node[above] at (V) {$V$};
		% --- VẼ CỘT ĐÈN TẠI N ---
		\begin{scope}[shift={(N)}]
			\draw[fill=black] (-0.2,0) rectangle (0.2,0.2); % Chân đế
			\draw[fill=black] (-0.05,0.2) rectangle (0.05,2.5); % Thân cột
			\draw[fill=black] (-0.3,2.5) -- (0.3,2.5) -- (0.4,3.2) -- (-0.4,3.2) -- cycle; % Đầu đèn
			\draw[fill=yellow!60] (-0.25,2.6) -- (0.25,2.6) -- (0.3,3) -- (-0.3,3) -- cycle; % Kính đèn
			\coordinate (D) at (0,3); % Điểm D trên đỉnh cột đèn
		\end{scope}
		% --- VẼ CÁC ĐIỂM CHỐT A, B, C, O, E, F, H ---
		\coordinate (A) at (1.5,1.4); % Ước lượng vị trí cột cổng
		\draw[fill=gray!30] (1.3,0) rectangle (1.7,1.4); % Cột cổng A
		\coordinate (B) at ($(I)+(-0.6,-2.3)$);
		\coordinate (C) at ($(J)+(0.5,-1.5)$);
		\coordinate (O) at ($(U)!0.5!(I)$); % Điểm O trên mái
		\coordinate (E) at (0,4.5);
		\coordinate (F) at (0,5.8);
		\coordinate (H) at (-1.2,6.5);
		% --- VẼ LỒNG ĐÈN TẠI H ---
		\draw[thick] (G) -- (H);
		\draw[fill=red, rounded corners=2pt] ($(H)+(-0.3,-1)$)-|($(H)+(0.3,0)$)-|cycle;
		\draw[yellow, line width=1pt] ($(H)+(-0.3,-0.2)$) -- ($(H)+(0.3,-0.2)$);
		\draw[yellow, line width=1pt] ($(H)+(-0.3,-0.8)$) -- ($(H)+(0.3,-0.8)$);
		\node at (H) [above] {H};
		% --- VẼ HỆ THỐNG DÂY ĐIỆN (NÉT ĐỨT VÀ LIỀN) ---
		\draw[red, thick, dashed] (A) -- (B);
		\draw[red, thick, dashed] (B) -- (C) node[midway, right, black, scale=0.7]{$3$ m};
		\draw[red, thick, dashed] (C) -- (O);
		\draw[red, thick, dashed] (O) -- (D);
		\draw[red, thick] (D) -- (E);
		\draw[red, thick] (E) -- (F) node[midway, left, black, scale=0.7]{$2$ m};
		\draw[red, thick] (F) -- (H);
		
		% --- ĐÁNH DẤU CÁC ĐIỂM ---
		\foreach \p in {A,B,C,O,D,E,F,G,M,N,P,Q}
		\filldraw[black] (\p) circle (1pt);
		\node[above] at (A) {$A$};
		\node[right] at (B) {$B$};
		\node[below] at (C) {$C$};
		\node[above] at (O) {$O$};
		\node[left] at (D) {$D$};
		\node[above right] at (E) {$E$};
		\node[right] at (F) {$F$};
		\node[right] at (G) {$G$};
		\node[below] at (N) {$N$};
		\node[below] at (P) {$P$};
		\node[right] at (Q) {$Q$};
		\node[left] at (M) {$M$};
	\end{tikzpicture}
\end{center}
	\shortans{$36{,}5$}
	\loigiai{
		Chọn hệ trục $Oxyz$ sao cho $M \equiv O$, $ M(0;0;0)$, tia $Ox$ đi qua $Q$, tia $Oy$ đi qua $N$.\\			
		Khi đó:	$Q(12;0;0)$, $N(0;8;0)$, trục $O z \perp (MNPQ)$, $G(0;0;10)$,  $P(12;8;0)$.\\			
		Để độ dài dây điện ngắn nhất thì các đoạn dây phải đi theo đoạn thẳng.\\
		Theo giả thiết ta có
		$$T = AB + BC + CO + OD + DE + EF + FH
		= AB + CO + OD + DE + FH + 5.$$
		Ta có:
		$A(3;0;3)$, $D(0;8;3)$, $H\left(0;\dfrac{1}{2};10\right)$,	mặt phẳng $(PQIJ)\colon x = 12$, mặt phẳng $(IJUV)\colon z = 4 $.\\		
		Gọi $A'$ là điểm đối xứng với $A$ qua mặt phẳng $(IJUV) \Rightarrow A'(3;0;5)$. \\			
		Gọi $B'$, $C'$ lần lượt là điểm đối xứng của $B$, $C$ qua $(IJUV)$.\\
		Khi đó: $AB = A'B'$. \\
		Gọi \( A'' \) là điểm đối xứng của \( A' \) qua mặt phẳng $(PQIJ) \Rightarrow A''(21;0;5)$, suy ra $A'B' = A''B'.$\\			
		Gọi $ S $ là điểm thỏa mãn $ \overrightarrow{A''S}  =\overrightarrow{B'C'} \Rightarrow A''B'C'S $ là hình bình hành
		$\Rightarrow A''B' = SC'.$\\
		Vì $ BC \parallel IJ \Rightarrow A''S \parallel IJ \Rightarrow S(21;3;5).$\\		
		Do đó, $AB + CO + OD$ nhỏ nhất bằng $SD$ khi và chỉ khi $S, O, C', D$ thẳng hàng, $SD = \sqrt{470}$.\\				
		Gọi $K$ là điểm thỏa mãn $\overrightarrow{HK} = \overrightarrow{FE} = (0; 0; -2) \Rightarrow K\left( 0; \dfrac{1}{2}; 8 \right)$ và $KHFE$ là hình bình hành.\\				
		Gọi $D'$ là điểm đối xứng với $D$ qua trục $Mz \Rightarrow D'(0; -8; 3)$. \\
		Lại có $ DE + FH = DE + EK = D'E + EK. $
		Để $DE + FH$ nhỏ nhất thì $D'$, $E$, $K$ thẳng hàng và $D'K = \dfrac{\sqrt{389}}{2}$.\\			
		Vậy $\min T = \sqrt{470} + \dfrac{\sqrt{389}}{2} + 5 \approx 36{,}5$.		
	}
	
\end{ex}

\begin{ex}%[Đề thi thử môn Toán TRƯỜNG THPT Cổ Loa - Hà Nội năm học 2025-2026]%[Phạm Quốc Toàn - Phạm Hoàng Việt]%[2H5H1-7]
	Tại một khu quảng trường, người ta muốn trang trí dải đèn led để chào mừng một sự kiện. Từ đỉnh $M$ của một cái cây cao (minh họa bằng đoạn $MK$), người ta kéo dây đèn led thẳng xuống một vị trí $H$ nào đó trên mặt đất, rồi tiếp tục kéo dây đèn led thẳng từ $H$ đến một vị trí $A$ cố định trên mặt đất sao cho $MH$ luôn vuông góc với $HA$. Sau đó, từ vị trí $H$, người ta sẽ đi dây điện thẳng đến công tắc nằm trên một bức tường vuông góc với mặt đất. Chọn hệ trục toạ độ $Oxyz$ với mặt phẳng $(Oxy)$ trùng với mặt đất (đơn vị trên mỗi trục tính theo mét), người ta xác định được $M(4;-2;3)$, $A(-2;6;0)$ và bức tường chứa công tắc nằm trên mặt phẳng $(\alpha) \colon 3x+4y+89=0$. Để độ dài đường dây điện phải đi là ngắn nhất, người ta đã tính toán khoảng cách nhỏ nhất từ vị trí $H$ đến bức tường bằng $a$ (m). Giá trị của $a$ bằng bao nhiêu?
	\begin{center}
		\begin{tikzpicture}[font=\footnotesize, line join=round, line cap=round, >=stealth,scale=.7]
			\coordinate (A) at (3,2);
			\coordinate (H) at (5,1);
			\coordinate (M) at (8,7);
			\coordinate (K) at (8,2);
			%%%%%%%%%%%%%%%%%%%%%%%%%%%%%%%%
			\draw (0,0)--(0,5) -- (2,8) -- (2, 3) -- (0,0) (0,0) -- (8,0) -- (10,3) -- (2, 3);
			\draw (A) -- (H) -- (M) -- (K);
			\foreach \i/\g in {A/90,H/-90,M/90,K/-90}{\draw(\i) circle (1pt) ($(\i)+(\g:3mm)$) node[scale=1]{$\i$};}
			\pic[draw,thin,angle radius=2mm] {right angle=A--H--M};
			%%%%%%%%%%%%%%%%%%%%%%%%%%%%%%%%
			\fill (0.2,0) node[above right]{$(Oxy)$}; 
			\fill (2,3) node[left]{$(\alpha)$}; 
		\end{tikzpicture}
	\end{center}
	\par\shortans[0]{15}
	\loigiai{
		Ta có $\widehat{AHM}=90^\circ$.\\
		Suy ra $H$ thuộc mặt cầu $(S)$ có đường kính $AM$.\\
		Mặt cầu $(S)$ có tâm $I\left(1;2;\dfrac{3}{2}\right)$, bán kính $R=\dfrac{\sqrt{109}}{2}$.\\
		Ta có $H=(S)\cap (Oxy)$.\\
		Suy ra $H$ thuộc đường tròn $(C)$ có tâm $J(1;2;0)$ là hình chiếu của $I$ lên mặt phẳng $(Oxy)$. \\
		Bán kính của đường tròn là $r=\sqrt{R^2-IJ^2}=5$.\\
		Ta có $\mathrm{d} \big(J, (\alpha)\big)=20$ và $(\alpha) \perp (Oxy)$ nên $\mathrm{d} \big(H, (\alpha)\big)_{\min}=\mathrm{d} \big(J, (\alpha)\big)-r=20-5=15$.
	}
\end{ex}

\begin{ex}%[2H5V1-5]%[TT-LienTruong-HaNoi-N25-26]%[Phạm Quốc Toàn - Lê Phúc]
	Cho hình chóp $S. ABCD$ có đáy là hình thoi cạnh $3$, góc $\widehat{ABC}=60^\circ$ và cạnh bên $SA$ vuông góc với mặt phẳng đáy. Gọi $M$, $N$ lần lượt là trung điểm của $SB$, $SD$. Biết góc nhị diện $[M, AC, N]$ bằng $120^\circ$. Tính khoảng cách từ điểm $C$ đến mặt phẳng $(SBD)$. (\textit{Không làm tròn kết quả các phép tính trung gian, chỉ làm tròn kết quả cuối cùng đến hàng phần mười})
		\par\shortans[0]{1{,}1}
	\loigiai{
		\immini[thm]{Đặt hệ trục tọa độ $Oxyz$ với gốc $O$ là giao điểm của hai đường chéo $AC$ và $BD$.\\
		Vì $ABCD$ là hình thoi cạnh $3$ và $\widehat{ABC}=60^\circ$, ta có tam giác $ABC$ đều, suy ra $AC=3$ và $OB=OD=\dfrac{3\sqrt{3}}{2}$.\\
		Tọa độ các đỉnh là $O(0;0;0)$, $A\left(-\dfrac{3}{2};0;0\right)$, $C\left(\dfrac{3}{2};0;0\right)$, $B\left(0;\dfrac{3\sqrt{3}}{2};0\right)$, $D\left(0;-\dfrac{3\sqrt{3}}{2};0\right)$. \\ 
		Vì $SA\perp (ABCD)$, gọi $SA=h>0$ là chiều cao của hình chóp $S.ABCD$ nên $S\left(-\dfrac{3}{2};0;h\right)$.\\ 
		Tọa độ trung điểm của $SB$ là $M\left(-\dfrac{3}{4};\dfrac{3\sqrt{3}}{4};\dfrac{h}{2}\right)$. \\ 
		Tọa độ trung điểm của $SD$ là $N\left(-\dfrac{3}{4};-\dfrac{3\sqrt{3}}{4};\dfrac{h}{2}\right)$. }
		{	\begin{tikzpicture}[font=\footnotesize, line join=round, line cap=round, >=stealth,scale=0.6]
				\coordinate (S) at (3,7);
				\coordinate (A) at (3,2);
				\coordinate (B) at (0,0);
				\coordinate (C) at (5,0);
				\coordinate (D) at (8,2);
				\coordinate (O) at (4,1);
				\coordinate (G) at (4,4.5);
				\coordinate (M) at (1.5,3.5);
				\coordinate (N) at (5.5,4.5);
				%%%%%%%%%%%%%%%%%%%%%%%%%%%
				\draw[->] (5,0)--(6,-1) node[below]{$x$};
				\draw[->] (B)--($(O)!1.3!(B)$) node[below]{$y$};
				\draw[->] (4,4.5)--(4,7.6) node[right]{$z$};
				%%%%%%%%%%%%%%%%%%%%%%%%%%%
				\draw (S) -- (B)  -- (C) -- (D) -- (S) -- (C);
				\draw[dashed] (S)--(A) -- (B)  (A) -- (D) -- (B) (A) -- (C) (O) -- (G);
				%%%%%%%%%%%%%%%%%%%%%%%%%%%
				\foreach \i/\g in {S/90,A/180,B/-90,C/-90,D/0,O/-90,M/180,N/90}{\fill(\i) circle (1.5pt) ($(\i)+(\g:3mm)$) node[scale=1]{$\i$};}
				\pic[draw,thin,angle radius=1.5mm] {right angle=G--O--B} 
				pic[draw,thin,angle radius=1.5mm] {right angle=G--O--C};
		\end{tikzpicture}}
	\noindent
		Góc nhị diện $\left[M, AC, N\right]$ bằng $120^\circ$ nên góc giữa hai mặt phẳng $(MAC)$ và $(NAC)$ bằng $60^\circ$. \\ 
		Vectơ pháp tuyến của mặt phẳng $(MAC)$ là $\overrightarrow{n}_1=\left[\overrightarrow{AC}, \overrightarrow{AM}\right]=\left(0;-h;\dfrac{3\sqrt{3}}{2}\right)$.\\
		Vectơ pháp tuyến của mặt phẳng $(NAC)$ là $\overrightarrow{n}_2=\left[\overrightarrow{AC}, \overrightarrow{AN}\right]=\left(0;h;\dfrac{3\sqrt{3}}{2}\right)$.\\
		Vậy 
		\begin{align*}
			& \cos 60^\circ =\dfrac{\left|\overrightarrow{n}_1\cdot{\overrightarrow{n}_2}\right|}{\left|\overrightarrow{n}_1\right|\cdot \left|\overrightarrow{n}_2\right|}=\dfrac{\left| \dfrac{27}{4}-h^2\right|}{\sqrt{h^2+\dfrac{27}{4}}\cdot \sqrt{h^2+\dfrac{27}{4}}}  \\
			& \Leftrightarrow \dfrac{1}{2}=\dfrac{\left| \dfrac{27}{4}-h^2\right|}{\sqrt{h^2+\dfrac{27}{4}}\cdot \sqrt{h^2+\dfrac{27}{4}}} \\ 
			& \Leftrightarrow \sqrt{h^2+\dfrac{27}{4}}\cdot \sqrt{h^2+\dfrac{27}{4}}=2\left| \dfrac{27}{4}-h^2\right| \\ 
			& \Leftrightarrow h^2+\dfrac{27}{4}=2\left| \dfrac{27}{4}-h^2\right| \\ 
			& \Leftrightarrow h=\dfrac{3}{2}. 
		\end{align*}
		Phương trình mặt phẳng  $(SBD)$ đi qua ba điểm $S\left(-\dfrac{3}{2};0;\dfrac{3}{2}\right)$, $B\left(0;\dfrac{3\sqrt{3}}{2};0\right)$, $D\left(0;-\dfrac{3\sqrt{3}}{2};0\right)$ là $x+z=0$.\\
		Khoảng cách từ điểm $C\left(\dfrac{3}{2};0;0\right)$ đến mặt phẳng $(SBD)$ là
		\[
		\mathrm{d}\big(C;(SBD)\big)=\dfrac{3}{2\sqrt{2}}=\dfrac{3\sqrt{2}}{4}\approx 1{,}1. 
		\]
	}
\end{ex}

\begin{ex}%[2H5V1-5]%[Dự án đề thi thử đợt 3-Nguyễn Đức Lợi]%[SGD tỉnh Lạng Sơn]%[GVPB - ThanhTa]
	Trong không gian $Oxyz$, cho điểm $M(1;-2;-5)$. Gọi $(P)$ là mặt phẳng đi qua $M$ và cách gốc tọa độ một khoảng lớn nhất. Biết mặt phẳng $(P)$ cắt $Ox$, $Oy$, $Oz$ lần lượt tại ba điểm $A$, $B$, $C$. Thể tích tứ diện $OABC$ bằng bao nhiêu?
	\par
	\shortans{450}
	\loigiai{
		\begin{center}
			\begin{tikzpicture}[scale=0.8, font=\footnotesize, line join=round, line cap=round, >=stealth]
				% Trục tọa độ
				\draw[->] (3,0,0) -- (4,0,0) node[below] {$x$};
				\draw[->] (0,4,0) -- (0,5,0) node[left] {$z$};
				\draw[->] (0,0,5) -- (0,0,6) node[right] {$y$};
				
				% Các điểm giao tuyến (ước lệ minh họa)
				\coordinate (A) at (3,0,0);
				\coordinate (B) at (0,0,5);
				\coordinate (C) at (0,4,0);
				\coordinate (O) at (0,0,0);
				\coordinate (M) at (1,1,1.5); % Minh họa điểm M
				
				% Vẽ mặt phẳng (P)
				\draw[thick, fill=blue!10, opacity=0.6] (A) -- (B) -- (C) -- cycle;
				
				% Vẽ đoạn OM
				\draw[dashed, thick, red] (O) -- (M) (O)--(A) (O)--(B) (O)--(C);
				\fill[red] (M) circle (1.5pt) node[right] {$M$};
				
				% Tên các điểm giao
				\fill (A) circle (1pt) node[below] {$A$};
				\fill (O) circle (1pt) node[below] {$O$};
				\fill (B) circle (1pt) node[left] {$B$};
				\fill (C) circle (1pt) node[left] {$C$};
			\end{tikzpicture}
		\end{center}
		Gọi $h$ là khoảng cách từ $O$ đến mặt phẳng $(P)$.\\\
		Ta luôn có $h = \mathrm{d}(O, (P)) \le OM$.\\
		Do đó, khoảng cách $\mathrm{d}(O, (P))$ lớn nhất khi và chỉ khi $(P)$ vuông góc với $OM$ tại $M$.\\
		Khi đó, phương trình mặt phẳng $(P)$ đi qua $M(1; -2; -5)$ và có vectơ pháp tuyến $\overrightarrow{OM} = (1; -2; -5)$ là
		\[1(x - 1) - 2(y + 2) - 5(z + 5) = 0 \Leftrightarrow x - 2y - 5z - 30 = 0.\]
		Tọa độ giao điểm của $(P)$ với các trục tọa độ
		\begin{itemize}
			\item Cho $y=0$, $z=0$ suy ra $x = 30 \Rightarrow A(30; 0; 0) \Rightarrow OA = 30$.
			\item Cho $x=0$, $z=0$ suy ra $-2y = 30 \Rightarrow y = -15 \Rightarrow B(0; -15; 0) \Rightarrow OB = 15$.
			\item Cho $x=0$, $y=0$ suy ra $-5z = 30 \Rightarrow z = -6 \Rightarrow C(0; 0; -6) \Rightarrow OC = 6$.
		\end{itemize}
		Tứ diện $OABC$ có $OA$, $OB$, $OC$ đôi một vuông góc nên thể tích là
		\[V = \dfrac{1}{6} \cdot OA \cdot OB \cdot OC = \dfrac{1}{6} \cdot 30 \cdot 15 \cdot 6 = 450.\]
	}
\end{ex}

\begin{ex}%[2H5V1-7]%[TT Cụm 5 - Ninh Bình-NH25-26]%[GVSB:Dương Công Tạo]%[GVPB1: Nguyễn Tấn Tài]
	\immini[thm]{Một hộ gia đình muốn xây dựng một kho chứa mini dạng khối tứ diện $OABC$ tận dụng góc tường nhà (với $O$ là góc tường, ba cạnh $OA$, $OB$, $OC$ nằm trên ba mép tường vuông góc). Để lắp đặt hệ thống thông gió, mặt phía trước của kho (mặt phẳng $ABC$) bắt buộc phải đi qua một van điều tiết đặt tại vị trí $M$ có tọa độ $M(1;2;1)$ (đơn vị tính là mét). Gia đình muốn thiết kế kho sao cho thể tích của kho là nhỏ nhất để không làm ảnh hưởng quá nhiều đến diện tích sinh hoạt của sân chung. Khi kho được thiết kế với thể tích nhỏ nhất, hãy tính khoảng cách từ góc tường $O$ đến mặt phẳng $(ABC)$.}
	{\begin{tikzpicture}[scale=.75,
			x={(220:1.3cm)}, 
			y={(10:1.1cm)}, 
			z={(90:1cm)},  
			>=stealth    
			]
			\coordinate (O) at (0,0,0);
			\coordinate (A) at (3.5,0,0); 
			\coordinate (B) at (0,4.0,0); 
			\coordinate (C) at (0,0,3.5);
			\fill[lightgray!40] (0,0,3.7) -- (0,0,0) -- (0,4.2,0) -- (0,4.2,3.7) -- cycle;
			
			\fill[lightgray!40] (3.7,0,0) -- (0,0,0) -- (0,0,3.7) -- (3.7,0,3.7) -- cycle;
			
			\draw[->] (O) -- (4,0,0) node[anchor=north east, pos=1] {$x$};
			\draw[->] (O) -- (0,4.5,0) node[anchor=west, pos=1] {$y$};
			\draw[->] (O) -- (0,0,4) node[anchor=south, pos=1] {$z$};
			
			\draw[thick] (A) -- (B) -- (C) -- cycle;
			
			\node[above left] at (A) {$A$};
			\node[below right] at (B) {$B$};
			\node[above left] at (C) {$C$};
			\node[below] at (O) {$O$};
			
			\coordinate (M) at (1.5, 2.5,1.2); 
			
			\fill[black] (M) circle (1pt); 
			\node[above=3pt] at (M) {$M(1; 2;1)$}; 
			
			\draw[dashed, black!60] (M) -- (0, 1, -0.3); 
			
	\end{tikzpicture}}
	\par  \shortans{2}
	\loigiai{
		Đặt hệ trục tọa độ $Oxyz$ có gốc $O(0;0;0)$, các tia $Ox$, $Oy$, $Oz$ lần lượt chứa các cạnh $OA$, $OB$, $OC$ của tứ diện. \\
		Gọi tọa độ các điểm là $A(a;0;0)$, $B(0;b;0)$, $C(0;0;c)$ với $a$, $b$, $c > 0$. \\
		Phương trình mặt phẳng $(ABC)$ có dạng mặt phẳng theo đoạn chắn $\dfrac{x}{a} + \dfrac{y}{b} + \dfrac{z}{c} = 1$. \\
		Vì mặt phẳng $(ABC)$ đi qua điểm $M(1;2;1)$ nên ta có hệ thức $\dfrac{1}{a} + \dfrac{2}{b} + \dfrac{1}{c} = 1$. \\
		Thể tích khối tứ diện $OABC$ là $V = \dfrac{1}{6}abc$. \\
		Áp dụng bất đẳng thức Cauchy (AM-GM) cho ba số dương $\dfrac{1}{a}$, $\dfrac{2}{b}$, $\dfrac{1}{c}$, ta có \[1 = \dfrac{1}{a} + \dfrac{2}{b} + \dfrac{1}{c} \ge 3\sqrt[3]{\dfrac{1}{a} \cdot \dfrac{2}{b} \cdot \dfrac{1}{c}} = 3\sqrt[3]{\dfrac{2}{abc}} \Rightarrow 1 \ge 27 \cdot \dfrac{2}{abc} \Rightarrow abc \ge 54.\]
		Thể tích nhỏ nhất $V_{\min} = \dfrac{1}{6} \cdot 54 = 9$ m$^3$ khi và chỉ khi $\dfrac{1}{a} = \dfrac{2}{b} = \dfrac{1}{c} = \dfrac{1}{3} \Rightarrow a = 3$, $b = 6$, $c = 3$. \\
		Khi đó, phương trình mặt phẳng $(ABC)$ là \[\dfrac{x}{3} + \dfrac{y}{6} + \dfrac{z}{3} = 1 \Leftrightarrow 2x+y+2z-6=0.\]
		Khoảng cách từ gốc tọa độ $O$ đến mặt phẳng $(ABC)$ là \[\mathrm{d}\left(O,(ABC)\right)= \dfrac{\left|2 \cdot 0+0+2 \cdot 0-6\right|}{\sqrt{2^2 + 1^2 + 2^2}} = \dfrac{|-6|}{\sqrt{9}} = \dfrac{6}{3} = 2\text{ (m)}.\]
	}
\end{ex}

\begin{ex}%[2H5V1-7]%[Dự án C đề thi thử TN Môn Toán NH25-26 - Dot 4 - Hector Tran]%[Liên trường THPT cụm 09 - Đợt 2 - Hà Nội]%[GVPB - Thanh Phát]
	\immini[thm]{
		Trong giờ thể dục học về kỹ thuật chuyền bóng hơi, Bình và An tập chuyền bóng cho nhau. Ở một động tác Bình chuyền bóng cho An, quả bóng bay lên cao nhưng lại lệch sang bên trái của An và rơi xuống vị trí cách chỗ An đứng $0{,}5$ m và cách chỗ Bình $4{,}5$ m. Chọn hệ trục tọa độ $Oxyz$ sao cho gốc tọa độ $O$ tại vị trí của Bình, vị trí của An nằm trên tia $Oy$ và mặt phẳng $(Oxy)$ là mặt đất. Biết rằng quỹ đạo của quả bóng nằm trong mặt phẳng $(\alpha) \colon x + by + cz + d = 0$ và $(\alpha)$ vuông góc với mặt đất. Khi đó, giá trị của $200b^2 - 100c^2 + 200d^2$ bằng bao nhiêu?
	}{
		\begin{tikzpicture}[line join = round, line cap=round,>=stealth,font=\footnotesize,scale=1]
			\draw[->] (2,2)coordinate (O)--(7.5,2)coordinate (y);
			\draw[->] (2,2)--(0.6,0.8)coordinate (x);
			\draw[->] (2,2)--(2,5.2)coordinate (z);
			\fill (2.3,3.1) circle(2pt) (6.55,1.35) circle(2pt);
			\draw[dashed] (2.3,3.1)..controls +(28:3.0) and +(90:1) .. (6.55,1.35)coordinate (H)node[below]{$A$};
			\fill (1.9,2.85) circle(3pt);
			\draw[line width=3pt] (1.9, 2.85) -- (2.0, 2.6) -- (2.0, 2.3) -- (1.8, 2.1)
			(2.0, 2.3) -- (2.05, 2.2) -- (2.0, 2.0)
			(2.0, 2.75) -- (2.2, 2.95)
			;
			\fill (6.65,2.7) circle(3pt);
			\draw[line width=3pt]
			(6.65, 2.7) -- (6.85, 2.4) -- (6.75, 2.2) -- (6.95, 2.0)
			(6.85, 2.4) -- (6.65, 2.3) -- (6.65, 2.0)
			(6.7, 2.6) -- (6.5, 2.5)
			(6.7, 2.6) -- (6.5, 2.4)
			;
			\draw (O)--(H)--($(O)!(H)!(y)+(0.13,0)$)coordinate (t);
			\draw pic[draw,angle radius=2mm] {right angle = H--t--y};
			\path
			($(t)+(0,-0.4)$) node[right]{$0{,}5$\,m}
			($(O)+(2,-0.75)$) node[]{$4{,}5$\,m}
			;
			\foreach \p/\r in {O/-90,y/90,x/180,z/180}
			\fill (\p) node[shift={(\r:3mm)}]{$\p$};
		\end{tikzpicture}
	}
	\par\shortans{2,5}
	\loigiai{
		Bình ở $O(0; 0; 0)$. Quả bóng rơi xuống tại điểm $A(x; y; 0)$ thỏa mãn $OA = 4{,}5$ và cách An một đoạn $0{,}5$.\\
		Suy ra $A\left( 0{,}5; \sqrt{20};0\right) $.\\
		Mặt phẳng $(\alpha) \colon x + by + cz + d = 0$ đi qua $O$ nên $d = 0$.\\
		Điểm $A\left( 0{,}5; \sqrt{20};0\right) $ thuộc $(\alpha)$ nên $0{,}5+ \sqrt{20} b = 0 \Leftrightarrow b = -\dfrac{\sqrt{5}}{20}$.\\
		Mặt khác $(\alpha)$ vuông góc với mặt đất $(Oxy)$ nên $\overrightarrow{n}_{\alpha} \perp \overrightarrow{k} \Rightarrow c = 0$.\\
		Vậy giá trị $200b^2 - 100c^2 + 200d^2 =200\left( -\dfrac{\sqrt{5}}{20}\right) ^2 =2{,}5$.
	}
\end{ex}

\begin{ex}%[Đề thi thử môn Toán Sở giáo dục tỉnh Tuyên Quang năm học 2025-2026]%[Đoàn Minh Tâm]%[2H5V1-7]%[GVPB: Thanh Phát]
	Ông An đang xây một ngôi nhà, trong quá trình xây phải đổ bê-tông cho một mái vát để lợp ngói. Ông tính toán việc ghép cốt pha đi qua điểm $B$ trên một chân tường và điểm $C$ trên cột góc nhà, đồng thời mặt ghép cốt pha phải đi qua điểm $A$ trên chân tường còn lại cách điểm $O$ ở góc giao hai chân tường một khoảng $5$ m, ông cũng tận dụng một chiếc cột có sẵn để chống mặt ghép (xem hình dưới). Biết rằng hai bức tường được xây vuông góc với nhau, mỗi bức tường đều vuông góc với sàn mái nhà, cột có chiều cao $1$ m và cách hai bức tường với cùng khoảng cách $1$ m (đỉnh cột là điểm $M$). Diện tích nhỏ nhất của khung ghép cốt pha $ABC$ là bao nhiêu mét vuông (\textit{làm tròn kết quả đến hàng phần trăm})?
	\begin{center}
		\begin{tikzpicture}[
			% Tùy chỉnh góc nhìn 3D
			x={({0.866cm,-0.15cm})},
			y={({-0.6cm,-0.5cm})},
			z={({0cm,1cm})},
			scale=1.2,
			font=\footnotesize,
			line join=round,line cap=round,line width=.6pt,>=stealth
			]
			\coordinate (O) at (0,0,0);
			\coordinate (A_ext) at (7.5,0,0);
			\coordinate (A) at (7,0,0);
			\coordinate (B_ext) at (0,4.5,0);
			\coordinate (B) at (0,3.5,0);
			\coordinate (C_ext) at (0,0,4.5);
			\coordinate (C) at (0,0,3.8);
			\coordinate (Mx) at (2, 0, 0);
			\coordinate (My) at (0, 2, 0);
			\coordinate (Mz) at (0, 0, 2);
			\coordinate (Mxy) at (2, 2, 0);
			\coordinate (M) at (2, 2, 2);
			\begin{scope}
				\clip (A) -- (B) -- (C) -- cycle;
				\fill[gray!10] (A) -- (B) -- (C) -- cycle;
				\foreach \x in {0.2, 0.4, ..., 7.4} {
					\pgfmathsetmacro{\factor}{1-\x/7.5}
					\draw[very thin, gray!60] (\x, {\factor*3.5}, 0) -- (\x, 0, {\factor*3.8});
				}
				\foreach \y in {0.2, 0.4, ..., 3.4} {
					\pgfmathsetmacro{\factor}{1-\y/3.5}
					\draw[very thin, gray!60] ({\factor*7.5}, \y, 0) -- (0, \y, {\factor*3.8});
				}
			\end{scope}
			
			\draw[dashed, thick] (O) -- (A);
			\draw[dashed, thick] (O) -- (B);
			\draw[dashed, thick] (O) -- (C);
			\draw[thick] (A) -- (A_ext);
			\draw[thick] (B) -- (B_ext);
			\draw[thick] (C) -- (C_ext);
			\draw[thick] (A) -- (B) -- (C) -- cycle;
			\draw[dashed, thick] (Mx) -- (Mxy) -- (My);
			\draw[dashed, thick] (O) -- (Mxy);
			\draw[dashed, thick] (Mxy) -- (M)--(Mz);
			\fill (M) circle (1pt);
			\fill (O) circle (1pt);
			\fill (A) circle (1pt);
			\fill (B) circle (1pt);
			\fill (C) circle (1pt);
			\node[above right ] at (A) {$ A $};
			\node[left ] at (B) {$ B $};
			\node[above right ] at (C) {$ C $};
			\node[above right ] at (O) {$ O $};
			\node[above right=-1pt ] at (M) {$ M $};
			\tikzset{lbl/.style={fill=white, inner sep=1.5pt}}
			\path (O) -- (A) node[midway, sloped, above= 1pt, lbl] {Chân tường $1$};
			\path (O) -- (B) node[pos=0.38, sloped, above= 1pt, lbl] {Chân tường $2$};
			\path (O) -- (C) node[pos=0.45, sloped, above= 1pt, lbl] {Cột góc nhà};
			\path (Mxy) -- (Mx) node[midway, sloped, below=1pt, lbl] {$ 1 $ m};
			\path (Mxy) -- (My) node[midway, sloped, below=1pt, lbl] {$ 1 $ m};
			\path (Mxy) -- (M) node[pos=0.7, right=1pt, lbl] {$ 1 $ m};
		\end{tikzpicture}
	\end{center}
	\par\shortans[oly]{9,38}
	\loigiai{
		\immini[thm]{
			Chọn hệ trục tọa độ $Oxyz$ với $O(0;0;0)$ là góc giao hai chân tường, tia $Oy$ nằm trên chân tường chứa điểm $A$, tia $Ox$ nằm trên chân tường chứa điểm $B$, tia $Oz$ trùng với cột góc nhà chứa điểm $C$.\\
			Theo giả thiết bài toán, ta có tọa độ các điểm $A(0;5;0)$.\\
			Giả sử $B(b;0;0)$ và $C(0;0;c)$ với $b>0$, $c>0$.\\
			Mặt phẳng $(ABC)$ chứa mặt ghép cốt pha có phương trình theo đoạn chắn là
			\[\dfrac{x}{b} + \dfrac{y}{5} + \dfrac{z}{c} = 1.\]
		}{
			\begin{tikzpicture}[
				x={({0.866cm,-0.15cm})},
				y={({-0.6cm,-0.5cm})},
				z={({0cm,1cm})},
				scale=.7,
				font=\footnotesize,
				line join=round,line cap=round,line width=.6pt,>=stealth
				]
				\coordinate (O) at (0,0,0);
				\coordinate (A_ext) at (7.5,0,0);
				\coordinate (A) at (7,0,0);
				\coordinate (B_ext) at (0,4.5,0);
				\coordinate (B) at (0,3.5,0);
				\coordinate (C_ext) at (0,0,4.5);
				\coordinate (C) at (0,0,3.8);
				\coordinate (Mx) at (2, 0, 0);
				\coordinate (My) at (0, 2, 0);
				\coordinate (Mz) at (0, 0, 2);
				\coordinate (Mxy) at (2, 2, 0);
				\coordinate (M) at (2, 2, 2);
				\draw[dashed, thick] (O) -- (A);
				\draw[dashed, thick] (O) -- (B);
				\draw[dashed, thick] (O) -- (C);
				\draw[thick,->] (A) -- (A_ext) node[right]{$y$};
				\draw[thick,->] (B) -- (B_ext) node[below]{$x$};
				\draw[thick,->] (C) -- (C_ext) node[above]{$z$};
				\draw[thick] (A) -- (B) -- (C) -- cycle;
				\draw[dashed, thick] (Mx) -- (Mxy) -- (My);
				\draw[dashed, thick] (O) -- (Mxy);
				\draw[dashed, thick] (Mxy) -- (M)--(Mz);
				\fill (M) circle (1.5pt);
				\fill (O) circle (1.5pt);
				\fill (A) circle (1.5pt);
				\fill (B) circle (1.5pt);
				\fill (C) circle (1.5pt);
				\node[above right ] at (A) {$A$};
				\node[left ] at (B) {$B$};
				\node[above right ] at (C) {$C$};
				\node[above left ] at (O) {$O$};
				\node[above right=-1pt ] at (M) {$M$};
			\end{tikzpicture}
		}
		\noindent Cột có chiều cao $1$ m, cách hai bức tường $1$ m nên đỉnh cột là $M(1;1;1)$.\\
		Vì mặt ghép cốt pha đi qua $M$ nên điểm $M \in (ABC)$.\\
		Thay tọa độ $M$ vào phương trình mặt phẳng ta được
		\[\dfrac{1}{b} + \dfrac{1}{5} + \dfrac{1}{c} = 1 \Leftrightarrow \dfrac{1}{b} + \dfrac{1}{c} = \dfrac{4}{5} \Leftrightarrow \dfrac{b+c}{bc} = \dfrac{4}{5} \Leftrightarrow b+c = \dfrac{4}{5}bc.\]
		Áp dụng bất đẳng thức AM-GM, ta có
		\[b+c \ge 2\sqrt{bc} \Leftrightarrow \dfrac{4}{5}bc \ge 2\sqrt{bc} \Leftrightarrow \sqrt{bc} \ge \dfrac{5}{2} \Leftrightarrow bc \ge \dfrac{25}{4}.\]
		Diện tích khung ghép cốt pha là diện tích tam giác $ABC$ là
		\[S = \dfrac{1}{2} \left| \left[ \overrightarrow{AB}, \overrightarrow{AC} \right] \right| = \dfrac{1}{2} \sqrt{b^2c^2 + 25b^2 + 25c^2}.\]
		Ta cần tìm giá trị nhỏ nhất của $S$, ta xét $S^2$ như sau
		\[S^2 = \dfrac{1}{4} \left[ b^2c^2 + 25(b^2+c^2) \right] = \dfrac{1}{4} \left[ b^2c^2 + 25 \left( (b+c)^2 - 2bc \right) \right].\]
		Thay $b+c = \dfrac{4}{5}bc$ vào biểu thức trên, ta có
		\[S^2 = \dfrac{1}{4} \left[ b^2c^2 + 25 \left( \dfrac{16}{25}b^2c^2 - 2bc \right) \right] = \dfrac{1}{4} \left( 17b^2c^2 - 50bc \right).\]
		Đặt $t = bc$ với $t \ge \dfrac{25}{4} = 6{,}25$.\\
		Xét hàm số $f(t) = 17t^2 - 50t$ trên $[6{,}25; +\infty)$.\\
		Ta có $f'(t) = 34t - 50$.\\
		Với $t \ge 6{,}25$ thì $f'(t) \ge 34\cdot 6{,}25 - 50 = 162{,}5 > 0$.\\
		Do đó hàm số $f(t)$ đồng biến trên $[6{,}25; +\infty)$.\\
		Suy ra $\min f(t) = f(6{,}25) = 17\cdot (6{,}25)^2 - 50\cdot 6{,}25 = 351{,}5625$.\\
		Khi đó $S^2 \ge \dfrac{1}{4} \cdot 351{,}5625 = 87{,}890625 \Rightarrow S \ge \sqrt{87{,}890625} = 9{,}375 \approx 9{,}38 \text{ m}^2$.\\
		Dấu \lq\lq $=$\rq\rq\ xảy ra khi $b=c=2{,}5$.
	}
\end{ex}

\begin{ex}%[2H5V2-8]%[Dự án C đề thi thử TN Môn Toán NH25-26 - Dot 4 - Vũ Hồng Toàn]%[THPT Mai Anh Tuấn - Thanh Hóa]%[GVPB - Lê Hữu Kiệt]
	Trong không gian một bức tường hình chữ nhật $ABCD$ có chiều cao $3$\,m được dựng vuông góc với mặt phẳng $(Oxy)$ với hai điểm $A(3;1;0)$ và $B(1;4;0)$. Giả sử rằng mặt phẳng $(Oxy)$ là mặt đất, trục $Oz$ hướng lên trên và mỗi đơn vị trên hệ trục tương ứng là $1$\,m. Một trạm phát sóng được đặt tại điểm $M(5;6;5)$. Trạm phát sóng phát ra các sóng thẳng và truyền về mọi hướng nhưng khi gặp bức tường thì bị chắn lại vì vậy có một vùng sau bức tường không có sóng. Một chất điểm bắt đầu chuyển động thẳng đều trong mặt phẳng mặt phẳng $(Oxy)$ từ điểm $N(-6;-4;0)$ qua gốc tọa độ $O$ đến chân tường với vận tốc $0{,}8$\,m/s. Hỏi sau bao nhiêu giây từ lúc chất điểm bắt đầu đi vào vùng mất sóng đến lúc chất điểm chạm vào chân tường và quay về vị trí ban đầu, giả thiết độ dày của bức tường là không đáng kể {\it(kết quả làm tròn đến hàng phần chục)}?
	\begin{center}
		\begin{tikzpicture}[scale=1, font=\footnotesize, line join=round, line cap=round, >=stealth]
			\tikzset{declare function={a=6;b=0.4*a;goc=-130;h=0.65*a;}}
			\path (0:0) coordinate (O)
			(goc+20:b-1) coordinate (A)++(0,b-0.5)coordinate (D)
			(goc+125:b) coordinate (B)++(0,b-0.5)coordinate (C)
			(a-1,-b+1)coordinate (Hx)++(0,h)coordinate (M)
			(-goc+40:h)coordinate (N)
			(intersection of M--C and Hx--B) coordinate (E)
			(intersection of M--D and Hx--A) coordinate (F)
			($(N)+(Hx)-(F)$)coordinate (Ix)
			(intersection of N--O and A--D) coordinate (No)
			(intersection of N--O and A--B) coordinate (H)
			(intersection of N--O and E--F) coordinate (I)
			(intersection of {O--$(goc:1)$} and D--A) coordinate (Ox)
			(intersection of {O--$(1,0)$} and B--C) coordinate (Oy)
			(intersection of {O--$(0,1)$} and D--C) coordinate (Oz) ;
			\fill[gray!10] ($(F) + (-1,-0.5)$) -- ($(N) + (-1,0.5)$) -- ($(Ix) + (0,0.7)$) -- ($(Hx) + (1,-0.5)$) -- cycle;
			\draw[gray, line width=2pt] (M)--(Hx);
			\draw[->] (Oy)--(a+0.5,0) node[shift={(-100:7pt)}]{$y$};
			\draw[->] (Ox)--(goc:b)node[shift={(0:7pt)}]{$x$};
			\draw[->] (Oz)--(0,h-1)node[shift={(200:7pt)}]{$z$};
			\fill[gray!60, opacity=0.6, pattern=north west lines, pattern color=red!30] (A) -- (B) -- (E) -- (F) -- cycle;
			\node at (-2,-0.7,0) [teal, font=\footnotesize, align=center] {Vùng mất sóng};
			\shade[shading=ball](M)circle(4pt);
			\draw[dashed] (Oy)--(O)--(Ox) (Oz)--(O);
			\fill[draw, pattern=bricks, pattern color=red!50] (A)--(B)--(C)--(D)--cycle;
			\draw[red] (N)--(No);
			\draw[red, dashed, thin] (H)--(No);
			\foreach \x /\gN in {A/-90,B/-10,C/90,D/-40,O/-80,N/90,M/90}
			\fill (\x) circle (1pt)($(\x)+(\gN:2.5mm)$) node {$\x$};
		\end{tikzpicture}
	\end{center}
	\par\shortans{21,1}
	\loigiai{
		\begin{center}
			\begin{tikzpicture}[scale=1, font=\footnotesize, line join=round, line cap=round, >=stealth]
				\tikzset{declare function={a=6;b=0.4*a;goc=-130;h=0.65*a;}}
				\path (0:0) coordinate (O)
				(goc+20:b-1) coordinate (A)++(0,b-0.5)coordinate (D)
				(goc+125:b) coordinate (B)++(0,b-0.5)coordinate (C)
				(a-1,-b+1)coordinate (Hx)++(0,h)coordinate (M)
				(-goc+40:h)coordinate (N)
				(intersection of M--C and Hx--B) coordinate (E)
				(intersection of M--D and Hx--A) coordinate (F)
				($(N)+(Hx)-(F)$)coordinate (Ix)
				(intersection of N--O and A--D) coordinate (No)
				(intersection of N--O and A--B) coordinate (H)
				(intersection of N--O and E--F) coordinate (I)
				(intersection of {O--$(goc:1)$} and D--A) coordinate (Ox)
				(intersection of {O--$(1,0)$} and B--C) coordinate (Oy)
				(intersection of {O--$(0,1)$} and D--C) coordinate (Oz) ;
				\fill[gray!10] ($(F) + (-1,-0.5)$) -- ($(N) + (-1,0.5)$) -- ($(Ix) + (0,0.7)$) -- ($(Hx) + (1,-0.5)$) -- cycle;
				\draw[gray, line width=2pt] (M)--(Hx);
				\draw[->] (Oy)--(a+0.5,0) node[shift={(-100:7pt)}]{$y$};
				\draw[->] (Ox)--(goc:b)node[shift={(0:7pt)}]{$x$};
				\draw[->] (Oz)--(0,h-1)node[shift={(200:7pt)}]{$z$};
				\fill[gray!60, opacity=0.6, pattern=north west lines, pattern color=red!30] (A) -- (B) -- (E) -- (F) -- cycle;
				\node at (-2,-0.7,0) [teal, align=center] {Vùng mất sóng};
				\shade[shading=ball](M)circle(4pt);
				\draw[dashed] (Oy)--(O)--(Ox) (Oz)--(O);
				\fill[draw, pattern=bricks, pattern color=red!50] (A)--(B)--(C)--(D)--cycle;
				\draw[red] (N)--(No);
				\draw[red, dashed, thin] (H)--(No);
				\draw[thin, gray] (E)--(M) --(F)--cycle;
				\foreach \x /\gN in {A/-90,B/-10,C/90,D/-40,O/-80,H/-30,M/90,E/130,F/180,N/90,I/90}
				\fill (\x) circle (1pt)($(\x)+(\gN:2.5mm)$) node {$\x$};
			\end{tikzpicture}
		\end{center}
		Ta có $C(1;4;3)$, $D(3;1;3)$.\\
		Gọi $E(x;y;0)$, $F$ lần lượt là giao điểm của $MC$, $MD$ với mặt phẳng $(Oxy)$.\\
		Khi đó
		\allowdisplaybreaks
		\begin{eqnarray*}
			\overrightarrow{MC} = k\overrightarrow{ME} \Leftrightarrow \heva{&1-5 = k(x-5)\\& 4-6 = k(y-6)\\& 3-5 = k(0-5)} \Leftrightarrow \heva{&x = -5\\& y = 1\\&k = \dfrac{2}{5}} \Rightarrow E(-5;1;0).
		\end{eqnarray*}
		Tương tự, ta cũng có $F(0;-\dfrac{13}{2};0)$.\\
		Gọi $I$, $H$ lần lượt là giao điểm của đường thẳng $NO$ với $EF$ và $AB$.\\
		Vì $I$, $H$, $O$, $N \in (Oxy)$ nên chỉ xét trên mặt phẳng tọa độ $(Oxy)$.\\
		Ta có $N(-6;-4)$; $O(0;0)$ nên phương trình đường thẳng $NO$ là $2x - 3y = 0$.\\
		$E(-5;1)$; $F\left(0;-\dfrac{13}{2}\right)$ nên phương trình đường thẳng $EF$ là $3x + 2y + 13 = 0$.\\
		Với $A(3;1)$; $B(1;4)$ nên phương trình đường thẳng $AB$ là $3x + 2y - 11 = 0$.\\
		Suy ra $I(-3;-2)$, $H\left(\dfrac{33}{13};\dfrac{22}{13}\right) \Rightarrow IH = \dfrac{24\sqrt{13}}{13}$, $IN = \sqrt{13}$.\\
		Tổng quãng đường chất điểm di chuyển cần tính thời gian là
		$$IH + HN = 2IH + IN = \dfrac{61\sqrt{13}}{13}.$$
		Vậy thời gian cần tìm là $t = \dfrac{\dfrac{61\sqrt{13}}{13}}{0{,}8}  \approx 21{,}1$\,(giây).
	}
\end{ex}

\begin{ex}%[2H5V3-4]%[KNTT - Lớp 12 - Dự án C - Đợt 2]%[Loc do]%[PB: Phan Anh]
	Trong không gian $Oxyz$ (đơn vị km), Trái đất được mô phỏng là mặt cầu tâm $O$ bán kính $R=6\,400$ và đường xích đạo được xem là đường tròn giao tuyến giữa mặt cầu mô phỏng trái đất và mặt phẳng $(Oxy)$. Tại Bắc Cực $I(0;0;6\,400)$, Bộ chỉ huy thiết lập một \lq\lq Vòm An Ninh\rq\rq.Vùng an toàn này là một mặt cầu tâm I có bán kính quét là $R_q =1\,280\sqrt{10} $. Bất kỳ tàu nào đi vào biên giới của vùng này sẽ được bảo vệ tuyệt đối. Siêu điệp viên Alpha đang ẩn náu tại tọa độ $M(0;5\,120;3\,840)$. Anh nhận được mật lệnh:  Phải di chuyển từ vị trí hiện tại xuống đường Xích đạo để kích hoạt thiết bị định vị, sau đó lập tức rút lui về ranh giới Vòm An Ninh tại điểm nhập $K$ theo lộ trình ngắn nhất để tàu thoát. Biết vận tốc di chuyển của Alpha là $v_A =60$ km/h. Cùng lúc đó, tàu truy kích Beta đang phục kích tại vị trí $N(0;-6\,400;0)$. Ngay khi Alpha xuất phát, Beta tính toán được chính xác điểm $K$ và lao đến để chặn đầu. Tàu Beta phải duy trì vận tốc tối thiểu là bao nhiêu km/h để có thể chặn được Alpha trước lúc điệp viên này vào ranh giới vùng an toàn? (xem vận tốc của Alpha và Beta là chuyển động đều và hai tàu di chuyển trên mặt cầu. Làm tròn kết quả đến hàng phần chục)
	\par\shortans[0]{$84,6$}
	\loigiai{
		\begin{center}
			\begin{tikzpicture}[>=stealth,line join=round,line cap=round,font=\footnotesize,scale=1]
				\draw[->] (-4,0)--(4.5,0) node[below]{$x$};
				\draw [->] (0,-3.5)--(0,5) node [right]{$y$};
				\path (0,0) coordinate (O);
				\draw[black] (O) circle (3cm);
				\coordinate (H) at ($(O)+(0:3)$);
				\coordinate (M) at ($(O)+(25:3)$);
				\fill (2.7,0) circle (1pt) node[below] {$5\,120$};
				\fill (0,2.7) circle (1pt) node[left] {$5\,120$};
				\fill (1.3,0) circle (1pt) node[below] {$3\,840$};
				
				\coordinate (K) at ($(O)+(65:3)$);
				\coordinate (I) at ($(O)+(90:3)$);
				\coordinate (N) at ($(O)+(180:3)$);
				\draw[black] (I) circle (1.3cm);
				\draw [dashed] (0,1.3)--(M)--(2.7,0) (0,2.7)--(K)--(1.3,0) (K)--(N) ;
				\fill (0,1.3) circle (1pt) node[left] {$3\,840$};
				\foreach \t/\g in {H/30,M/20,K/30,I/60,N/-120}{
					\fill (\t) circle (1pt) node[shift={(\g:7pt)},font=\scriptsize]{$ \t $};
				}
				
			\end{tikzpicture}
		\end{center}
		Ta có $\overrightarrow{OI}=(0;0;6\,400)$, $\overrightarrow{OM}=(0;5\,120;3\,840)\Rightarrow \left[\overrightarrow{OI},\overrightarrow{OM}\right]=(-2\,048\,000;0;0)$.\\
		Suy ra $(OIM)\colon x=0$ hay $O$, $I$, $M$ thuộc mặt phẳng $(Oyz)$, ta có độ dài cung tròn $l=R\cdot\alpha $.\\
		Alpha đi theo hai cung $\wideparen{MH}$ và $\wideparen{HK}$. \\
		Ta có phương trình đường tròn 
		\begin{itemize}
			\item $(O)\colon y^2 +z^2 =6\,400^2$; \hfill$(1)$
			\item $(I)\colon y^2 +(z-6\,400)^2 =\left(1\,280\sqrt{10}\right)^2$. \hfill$(2)$
		\end{itemize}
		Lấy phương trình ($1$) trừ phương trình ($2$) ta được $z=5\,120\Rightarrow y=3\,840$. Suy ra $K(0;3\,840;5\,120)$.\\
		Độ dài cung Alpha đi là 
		\[l_{\wideparen{MH}} +l_{\wideparen{HK}} =\tan ^{-1} \left( \dfrac{3\,840}{5\,120} \right)\cdot 6\,400+\tan ^{-1} \left( \dfrac{5\,120}{3\,840} \right)\cdot 6\,400=3\,200\pi.\]
		Do đó thời gian Alpha đi là $t_{Alpha} =\dfrac{3200\pi}{60} $.\\
		Beta đi theo cung $\wideparen{NK}$, nên có $l_{\wideparen{NK}} =6\,400\left(\pi -\tan ^{-1} \dfrac{5\,120}{3\,840}\right) =14\,171{,}50359$.\\
		Vận tốc của tàu Beta là $v_{Beta} =\dfrac{14\,171{,}50\,359}{\dfrac{3\,200\pi}{60}} \approx 84{,}57993178$.\\
		Vậy $v_{Beta} \approx 84{,}6$ (km/h).
	}
>>>>>>> 33890a9ce1cfbd3079ed36a93e836cf3ef1fde30
\end{ex}